\documentclass[10pt,a4paper,twoside,bg=print,twocolumn,openany,nodeprecatedcode]{dndbook}
\usepackage[utf8]{inputenc}
\usepackage{tabulary}
\usepackage[normalem]{ulem}
\usepackage{color,soul}
\usepackage{float}
\usepackage{caption}
\usepackage{wrapfig}
\usepackage{tocloft}
\usepackage[toc]{multitoc}
\usepackage[hidelinks]{hyperref}
\raggedbottom
\extrafloats{2000}
\cftsetindents{section}{0em}{0em}
\cftsetindents{subsection}{0em}{0em}
\cftsetindents{subsubsection}{0.2em}{0em}
\setcounter{tocdepth}{1}
\renewcommand*{\multicolumntoc}{3}
\setlist[description]{nolistsep,listparindent=3pt,topsep=6pt}
\date{}
\begin{document}
\frontmatter
\tableofcontents
\mainmatter
\addcontentsline{toc}{chapter}{Chapter 1: Introduction to Drakkenheim}
\markboth{CHAPTER 1: INTRODUCTION TO DRAKKENHEIM}{CHAPTER 1: INTRODUCTION TO DRAKKENHEIM}
\chapter*{Introduction to Drakkenheim}
\textbf{Drakkenheim is no more}. An eldritch storm of falling stars devastated the city on that woeful eve, leaving behind a meteorite that would have lasting effects. Fifteen years later, Drakkenheim is a dismal urban wasteland contaminated by otherworldly magic and haunted by hideous monsters. Fantastic wealth, lost knowledge, and powerful artifacts lie ready for the taking by adventurers brave or foolish enough to venture into the ruins. However, even those who survive the dangers of Drakkenheim may not return unmarked by its terrors!

\includegraphics[width=0.4\textwidth]{images/../5e.tools/img/adventure/DoDk/000-00-001.drakenheim-city.png}
\section{Campaign Overview}
\textit{Dungeons of Drakkenheim} is a fantastic adventure set in a gothic renaissance ruined city. The campaign is designed with four player characters in mind, so adjust accordingly for smaller or larger groups. You can play \textit{Drakkenheim} as a series of loosely-connected exploration quests or an unpredictable narrative filled with faction conflict. The material herein will challenge characters from 1st to 13th level.


\begin{description}
\item[Exploration.]
Characters embark on dangerous expeditions into monster-filled ruins, haunted streets, decrepit gardens, shattered mage towers, baroque cathedrals, and ancient castles to plunder treasure and solve mysteries.
\item[Interaction.]
Five rival factions clash over the ruins of Drakkenheim. A tangled web of secrets, subterfuge, and sabotage awaits when the factions' opposing objectives and your characters' personal goals come into conflict!
\item[Combat.]
Beyond fierce warriors dispatched by the factions, all manner of otherworldly horrors inhabit Drakkenheim. Characters battle undead husks, predatory beasts, mutated monstrosities, and terrifying abominations produced by the mysterious meteor.
\end{description}

\DndDropCapLine{Y}{}ou'll need a copy of the core rules for Fifth Edition to properly run this adventure. When reading this book, text in \textbf{bold} usually references a creature or non-player character stat block, while text in \textit{italics} often indicates spells or magic items. Many new monsters, items, and spells are described in the appendix sections of this book; otherwise, consult the Fifth Edition core rules.


\subsection{Delerium}
Iridescent crystals of vast magical potential are found throughout Drakkenheim. Known as delerium, these lambent stones emanate unnatural energies that induce madness and monstrous transformations. Despite these dangers, delerium is ideally suited to crafting magic items and fuelling mighty new spells. Sorcerers, warlocks, wizards, and all manner of occult magicians crave its supernatural power, thus delerium commands a high price within arcane circles and underground markets alike. Many prospectors risk everything to collect a few fragments, but the true origins of the strange mineral remain unknown. The rival factions stand divided over whether the crystals should be destroyed, harnessed, or worshipped, and their simmering disagreements threaten outright war. Appendix D: Delerium, Magic \& Spells covers the crystals in detail.

\subsection{Broken Realm}
Drakkenheim was the cosmopolitan capital of Westemär, ruled by the imperious House von Kessel. A decade of bloody civil war erupted after the city's destruction. The realm was left shattered and the royal line broken. Now, the political order of the wider world lies in shambles, torn apart by religious schism and military conflict, and only a faint hope remains that the city can be rebuilt and the nation restored.

\subsection{Rival Factions}
Five factions have arrived in Drakkenheim to advance their ambitious agendas. Though their ideologies and methods differ, each seeks valuable resources and lost secrets within the city. However, the ruins present many confounding obstacles for their agents. As such, all the factions hope to attract talented adventurers to their cause, drawing would-be heroes into a fierce political intrigue that may boil over into violent struggle. Those who wish to decide the fate of Drakkenheim must choose their allies well and their enemies carefully! Chapter 3: Factions contains complete details for running each faction, but they are summarized below:

\subsubsection{Hooded Lanterns}
Formally known as the 4th Provisional Expeditionary Force to Reclaim the Capital, the Hooded Lanterns are an irregular military drawn up from veterans of the Civil War and remnants of the old Drakkenheim City Watch. Led by the dour \textbf{Lord Commander Elias Drexel}, they wage a guerrilla war against the monsters, bandits, and scavengers who have taken root in their nation's capital. Each has pledged their swords and lives to rebuild Drakkenheim and restore the Realm of Westemär. Several noble houses have set aside their feuds to support the Hooded Lanterns' efforts. These families believe once the city is restored, evidence will emerge revealing the true heir to the throne of Drakkenheim... or enable a cunning house to seize the crown.

\subsubsection{Queen's Men}
The Queen's Men are a loose affiliation between a hundred gangs of brigands, outlaws, and scoundrels, all of whom owe fealty to the enigmatic \textbf{Queen of Thieves}. These reprobates prey upon the pilgrims, prospectors, and explorers drawn to the city, extort and rob adventurers, smuggle delerium to disreputable clients in distant lands, and plunder the fantastic treasures and incredible wealth left behind in the city. The Queen of Thieves dreams of building an influential and powerful criminal empire forged in the anarchy and lawlessness of the ruins, and conspires against the other factions at every turn.

\subsubsection{Knights of the Silver Order}
The Knights of the Silver Order are an organization of oath-bound paladins and devout warriors sworn to combat supernatural evil, dark magic, and otherworldly incursions. The questing knights normally travel the realms in small bands and companies, and adhere to the Faith of the Sacred Flame, the dominant religion across the continent. They believe delerium is a contaminating blight borne of the darkest chaos, and have resolved to annihilate the corrupted crystals, slay every monster it has created, and hunt down diabolical mages who would wield its power to work evil magic. Now an entire regiment has been dispatched to Drakkenheim to carry out this righteous task, led by \textbf{Knight-Captain Theodore Marshal}.

\subsubsection{Followers of the Falling Fire}
The annihilation of Drakkenheim incited a serious religious schism within the Faith of the Sacred Flame. A new sect has emerged known as the Followers of the Falling Fire, who adhere to the prophecies of former Flamekeeper \textbf{Lucretia Mathias}. She attests that delerium is a divine gift, not a blight, and that these sacred stones will offer salvation against a greater darkness yet to come. The mainstream clergy have excommunicated Lucretia Mathias and branded her followers as an insane cult with heretical beliefs and blasphemous practices. Even still, masses of devout commonfolk now embark on a dangerous pilgrimage to Drakkenheim to take their place in a divine plan.

\subsubsection{Amethyst Academy}
The Amethyst Academy is a magical school for sorcerers and wizards, where students learn magic in remote castles and secluded universities. Directed by a shadowy cabal of powerful archmages, the Academy also operates an enterprising mage guild that controls the manufacture of magical wares, and orchestrates an influential arcane syndicate that provides occult counsel to the nobility of the continent. However, the Academy was dealt a terrible blow following Drakkenheim's destruction, as the city was home to their greatest stronghold: the Inscrutable Tower. Archwizard Eldrick Runeweaver leads the Academy's efforts to access the tower and safeguard the secret lore within. Meanwhile, their survey teams conduct ongoing research regarding delerium's supernatural qualities and unknown dangers. The mages plan to harness the crystals' vast energies to craft magical items and power new arcane spells.

\DndQuote%
{}
{''A dangerous foe, moving almost unseen, it invades your wits and saps your strength when it's needed most.''}
{Pluto Jackson}
\subsection{The Haze}
There are more hazards in Drakkenheim than monsters and rival factions. A billowing mist filled with glinting multi-coloured motes has settled over the entire city, which adventurers call the \textbf{Haze}. It's like the dust never really settled after the meteor fell.

\textbf{Characters gain no benefit from taking a long rest within the Haze.} Adventurers should be well-rested before they head into the ruins, as they'll need to escape the city to regain their hit points and spell slots. This effect cannot be circumvented by any mundane equipment, nor bypassed with spells such as \textit{rope trick}, \textit{Leomund's Tiny Hut}, or similar character abilities. These effects are described in detail in Chapter 5: Exploring Drakkenheim.

\subsection{Contamination}
Characters will also encounter deadly \textbf{eldritch contamination} caused by creatures, delerium, and other magical phenomena in Drakkenheim. Abilities, equipment, and spells that protect against diseases, magic, or poisons do not work against contamination. It cannot be healed naturally, nor readily removed with low-level spells. Player characters will need to work with the factions and find creative solutions to manage these problems during their adventures. The full rules and debilitating effects of contamination are outlined in Appendix C: Contamination.

\subsection{Dangers of the Dark City}
Finally, be warned that in Drakkenheim, characters will often stumble upon creatures and other perils far beyond what they can defeat on their own. Discretion and cunning may prevail when strength of arms fails. Adventurers must seek out like-minded allies, and be wary of making new enemies. They must have all the equipment they need before setting out on an expedition, use their resources carefully, and be prepared to retreat at a moment's notice!

\section{Content Warning}
Drakkenheim portrays a dark fantasy world. We recommend players and Game Masters have an open conversation at the start of the campaign to discuss their lines and veils. In support of this conversation, please note references to the following are common in this adventure:


\begin{itemize}
\item Violence, murder, blood, gore, cannibalism, and body horror.
\item Degenerative mutations that cause physical disfigurement and madness/insanity.
\item Rats, spiders, insects, demons, undead, ghosts, and other monsters.
\item Natural disasters, large-scale loss of life, civil war, displaced persons, and refugees.
\item Moral ambiguity, social/political manipulation, religious zealotry and military nationalism.
\end{itemize}

The campaign does not contain explicit sexuality, sexual assault, racial prejudice, systemic racism, homophobia, or transphobia. Any inferences to such are wholly unintentional.

\section{Character Creation}
In a \textit{Dungeons of Drakkenheim} campaign, the player characters are the heart of the story. Their choices, conflicts, and personal goals could decide the fate of Drakkenheim, if they survive! Below are guidelines for creating a player character ready to take on the dark city. We recommend players and Game Masters collaborate closely during character creation for the best possible results.

\subsection{Character Options and Backgrounds}
Adventurers of any race or class in the Core Rules can be found in Drakkenheim.


\begin{itemize}
\item Appendix E: The World of Drakkenheim describes how many character ancestries, religions, martial arts, and magical traditions fit into the wider world.
\item Appendix F: Character Backgrounds includes new backgrounds players may use to closely connect their characters to the events, people, and mysteries in Drakkenheim.
\end{itemize}

\subsection{Personal Quests}
A character's core motivation for adventuring in Drakkenheim is represented by their personal quest. This is a specific individual objective for player characters to achieve during the campaign, but players need not reveal their personal quest to each other.

When a character completes their personal quest, they choose one of the following rewards:


\begin{description}
\item[Ability Score Improvement.]
One ability score of your choice increases by 2, as does your maximum for that score.
\item[New Feat.]
You gain a bonus feat of your choice, subject to the GM's approval.
\end{description}

Player characters may select from the quests below, or work with the Game Master to design one:

\subparagraph{Apocalyptic Vision}
I have been haunted by visions of the meteor for fifteen years. Every time I sleep, I see it crashing down upon me. My personal quest is to see for myself what lies at the heart of the crater. I must overcome whatever I find there in order to stop these nightmares.

\subparagraph{Arcane Secrets}
Forbidden lore lies in Drakkenheim. Several dark wizards living in the ruins have mastered terrible and powerful new magic. The Inscrutable Tower library stands unguarded, filled with spellbooks containing secret arcana. My personal quest is to learn these spells.

\subparagraph{Blighted Landscape}
A blight on the natural order envelops Drakkenheim, but the contamination could spread beyond the city. My personal quest is to study the unnatural phenomena found within the most corrupt places in the city: Queen's Park Garden, the underground waterways, and the Crater's Edge. Perhaps then I can develop a way to reverse the contamination.

\subparagraph{Claim the Throne}
Though my family has fallen on hard times, we were distantly related to the royal House von Kessel. I may have a claim to the throne of Drakkenheim. My personal quest is to search the cathedral, noble estates, and the castle for sufficient evidence to conclusively prove my inheritance. Then I will take the crown for myself.

\subparagraph{Faction Aspirant}
I believe strongly in the ideals and goals of one of the five factions: either the Amethyst Academy, the Followers of the Falling Fire, the Hooded Lanterns, the Knights of the Silver Order, or the Queen's Men. My personal quest is to join the faction, aid their mission, climb the ranks, and rise to a high leadership position within their organization.

\subparagraph{Family Matter}
One of my family members or close friends has become mixed up with one of the factions. Perhaps they have foolishly decided to join the organization. Alternatively, they've been taken prisoner by the faction due to some trespass or misunderstanding. Either way, my personal quest is to bring them back home.

\subparagraph{Find a Cure}
I have a friend, family member, or loved one who was contaminated by delerium and turned into a monster. They have gone mad, and are lost somewhere in Drakkenheim. My personal quest is to locate them, then discover a way to reverse their horrific transformation.

\subparagraph{Lost Heirloom}
I used to live in Drakkenheim. My personal quest is to recover a lost heirloom from my old home, which was located somewhere in the Inner City. Perhaps the item is sentimental, valuable only to members of my family. Alternatively it may have magical properties, and I wish to use these powers to aid my own ambitions.

\subparagraph{Monster Slayer}
I wish to leave behind a heroic legacy by defeating dread beasts. My personal quest is to slay three of the greatest monsters rumoured to dwell in Drakkenheim, and bring back their heads as trophies. Each must be a notable creature with a challenge rating greater than my own level when I defeat them. My allies can assist me, but I must strike the decisive blow myself!

\subparagraph{Overwhelming Debt}
I owe a tremendous debt of 10,000 gold pieces, maybe more. My personal quest is to pay it back by making a fortune in Drakkenheim collecting delerium, plunder, and completing work-for-hire. Meanwhile, I must hide from bounty hunters my creditor send to kill me.

\subparagraph{Personal Pilgrimage}
As an acolyte of the faith, my personal quest is to visit several holy places in the city and perform a sacred ritual before the remains of the martyrs interred within each: the Cathedral of Saint Vitruvio, Saint Selina's Monastery, and the Chapel in Castle Drakken. Through these devoted acts I will understand my place in the divine order.

\subparagraph{Score to Settle}
I have come to Drakkenheim to kill one of the faction leaders: either Archwizard Eldrick Runeweaver, \textbf{Knight-Captain Theodore Marshal}, \textbf{Lord Commander Elias Drexel}, \textbf{Lucretia Mathias}, or the \textbf{Queen of Thieves}. Perhaps I seek revenge for a past wrong. Alternatively, I may covet their position and wish to replace them. Either way, my personal quest is to slay them.

\subsection{Character Relationships}
The campaign begins with the player characters setting out on their first foray into the ruins, but how did they meet? Consider establishing a personal connection between two or more player characters, such as one of the following:

\textbf{...siblings, cousins, or childhood friends}.

\textbf{...a shared origin, religion, or teacher}.

\textbf{...former rivals turned unlikely allies}.

\textbf{...one saved the life of the other}.

\textbf{...saw each other in a dream or vision}.

\section{Four Quick Facts about the World of Drakkenheim}
Appendix E: The World of Drakkenheim explains these background details further, but the setting of Drakkenheim diverges from the default assumptions of Fifth Edition in four key ways:


\begin{description}
\item[The gods are silent and distant.]
They don't manifest physically nor speak with worshippers, and do not interfere in earthly matters. Planar cosmology is mysterious and unknowable, but sages have still developed many conflicting (and wildly incorrect) theories.
\item[Divine spellcasting powers aren't granted by gods.]
Instead, clerics, druids, and paladins tap into sacred energies through devotion, meditation, and resolve. Violating one's religious tenets won't result in a loss of this ability, but a crisis of faith might.
\item[Individuals can't become wizards through study alone.]
The arcane magic of both sorcerers and wizards alike is borne in the blood, which may be harnessed via practice and study. The laws of the land bar arcane spellcasters from holding noble titles.
\item[High-level NPCs and spellcasters are exceptionally rare in the wider world.]
Despite their few numbers, such people still have a profound impact on society. However, Drakkenheim has attracted an uncommonly high concentration of such individuals.
\end{description}

\section{The Original Campaign}
The adventure in this book is based on the \textit{Dungeons of Drakkenheim} live stream campaign broadcast by the \textbf{Dungeon Dudes} which ran from October 2018 to December 2019. It featured Kelly McLaughlin as Sebastian Crowe, Jill Danaitis as Veo Sjena, Joe O'Gorman as Pluto Jackson, and Monty Martin as Game Master.

You can watch it on the Dungeon Dudes YouTube Channel (www.youtube.com/dungeondudes) and find it on most major podcast platforms. Check it out if you're looking for inspiration for how we ran the adventures in this book, but be warned: spoilers abound! We've also created a wealth of instruction videos for players and Game Masters covering everything about tabletop roleplaying games, from character building to designing campaigns.

Every major event, location, and non-player character is detailed in this book, and we've included many more completely new ones for you to explore in your own game. In many ways, the original campaign was a playtest for this book, and we've refined and updated several sections since. If you've watched the original series and notice an inconsistency between this book and the show (there are many), the version in this book takes precedence. On the flipside, you should feel empowered to alter, change, and expand the material here, especially if doing so builds stronger connections or better stories for you and your players.

The events of our original campaign are non-canonical with regards to this book, so several elements of Drakkenheim strongly connected to the original characters have been altered or removed. If you wish, the original "Drakkenforce" characters may appear in your campaign as rival adventurers, but their stories have been set aside herein so your players may take centre stage.

\chapter{Running the Campaign}
\textit{Dungeons of Drakkenheim} is a nonlinear adventure that combines urban exploration and faction conflict, bound together by the personal quests of the player characters. Early in the campaign, your players might closely follow the story overview described below. As they advance further into the city and interact with the factions, their decisions will reshape the story in unpredictable ways.

\DndDropCapLine{T}{}his book contains inspiration, guidelines, and tools to help you adjudicate and respond to your player characters' actions as they explore Drakkenheim. Running an exciting and dynamic nonlinear adventure requires sound judgement and good use of logic on your part. You should feel empowered to make changes so your players have a rewarding and challenging experience. Occasionally, you may even need to devise your own material or repurpose adventure hooks, characters, and locations to better suit the changing conditions in your campaign.


\includegraphics[width=0.4\textwidth]{images/../5e.tools/img/adventure/DoDk/001-01-001.outskirt-of-drakkenheim.png}
\section{Story Overview}
The characters arrive in Drakkenheim to pursue their personal quests or a few freelance missions from friendly contacts in Emberwood Village. They face ravenous ratlings, malfeasant wizards, and rival adventurers in the Outer City Ruins. During their first forays into the city, the characters witness the strange magic of delerium and contend with the mysterious Haze. The threat of contamination makes survival in Drakkenheim difficult and dangerous. The characters may seek assistance and resources from the five factions to accomplish their goals and find a way deeper into the city.

As the characters learn more about the factions' differing ideologies, the various faction leaders take notice of their progress. Once the characters successfully traverse the City Walls and begin exploring the Inner City Ruins, each leader approaches the characters for help with their objectives. They appeal to the characters' Ideals and Bonds, offer resources and support, and promise crucial assistance with the characters' personal quests.

However, as the characters collaborate with the factions to advance their mutual goals, they become increasingly entangled in the political machinations between the competing groups. Each organization has a different vision for the future of Drakkenheim, and the tenuous peace won't last much longer. Nevertheless, with the support of one or more of these groups, the characters can explore dangerous Inner City locations such as the Cathedral of Saint Vitruvio, the Inscrutable Tower of the Amethyst Academy, or perhaps even mount an expedition into the crater. Characters may unearth stunning revelations from Drakkenheim's past and recover powerful artifact which threaten the balance of power between the factions. They might even discover how the Haze threatens to engulf the world itself.

The factions are deeply divided over what must be done to solve these problems. As a result, simmering tensions boil over into outright conflict. As the city streets become a battlefield between the organizations, the faction leaders each demand the characters' loyalty. Now, the characters must decide which factions are their allies, and which are their enemies. Eventually, the unraveling situation may require the characters to act directly against one or more of the factions, who respond with counterstrikes of their own. In turn, the characters may launch major assaults against enemy factions to destroy their strongholds, defeat their leaders, and drive them out of Drakkenheim. As the conflict reaches its zenith, the surviving factions set their master plans into motion, requiring an expedition into Castle Drakken. There, a final showdown will decide the fate of Drakkenheim.

\section{Character Advancement}
Character advancement in Drakkenheim may be handled by using traditional experience points. In addition to earning XP for defeating monsters, you should award additional experience points (equal to the reward for defeating a level-appropriate monster) whenever characters:


\begin{itemize}
\item Finish a personal quest
\item Complete a faction mission
\item Slay a legendary monster (either a monster or faction leader)
\item Obtain a \textit{Seal of Drakkenheim} or a \textit{Relic of Saint Vitruvio}
\item Discover secret information about the meteor or Drakkenheim
\end{itemize}

Alternatively, you can choose to dispense with tracking experience points and simply decide when the characters gain a level.

The flowchart shows a possible path characters may take through their adventures in Drakkenheim, along with suggested levels to provide an appropriate challenge for the characters. Use this as a guideline for how fast characters should advance, and award character levels when they complete a personal quest or major faction mission.

\subsection{Adventure Flowchart}
\includegraphics[width=0.4\textwidth]{images/../5e.tools/img/adventure/DoDk/adventure-flowchart.png}
\subsection{Deadly Challenges}
Drakkenheim is a dangerous place. Not every combat encounter is fairly balanced in a level-appropriate manner. Deadly threats lurk down every street, and many scenarios intentionally stack the odds against the player characters. They must use cunning, ingenuity, and teamwork to survive and seek the aid of the factions to accomplish their most important goals.

\subsection{Character Death}
Characters looking to revive a slain companion have limited options. Few non-player characters are capable of casting the necessary spells, and the expensive material components are rare. Prominent members of the Silver Order and the Followers of the Falling Fire can cast \textit{raise dead} and other higher level spells. Consult the Faction Boons section in chapter 3 regarding each respective faction for details on what each faction requires. Characters might also seek out the druid of the Shrine of Morrigan (see chapter 6). Finally, Spell Scroll of \textit{revivify}, \textit{raise dead}, and \textit{resurrection} may be found as treasure in various locations in the city, usually the chapels and cathedrals. At your option, the characters may find another scroll elsewhere in the ruins.

Regardless, the players may need to create a replacement character, or play one temporarily. A player may easily create a new arrival to Emberwood Village to replace a slain character, or perhaps be encountered as the lone survivor of another adventuring party. Alternatively, they might play a faction member assigned to the party as a reinforcement.

\DndQuote%
{}
{''My best tip for staying alive in Drakkenheim is to just kill everything that moves. No time to ask questions. Hesitation is dangerous.''}
{Sebastian Crowe}
\section{Origins of the Meteor}
\DndQuote%
{}
{''Magic attracts more magic. Under extraordinarily rare circumstances, magical energies crystallize into delerium. The stones slowly draw more raw magic into the world from distant planes, possible realities, and alternate dimensions. The rate is exponential. Considerably more leakage occurs around vast quantities of delerium. These chaotic energies are inherently destabilizing. Unfocused magic conjures up paradoxical improbabilities at random, leaving behind lingering eldritch contamination. Uncontained, these magical forces manifest as the Haze. More delerium forms within, and so the process accelerates.''}
{Archwizard Ryan Greymere}
The meteor left behind vast deposits of delerium, enough to manifest an erratic magical field called the \textbf{Haze}. New delerium deposits gradually emerge within the Haze, which causes the affected region to expand slowly over time. Left unchecked, the Haze can overtake an entire world within a few centuries, leaving behind a chaotic landscape of churning madness and wild magic. The infected world cataclysmically shatters when more destructive magic inevitably erupts, sending a hail of delerium-filled comets throughout the cosmos. After hurtling through the Astral Void for eons, some collide with new worlds where the process begins again. Countless worlds have met this fate, the rest are wholly ignorant of the doom that might one day descend upon them.

If nothing is done to stop the spread, the Haze will fully engulf Westemär in roughly one hundred and fifty years. Once this happens, little hope for reversal remains, and the dread mists will envelop the world two centuries thereafter. Civilization and society collapse into a monster-filled wasteland that catastrophically ruptures about one thousand years later.

None know the complete truth about the Haze at the outset of the campaign, and little knowledge is revealed via divination magic save mad whispers. Throughout the adventure, player characters may find clues about delerium from which they may develop their own theories and speculations regarding the origins of the meteor. They can learn much by exploring deep into the ruins, especially locations such as the Inscrutable Tower, the crater, and Castle Drakken. Faction leaders and malfeasant wizards might also offer their research and insights to help players along the way.

\section{Fall of House von Kessel}
House von Kessel ruled Westemär from the capital of Drakkenheim for one hundred fifty years. The last monarch, Ulrich IV, inherited a period of relative political stability, and so pursued interests in architecture, education, and carnal affairs instead of politics and warfare. He married a Caspian princess, Lenore, and they had three children: Leonard, Katarina, and Eliza. Ulrich IV was ten years into his reign when the meteor struck Drakkenheim. The King, his wife, and their children were never seen again after that day, but their true fate remains a mystery.

However, many members of the royal household were not in Drakkenheim when the meteor struck, including the King's younger siblings, Mannfred and Cecilia von Kessel. The two spent vast sums commissioning the clergy of the Sacred Flame and the mages of the Amethyst Academy to use divination magic to discern what happened, but to no avail. They nearly bankrupted themselves by launching multiple large-scale military expeditions attempting to reclaim Drakkenheim, all of which failed. Eventually, Mannfred became convinced further efforts were futile, and proposed he would take up the crown and relocate the capital. However, Cecilia contested the claim, arguing there was no conclusive evidence that confirmed their brother and his children were dead. The nobility was divided over the issue, and a civil war ensued that raged for nearly a decade. As Westemär itself descended into political and financial ruin, Mannfred von Kessel and his children were assassinated in a dramatic betrayal. Cecilia von Kessel herself died unexpectedly only a few days later, leaving behind no heirs of her own. Lacking a clear successor, the Civil War ended with a whimper.

\subsection{The Royal Family}
Since the royal household were recognizable public figures in their lifetimes, any character proficient in History instantly recognizes them in artwork or portraits.


\begin{description}
\item[King Ulrich von Kessel IV.]
was a portly man of average height with a wiry beard and rather unremarkable features, artwork often exaggerates his appearance with a heroic barrel chest, a lush curly beard, and a regal countenance. Stately artwork shows him clad in a royal mantle wearing the \textit{Crown of Westemär}. Those who knew the King in life remember him as a capable if overindulgent ruler who left most matters of state to his administrators.
\item[Queen Lenore von Kessel.]
was a woman in her late thirties; six feet tall with graceful limbs, hair dyed bright orange, and bleached white skin. Her face had chiseled and angular features, and with every gesture she struck a pose. She wore outrageously ornamented gowns and jewelry, including a famous golden necklace set with eleven emeralds.
\item[Leonard von Kessel.]
was eighteen years old when the meteor struck, but would be thirty three years old now. Leonard took to military interests at a young age. Artwork often depicts him as an aristocratic young soldier with sandy brown hair wearing a dragon-decorated breastplate and wielding a longsword.
\item[Katarina von Kessel.]
was an inquisitive and acrobatic eleven-year-old with wide eyes and auburn hair. She typically wore a frilled shirt, fitted trousers, and cape instead of dresses, and had a penchant for stealing hats. She would now be twenty six years old.
\item[Eliza von Kessel.]
was only six years old at the time of the disaster. The button-nosed child was usually dressed in a bonnet and patterned dress. Assuming she survived, she would be twenty one years old today.
\end{description}

\subsection{A Long Lost Heir}
None have yet discovered that King Ulrich von Kessel IV and Queen-Consort Lenore von Kessel were both transformed into monsters. Ulrich is now part of the hideous \textbf{Amalgamation} in Castle Drakken, and Lenore has become a mad creature hiding in Queen's Park. Their children haven't been seen in fifteen years. It's up to you to decide whether or not they survived, and where they might be found.

\includegraphics[width=0.4\textwidth]{images/../5e.tools/img/adventure/DoDk/002-01-002.family-portrait.png}
Beyond the immediate offspring of Ulrich IV, however, many relatives and cadet branches exist. As such, the House von Kessel family tree has been left intentionally vague so you and your players may freely insert your own player and non-player characters as potential heirs. They need not be one of Ulrich's known children, but instead illegitimate offspring, a grandchild, niece, nephew, or other family relative. The most dramatically compelling possibilities include:


\begin{itemize}
\item One or more player characters (known or unbeknownst to them)
\item A close relative of a player character (offspring, half-sibling, cousin, or niece/nephew)
\item The \textbf{Queen of Thieves} or a faction lieutenant (see the NPCs entries for specific ideas)
\end{itemize}

Furthermore, there is no need for there to be a single surviving heir, or even any at all. Anyone can say they have royal blood. It's another matter to prove it, and gain the support to actually take the crown. Finally, it's possible another noble household could gain sufficient political clout and military power to turn a tenuous claim into a credible one, and simply supplant House von Kessel entirely.

\subsubsection{Complications}
At the outset of the campaign there is no heir apparent to the throne, nor have any credible claimants to the throne stepped forward. There are many reasons why a potential claimant to the throne hasn't emerged yet. Decades of affairs, disaster, and war have ravaged the public records, and conclusively determining heredity via magical means requires the blood or remains of Ulrich IV. While an heir could possess a locket, signet ring, or some other document, this must be substantiated with something from the ruins to prove their identity and ancestry. Consider introducing one of the following additional complications:

\includegraphics[width=0.4\textwidth]{images/../5e.tools/img/adventure/DoDk/003-01-003.meteorite-delerium.png}

\begin{itemize}
\item Kept as a long-forgotten prisoner in an iron helmet in the Queen of Thieves' dungeon.
\item Grew up adopted; no idea about their true parentage.
\item Working as a humble stablehand living in Emberwood Village.
\item Secretly hidden and protected by a distant noble house.
\item Petrified into a statue in Castle Drakken.
\item Accidentally banished to another plane.
\item Succumbed to contamination; transformed into a monster.
\item Secretly an arcane spellcaster; disinherited under the Edicts of Lumen. (See Appendix E)
\item Target of an assassination attempt by faction agents; have been hiding since.
\item Born from an affair; their birth was kept secret to conceal a scandal.
\end{itemize}

\section{Embracing the Unknown}
What has happened to Drakkenheim is simply a cosmic accident. It is not the intentional work of an inscrutable extraplanar being, nor did any chortling demon lord or scheming evil deity set these events in motion. Instead, it is more akin to a natural disaster, albeit one caused by the magical forces that underscore the cosmos. Much like an earthquake, hurricane, or volcano, the calamitous forces that cause such magical catastrophes are entirely beyond mortal means to control or reverse. This fact may come as a surprise to players who expect they can simply defeat some tangible evil entity and save the world. However, the dark fantasy and cosmic horror literature that inspired this adventure rarely ends in such clear-cut ways. Instead, the heroic actions of the player characters may avert further disaster, but will never erase the damage that has already been done.

While outwardly presenting a tantalizing mystery, the meteor functions more as a catalyst for conflict in Drakkenheim. My players failed to fully uncover the origins of the meteor during their adventures, but our campaign still had an immensely satisfying conclusion. In fact, themes of cosmic horror and madness are greatly enhanced when the real answers are never fully understood.

\subsection{The Seals of Drakkenheim}
These six wondrous items were entrusted to the royal council of Westemär. Each identifies the bearers as important administrative and political figureheads. In addition to their potent magical powers (described in Appendix D), the seals have an important purpose in a special coronation ritual that attunes a new monarch to the \textit{Crown of Westemär} itself (described in chapter 9). As such, many factions seek these seals, which are thought lost in the ruins. A few are already in the hands of several faction agents however! Here's where you can find them though:

\subparagraph{Chancellor's Crest}
Balrin Barthes failed to escape Drakkenheim, and was transformed into a monster. This mutated wretch still wears the Chancellor's Crest, and serves the horrific Duchess deep in its lair below the cistern.

\subparagraph{Inscrutable Staff}
Adriana Modiera, the Archmage of Drakkenheim, died within the Inscrutable Tower. This mighty magical staff remains there today. Traditionally, the Amethyst Academy has the right to determine who holds this staff.

\subparagraph{Lord Commander's Badge}
Elias Drexel survived the fall of Drakkenheim and spent years fighting in the Civil War before taking charge of the Hooded Lanterns. He still holds his original badge of office.

\subparagraph{Saint Vitruvio's Phylactery}
High Flamekeeper Gregora led the faithful in the Cathedral of Saint Vitruvio. Her true fate is unknown, but now a monstrous garmyr priest carries the phylactery in the cathedral. It is traditionally entrusted to the highest ranking cleric of the Sacred Flame in Drakkenheim, as this item is also one of the Relics of Saint Vitruvio.

\subparagraph{Spymaster's Signet}
Spymaster Magdalena Vasari was a half-elf assassin. A loyal operative through-and-through, Vasari took advantage of her assumed death to secretly seek any surviving heirs of the royal family during the Civil War. She may have succeeded had she not been killed and her signet ring stolen by the \textbf{Queen of Thieves} several years ago.

\subparagraph{Steward's Seal}
Johann Eisner, the royal steward, became trapped in Castle Drakken after failing to rescue Ulrich IV. However, he entrusted his young daughter with hiding his seal as a precaution for this exact predicament. While her fate is now unknown, the seal is now amongst the treasures of the \textbf{Crimson Countess} in the Cosmological Clocktower.

\section{Running Personal Quests}
These adventure hooks give the players their own personal objectives for venturing into the ruins. Players might choose one of the quests presented in chapter 1, or you may work with the player to invent one of your own design.

The personal quests have been kept open-ended so you can work with your players to deeply customize them. You should feel free to integrate unexpected twists or complications along the way, and establish strong connections to the city, events, and factions in the process. Here's a few suggestions for how these quests might proceed:

\subparagraph{Find a Cure}
The character's relative may be found in an area of the Deep Haze: near the Inscrutable Tower, Queen's Park, the crater, or perhaps in thrall to the Duchess or the \textbf{Pale Man}. Unfortunately, they are indeed mad, and once found resist capture and constantly try to escape. A monstrous transformation can be reversed in three ways: using a \textit{wish} spell, by \textbf{Lucretia Mathias} performing a divine intervention, or by researching the \textit{siphon contamination} spell.

\subparagraph{Claim the Throne}
This quest ties the character strongly to a central conflict in the campaign. A fun twist is introducing a rival adventurer with an equally valid claim. Consult that section for more details. This character could even be one of the children of Ulrich IV.

\subparagraph{An Apocalyptic Vision}
This character will need to seek out the means to survive the most contaminated area in the city. Once they reach the heart of the crater, the character has another apocalyptic vision when they touch the Delerium Heart. The character must now decide what they will do to save the world.

\subparagraph{Faction Aspirant}
Advancement in a faction requires a full commitment from the character, exclusive of any other faction. Achieving this quest happens when a faction's leader regards the character as a capable and trusted confidant, and appoints them as their right-hand lieutenant. The most challenging aspect of this quest are the conflicts of interest it may cause with the other player characters, but embrace this conflict as a sticking point between the characters!

\subparagraph{A Score to Settle}
This is a very difficult personal quest, but a straightforward one. The faction leaders are all well-protected and formidable in their own right, though Elias Drexel and Theodore Marshal are significantly easier to encounter in person. Characters targeting the \textbf{Queen of Thieves} or \textbf{Eldrick Runeweaver} are in for a vexing challenge!

\subparagraph{Arcane Secrets}
These spells can be learned from \textbf{Oscar Yoren}, the \textbf{Pale Man}, Ryan Greymere and spellbooks contained within the Inscrutable Tower; the character must seek knowledge from each. Collecting these spells raises the ire of the Amethyst Academy, who wish to keep knowledge of this magic secret. Characters can complete this quest by researching the \textit{contaminated power} spell.

\subparagraph{The Blighted Landscape}
After visiting each location, the character can research the Aqua Expurgo, \textit{neutralizing field}, and \textit{contamination immunity} spells respectively, even if they haven't met the other prerequisites to research these spells.

\subparagraph{A Family Matter}
Captive or convert, the relative is found at a faction stronghold. They could even be one of the faction lieutenants! It is difficult to convince them to leave; they have to see the faction do something horrible with their own eyes.

\subparagraph{Overwhelming Debt}
Rival adventurers, a Queen's Men strike team, or an assassin might be pursuing a bounty on the character. You may want to have the Queen of Thieves or the Amethyst Academy "acquire" the character's debt and use that as leverage over them!

\subparagraph{A Lost Heirloom}
Any rare magic item is an appropriate heirloom, and you can choose any location in the city where the item is lost, although Reed Manor, the Kleinberg Estate, and Slaughterstone Square are good possibilities. If you'd really like to create complications, this item could be stolen by the Queen of Thieves or a rival adventurer, or could even be one of the Seals of Drakkenheim.

\subparagraph{Personal Pilgrimage}
The ritual the character wishes to perform at each location need not be simple. Instead, it might be a time-consuming process that requires a rare offering and takes several hours to complete, or perhaps can only be performed at a certain time of day (either dawn or dusk). As such, the character will need the protection and support of their allies to stay safe during the rite.

\subparagraph{Monster Slayer}
Suitable foes might include the \textbf{Rat Prince}, the \textbf{Crimson Countess}, the Duchess, the \textbf{Lord of the Feast}, the Prince of Carnage, a \textbf{grotesque gargant}, or a \textbf{crater wurm}.

\subsection{Finishing a Personal Quest}
Upon completing a personal quest, a player character gains rewards as described in chapter 1. In addition, you should reward each character in the party experience points based on the advancement method you are using.

At your option, once a character completes their personal quest, you might work with the player to create a follow-up personal quest for them to reflect the character's ongoing and evolving story. Furthermore, if conditions emerge where a character either can't complete or decides to abandon their personal quest, you should work with the player to develop a new one.

\subsection{Creating a Custom Personal Quest}
When designing a custom personal quest, keep the following principles in mind:


\begin{itemize}
\item The quest should require characters to explore one or more locations within the Inner City ruins. A character shouldn't be able to complete their personal quest too early. Conversely, a personal quest involving Castle Drakken likely can't be completed until late in the campaign.
\item It should present some sort of complication with one of the factions, or perhaps require the characters to seek help from the faction to complete their quest.
\item The quest should be specific and achievable, and not related to character advancement (such as reaching 10th level or taking a specific feat).
\item Be careful about using a "ticking clock" for a personal quest. While this can introduce a compelling complication, focus on long-term goals rather than extremely urgent matters.
\end{itemize}

Players often suggest personal quests involving missing persons, lost objects, or revenge. Many example personal quests provided already cover such motivations, such as "A Family Matter'', "A Lost Heirloom", and "A Score to Settle". Ask players to consider these options first. Perhaps their idea can be easily adapted to work with one of these quests (or vice versa).

\section{Running the Factions}
Five factions are struggling to control the ruins of Drakkenheim: the Hooded Lanterns, the Queen's Men, the Knights of the Silver Order, the Followers of the Falling Fire, and the Amethyst Academy. Each is commanded by a prominent faction leader, but the faction as a whole behaves like a major recurring NPC, except they are composed of many individuals working together who share the same ideology, goals, methods, and resources. Chapter 3: Factions describes the background, objectives, roleplaying traits, major NPCs, and assets of each faction. Using this information, it's up to you to determine how the factions might help or hinder the player characters during their adventures in Drakkenheim.

Rather than uniting the factions as one against a common foe, the player characters must navigate a fractious web of intrigue, negotiation, sabotage, and violent conflict if they wish to uncover the greatest treasures and secrets in Drakkenheim. Their relationships with each faction will naturally develop as a consequence of the character's choices, but it's ultimately up to them to decide who they want to join and who they will oppose. Thus, each individual faction might become a supportive patron, fierce competitor, close ally, bitter enemy, or anything in between.

Exploring the most dangerous locations in Drakkenheim is extraordinarily difficult without gaining the support (or inciting the ire) of one or more factions. Correspondingly, the factions need the player characters' help to accomplish their own goals in the city, and offer rewards to their allies and assistance with the characters' personal quests. Conversely, factions that become antagonistic create dynamic problems and new complications throughout the campaign when they use their resources against the characters. Either way, the factions are best played in a reactive manner. Their leaders make firm demands when characters try to negotiate, orchestrate counterstrikes when characters foil their plans, and conduct sabotage when characters oppose their goals.

Eventually, these interactions may blossom into a full-scale hostility between the factions. Once this occurs, aim to keep each faction in play as long as possible. Don't be afraid of letting the players get in over their heads if they make enemies with more than one faction, either. Your players will find surprising ways to foil the factions, and develop cunning plans when they collaborate with them. Embrace their creativity by responding with new challenges, and rewarding players when they succeed by having their allies recognize their accomplishments. The factions are equipped with the tools to help you escalate tension and raise the stakes of the conflict in interesting ways.

Many adventure sites provide hooks and missions connected to the factions you may use to prompt the characters to visit each area. Each location also details how the factions could respond to possible events and outcomes that may occur.

\section{Cast of Characters}
\textbf{Lord Commander Elias Drexel} (grim leader of the Hooded Lanterns)

\textbf{Captain Ansom Lang} (determined Hooded Lanterns Lieutenant)

\textbf{Lieutenant Petra Lang} (hopeful Hooded Lanterns Lieutenant)

\textbf{Knight-Captain Theodore Marshal} (valiant leader of the Silver Order)

\textbf{High Flamekeeper Ophelia Reed} (devoted Silver Order Lieutenant)

\textbf{Archwizard Eldrick Runeweaver} (calculating leader of the Amethyst Academy)

\textbf{River} (pragmatic Amethyst Academy Lieutenant)

\textbf{Lucretia Mathias} (mystical prophet of the Falling Fire)

\textbf{Nathaniel Flint} (enthusiastic Falling Fire Lieutenant)

\textbf{Queen of Thieves} (enigmatic criminal mastermind)

\textbf{Blackjack Mel} (sleazy Queen's Men Lieutenant)

\textbf{King Ulrich von Kessel IV} (monarch of Drakkenheim, assumed deceased)

\textbf{Queen-Consort Lenore von Kessel} (royal consort, assumed deceased)

\textbf{Leonard, Katarina, and Eliza von Kessel} (royal heirs, fate unknown)

\subsection{Negotiating an Alliance}
Despite their conflicts, the faction leaders know they can't achieve their goals on their own. Not only do they seek the aid of the player characters, many factions see the benefits to collaborating with each other. Diplomatic characters could negotiate an alliance, truce, or ceasefire between two (or more) factions. Convincing any two factions to set aside their differences and work together should be difficult, but achievable. In order to broker such a deal, characters will need to demonstrate how the alliance will help both factions achieve their objectives, but also make an appeal to the roleplaying traits, ideals, bonds, and flaws of the factions and their leaders. These tense negotiations can make for some of the most memorable moments in the campaign, so potential alliances are best arranged in dramatic locations such as the faction strongholds and the Cathedral of Saint Vitruvio.

The factions have been designed such that a long-term alliance between any two factions is achievable, and a short-term truce involving three is difficult but entirely possible. Countless situations in your campaign may emerge that could cause the factions to re-evaluate their priorities, and clever and convincing players might present compelling arguments. Bear in mind of course, that the factions are composed of people: and people don't always act in logical ways. These alliances are always unstable, however, as each faction presses to have their own goals prioritized. Unless the player characters make an ongoing effort to keep a coalition together, any alliance between factions will inevitably collapse.

Factions such as the Amethyst Academy and the Queen's Men may even use alliances as a convenient opportunity to manipulate the player characters and the other factions, and possibly even make a dramatic reversal at a critical moment. Especially cunning characters might help the factions weave such deceptions and plots, resulting in brutal betrayal against an unsuspecting ally which allows the faction to seize full control of Drakkenheim for themselves.

Since the threat of betrayal is so real, player characters who build a coalition between multiple factions may compromise the trust they earn with an individual faction. They won't receive the greatest boons the faction can offer, and faction agents don't trust characters with sensitive information. Faction leaders are wary of double-agents, and would much rather see the player characters make a commitment to them. As a result, the player characters should constantly be called upon to demonstrate their loyalty.

\includegraphics[width=0.4\textwidth]{images/../5e.tools/img/adventure/DoDk/004-01-004.assassination.png}
\subsection{Resolving the Conflict}
Unfortunately for the player characters, discovering the full properties of delerium and facts regarding the meteor's origins will not unite the factions. Instead, these revelations will drive further wedges between the fractious groups since they deeply disagree on the best way to address the problem. Given that The Haze threatens to overtake the entire world, it may at first appear the Silver Order is completely correct that the crystals must be destroyed. Throughout the campaign however, the feasibility of accomplishing such a task should be consistently cast in doubt. You should strive to contrast this viewpoint against the rational yet questionable stance of the Amethyst Academy which insists arcane means may harness delerium to prevent further catastrophe, and the zealous and apocalyptic prophecies of the Followers of the Falling Fire. Whenever it seems like the players have made up their minds, rival factions should present a gnawing sense that some alternative might still be possible. You should avoid presenting any "correct" answer: It's up to the characters to piece it all together, draw their own conclusions, and decide how to act.

As such, there is no fixed resolution to the faction conflict. By the end of the campaign, several factions could be routed from Drakkenheim entirely, their leaders slain and strongholds toppled. While a showdown in Castle Drakken is inevitable, who exactly is present for that battle and the outcomes that may arise are unforeseeable. Chapter 10: The Fate of Drakkenheim describes a few possibilities.

\DndQuote%
{}
{''A lot of people seem hopeful to find the royal family alive in the city, as if that's going to fix everything. First off, they weren't all that great, second of all, the chances they are alive are.... dismal.''}
{Sebastian Crowe}
\section{Rival Adventurers}
Drakkenheim attracts all manner of would-be heroes and fortune-seekers. Characters may often encounter these rival adventurers during the campaign as recurring NPCs, both within the ruins and in Emberwood Village. At any given time, there might be as many as a dozen other groups exploring Drakkenheim, perhaps more.

\includegraphics[width=0.4\textwidth]{images/../5e.tools/img/adventure/DoDk/005-01-005.lock.png}
Rivals may be competitive and quarrelsome, or friendly and sporting. Either way, good rivals spur the characters into action and create unique roleplaying opportunities. You might even invite a player to establish a pre-existing relationship between their character and a rival during character creation. Friendly rival adventurers might interact with characters in the following ways:


\begin{itemize}
\item Prompt the characters to recount their own stories and deeds
\item Share rumours and gossip with the characters
\item Suggest the characters visit a specific adventure site or take on a new mission
\item Challenge the characters to a contest (such as archery, wrestling, magical displays, drinking, slaying a monster, finding a treasure, or a duel)
\item Team up with the characters to take on a difficult objective
\end{itemize}

Alternatively, some rival adventurers might be more hostile towards the characters:


\begin{itemize}
\item Trash-talk and humiliate the characters in Emberwood Village
\item Belittle the characters' accomplishments and insult their skills
\item Complete a quest or explore a location the characters have ignored
\item Act as reinforcements for an enemy faction
\item Take up a bounty on the characters
\end{itemize}

The factions often hire adventuring parties to help achieve their goals. If the player characters don't take up a mission or offer, their leaders might turn to rival adventurers instead! It is also possible that rival adventurers could start as a friend or enemy, and end up the opposite!

\subsection{Rival Adventuring Parties}
You can use the sample adventuring parties below, or create one of your own design.

When creating your own adventuring party, use pre-generated stat blocks from the Core Rules rather than spending time making character sheets for these NPCs like you would a player character. When it comes to making a good rival, interesting roleplaying traits are far more important than game statistics. A rival works well when they reflect the personalities and flaws of the player characters themselves, then exaggerate their best and worst qualities to the extreme while sharing some common ground. Diametrically opposing those traits tends to result in antagonists, not rivals. Flesh out such characters by pulling in one or two roleplaying traits from one of the backgrounds or personal quests presented in this book.

\DndQuote%
{}
{''My favourite game is to tell a group of adventurers where great treasure in the city is. Frankly, I have no idea where treasure is, but now they will stay out of my way.''}
{Sebastian Crowe}
\subsubsection{Caspian Expedition}
This expedition was dispatched by a Caspian noble house.


\begin{itemize}
\item Jupiter Jones is a barrel chested human \textbf{gladiator} with a chiseled jaw and herculean features. Though he claims he's adventuring in Drakkenheim to bring glory to his name, he's actually on a personal quest to prove his family's connection to House von Kessel, as he is in fact Queen Lenore's nephew.
\item Lemminton Withalf, a greasy, misanthropic, and sniveling gnome bard (use the game statistics for a \textbf{scoundrel} with bard spells instead of wizard spells) serves as herald.
\item A personal household guard of a dozen thugs.
\end{itemize}

\subsubsection{Heroic Veterans}
This adventuring team has a small reputation in Drakkenheim. Its members include:


\begin{itemize}
\item Pluto Jackson, a human \textbf{knight}. This brash Caspian prince has come to Drakkenheim ostensibly to slay the greatest monsters in the city to make a name for himself. In truth, he's searching for proof his nephew is a legitimate heir to the throne.
\item Veo Sjena, a catfolk \textbf{urban ranger}. This mangy alleycat has survived in Drakkenheim for fifteen years and knows the city better than anyone. She's looking for any sign of her missing father.
\item Sebastian Crowe, a half-elf \textbf{mage}. This meddlesome and overconfident Academy dropout is haunted by nightmares of the meteor.
\end{itemize}

\subsubsection{Intrepid Explorers}
These longtime friends have been adventuring together for years. They include a \textbf{berserker}, a \textbf{hedge mage}, a \textbf{veteran}, a \textbf{priest}, and a \textbf{spy}.

\subsubsection{Foolhardy Rookies}
These new arrivals might not last a day in Drakkenheim, but they've got spunk! The group consists of an \textbf{acolyte}, a \textbf{guard}, a \textbf{noble}, and a \textbf{scout}.

\section{Starting the Campaign}
You can begin characters at 1st or 3rd level. We strongly recommend holding a Session Zero before starting the campaign. During this special session, make sure you achieve the following:


\begin{itemize}
\item Generate the player characters. Present players with the information included in the introduction to this book, especially the "Character Creation" section.
\item Share the appendices "The World of Drakkenheim" and "Character Backgrounds" with players who want to create characters with well-established connections to the setting.
\item Allow each player to choose a personal quest by sharing the table or cards with them. Remember, they need not disclose these to each other.
\item Ask the players to decide how their characters met and why they're working together. You can use the "Road to Drakkenheim" adventure below to bring the characters to the city itself.
\item Distribute one rumour to each player character using the tables or cards.
\item Share the Poster Map of Drakkenheim with the players.
\end{itemize}

In addition, Session Zero is an important opportunity to establish a few expectations with your players about the campaign. We've created a video guide to Session Zero on our YouTube channel! Beyond scheduling and table etiquette, you should cover a few topics specific to this campaign:


\begin{itemize}
\item Discuss the themes of dark fantasy with your players by reviewing the content warning provided in the introduction of this book.
\item Highlight the warnings in the "Dangers of the Dark City" section in the introduction. You may wish to present the full rules for contamination to your players from Appendix C.
\item Players might react to the rules for contamination and the Haze by asking questions about whether or not their character abilities can protect them from these effects. Even if they propose something inventive, your answer should be a firm "no". Instead, explain they may discover new ways to overcome these problems during the adventure.
\item Remind your players the adventure occurs entirely within the bounds of Drakkenheim and its immediate environs. While the wider world gives context to the story, travelling outside Drakkenheim is beyond the scope of the campaign. Make a friendly agreement with your players that no matter what happens, their characters won't leave the area. If conditions emerge where it would make sense for a player character to depart the city, the player should play a new character until their previous one returns to Drakkenheim.
\end{itemize}

After you've finished your Session Zero, begin your first game session with the introductory adventure, "The Road to Drakkenheim".

\subsection{Alternate Opening: Faction Contact}
If you're beginning at 3rd level, you may instead begin the campaign with the player characters in direct contact with a faction lieutenant, who has hired them for a job in the Outer City described in chapter 6. We began the original \textit{Dungeons of Drakkenheim} campaign with characters meeting River of the Amethyst Academy at Eckerman Mill, who hired them to recover a delerium crystal from the ruins of the Rat's Nest Tavern.

\section{The Road to Drakkenheim}
\DndQuote%
{}
{''You couldn't pay me any amount of money to go into that city, but there are people willing to risk it for the right price. Those people gotta eat, and honest folk can make good profit feeding them. Still, there's bastards who just wanna take it. That's where you come in.''}
{Eren Marlowe, Provisioner}
If you wish to start with characters at 1st level and ease them into the campaign as a whole, begin with this short prelude. In this introduction, the characters are making their way towards Drakkenheim for the first time, accompanying a supply caravan bound for Emberwood Village. You can also use this opening to start characters at 3rd level, but you may wish to raise the difficulty of these encounters by adding additional creatures.

\subsection{Adventure Hook}
\subparagraph{Supplies for Emberwood Village}
Though each has their own reasons for venturing into Drakkenheim, the characters have been hired by Eren Marlowe (a human \textbf{commoner}) to guard their caravan after the last group they hired stiffed them at the last minute. They'll pay each character 25 gold for the trip upon arrival in Emberwood Village. It's the perfect way to make some money as characters trek towards the city. The characters meet Eren Marlowe in the city of Altberg, which is the closest large city to Drakkenheim. It's a three week journey to the ruins.

\subsection{Eren Marlowe}
Eren Marlowe is a lean person in their mid-forties with shoulder-length brown hair under a leather cap. They wear a slightly shabby woolen traveller's jacket and well-worn slacks. Chatty yet professional and with an opinion about everything, they enjoy having company on the road as much as they enjoy having the extra protection. As the characters travel along the road, Eren Marlowe asks each of them about their reasons for travelling to Drakkenheim, and how the characters met each other. This is a perfect opportunity to introduce the player characters, their relationships, and personal quests. In turn, Marlowe can share the following information with characters:


\begin{itemize}
\item They are bound for Emberwood Village, which is the only settlement within two weeks' travel of Drakkenheim that hasn't been abandoned. It is about five miles south of the dark city.
\item The increasing number of fortune-seekers, mercenaries, and prospectors exploring the ruins has turned the sleepy village into a boom town again.
\item Most farms and villages within one hundred miles of Drakkenheim have been abandoned. Crops and livestock don't grow anymore as the water and soil are contaminated by foul magic.
\item The folk of Emberwood Village rely on outside supplies entirely. Flamekeeper Hanna, a priest of the Sacred Flame, provides seed funds to attract the caravanners.
\item Once they arrive, Eren Marlowe will stay in Emberwood Village for a few weeks selling goods.
\item If characters are interested in heading into Drakkenheim, Marlowe offers to introduce them to their friend and frequent business partner, Armin Gainsbury, who runs a supply company in Emberwood Village. Armin sells useful adventuring equipment and can provide good advice about exploring the ruins.
\end{itemize}

\subsection{Events on the Road}
While most of the journey is uneventful, once characters are a few days from Drakkenheim, read:

\begin{DndReadAloud}{}
\textit{After two weeks travelling on the road, the countryside turns dreary and desaturated. You pass many abandoned villages with dilapidated homes, crumbling farmsteads with sallow empty fields, and forests of dead trees. Thorny brambles and harsh scrubgrass are all that grow along the mud-slick road. Drizzling rain fills most days, and the wind howls at night. Aside from the creaking wagon wheels, you only hear the cawing of crows}.

\end{DndReadAloud}

\subsubsection{Marlowe's Caravan}
Eren sits atop a large covered wagon pulled by two draft horses. It's packed with common provisions and trade goods: flour, salt, dried fruits, preserved meats, several jars of sauerkraut, a few casks of mead, a case of cheap wine, and a barrel of fresh water. The wagon also packs a few sets of common artisan's tools, \textit{cook's utensils}, empty waterskins, Climber's Kit, rope, and mining tools. Marlowe carries two Potion of Healing on their person for emergencies, and an extra \textit{light crossbow}. Beneath a hidden plank in the wagon is a locked chest that contains 200 gold.

\subsubsection{Bandit Shakedown}
\begin{DndReadAloud}{}
\textit{"Oi, oi! Hold there, friends... you're travelling on a royal road, eh? In these troubled times it falls to loyal servants like ourselves to keep these highways free of reprobates and outlaws, eh? I'm afraid we'll need to collect a fee to ensure your safe passage to the capital, eh?"}

\end{DndReadAloud}

During one afternoon on the road, several cutthroats emerge from the ditches and rubble of a nearby farmhouse with crossbows at the ready.

Bogdan, a \textbf{thug} mounted on a \textbf{riding horse}, leads four bandits: Hosh, Mildred, Sanna, and Mudface. They demand the horses, all the booze in the wagon, a pound of food each, and a payment of 25 gold per character (including Eren Marlowe).

\subparagraph{Negotiation}
The bandits can be talked down with three successful Charisma (Deception, Intimidation, or Persuasion) checks (in any combination). It's DC 10 if the characters offer to pay the bandits something, but DC 15 if they try to convince the bandits to walk away empty-handed but alive. However, after three failed checks, the bandits are convinced the characters are full of hot air and won't entertain any further negotiations.

\subparagraph{Combat}
If it comes down to a fight, the bandits try to gang up on the most vulnerable target. Then, Bogdan threatens to finish off an unconscious character unless his demands are met. Otherwise, the brigands disengage when wounded. They flee if one-third of their number are killed and none of the player characters are downed.

\subparagraph{Developments}
These bandits are only loosely associated with the Queen's Men and don't venture into Drakkenheim themselves. If captured or questioned, they don't know anything about the Queen's Men or Drakkenheim that isn't common knowledge. They have a hideout in an old barn a few miles off the road. If the characters follow up, the barn is guarded by six bandits. It contains their spoils: 75 gold in trade goods and coins and a \textit{potion of healing}.

\DndQuote%
{}
{''Sneaking and skulking is the preferred method, second best is having someone like me around.''}
{Pluto Jackson}
\subsubsection{A Night on the Road}
A trio of harrowed adventurers approach the characters' campsite on the last night of their journey. They are tired, hungry, and hoping for a place to rest. They are leaving Drakkenheim after failing their quests and losing two of their other companions. They are bitter, defeated, cynical, and never plan to return, but share one or more rumours about the city as they sit and share stories by the fire. They are hesitant to talk about their fallen friends, but suggest watching out for the rats in Drakkenheim. They include:


\begin{itemize}
\item Ludwig von Graff, a penniless human \textbf{noble} who was trying to find the deed to his ancestral estate.
\item Endra Jansen, a human \textbf{scout} who was seeking fortune to pay off an old debt.
\item Rikard Vos, a human \textbf{veteran} who survived the civil war and had returned to Drakkenheim searching for lost family members.
\end{itemize}

Vos has accumulated five levels of \textbf{contamination} and is now incapacitated. He follows his companions in a daze, and cannot speak as his tongue has rotted out from a \textbf{rasping mutation}. His friends have covered him in a heavy cloak to conceal his \textbf{wasting mutation} and the \textbf{tentacled limb} growing in place of his left arm.

If the characters detect Vos' mutations or pry the adventurers about his illnesses, they confess that he became sick after picking up one of the strange crystals in Drakkenheim. They insist the illness is not contagious, but rather caused by overexposure to the contaminated magic in Drakkenheim. They've heard expensive magical medicine might be able to help, but they can't afford it. Instead, they're hoping he'll get better once they bring him far away from the city.

Unfortunately, during the night, Rikard Vos succumbs to contamination and transforms into a tentacled delerium dreg. The monster flies into a rage and attacks one of its former companions. It's up to the characters to decide how to intervene. The creature is beyond reason and attacks unceasingly until killed. If von Graff and Jansen survive, they are deeply troubled by this incident. They solemnly warn the characters about the horrors of the city as they set off in the morning.

Marlowe has heard stories about contamination and delerium, but never witnessed it first hand. They note that handling delerium should be done with extreme care, and that far more terrible creatures are said to dwell in Drakkenheim, many the result of those who suffered a similar fate as Rikard Vos.

\subsection{Rumours}
Rumours offer characters clues about Drakkenheim and prompts which encourage them to freely explore the ruins to see what they can find. A few rumours are true, some of them are somewhat misleading, and others are entirely false. Regardless, there is always value in looking into a rumour, as each will lead to great adventures! Here are a few suggestions for when characters might find a rumour:


\begin{itemize}
\item During downtime in Emberwood Village
\item Written on the walls of Eckerman Mill
\item Clues found in a random encounter
\item Interacting with rival adventurers
\end{itemize}

You can choose one of the rumours or roll one randomly on the table below:


\begin{DndTable}{cX}

d100 & Rumour\\
1-4 & Don't go to Slaughterstone Square.\\
5-8 & The Trolls of King's Gate will let you pass for a price, but it will cost you an arm and a leg.\\
9-12 & Queen's Park Garden is filled with flowers unlike any you've ever seen, they glow with a strange colour to them and some people said they have immense healing properties.\\
13-16 & Someone has been selling potions out of the old Reed Manor.\\
17-20 & The Black Ivory Inn is up and running, and it sounds like business is booming too!\\
21-24 & There is an invisible wall surrounding the old Rose Theatre.\\
25-28 & The Queen's Men run vicious fighting pits beneath Buckledown Row, which lead to their secret stronghold in the sewers. The \textbf{Queen of Thieves} will grant any request within her power to whomever can survive a bout with her monstrous champion, but none so far have claimed the prize.\\
29-32 & It's said the many gargoyles perched atop the City Walls come to life and attack any who dare attempt to scale the walls.\\
33-36 & The strange fishlike folk in the sewers are quite friendly, and have saved many desperate explorers.\\
37-40 & The Followers of the Falling Fire have gathered at Saint Selina's Monastery. It's deep within the Haze, so how do they manage to survive there?\\
41-44 & Elias Drexel was a mighty general during the civil war, leading the Hooded Lanterns isn't even the first time he's tried to retake Drakkenheim.\\
45-48 & The Archmage of Drakkenheim died when the meteor fell, and the Amethyst Academy has yet to appoint a true successor.\\
49-52 & The Knights of the Silver Order have camped outside the city, I hear they have weapons that can torch a village in moments.\\
53-56 & The Noble Man, in his noble home. Withered, old, and all alone. His stare can chill you to the bone, a \textbf{pale man} on a pale throne.\\
57-60 & I heard there are delerium crystals the size of a horse within the crater. Imagine how much they'd be worth!\\
61-64 & Deep beneath Drakkenheim are the remains of an ancient dragon. An old legend claims they may be resurrected in a time of need.\\
65-68 & The \textit{Crown of Westemär} has the power to grant the monarch's every wish.\\
69-72 & A shrine to the Old Gods lies north of the castle, tended by an elfish druid with power over life and death.\\
73-76 & A dwarven clan from the Glimmer Mountains has laid claim to a rich field of delerium south of the Crater.\\
77-80 & Exploring the inner city is dangerous, and it's not just because of the monsters. Anywhere big chunks of the meteor landed are covered in the Deep Haze—the tower, the cathedral, the castle, or of course, the crater. It's a deathly fog that saps the body and clouds the mind... and that's only the beginning. What happens after—it's far worse.\\
81-84 & Queen Lenore's favourite place to spend time was the grotto in the garden, where she housed many exotic flowers.\\
85-88 & The inner city has become the hunting ground for a fearsome monstrosity known only as the \textbf{Lord of The Feast}. All who dare enter become his prey.\\
89-92 & A flock of harpies dwells in the Cosmological Clocktower. They viciously attack any who trespass on their territory, but they are deathly afraid of cats.\\
93-96 & Whispers say the ratlings serve a demonic god.\\
97-00 & Many of the monsters lurking in Drakkenheim are in fact its former inhabitants, twisted by fell magic. They can be made docile by showing them a token of their former life.\\
\end{DndTable}

\section{Conclusion}
Marlowe has no intention of turning back, and insists the caravan press on towards Emberwood Village (described in chapter 4). In town, they will meet agents of the five factions, detailed in the next chapter.

\chapter{Factions}
\DndDropCapLine{T}{}his chapter describes the roleplaying traits, resources, and major NPCs for each of the five factions. Each dossier is broken down into the following sections:


\includegraphics[width=0.4\textwidth]{images/../5e.tools/img/adventure/DoDk/006-02-001.large-argument.png}
\section{Background}
This details a brief history of each faction and its recent activities in Drakkenheim.

\subsection{Objectives}
These describe the faction's major goals in Drakkenheim. Most concern two core conflicts that define the campaign: \textit{Who should control Drakkenheim? What should be done about delerium}? Conflict is inevitable because every faction desires very different outcomes with regard to these questions. All wish to control the ruins and decide how delerium will be used (or destroyed). In addition, some factions have additional objectives which may directly oppose those of a rival faction.

Each faction's \textbf{Roleplaying Traits} resonate strongly with their \textbf{Objectives}. However, a faction may tolerate setting aside or even adjusting its objectives if it perceives a chance for greater gain in the future; if acting on its objectives would incite overwhelming retribution from political forces beyond Drakkenheim; or if circumstances emerge where achieving its objectives would contradict the faction's own ideals. In such a situation, a faction might even modify its objectives. Player characters might create alliances or expose weaknesses in the factions by taking advantage of this fact.

\subsection{Key Figures}
Each faction has two major NPCs: the faction leader and their lieutenant(s). The stat blocks for these characters are detailed in Appendix B, and their roleplay traits are provided here. Use these NPCs as the player characters' primary points-of-contact during faction interactions such as:


\begin{itemize}
\item Briefing or hiring the characters for a mission,
\item Meetings and diplomatic negotiations,
\item Sharing information, advice, and creating mutual plans,
\item Offering rewards or resources to the player characters on behalf of the faction.
\end{itemize}

\subsubsection{Faction Leader}
A formidable legendary non-player character commands each faction. These faction leaders issue orders to faction agents, direct faction resources, and make the final decision on how the faction acts in all matters. Each has their own individual Personality Traits, Ideals, Bonds, and Flaws, and perhaps a few dark secrets too! As conflict in Drakkenheim escalates, the player characters should interact with the faction leaders often, and might even bring several faction leaders together for negotiations.

The faction leaders spend most of their time coordinating strategy and logistics within their faction stronghold. They are the organizational lynchpin for their entire faction, and rarely take to the field. A faction leader should only personally venture into Drakkenheim for a major gathering or offensive involving the cathedral, Castle Drakken, or a faction stronghold. The risk to the faction is simply too great to risk their lives in the field otherwise.

If a faction leader is slain or permanently incapacitated, the faction must withdraw from Drakkenheim shortly thereafter unless there is someone else to take up their mantle. Surviving faction agents may seek immediate retribution, but the death of a faction leader seals their defeat in Drakkenheim for the foreseeable future.

Under the right circumstances, a player character who demonstrates unwavering belief in the faction and great leadership could rise through the ranks and replace a faction leader.

\subsubsection{Faction Lieutenant}
The faction lieutenant is the right-hand of the faction leader. As such, the lieutenant is usually the first point of contact the player characters have with a faction, and most can be encountered in Emberwood Village. Once the lieutenant is confident the players are capable and trustworthy collaborators, they'll facilitate an introduction with the faction leader.

Many lieutenants also lead elite faction strike teams. If the lieutenant is slain, the faction leader will appoint another faction member as a replacement.

\subsubsection{Supporting Characters}
Each faction has numerous rank-and-file members, support staff, craftspeople, recruits, and other followers. Use these lists to generate quick NPCs who might be guarding a stronghold, relaying a message, maintaining the supply lines, or fighting alongside (or against!) the player characters. You can add additional NPCs to the factions if you wish, but be careful when adding high-level NPCs or spellcasters capable of casting spells above 3rd level. The factions have specific and limited access to high-level magic. Powerful NPCs are meant to be rare so the player characters may rise to a position of importance and prominence during the campaign.

\subsection{Roleplaying Characteristics}
Much like a player character, every faction has Personality Traits, Ideals, Bonds, and Flaws. Rather than qualities of any individual person, these describe ideological beliefs shared by all faction members. Use these characteristics to guide roleplaying interactions with faction agents, determine the way the faction behaves, and judge how their leaders might react to changing circumstances or unexpected events.

Faction NPCs often possess these characteristics, but may have their own personal interpretations. However, a faction that betrays its ideals or bonds has betrayed its own members. When a faction repeatedly fails to uphold its own ideals and bonds, its agents might start deserting Drakkenheim or refuse to comply when called to action.

\subsection{Strongholds}
Most factions have several holdings in Drakkenheim, but one major location is the faction stronghold. Full details on each faction stronghold are described in chapter 8. This is where most faction members reside and where many faction resources are kept. If a faction loses its stronghold, the rank-and-file members are routed and the power of the faction is broken. If the faction leader survives, they may still be able to muster together a few remaining diehard agents to cause a few more complications for the player characters, but the faction teeters on the brink of defeat.

\DndQuote%
{}
{''As useful as magicians and wizards can be when organized, You should always have a bag of tricks...for sale of course.''}
{Pluto Jackson}
\subsection{Missions}
The faction leader or faction lieutenant may present these missions to the player characters, which further the faction's objectives. The missions are presented in a rough progression: those presented first may lead to later ones. Characters who complete these missions gain greater respect and trust from the faction, and are rewarded with faction boons (described below). In addition, when the player characters complete a faction mission they receive experience points based on the advancement method you are using (as described in Character Advancement in chapter 2).

Many missions are directly tied to the \textbf{City Locations} described in chapters 6 and 7. In such cases, further details for running that mission are found in the "Adventure Hooks" section of the location entry. Other missions are open-ended adventure seeds to spur interactions with other factions. Nevertheless, you should feel free to invent additional missions for the factions based on the player characters' choices. Especially when the factions clash against each other, the player characters themselves collaborate with their faction allies to devise their own missions and specific goals.

\subsection{Boons}
Every faction possesses valuable resources and knowledge they use to advance their agenda, including equipment, information, magic items, mounts, spellcasting services, and more. These are represented by faction boons, which may be awarded to the player characters for completing faction missions, but also offered to player characters \textit{before} a mission to help them complete it.

These boons are presented in a roughly escalating order: latter boons should only be offered later in the campaign to characters who have demonstrated their commitment to the faction. Unless otherwise indicated, a faction can only offer each boon once. Beyond this, how often you reward characters with a faction boon is up to you.

\subsubsection{Other Resources}
The factions' exact assets and wealth are intentionally abstract so you may exercise your judgement and decide what's appropriate given the circumstances. Regardless, the factions do not have infinite resources. Every faction is stretched thin, and must cope with severe limits and constraints on their finances, personnel, materials, equipment, and time. The nearest large city is several weeks' travel away and the surrounding villages have been abandoned. More troops require exponentially more supplies since there are so few uncontaminated food and water sources. For many of the factions, simply maintaining the ground they currently hold requires everything they've got. \textbf{This is why the factions need the player characters to help them.}

The characters may occasionally demand more from the factions; citing dire need or urgent circumstances. If the characters push hard for more resources and appeal to the factions' ideals, objectives, and bonds, you might grant them an additional boon for that mission. However, the next time the characters need help from the faction or would receive a boon, the faction can't offer them anything as they are recuperating their stockpile!

\subsection{Strike Teams}
While the factions seek assistance from the player characters, they aren't without their own agents. When they must dispatch their own warriors directly into the ruins, they send a \textbf{Strike Team}.

\includegraphics[width=0.4\textwidth]{images/../5e.tools/img/adventure/DoDk/007-02-002.queens-men-thug.png}
An allied faction might send strike teams as reinforcements to help the player characters on a critical mission. Consequently, an enemy faction can dispatch a strike team against the player characters to disrupt missions, capture them, or even assassinate them. Each faction keeps tabs on the player characters, and might even send a strike team to interfere during a mission that wouldn't otherwise involve them just to sabotage the characters' efforts. More ruthless factions might take the fight directly to the characters' doorsteps, hitting them while they are retreating, recuperating, or simply unprepared for battle. Especially devious GMs may opt to have faction strike teams operate behind-the-scenes by carrying out attacks while the player characters are occupied elsewhere!

\subparagraph{Mustering Forces}
Any faction can call up a strike team within a few hour's notice. However, mustering up larger forces for major offensives takes time. When a faction needs to assemble multiple teams at once to perform a specific task, it takes at least one day per strike team involved. This can give characters the chance to notice something big is being planned.

\subparagraph{Reinforcements}
The factions regularly bring in reinforcements and new recruits. However, increasing casualties amongst the rank-and-file devastate morale within all the factions, and it's a problem when its more elite strike teams are captured, defeated, or killed. When a faction strike team is destroyed, it takes roughly five days for the faction to assemble a replacement team. Under desperate circumstances, a faction can put together another group immediately. Should this group also be taken out, it takes ten days to draw up further replacements thereafter.

\subparagraph{Retribution}
Faction leaders often investigate when their agents go missing. Each faction possesses a web of informants, scouts, and spies, and can possibly use magical means to find out what happened. Unless an active effort is made to conceal what occurred, most factions can discover the fate of their agents within seven days and possibly even determine who was responsible. Skirmishes between adventurers and low-level faction agents are relatively common in the ruins, so such events may not incite immediate retribution, but rather careful observation. If the situation continues to escalate, the faction responds.

\subparagraph{Major Offensives}
Every faction knows that sending large numbers of soldiers into Drakkenheim is a foolish risk: several attempts have been made by much larger forces and failed. The biggest issue is contamination: there simply aren't enough spellcasters and materials to remove it quickly enough. Generally speaking, the factions will only consider launching major offensives if it means taking the cathedral, Castle Drakken, or another faction stronghold.

\subsection{Schemes}
Should your player characters become enemies with one of the factions, the factions enact these dastardly plots against them. These schemes are presented in escalating hostility; so latter schemes are only used in dire conflict. Set one of these schemes into motion the next time the player characters venture into the ruins. If the players manage to foil the faction's schemes, the defeated faction will need at least a week before they can muster up another. On the other hand, the factions don't do things "just to get revenge". They instead consider their long-term goals first. They retreat, regroup, and find a way to attack from a position of strength.

You can invent your own schemes. Try using a strike team in a creative or unexpected way, or consider the personal capabilities of each faction leader: especially the spells, magic items, and unique knowledge at their disposal. Don't be afraid to turn the player characters' own motivations, goals, and flaws against them either!

\subsection{Faction Reputation}
\subparagraph{Champions}
The player characters have demonstrated themselves to be staunch supporters of the faction's cause. They've successfully completed multiple missions for the faction that prove their loyalty, including a mission where they acted directly against another faction. The faction offers the characters complete access to their resources, support, and information. The characters receive faction boons to help them complete missions.

\subparagraph{Allies}
The party and the faction are working together to advance their mutual goals, but they have some minor differences and disagreements. Characters might be working with multiple factions in this manner. The characters have helped the faction in several ways by completing a few missions, and proven themselves reliable and capable. The faction rewards characters with boons after they complete these missions.

\subparagraph{Indifferent}
They are neutral. Perhaps they aren't known to each other yet. The faction is willing to hire the party on a freelance basis to help fulfil their objectives—especially ones that require arms' length interaction (such as working against another faction), but the faction doesn't trust the players yet. They may have had some rough interactions with low-level faction members, so the faction won't offer the players resources, support, or information beyond what is required to fulfil any missions they hire them to do. The faction might offer to trade for services and pay the characters in gold rather than grant them boons.

\subparagraph{Adversaries}
There isn't any trust here. The players show disregard for the faction's objectives. Minor skirmishes may have broken out with low-level faction members. The faction won't let the players use their resources or help. If the players want to repair the relationship, the faction offers the players a chance to complete a mission to demonstrate their commitment, possibly by striking directly at a rival faction. The faction occasionally sends strike teams to impede the party, disrupt their goals, and dissuade them from entering the city. However, these teams don't try to kill the characters. If violence breaks out, these strike teams either knock the characters out to try to capture them, deescalate, or flee.

\subparagraph{Enemies}
The characters have betrayed the faction's trust, murdered its members, and are working against the faction's most important goals in a profound way. The faction actively tries to thwart the characters using its schemes, and sends its deadly strike teams against the the characters.

\includegraphics[width=0.4\textwidth]{images/../5e.tools/img/adventure/DoDk/008-02-003.bounty-hunt.png}
\section{The Amethyst Academy}
\DndQuote%
{}
{''Suffice to say, the Amethyst Academy has developed the means to harness the incredible arcane energies contained within delerium. We shall use it to create unfathomable magical wonders on a scale hitherto undreamt of, for the betterment and protection of all people. Only the Amethyst Academy has the knowledge to accomplish such a feat, and the wisdom to temper the wild and destructive potential of this volatile resource.''}
{Archwizard Eldrick Runeweaver}
\includegraphics[width=0.4\textwidth]{images/../5e.tools/img/adventure/DoDk/009-02-004.amethyst.academy.png}
The Amethyst Academy is a research institute, educational organization, and mage guild pursuing the arcane mysteries of delerium. The faction is small and outwardly aloof, yet subtly influential and tremendously wealthy. They are led by a shadowy directorate of powerful archmages who dispatched Archwizard \textbf{Eldrick Runeweaver} to oversee their operations in Drakkenheim. Most of its members are potent arcane spellcasters with access to rare magical resources, but only a handful of agents are active in Drakkenheim.

\subsection{Background}
The Amethyst Academy began as a secret school for young sorcerers during a time when arcane magic faced extreme persecution from the clergy of the Sacred Flame. The mages located children who manifested magical talents, rescued them from their fearful and paranoid communities, and educated them in hidden strongholds and far-away castles so they might wield their abilities with trained and refined control. The secret school gradually grew into a secret society in its own right.

Sequestered in unseen universities to hide from the mage-slaying Knights of the Silver Order, the mages pooled their knowledge. They conducted research into arcane and cosmological phenomena, and developed the modern practice of wizardry. Over time, the Amethyst Academy secured allies amongst the nobility by providing them clandestine magical services and occult counsel, crafting enchanted heirlooms, protecting their castles with arcane wards, and teaching their heirs magic.

\includegraphics[width=0.4\textwidth]{images/../5e.tools/img/adventure/DoDk/010-02-005.eldrick-runeweaver.png}
When these schemes were discovered by the Faith of the Sacred Flame, a renewed series of brutal inquisitions and vengeful witch-hunts began. What had been a shadow conflict between paladins and sorcerers for centuries spilled into outright war between noble houses who championed the Faith, and those that supported the Amethyst Academy. These conflicts ended with the Edicts of Lumen (detailed in Appendix E). Enacted three hundred years ago, this landmark treaty stopped the murderous suppression of arcane magic by the religious ministry, and established a new balance of power between mages, the Faith, and the nobility throughout the continent.

The Edicts of Lumen bestow extensive protections, full autonomy, and economic advantages to the Amethyst Academy, but also require the Academy to observe strict standards of mercantile neutrality and political non-intervention. Most importantly, it grants the Amethyst Academy guardianship of any child who demonstrates magical aptitude until they come of age and master their abilities. In exchange, the agreement stipulates that all arcane spellcasters are disinherited from any land, rank, or titles, including the scions of noble and royal families. Since the Edicts, Academy mages may work magic without fear of oppression from church or state, though the organization remains withdrawn as it increasingly turns its attention to global affairs and extraplanar happenings. However, the direction and focus of the Amethyst Academy has dramatically shifted following the destruction of Drakkenheim, when their greatest stronghold was cut off from them.

While the Amethyst Academy outwardly upholds the Edicts of Lumen, in recent years some have claimed the mages furtively flaunt its restrictions and disregard the terms that forbid dark magical practices. Others suspect the mages are weaving an arcane conspiracy to control commerce and influence politics while shielding themselves from the wrath of the Sacred Flame.

\subsection{Attire}
Agents typically wear purple robes, cloaks, or hoods.

\subsubsection{Academy Rings}
Members of the Amethyst Academy each wear a set of rings made from exotic arcane metals. The number of rings worn corresponds to the highest level spell they can cast. Known as \textit{academy rings}, these are magic items that identify them as an Academy graduate, and often permit access to Academy strongholds. A set of five or more functions as a \textit{ring of spell storing}, and some allow control over a \textbf{shield guardian}. However, each set is made for a specific individual, and only that person may attune that set of academy rings.

\subsection{Strongholds}
The Amethyst Academy does not currently have a dedicated base of operations in Drakkenheim. Indeed, one of their most important objectives is reclaiming the Inscrutable Tower, so it can be used for such purposes!

\subparagraph{Emberwood Accommodations}
For the time being, most Academy members simply make use of local accommodations: Eldrick and River have booked a permanent suite in the Red Lion Hotel. However, the Amethyst Academy has secretly created a permanent \textit{teleportation circle} in the cellar beneath the abandoned Eventide Manor in Emberwood Village. Academy agents come and go from the cellar using \textit{invisibility} spells.

\subparagraph{Distant Fortresses}
Far away from Drakkenheim, the Amethyst Academy has many great fortresses, libraries, and universities. Most notable are the Enigma Ziggurat in the Free City of Liberio, and Paradox Castle in the distant Eastern Vale. Both have well-guarded \textit{Teleportation Circle}, and Academy agents working in Drakkenheim usually teleport back and forth between these strongholds and the secret circle in Emberwood Village. The sigil sequences of these teleportation circles is a closely guarded secret known only to Academy mages who can cast 5th level spells.

\section{Roleplaying Characteristics}
\subparagraph{Symbol}
A lidless arcane eye set within an octagram. Their colours are gold and purple.

\subsection{Personality Traits}

\begin{itemize}
\item We act businesslike and professional when dealing with the wider world.
\item We put arcane theory into magical practice. Innovation, creativity, and invention must have purpose.
\item We do not trust outsiders, and don't volunteer information that is not requested or required. Non-magical folk rarely grasp arcane matters anyways.
\end{itemize}

\subsection{Ideals}

\begin{itemize}
\item We rigorously study and carefully teach arcane magic. We value eldritch knowledge, through which we may harness the power of the cosmos itself.
\item Mages can accomplish incredible wonders by working together effectively. We are protected and strengthened by our disciplined hierarchy.
\item There is no such thing as evil magic, but some mages misuse magic for evil purposes. Unfortunately, those who do not practice the arcane arts rarely grasp such nuance.
\end{itemize}

\subsection{Bonds}

\begin{itemize}
\item The Amethyst Academy safeguards occult secrets that would be abused and feared by the wider world. What others view as dangerous, we seek to understand.
\item We cannot allow the brutal suppression of arcane magic to ever occur again, so we respect the Edicts of Lumen. We must maintain the trust of the nobility and avoid arousing suspicion and hostility from the Faith of the Sacred Flame.
\item We use wealth accumulated by selling magical services to educate young mages and support further arcane research. Our students are our future, and they must be kept safe.
\end{itemize}

\subsection{Flaws}

\begin{itemize}
\item Since we are so often misunderstood, we act in unseen and secret ways.
\item We are researchers and scholars, not warriors. Risking our talented members is wasteful and reckless when their minds can be put to much better use. Other lives are disposable.
\item We believe only Academy mages possess the intelligence and wisdom to wield arcane magic in a responsible way. Magic users who aren't Academy members are undisciplined amateurs at best, and more often become unhinged villains.
\end{itemize}

\DndQuote%
{}
{''The Academy proves to be mysterious and secretive at every turn, but I know not to ask too many questions if I want to bring home the bacon. I'd be a fool to bite the hand that feeds me.''}
{Veo Sjena}
\subsection{Objectives}
\subsubsection{Collect and Study Delerium}
Far away from Drakkenheim in secret laboratories, Amethyst Academy wizards relentlessly research delerium samples recovered from the ruins. Their brilliant minds have crafted potent arcane artefacts and invented mighty new spells powered by the vast energies contained within the crystals. Though the true origin of delerium remains unknown to them, many Academy assets are now dedicated to the study and refinement of delerium. They procure enormous quantities to fuel their esoteric experiments and eldritch industries. Rather than risk their own members' lives in the ruins of Drakkenheim, the Academy usually obtains delerium through various grey market sources, typically by hiring adventurers.

The mages do not disregard the dangers presented by delerium's contaminating properties. In fact, it is precisely because of these hazards that the mages believe only they can use the crystals responsibly. Thus, the Academy seeks complete control over this critical resource both to protect the wider world and to secure their magical hegemony for the foreseeable future. However, other factions are now moving to reclaim Drakkenheim and potentially install a new ruling dynasty. Such political upheavals could affect the Academy's ability to research and harvest delerium. The Academy does not want a ruler who will destroy delerium, force the Academy to pay exorbitant sums to obtain it, or impose conditions on how the Academy may use the crystals.

Ultimately, the Academy has zero interest in seeing Drakkenheim restored and rebuilt, and view this as a folly doomed to fail. Rather, they have an ambitious plan to transform the city ruins into a quarry for their research and industry. Once they fulfil their other goals, the Amethyst Academy wants the other factions to totally abandon Drakkenheim so they can have sole control over the ruins.

\includegraphics[width=0.4\textwidth]{images/../5e.tools/img/adventure/DoDk/011-02-006.inscrutable-staff.png}
\subsubsection{Reclaim the Inscrutable Tower}
Drakkenheim was home to one of the Amethyst Academy's most important strongholds: the Inscrutable Tower. This immense structure houses vast repositories of magical lore, incomplete magical experiments, hidden vaults containing unspeakable extraplanar entities, and powerful magic items sheltered from the world at large—all rendered completely inaccessible to the Academy since the meteor fell. The Deep Haze has settled within the tower, scrambling the Academy's own arcane wards and foiling every attempt to enter or observe the Tower using teleportation or divination magic.

The mages have already expended considerable effort attempting to retake the Tower on their own. These attempts catastrophically failed. The pragmatic wizards cannot afford to exhaust the Academy's resources in such a fruitless manner any longer, so now they turn to research and patient methods to advance their goals in Drakkenheim. They seek a suitably talented, trustworthy, and expendable adventuring party to conduct survey and recovery missions into the Inscrutable Tower. If the Tower were to be secured and repaired, the Amethyst Academy would have the ideal staging ground to begin the mass industrial exploitation of delerium, and would be able to bring its full resources to bear in Drakkenheim.

\subsubsection{Obtain the Inscrutable Staff}
Many talented mages were lost or killed by the meteor strike—including the former Archmage of Drakkenheim, Adriana Modiera. Her position on the Archmage's Directorate has remained vacant ever since. The Academy desperately wishes to discover her ultimate fate, recover her research, and obtain the \textit{Inscrutable Staff}, the symbol of office for the Archmage of Drakkenheim, which is also one of the \textit{Seals of Drakkenheim}. Traditionally, the Archmage of Drakkenheim acts as an advisor on the royal council, so possessing this \textit{Inscrutable Staff} ensures the Amethyst Academy maintains their influence in the political destiny of Westemär.

\subsection{Key Characters}
Archwizard \textbf{Eldrick Runeweaver} heads up the Amethyst Academy operations in Drakkenheim and serves as the faction leader. The tiefling mage River works alongside him as lieutenant.

Eldrick Runeweaver serves as the acting Archmage of Drakkenheim within the Academy Directorate, but this is a provisional position. He will not be confirmed to a lifetime appointment until he successfully secures the Inscrutable Tower. He is accountable to the Academy Directorate, but only reveals his true position to close collaborators. Although he does not lead the Amethyst Academy as a whole, his death or defeat would cause the Directorate to seriously reevaluate their position in Drakkenheim.

\subsubsection{Eldrick Runeweaver}
One of the finest professors in the Amethyst Academy, \textbf{Eldrick Runeweaver} has personally instructed a generation of talented arcane spellcasters. He is the foremost living expert on abjuration magic, but is also highly regarded within the Academy as a dedicated administrator and calculating leader. If Eldrick Runeweaver successfully retakes the Inscrutable Tower, he will likely take up the vacant spot in the Academy Directorate and assume the mantle of Archmage of Drakkenheim.

Eldrick Runeweaver is a towering and broad-shouldered man in his mid-sixties. His amber eyes glow softly behind his circular glasses, and he has a square-cut, greying black beard. He wears nine rings of silver, gold, and electrum: one on each finger except his left thumb. He wears flowing purple garb with golden embroidery, and carries an obsidian staff tipped with a delerium crystal.


\begin{description}
\item[Personality Trait.]
I carefully calculate my every action, and I speak using a few well-chosen words. I'm a teacher at heart, so I often phrase my responses in the form of a pointed question.
\item[Ideal.]
I will seek out more information whenever I encounter an unknown situation or problem. Through reason and logic, the facts always bear out the truth of the matter.
\item[Bond.]
I hold a position of responsibility and authority within the Amethyst Academy, and it is my duty to protect the Academy's interests, carry out its agenda, and keep its many secrets.
\item[Flaw.]
I care deeply for my students, but tend to overlook their faults and failures. I've had several wayward apprentices over the years, and I've often had to cover up their mistakes.
\end{description}

\textbf{Eldrick Runeweaver} has a \textbf{cat} familiar, and travels with an \textbf{iron golem} bodyguard and a \textbf{homunculus} servant. He most often communicates via \textit{sending}, \textit{Rary's Telepathic Bond}, and \textit{project image} spells, or by sending his \textit{simulacrum} instead. Eldrick always has a \textit{contingency} spell in place, and has a \textit{clone} kept in a distant stronghold with extra equipment should the unthinkable occur.

\subparagraph{Eldrick's Simulacrum}
This special \textit{simulacrum} uses the game statistics for an \textbf{Archmage}.

\subparagraph{Eldrick's Spellbooks}
Eldrick keeps many spellbooks in secret places, which contain every spell on the Wizard spell list except \textit{wish} (knowledge of this spell is restricted to the Academy Directorate).

\subsubsection{River}
\includegraphics[width=0.4\textwidth]{images/../5e.tools/img/adventure/DoDk/012-02-007.river.png}
River is a tiefling \textbf{mage}. She is a short woman in her late twenties. Her attire is practical and professional, a hip-length purple jacket with a high collar, and belted at her waist are several pouches and pockets. She wears five rings on her right hand. Two oddly-shaped dark brown horns sprout backwards from her flowing blue hair, creating the impression of two boats at sea. River has spotted brown skin and impish features: her teeth are too sharp, her tongue too long, and it's slightly unsettling when she smiles.


\begin{description}
\item[Personality Trait.]
I converse using precise academic terminology, but amongst those I trust as competent and capable, I inject a wry and sardonic sense of humor.
\item[Ideal.]
I like people who are professional, deliver results, and don't ask questions. I cannot stand fools who waste my time with buffoonery, haggling, and excuses.
\item[Bond.]
Eldrick Runeweaver is my mentor, teacher, and friend. I've worked alongside him for years, and any success he has will be mine as well.
\item[Flaw.]
I grew up surrounded by the magical amenities of the Amethyst Academy, so I don't understand what life is like for those who had a more mundane upbringing.
\end{description}

River always has \textit{sending} and \textit{teleportation circle} spells prepared, and is accompanied by her \textbf{imp} familiar and a \textbf{shield guardian} bodyguard. She wears a \textit{ring of spell storing} that also acts as her Shield Guardian Amulet. Only River can attune to this ring. She also carries 2d6 uncommon Spell Scroll or potions.

\subsubsection{Academy Mages}
A handful of other mages lead survey teams in the ruins. Here are several NPCs to fill out the various Academy agents the player characters may encounter in Drakkenheim:


\begin{description}
\item[Rufus Stonewall.]
A grumpy old professor who is looking forward to retiring.
\item[Nimble Quill.]
An eager new graduate hoping to prove themselves.
\item[Cogsworth B. Copperpot.]
An addled old diviner confused about his own predictions.
\item[Murdok the Caller.]
A know-it-all conjurer with a knack for having the perfect spell prepared.
\item[Corvus Greysky.]
A grim necromancer with a morbid sense of humor.
\item[Elvira Evermore.]
A stern abjurer who takes magic use very seriously.
\item[Luna Mumph.]
A young gnome illusionist who seems lost in her own world.
\item[Harmon Munk.]
An elf transmuter who is always misplacing his alchemical reagents.
\item[Neptune Lee.]
A charismatic enchanter who likes to boast about his accomplishments.
\item[August "Gus" Elderfire.]
An experienced evoker who never does things by the book.
\end{description}

\subsection{Missions}
\subparagraph{Crystal Sample}
Recover a delerium crystal from the Rat's Nest Tavern. (See chapter 6)

\subparagraph{Exclusive Contract}
Secure the cooperation of the Ironhelm dwarves prospecting in the Scar. (See chapter 6)

\subparagraph{Stolen Homework}
Steal the research and findings of malfeasant wizard \textbf{Oscar Yoren} from his lair in Reed Manor. (See chapter 6)

\subparagraph{Recover an Eldritch Lily}
If the Amethyst Academy learns about the formula for \textit{aqua expurgo}, they'll dispatch characters to gather samples from Queen's Park. (See chapter 7)

\subparagraph{Broker an Alliance}
The Amethyst Academy asks characters to broker a truce with either the \textbf{Queen of Thieves}, the Hooded Lanterns, or possibly, the Followers of the Falling Fire so they can secure their supply of delerium and access to the ruins.

\subparagraph{Plight of the Pale Man}
Slay the \textbf{Pale Man} and acquire his spellbooks. (See chapter 7)

\subparagraph{End the Iron Banshee}
Seek out and destroy Ryan Greymere, the Iron Banshee. (See chapter 7)

\subparagraph{Reclaim the Inscrutable Tower}
The Amethyst Academy asks characters of 9th level or higher to participate in a raid mission to reclaim the Inscrutable Tower. (See chapter 8)

\subparagraph{Expedition to the Crater}
The Amethyst Academy commissions the characters to mount an expedition to the Crater Basin and see what's there. (See chapter 7)

\subparagraph{The Sequestered City}
Once the Amethyst Academy completes its goals of obtaining the \textit{Inscrutable Staff} and claiming the Tower, they ask the characters to help them subtly manipulate or sabotage the other factions into leaving Drakkenheim. (See chapter 10)

\subsection{Boons}
\subparagraph{Arcane Supplies}
The Amethyst Academy gives each character one uncommon potion or \textit{spell scroll} (of a spell from the Wizard spell list) of their choice when they begin a mission that directly advances the Academy's goals. Higher-level characters with proven reputations may be offered items of rare or very rare quality instead.

\subparagraph{Arcane Spellcasting Services}
River and Eldrick Runeweaver can offer their services as spellcasters to the characters as needed (they won't personally travel to the ruins, however).

\includegraphics[width=0.4\textwidth]{images/../5e.tools/img/adventure/DoDk/013-02-008.living-deep-haze.png}
\subparagraph{Material Components}
The Amethyst Academy can provide material components for spells of 5th level or lower, whose sum total value is no more than 500 gold pieces, prior to any mission.

\subparagraph{Occult Lore}
The Amethyst Academy is the foremost authority on magical lore and arcane phenomena. Their sages can answer questions about contemporary magic, magical creatures such as constructs, elementals, fey, fiends, and undead, any spell or magic item in the Core Rules, and historical facts about the past five hundred years. If you feel the matter is something that might require some deeper research, it takes the Academy one week to unearth relevant information. The Academy uses its discretion when revealing details about sensitive topics—especially information about the meteor or delerium, and never reveals its own secrets or details about its internal politics to outsiders.

\subparagraph{Elemental Summons}
If the player characters have a reasonable need for reinforcements during a mission, the Amethyst Academy offers the characters an \textit{elemental gem}.

\subparagraph{Magical Artifice}
As a reward for completing an Academy mission, the mages offer the party one uncommon magic item of their choice, subject to GM approval. The Amethyst Academy grants this boon once per player character. Once each character has received an uncommon item, the Academy will begin granting them one rare item each, then one very rare item each, as appropriate for completing missions of commensurate importance.

\subparagraph{Guild Membership}
A sorcerer, warlock, or wizard who can cast spells of 5th level or higher is offered guild membership in the Amethyst Academy. If they accept, they are given a purple Robe of Useful Items, a \textit{wand of magic missiles}, and a \textit{ring of spell storing}. These items bear Academy iconography, and clearly identify the character as an Academy member when worn.

\subparagraph{Library Gift}
Characters instrumental to retaking the Inscrutable Tower who demonstrate a complete commitment to the Academy's ideals and goals are rewarded with their choice of a \textit{manual of bodily health}, a \textit{tome of clear thought}, or a \textit{tome of understanding}.

\subsection{Strike Teams}
\subsubsection{Summoned Spies}
The Amethyst Academy agents send summoned creatures and magical constructs into the ruins to gather information: either 7 (2d6) \textbf{mephits}, 2 (1d4) imps, or a single \textbf{homunculus}. As magical communications are disrupted in the Haze, these creatures must survive and escape the city to report their findings, so they rarely engage in combat. These creatures often do not survive these missions, and some get lost and never return. Nevertheless, they're an essential tool for the Academy to have eyes and ears inside Drakkenheim.

\subsubsection{Academy Survey Team}
Though rare, a handful of mages lead survey teams into the ruins to carry out field research. They are typically encountered conducting experiments, observing magical phenomena, facilitating communications between other Academy agents, and providing magical support as needed. However, if the situation in Drakkenheim escalates, these agents help by casting spells, orchestrating traps, and summoning creatures.

A survey team is led by a \textbf{mage} with a \textbf{pseudodragon} familiar and a \textbf{shield guardian} bodyguard. The \textbf{mage} wears a \textit{ring of spell storing} that also acts as a Shield Guardian Amulet, but only the mage can attune to this ring. They also possess 1d6 uncommon Spell Scroll or \textit{potions}. They always have the \textit{sending} and \textit{teleportation circle} spells prepared.

Most travel with 1d4 research assistants: either hedge mages or academy apprentices. An Academy Apprentice uses the game statistics of a \textbf{commoner} with Intelligence 13 and the ability to cast two cantrips and one 1st level-spell from the Wizard spell list. (Spell Attack +3, Spell Save DC 11)

\subsubsection{Bound Elementals}
When the Amethyst Academy requires extra muscle, they conjure elemental creatures and press them into their service using \textit{planar binding} spells. Air, earth, fire, and water elementals are popular additional bodyguards for Academy Surveyors, and the Amethyst Academy will conjure invisible stalkers if they need to dispatch a dangerous enemy.

\subsection{Schemes}
The Amethyst Academy is a surreptitious foe. When the Amethyst Academy acts, it does not appear to be doing anything at all. They conduct their affairs like a business, and orchestrate events behind the scenes to turn out in their favour.

\subparagraph{Magical Surveillance}
If the player characters pose a threat to Amethyst Academy's goals, the mages begin using \textit{scrying} spells to observe the characters when they are outside the Haze, and send summoned spies to tail them within the city. At first, their aim is to simply learn as much about the characters and their plans as possible, then use this knowledge to dispatch their own strike teams to disrupt the players. Eventually, they may begin using \textit{dream} spells to interrupt their rests!

\includegraphics[width=0.4\textwidth]{images/../5e.tools/img/adventure/DoDk/014-02-009.amethyst-academy-wax-seal.png}
\subparagraph{Well-equipped Rivals}
If the player characters do not work with them, the Amethyst Academy invests its resources in rival adventuring teams instead. In tandem with their magical surveillance, the Amethyst Academy sends these rival adventurers to sabotage the characters missions. In many cases these rival adventurers are given a false premise and don't actually know disrupting the player characters is their primary purpose. The Amethyst Academy might even charge rival adventurers to accomplish one of the player characters' goals before they can do it (such as taking the Clocktower, recovering a Relic of Saint Vitruvio, finding a Seal of Drakkenheim, or something important to a personal quest). They keep the item as a bargaining chip.

\subparagraph{Shadow Assassins}
Should the Amethyst Academy need to dispose of the characters entirely, they dispatch one \textbf{invisible stalker} per character to carry out the deed while the characters are sleeping. They use a \textit{scrying} spell to confirm the characters are resting, and a \textit{teleport} spell to move their agents into position nearby. If the characters are using magical means to protect themselves, such as a \textit{Leomund's Tiny Hut} spell, they send one \textbf{mage} who casts \textit{greater invisibility} upon themselves and then casts \textit{dispel magic} to destroy the barrier. The invisible stalkers are instructed to coordinate their attacks on one character at a time, starting with healers and spellcasters, and will attack downed characters to confirm a kill.

\subparagraph{Stark Demonstration}
Under the flag of truce, the Amethyst Academy invites the player characters and faction leaders to attend a demonstration of their latest creation with delerium. Known as the Disjunction Bomb, this devastating weapon can annihilate an entire city district. They matter-of-factly detonate a test weapon in a cleared wilderness area, and then offer to discuss terms with the characters and the factions about how to best proceed.

\DndQuote%
{}
{''They teach you magic, then you get expelled for using magic too much. Relationships are... complicated.''}
{Sebastian Crowe}
\section{The Followers of the Falling Fire}
\DndQuote%
{}
{''One day soon, we will near the end of our world—an age where great heroes will be needed. These holy stones were heaven-sent so we might prepare our spirits to meet the all-consuming shadows. In our darkest moment, the Sacred Flame will set alight our blazing hearts as one.''}
{Lucretia Mathias}
The Followers of the Falling Fire are a religious sect who believe the meteor signalled a celestial prophecy, and consider delerium a holy sacrament. Called to pilgrimage by the prescient testament of their leader, \textbf{Lucretia Mathias}, zealous masses of devout common-folk now journey to the crater's edge to take their place in a divine plan.

\subsection{Background}
\includegraphics[width=0.4\textwidth]{images/../5e.tools/img/adventure/DoDk/015-02-010.falling-fire.png}
The Faith of the Sacred Flame has long looked to the stars for guidance. The flickering stars cast within the Astral Void symbolize how the faithful see their place in the cosmos: points of light holding darkness at bay. Scripture tells how a shooting star inspired Saint Tarna to seek redemption before a pair of angels. The story of a blackguard turned paladin who founded the Faith itself established an important message within the religion: that from the deepest darkness may arise the brightest light.

\textbf{Lucretia Mathias} was one such Flamekeeper who read omens in the night sky, and was regarded as one of the most learned and wise amongst the clergy. She channeled immense divine power, and invoked miracles that have not been witnessed in a generation. A decade before the meteor destroyed Drakkenheim, she beheld a vision of a falling star that would arrive on earth and signal a great upheaval for the faithful. Despite her reputation, her visions were dismissed as an apocalyptic fantasy and largely suppressed at the time.

Her visions proved true; right down to the hour. However, even Lucretia Mathias did not expect the actual destruction the meteor would cause, nor did she anticipate it would annihilate her birthplace, Drakkenheim. She spent several years afterwards contemplating the heavens and several strange stones recovered from the city. Her conclusion was dire indeed: the world stands on the precipice of darkness unending. Yet within delerium, she saw the spark of salvation. Though the stones are the stuff of chaos, black charcoal and midnight oil may fuel a great flame. Perhaps if taken up by righteous souls, the crystals might become an earthly vessel for the Sacred Flame itself. Rather than fall to all-consuming shadows, Lucretia Mathias prophesied a New Age of Heroes that could meet the coming night. She spoke widely about her beliefs and discoveries, and her fiery sermons galvanized a throng of loyal supporters who proselytized her new creed across Elyria, Westemär, and Caspia. As her support grew, she resolved to make a pilgrimage to the crater itself to confirm her revelations.

\includegraphics[width=0.4\textwidth]{images/../5e.tools/img/adventure/DoDk/016-02-011.lucretia-mathias.png}
Meanwhile, the mainstream clergy reacted to delerium with disgust and horror. When Lucretia Mathias returned from her pilgrimage bearing a delerium crystal embedded in her own body, she was declared a heretic and excommunicated. However, she refused to be silenced. Escaping Elyria, she expanded her original writings and published them as the \textit{Testament of the Falling Fire}. This religious essay proclaimed a radical new doctrine for the faith, and called upon all devout believers to make pilgrimage to the crater in Drakkenheim. Those who finish this arduous journey partake in a controversial rite known as the \textit{Sacrament of the Falling Fire}.

\subsubsection{Attire}
The Followers of the Falling Fire give up most of their worldly possessions. Many are clad only in sackcloth jerkins or monastic robes, and often carry salvaged weapons adapted from everyday equipment and tools.

\subsubsection{Strongholds}
While the Falling Fire is now sending missionaries out all over the world, they have four important holdings in Drakkenheim.

\subparagraph{Hendrix Farm}
Pilgrims camp out here in Emberwood Village before embarking on the last stage of their journey to the crater.

\subparagraph{Champion's Gate}
Not a stronghold in any true sense, the pilgrims use Champion's Gate as a waypoint on their way to the crater's edge or Saint Selina's Monastery.

\subparagraph{Chapel of Saint Gresha}
Located on the edge of the crater, this is the place where Lucretia Mathias performs the \textit{Sacrament of the Falling Fire} each day at dawn.

\subparagraph{Saint Selina's Monastery}
Located in the Deep Haze, this is the sanctuary for members of the Falling Fire who have completed their pilgrimage and taken the Sacrament.

\section{Roleplaying Characteristics}

\begin{description}
\item[Symbol.]
A soaring comet. Their colours are black and white.
\end{description}

\subsection{Personality Traits}

\begin{itemize}
\item We often speak prayers and quote the writings of Lucretia Mathias. We enthusiastically share the message of our new faith with everyone.
\item We are fascinated by the stars, and see omens and prophecy in everything.
\item We see hope for the downtrodden and absolution for the lost. We offer deliverance and mercy for any who seek the truth and the light, but we condemn those who do not.
\end{itemize}

\subsection{Ideals}

\begin{itemize}
\item A New Age of Heroes shall meet the Creeping Darkness that threatens the world. Through our devotion, we shall light the Sacred Flame.
\item Every soul is worthy of redemption, you need only see the good within. From the darkest shadow comes the brightest light.
\item Material possessions and monetary wealth hold no value anymore. All that matters is our faith.
\end{itemize}

\subsection{Bonds}

\begin{itemize}
\item Delerium shards have a sacred purpose as the key to our salvation. They must only be used in our religious ceremonies. It is utter profanity to turn the stones into arcane trinkets, and total sacrilege to destroy the holy crystals.
\item The ruins are a holy site now. We must make a pilgrimage to the crater to take part in the \textit{Sacrament of the Falling Fire}. This test of faith will demonstrate our devotion, and we will know our place in the divine plan.
\item We are disillusioned by the hypocrisy of the mainstream religious ministry of the Sacred Flame. The Flamekeepers have tarnished our faith with political ambitions and affluent vices.
\end{itemize}

\subsection{Flaws}

\begin{itemize}
\item Lucretia Mathias and her disciples who have completed the \textit{Sacrament of the Falling Fire} have become living vessels for the Sacred Flame. We follow them faithfully and heed their every word unquestioningly.
\item The divine prophecy of Lucretia Mathias is irrefutably true. Any evidence to the contrary is merely a deception meant to dissuade us from our righteous path.
\item The old world shall inevitably end in darkness, petty political squabbles and earthly concerns are meaningless compared to our holy mission.
\end{itemize}

\DndQuote%
{}
{''What strange manners to travel into the dangers of this city with nothing but faith as a shield. The border between brave and foolish is thin among them.''}
{Pluto Jackson}
\subsection{Objectives}
\subsubsection{Transform Drakkenheim into a Holy Site}
\textbf{Lucretia Mathias} has proselytized that hundreds of thousands, perhaps even millions of devoted followers will need to complete the pilgrimage to fulfil her prophecy. Once they do, the Sacred Flame will set alight every righteous mortal heart, and a glorious new age of heroes shall begin. Until then, they must ensure many more can make the journey, and preserve the delerium, the sacred charcoal needed to enact their rituals. Thus, they wish to make Drakkenheim itself a holy site.

\includegraphics[width=0.4\textwidth]{images/../5e.tools/img/adventure/DoDk/017-02-012.nathaniel-flint.png}
Those who complete their pilgrimage to the crater's edge participate in a ritual called the \textit{Sacrament of the Falling Fire}. This controversial rite requires supplicants to impale a delerium fragment into their hearts. Those who survive the pilgrimage are spiritually transformed and physically empowered, rendered immune to contamination and able to survive in the Haze. Many who complete the pilgrimage suddenly discover they are able to tap into powerful spiritual energies, which they claim are the Sacred Flame. Thus many followers channel divine magic who had not before exhibited the necessary willpower to master such abilities.

\subsubsection{Spread and defend their new faith}
The Followers of the Falling Fire seek to convert others to take up their cause and their beliefs, no matter their former lives. They believe members of the Hooded Lanterns and the Silver Order could be convinced to abandon their organizations to join the Followers of the Falling Fire.

Should a new ruler of Westemär emerge, the Followers of the Falling Fire wish to convert this individual to their faith. Doing so will legitimize the Falling Fire politically, and make sure that they can continue their pilgrimages.

At the same time, they must defend their faith. The Followers of the Falling Fire have been declared heretics and excommunicated by the clergy of the Sacred Flame, and must protect themselves from attacks (both physical and political) by the Knights of the Silver Order.

\subsubsection{Claim the lost relics of Saint Vitruvio}
Drakkenheim was home to many saints of the Sacred Flame: most notably, the heroic Saint Vitruvio, who fought astride the mighty gold dragon Argonath centuries ago. Both were buried beneath the cathedral that bears his name. Many relics of this mighty paladin are lost in Drakkenheim, such as \textit{Ignacious, the Sword of Burning Truth}. If the Followers of the Falling Fire were to obtain these relics and resurrect the dragon it would significantly legitimize their beliefs to many laypeople.

They wish to reclaim these artefacts and reconsecrate the altar of the Sacred Flame in the Cathedral of Saint Vitruvio. The Saint Vitruvio's Phylactery is also one of the \textit{Seals of Drakkenheim}. Traditionally held by a Flamekeeper of the Sacred Flame, possessing it ensures the Followers of the Falling Fire gain influence over the political destiny of Drakkenheim and Westemär.

\subsection{Key Characters}
The vast majority of the Followers of the Falling Fire are commoners, mere pilgrims seeking salvation in the ruined city. Many bring children, infirm, or sick family members seeking miracle healing. Hundreds have arrived already, and thousands more are coming.

\subsubsection{Lucretia Mathias}
\textbf{Lucretia Mathias} is a willowy and wizened woman with piercing eyes and pursed lips. She is in her early nineties, and keeps her greying black hair in a coiled bun beneath a linen shawl. She wears simple grey robes, and clutches a heavy leather-bound tome. A glowing golden crystal dangles from a silver chain around her neck. Her garb lightly conceals the delerium shard embedded in her heart.


\begin{description}
\item[Personality Trait.]
I speak in proverbs and quote scripture when I speak, inspiring others to interpret my cryptic words and come to their own revelations.
\item[Ideal.]
The brightest lights shine in the darkest of places. None are beyond redemption, and all may find their inner truth and purpose by following the Falling Fire.
\item[Bond.]
I have foreseen the inevitable end of our world, but through our sacred self-sacrifice we shall set alight the Sacred Flame to bring enlightenment and salvation to the universe.
\item[Flaw.]
I only reveal the greatest truths I have seen to those I am certain can grasp them fully, and will understand them as I do.
\end{description}

\subsubsection{Nathaniel Flint}
Nathaniel Flint is a human \textbf{chaplain} who completed the \textit{Sacrament of the Falling Fire}. He is an animated and jovial man in his early thirties, with medium length curly brown hair, mutton chops, a bushy mustache, and a flushed face. He does not wear armour, and instead wears pilgrims' robes with a roped belt and a tabard emblazoned with the sign of the Falling Fire. A scarfed hood conceals the delerium shard embedded in his chest.


\begin{description}
\item[Personality Trait.]
I approach every situation with a carefree eagerness.
\item[Ideal.]
I see the prophecies of Lucretia Mathias as completely reasonable and true, and don't understand why anyone wouldn't embrace them.
\item[Bond.]
There is no greater purpose to my being than to help others find their divine light and true purpose in the crater.
\item[Flaw.]
Everything that happens is part of a divine plan, no matter how horrible or unfortunate, it was always meant to be.
\end{description}

\subparagraph{Secret Identity}
At your option, Nathaniel Flint may in fact be Leonard von Kessel, but he has no memory of this fact. Alternatively, he could be connected to the royal family in some other way.

\subsubsection{Pilgrims of the Falling Fire}
Here's a few readymade NPCs who are travelling to Drakkenheim to complete their pilgrimage:


\begin{description}
\item[Brother Abraham Schaefer.]
A faithful old preacher naïve to the horrors of Drakkenheim.
\item[George Hopper.]
A courageous fighter who dutifully protects other pilgrims.
\item[Anita Winters.]
A strong and sturdy warrior who had visions of the meteor.
\item[Dekota Jones.]
A young explorer who questions everything.
\item[Cosmos.]
An optimistic tiefling who is obsessed with planar travel.
\item[Rufus Apollo.]
An old dwarven cleric who studies astrological signs and patterns.
\item[Ingrid and Mara Stumer.]
Sisters who followed dreams and visions to the city.
\item[Silvie Roseshot.]
A former Flamekeeper inspired by the Testament of the Falling Fire.
\item[Lucy Wainwright.]
A creative and intelligent woman with a child who came to the city for salvation.
\item[Jared Walsh.]
A young paladin who believes he will fulfil his oath by taking the Sacrament.
\end{description}

\subsection{Missions}
\subparagraph{Sceptre of Saint Vitruvio}
Retrieve the holy relic from the Chapel of Saint Brenna. (See chapter 6)

\subparagraph{The Lost Disciple}
A prominent disciple went missing near the Black Ivory Inn. (See chapter 6)

\textbf{Dwarven Defilers} Drive out the Ironhelm Dwarves who are mining delerium in the Scar. (See chapter 6)

\subparagraph{Pilgrims' Plight}
The vast majority of the Followers of the Falling Fire are common folk. There are few warriors amongst them. Protect a pilgrimage from Emberwood Village to Champion's Gate.

\subparagraph{Pilgrimage of Falling Fire}
Characters of 6th level or higher who faithfully serve the Falling Fire are invited to take the pilgrimage to the crater. (See chapter 7)

\subparagraph{Lighted Lanterns}
The characters are tasked with convincing the Hooded Lanterns to convert to the Falling Fire and serve their cause.

\subparagraph{Saint Vitruvio's Cathedral}
Destroy the \textbf{Lord of the Feast} and retake the cathedral. Afterwards, help us consecrate the Brazier for the Falling Fire and retrieve Ignacious, the Sword of Burning Truth.

\subparagraph{A Messenger of Faith}
Entreat members of the Silver Order to back down and join our cause - otherwise, destroy them!

\subparagraph{Drive out the Faithless}
We shall put an end to the Amethyst Academy's experiments with the holy delerium by invading the Inscrutable Tower. (See chapter 8)

\subparagraph{Shield of the Sacred Flame}
An expedition to Castle Drakken is needed to recover the fabled shield of Saint Vitruvio.

\subparagraph{Awaken Argonath}
We will need a pristine delerium geode from the Delerium Heart at the centre of the crater to complete the resurrection of Argonath.

\subparagraph{A Worthy Disciple}
Lucretia Mathias is wizened, and though her divine powers imbue her with vigour, she is not long for this world. She has not yet found a true disciple who could take up her mantle after she passes, and she knows that her followers will need to rally behind a new leader when she is gone.

\includegraphics[width=0.4\textwidth]{images/../5e.tools/img/adventure/DoDk/018-02-013.worshiper-andreiaugrai.png}
\subsection{Boons}
\subparagraph{Divine Spellcasting Services}
\textbf{Lucretia Mathias} can offer her services as a spellcaster to the characters. She can help them recover via spells such as \textit{lesser restoration}, \textit{prayer of healing}, \textit{purge contamination}, \textit{remove curse}, \textit{greater restoration}, and \textit{heroes' feast}. While characters must still provide any costly material components, she does not charge them to cast the spells.

\subparagraph{Celestial Lore}
Lucretia Mathias is an unparalleled expert on celestial phenomena, cosmology, religion, and divine magic. In addition, she possesses unique insight about delerium and the meteor, and has prescient divine visions. While she could conceivably answer questions about anything, she speaks in prophecy and mystery, and often encourages others to "seek the truth for themselves." She's more likely to guide characters towards an answer than outright reveal it. For example, while Lucretia Mathias has personally ventured to the centre of the crater and knows what lies there, she believes others must witness it firsthand. Furthermore, she has her own subjective point of view about events coloured by her faith, though she sincerely believes them to be true.

\subparagraph{Prophecy}
Lucretia Mathias can cast a wide range of divination spells, including \textit{commune}, \textit{divination}, \textit{legend lore}, and \textit{scrying}. She will cast these spells to gain knowledge and insight about future events.

\subparagraph{Revive the Fallen}
Lucretia Mathias is willing to cast \textit{raise dead}, \textit{resurrection}, or even \textit{true resurrection} upon a faithful champion. She will use delerium as a material component for these spells, and so any character raised by her gains a form of indefinite madness.

\subparagraph{Sanctified by the Falling Fire}
Characters who demonstrate full faith and devotion to the Falling Fire are invited to make a pilgrimage to the crater's edge and partake in the Sacrament of the Falling Fire. This ritual imbues characters with powerful protection against the Haze, but has a serious cost.

\subparagraph{Relics of the Faithful}
As a reward for completing an important mission, the Followers of the Falling Fire will grant the party one magical rod, staff, or wondrous item from their reliquary. This will be an uncommon item for an Outer City mission, a rare item for the Inner City or a faction stronghold, or a very rare magic item for a successful mission to Castle Drakken or the cathedral. This boon can be granted multiple times, once for each character in the party. The Game Master can either choose the item, or allow the player to choose it.

\subparagraph{Testament of the Falling Fire}
Lucretia Mathias allows characters to read one of her original manuscripts. These take the form of a \textit{Tome of Clear Thought}, a \textit{Tome of Understanding}, and a \textit{Tome of Leadership and Influence}.

\includegraphics[width=0.4\textwidth]{images/../5e.tools/img/adventure/DoDk/019-02-014.incense-brazier.png}
\subsubsection{Sanctified by the Falling Fire}
The \textit{Sacrament of the Falling Fire} is known only to \textbf{Lucretia Mathias} and \textbf{Saint Gresha}, who cast the spells each day at dawn and dusk respectively at the Chapel of Saint Gresha on the crater's edge. A character who participates in the \textit{Sacrament of the Falling Fire} gains the following traits, so long as they adhere to the tenets and faith of the Falling Fire:


\begin{itemize}
\item You gain advantage on saving throws against being charmed or frightened, and resistance to \textit{necrotic} and \textit{radiant damage}. You are immune to contamination.
\item If you die, your soul enters the delerium fragment embedded into your chest. You can't be raised from the dead while your soul is in the fragment.
\end{itemize}

The embedded delerium fragment does not contaminate or otherwise harm a character who keeps their conviction.

\subsubsection{Loss of Faith}
A character who violates or loses their faith in the Falling Fire loses the protections of sanctification and the beneficial traits above. Instead, their body is directly exposed to the delerium in their chest and they immediately gain one level of contamination. Until the fragment is removed or they atone for their sins, their hit point maximum is reduced by 10 (3d6) and they gain a level of delerium contamination each time they finish a long rest.

\subsubsection{Sanctified Delerium}
A delerium fragment filled with a sanctified soul glows with golden light. It becomes inert and immune to all damage except bludgeoning damage from magical weapons. Sanctified delerium possesses no other magical properties and is worthless for all arcane applications. A creature holding the fragment can telepathically communicate with the soul held within. If a sanctified fragment is destroyed, the soul inside is utterly annihilated, but the fragment can be used in a spell to return the dead soul within to life.

\subsubsection{Removing the Fragment}
The sacrament ritual fuses delerium to the body. The fragment can only be physically removed through great effort, requiring one hour of work and a DC 25 Medicine check. On a successful check, the fragment is removed but the sanctified character suffers 55 (10d10) damage. The fragment can be harmlessly removed from a willing creature using a \textit{heal} or \textit{wish} spell.

\section{What has Lucretia Mathias foreseen?}
Within the quotation at the start of this section, the entire truth of Lucretia Mathias' visions are laid bare. While it might \textit{seem} like she is speaking in metaphor, she is actually being quite literal. Lucretia Mathias has seen how the Haze will consume the world, but also knows that the spread of delerium through comets has continued for untold eons. She intends to break this cycle.

"From the deepest darkness comes the brightest light." is a core doctrine for the Sacred Flame, and Lucretia Mathias. believes the chaotic and destructive energies of the Haze and delerium can be re-made into a vessel for light and order.

One must only plant the seed of righteousness within. Thus the purpose of the \textit{Sacrament of the Falling Fire} is only realized when participants die and their souls are drawn into the delerium crystals in their chests. Lucretia Mathias believes that if enough of these sanctified stones are gathered by the time the world ends, a divine apotheosis will result in a new generation of harbingers of light and salvation, rather than chaos and doom. Whether or not this will actually happen is impossible to say. Since Drakkenheim is a world of distant and unknowable gods, there is no way to prove or disprove her claims. It is entirely a matter of faith.

\subsection{Strike Teams}
The Followers of the Falling Fire lack military training and possess few resources. What they do have, however, is great faith and determination. In battle, the Followers of the Falling Fire surge forward in a mob to engage in zealous combat. However, they staunchly defend one another and look out for their fellow believers. They have little regard for tactics or strategy, believing instead that their actions are guided by divine providence instead. They rarely retreat.

\subsubsection{Pilgrimage}
A mob of 21 (6d6) commoners led by a \textbf{priest}. Such groups may also include several noncombatants, as well as wagons pulled by draft animals such as oxen or mules.

\subsubsection{Militant Faithful}
10 (3d6) cultists led by 2 (1d4) cult fanatics.

\subsubsection{Sanctified Zealots}
An assorted group of 7 (2d6) berserkers, cult fanatics, and priests. These individuals have completed the Sacrament, and have the Sanctified by the Falling Fire traits.

\subsubsection{Converted Clergy}
Any other strike team might be joined by a \textbf{knight}, \textbf{chaplain}, or \textbf{cavalier} originally a member of the mainstream clergy who has shifted their faith to the Falling Fire. Many might even be former members of the Silver Order.

\subsubsection{Angelic Avenger}
\textbf{Lucretia Mathias} can summon 1d4 devas or a \textbf{planetar} to lead another strike team.

\subsection{Schemes}
\subparagraph{Celestial Visions}
If the player characters pose a threat to Followers of the Falling Fire, Lucretia Mathias begins using \textit{scrying} spells to observe the characters when they are outside the Haze, and using \textit{divination} and \textit{commune} spells to learn about their personality traits, ideals, bonds, and flaws. Her aim is to simply learn about the characters, their motivations, and actions.

\subparagraph{Compelling Plea}
Lucretia Mathias uses her knowledge of the player characters to implore them to join the Followers of the Falling Fire. She uses information gleaned about the characters' ideals, bonds, and flaws to impress the importance of her own visions. She makes prophecies and delivers omens that the characters will not succeed without faith, and perhaps even exposes their vices, flaws, and past sins. However, in bringing up the characters' failures, secrets, and flaws, she offers them hope for atonement and absolution by joining her.

\subparagraph{Salvation and Penitence}
If the characters are in dire circumstances, Lucretia Mathias receives a vision of their plight and dispatches a Falling Fire strike team to rescue them.

The strike team brings the characters back to Saint Selina's and offers them healing and help, but also takes their possessions and tries to forcefully convert them. They may even imprison the characters until they repent.

\subparagraph{Conversion}
If the player characters have any family members or loved ones, the Followers of the Falling Fire seek them out and attempt to convert them. Alternatively, they convert members of the Queen's Men, Silver Order, or Sacred Flame to join them.

\subparagraph{Truth of the Crater}
Lucretia Mathias invites the player characters to accompany her to the centre of the crater so they can see for themselves the truth of what is to come. If they do not agree to join her, she incites the monsters therein to attack, relying on her own divine power to protect herself.

\DndQuote%
{}
{''They may have strong faith and good intentions, but you've gotta be out to lunch if you're willing to stab delerium into your chest!''}
{Veo Sjena}
\section{The Hooded Lanterns}
\DndQuote%
{}
{''There cannot be a Westemär without Drakkenheim.''}
{Lord Commander Elias Drexel}
The Hooded Lanterns are an \textbf{irregular military force} drawn up from remnants of the old Drakkenheim City Watch and veterans of the Westemär Civil War. Led by the grim \textbf{Lord Commander Elias Drexel}, they wish to drive out the monsters and outlaws who infest their nation's capital and restore the realm of Westemär.

\subsection{Background}
\includegraphics[width=0.4\textwidth]{images/../5e.tools/img/adventure/DoDk/020-02-015.hooded-lanterns.png}
Multiple military operations have attempted to retake Drakkenheim over the past fifteen years. Each has ended in miserable failure, and the contaminated ruins have swallowed thousands of soldiers' lives. What remains of these assaults are merely twisted monsters and undead monstrosities formed from the risen corpses of these regiments. After these disastrous attempts, Drakkenheim was left abandoned while civil war ripped Westemär apart instead. These conflicts also ended in vain, with the royal line of House von Kessel undone by fratricide.

Westemär is slowly disintegrating into a realm of squabbling fiefdoms, yet some amongst the fractured nobility hold a glimmer of hope in restoring the nation and rebuilding the capital. Three years ago, they marshaled up the remnants of the royal guard, the city watch, and patriotic veterans of the civil war to make one last-ditch attempt to retake the city. Although officially dubbed the 4th Provisional Expeditionary Force to Reclaim the Capital, this guerilla regiment was nicknamed after oil lamps carried by the old city watch: the Hooded Lanterns.

The core members of the Hooded Lanterns are warriors fighting for their homeland. Few honest soldiers survived the civil war, so many veterans have a life of bloody deeds behind them already. On the other hand, new recruits are seldom old enough to have any strong memories of Drakkenheim before it was destroyed. As such, many possess naïve and idealistic notions honed into simmering rage and fury at a golden age stolen from them.

Under the leadership of \textbf{Lord Commander Elias Drexel}, the Hooded Lanterns have finally broken the losing streak of previous military operations. Rather than making an all-out assault on the ruins, they have taken a methodical approach. After spending nearly a year carefully scouting out the ruins, they began clearing out important fortifications in the city to use as staging grounds for future operations. This strategy has paid off: the Hooded Lanterns have managed to establish countless safehouses and supply caches within the city. They captured Shepherd's Gate and the Drakkenheim Garrison last year. However, in recent months their efforts have been stymied by sabotage enacted by the \textbf{Queen of Thieves}, and the emergence of a horrific garmyr warlord in the Inner City known as the \textbf{Lord of the Feast}. While the Hooded Lanterns have imposed the mantle of law-and-order in the city, these setbacks have revealed their actual authority is limited to their strength of arms. The regiment maintains only a fragile foothold in Drakkenheim, and their supply lines are now exceptionally difficult to maintain. It's taking almost everything the Hooded Lanterns have to hold their current position.

\includegraphics[width=0.4\textwidth]{images/../5e.tools/img/adventure/DoDk/021-02-016.elias-drexel.png}
The Hooded Lanterns are supported by surviving loyal vassals of the royal family who wish to see the old monarchy of Westemär restored. Many of these wealthy nobles lost estates, treasure vaults, family heirlooms, and personal documents when the city was destroyed, and expect the Hooded Lanterns to retrieve them. In return, these noble houses provide soldiers, material, and other supplies to the Hooded Lanterns on a periodic basis, but their financial resources are not inexhaustible. Indeed, the Hooded Lanterns have often resorted to selling delerium themselves to make ends meet.

\subsubsection{Attire}
The Hooded Lanterns are a professional military unit. Every member uses standardized equipment: their uniform consists of a forest green cloak with yellow trim, the traditional colours associated with the Drakkenheim coat-of-arms. They bear the Hooded Lanterns' symbol on their chest pieces, and shoulder stitches denote their rank and company. Scout units wear light armour, and often wield Longsword, Shortsword, and Longbow. Guard units are equipped with Chain Mail, Shield, Spear, and crossbows. Each member also wears a metal chain with an inscribed plate bearing their name, hometown, and next of kin. Most Hooded Lanterns cut their hair when they sign up, but since many spend weeks at a time fighting, they end up a disheveled mess by the end of their assignments.

\subsubsection{Strongholds}
The Hooded Lanterns hold three key positions in Drakkenheim and its environs: the Emberwood Watchtower, Shepherd's Gate, and the Drakkenheim Garrison. They have converted the stockade beneath the garrison into a bunker insulated from the Haze. Beyond these key locations, the Hooded Lanterns have several minor bases in old cottages and watchtowers a few days outside the city where their soldiers rest and recuperate after being posted to the ruins.

\section{Roleplaying Characteristics}
\subparagraph{Symbol}
A hooded lantern held by a dragon's claw.

\subsection{Personality Traits}

\begin{itemize}
\item We're not looking for glory or wealth, we're fighting for our home. We expect others to respect the seriousness and importance of our goals.
\item We are resolute in facing dangerous monsters and terrible horrors head on. These mutated creatures aren't human anymore.
\item If you're willing to join our cause, we'll embrace you as a loyal comrade-in-arms.
\end{itemize}

\subsection{Ideals}

\begin{itemize}
\item We all must do our duty to restore our nation. When your commander gives you an order, you follow it without question. You must be prepared to give your life for the cause.
\item In battle, you have to seize every advantage. A plan is worth more than any weapon.
\item We will uphold the old ways of Westemär. This must be a city of laws and order. They mean no less today than they did twenty years ago.
\end{itemize}

\subsection{Bonds}

\begin{itemize}
\item This city was our home, and it will be again. It's the heart of the realm. There cannot be a Westemär without Drakkenheim.
\item We trust our comrades-in-arms with our lives and are deeply loyal to one another. We staunchly believe that together we can accomplish what none of us could do alone.
\item The old nobility of Westemär provide us with material and supplies. We have a responsibility to ensure their investments and interests are protected, otherwise we won't have the resources to continue our fight.
\end{itemize}

\subsection{Flaws}

\begin{itemize}
\item After a decade of civil war and years fighting for survival in the ruins, most of us know nothing but war. Our young recruits don't even have any memory of the old Drakkenheim.
\item We all knew the risks when we signed up. If we are to win the battles ahead, we must be willing to sacrifice the lives of our comrades.
\item Drakkenheim must be rebuilt the way it used to be. The crown, the castle, the city must all be retaken and rebuilt. It wouldn't be home otherwise.
\end{itemize}

\DndQuote%
{}
{''Adapative and cunning, they are true survivalists. Some of the bravest warriors I've met in Westeimär.''}
{Pluto Jackson}
\subsection{Objectives}
\subsubsection{Restore Drakkenheim}
The patriotic warriors of the Hooded Lanterns desperately want to rebuild their former home. The rank-and-file soldiers fight a guerilla war to reclaim Drakkenheim street-by-street, slaying monsters, driving out outlaws, and arresting scavengers. \textbf{Lord Commander Elias Drexel} plans on establishing strongholds at key sites of strategic importance: the City Gates, Saint Selina's Monastery, Saint Vitruvio's Cathedral, the Clocktower, and eventually, Castle Drakken. Thus far, the Hooded Lanterns have only succeeded in taking Shepherd's Gate and the Drakkenheim Garrison.

\includegraphics[width=0.4\textwidth]{images/../5e.tools/img/adventure/DoDk/022-02-017.the-master-of-the-forge.png}
\subsubsection{Obtain the Seals of Drakkenheim}
The ruling council of Westemär wore enchanted badges of office attuned to the magical defences and arcane wards that once protected Drakkenheim, and also served as the keys to Castle Drakken itself. Known as the \textit{Seals of Drakkenheim}, all are lost within the ruins save the \textit{Lord Commander's Badge}, carried by Elias Drexel himself. Until all the seals are assembled, the city's magical defences are subject to the chaotic whims of the Haze. The Hooded Lanterns are constantly searching the ruins for any sign or clue of the whereabouts of these items. Unfortunately, several are already in the hands of rival factions. See Appendix D for their game statistics.

\subsubsection{Re-establish a rightful monarchy}
Many Hooded Lanterns hope a legitimate heir to the throne still survives somehow. Others believe that even if House von Kessel is truly extinct, a new monarchy could be formed by a prominent noble house with some connection to the royal line. Both Saint Vitruvio's Cathedral and Castle Drakken are believed to contain archives detailing the line of succession, and magically preserved vials of royal blood that can be used in a special ritual to confirm legitimate royal ancestry. A potential claimant with access to these documents could galvanize the support needed to ascend to the throne. However, proof of royal ancestry alone is not enough - any claim is spurious at best without the \textit{Seals of Drakkenheim}, the \textit{Crown of Westemär}, and Castle Drakken. Even a true heir to House von Kessel will have immense difficulty holding the Realm of Westemär together without these keys to power.

\subsection{Key Characters}
\textbf{Lord Commander Elias Drexel} leads the Expeditionary Force as the Faction leader. He has two close confidants who interchangeably act as the faction lieutenant: Captain Ansom Lang and his sister, Lieutenant Petra Lang.

\subparagraph{NOTE}
The Hooded Lanterns use a formal military rank structure. Although Ansom and Petra are ranked as "captain" and "lieutenant", this doesn't change their role as faction lieutenantsfor the purposes of their interactions with the player characters.

\subsubsection{Lord Commander Elias Drexel}
Elias Drexel is a man in his early fifties. He is stocky and imposing with a furrowed brow and a grim frown on his wide and angular face. His harsh features are framed with a thick moustache, side chops, and a mane of long braided brown hair. He is garbed in chainmail armour with heavy leather gloves and boots, worn under a fur-lined dark green cloak and iron pauldrons. He wears a pinned brooch - the \textit{Lord Commander's Badge}, the symbol of his office. A quiver, a longbow, and two longswords are strung on his back.


\begin{description}
\item[Personality Trait.]
I am blunt and direct, and expect my soldiers to obey my commands and deliver results promptly. I find no humor in anything.
\item[Ideal.]
Loyalty and trust are priceless. They must be carefully given, but staunchly held. There is nothing worse than betrayal and treason.
\item[Bond.]
During the Civil War, I switched sides and killed several scions of House von Kessel who could have claimed the throne. I once believed it was for the right cause, but now I feel personally responsible for breaking the royal line. My mission in Drakkenheim is penance and redemption.
\item[Flaw.]
I'm filled with rage and conflict because of my contradictory actions and beliefs, and I don't take defeat, setbacks, or even missteps lightly.
\end{description}

Elias Drexel was appointed to the office of Lord Commander before the fall of Drakkenheim, and was outside the city leading training drills when the meteor fell. Afterwards, he took part in three military expeditions to reclaim Drakkenheim, and saw firsthand their horrors. He then went on to fight in the Westemär Civil War alongside Mannfred von Kessel. However, the idea that his soldiers died for nothing in Drakkenheim festered in his mind, and made him doubt Mannfred's ability to be a good ruler. He secretly participated in a successful conspiracy to assassinate Mannfred von Kessel. Little did he know that shortly thereafter, tragedy would soon befall Cecilia von Kessel. He isn't exactly proud of his choice to betray Mannfred, and that causes him considerable guilt, regret, and internal self-loathing.

\subsubsection{Command Staff}
The Lord Commander has several important aides who assist with the day-to-day operations of the Hooded Lantern's forces:


\begin{description}
\item[The Warden.]
a \textbf{knight} who oversees the defences and dungeons of the City Watch Barracks.
\item[The Quartermaster.]
a \textbf{spy} who attends to the supply cache.
\item[The Apothecary.]
a human \textbf{druid} who tends the wounded and manages their hospice.
\item[The Master of the Forge.]
a dwarf \textbf{priest}.
\item[Paige.]
a \textbf{commoner} and personal attendant to the Lord Commander.
\end{description}

\subsubsection{Petra Lang}
Lieutenant Petra Lang is a human \textbf{urban ranger}. She is a fit and nimble teenager with long black hair, a narrow pointed chin, and bold features. She wears a hooded green cloak and light chain armour of the Hooded Lanterns with a lieutenant's insignia on her shoulder pauldron.


\begin{description}
\item[Personality Trait.]
I see the value in those willing to risk their lives in the ruins. I have great insight and trust in my own judgement of character.
\item[Ideal.]
The right path is often the harder path, but it is always one worth walking.
\item[Bond.]
I was adopted by Elias Drexel, who raised and trained me. I don't want to let him down.
\item[Flaw.]
I am too open and trusting of people who claim to be on the same side as me, and I put a lot of weight in counting on people who I assume are able to handle themselves.
\end{description}

\subsubsection{Ansom Lang}
Captain Ansom Lang is a human \textbf{urban ranger}. He is an athletic teenager with long black hair and a short neat beard, angular jaw, and rugged features. He is garbed in a hooded green cloak and light chain armour of the Hooded Lanterns that prominently displays a captain's insignia on his shoulder.


\begin{description}
\item[Personality Trait.]
I have been in Drakkenheim long enough to know not to trust anyone. I am extremely wary of adventurers, and anyone who might cause problems for the Lanterns.
\item[Ideal.]
I believe that if my sister and I work hard enough, with the right people, that maybe we can change things.
\item[Bond.]
Myself and Petra have had each other's backs since we were children. That will never change.
\item[Flaw.]
I find reasons to not like people, and find reasons to doubt the commitment of the people around me.
\end{description}

\subparagraph{Royal Bastards}
Ansom and Petra Lang were adopted by Elias Drexel during the Civil War. At your option, you may determine they are secretly the illegitimate offspring of King Ulrich von Kessel IV and an unknown mother. This fact was concealed to prevent a scandal, and even Elias Drexel doesn't know it. A letter revealing this could be found in the Steward's Archive in Castle Drakken.

\subsubsection{Company Roster}
The Hooded Lanterns are organized into a regiment of five companies, making up a total force of approximately one thousand soldiers plus their support staff, orderlies, and crew. Each company is led by a captain and two lieutenants, and each has a distinct role:

\subparagraph{Scout Company}
Intrepid infiltrators who undertake dangerous reconnaissance operations into the city centre. Led by Captain Raine Highlash.

\subparagraph{Garrison Company}
Protects the established holdings of the Hooded Lanterns such as the Watchtower, Shepherd's Gate, and the Drakkenheim Garrison. Led by Captain Jakob Slovak, the Garrison company operates at double strength, since these soldiers must be constantly rotated to prevent them from becoming contaminated while on duty.

\subparagraph{Vanguard Company}
These experienced shock troops carry out offensive operations against priority targets, led by Captain Ansom Lang.

\subparagraph{Reserve Company}
A reserve force maintained outside the city composed of new recruits and aging veterans charged with maintaining supply routes.

\subparagraph{Royal Guard}
A special unit activated only if the Hooded Lanterns discover a royal heir. This small unit is composed of veterans hand-picked by the Lord Commander for trustworthiness and skill.

\subsubsection{Hooded Lanterns Soldiers}
Below are several NPCs to fill out the various Hooded Lanterns guards and scouts the player characters may encounter in Drakkenheim:


\begin{description}
\item[Chud Hopkins.]
A halfling \textbf{scout} who occasionally helps with patrols or watches, or works as a chef in the barracks kitchen.
\item[Sten Livingston.]
A human \textbf{scout} often positioned out in front of the City Watch Barracks as one of the head watchmen.
\item[Raine Highlash.]
A human \textbf{urban ranger} who runs patrols around the Outer City.
\item[Jakob Slovak.]
A human \textbf{urban ranger} of Shepherd's Gate.
\item[Luther Hess.]
A rotund \textbf{guard}, lazy and unwilling to do above the bare minimum.
\item[Alena Kruger.]
A calm human \textbf{veteran} who prides herself in her observational skills.
\item[Tresa Metterling.]
An optimistic \textbf{guard} with a hopeful outlook for Drakkenheim.
\item["Old" Symon Vasilisk.]
An aging human \textbf{guard} who rarely leaves the barracks and is on his way to retirement.
\item[Rolf Wagner.]
A young man of \textbf{noble} birth, sent to the city on his father's behalf to fight for the Lanterns.
\item[Greta Braun.]
A young \textbf{guard} with aspirations of climbing the ranks of the Lanterns.
\item[Alycia Martell.]
Afraid of almost everything in Drakkenheim but determined to fight for her beliefs.
\end{description}

\subsection{Missions}
\subparagraph{Missing Lantern}
Petra Lang went missing on a recent supply run. Her brother, Ansom, has resolved to rescue her from the Rat's Nest Tavern. (See chapter 6)

\subparagraph{Potions Inspectors}
The recent shipment of health potions from \textbf{Oscar Yoren} never showed up. The Hooded Lanterns need someone to look into this matter and make sure Oscar is still upholding his end of their deal. (See chapter 6)

\subparagraph{Dwarven Prospects}
A group of intrepid dwarves has a delerium mining operation in the Smithy on the Scar, and have brought incredible equipment with them. The Hooded Lanterns want to make a deal with the dwarves to secure their equipment and a share of the delerium. (See chapter 6)

\subparagraph{Resupply Mission}
The Hooded Lanterns dispatch the characters into the Inner City to set up and resupply several hidden caches and safehouses.

\subparagraph{Time for a new outpost}
The Cosmological Clocktower has long been an important objective. Clear it out so the Hooded Lanterns can establish it as a forward base at the centre of the city. (See chapter 7)

\subparagraph{Sizing Up the Silver Order}
Elias Drexel commissions the characters to approach the Silver Order as a neutral third party. He wants information: why they're here and what their plans are, and if they are open to an alliance or truce.

\subparagraph{Hazy Solutions}
The Hooded Lanterns don't possess a reliable defence against the Haze or a remedy to contamination. The Hooded Lanterns will need to secure the assistance of another faction to develop a means of dealing with the Haze: either the Amethyst Academy, the Followers of the Falling Fire, Knights of the Silver Order, or perhaps even the malfeasant wizards rumoured to dwell in the city. (See chapter 7)

\subparagraph{Noble Backer}
A wealthy \textbf{noble} who supports the Hooded Lanterns, used to own the Kleinberg Estate. They demand the Hooded Lanterns obtain a family heirloom and personal documents left behind in the estate vault. Otherwise, they'll pull their crucially needed funding.

\subparagraph{Audience with the Queen of Thieves}
The Hooded Lanterns believe the \textbf{Queen of Thieves} has information and spies within the regiment, which have made it difficult for them to plant their own agent inside the Queen's Men. They commission the adventurers to try to infiltrate the Queen's Men, get close to the Queen of Thieves, figure out her plans, and find a way to bring her down.

\subparagraph{Capture the Cathedral}
The Hooded Lanterns ask characters to coordinate an assault on the cathedral to slay the \textbf{Lord of the Feast} and open the cathedral vaults. This will hopefully have crucial evidence for locating a royal heir. (See chapter 7)

\subparagraph{Expedition to Castle Drakken}
Once the Hooded Lanterns have assembled the Seals of Drakkenheim and found a new heir, they plan an attack on Castle Drakken to claim the castle, crown the new monarch, and use Minazorond to defend the city!

\includegraphics[width=0.4\textwidth]{images/../5e.tools/img/adventure/DoDk/023-02-018.shipment-box.png}
\subsection{Boons}
\subparagraph{Supply Caches}
The Hooded Lanterns have hidden stockpiles of supplies in the city for emergencies. They'll teach characters how to find them quickly. Each contains 1d6 Potion of Healing or a Keoghtom's Ointment, Healer's Kit, ammunition, specially preserved rations, and the location is usually a safe place to take a short rest.

\subparagraph{Greencloaks}
The Hooded Lanterns offer each character either a \textit{cloak of elvenkind} or a \textit{cloak of protection} bearing the insignia of the Hooded Lanterns.

\subparagraph{Armoury}
As a reward for completing an important mission, the Hooded Lanterns grant the party one magical weapon from their armoury. This will be an uncommon magic weapon for an Outer City mission, a rare magic weapon for the Inner City or a faction stronghold, or a very rare magic weapon for a successful mission to Castle Drakken or the cathedral. The Hooded Lanterns will grant this boon once per player character.

\subparagraph{Call for Aid}
The Hooded Lanterns give the characters a flare gun and a horn that can be used as an action to summon emergency help. When used, help arrives 10 (3d6) minutes later in the form of a randomly determined Hooded Lanterns strike team. Player characters who call for help recklessly or without appropriate need lose this boon.

\subparagraph{Intelligence Reports}
No one knows the layout and the defences of Drakkenheim better than the Hooded Lanterns.

If asked, the Hooded Lanterns can provide detailed descriptions, directions, and information on the prominent structures in the city, maps of the city streets and public buildings, as well as information about the movements of known monster groups and other factions in Drakkenheim.

\subparagraph{Urban Survival Training}
Characters can spend one month training with veteran Hooded Lanterns soldiers to gain one of the following:


\begin{description}
\item[Ability Score Improvement.]
The character can increase one ability score by 2 or increase two ability scores by 1 each. The ability score can be increased above 20, up to a maximum of 30.
\item[New Feat.]
The character gains a new feat of their choice (subject to GM approval).
\end{description}

\subsection{Strike Teams}
The Hooded Lanterns are urban guerilla fighters honed by training, discipline, and experience. Their combat doctrine emphasizes organized teamwork, scouting ahead, stealth ambushes, hit-and-run tactics, seizing elevated ground, and heavy use of ranged attacks and traps. Once engaged, squads communicate using whistles, code-signs, signal flares, and horn calls to execute complex and practiced maneuvers. While the Hooded Lanterns don't have access to much magic, they're prepared to fight spellcasters in battle. They know much about the monsters they fight on a regular basis, and they always confirm their kills. The Hooded Lanterns know the layout of the city streets well, and stash caches of healing potions and emergency supplies throughout the city.

\subsubsection{Supply Convoys}
These supply convoys are often encountered in Emberwood Village and the northwestern side of the outer city, usually bound for one of the Hooded Lanterns' strongholds. Each consists of 10 (3d6) guards, 7 (2d6) scouts, 2 (1d4) veterans, and 5 (2d4) mules.

\subsubsection{Scout Squad}
A typical Hooded Lanterns unit consists of 10 (3d6) scouts led by an \textbf{urban ranger}. Smaller teams act as dispatch riders on warhorses.

\subsubsection{Garrison Squad}
10 guards led by a \textbf{veteran}. Occasionally, the garrisons have a \textbf{chaplain}, either a \textbf{priest} or a \textbf{druid}. Such units may also have brought traps, field defences, or even crew siege equipment.

\subsubsection{Vanguard Commando Squad}
A Hooded Lanterns commando squad consists of 5 (2d4) urban rangers led by an \textbf{assassin}. These units employ the full playbook of Hooded Lanterns combat tactics.

\subsubsection{Royal Guard}
A group of 10 (2d6) veterans led by a \textbf{gladiator}.

\subsection{Schemes}
\includegraphics[width=0.4\textwidth]{images/../5e.tools/img/adventure/DoDk/024-02-019.apothecary.png}
\subparagraph{Search and Seizure}
A squad of Hooded Lanterns confront the characters and demand the characters turn over their equipment and vacate the city immediately. If they are seen again, they will be brought to justice.

\subparagraph{Blockade}
The Hooded Lanterns forbid characters from entering the city via Shepherd's Gate. They step up patrols around the other gates to block characters from using them.

\subparagraph{Stalkers}
A Hooded Lanterns strike team tails the characters as they move through the streets. They assault the characters after they are weakened by a difficult encounter and attempt to arrest them.

\subparagraph{Wanted}
A series of wanted posters are put up around Emberwood Village with a reward for bringing the characters to justice. This draws the attention of Queen's Men bounty hunters as well as rival adventurers, and anyone else keen to make money to bring the characters to justice.

\subparagraph{Deadly Ambush}
The Hooded Lanterns set up extensive traps down a city street and lure the characters into a deadly ambush where two strike teams await on rooftops and in covered windows to surprise the characters. They attempt to encircle the characters in a pincer maneuver.

\subparagraph{Crown Authority}
The Hooded Lanterns banish the characters from Emberwood Village: local merchants refuse to do business with the characters as a result, and the characters won't be allowed to stay at any of the inns or taverns. The Hooded Lanterns place additional patrols in Emberwood Village at the watch tower and Caravan Court to drive out the characters when they come into town.

\subparagraph{Strange Bedfellows}
In desperate situations, the Hooded Lanterns ally themselves with another faction who also has made an enemy of the player characters.

\subparagraph{Execution Sentence}
The Hooded Lanterns orchestrate a deadly plan to draw the characters into battle against the Executioner of Slaughterstone Square. They blockade the roads and take positions on the roofs, preventing the characters from escaping the square while the Executioner kills them.

\section{Knights of the Silver Order}
\DndQuote%
{}
{''These evil crystals are the stuff of dark magic; anathema to order and life itself. No goodness can come from their corrupt power, and those foolish enough to believe delerium can be turned to useful purposes are truly lost.''}
{Knight-Captain Theodore Marshal}
The Silver Order is an \textbf{elite church-sanctioned regiment} of paladins and ordained knights oath-sworn to combat supernatural evils, otherworldly incursions, and dark magic. As faithful devotees of the Sacred Flame, they believe delerium is a dangerous and contaminating blight. Led by the valiant \textbf{Knight-Captain Theodore Marshal}, the knights have been dispatched to eradicate the corruption at its very source.

\subsection{Background}
\includegraphics[width=0.4\textwidth]{images/../5e.tools/img/adventure/DoDk/025-02-020.silver-order-symbol.png}
The Silver Order was founded several centuries ago as a militant arm of the Faith of the Sacred Flame. These knights and paladins act by the decree of the Divine Matriarch to defend the faithful from demons and monsters, and are called up as witch-hunters when wayward mages dabble in forbidden arcana. Their holy missions transcend sovereign borders, but the order is headquartered in the Holy City of Lumen in Elyria.

Paladins are revered figures in the Faith of the Sacred Flame, as the founder of the faith was herself a paladin. Though not every warrior in the Silver Order is a paladin, each strives to emulate Saint Tarna's valiant example. Though they usually operate as individual questing knights or in small warrior-companies, occasionally they are deployed as a regiment. Members of the Silver Order have been traditionally drawn from all peoples and all levels of society. Nobles and commoners are found within the Order, and one's station in life does not reflect one's position in the Order.

In their role as witch-hunters and mage-slayers, the Silver Order was complicit in the persecution and oppression of arcane spellcasters. To this day, the issue remains a deep stain on the order's honour, as many observed that the injustices and persecution visited upon spellcasters stand in stark contrast against the tenets of the faith of the Sacred Flame, a faith that calls out for justice for the oppressed. Nevertheless, staunch supporters of the Silver Order argue that their deeds were indeed justified. Decisive action is necessary to defend the innocent and the faithful from wicked sorcerers who sought to dominate others with magical might, or cut down foul warlocks who spread wickedness through demonology and necromancy. The chronicle of the Silver Order records countless heroic missions that have prevented a dire threat from growing into a far greater crisis.

\includegraphics[width=0.4\textwidth]{images/../5e.tools/img/adventure/DoDk/026-02-021.theodore-marshal.png}
The Silver Order answers to the Faith of the Sacred Flame, not to the command of any sovereign nation. However, the Order remains strongly connected to its roots, and most of its members hail from Elyria. Ever since the realm of Elyriacame under the de facto rule of the religious ministry, the order has become increasingly entangled in state affairs. Their knights have been deployed on inquisitional crusades against heretics, heathens, and other political enemies of the Sacred Flame. The distinction between an "enemy of the Sacred Flame" and an "enemy of Elyria" is not entirely clear anymore. As such, some worry the Order's presence in Drakkenheim is an invasive ploy by Elyria to wrest control of Westemär, rather than a righteous expedition to purify the city of eldritch contamination.

\subsubsection{Attire}
As befits their name, the knights are clad in silver plate with gold filigree. They wear flowing tabards of white, prominently decorated with symbols of the Sacred Flame, rather than any individual heraldry. The Flamekeepers who accompany the knights wear religious vestments of white and gold.

\subsubsection{Strongholds}
The Knights of the Silver Order have built a military encampment known as Camp Dawn about two miles outside Drakkenheim, just on the outside edge of the Haze. Outside Drakkenheim, the Silver Order can often find shelter in any chapel or cathedral of the Sacred Flame, though not every Flamekeeper in Westemär now supports their mission. Their headquarters is the impregnable Vigilant Keep in the far city of Lumen in Elyria, where the Grandmaster oversees the entire order at the behest of the Divine Matriarch.

\section{Roleplaying Traits}
\subparagraph{Symbol}
A burning chalice upon a shield. Their colours are silver and gold.

\subsection{Personality Traits}

\begin{itemize}
\item We conduct ourselves with honour and respect. We do not speak falsehoods, and we show compassion and courteousness for all, even our enemies.
\item We bravely stand up for what is right, and expect others to do the same.
\item We profess our faith openly, so all may bask in the warmth of the Sacred Flame.
\end{itemize}

\subsection{Ideals}

\begin{itemize}
\item We are united by faith in the Sacred Flame. All are called to uphold its tenets, and all must heed its wisdom. We cannot be divided by ancestry, nobility, or national creed.
\item We will never abandon the righteous to torment and death. Not while we have strength left to act.
\item Elyria has become a mighty realm through its faith in the Sacred Flame. It is a shining example that all other monarchs and rulers should emulate.
\end{itemize}

\subsection{Bonds}

\begin{itemize}
\item We are oath-bound to protect those who cannot protect themselves from supernatural evil, cosmic terrors, and malevolent sorcerers. We are the shield of the faithful and the torch against the darkness that comes from beyond. Evil must be uprooted before it can spread.
\item We are called to act by the Divine Matriarch of the Sacred Flame, and will always follow her decree.
\item It is our duty to protect the relics and holy places of our faith, and ensure they do not fall into wicked hands who would use them for vile purposes.
\end{itemize}

\subsection{Flaws}

\begin{itemize}
\item We never back down from a fight. If you challenge us, we will never surrender.
\item Those who do not follow the Faith of the Sacred Flame are misguided, especially those who worship other gods or twist the message of our faith. They have been turned away from the light through deception and falsehoods spread by charlatans.
\item Power corrupts, and nothing corrupts like arcane power. We don't trust arcane spellcasters, and we are prepared to strike against them should they work evil magic or traffic with fiends.
\end{itemize}

\DndQuote%
{}
{''Metal headed, stubborn, and overly committed to falsities justified by holy laws. If I praised the Light whenever I blew people up, would I win medals?''}
{Sebastian Crowe}
\subsection{Objectives}
\subsubsection{Destroy delerium at its source}
The Divine Matriarch of the Sacred Flame has decreed delerium an unholy abomination. These chaotic seeds of evil magic must be wholly destroyed before they can take root. By her command, the Silver Order has come to cleanse Drakkenheim of this foul taint: even if it means burning the city to the ground. Furthermore, the knights have resolved to stop others from harvesting and using delerium, especially the Amethyst Academy.

The Silver Order insists that Drakkenheim - and by extension, the nation of Westemär - cannot be rebuilt until the threat of delerium is entirely eliminated. While the Silver Order does not intend to establish Elyrian rule over Westemär, they want to ensure any new ruler honours the Sacred Flame, commits to destroying delerium, and helps them dissolve the Followers of the Falling Fire.

\subsubsection{Stamp out the heresy of the Falling Fire}
The demagogue \textbf{Lucretia Mathias} has fractured the Faith of the Sacred Flame, and twisted its noble and pure religious teachings into the debased worship of a blasphemous and evil rock from the darkest depths of the outer void. Lucretia Mathias must be captured and brought back to Elyria to stand trial for her grave sins, where she may find repentance in a burning pyre. Her misguided followers must be driven from the city.

\subsubsection{Preserve the holy relics in Drakkenheim}
Drakkenheim was home to many saints of the Sacred Flame: most notably, the heroic Saint Vitruvio, who fought astride the mighty gold dragon Argonath centuries ago. Both were buried beneath the cathedral that bears his name. Many relics of this mighty paladin are lost in Drakkenheim, such as \textit{Ignacious, the Sword of Burning Truth}. The Silver Order wishes to reclaim these artefacts and reconsecrate the altar of the Sacred Flame in the Cathedral of Saint Vitruvio, so it may serve as a beacon of light from which the rest of the city may be purged.

It is of vital importance to the Silver Order that these relics and places do not fall into the hands of the Followers of the Falling Fire, who they fear may use them for heretical purposes.

One of the relics of Saint Vitruvio is also one of the \textit{Seals of Drakkenheim}: the Saint Vitruvio's Phylactery. Traditionally, this item is entrusted to the High Flamekeeper of Drakkenheim. If the Silver Order obtains it, they can ensure the Faith of the Sacred Flame has some influence over the political destiny of Drakkenheim and Westemär by demanding who is appointed the new High Flamekeeper.

\subsection{Key Figures}
\textbf{Knight-Captain Theodore Marshal} is faction leader of the Silver Order in Drakkenheim. He keeps close counsel with High Flamekeeper Ophelia Reed, the faction lieutenant.

\subsubsection{Knight-Captain Theodore Marshal}
Theodore Marshal is a man in his prime. He is athletic and tall with boldly sculpted features, bronzed skin, and unkempt sandy brown hair. Unshaven stubble has grown into a beard over the criss-cross of scars on his face. A leather patch covers his right eye. He wears gleaming silver plate armour etched with golden trim, and carries a shimmering sword and a broad shield engraved with the symbol of the Sacred Flame.


\begin{description}
\item[Personality Trait.]
I act with determined calm and respect. Little rouses my anger, instead I respond with disappointment when others incite my ire.
\item[Ideal.]
I believe that good-hearted people can set aside their differences to work together.
\item[Bond.]
I am willing to lay down my life for a just and righteous cause.
\item[Flaw.]
I will never give up on my comrades. I do not trade lives.
\end{description}

\subsubsection{Ophelia Reed}
High Flamekeeper Ophelia Reed is a human \textbf{chaplain}. She is a stout black woman in her late forties with a warm smile, rounded features, and braided hair. She wears the flowing yellow and white robes of a cleric of the Sacred Flame. She carries a silvered staff and a censer with a \textit{continual flame} spell cast upon it.


\begin{description}
\item[Personality Trait.]
I attentively listen to others, and respond from a place of warmth and kindness.
\item[Ideal.]
I will always seek the right remedy for any ailment or problem, and offer my patient guidance to those who ask for it.
\item[Bond.]
I have taken a vow of poverty and devoted my life to the service of the Sacred Flame.
\item[Flaw.]
Darkness is all-consuming. Once it has taken root. For some, death is the only salvation.
\end{description}

\subsubsection{Silver Order Crusaders}
The Knights of the Silver Order have marched the long road from distant Elyria. They have mustered a great strength of the Order, consisting of a hundred knights and their retinues, men-at-arms, and retainers. Dozens of clerics of the Sacred Flame have joined their forces as well.

Below are several concepts for quick NPCs to fill out the ranks of the Silver Order knights, paladins, and clerics.


\begin{description}
\item[Knight-Lieutenant Cassandra Wyatt.]
a human \textbf{cavalier} who leads the aerial units.
\item[Ser Virgil Underwood.]
a human \textbf{knight}.
\item[Austin Edwards.]
A young \textbf{knight}, eager to leap into battle for his cause.
\item[Adam.]
A stoic and thoughtful paladin with a cautious attitude towards Drakkenheim.
\item[Claudia.]
A determined \textbf{knight} with the qualities of great leadership.
\item[Delilah.]
An aspiring Flamekeeper with a fiery attitude and strong will.
\item[Quinn.]
A new recruit, optimistic and eager to impress her team.
\item[Elijah.]
A paladin who takes great pleasure in defeating monsters in combat.
\item[Buddy Knox.]
A guitar strumming bard who sings the deeds of the Knights of the Silver Order.
\item[Noah.]
An aging man with a big white beard who takes things in stride.
\item[Ruth.]
a Flamekeeper who takes her job as a healer very seriously.
\item[Zacharias.]
An arrogant \textbf{knight} who fears no monster or abomination in Drakkenheim.
\end{description}

\subsection{Missions}
\subparagraph{Sceptre of Saint Vitruvio}
Reclaim the lost relic in the chapel of Saint Brenna. (See chapter 6)

\subparagraph{Dwarven Heathens}
Drive out or destroy the dwarven smithy. (See chapter 6)

\subparagraph{Witch Hunt}
Destroy one of the malfeasant wizards in Drakkenheim: \textbf{Oscar Yoren}, the \textbf{Pale Man}, or the Iron Banshee.

\subparagraph{Make an Alliance}
The Silver Order wishes to broker an alliance with the Hooded Lanterns if possible. However, we'll need to convince the Hooded Lanterns that it's in their best interests to burn down the city entirely, then rebuild.

\subparagraph{Battle of Temple Gate}
We need our own means to access the Inner City. We shall assault Temple Gate and wrest control of it from the garmyr.

\subparagraph{Saint Vitruvio's Cathedral}
Destroy the \textbf{Lord of the Feast} and retake the cathedral. Afterwards, help us relight the Brazier of the Sacred Flame.

\subparagraph{Root out the Heretic}
Capture \textbf{Lucretia Mathias}!

\subparagraph{Shattered Tower}
If the Amethyst Academy will not be dissuaded to give up their pursuit of delerium, then we will cut them off from it entirely. Invade the Inscrutable Tower and destroy it from within!

\subparagraph{The Resurrection of Argonath}
Assembling the relics of Saint Vitruvio, the Silver Order attempts to resurrect the gold dragon Argonath, but they'll need the Diamond of Mount Kadath from the vaults of Castle Drakken to complete the spell.

\subparagraph{Ashes to Ashes}
Burn Drakkenheim to the ground and destroy the Delerium Heart.

\includegraphics[width=0.4\textwidth]{images/../5e.tools/img/adventure/DoDk/027-02-022.cassandra-wyatt.png}
\subsection{Boons}
\subparagraph{Supplies}
The Silver Order gives the party 1d4 Potion of Healing or Spell Scroll (of a spell from the cleric spell list) of their choice whenever they begin a mission that directly advances the Silver Order's goals. These will be uncommon items for missions to the Outer City, rare items for the Inner City or faction strongholds, and very rare items for missions involving the crater, Castle Drakken, or the cathedral.

\subparagraph{Stables}
The Silver Order offers each player character a \textbf{warhorse}. If the horse is killed, the Silver Order offers the characters another one.

\subparagraph{Armoury}
As a reward for completing an important mission, the Silver Order will grant the party one magical melee weapon, a magical suit of heavy armour, or a magical shield from their armoury. This will be an uncommon item for an Outer City mission, a rare item for the Inner City or a faction stronghold, or a very rare magic item for a successful mission to Castle Drakken or the cathedral. This boon can be granted multiple times, once for each character in the party. The Dungeon Master can either choose the item, or allow the player to choose it.

\subparagraph{Divine Spellcasting Services}
High Flamekeeper Ophelia Reed can offer her services as a spellcaster to the characters. She can help them recover via spells such as \textit{lesser restoration}, \textit{prayer of healing}, \textit{purge contamination}, \textit{remove curse}, and \textit{greater restoration}. While characters must still provide any costly material components, she does not charge them to cast the spells.

\subparagraph{Religious Lore}
The Faith of the Sacred Flame believes knowledge is a form of illumination, and keeps historical texts and religious lore. Flamekeepers can answer questions about religious practices, magical creatures such as celestials, fiends, and undead, historical facts about the past five hundred years, and mythological stories about events that happened up to one thousand years ago. If you feel the matter is something that might require some deeper research, it takes the Flamekeepers one week (or 2d6 days) to unearth relevant information.

\subparagraph{Take to the Sky}
The Silver Order grants the player characters access to their griffons and training on how to fly them. Each character gains a griffon mount. If this mount is killed it cannot be replaced.

\subparagraph{Revive the Fallen}
Ophelia Reed will cast \textit{raise dead} on a fallen player character of 9th level or higher who falls in battle fulfiling objectives of the Silver Order. They can provide the material components for the spell. The Silver Order will offer this boon once per player character. However, player characters of any level who are trusted by the Silver Order, provide the material component for \textit{raise dead}, and make an offering of 1,000 gold pieces to the Sacred Flame may also receive this boon.

\subparagraph{Teachings of the Faith}
The Silver Order allows characters to read from relic holy books of the Sacred Flame. They present them with a \textit{Manual of Gainful Exercise}, a \textit{Tome of Understanding}, or a \textit{Tome of Leadership and Influence}.

\subsection{Strike Teams}
The Silver Order favours aggressive charges and defensive bastions, employing a "hammer and anvil" approach whether on the battlefield or in skirmishes. Their inspiring commanders drive their rank-and-file soldiers to greater acts of boldness and heroism, while the magical support of their clergy means they can tend the wounded and unravel evil magic levied against them. The Silver Order have superlative training in combating spellcasters and fiends, and know how to recognize the magic and tactics of such foes. They target spellcasters aggressively, and seek to break their concentration or end their spells using \textit{dispel magic}.

Before a battle, the Silver Order draws up a straightforward battle plan, which they execute with discipline and focus. In all matters, they operate with honor. They will often reach out in parley before engaging, and even in battle they will often implore their foes to retreat or surrender once they seize the upper hand. On the contrary, members of the Silver Order expect no quarter themselves, and offer none to the wicked and unrepentant.

\subsubsection{Questing Knight}
\includegraphics[width=0.4\textwidth]{images/../5e.tools/img/adventure/DoDk/028-02-023.ophelia-reed.png}
A lone Knight of the Silver Order often carries out missions accompanied by a personal retinue of 10 (3d6) guards. This \textbf{knight} is often mounted on a \textbf{warhorse}. Knights are occasionally accompanied by a smaller number of retainers who are also on horseback.

\subsubsection{Silver Order Cavalry}
This company of heroic knights consists of a group of 5 (2d4) knights all mounted on or warhorses. Such units are mustered for offensive operations on when force is certainly expected.

\subsubsection{Silver Order Infantry}
The knightly retainers sometimes form large infantry units equipped with spears and shields, pikes, or halberds: such units may include 35 (10d6) guards, perhaps more. They may be commanded by a single \textbf{knight}. The Silver Order deploys such units to hold ground while their cavalry are charged with seizing it.

\subsubsection{Clergy}
Any Silver Order strike team may include support from a cleric of the Sacred Flame: either a \textbf{priest} or \textbf{chaplain}. They may be accompanied by 3 (1d6) acolytes. Sometimes several Flamekeepers gather together on their own for special circumstances when a great deal of divine magic is needed, or to perform religious rites.

\subsubsection{Sky-Cavaliers}
The elite air force of the Silver Order are units of 3 (1d6) cavaliers mounted on griffons. These riders strike like heavenly thunderbolts against ferocious monsters, flying demons, and other airborne threats, or descend from the sky upon wretched spellcasters and enemy leaders.

\subsection{Schemes}
As an honorable military regiment, the Silver Order does not engage in subtle or stealthy schemes. Instead, they enact traditional military strategies with commensurate discipline and effectiveness.

\subparagraph{Loyal to the Flame}
The Silver Order petitions Flamekeeper Hanna in Emberwood Village to withhold her spellcasting services from the player characters.

\subparagraph{Military Blockade}
The Silver Order creates fortified outposts on the main roads leading into Drakkenheim. Each consists of a fortified watchtower and a palisade wall, and is garrisoned by 2 (1d4) knights, 10 (3d6) guards, and a \textbf{priest}. They have warhorses and ballistas. These blockades prevent outside supplies from reaching Emberwood Village or other strongholds within the city - most crucially, food.

\subparagraph{Occupy Emberwood Village}
The Silver Order mobilizes in force to take Emberwood Village. Unless the Silver Order is stopped, they drive out the other factions entirely and install a strict curfew on the villagers, effectively ending all trade and services within the town. The knights fortify the town and relocate their base of operations here, retaining a small garrison at Camp Dawn.

\subparagraph{An Honourable Challenge}
\textbf{Knight-Captain Theodore Marshal} challenges the most martially skilled player character to a one-on-one duel to the death. This honorable combat must be conducted without any external interference. If he wins, he immediately challenges another character to a duel under the same conditions. He repeats this process until either himself or the entire party is killed.

\includegraphics[width=0.4\textwidth]{images/../5e.tools/img/adventure/DoDk/029-02-024.sacred-flame-commands.png}
\section{Queen's Men}
\DndQuote%
{}
{''Drakkenheim is mine. Nothing moves in Drakkenheim without me allowing it to happen. I have eyes on every street, a fix on every prize, and tabs on every would-be band of fortune-seekers. Only one can rule this city—me.''}
{The Queen of Thieves}
The Queen's Men are a loose affiliation between a hundred gangs of brigands, outlaws, and scoundrels, all who swear fealty to the enigmatic \textbf{Queen of Thieves}. These reprobates prey upon pilgrims, prospectors, and explorers drawn to the city, extort and rob adventurers, smuggle delerium to disreputable clients in distant lands, and plunder the fantastic treasures and incredible wealth left behind in the city.

\subsection{Background}
\includegraphics[width=0.4\textwidth]{images/../5e.tools/img/adventure/DoDk/030-02-025.queen-of-thieves.png}
Drakkenheim is a haven for criminals, deserters, and fugitives. The ever-present dangers of the dark city make it the perfect place to hide for someone on the wrong side of the law with nothing left to lose. Opportunities for plunder abound within this new frontier for banditry. Many more ply their trade raiding the supply caravans making their way to Emberwood Village. Sly cutthroats know that simply taking the spoils from a weakened but successful adventuring party at swordpoint makes for a highly profitable enterprise, and a far less risky one to boot!

When the Hooded Lanterns arrived in Drakkenheim a few years ago, they resolved to restore order. Presenting a united front against the fractious outlaws, the Hooded Lanterns met with early success in driving out the more violent and bloodthirsty among them. Then, an enigmatic figure emerged. This self-proclaimed "Queen of Thieves" used assassination, trickery, intimidation, and guile to reshape the disparate gangs into a budding criminal empire. Her plots have sabotaged many important Hooded Lanterns operations, and the Lord Commander believes the Queen of Thieves blackmailed several of his important backers into withdrawing their support.

\includegraphics[width=0.4\textwidth]{images/../5e.tools/img/adventure/DoDk/031-02-026.queen-of-thieves.png}
More importantly, the Queen of Thieves has established a general truce between the various bandit leaders occupying the city to share territory, resources, and contacts. In exchange, every gang must pay a regular tribute in gold and plunder to the Queen of Thieves. Those who don't become a target for all the others, while loyal and successful scoundrels enjoy extra benefits and kickbacks. Though the Queen of Thieves is now infamous for her schemes and skill, her web of black market smugglers secures her rule over the underworld. Her agents funnel treasures and delerium purloined from Drakkenheim to skilled fences and unscrupulous merchants in the world beyond. This ensures ample gold and booze flow back into Drakkenheim to keep her army of scoundrels happy and well-supplied.

\subsubsection{Strongholds}
\subparagraph{The Skull \& Sword Taphouse}
The Skull \& Sword Taphouse in Emberwood Village is a favoured haunt for many outlaws, and the best place to find Blackjack Mel. This location is detailed in Chapter 4: Emberwood Village

\subparagraph{Buckledown Row}
In Drakkenheim itself, the gangs often meet in several run-down pubs along this dingy street. The sewers beneath connect to a series of smugglers' tunnels that lead to the underground stronghold where the Queen of Thieves holds court. Though the way is littered with deadly traps, it leads nowhere—a cruel joke the Queen of Thieves plays upon would-be hunters.

\subparagraph{The Court of Thieves}
Only the Queen of Thieves' most trusted agents know the true location of her personal lair. See the Faction Strongholds chapter for details.

\section{Roleplaying Characteristics}
\subparagraph{Symbol}
The Queen of Thieves has taken the Queen of Diamonds as her personal calling card. Her colour is red. The individual gangs wear their own colours and symbols.

\subsection{Personality Traits}

\begin{itemize}
\item We're willing to lie, cheat, steal, threaten, and even kill to get what we want.
\item You can't trust anyone, but you can trust their greed Everyone's got a price.
\item We get even whenever someone crosses us.
\end{itemize}

\subsection{Ideals}

\begin{itemize}
\item We're outcasts and outlaws. The laws written by nobles and priests turned us into criminals. We're glad the old world is falling apart. It never did right by us anyways.
\item Anything can be yours if you're strong enough or smart enough to take it.
\item Always stay one step ahead of the game. There's no heroes here. If you won't do what it takes to survive, you're as good as dead.
\end{itemize}

\subsection{Bonds}

\begin{itemize}
\item You came to the wrong neighbourhood, friend. We rule this city now!
\item We must pay tribute to the Queen of Thieves and respect the general truce between all the outlaws in Drakkenheim. Everyone else is fair game!
\item A favour, a secret, or a good contact can be just as valuable as money to the right person.
\end{itemize}

\subsection{Flaws}

\begin{itemize}
\item Take any chance you get to indulge your vices. After all, it might be the last time you get to enjoy it!
\item Always look out for yourself. There's no loyalty amongst thieves. When the chips are down, it's everyone for themselves.
\item Whatever it is, wherever it came from, however we get it, whoever wants it - we don't care. If someone wants to pay for it, that means money for us!
\end{itemize}

\DndQuote%
{}
{''You won't know she is there unless she wants you to. She is not to be trusted and never underestimated. The most dangerous royalty I've ever met.''}
{Pluto Jackson}
\subsection{Objectives}
The Queen of Thieves dreams of building an influential and powerful criminal empire forged in the anarchy and lawlessness of the ruins, while conspiring to sabotage the other factions' operations at every turn. However, she does have her eyes on one special prize.

\subsubsection{Transform Drakkenheim into an Outlaw's Haven}
The Queen's Men see no reason for Drakkenheim to change. As it is, the city is a bandit's paradise. They scour the ruins for treasure, delerium, and lost magic items to sell to the highest bidder on the black market. Most outlaws don't have any loftier goals, but the Queen of Thieves wants to expand her criminal empire.

However, other factions are trying to gain ground in Drakkenheim. Driving them off isn't a simple matter, either. Several of these new rivals are also potential customers: the Amethyst Academy pays top dollar for delerium. The pilgrims flooding into the city to join the Followers of the Falling Fire are easy pickings for shakedown and robbery as well.

Others pose a more existential risk to the Queen of Thieves. If she pushes too hard against the Hooded Lanterns or the Silver Order, a much more aggressive military crackdown could ensue. To preserve the delicate balance, the Queen of Thieves has resolved to deceive, blackmail, exploit, and outwit the other factions, but not to destroy them—that would be bad for business! Therefore, she aims to disrupt and sabotage the other factions' goals. Perhaps she might even install a puppet on the throne of Westemär!

\subsubsection{Obtain the Crown of Westemär}
\DndQuote%
{}
{''Stealing it would be the crowning achievement for a Queen of Thieves, don't you think?''}
{Sebastian Crowe}
This objective is only known to the most trusted associates of the Queen of Thieves. Most Queen's Men don't know they are working towards this goal. The \textit{Crown of Westemär} is more than just a symbol of royal authority: it is a powerful artefact in its own right that may grant several \textit{Wish} to whomever wears it. There is no greater prize to the Queen of Thieves, and she wants to steal the crown so she may use its power for herself. While the Queen of Thieves has personal desires for this power, the Crown is also an effective trump card in securing control of Drakkenheim. Having the ability to make wishes is a massive deterrent for anyone who would move against her, cementing her power over the ruined city. All reliable signs indicate the crown is lost in Castle Drakken, but that's not the only complication. In order to attune to the Crown of Westemär to use its powers, one must also possess all six \textit{Seals of Drakkenheim}.

\subsection{Key Figures}
The \textbf{Queen of Thieves} is the criminal overlord of Drakkenheim. She doesn't trust any of her underlings as far as she could throw them, but Blackjack Mel has proven a capable toadie and serves as faction lieutenant.

\subsubsection{The Queen of Thieves}
\DndQuote%
{}
{''Who I am is not as important as what I am. I am quicker than thought, more silent than contemplation, more slippery than a fading memory, and more unshakable than a guilty conscience. I deal in secrets and mendacity as lesser thieves deal in coins and larceny. I can fulfil your wildest dreams, or I can be your worst nightmare.''}
{The Queen of Thieves.}
The \textbf{Queen of Thieves} is a scheming mastermind who controls the Drakkenheim underworld. She has transformed a hundred squabbling bands of brigands and outlaws into a burgeoning criminal empire. However, none know her true identity or origin, for she wears many faces. Though she is most infamously encountered as a lithe masked woman garbed in swashbuckling attire with a wide-brimmed hat and jewelled rapier, the Queen of Thieves adopts a variety of costumes, facial features, skin and hair colours, body shapes, ages, and genders depending on her mood and circumstances. As such, she typically confirms her identity by revealing some secret about whoever she is speaking with or demonstrating her uncanny ability to change appearance seemingly at-will.


\begin{description}
\item[Personality Trait.]
I am never caught unprepared, and have contingencies for my contingencies. I speak with rakish confidence while casually exposing everything I know about others, and flaunt my ability to be anybody or anywhere.
\item[Ideal.]
I indulge in playing mind games and manipulating people with the promise of profit. Instead of spilling blood, I spill secrets. My enemies are worthless as corpses. When I have to steal lives, I spend them dearly first.
\item[Bond.]
Law is powerless in the chaos of the ruins. As long as anarchy reigns, so shall I.
\item[Flaw.]
Anyone who trusts me is a terrible fool. Double-crossing is my nature. I expect everyone to betray me, and I'll relish the opportunity to defeat them when they try.
\end{description}

Few amongst the Queen's Men have met the Queen of Thieves in person, and even fewer have attended the Court of Thieves. Most only know if they get a note attached to a Queen of Diamonds playing card, you better do what it says.

\subsubsection{Who is the Queen of Thieves?}
The true power of the Queen of Thieves is her apparent anonymity. Once her true identity is revealed, she is effectively defeated. In this regard, it's important to separate the "persona" that is the Queen of Thieves from the actual \textit{person} behind her. While this book includes one possible suggestion, you may decide the true identity of the Queen of Thieves in your campaign, and decide her true motivations for yourself. You might choose one of the missing heirs to the monarchy, or a long-lost friend or relative connected to one of your player characters. The Queen of Thieves does not disclose her identity readily, but foreshadowing various possibilities with uncanny remarks and subtle hints makes the final reveal that much more dramatic!

In the original \textit{Dungeons of Drakkenheim} campaign, the \textbf{Queen of Thieves} was none other than \textbf{Katarina von Kessel}, one of the lost scions of House von Kessel. Unfortunately for her, she is an arcane spellcaster. Under the Edicts of Lumen, she cannot inherit the throne.

An "incident" involving an upstart wizard resulted in her sister, Eliza, being banished to the Abyss. Unwilling to accept her sister's death, Katarina von Kessel fled the Amethyst Academy and travelled the world to find a way to bring her sister back. In the process, she learned to weave her proclivities for enchantment and illusion magic with her noble-born fascination for swordplay and scheming. Eventually, she hatched a plan to return to Drakkenheim under a false identity and claim her birthright. She does not want to be used as a pawn of the Amethyst Academy or the nobility. She wants to do things \textit{her way}, and does not believe such organizations have \textit{her} best interests in mind—nor the interests of the common folk, for that matter. Katarina von Kessel wants to obtain the \textit{Crown of Westemär} and use its power to restore her lost sister's life, then erase all public knowledge of her magical abilities so she may rule Westemär legitimately. She hasn't yet decided on her last wish, however.

\DndQuote%
{}
{''Sometimes it's the things you love that end up hurting you the most.''}
{Sebastian Crowe}
\includegraphics[width=0.4\textwidth]{images/../5e.tools/img/adventure/DoDk/032-02-027.convenient-trade.png}
\subsubsection{Blackjack Mel}
Blackjack Mel is a human \textbf{scoundrel}. He is a short and scrawny man in his late thirties with a funny looking but clean-shaven mug and greasy slicked-back hair. He wears a black padded jacket and buckled slacks that appear distinguished at first glance, but are actually quite shabby and worn.


\begin{description}
\item[Personality Trait.]
I can always make a deal with the right angle, I just gotta pitch it the right way. I don't give up when the mark doesn't take the bait, just try something new!
\item[Ideal.]
We're all in this for ourselves. You scratch my back, I'll scratch yours.
\item[Bond.]
Listen, I got some close connections with some powerful people. You wanna meet the Queen of Thieves? Stick with me. We'll go places together. Just uh, watch out you don't piss anyone off, okay?
\item[Flaw.]
I'm a compulsive liar, so what? Yeah, sometimes even I get confused by my own lies, but the truth can always do with a little exaggeration, it makes life more interesting!
\end{description}

\subsubsection{Outlaws and Scoundrels}
The Queen of Thieves upholds a general truce between a hundred individual gangs of brigands. Each outfit holds a claim to its own territory within the ruins that is their exclusive right to plunder. In return, each must pay a share to the Queen of Thieves. They must also heed her summons, which takes the form of a playing card showing the Queen of Diamonds.

She arbitrates disputes and enforces dissent in a ruthlessly efficient manner. Ultimately, the gangs are disposable pawns to the Queen of Thieves, and she \textit{doesn't care} if the player characters get into fights with them.

Below are a few of the most prominent gangs that work under the Queen of Thieves. Feel free to expand or use these as you see fit to help flesh out the myriad outlaws encountered in the city:

\subsubsection{Pins \& Clubs}
The Pins \& Clubs are general ruffians. They often have iconic colourful mohawks or other wild coloured hairstyles. They wear a lot of spikes and chains and notoriously break the legs of people who cross them.

\subparagraph{Leader}
Leon The Swine, a scrawny middle aged \textbf{bandit captain} with a scar over his left eye and a mohawk. He carries a saber and a dagger, and chains hang from his leather vest and black trousers.

\subparagraph{Symbol}
A club logo from a card deck with several pins sticking out of it.

\subparagraph{Quote}
"If ya don't like the way they looked at ya, if they seem too pompous, too arrogant, too wise, or too righteous, you club them down and let them know where they really stand."

\subsubsection{Rose \& Thorns}
This group is known for their confidence. They often are known for using daggers, spears, and shortswords and always have high quality leather armour and steel-toed spiked boots. A lot of the most well known members of this gang are a group of women renowned amongst the thieves as the Bouquet, led by the unwavering Rose Carver.

\subparagraph{Leader}
Rose Carver, a striking red-haired \textbf{veteran} with piercing eyes and studded leather armour. She has a belt filled with daggers, and steel-tipped boots with several spikes on the toes.

\subparagraph{Symbol}
A thorn covered rose.

\subparagraph{Quote}
"All beautiful things have their sharp edges. Let's show them ours!"

\subsubsection{The Wounded Hearts}
The Wounded Hearts use cunning, deceit, and charm to get what they want. They are masters of swindling, scheming, and double crossing. Nowhere else will you find better liars. They often wear elegant noble's garb and act stuck up, as if they deem themselves superior to the other gangs.

\subparagraph{Leader}
Christain Lament, a tall blonde elven \textbf{spy}. He has blue eyes and wears a flowing cape and black noble's clothing with gold filigree. He carries a rapier and daggers.

\subparagraph{Symbol}
A heart with a blade through it.

\subparagraph{Quote}
"Chivalry and cordiality will get you far in this world. Kindness is the great mask of all evils."

\subsubsection{The Howlin' Dogs}
These bandits are all about intimidation. They have size - usually only the strongest get into this gang. They often wear vests and very little else, letting their muscles speak for themselves. They often carry axes or greatswords. Many of them dual wield. They are rude, mean, and ferocious. If you see one looking at you, it's best to look away.

\subparagraph{Leader}
Baskerville, a large grey-skinned half-orc \textbf{gladiator} with muscles bulging out of his vest. He has several piercings and carries two axes. His body is covered in various scars and wounds.

\subparagraph{Symbol}
A growling dog skull.

\subparagraph{Quote}
"Most people think they can handle a dog cornered and rabid, but they can't, they don't know how sharp those teeth can be."

\subsubsection{The Catacomb Cobras}
Dressed all in black, and often adorned with capes or masks, these bandits are the stealthiest of the lot. They use a lot of poisons and sleight of hand. It's more often that those who double-cross the Cobras never see them coming. A lot of mysterious deaths are assumed to be linked back to this gang.

\subparagraph{Leader}
Veronica Venom, a human \textbf{assassin} who wears a black hood and dark leather armour. She has a bandolier of poison vials and 4 daggers fastened in her belt.

\subparagraph{Symbol}
A snake coiled around a skull.

\subparagraph{Quote}
"Sometimes, if you don't pay attention, you wind up in a snake pit by simply taking a wrong turn, and those snakes don't care if you're lost, they still bite."

\subsection{Missions}
\subparagraph{Tribute for the Queen}
If you want the Queen's protection, you need to pay your dues. Deliver a payment of 100 gold to the well in Buckledown Row. We don't care how you get it.

\subparagraph{Shakedown}
Go into the city and steal the spoils from a rival adventuring party.

\subparagraph{Grave Robbers}
Steal a valuable sceptre from the Chapel of Saint Brenna. (see Chapter 6)

\subparagraph{Entry Points}
Secure the smuggler's passage that ran under the Black Ivory Inn.

\subparagraph{Hired Goons}
Rough up \textbf{Oscar Yoren} and convince him to stop selling potions to the Hooded Lanterns, and start giving them to the Queen's Men. (see Chapter 6)

\subparagraph{Explosive Plans}
Purchase a \textit{wand of fireballs}, \textit{necklace of fireballs}, and ten scrolls of \textit{fireball} from the Amethyst Academy. We'll need them later...

\subparagraph{Browbeat}
The Dwarves of the Scar think they can stiff their payments to us. Show them the error of their ways.

\subparagraph{The Steward's Seal}
Find this missing Seal of Drakkenheim. (see Appendix D)

\subparagraph{Double Agents}
Ingratiate yourselves with the Hooded Lanterns or the Amethyst Academy. If you have to bump off a few of my boys to make it look convincing, do it. Once you've got their trust, we'll move on to the next steps.

\subparagraph{Rooked}
Infiltrate Saint Selina's Monastery and steal a Seal of Drakkenheim from the Followers of the Falling Fire.

\subparagraph{The Academy Job}
Steal the Inscrutable Tower. Yes, the whole tower. Bring back the \textit{Inscrutable Staff}. (See Chapter 8)

\subparagraph{The Cathedral Vaults}
Sneak into the Cathedral of Saint Vitruvio. Bring back everything in the House von Kessel Vault, as well as the Saint Vitruvio's Phylactery.

\subparagraph{Jailbreak}
Infiltrate the Hooded Lanterns barracks to rescue or kill someone held within.

\subparagraph{Sabotage and Misdirection}
Infiltrate Camp Dawn, and destroy their griffon tower, stables, hospice, and armoury. Make it look like the Hooded Lanterns were responsible.

\includegraphics[width=0.4\textwidth]{images/../5e.tools/img/adventure/DoDk/033-02-028.dice-throw.png}
\subparagraph{Such a Tragic Loss}
Arrange the death of Elias Drexel and Theodore Marshal, but make it look like it was their own mistaken folly, rather than an assassination.

\subparagraph{The Queen's Gambit}
Once the Queen of Thieves has assembled all six Seals of Drakkenheim, she asks the player characters to gather their best assault team. It's time to take the ultimate prize: the Queen of Thieves will seize the crown for herself!

\subsection{Boons}
The Queen of Thieves provides many favours to those who serve her well.

\subparagraph{Who Says Crime Doesn't Pay}
The Queen of Thieves makes sure her agents are well-paid. She rewards characters with 100 gold pieces each for successful missions to the Outer City, 250 gold each for missions to the Inner City, and 1,000 gold each for missions to Castle Drakken or faction strongholds.

\subparagraph{Poisons}
Trusted associates of the Queen of Thieves can procure any \textit{poisons} described in the \textit{Core Rules} from her quartermasters. The Queen of Thieves will give characters an allowance of up to 500 gold pieces to procure these poisons before any mission.

\subparagraph{Purloined Equipment}
If the player characters need a mundane item for a mission, the Queen of Thieves can use her web of connections to obtain it for them. This may be any mundane item in the Core Rules. How valuable the item can be depends on how long the characters are willing to wait to get it: 100 gold pieces per day spent seeking it out, up to a maximum of three weeks (or 2,100 gold)

\subparagraph{Potions and Scrolls}
The Queen of Thieves can provide the party with 1d4 Spell Scroll from the enchantment and illusion school of magic, and potions that aid stealth or trickery, such as a \textit{potion of invisibility} or a \textit{potion of gaseous form}. Determine these randomly from items up to very rare quality, or choose items appropriate for the mission at hand.

\subparagraph{Fences}
Access to smugglers and dealers of ill-gotten goods. There's many items stolen from noble houses, guild vaults, or churches that normal dealers simply won't purchase. These merchants have few scruples, and can find a buyer for almost anything.

\subparagraph{Smuggler's Passages}
The Queen's Men have maps of the smuggler's passages and sewers. Characters can use these maps to move through the city sewers at twice the normal speed, and make checks for random encounters while doing so with advantage.

\subparagraph{Stolen Magic Items}
As a reward for completing an important mission, the Queen of Thieves bestows the party with a magic item she has stolen. The Game Masters can choose this item or determine it randomly. The item can be consumable or permanent. This will be a common item for an Outer City mission, an uncommon item for the Inner City or a faction stronghold, or a rare item for a successful mission to Castle Drakken, the cathedral, or the Inscrutable Tower. As long as the items are randomly determined, there's no limit to how many times she might grant this boon.

\subparagraph{Cloak and Dagger Training}
The Queen of Thieves allows the characters access to her training rooms. Her trusted associates Lies \& Deceit can spend time showing the players new skills and combat techniques. Each character can gain one of the following of their choice:


\begin{description}
\item[Ability Score Improvement.]
The character can increase one ability score by 2 or increase two ability scores by 1 each. The ability score can be increased above 20, up to a maximum of 30.
\item[New Feat.]
The character gains a new feat of their choice (subject to GM approval).
\end{description}

\subparagraph{Bandit Backup}
If the player characters have a reasonable need for reinforcements during a mission, the Queen of Thieves can dispatch one of her gangs to help them out. The gang sends 1d6 bandits and a \textbf{bandit captain} to help out. They demand an equal share of any loot the characters find on the mission, and a share of any payment. If all bandits are killed on the mission, the player characters will not receive this boon again. If circumstances with other factions escalate into open conflict, the Queen of Thieves might dispatch 1d4 of her assassins, mages, or gladiators to join the characters.

\subparagraph{Informants}
The Queen of Thieves has a wealth of spies and informants. The player characters can ask her to spy on or investigate a location, character, or event. She can provide a clue or details on that location. Her information is reliable, but not always complete.

However, the Queen of Thieves always uses her discretion and judgement in giving out this information, and she's cautious to reveal anything that could be used against her or be used to predict her own plans. She knows the value of revealing information at the right moment, so she's careful about sharing information she might use as blackmail later.

\includegraphics[width=0.4\textwidth]{images/../5e.tools/img/adventure/DoDk/034-02-029.wooden-toy.png}

\begin{itemize}
\item She never reveals the locations of the \textit{Seals of Drakkenheim}, even if she's sending the characters to claim one, she doesn't reveal that the item in question is a \textit{Seal of Drakkenheim}.
\item She does not tell anyone about the powers of the \textit{Crown of Westemär}, what she plans to do with it, or any personal information about herself. If asked, she laughs.
\item The Queen of Thieves knows the ideals, bonds, and flaws of almost any non-player character.
\item She can explain the background, objectives, and goals of the other factions, and the complications they face.
\item She can provide a clue about an alternative entrance to a location, or details about the monsters who might dwell there or the treasures hidden within.
\item She knows what happened to the royal family, but never discloses this information to the player characters.
\item Furthermore, she isn't privy to the inner workings of the Amethyst Academy, she doesn't know what's in the crater or where delerium came from, or if there is any truth to \textbf{Lucretia Mathias}' claims.
\end{itemize}

\subsection{Strike Teams}
\subsubsection{The Friendly Barmaid}
Cora is a human \textbf{spy} who sometimes works at the Bark and Buzzard in Emberwood Village. She's an informant for the Queen of Thieves and a notorious gossip. She likes to talk to drunk adventurers, flirt with them, and find out what they're planning.

\subsubsection{Outlaw Gang}
A typical gang consists of 24 (7d6) bandits or thugs, 2 (1d4) bandit captains, and 4 (1d6) spies. Special members include:


\begin{DndTable}{cX}

1d6 & Special Member(s)\\
1 & None\\
2 & 7 (2d6) worgs\\
3 & 4 (1d6) bugbears\\
4 & 2 (1d4) ogres\\
5 & 1 (1d3) doppelgangers\\
6 & 1 \textbf{hedge mage}\\
\end{DndTable}

When called upon, the gangs do not necessarily send their entire outfit. You should determine an appropriate number of members who carry out the strike mission. They always send enough to outnumber the characters at least two to one, however.

\subsubsection{Enforcers}
When the Queen of Thieves needs to send a strong message, she sends a group of 4 (1d6) bugbear gladiators to start cracking skulls and breaking legs. They are also assigned to guard any prisoners.

\subsubsection{Infiltrators}
This strike team of 2 (1d4) doppelgangers is led by one known as Thought-stealer. The Queen of Thieves herself oversees this strike team, using \textit{Rary's Telepathic Bond} and \textit{scrying} spells. This strike team only conducts infiltration, information gathering, and clandestine asset recovery. They're rarely used for direct attacks such as assassination or kidnapping: they are her eyes and ears.

\subsubsection{Saboteurs}
Lies (a tiefling \textbf{mage}) and her twin brother, Deceit (a tiefling \textbf{assassin}) coordinate a strike team of 5 scoundrels. This strike team fulfils a wide variety of roles, and is sent in when the Queen of Thieves needs exceptionally skilled agents for an important mission. Occasionally, the Queen of Thieves uses \textit{modify memory} on members of this strike team so they can't reveal secrets about their missions when captured, or even casts \textit{dominate person} on one of them so she can "personally" oversee the operation.

\subsubsection{Royal Assassins}
This elite strike team consists of one \textbf{assassin} per player character. It's led by Dell Talquamar and the Boogeyman (both assassins) alongside Jerry Rigz, a goblin \textbf{mage} to provide logistics and magical support. This strike team is only deployed to steal a \textit{Seal of Drakkenheim}, to take out a rival faction leader, or to destroy the player characters.

\includegraphics[width=0.4\textwidth]{images/../5e.tools/img/adventure/DoDk/035-02-030.blackjack-mel.png}
\subsection{Schemes}
\subparagraph{Shakedown}
The Queen of Thieves sends one of her gangs to go beat the party up after a mission in the ruins. She doesn't have them killed, but she steals their equipment and valuables.

\subparagraph{Blackmail}
The Queen of Thieves uses her agents and magic to discover a terrible secret about one of the player characters, or some crime or trespass they committed in the ruins. She threatens to reveal the secret if the characters don't work for her.

\subparagraph{Bribery}
Using the intel she has gathered, the Queen of Thieves caters to the goals of the party and offers to pay them an exorbitant sum of money (as you see fit) to buy the characters' loyalty—or at least their services.

\subparagraph{Eyes and Ears}
Spies tail the party and gather information on their next moves and plans, along with any factions they associate with for the next 1d4 days. If undiscovered they relay all information to the Queen of Thieves for her to use to her advantage.

\subparagraph{Backstabbing}
The Queen of Thieves expects the player characters will betray her. She is never surprised when this happens, and should always have a plan in mind for if the player characters' decide to double-cross her.

\subparagraph{Puppetmaster}
The Queen of Thieves kidnaps a faction member. Thereafter, she controls the faction agent using \textit{dominate person} and sends them back as a spy. If caught, the Queen of Thieves might try to create the impression that the puppet is actually a \textbf{doppelganger} or magical disguise, rather than a dominated person. Thanks to the Queen's \textit{Modify Memory} power, this faction member won't have any knowledge of what happened.

\subparagraph{An Audience with the Queen}
The Queen of Thieves invites the characters to the Court of Thieves to negotiate. Once they arrive, she lures them into the arena to face her gladiators, and then her favoured pet: Big Linda.

\subparagraph{Captured by the Queen}
The Queen of Thieves prefers taking the characters prisoner over killing them. If captured, she'll interrogate them using \textit{suggestion} and \textit{detect thoughts} spells to learn their secrets and gain leverage over them, so she can force them to do her bidding.

\subparagraph{Imposter Syndrome}
The Queen of Thieves sends doppelgangers disguised as the party to Emberwood Village where they cause a commotion and set fire to a building. They flaunt their false identities and make sure to tarnish the names of the characters.

\subparagraph{Sabotage}
The Queen of Thieves sends a strike team to sneak into a faction stronghold with explosives (perhaps a Spell Scroll of \textit{fireball} or \textit{necklace of fireballs}) and blow up a critical resource.

\subparagraph{Theft}
The Queen of Thieves dispatches one of her elite strike teams to steal magic items or valuable information from the player characters or their faction allies.

\subparagraph{An Offer You Can't Refuse}
The Queen of Thieves secures an asset valuable to completing a character's personal quest. She uses it as a bargaining chip.

\DndQuote%
{}
{''No matter how she may look on the outside, I know the Queen of Thieves is rotten to the core.''}
{Veo Sjena}
\chapter{Emberwood Village}
\begin{DndReadAloud}{}
\DndDropCapLine{A}{} few miles south of Drakkenheim lies Emberwood Village, a humble crossroads along Champion's Way surrounded by forsaken farmsteads and dark pine forests. Many residences are abandoned and slowly deteriorating with their windows and doors nailed shut. However, the village centre is alive with activity. Conversation, laughter, and music resound from several taverns along the main thoroughfare. Amidst a bustling market square, a small crowd of adventurers, merchants, and prospectors busily traffic all manner of weapons, equipment, coins, recovered treasures, and a few strange glimmering crystals. Even the light of the Sacred Flame still glows within the modest stone chapel at the heart of town.


\end{DndReadAloud}

\section{Overview}
This chapter details Emberwood Village, a safe haven where player characters may rest, recuperate, and resupply between their adventures in Drakkenheim. Equipment, lodgings, spellcasting services, a few magic items, and other amenities are available from the many folk dwelling in the village.

By interacting with the various townsfolk, faction agents, and other non-player characters in Emberwood Village, the player characters will discover many adventure hooks and rumours that direct them towards key locations within Drakkenheim. The taverns and inns are the perfect place to make initial contact with various faction lieutenants who use the establishments as neutral ground for meetings. Many rival adventurers can be encountered in Emberwood Village as well, ready to swap stories and stoke friendly competition.

\includegraphics[width=0.4\textwidth]{images/../5e.tools/img/adventure/DoDk/036-03-001.emberwood-village.png}
\subsection{Emberwood Village Background}
Two decades ago, Emberwood Village was little more than a farmers market, blacksmith, a pub, a small chapel and a few dozen cottages. The local businesses served the farmlands and logging camps surrounding the capital, and the town was an occasional rest stop for merchants to stop off on their way to Drakkenheim.

As the town lies some five miles south of the capital city, it was not directly damaged by the meteor shower, though the impact shattered the windows of every building, caused some weaker structures to collapse, and left some in the village blinded or deafened. Unfortunately, the aftermath was all the more terrible. Over the past fifteen years, the contamination in Drakkenheim has affected the countryside for miles around. The once fertile fields and dense woods are now barren and dying. Gradually, most folk gave up, departing for safer and more prosperous places. A few people remained behind, determined to hold onto their family homes and salvage their livelihoods.

Emberwood Village persisted despite the odds. After the civil war, adventurers began showing up in the town. Explorers heading towards the dark city needed a safe place to tend their wounds and rest, and so Emberwood Village gradually became a sort of boom town on a new frontier for adventure.

\section{Interactions}
Emberwood Village is home to only about two hundred permanent residents, almost all are commoners who endured years of hardship after Drakkenheim was destroyed. However, the recent flow of coin from explorers scouring the ruins allows many to live modest or comfortable lifestyles. The common folk tend to keep their dealings with adventurers friendly but professional. They give rough-looking warriors a wide berth, but enjoy hearing news, gossip, and stories about daring exploits in the ruins. As long as interactions stay peaceful, the villagers will deal with anyone regardless of their faction allegiances. The villagers live in strange times, and are used to seeing odd-looking folk and weird magic. They can recognize delerium, severe contamination, and mutations, and encourage such individuals to seek magical aid.

\subsection{Arrival by Caravan}
If the characters arrived in Emberwood Village in the company of Eren Marlowe (see "The Road to Drakkenheim"), they are paid the agreed rate of 25 gp per character. While Marlowe doesn't sell goods commonly needed by adventurers, they can offer some friendly advice about town:


\begin{itemize}
\item They invite characters to meet them again at the market square in the village. Marlowe will introduce characters looking for adventuring equipment to Armin Gainsbury.
\item Marlowe suggests characters visit the Bark and Buzzard if they're looking for a place to relax and get a good meal. Characters can ask the owner, Karin Alsberg, about setting them up in one of the empty cottages around town should they need accommodations.
\end{itemize}

\subsection{A Curious Crowe}
Emma Crowe, a ten-year-old red-headed and freckle-faced human \textbf{commoner} keeps lookout from a dead tree on the way into town. She skips up to adventurous looking newcomers and introduces herself as the daughter of the local blacksmith, Tobias Crowe. She offers them a tour around Emberwood Village for a gold piece. At each location, the precocious youngster tells a story about an adventurer who once frequented the business, then gives a lighthearted account of the grisly fates they met in Drakkenheim. The tour ends at the Crowe and Sons Smithy.

\subsection{Adventuring Equipment}
Most weapons, armour, and adventuring equipment described in the \textit{Core Rules} are available in Emberwood Village. The various merchants and their wares are described in the Area Details below.

Emberwood Village relies entirely on trade for goods and supplies, since all other settlements for miles around are now abandoned and local food production has become impossible. As a result, weapons, armour, and adventuring equipment are sold for double the normal price found in the Core Rules. Food costs five times as much, and water isn't given freely. Instead, it is priced like alcohol.

Exotic equipment, vehicles, mounts, or other mundane items that cost over 500 gp are especially rare and must be brought in from outside Drakkenheim. It takes 21 (6d6) days to deliver such items, which cost five times the typical price.

The folk of Emberwood Village have grown accustomed to adventurers trying to haggle and swindle them. They don't give discounts or make deals, and tell adventurers who push for better prices to pay up or go elsewhere (there isn't anywhere else to go).

\subsection{Selling Treasure}
Most merchants in Emberwood Village will trade or purchase found treasure, salvaged equipment, gems, or art objects at half their listed value.

\subsection{Trading Delerium}
Only a few intrepid and wealthy purveyors can afford to trade in delerium. These include Aldor the Immense, Orson Fairweather, and River of the Amethyst Academy. Typically, most delerium is bought and sold for about half its listed value in Appendix D. However, some traders may offer more preferable rates for proven associates. Due to their rarity, crystals and geodes typically command their full market value.

\subsection{Lodgings}
Though there are many taverns and public houses in Emberwood Village, the only actual inn is the Red Lion Hotel (see description below). However, the Alsbergs of the Bark and Buzzard tavern can make arrangements for characters to bunk in one of the abandoned homes and cottages in town.

\subsection{Magic Items and Spellcasters}
A few NPCs in Emberwood Village sell a limited selection of magic items, mostly potions and scrolls:


\begin{itemize}
\item Aldor the Immense sells a small selection of uncommon and rare magical items, and will offer to purchase found magic items as well.
\item Flamekeeper Hanna or Old Zoya offer spellcasting services of spells up to 3rd level.
\end{itemize}

These characters also sell expensive spell components for spells of 3rd level and lower, and occasionally have one or two components for spells of 4th or 5th level.

Characters seeking more powerful magic items, higher level spell components, or the services of higher-level spellcasters must seek assistance from the factions.

\DndQuote%
{}
{''A lot of fires happen in this town. Oh wait, none have happened since I left? That's strange.''}
{Sebastian Crowe}
\subsection{Law and Order}
Although the Hooded Lanterns claim otherwise, there is no established law and order in Emberwood Village. Most businesses provide their own protection and security, and the villagers look out for each other. Word travels fast, though. Services, supplies, or shelter aren't available to adventurers with a violent and antagonistic reputation—for any price.

Belligerent curs who disturb the peace are beaten. Those who commit robbery, fraud, and property damage are driven out of town and never welcomed back. Murderers and arsonists are executed. Crimes of any sort involving magic are punished by hanging. On the flipside, there's little the villagers can do when perpetrators aren't caught red-handed. There's no authority willing to investigate crimes or search for missing people, so the villagers must petition adventurers or the factions for help. There are no trials; the mob becomes judge and jury.

In this respect, characters who make themselves villains may be targeted by the Hooded Lanterns, the Silver Order, or even the \textbf{Queen of Thieves} should they directly threaten the immediate safety of Emberwood Village. Alternatively, a rival adventuring party might even take up the bounty or intervene in a crisis. However, if any of the villagers are killed, no one comes to replace them.

\subsection{Making Contact with the Factions}
During their first visit to Emberwood Village, characters may encounter several faction agents. Each is a key contact between the characters and the factions during their adventures in Drakkenheim:


\begin{description}
\item[Ansom.]
and \textbf{Petra Lang} are human urban rangers of the Hooded Lanterns. They are frequently relaying messages and leading \textbf{scout} teams to the Old Watchtower, and can be encountered there most days.
\item[Blackjack Mel.]
is a human \textbf{scoundrel} working for the Queen's Men. He can always be found at the back tables of the Skull & Sword with a bodyguard of 2d6 thugs. He's happy to talk shop with mercenaries.
\item[Nathaniel Flint.]
is a human \textbf{chaplain} of the Followers of the Falling Fire. He is regularly seen at the Hendrix Farm preparing the latest group of pilgrims for the last leg of the journey towards Drakkenheim. He uses the old barn as a meeting place for important matters.
\item[Ophelia Reed.]
is a human \textbf{chaplain} of the Silver Order. She is visiting the Chapel of Saint Ardenna where Flamekeeper Hanna has granted her permission to hold meetings and discussions.
\item[River.]
is a tiefling \textbf{mage} of the Amethyst Academy. She can be met at the Red Lion Hotel. She rents a private suite and is often reading in the parlour, where she might chat with adventurers.
\end{description}

Roleplaying traits for each of these characters are described in detail in chapter 3 but all act mostly neutral during their first encounter with the player characters. Instead, they try to size up their potential. Each explains some information about their faction, such as their ideals and general mission (see the faction dossiers in chapter 3 for more information). These NPCs might share a drink or meal with the characters, ask them if they've been to the ruins yet, or inquire about why they've come to Drakkenheim. However, they won't offer missions or quests until the characters have survived at least one expedition of their own to the ruins.

\section{Area Details}
\includegraphics[width=0.4\textwidth]{images/../5e.tools/img/adventure/DoDk/037-map-3.01-emberwood-village-map.png}
\includegraphics[width=0.4\textwidth]{images/../5e.tools/img/adventure/DoDk/038-map-3.01-emberwood-village-map-player.png}
\subsection{Bark and Buzzard}
Nestled amongst the residential cottages of Emberwood Village, this small one-level public house is a popular spot for locals and adventurers seeking a peaceful and laid back place to unwind. The sign hanging above the door depicts a cartoonish vulture feasting on the entrails of a bloodhound, both smiling pleasantly. The public house has a cozy and familial character: there's no proper bar, nor any walls separating the kitchen from the dining room. Instead, a few benches are arranged beside the main cask, and about a dozen round tables surround an inviting hearth. The Alsberg family has owned and maintained the establishment for a generation. Karin Alsberg and her husband, Holger (both human commoners) make fine home-brewed beer known as Ember Ale, and serve hearty braised lamb stew with fresh-baked bread.


\begin{itemize}
\item A mug of fine ale costs 2 sp. A meal costs 1 gp.
\item The Alsbergs keep half a dozen one-room cottages in the village that they will rent to travellers. The price is 2 gold per night.
\end{itemize}

\subsubsection{Old Rattlecan}
Sitting behind the Bark and Buzzard like a scarecrow, Old Rattlecan is an intelligent suit of animated armour which was found in the ruins of Drakkenheim. A mysterious magical enchantment causes it to rebuild itself when destroyed:


\begin{itemize}
\item Rival adventurers drinking in the Bark and Buzzard may invite player characters to "come kick the can out back". They challenge them to see which character can defeat Old Rattlecan the fastest, or who can survive the longest before being knocked out. Adding handicaps or bizarre "house rules" are popular ways to spice up the contest, such as being blindfolded, using no weapons, or taking drinks between hits.
\item Rattlecan rejuvenates after one minute when reduced to 0 hit points unless it is \textit{Disintegrate} or \textit{dispel magic} is cast on its remains. Everyone in town is terribly upset if this happens.
\item A creature reduced to 0 hit points by Old Rattlecan is unconscious but stable and wakes up one minute later with 1 hit point.
\item Old Rattlecan can occasionally be heard making remarks to those it fights. It yells things like "Have at thee cur!", "'Tis a flesh wound!" or "Was that the best you can do?" as it seems to revel in battling any challenger.
\end{itemize}

\subsection{Gilded Lily}
This lavish two-storey building is surrounded by a veranda decorated with a swirling lily motif and painted purple. It has ornate golden glowing lanterns hanging between the veranda arches and within its curtained windows. Inside, four large golden lamps light a small stage set before rows of candlelit tables. Floral arrangements line the walls. The well-stocked bar displays wines and spirits from across the world. The Gilded Lily is owned by a human \textbf{noble} named Mr. Killian Vandire, but the exquisite hostess Madam Rochelle (a human \textbf{spy}) runs the business and organizes the nightly performance acts. The tavern is a popular place to share company, gossip, and have a few drinks, and many adventurers make their last stop here for some good memories before trudging into the dismal ruins.


\begin{itemize}
\item Drinks cost 5 silver: wine, spirits, and coffee only, no beer or food is served here.
\end{itemize}

\subsubsection{Open Mike Night}
"Mike" was the master of ceremonies who went missing shortly after the meteor fell. One of Mike's final wishes was to open the stage to any performer for magic shows, musical numbers, poetry readings, performances, or storytelling. The Gilded Lily is willing to host any performer who wishes to take the stage and though they do not pay their acts, they do put a hat out so patrons can toss a coin to their performer if they like the show.


\begin{itemize}
\item Any sort of act is allowed in the Gilded Lily. Performers may tell a story, recite a poem, sing, or play an instrument. Spellcasters might put on a magic show, and even demonstrations of acrobatic or athletic skill are welcome. Player characters are invited to be as creative as they want to earn some extra coins.
\item Player characters who take to the stage may make a skill check related to their performance: acrobatics, arcana, athletics, history, and performance are all suitable ways to represent their act.
\item Characters are paid in tips: they receive an amount of gold based on their skill check result:
\end{itemize}


\begin{DndTable}{cX}

9 or less & The character is booed off the stage, and receives no gold for their performance. If this happens three or more times, they are banned from ever performing again.\\
10 to 14 & Characters collect 2d6 gp in coins from tips.\\
15 to 19 & Characters earn 5d6 gp from the impressed audience\\
20 or higher & The wonderful performance earns 10d6 gold pieces, and characters are invited to perform an encore performance. However, characters must exceed their previous check result to win further encores.\\
\end{DndTable}

\subsection{Skull \& Sword Taphouse}
This dingy saloon was converted from an old farmstead. A jawless skull impaled upon a rusty sword is embedded beside the main doors: no one knows who it was or who did them in, but it's tradition to make a macabre greeting to the skull on your way into the bar. Inside the establishment is a rowdy environment - brawls and fist fights happen daily, and several ruffians are thrown out nightly.

The Skull \& Sword is a popular place for scoundrels and misfits to gather thanks to the cheap mead and liquor served here. It may very well be the best place in Emberwood Village to pick up rumours from the city, but you may want to keep a low profile hanging out around there, or be ready to draw swords for looking at someone the wrong way.

The establishment is run by a half-orc \textbf{gladiator} named Frida Longhorn, who keeps most people in check. Despite the rough reputation of the establishment, Frida is pretty good at keeping the peace and making sure there are no killings or broken windows in her bar. Her business partner, a halfling \textbf{spy} named Cuff Goldburg, is usually sitting at one of the back tables with Blackjack Mel. Cuff typically excuses himself when Blackjack Mel deals with mercenaries and adventurers.

The \textbf{Queen of Thieves} bankrolls the bar, ensuring the flow of cheap booze (which is almost all stolen). While commonfolk in Emberwood Village believe the Skull \& Sword is a hub for unlawful activity, most of what goes on here is above board aside from a few sketchy meetings. The real dirty deals and smuggling happen in Buckledown Row where the other factions can't easily interfere.


\begin{itemize}
\item A mug of cheap mead is 5 copper pieces. Brandy, rum, vodka, and gin are 1 silver.
\item Blackjack Mel is here every night, where he has a favourite spot amongst the back tables with 7 (2d6) thugs. He's eager to meet potential recruits in a private booth in the back corner.
\end{itemize}

\subsection{Red Lion Hotel}
Upon a small hill overlooking the village is the impressive Red Lion Hotel. Before the meteor struck it was known as Stavros Manor, home to a human \textbf{noble} named Kosta Stavros who served as Reeve of Emberwood Village. After narrowly avoiding destitution when Drakkenheim was destroyed, Stavros converted his home into a hotel that serves wealthy explorers, merchants, and nobles conducting business near Drakkenheim. He is an older man with a sallow face, thinning hair, and monocle, always garbed in a rich blue vest and doublet.

The three-storey manor is made of red brick. It features lion-shaped gargoyles, large glass-paned windows, and embossed wooden doors painted green. The estate includes a dozen well-appointed bedrooms, each with their own private bath and common room. A full service staff provides patrons' every need. Meals are taken in a sumptuous dining room outfitted with an impressive crystal chandelier, and drinks are served in a vast parlour furnished with fine antiques and decorated with various sculptures, paintings, and bookshelves. Each room features a great hearth and lovely views which overlook the village or the surrounding countryside.

The Reeve has abandoned his civic duties, and is now little more than another business owner in the town. However, he still lords his imagined authority and privilege over others in the village, and prefers to share wine and gossip with travellers and merchants than deal with common folk. He regards the idea of adventuring in Drakkenheim as a sort of sport. He often makes wagers on who will die first amongst any adventuring groups who visit the Red Lion Hotel.


\begin{itemize}
\item Kosta Stavros loves the fine arts, especially sculpture, paintings, and literature. He's assembled a vast collection of art objects recovered by adventurers from the ruins, and exchanges coin or credit towards accommodation for works in good condition.
\item River of the Amethyst Academy stays at the Red Lion Hotel. She's happy to chat, but prefers to conduct business privately. She asks any potential new associates of the Amethyst Academy to meet with her later at Eckerman Mill (see chapter 6).
\item The Ironhelm Dwarves (see Smithy on the Scar in chapter 6) keep rooms here, along with many merchants who operate in Caravan Court and one or two groups of wealthy rival adventurers.
\item Rooms are 25 gp per night and include meals.
\item Wine and liquors are sold at 100 gold per bottle.
\end{itemize}

\subsection{Crowe and Sons Smithy}
This smithy is little more than a one-level log cabin with an overhanging thatched roof covering a workbench, stone forge, anvil, and grinding wheel. Various tools hang along the walls beside several finished weapons.

Master smith Tobias Crowe (a human \textbf{veteran}) owns and operates the smithy with his son and apprentice, Peter (a \textbf{commoner}). Crowe is a quiet family man who stands six feet tall with fiery red hair and beard. He's often drenched in sweat and soot, and wears a greasy pair of old goggles. Crowe proudly demonstrates the high quality of his work to customers, and openly criticizes the shoddy workmanship of their existing equipment. He spends most of his days repairing broken gear, and his evenings at the Bark and Buzzard, before returning home to his wife Moira and daughters Emma and Sybil. Although he lacks the spellcasting abilities to work with delerium, he dreams of one day creating a masterpiece out of the meteoric iron found in Drakkenheim.


\begin{itemize}
\item Tobias Crowe sells weapons and armour at twice the price found in the Core Rules.
\item Plate armour must be made-to-order. Payment is triple the list price, paid up front, and is available in four weeks.
\item He accepts payment in coins, gems, or precious metals.
\item He pays half-price for salvaged equipment if he thinks he can repurpose or melt it down.
\end{itemize}

\includegraphics[width=0.4\textwidth]{images/../5e.tools/img/adventure/DoDk/039-03-002.approaching-emberwood.png}
\subsection{Chapel of Saint Ardenna}
Located at the heart of Emberwood Village is a small chapel of the Sacred Flame. The circular fieldstone building has a copper domed rooftop, thin windows set with stained glass, and simple wooden double doors. Outside is a small crematory garden and scattering plot for the ashes of the dead. Within the chapel is a bronze \textit{brazier of the Sacred Flame}, alight with golden fire that burns all night and day. It serves as a peaceful place of quiet worship in the harsh wasteland surrounding Drakkenheim.

The chapel is maintained by Flamekeeper Hanna (a human \textbf{priest}) and 2 (1d4) human acolytes. Hanna is a young woman who wears a simple white and yellow frock and keeps her black hair tied in a braid. Hanna was merely an acolyte when the last Flamekeeper died several years ago, but opted to take up the position herself rather than abandon the folk of Emberwood Village. She believes the Faith of the Sacred Flame is needed here now more than ever, and offers hospitality, kindness, and prayer to any who visit the chapel.

\textbf{High Flamekeeper Ophelia Reed} of the Silver Order regularly visits the Chapel of Saint Ardenna, protected by Ser Virgil Underwood and Ser Cassandra Wyatt, two Silver Order cavaliers. Ophelia Reed has been trying to convince Flamekeeper Hanna to join the Silver Order, but thus far Hanna has kept her neutral allegiance to the folk of Emberwood Village. Still, Flamekeeper Hanna enjoys conversation with Ophelia Reed, and allows the Silver Order emissary to take shelter in the chapel and conduct meetings with adventurers here.

\subsubsection{Spellcasting Services}
Flamekeeper Hanna can cast cleric spells of 3rd level and lower, and knows the \textit{purge contamination} spell and always keeps it prepared. She charges 25 gp for a 1st-level spell, 50 gp for a 2nd-level spell, or 150 gp for a 3rd-level spell, plus the cost of any expensive material components.

She takes payment in coins, gems, or relics. She uses the money to help fund the caravans that supply Emberwood Village with basics for the townsfolk. Flamekeeper Hanna will occasionally waive these fees to help innocent folk in dire situations, but not adventurers who went looking for treasure and trouble in Drakkenheim.

\subsubsection{Spell Scrolls}
Paid in advance, once per day Flamekeeper Hanna can create Spell Scroll of any cleric spell she can cast. She charges 75 gp for a 1st level \textit{spell scroll}, 150 gp for a 2nd-level \textit{spell scroll}, and 450 gp for a 3rd-level \textit{spell scroll}, plus the cost of any expensive material components. There's a single exception: a \textit{scroll} of \textit{revivify} costs 1,000 gp.

\subsection{Watchtower}
This fortified stone watchtower is the tallest structure in Emberwood Village, and lies upon the road on the north edge of town. A two-storey wooden house abuts the base. Inside is a small kitchen, a mess hall, and about two dozen bunk beds. The tower has three main levels, one serving as a meeting room, another as armoury, and the final as a private chamber for the watchmaster. Outside the tower the Hooded Lanterns keep a small training yard and a modest stable large enough for six horses. Nearby, a covered wooden bridge crosses the narrow Timberwash stream and leads out of Emberwood Village towards Drakkenheim.

This location once served as the initial outpost for the Hooded Lanterns until they successfully established a base of operations within the city proper, and the remnants of a much larger military camp are scattered about the nearby field. Now, only a small detachment keeps watch in Emberwood Village.


\begin{itemize}
\item A small force of a dozen Hooded Lantern scouts are posted here, with one or two other squads out on patrol in the countryside.
\item Raine Highlash, an \textbf{urban ranger} with a predatory look about her, serves as watch captain.
\item Ansom or Petra Lang report in here regularly at the behest of the Lord Commander, and use the outpost to meet with mercenaries and adventurers.
\end{itemize}

\DndQuote%
{}
{''Warm fires, great company, home cooked meals! What more could you ask for in a city like Drakkenheim? Did I mention the home-cooked meals?!''}
{Veo Sjena}
\subsection{Shrine of the Old Gods}
Seven standing stones comprise this ancient shrine that rests within a small grove of dying trees on the outskirts of town. Inscribed on each stone are the legends of primordial gods, and between them rests a stone pedestal adorned with various animal bones, plants, and candles that are lit every night. A crude old hut lies down a dirt path near the shrine, smoke billowing from a hole in the thatched roof chimney, and an odd scent of herbs wafting from within.

The hut is home to Old Zoya, an elderly human \textbf{druid} who tends the shrine. Many in town wrongfully refer to her as a witch, but hold a deep respect for the ornery mystic. She claims the shrine protects Emberwood Village through these dark times. Zoya holds no love for the Sacred Flame, and keeps faith for the old, true gods.

\subsubsection{Spellcasting Services}
Old Zoya can cast druid spells of 3rd-level and lower, and knows the \textit{purge contamination} spell and always keeps it prepared. She charges 25 gp for a 1st-level spell, 50 gp for a 2nd-level spell, or 150 gp a 3rd-level spell, plus the cost of any expensive material components.

She takes payment in coins, gems, or relics. She refuses to speak of what she does with the money, and grows impatient with anyone who gets too nosy.

\subsubsection{Potions of Healing}
Each week, Zoya brews up 1d6 Potion of Healing which she sells for 50 gold each. If paid in advance, once per day she can create a \textit{potion of greater healing} for 150 gp.

\subsection{Hendrix Farm}
The Hendrix Farm is one of the only occupied farmsteads remaining in Emberwood Village. It was once run by an old human \textbf{commoner} named Elijah Hendrix, until his crops became contaminated, his livestock perished, and his wife grew ill and passed away after the meteor fell. He gave up on farming, and nearly gave up on life as well. However, several years ago, \textbf{Lucretia Mathias} stayed at his farmhouse during her first pilgrimage to Drakkenheim. Her words inspired him with renewed purpose, and he now offers shelter to any future pilgrims of the Falling Fire who are making the trek into Drakkenheim.

The fallow fields are devoid of crops and no longer used for agricultural production. Instead, they are littered with small campfires and tents for the passing pilgrims. The large wooden barn and nearby farmhouse are kept in good repair, though all the farm equipment and ploughs were long ago dismantled for parts.

\textbf{Nathaniel Flint} and many members of the Falling Fire gather here as their last stop before heading further into the ruins. At any given time there may be as many as fifty commoners congregating here with their wagons. Though he has been invited to come take the Sacrament himself many times, Farmer Hendrix believes it is his destiny to wait here until the very last pilgrim comes.

\subsection{Eventide Manor}
This manor house on the south end of town burned down many years ago. It was once owned by a mysterious elven woman who had a child with a local craftsman, but was abandoned after she passed away and the boy was taken by the Amethyst Academy. Since it caught fire, none have come to claim or repair it. Children now tell tales of the strange noises coming from the cellar of the house, and how the charred doors and windows open and close on their own.


\begin{itemize}
\item Unbeknownst to the folk of Emberwood Village, this house was owned by a former Archmage of the Amethyst Academy. Behind an illusory wall in the cellar protected by a \textit{glyph of warding} is a \textit{teleportation circle} used today by members of the Amethyst Academy to come and go from their distant strongholds. Academy members only enter or exit the manor under the effects of an \textit{invisibility} spell, hence the children's rumours.
\end{itemize}

\subsection{Caravan Court}
A marketplace has sprung up surrounding the boarded-up water well in the village square. This colourful collection of caravans, wagons, canvas tents, and makeshift stalls is the hub of commerce for adventurers, scavengers, and prospectors looking to make their fortune in Drakkenheim. Dozens of canny merchants have set up shop here to cater to the needs of the factions and explorers.

Traders and new goods arrive here daily, and the market is constantly abuzz with explorers looking to sell relics and treasures unearthed in Drakkenheim, or buy equipment needed for a dangerous voyage into the city ruins. Amongst the bustle, heavily armed and armoured thugs load ponderous wagons with locked and leaden chests filled with delerium shards for trade in distant cities. Activity carries on through the afternoon, but most pack up shortly after sunset to head to the taverns and bars for drinks.

\subsubsection{Marlowe's Provisions}
\DndQuote%
{}
{''Everyone's gotta eat! Don't want to go trouncing through dangerous ruins on an empty stomach! Oh, and please drink some water, you look parched.''}
{Eren Marlowe}
Eren Marlowe, a human \textbf{commoner}, sets up a small makeshift stall with a few wooden tables displaying jars of dried fruits, packages of cured meats, jugs of water, mead, and bags of assorted nuts. Adventurers often come here for provisions that won't rot or spoil during their treks into the contaminated ruins.


\begin{itemize}
\item Any food, water, and rations can be purchased here for five times the market price
\end{itemize}

\subsubsection{Gainsbury Expeditionary Supply Company}
\DndQuote%
{}
{''I guarantee you won't survive long in Drakkenheim without a good rope and a full lantern! Best buy them while supplies last!''}
{Armin Gainsbury}
This caravan is loaded high with boxes and barrels of goods. Laid out on a small table are samples of climbing gear, cartographer's instruments, goggles, heavy leather gloves, mining picks, torches, tongs, and assorted tools for navigation and survival in the ruins.


\begin{itemize}
\item Armin Gainsbury, a bespectacled and beanpole-shaped human \textbf{commoner}, carries all manner of adventuring gear which they sell for twice the market value.
\item Gainsbury is especially well stocked on mining and prospecting equipment, and sells a pamphlet detailing the correct methods for safely handling and extracting delerium crystals. A copy costs 1 gp. The information inside is filled with hyperbolic (but correct) warnings. The instructions are overly detailed and rather long-winded, but reasonably accurate.
\end{itemize}

\subsubsection{Fairweather Trades and Exports}
\DndQuote%
{}
{''When you've been around the world as many times as me, you see all manner of treasure. In all my years, I ain't never seen anything like these rocks. I'm happy to make a barrel or two of gold delivering it to whomever wants to empty their pockets for 'em.''}
{Orson Fairweather}
An iron-bound stagecoach is parked here flanked by two massive flat-bedded wagons with stacks of leaden and steel-studded chests, overseen by a dozen armoured men.


\begin{itemize}
\item Orson Fairweather (a human \textbf{gladiator}) is among the most prolific exporters of delerium operating in Emberwood Village. The stout man has a heavily scarred and aged face, a leather eyepatch, and a mouth of gold teeth. He speaks little of his own history or exploits, but entertains rumours about his past life as a pirate and adventurer.
\item Fairweather has a security force of a dozen veterans, and has paid the Amethyst Academy to create \textit{Arcane Lock} and \textit{Glyph of Warding} to protect his wares from thieves.
\end{itemize}

Fairweather is among the few merchants in town willing and able to purchase delerium in vast quantities. He offers half the list price in Appendix D and doesn't make deals. He furiously rebukes anyone who questions his motives or asks prying inquiries about his buyers. Countless folk in distant lands purchase delerium for unknown reasons, and in the grand scheme of things, Fairweather is just a middle man.

\subsection{Aldor the Immense}
\includegraphics[width=0.4\textwidth]{images/../5e.tools/img/adventure/DoDk/aldor-the-immense.png}
\DndQuote%
{}
{''Anyone who says money can't buy happiness doesn't know where to shop.''}
{Aldor the Immense}
This massive and heaving covered carriage is decorated with gold filigree and glimmering lanterns. Hitched beside it are two great steel-skinned oxen that exhale green fumes from their snouts. Three lanky attendants garbed in light armour and carrying crossbows patrol the grounds.

Seated before the carriage on a buckling wooden chair is a massive human man who must be nearly seven feet tall and almost as wide. He is garbed in masterfully tailored jet-black silk that flatters his obese frame: a frilled doublet, padded breeches, and riding gloves, all worn under a great fur-lined cloak, topped with a feathered cap. Each piece is delicately embroidered with elegant silver trim. The smell of fine tobacco mixed with perfume hangs in the air around him. He smokes a fine black pipe and drinks from a bottle of vintage wine, a rakish grin on his flabby face. He greets characters warmly then rubs his hands together with anticipation. "Let's get down to business, shall we?"


\begin{itemize}
\item Aldor the Immense is a \textbf{djinni} who delights in mortal pleasures and freely travels the world as a dealer of exotic wares. He possesses a sharp wit, a sophisticated sense of humor, puissant skill with money, and a gluttonous appreciation for material indulgences.
\item His elemental nature is completely indiscernible, and he appears in all respects to be a grotesquely large human man. He wears an \textit{amulet of proof against detection and location}.
\item Aldor the Immense deals in magic items and delerium. He is accustomed to small-time heroes attempting to browbeat him, would-be seducers looking to cut a deal, and two-bit thieves trying to rip him off. He is always one step ahead in the mental bargaining game.
\item His attendants and bodyguards are three human assassins named Brill, Meyers, and Vandeerbeek. His oxen are both gorgons.
\item Aldor the Immense has commissioned the Amethyst Academy to create \textit{Arcane Lock} and \textit{Glyph of Warding} to protect his wares. In addition, he has a \textit{tiny chest} in which he keeps his own private treasures.
\end{itemize}

\DndQuote%
{}
{''A fascinating place. Thriving in the shadow of a dangerous backdrop it has become a place to let your guard down and relax even if just a short while...''}
{Pluto Jackson}
\subsubsection{Magic Items for Sale}
Aldor has an enigmatic and ever-shifting stock of magical items that replenishes each week.

\includegraphics[width=0.4\textwidth]{images/../5e.tools/img/adventure/DoDk/040-03-003.obese-toad.png}

\begin{itemize}
\item 2d6 Potion of Healing (100 gp each) and 1d6 Potion of Greater Healing (300 gp each)
\item Spell Scroll from the wizard spell list: 1d10 1st level \textit{spell scrolls} (75 gp each), 1d8 2nd level \textit{spell scrolls} (150 gp each), and 1d6 3rd-level \textit{spell scrolls} (450 gp each). Add twice the cost of any expensive components to the price of the scroll.
\item 1d4 randomly determined uncommon magic items sold for 3d6 x 100 gp each.
\item Material components for spells of 5th level and lower.
\item Use your discretion when generating these magic items. If you wish, occasionally Aldor the Immense might carry a single rare magic item, which he sells for 10,000 gp.
\end{itemize}

\section{Concluding the Chapter}
Allow player characters ample time getting familiar with Emberwood Village. Ensure characters gather a few rumours, make contact with at least two or three of the major faction agents, and encounter one or two groups of rival adventurers. In addition, player characters can take this opportunity to get to know one another more, and perhaps snoop around town for any information related to their personal quests.

Eventually, several townsfolk may begin asking characters if they've been to the ruins yet or tried their hand at collecting delerium shards. Boastful rival adventurers could even set up a wager. If characters are looking for advice, NPCs might recommend characters seek out Gainsbury Expeditionary Supply Company for equipment and information on harvesting delerium, or chat with Fairweather Trades and Exports about selling whatever they manage to dig up. The next chapter, Exploring Drakkenheim, details how characters can explore the ruins, and includes a "Delerium Hunt" to help run their first foray!

\section{Character Advancement}
If your player characters began at 1st level and accompanied Eren Marlowe's caravan, after they've interacted with folk in Emberwood Village and met a few faction members, award them sufficient experience to advance them to 2nd level before they embark on their first foray into Drakkenheim.

\chapter{Exploring Drakkenheim}
\DndDropCapLine{F}{}ollowing an overview of the city itself, this chapter introduces the core mechanics for running urban exploration in Drakkenheim: navigating the city streets and sewers, generating random encounters, and the properties of the mysterious Haze. You'll want to refer to these mechanics throughout the campaign, alongside the appendix sections of this book that detail the rules for Contamination, Delerium, the Random Encounter Tables, and of course, the City Map.


Finally, this chapter closes with an introductory exploration adventure: a Delerium Hunt!

\includegraphics[width=0.4\textwidth]{images/../5e.tools/img/adventure/DoDk/041-04-001.inside-the-crater.png}
\section{City Overview}
Drakkenheim was a grandiose capital city famous for its many architectural wonders. The even-handed rule of House von Kessel fostered a vibrant and industrious spirit that welcomed folk throughout the realm and beyond, and the city was an international hub for trade and culture. Nearly a hundred thousand people called Drakkenheim home, and thousands more lived in small villages and hamlets a few days travel from the city. Only a few hundred people survived the cataclysm that devastated the city, and the folk who lived in the surrounding countryside have abandoned the area for safer regions.

The recorded history of Drakkenheim stretches back about six hundred years, but a settlement and fortification have stood here since antiquity. The city was named after its founders, the now-extinct House von Drakken. Four distinct dynasties have ruled the city since, with House von Kessel claiming the throne one hundred fifty years ago.

Drakkenheim is divided into three main regions:


\begin{description}
\item[The Outer City,.]
estates, suburbs, and slums surrounding the city walls.
\item[The Inner City,.]
protected by the City Walls and accessible through five Gates.
\item[Castle Drakken,.]
a palace and military fortification that lies atop a high crag.
\end{description}

The city itself is further divided by the \textbf{Drann River}. Common parlance calls the two halves the North Ward and the South Ward.

\subsection{Streets}
Five major highways lead to Drakkenheim. The roads transition into the city's main streets: King's Road, Temple Road, Shepherd's Way, Champion's Way, and College Street. Each passes through a correspondingly-named gate, and ranges from twenty to forty feet wide with broad sidewalks along either side. Wrought-iron lamp posts and sewer drains are placed throughout. They converge at Market Square. Elsewhere in the city, the common cobblestone-paved roads are twelve to eighteen feet wide. Most have a narrow two-foot wide brick sidewalk. Sidestreets are often no more than ten feet wide, and in other areas only a narrow alley separates the buildings. In the outer city slums, many roads are merely gravel-covered or raw dirt, though there are some stone-paved streets.

\subsection{Buildings}
Drakkenheim is a city of stone buildings and slate spires. Many in the Inner City may be three or four storeys tall, sometimes taller. Most are tightly packed townhouses built on stone ground levels, although the upper levels are often made from wood, and boast fine glass windows and sturdy chimneys. Wealthier homes even feature indoor plumbing directly connected to the city's water system. Dozens of smaller chapel domes and spires mark distinct boroughs in the city. In contrast, the outer city is mostly a ramshackle slum consisting of crumbling tenements, wattle and daub cottages, and fieldstone shacks. However, the northern estates and the sections along major roads boast fine construction similar to that of the inner city.

\subsection{Drann River}
The Drann River is roughly four hundred feet wide and flows southwest. The water is deep, slow-moving, and thoroughly contaminated. Four stone causeways span the water. Stone quays and piers surround the river within the Inner City.

\subparagraph{Contaminated Water}
All water in Drakkenheim is contaminated, including the Drann River. Characters who end their turn in contaminated water must succeed on a DC 10 Constitution saving throw or gain one level of contamination.

\subsection{City Walls}
Stone masonry walls forty feet high and five to ten feet thick separate the Inner City and the Outer City, with entrance via five gates. These fortifications present a major obstacle when exploring the city: the five city gates are occupied by faction forces or monsters, and hostile gargoyles attack any who scale the walls. Refer to the "Walls of Drakkenheim" in Chapter 7: Inside the Walls of Drakkenheim for details.

\subsection{Castle Drakken}
Atop a high cliff on the north side of town perches Castle Drakken. The fortification and palace is described in chapter 9.

\section{The Haze}
The Haze is magical radiation emitted by the vast concentration of delerium in Drakkenheim. It permeates the ruins and extends about two miles outside the walls.

\subsection{Effects on Characters}
Creatures do not gain any benefits from finishing a long rest within the Haze.

Creatures may remain in the Haze for up to twenty-four hours before they must finish a long rest to recuperate from exposure to its eldritch energies. For each additional hour spent within the Haze beyond twenty four hours, a creature must succeed on a DC 15 Constitution saving throw or gain one level of contamination.

Most monsters encountered in Drakkenheim possess the \textit{fully contaminated} trait (see Appendix A) and ignore these effects.

\subsection{Environmental Effects}
At dawn, the Haze magically manifests mist throughout the ruins. Characters can see normally through the mist up to 150 feet. Vision beyond is \textit{lightly obscured}, but vision past 300 feet is totally obscured. Any mists dispersed during the day reform ten minutes later.

The mists dampen sunlight. Creatures with \textbf{sunlight sensitivity} or similar traits do not suffer these penalties during the day while within the Haze.

The mist disperses shortly after sunset, but the Haze appears as a shifting corona of octarine light shimmering over the city during the night. Floating motes and hazy particles hang over the city streets providing illumination akin to moonlight, and delerium crystals emit a bright glow to a range of five feet.

The corpses of beasts and humanoid creatures do not rot within the Haze. However, most unattended corpses left in the streets are devoured within a day by ravenous ratlings. Other food and drink spoils after 2d6 days within the Haze.

\subsection{Effects on Spells}
The erratic magic of the Haze disrupts many spells cast within it. Any teleportation spell cast by a creature outside the Haze targeting a creature or location within the Haze automatically fails and the spell slot is wasted. Creatures inside the Haze may teleport normally within, but any attempt to leave the Haze via teleportation fails. Telepathic communications or effects such as the \textit{sending} spell transmitted from outside the Haze fail to contact a character within it, and vice-versa. Sensors created by divination spells can't appear inside the Haze, and those created outside the Haze can't enter it. Finally, other divination spells fail to reveal useful information about delerium, the Haze, the origins of the meteor, or any events that occurred within the Haze.

When a character attempts any of the above using a spell or other effect, they must succeed on a DC 15 Intelligence saving throw. On a failure, they take 6d6 psychic damage and become incapacitated with madness until they finish a long rest. During this time they speak only in gibberish. A \textit{greater restoration} spell ends this effect.

\textit{Rope trick}, \textit{Leomund's Tiny Hut} and similar spells or abilities do not provide shelter or protection from the Haze. However, characters may rest within a \textit{Mordenkainen's Magnificent Mansion} spell cast within the Haze.

\subsection{Deep Haze}
Certain areas are so suffused by the Haze that mere exposure is extremely dangerous and potentially fatal. Known as the Deep Haze, these areas have the following additional traits:


\begin{itemize}
\item Deep Haze is instantly recognized by a thick prismatic fog. Characters can see normally through this fog up to 60 feet. Vision beyond is \textit{lightly obscured}, but vision past 120 feet is totally obscured. If dispersed during the day, the fog reforms one minute later.
\item Characters entering the Deep Haze must succeed on a DC 15 Constitution saving throw, and again for each full hour spent within. On a failed saving throw, they suffer 10 (3d6) \textit{necrotic damage} and gain one level of contamination.
\end{itemize}

The Haze map indicates which areas of the city are covered by the Deep Haze. Other locations suffused by the Deep Haze are noted in location descriptions.

\subsection{Delerium Sludge}
Gargantuan delerium geodes perspire an opalescent sludge. Characters submerged in delerium sludge make a DC 18 Constitution saving throw at the start of their turn. They take 42 (12d6) \textit{necrotic damage} on a failed saving throw, and half as much on a successful one. In addition, characters who fail the saving throw gain one level of contamination.

\section{Navigating the Ruins}
Drakkenheim behaves more like a dangerous wilderness than an urban location. Characters risk perilous random encounters with hostile monsters and mysterious magical phenomena as they traverse the ruins, but may also stumble upon delerium, lost treasures, and faction agents.

\includegraphics[width=0.4\textwidth]{images/../5e.tools/img/adventure/DoDk/042-map-4.01-city-street-and-sewers--map.png}
\subsection{Getting to the City}
Aside from Emberwood Village, all towns and villages within one week's travel of Drakkenheim have been wholly abandoned. The nearest large cities are a three-week journey away.

Travel times from safe havens are indicated below. When travelling to the city, characters may arrive anywhere on the city outskirts.


\begin{DndTable}{Xcc}

Origin Point & Distance & Travel Time on Foot\\
Eckerman Mill & 1 mile & 20 minutes\\
Camp Dawn & 2 miles & 1 hour\\
Emberwood Village & 5 miles & 2 hours\\
\end{DndTable}

\subsection{Urban Exploration}
As the player characters traverse the ruined city, you can handle urban exploration in an abstract manner as follows:


\begin{itemize}
\item Using the city map, identify approximately where the characters are currently located.
\item Let the players determine what direction they want to go and choose their travel pace (see "Moving Through the Streets" below).
\item Check for random encounter each hour (see below for tables).
\end{itemize}

It's not necessary to precisely track which streets the characters take; a rough route is sufficient. Correspondingly, navigation checks are not required. Characters can locate key adventure sites without much difficulty, as long as they have approximate directions and a general description of their destination. Instead, the random encounter tables include several situations in which the characters might make a wrong turn or become lost.

\subsubsection{Moving through the Streets}
Ruined buildings and maze-like streets are exceptionally \textit{difficult terrain}. How fast characters can move on foot depends on their priority:


\begin{description}
\item[Fast.]
At this pace, characters advance one mile per hour. They take the most direct route to their destination and move openly on major streets. Characters cannot use stealth, and make Perception checks with disadvantage. When checking for random encounters (see "Random Encounters" later in this chapter), each character rolls twice and takes the lower result.
\item[Normal.]
At this pace, characters advance one-half mile per hour. They proceed carefully through the ruins staying vigilant for threats and landmarks. Characters cannot use stealth, but may make Investigation and Perception checks with advantage.
\item[Slow.]
At this pace, characters advance one-quarter mile per hour to evade hostile threats. They sneak down back alleys and avoid major streets. Characters may use stealth, and make Perception checks normally.
\end{description}

\subsubsection{Searching the Ruins}
Use these guidelines whenever the characters deliberately search for something in the ruins, such as scavenged supplies, lost treasures, unknown locations, a hidden foe, or a safe place to rest. Conducting a good search takes at least one hour; during which the characters may survey an area with up to a ¼ mile diameter. After one hour of searching, each character makes an Investigation or Wisdom (Survival) check (DC 10 for the Outer City, DC 15 for the Inner City). Exceeding the DC by 5 or more grants an additional success. Whether or not the characters find anything is determined by the total number of successes:


\begin{description}
\item[3 or more successes.]
The party finds what they are looking for if it can be found within the search area, otherwise, they conclusively rule out any possibility that what they were seeking was in the area of their search.
\item[2 or more failed checks.]
The characters stumble into a Random Encounter, but if they still had enough successes, they also find whatever they sought.
\end{description}

If your party is larger or smaller than four characters, increase or decrease the number of required successes proportionately. Alternatively, characters might use special skills when searching for something specific:


\begin{itemize}
\item Arcana. Delerium deposits or magical phenomena.
\item \textit{Cartographer's Tools}. Searching an area the characters have visited at least once before themselves, or for which they have a map.
\item History. Specific and notable monuments, buildings, or residences.
\item Nature. Exploring natural environments such as Queen's Park.
\item Survival. Tracking monsters or people.
\item Stealth and Perception checks do not contribute to successes, but may be used to evade or detect hostile threats while the characters search the ruins.
\end{itemize}

\subsubsection{Finding Delerium Deposits}
When characters seek out delerium by \textbf{Searching the Ruins} (described above), success requires a sucessful Intelligence (Arcana or Investigation) or Wisdom (Survival) check with a DC indicated by the table below. A check result which exceeds the DC by 5 or more grants an additional success.

\begin{DndTable}[header=Outer City Delerium Deposits (DC 15)]{cX}

Number of Successes & Result\\
2 success or fewer & Nothing\\
3 successes & 10 (3d6) delerium chips worth 10 gp each\\
4 success & 3 (1d6) delerium fragments worth 100 gold each\\
5 or more successes & 1 delerium shard worth 500 gp\\
\end{DndTable}

\begin{DndTable}[header=Inner City Delerium Deposits (DC 20)]{cX}

Number of Successes & Result\\
1 success or fewer & Nothing\\
2 successes & 10 (3d6) delerium chips worth 10 gp each\\
3 successes & 7 (2d6) delerium fragments worth 100 gold each\\
4 successes & 3 (1d6) delerium shards worth 500 gp\\
5 or more successes & 1 delerium crystal worth 1,000 gp\\
\end{DndTable}

Results are cumulative, so a group that makes four successes finds 3d6 delerium chips and 1d6 fragments. As normal, generate a random encounter if characters have two or more failed check results while searching for delerium.

Characters searching within a quarter mile of the Crater's Edge receive one automatic additional success due to the rich quantities of delerium found in this exceptionally dangerous region.

\subsubsection{Drakkenheim Sewers}
The Drakkenheim sewers were a sophisticated engineering marvel. The system combined underground gravity-based aqueducts and drainage canals with carefully placed locks and floodgates to deliver clean and fresh water to the city population and draw away the massive quantities of sewage and wastewater. Major canals run under the city streets, and can be accessed from the surface by numerous maintenance hatches.

After over a decade of disuse, most tunnels are flooded with contaminated water and many other passages have collapsed. Other sections are now monster lairs. As a result, travelling in the sewers is extraordinarily hazardous and ill-advised. However, characters might consider taking the sewers to locate a route under the city walls, or to gain access to hidden underground dungeons such as the cistern (see chapter 7) or the Court of Thieves (see chapter 8).

\subparagraph{Sewer Navigation}
Handle sewer navigation abstractly using the same mechanics for navigating the streets above, except characters move at half their normal pace through the twisting tunnels. Sewer passages and access points could be almost anywhere in the city, so use the surface map as an approximate guide to track the character's progress. In addition, after each hour spent in the sewers characters must succeed on a DC 10 Constitution saving throw or gain one level of contamination from contact with sewer water.

\subparagraph{Searching the Sewers}
Just as characters may search the streets, they might search the sewers for hidden entrances to some adventure sites. Other secret locations may be found in this manner. A DC 20 check result is needed for success when searching the sewer passages.

\subsubsection{Other Exploration Activities}
\subparagraph{Taking a Rest}
Before taking a short rest in the ruins, characters may spend one hour locating a safe spot or hiding their campsite. If they do, don't check for random encounters during their rest. Note that the Haze (described above) prevents characters from benefiting from a long rest in the city. Characters who attempt to do so anyways check for random encounters each hour.

\subparagraph{Flying}
Flying characters may soar across the city rapidly. At your discretion, they may risk hostile attention from the gargoyles defending the walls or the harpies inhabiting the Inner City.

\subparagraph{Vehicles}
Land vehicles are largely impractical in the ruins, although characters could use a boat to travel via the river normally (though at your discretion characters may risk attack from aquatic or sewer-dwelling monsters!) Characters moving with a land vehicle or portaging a water vehicle travel ¼ mile per hour. Land and water vehicles can't be hidden without magic, but individual characters might be able to use stealth.

\subparagraph{Finding Sustenance}
Finding suitable sustenance within Drakkenheim is difficult since most food and water within the city is either spoiled or contaminated. Instead, characters should bring their own supply of food and water from outside the city. Nevertheless, a character who eats or drinks anything found in Drakkenheim that has not been magically purified makes a DC 15 Constitution saving throw. On a failed save, they gain one level of contamination.

\subsubsection{A Combined Approach}
Drakkenheim is dangerous and unpredictable. There is no such thing as a foolproof plan for safely exploring the ruins. Nevertheless, characters using special resources, spells, magic items, or with a very inventive plan might greatly improve their navigation. For example, characters could use spells such as \textit{invisibility} and \textit{pass without trace} together to travel quickly while remaining unseen, or perhaps \textit{water walk} helps characters move over the river. If they aren't expending resources, you should call for a few difficult skill checks. On a failure, they stumble into a difficult random encounter with extra complications.

\subsubsection{Escaping the Ruins}
Characters must use their resources carefully while exploring Drakkenheim, as there is always a chance they'll face a random encounter or two when \textit{leaving} the ruins as well. At some point during their adventures, characters might be heavily wounded and out of resources deep in the city ruins. As they stagger back to Emberwood Village, they roll up a difficult random encounter. This could be the one that ends it for them! Be lenient the first time this happens, but merciless the second.

\includegraphics[width=0.4\textwidth]{images/../5e.tools/img/adventure/DoDk/043-04-002.entering-the-sewers.png}
\section{Random Encounters}
Dangerous monsters lurk everywhere in Drakkenheim.

While characters navigate the ruins, ask each player to roll 1d20...


\begin{itemize}
\item For each hour spent exploring the city (as directed above)
\item Any time you feel the players are being indecisive or idle in a dangerous environment
\item Whenever the players are having an easy time and you need to shake things up
\item When directed by specific circumstances described within an adventure site
\end{itemize}

There is an encounter if \textit{any} player rolls a natural one. Consult the random encounter tables below to determine what the characters encounter based on their location in the city. Some locations have their own unique random encounter tables.

You may wish to invent additional complications to the encounter for each additional one rolled, such as an ambush, additional monsters, unfavourable terrain, or erratic magical effects.

Conversely, rolling a natural 20 might add a bit of good fortune for the players, such as an opportunity to get the drop on foes or avoid them, a delerium deposit, monsters carrying treasure or a magic item, nearby allies who can help, or other advantageous circumstances for the characters.

\subparagraph{If every player rolls a one}
What's the worst that can happen? It happens!

Finally, you can adjust the die size rolled when players check for random encounters, such as using a 1d10 or even 1d6. The odds of a random encounter happening increase with lower die sizes, so this is ideal if you have a smaller group or want the characters to face more random encounters.

\subsection{Random Encounter Tables}
\begin{DndTable}[header=Sewers]{cX}

1d20 & Encounter\\
1 & Horribly Lost\\
2 & Going in Circles\\
3 & Wrong Turn\\
4 & Gloaming Ray\\
5 & Living Ruins\\
6 & Troll Traveller\\
7 & Sewer Monster\\
8 & Drowned Dead\\
9 & Ratling Scavengers\\
10 & Stalking Vermin\\
11 & Hateful Dead\\
12 & Lost Ones\\
13 & Gibbering Flesh\\
14 & Dozing Hulk\\
15 & Shadows of Drakkenheim\\
16 & Fallen Heroes\\
17 & Cleaning Cube\\
18 & Deep Ones\\
19 & Rival Adventurers\\
20 & Double Trouble\\
\end{DndTable}

\begin{DndTable}[header=Outer City]{cX}

d100 & Encounter\\
1-4 & Run out of Town\\
5-8 & Uninvited Guests\\
9-12 & Wrong Turn\\
13-16 & Troll Traveller\\
17-20 & Menacing Manticore\\
21-24 & Shadows of Drakkenheim\\
25-28 & Phase Webs\\
29-32 & Watching Gargoyles\\
33-36 & Ratling Scavengers\\
37-40 & Garmyr Hunters\\
41-44 & Hateful Dead\\
45-48 & Gibbering Flesh\\
49-52 & Ghost Lights\\
53-56 & Haze Haunt\\
57-60 & Angry Furniture\\
61-64 & Lost Ones\\
65-68 & Shambling Husks\\
69-72 & Old Alchemist's Shop\\
73-76 & Hooded Lantern Patrol\\
77-80 & Queen's Men Looters\\
81-84 & Questing Knight\\
85-88 & Academy Surveyor\\
89-92 & Pilgrims of the Falling Fire\\
93-96 & Rival Adventurers\\
97-100 & Double Trouble\\
\end{DndTable}

\begin{DndTable}[header=Inner City]{cX}

d100 & Encounter\\
1-3 & Horribly Lost\\
4-5 & Executioner's Summons\\
6-7 & Going in Circles\\
8-9 & Lord of the Feast!\\
10-11 & Uninvited Guests\\
12-13 & Crimson Countess\\
14-15 & Wandered into the Garden\\
16-17 & Sewer Monster\\
18-20 & Stalking Vermin\\
21-23 & Ratling Raiders\\
24-26 & Hateful Dead\\
27-30 & Lost Ones\\
31-32 & Troll Traveller\\
33-34 & Phase Webs\\
35-37 & Gibbering Flesh\\
38-39 & Old Alchemist's Shop\\
40-41 & Menacing Manticore\\
42-43 & Shadows of Drakkenheim\\
44-45 & Ghost Lights\\
46-47 & Haze Haunt\\
48-50 & Watching Gargoyles\\
51-54 & Garmyr Hunters\\
55-57 & Garmyr Ravagers\\
58-59 & Overgrown Ruin\\
60-62 & Living Ruins\\
63-64 & Chimera Nest\\
65-68 & Harpy Flock\\
69-70 & Lurking Wraiths\\
71-72 & Fallen Heroes\\
73-74 & Gloaming Ray\\
75-77 & Deep Ones\\
78-81 & Rival Adventurers\\
82-84 & Pilgrims of the Falling Fire\\
85-88 & Hooded Lantern Patrol\\
89-92 & Queen's Men Looters\\
93-94 & Questing Knight\\
95-96 & Academy Surveyor\\
97-100 & Double Trouble\\
\end{DndTable}

\subsubsection{Common Locations}
Use these descriptions to set the scene for a random encounter.


\begin{DndTable}{cX}

1d10 & Description\\
1 & Breakaway meteors demolished this area, leaving only broken cobblestones and toppled ruins riddled with small craters.\\
2 & A makeshift encampment. There's a canvas tent, sleeping bags, and a few improvised barricades or defences. Embers still burn in the campfire.\\
3 & A deserted alleyway filled with trash, access point to the sewers.\\
4 & A devastated street filled with broken wagons, abandoned vendor's stalls, and bones.\\
5 & An open plaza or park around a small monument, fountain, or well.\\
6 & A courtyard leading up to a civic building, shrine, chapel, or estate.\\
7 & A cluster of ramshackle taverns.\\
8 & Seedy tenement buildings and rundown shacks.\\
9 & Several workshops surrounded by warehouses.\\
10 & Rows of dilapidated and heavily damaged townhouses.\\
\end{DndTable}

\subsubsection{Warped Ruins}
Many buildings in Drakkenheim are not exactly normal anymore. Use these and the Arcane Anomalies in Appendix D to inject weirdness and horror as characters explore the city.


\begin{DndTable}{cX}

1d10 & Description\\
1 & Stonework and timbers have transformed into solid glass.\\
2 & Cloaked in eldritch fire, but it's not burning down.\\
3 & Blown apart, but frozen in time at the middle of the explosion.\\
4 & Flesh, limbs, eyes, mouths, and faces are merged into the walls. They babble incoherently.\\
5 & The side facing the crater is melted like sticky wax.\\
6 & Disembodied screams and whispers resound from within.\\
7 & All colour within becomes black and white\\
8 & An illusion of a childhood memory briefly appears in the front door, then vanishes.\\
9 & Looking at the building causes déjà vu.\\
10 & Scorch marks in the shape of human silhouettes are found on the walls and floors.\\
\end{DndTable}

\subsubsection{Lucky Finds}
After a random encounter, roll 1d20 to see if the characters pick up any useful loot.


\begin{DndTable}{cX}

1d20 & Loot\\
1-8 & Nothing but broken refuse and ruined scraps of weapons and armour.\\
9 & 1d4 sets of tools in an abandoned workshop.\\
10 & 5d6 gold pieces in a coin purse clutched by a severed hand.\\
11 & 1d6 Potion of Healing in a hidden cache.\\
12 & 1d4 uncommon Spell Scroll tucked in a scroll case by a wizard's corpse.\\
13 & 1d4 art objects worth 25 gp each are found in a modest townhouse.\\
14 & 2d6 delerium chips haphazardly stored in a backpack.\\
15 & 2d6 x 10 gp in a lockbox in an old merchant's tradehouse.\\
16 & 1 Keoghtom's Ointment upon an otherwise empty shelf.\\
17 & 1d4 delerium fragments carefully placed in glass containers.\\
18 & 1d4 Potion of Greater Healing in a forsaken shrine.\\
19 & A crumbling estate has a piece of artwork worth 250 gp inside.\\
20 & 1 rare \textit{spell scroll} in a pile of books and loose pages fluttering down the street.\\
\end{DndTable}

\subsection{Random Encounter Descriptions}
\subparagraph{Academy Surveyor}
An Amethyst Academy \textbf{mage} and 1d4 \textbf{hedge mage} apprentices are conducting research in the ruins.

\subparagraph{Angry Furniture}
2 (1d4) mimics are disguised as perfectly normal pieces of furniture in a nearby house.

\subparagraph{Chimera Nest}
A \textbf{chimera} is nesting in a tall spire. Aggressive and territorial, it swoops down to rip apart any who wander around its turf.

\subparagraph{Cleaning Cube}
A \textbf{gelatinous cube} awaits down a narrow corridor in the sewers. It stays completely still until a creature enters its space.

\subparagraph{Crimson Countess}
The \textbf{Crimson Countess} hunts above with a retinue of 5 (2d4) harpies.

\subparagraph{Dozing Hulk}
A single \textbf{haze hulk} lays sleeping across the narrow road. Squeezing past without waking it may prove difficult, but it clutches a delerium shard in its hand.

\subparagraph{Deep Ones}
10 (3d6) aquatic delerium dregs and 1 \textbf{chuul} occupy a sewer junction or are emerging from a sewer exit.

\subparagraph{Double Trouble}
Roll twice, ignoring this result on subsequent rolls. If two groups of monsters or NPCs are encountered, they're fighting each other.

\subparagraph{Drowned Dead}
A one-hundred-foot long corridor in the sewers is completely flooded. The tunnel has 10 (3d6) haze husks floating in it that appear lifeless unless attacked or until characters are halfway down the tunnel. The undead attempt to hold characters in place until they drown.

\subparagraph{Executioner's Summons}
The characters stumble into Slaughterstone Square.

\subparagraph{Fallen Heroes}
5 (2d4) haze wights are all that remains of this former adventuring party. They believe any living humanoids they encounter are mutated monsters.

\subparagraph{Garmyr Hunters}
10 (3d6) garmyr leading 1d4 worgs are stalking the streets for fresh meat. They are on keen lookout and watch for any signs of movement. Double the number encountered in the Inner City.

\subparagraph{Garmyr Ravagers}
7 (2d6) garmyr berserkers with 2 (1d4) hell hounds are rampaging through the streets, howling loudly and starting fires.

\subparagraph{Ghost Lights}
2 (1d4) will-o-wisps attempt to lure passing characters towards a mirage of treasure or delerium in the Haze. In the inner city, there are 2d6 wisps instead.

\subparagraph{Gibbering Flesh}
The walls here are covered in thick fleshy growths, parts of the walls, streets, and rubble have grown eyes and mouths that break away and form 2 (1d4) gibbering mouthers. In the Inner City, it has become a \textbf{protean abomination}.

\subparagraph{Gloaming Ray}
A \textbf{cloaker} flies through the dark rooftops above looking for prey.

\subparagraph{Going in Circles}
The characters lose their position and get lost. They no longer know where they are anymore or what direction they are facing. They must wander the city streets for the next 1d4 hours, after which they regain their original position. Check for random encounters as normal each hour. A random encounter could result in circumstances where the characters no longer become lost.

\subparagraph{Harpy Flock}
A pack of 10 (3d6) harpies flies overhead surveying the ground for anything worth hunting. They perch on rooftops and watch the streets below, trying to lure wandering characters up to the roofs to push them off to their deaths

\subparagraph{Hateful Dead}
A lone \textbf{haze wight} marches with 10 (3d6) haze husks. It commands its minions to attack any living creature they find, and raises any nearby corpses as reinforcements. Add 1d4 additional haze wights when encountered in the Inner City.

\subparagraph{Haze Haunt}
A \textbf{warp witch} moans and wails in the nearby building, scornfully mourning its miserable existence as it tries to remember its past. Add an additional 1d4 warp witches when encountered in the Inner City.

\subparagraph{Hooded Lantern Patrol}
5 (2d4) scouts led by an \textbf{urban ranger} are on a recon mission. If the characters are in trouble they step in to help them, but then demand the characters turn over half of whatever plunder they've found in the ruins, citing their authority under the "law" of Westemär and that "scavenging" in the ruins is technically prohibited without the assent of the Lord Commander.

\subparagraph{Horribly Lost}
The characters are badly turned around and become hopelessly lost in the ruins. They no longer know where they are anymore or what direction they are facing. They wander the city streets indefinitely, checking for random encounters as normal each hour. They don't regain their bearings unless the circumstances of a random encounter lead them to clues or friendly NPCs who can help them.

\subparagraph{Living Ruins}
Several inanimate parts of the ruins suddenly spring to life as 2 (1d4) contaminated elementals. Choose whichever type is most appropriate for the area.

\subparagraph{Lord of the Feast}
The \textbf{Lord of the Feast} leads a warpack of 10 (3d6) garmyr and 2 (2d4) hell hounds.

\subparagraph{Lost Ones}
10 (3d6) delerium dregs wander the streets ahead. They sorrowfully mutter nonsensical gibberish, but wail and screech when they encounter humanoids. Add 2 (1d4) haze hulks when encountered in the Inner City.

\subparagraph{Lurking Wraiths}
2 (1d4) arcane wraiths flutter about seeking erratic magic, and fixate on whichever character is carrying the most delerium or magic items.

\subparagraph{Menacing Manticore}
A \textbf{manticore} circles overhead looking for an easy meal. If it spots the characters, it swoops in to attack the most vulnerable member. It flies off towards the inner city if reduced to less than half its hit points. Manticores in the Inner City may hunt in packs of 5 (2d4).

\subparagraph{Old Alchemist's Shop}
A decrepit alchemist's shop that reeks of a chemical odor stands on a street corner. A pool of spilled chemicals has become an aggressive \textbf{ochre jelly}. In the Inner City, 2 (1d4) black puddings are encountered instead.

\subparagraph{Overgrown Ruin}
Strange alien plants and oddly-shaped vines creep up a crumbling ruin. A \textbf{shambling mound} and 2 (1d4) hypnotic eldritch blossoms grow in the tangled mass.

\subparagraph{Phase Webs}
Strange webs cover this section of the ruins and fill nearby buildings - some even span the streets themselves. 2 (1d4) phase spiders resting in the dark corners come to investigate any disturbance of their webs. The webs are especially thick in regions of the Inner City, where there are 7 (2d6) spiders instead.

\subparagraph{Pilgrims of the Falling Fire}
A group of 10 (3d6) cultists led by a \textbf{cult fanatic} and a \textbf{berserker} are heading towards the crater.

\subparagraph{Queen's Men Looters}
7 (2d6) bandits led by a \textbf{bandit captain} are hoping for easy pickings. The bandits try to hide from the players until they appear wounded or in need of a rest, so canny characters might spot them before this happens.

\subparagraph{Questing Knight}
A Silver Order \textbf{knight} on a \textbf{warhorse} with a retinue of 7 (2d6) guards are searching for lost relics and holy sites in the ruins.

\subparagraph{Ratling Raiders}
A \textbf{ratling} \textbf{warlock of the rat god} leads a great horde of 21 (6d6) ratlings.

\subparagraph{Ratling Scavengers}
10 (3d6) ratlings hide down a nearby alley or passage awaiting unsuspecting prey.

\subparagraph{Rival Adventurers}
Characters encounter a group of rival adventurers (see chapter 2).

\subparagraph{Run out of Town}
Characters get turned around badly, and arrive at the nearest edge of town.

\subparagraph{Sewer Monster}
An \textbf{otyugh} feasting on offal lurks near a sewer access point.

\subparagraph{Shadows of Drakkenheim}
5 (2d4) shadows slowly follow the characters. They try to remain hidden and follow the party until one of the player characters is more than ten feet away from the others, then the shadows strike in an attempt to overwhelm the straggler. An \textbf{arcane wraith} leads these shadows when encountered in the Inner City.

\subparagraph{Shambling Husks}
7 (2d6) haze husks stumble and shuffle around the streets. In this unsettling gait, they play out scenes of their former everyday lives.

\subparagraph{Stalking Vermin}
14 (4d6) ratlings and 5 (2d4) ratling guttersnipes have set up an ambush. These ratlings use bait to lure prey towards a dead-end street or sewer passage where the guttersnipes lurk on high ground to snipe their quarry with ranged attacks.

\subparagraph{Troll Traveller}
A \textbf{troll} is heading through the ruins to join the trolls at King's Gate. He wears a large pack containing raw contaminated meat, a tankard of rancid ale, 11 (2d10) gold, and a few sets of tools and trinkets. If encountered near a bridge in the Inner City, there are 3 trolls instead.

\subparagraph{Uninvited Guests}
The characters arrive at the nearest city gate.

\subparagraph{Wandered into the Garden}
The characters arrive at the nearest edge of Queen's Park Garden.

\subparagraph{Watching Gargoyles}
2 (1d4) wall gargoyles roaming overhead swoop down to attack. If encountered in the Inner City, there are 7 (2d6) instead.

\subparagraph{Wrong Turn}
The characters find themselves back where they started an hour ago. I guess you took a wrong turn somewhere?

\DndQuote%
{}
{''Like trying to breathe on a scorching day, the air burns, eyes water. I constantly wish I was taking a bath, I feel dirty just being in it.''}
{Sebastian Crowe}
\section{Delerium Hunt}
\DndQuote%
{}
{''Prove to us you can survive in the ruins, and then maybe we'll see if we can work together.''}
{}
In this short introduction to exploring Drakkenheim, player characters search the ruins for delerium and other plunder they can sell in Emberwood Village. Along the way, they'll face a few random encounters in the ruins against monsters, rival adventurers, and faction patrols.

\includegraphics[width=0.4\textwidth]{images/../5e.tools/img/adventure/DoDk/044-04-003.drakkenheim-destroyed.png}
\subsection{Adventure Hook}
\subparagraph{Prospecting and Plunder}
This fundamental hook is simple: the characters head into the Outer City Ruins on their own initiative to seek out whatever treasure and delerium they can find. Allow characters to take the reigns and plan their route. If for some reason characters don't think to do this on their own, a friendly NPC contact in Emberwood Village or even a faction lieutenant suggests the characters experience the ruins firsthand and prove they can survive the dangers there. It's risky, but there's a chance to make good money!

\subsection{Events}
Use the guidelines throughout this chapter as characters explore the ruins. Finding a delerium deposit in the Cratered Street requires characters successfully search the ruins once as described above. Check for random encounters normally, but use these events along the way:

\subsection{Approaching Drakkenheim}
If this is the first time the characters are venturing into Drakkenheim, read:

\begin{DndReadAloud}{}
Nestled within a pine-covered valley, Drakkenheim is a walled metropolis that straddles the wide and sluggish Drann River. A high rocky hill rises above the city, upon which perches Castle Drakken. Below the imposing fortress, cobblestone streets weave around densely-packed stone buildings and slate-shingle spires. The city spills from the outer walls into a ruined urban sprawl.

The great dome of Saint Vitruvio's Cathedral stands prominently southwest of the castle. At the heart of the city lies the Cosmological Clocktower of Market Square, and to the northwest is the Inscrutable Tower of the Amethyst Academy. The soaring obsidian-glass building is broken in half at the middle, but the top still hovers above in defiance of gravity.

Four bridges reach across the Drann River like desperate limbs; one is covered in massive stone statues of the heroes of old, and another lurches under the weight of the ramshackle homes and other buildings constructed upon its length.

Billowing clouds of midnight purple and dusky violet mask the sky overhead, casting the city in a dismal gloom. In the rare moments where the clouds part, stray sunbeams scornfully shine upon the city. There hasn't been a clear day in Drakkenheim in fifteen years.

\end{DndReadAloud}

\subsubsection{Entering the Ruins}
\begin{DndReadAloud}{}
The mist-shrouded countryside gives way to crumbling buildings and twisting mud-filled streets. Entering the desolate city ruins, you pass the bloated corpse of a long-dead horse. The smell of ozone fills the cold air as you walk through the tattered rubble of the Outer City. Once in a while you hear a barking dog, then perhaps a scream carried on the wind.

\end{DndReadAloud}


\begin{itemize}
\item If the player characters trigger a random encounter, instead of rolling on the Outer City table, use the Ratling Scavengers or Lost Ones results.
\end{itemize}

\subsubsection{Cratered Street}
\begin{DndReadAloud}{}
This blasted street is filled with small craters.

Tiny crystals glow with prismatic light within each.

\end{DndReadAloud}


\begin{itemize}
\item Characters find this location once they successfully search the ruins and score two or more successful check results.
\item There are 6 haze husks stumbling aimlessly around the craters.
\item Characters can extract 10 (3d6) delerium chips worth 10 gp each and 3 (1d6) delerium fragments worth 100 gp each. This takes about five to ten minutes.
\end{itemize}

\subsubsection{Disputed Claims}
A group of rival adventurers arrives on the site as the characters extract the delerium fragments.


\begin{itemize}
\item They are a \textbf{noble}, a \textbf{scout}, and a \textbf{hedge mage}. Tailor their identities and personalities to appropriately reflect your player characters. This opportunistic group tries to browbeat characters into giving up the claim.
\item If characters take too long deciding what to do about the haze husks, or take more than four hours to find the deposit, you may decide the rival adventurers have already claimed the first deposit (there are no haze husks present in this case). When the characters arrive, the rivals are extracting the delerium and gloat over the find. Characters who decide to move on to a second deposit encounter the haze husks as normal.
\item The rival adventurers have had some minor success so far. They have a \textit{potion of healing}, 35 (10d6) gold, 1d6 delerium chips, and their equipment.
\end{itemize}

\subsubsection{Escaping the Ruins}
Once the characters have found their prize, don't forget to check for random encounters as they leave the city! Only then may they pawn their finds in Emberwood Village.

\subsection{Character Advancement}
Once your characters complete their first successful Delerium Hunt, award them sufficient experience to advance them to 3rd level.

\subsection{Expeditions Onward}
Take note of which faction agents the player characters interacted with during their early interactions into Emberwood Village.

Who did they get along with? Who did they impress? Who did they rub the wrong way? Whichever faction lieutenant(s) had a good rapport with your player characters approach them shortly after the characters return from their first successful expedition to the ruins. The faction lieutenant then offers the characters a mission to one of the locations detailed in chapter 6 - we recommend River or Ansom send characters to the Rat's Nest, and Nathaniel, Ophelia, or Blackjack Mel can send characters to the Chapel of Saint Brenna.

Beyond this point, the paths the characters might take are in their hands. Use the rumours, Adventure Hooks provided for each location, and the dossiers of the factions to navigate their choices through the Outer City, advancing to the Inner City, and eventually to Castle Drakken.

\chapter{Outside the Walls}
\DndDropCapLine{A}{}ll manner of treasure-seekers have made their fortunes and met their fates in the Outer City. These crumbling ruins surround the city walls like huddled masses desperate for shelter. Describe the abandoned feeling of the ruined Outer City as follows:


\begin{DndReadAloud}{}
Travelling down the decrepit thoroughfares of Drakkenheim, you pass several buildings which have completely collapsed into rubble. Those which still stand are little more than condemned wretches teetering upon the brink of ruin. A rare few are undamaged, but still bear obvious signs of disuse and exposure: shattered windows, peeling paint, doors rotted off the hinges, and wet sagging rooftops. The cracked cobblestones are littered with abandoned refuse; broken wagons, litter, and an occasional bloodstain or scattered bones. Here and there, tiny motes of octarine glint within the light mist which hangs over the streets. For now, the streets are cold and deadly quiet.

\end{DndReadAloud}

\DndQuote%
{}
{''Magic is all fun and games until someone randomly turns into a potted plant...actually that is kind of fun too.''}
{Veo Sjena}
\section{Overview}
The locations in the Outer City make ideal early explorations into Drakkenheim. However, characters who have not reached at least 3rd level may find these locations quite challenging. The party need not visit every location detailed herein, but those who do may reach up to 6th level exploring the Outer City. Characters may visit these locations in any order, but the "Rat's Nest" and the "Chapel of Saint Brenna" are ideal locations to begin.

\subparagraph{Adventure Hooks}
Each location in this section describes a variety of adventure hooks you can use to prompt characters to explore the site: most involve one of the faction lieutenants approaching the party. See "Making Contact with the Factions" in Chapter 4: Emberwood Village for a summary of the faction lieutenants and where they can be found within the town.

\subparagraph{City Walls and Gates}
The party should explore the Outer City before taking on the Inner City. Nevertheless, if characters are determined to make their way into the Inner City, they can certainly try! Should the party approach any of the city gates or attempt to scale the city walls, consult Chapter 7 "Inside the Walls of Drakkenheim".

Once characters have gained experience exploring the ruins and developed a relationship with one or more of the factions, they'll be ready to tackle the Walls of Drakkenheim and begin exploring the Inner City locations (detailed in chapter 7).

\section{Black Ivory Inn}
\includegraphics[width=0.4\textwidth]{images/../5e.tools/img/adventure/DoDk/045-05-001.black-ivory-inn.png}
\DndQuote%
{}
{''Listen, I know a guy, my cousin Mike used to run a smugglers tunnel underneath an old inn on the outskirts of town, leads ya right into the city. Now, I hear some people saying the old Black Ivory place is up and running again? Maybe worth seeing if that smugglers tunnel is still operational for an easy way in and out of the city!''}
{Blackjack Mel}
The Black Ivory Inn stands utterly pristine; apparently spared from the ravages of the meteor. Warm light, sweet-smelling wine, mirthful laughter, and stunningly beautiful piano melodies emanate from within and waft through the ruined streets nearby. However, the legendary music hall and hotel is the site of a bizarre arcane anomaly that has warped the fabric of time and space. Until the malignant abomination anchoring this twisted magic is destroyed, visitors to the Black Ivory Inn are trapped within a time loop: doomed to forever play out the last four hours and thirty-three minutes before the meteor struck.

\DndQuote%
{}
{''Magic is hard to control at the best of times. Sometimes you wish you could throw fire, and you end up setting yourself on fire. I know from experience. Look at this scar... arcane anomaly.''}
{Sebastian Crowe}
\subsection{Overview}
In this adventure, the characters must escape the time-cursed Black Ivory Inn with their lives and sanity intact. They may be drawn in by the mysterious melody coming from the disturbingly normal looking building, or they might seek out the Black Ivory Inn deliberately to find a missing contact or the smugglers passage rumoured to run underneath the hotel.

\subsubsection{Adventure Hooks}
\subparagraph{A Haunting Melody}
While travelling through Kingside, the player characters may hear piano music in the air, and the sounds of mirthful laughter. Following the music leads them to a pristine building, spared from the ravages of time and the meteor itself. Alternately the characters may hear about the inn and its music from a rumour, and may choose to investigate it on their own accord.

\subparagraph{The Smuggler's Passage}
In its heyday, the Black Ivory Inn was a notable speakeasy known for having an excellent smuggler's passage to move illegal goods into Drakkenheim under the walls. If this passage can be secured it may prove useful for navigating the more dangerous areas of the city.

\subparagraph{Missing Prophet}
A missing priest of the Falling Fire was sent to aid and recruit the rumoured people staying at the Inn, but has not returned. The party is asked to go investigate this disappearance and, if possible, bring the priest back to his people.

\DndQuote%
{}
{''Magic itself is unpredictable, not wielded but conjured with surprise. It all looks as wild as a boar.''}
{Pluto Jackson}
\subsection{Interactions}
The time loop is the result of a curious interaction between the Haze and a peculiarly worded \textit{wish}.

\subsubsection{Time Loop}
A renowned musician performed here named Miss Charlotte. In her youth, Charlotte encountered an extraplanar entity that granted her a single \textit{wish}. Her wish was to be the greatest living pianist of all time. However, Miss Charlotte was already an unparalleled musician. Instead of improving her talents, her wish granted her immortality by binding her life and talents to her piano, not unlike the relationship between a lich and its phylactery. As a result, she would perpetually remain the greatest living piano player \textit{of all time}, so long as her piano remained intact.

Miss Charlotte was impaled by a delerium crystal during the destruction of Drakkenheim. She would have died if not for her wish. Caught between life and death, Miss Charlotte's wish has been twisted by the magic of the Haze, throwing the entire Black Ivory Inn into a temporal pocket that traps anyone who enters. Her body was transformed into a \textbf{protean abomination}, which now lurks below the temporal reflection of the Black Ivory Inn. Her mind projects a simulacrum of her original form that can continue playing her piano for all time.

\subsubsection{Becoming Trapped}
A character becomes trapped in the time loop the moment they wholly step inside the Black Ivory Inn. There is no saving throw.


\begin{itemize}
\item Trapped creatures cannot perceive or interact with creatures and objects outside the Black Ivory Inn. This effect begins 10 (3d6) minutes after the creature becomes trapped. During this period, victims can still perceive and interact with creatures outside Black Ivory Inn, but such creatures and objects appear hazy and faded before disappearing entirely. After this point, when a trapped creature looks outside the Black Ivory Inn, they see a time-lost version of Drakkenheim instead of the present-day ruins.
\item Characters trapped within the time loop are visible to outsiders only while they remain inside the Black Ivory Inn; they can't be perceived otherwise.
\end{itemize}

\subsubsection{Effects of the Time Loop}
The loop itself behaves like a pocket dimension or demiplane. The anomaly constantly plays out the last four hours and thirty-three minutes before the meteor struck in a time loop capturing the events from 3:40 PM to 8:13 PM on September 16th.


\begin{itemize}
\item Characters are instantly killed by the meteor at the end of each loop at 8:13 PM.
\item After the death or incapacitation of the entire party (whether at the end of the loop or by some other means) they awaken again at 3:40 PM.
\item The characters are restored to full hit points and receive the benefits of a long rest. Any bodily damage is restored, no matter how severe. The characters retain whichever items were in their possession the moment they entered the Black Ivory Inn, but lose any items or equipment they picked up since then. As such, items that characters brought with them but used during the previous loop are restored!
\item Characters who consumed any food or drink within the Black Ivory Inn during the previous loop must succeed on a DC 10 Wisdom saving throw. On a failed saving throw, they temporarily lose their awareness of being stuck in a time loop until reminded by someone else what has happened to them.
\item Each time the time loop repeats or the characters are killed, each character must succeed on a DC 5 Intelligence saving throw or gain a random form of indefinite of Drakkenheim Madness.
\item Time continues to pass for those outside the loop. A creature trapped in the loop for over a year of real-world time loses their memories and becomes bound to the loop forever. They die and crumble to dust if the loop ends, and then may only be restored via a \textit{wish} spell.
\end{itemize}

\subsubsection{Escaping the Loop}
A character may escape the loop via \textit{plane shift} or \textit{wish}, but teleportation spells directed outside the demiplane fail. Characters might also be able to convince Miss Charlotte to let them leave.

\subsubsection{Breaking the Loop}
The time loop ends if the \textbf{protean abomination} and the \textbf{piano} are both destroyed. When this happens the Black Ivory Inn and its patrons slowly erode into black ash. All that are left behind are the characters, and the five other visitors (or their corpses, if they did not survive the final loop). Little else remains of the Black Ivory Inn save rubble and ruin - it now appears as it should for a building struck by a meteor. There is a path through the rubble to the Smuggler's Warehouse below. The debris can be cleared away to re-open the underground tunnel.

\subsubsection{Key Events}
\begin{DndReadAloud}{}
"Miss Charlotte is the greatest pianist of all time."

\end{DndReadAloud}


\begin{itemize}
\item The time loop begins at 3:40 PM
\item Miss Charlotte plays all afternoon into the evening.
\item At 7pm she takes a break from playing. She sits at the bar and chats with patrons, enjoying a glass of wine and a small meal for thirty minutes.
\item The meteor strikes at 8:13 PM, shortly after dusk. In this moment, a meteor chunk impales Miss Charlotte, no matter where she is located at that time.
\end{itemize}

\begin{DndReadAloud}{}
There is a wash of dust and ash, and you open your eyes to thunderous applause as Charlotte finishes her first set.

\end{DndReadAloud}

Once inside the Black Ivory Inn, characters see Drakkenheim as it was fifteen years ago, in the late afternoon of September 16, 1111.

\subsubsection{Miss Charlotte}
Miss Charlotte is a willowy woman with slender limbs and delicate fingers. She wears a simple green gown, and keeps her hair tied in a bun behind her head. She gracefully plays the piano in an enraptured trance.


\begin{description}
\item[Ideal.]
To play music is all I've ever wanted, it's all I've ever had, it's everything I'll ever need, and the only thing I have left. Please don't take that from me.
\item[Bond.]
Artists ought to be free to share their gifts with the world. I admire any artist who has the same passion and skill for their craft as I do for mine.
\item[Flaw.]
I cannot bear what has happened, and the truth fills me with rage and anger. Instead, I will surround myself with music and mirth until the end of time.
\end{description}

Miss Charlotte and the piano are immune to all damage as long as the abomination has one hit point or more. In addition, as long as the abomination has at least one hit point, when Miss Charlotte fails a saving throw, she can choose to succeed instead.

When Miss Charlotte plays a chord on the piano (as an action), the abomination rejuvenates with all its hit points.

\begin{DndReadAloud}{}
"I remember you... you caused a bit of trouble last time you were here. Just try to relax and enjoy the music."

\end{DndReadAloud}

\subsubsection{Tavern Patrons}
Most patrons encountered within the Black Ivory Inn are hopelessly lost and cannot comprehend their current predicament. They are dimly aware of what is going on, but have endured so much madness it has crushed their actual will. Tales about meteors and time loops are met with laughter or bewilderment, which quickly turns to outrage and disgust if the player characters pursue it. Even proof such as delerium shards or uncanny foreknowledge fails to rouse them from their stupor.


\begin{itemize}
\item There's roughly 35 (10d6) human commoners or nobles enjoying music, coffee, food, and accommodations at the Black Ivory Inn, with a staff of 7 (2d6) more commoners who work as cooks, porters, and housekeepers alongside the notable figures below.
\item These people are forever lost to the time loop, and die when it ends. They never leave the inn, nor have any recollection or understanding of the time loop.
\item The tavern patrons defend Miss Charlotte if she's attacked. They can't perceive the protean abomination.
\end{itemize}

\subparagraph{"Open" Mike Connolly}
The owner of the Black Ivory Inn is a rotund human \textbf{scoundrel} with blonde slicked back hair, a curly moustache and a set of round framed glasses, dressed in a red velvet coat. He is often found sitting close to the stage admiring the musicians. Though he knows about the illicit activities happening beneath his bar, he doesn't discuss the smuggler's tunnel openly with anyone he hasn't met before, but mentioning his young nephew Mel sparks his attention. Mike has the key to both the cellar doors and the smuggler's tunnel, which he keeps on a key ring on his belt.

\subparagraph{Rogan Blum}
The well-dressed bartender of the Black Ivory Inn is a human \textbf{spy}. Rogan will listen to any tale of woe and regret from a paying customer, and his advice seems staggeringly insightful two or three drinks in.

\subparagraph{Phifer Winters}
A young human \textbf{commoner} serves drinks and coffee with a warm smile and a friendly compliment. Though she's often too busy working to chat, she'll excitedly talk about her aspirations to one day sail the Diamond Sea and travel to distant lands.

\subparagraph{Reginald Grimes}
This human \textbf{commoner} is approaching his seventies, but still cooks in the kitchen of the Black Ivory Inn. He takes great pride in his culinary work and is completely unable to tolerate criticism of his cooking. He deems each meal he makes a work of art.

\subsubsection{Smuggler's Ring}
A group of smugglers.


\begin{itemize}
\item 3 (1d6) bandit captains and 14 (4d6) bandits.
\end{itemize}

Captain Wicket is a \textbf{bandit captain} who has an old pirate hat and blue coat on. He is holding a meeting for his men about the next job they are doing in the basement meeting room in the Black Ivory Inn. He has three trusted commanders and a gang of thieves who have come to meet him here. He has requested that no one disturb their meeting. Detailed below is his crew:


\begin{itemize}
\item Left Hand Lila: a \textbf{bandit captain} with a hook for a right hand
\item Mr. Big: A gnome \textbf{bandit captain} with a devilish smirk and a knack for dramatic plans
\item Wind: an albino tiefling \textbf{bandit captain} with bright red eyes, black hair and huge horns
\item Ted Oats: a budding criminal with short black hair, eager for this job
\item Fred Oats: the twin brother of Ted with long black hair, defensive of his brother
\end{itemize}

\subsubsection{Trapped Visitors}
There are five people who have become trapped here recently whom the characters' actions may jolt back into awareness. Though they have forgotten themselves, they can be roused from their stupor by reminding them of present events.

\subparagraph{Osiris Chronolly}
A long-bearded and purple-robed Amethyst Academy \textbf{mage} who came to the Black Ivory Inn following a long-lost rumour that the legendary musician had gained her talents by making a wish. He, however, begrudgingly refuses any idea that he is trapped in a time loop unless some sort of irrefutable evidence can be placed before him. Once presented with proof, he speculates that breaking the loop will set recent arrivals free, but suspects those who were here when the time loop was created are lost forever.

\subparagraph{Aaron Grint}
This middle-aged \textbf{urban ranger} of the Hooded Lanterns used to come to the Black Ivory Inn during his youth, and was a big fan of Miss Charlotte. He sits in a corner with his hood up smoking a pipe. He blissfully believes he has returned to his lost youth and has forgotten the horrors he witnessed in the Civil War. The return of his memories drives him into a rage.

\subparagraph{Anika Patel}
This human \textbf{knight} of the Silver Order devours meal after meal while enjoying the music. She has forgotten that she came to Drakkenheim to fight with the Silver Order, and now believes her oath is to protect Miss Charlotte. Anika is willing to help, but is bewildered by the idea of a time loop. It takes extreme convincing for her to accept the concept. Regardless, she stands up to quell any violence that breaks out in the tavern.

\subparagraph{Balthazar Adamos}
This human \textbf{priest} of the Falling Fire came to the Black Ivory Inn to investigate the haunting music. He doesn't drink, so he's maintained his memories through several loops already, and his inability to escape is driving him mad. Balthazar has now claimed a room upstairs and locked himself inside. He becomes crazed if someone mentions the time loop to him, but is terrified that ending the loop will mean his death. He came to Drakkenheim to participate in the Sacrament of the Falling Fire, and believes if he dies his soul will forever wander in darkness.

\subparagraph{Fate Hope}
A human \textbf{scoundrel} clad in the colours of the Queen's Men, Fate arrived at the Black Ivory Inn looking for the smugglers tunnel. After talking to a few of the patrons, she sat at the bar to have a drink and lost her memories. She now spends each loop knocking back a few pints at the bar.

\subsection{Area Details}
\includegraphics[width=0.4\textwidth]{images/../5e.tools/img/adventure/DoDk/046-map-5.01-black-ivory-inn.png}
\includegraphics[width=0.4\textwidth]{images/../5e.tools/img/adventure/DoDk/047-map-5.01-black-ivory-inn-player.png}
\subsubsection{Approach}
\begin{DndReadAloud}{}
From the outside, the Black Ivory Inn appears perfectly normal. The sounds of beautiful music can be heard from a block away and the lights are all on. The sounds of revelry and clanging glasses are heard within, drawing those who wander by towards the inn. The inn itself is a small square building two storeys tall with a porch that wraps fully around it. It has a sign hanging above the door that reads "Black Ivory Inn, Come for a drink, Stay for the company."

\end{DndReadAloud}

\subsubsection{Inn}
\subparagraph{The Piano Bar}
Richly appointed with velvet curtains and gleaming varnished tables, each set with a candle in a glass bowl. Two well-dressed wait staff and the well-groomed bartender smile warmly.

\subparagraph{Stage}
Upon a candlelit stage in the far corner of the room stands a jet black piano with glistening white keys. Miss Charlotte is found here playing most of the time.


\begin{itemize}
\item The piano has AC 15 and 30 hit points. It is immune to non-magical damage.
\end{itemize}

\subparagraph{Kitchens and Pantries}
Most of the meals served at the Black Ivory Inn were prepared foods purchased from local bakers, farmers, and butchers.

\subparagraph{Scullery}
Washing rooms for laundry.

\subparagraph{Inn Rooms}
Several comfortable inn rooms each with two beds, a chest of drawers, a small table, and an armchair. Several also include a wardrobe or baths. A common privy on each floor is kept clean and fresh.

\subsubsection{Basement}
\subparagraph{Storage Cellar}
This mundane cellar is stocked with a few kegs of ale, and has a large wall filled with fine wines of various ages and assortments. Most of them are Crownland wines from the north vineyard of Drakkenheim. Many jars and crates of preserved foods and dried fruits and berries can also be found, along with sacks of grains and flour, and a few household tools and appliances. At the back of the cellar is an old oak door with a padlock on it.

\subparagraph{Gambling Den}
This room is clearly appointed for preparing smuggling operations and as a place for secret communication, planning, or illegal gambling to take place. The room is sparse of decorations save for a few round tables and some of the older stools from upstairs. The tables are laid out with decks of cards, bottles of booze, and even a map pinned on the wall that details the underground passages, smugglers' tunnels and sewer systems of Drakkenheim. Currently there is a private meeting happening down here between Captain Wicket and his squad. One of the bandit captains stands at the door to make sure they aren't disturbed as they lay out plans for smuggling a rare painting into Drakkenheim to sell to a fence procuring rare art for a local exhibit. The crew are hostile towards unwanted guests and tell any who come near that they are having a private meeting and they best go back upstairs before they start trouble.

\includegraphics[width=0.4\textwidth]{images/../5e.tools/img/adventure/DoDk/048-05-003.card-box.png}
\subparagraph{Smuggler's Warehouse}
The Smuggler's Warehouse is in disarray. There are smashed crates and a spray of gold coins across the floor. A shelf along the wall has toppled over and splintered along the ground. Along the far wall is an inconspicuous set of thick wooden doors that are locked with a large padlock. In the corner of the room is the \textbf{protean abomination}. It is a massive mound of flesh occupying the upper corner of the room, with sticky fleshy tendrils clinging to the various walls and surfaces around it. It looks like it's spreading outwards towards the door - it has grown and caked over the smuggler's tunnel doors almost completely, and one can make out bits of debris, wood, and gold that have been enveloped by the spreading horror.


\begin{itemize}
\item Stuck in the side of the creature is the remnant of a grand piano, and on the front of the creature one can make out a feminine face, twisted in horror, with glowing octarine eyes and a gaping mouth. The abomination has several other eyes and mouths around it that all moan and cry as one approaches. Its' eyes become fixed on anyone who enters the room.
\end{itemize}

\subparagraph{Underground Tunnel}
The smuggler's tunnel is a long straight tunnel that goes on as far as one can see. It is lined with occasional torches.


\begin{itemize}
\item A small satchel hanging on the wall has a Spell Scroll of \textit{dancing lights} and a Spell Scroll of \textit{knock}.
\item The tunnel connects to the Old Town cistern beneath Slaughterstone Square.
\end{itemize}

\DndQuote%
{}
{''In the deepest parts of the city, the Haze is as thick as pea soup...if only it tasted like it.''}
{Veo Sjena}
\subsubsection{Time-Lost Streets}
A trapped victim exiting the Black Ivory Inn enters into an illusory version of Drakkenheim as it was fifteen years ago. Once victims travel out of line of sight of the Black Ivory Inn, they immediately appear back before the inn again as if the streets themselves looped back around. All the doors and windows of all the buildings seem to open into a crudely rendered room, but when entered, a character appears back in the taproom of the Black Ivory Inn.

The folk walking the streets are merely phantasms with dim awareness. They repeat only a few short phrases such as "Lovely day today, isn't it." "Strange thing in the sky, that." or "I'll be on my way, then..." If attacked or damaged, they die, and the others scream and run.

\begin{DndSidebar}[float=!b]{Developments}
\subsection{Check out anytime you'd like}
With the protean abomination destroyed, the inn returns to the state it was meant to be in. The creature dissolves into nothing but a pile of goo and within it rests a chunk of meteoric iron with 2d4 fragments of delerium embedded in it. The ceiling of the inn has mostly collapsed from the impact and a clear hole in the floor leads directly into the basement where the abomination lay. The inn lies in shambles, another nondescript ruin of Drakkenheim. But, the entrance to the smugglers tunnel is still very much intact, and Open Mike's corpse can be found amongst the bodies of the patrons. The key to the smuggler's tunnel is safely fastened on his belt.



\subsection{A Safe Passage}
The smuggler's tunnel now lays open to the players. The players can choose to use this tunnel that leads them directly to the cistern below Slaughterstone Square without provoking a random encounter. They come out from an inconspicuous grate in the floor of the cistern that the monsters of Drakkenheim seem to have overlooked. Once there, the sewers function as normal.



\subsection{Thankful Pilgrim}
If he survives, the priest is endlessly thankful to be rid of the Black Ivory Inn and to still be alive and offers the party a \textit{bag of holding} or an \textit{alchemy jug}. He asks if they would accompany him to Saint Selina's Monastery. He would most like to introduce them to Lucretia Mathias and speak highly of their courageous and heroic deeds. He latches on to the characters, believing them to be the prophesied heroes that the book of Falling Fire spoke of.



\subsection{Friendly Factions}
The other survivors are all a little dazed from their time spent in the inn, confused and in shock of the grim reality they have just been awoken into. They all thank the characters and agree to speak highly of them to their leaders and commanders. If the characters made enemies of any of the survivors, or if they boldly helped them or rescued them, those stories are brought back to the factions they represented.



\end{DndSidebar}

\section{Buckledown Row}
\DndQuote%
{}
{''Listen, if you wanna make it in this town, you gotta stick with me, right? I know people. I can make connections, introductions! You got some fight in you. You wanna meet the Queen of Thieves? I got your ticket right here!''}
{Blackjack Mel}
This infamous row of dive bars and dingy taverns lies on the southeastern outskirts of Drakkenheim. Decades ago, these were the choice watering holes for lower-class labourers, sweat-soaked craftspeople, and assorted scoundrels who inhabited Drakkenheim. Today, Buckledown Row is a hub for outlaws, outcasts, scavengers, and worse. The taverns are mostly intact, and close enough to the edge of town for a band of rough-and-tumble ruffians to make the trek easy enough, but still far enough into the city to dissuade others from following. Ever since the Queen of Thieves brought the gangs under her influence, the street has been controlled by the Queen's Men: bandits and brigands meet here on neutral ground, join in revelry together, trade with fences and smugglers, and only occasionally slit each other's throats. There's few rules and even less order. Flowing ale and gold tend to calm the most violent disputes, but for anything else the hidden fighting pit is the perfect place to settle feuds, make some bets, and lose a few teeth.

\includegraphics[width=0.4\textwidth]{images/../5e.tools/img/adventure/DoDk/049-05-002.buckledown-row.png}
\subsection{Overview}
On the surface, Buckledown Row might appear to be the stronghold for the Queen's Men in Drakkenheim. It's certainly a gathering point for the gangs and a hub for illicit activity, and the characters just might be able to make contact with the Queen of Thieves here. Alternatively, they might just want to crack some skulls in the fighting ring!

\subsubsection{Adventure Hooks}
\subparagraph{Rules of Buckledown Row}
Characters might hear rumours about Buckledown Row, especially if they frequent the Skull \& Sword Taphouse in Emberwood Village. If they have been working with Blackjack Mel, he'll even suggest the characters test their might in the fighting pits. He offers to sponsor them as a fighting team and act as their manager (for an equal cut of the winnings, of course!). If they're interested, he encourages them to travel there with him. He'll be sure to set up a good matchup! Blackjack Mel explains that anyone who proves themselves in the Buckledown Row fighting pit wins the honor to face the Queen of Thieves' personal champion in her secret arena (located in the Court of Thieves). It's the only surefire way to meet the Queen of Thieves, and she's even proclaimed that she'll grant one personal favour to anyone who can defeat her champion. Of course, no one's actually managed to beat them yet...

\subsection{Interactions}
The Queen of Thieves makes an appearance in Buckledown Row, usually appearing for a few hours on a Friday or Saturday evening to watch the best fights.

The bandits that fill the streets and taverns might not notice strangers, but also are not afraid to stab or steal from anyone who seems like they don't belong.

Asking around the taverns and streets, it's easy to gather that the fighting pits below the taverns are the way people around here get known. Many of the fighters in the pits are renowned legends amongst the gangs. Anyone who can defeat the other fighters gets to take a crack at the Queen's champion in the royal fighting pits. Beat the champion, and you get a private audience with the Queen.

The five taverns in Buckledown Row are all filled with various gangs of the Queen's Men. Around every corner is gambling, cards, dice, and occasionally people brawling openly in the streets. A body being thrown out a window, or the sound of shattering glass or splintering wood are so common that they become almost part of the backdrop of the place. Everyone here is a little rough around the edges and if you stick out too much you'll draw a lot of unwanted attention.

Consult the Outlaws and Scoundrels from the Queen's Men Dossier for some of the symbols and characters to fill out the taverns. The gangs all mingle here but usually each gang has a preferred tavern that their leader can be found in, and a few characteristics to help bring any NPCs to life.

If the party manages to fit in without drawing unwanted attention, the gangs pay them no mind above anyone else. The gangs are gaining so many members that if you are in Buckledown Row, most assume you ought to be there. Discussing other factions, or any plans or adventures that seem suspicious, might arouse the attention of people nearby. Also if the party is acting awkward or nervous, this may draw attention - it may also help their cause to dress in appropriate gang attire or to have symbols of a gang on their gear. Again, the Outlaws and Scoundrels detailed in the Queen's Men Dossier can provide information the party might gather to help them disguise themselves or hatch a plan to fit in. If the party openly approaches a gang leader without invitation, they may gain some attention as well, but most of the leaders are eager to hear out anyone bold enough to approach them and will engage with the party in amused curiosity. Otherwise the gangs accept the characters as more new recruits and try and encourage them with drinking, games, sometimes a fight, or a game of cards. Or maybe trying to push them to join the fighting pits below.

\subsection{Area Details}
\includegraphics[width=0.4\textwidth]{images/../5e.tools/img/adventure/DoDk/050-map-5.02-buckle-down-row.png}
\includegraphics[width=0.4\textwidth]{images/../5e.tools/img/adventure/DoDk/051-map-5.02-buckle-down-row-player.png}
\subsubsection{Buckledown Row}
\begin{DndReadAloud}{}
Roughshod fences and makeshift barricades block off this muddy street of ramshackle taverns. Crudely made bridges stretch between the rooftops, and crooked street lamps glow brightly in the mists. Laughter, cheering, and a general uproar echo down the alleys from a motley assortment of thugs, ruffians, cutthroats, and scoundrels who drink, gamble, and traffic down the street. Keen-eyed sharpshooters look down from the lookout posts, and a pair of muscled bruisers stand watch outside the ramparts.

\end{DndReadAloud}


\begin{itemize}
\item At any time, 21 (6d6) bandits stroll the streets, and 5 (2d4) scoundrels lurk upon the rooftops. A pair of bugbear veterans are posted at each entrance. They'll drive off anyone who looks like a member of the Hooded Lanterns or the Silver Order, but they let anyone else through.
\end{itemize}

If conflict with the factions is escalating, 1d4 Queen's Men strike teams might be posted around Buckledown Row in case of trouble.

\subsubsection{Wishing Well}
\begin{DndReadAloud}{}
An ancient moss-covered fieldstone well rests on a terrace behind one of the taverns. Upon a skeleton nailed to a gnarled tree nearby is a crudely-painted sign that reads "Make a Wish". Peering down the well reveals an impenetrable gloom.

\end{DndReadAloud}


\begin{itemize}
\item Each week, representatives from each of the hundred gangs of the Queen's Men deposit their tribute to the Queen of Thieves here by dropping a marked satchel of gold or gems down the well. Each outfit pays a sum of 100 gold or more, depending on their arrangements with the Queen.
\item The gangs don't know what happens afterwards, but the Queen always knows when someone misses a payment.
\item The bottom of the well is concealed by a permanent \textit{major image} and a silent \textit{alarm} spell, but their magic has been hidden with the \textit{Nystul's Magic Aura} spell. Below is a \textit{portable hole} that collects the tribute.
\item In secret, the \textbf{Queen of Thieves} herself arrives invisibly throughout the week to collect her tribute. She always leaves behind a sack of gold. This "float" is used by the Queen to track down anyone who robs her: since she once possessed the gold, she will \textit{Scrying} upon the coins to locate any thief who rips her off.
\end{itemize}

\subsubsection{Sweaty Bugbear}
\begin{DndReadAloud}{}
This small ramshackle taphouse is little more than a dank hole in the wall. Almost literally. The door is missing and the wall where the entrance was has been blasted away to leave just a gaping hole that people come and go through. No signage seems to exist but in the large window next to the blasted entrance is the cook of the place. There is a smoking grill where a large bugbear named Brian is cooking some charred-looking meat; he is wet with sweat caused by the heat of the grill. Baskerville, the leader of the Howlin' Dogs gang, sits at a table in the corner with cushioned seats, a bottle of mead in hand and talking with some scary-looking thieves. There is a game of knife throwing happening in the back. There also seems to be a gnome in a top hat strapped to the target they are throwing knives at.

\end{DndReadAloud}


\begin{itemize}
\item The knife throwing thugs are willing to let anyone join. If you hit the gnome, you lose. If you hit between his legs it's 20 points. Between an arm and leg is 10, and between an arm and his head is 30. They also explain that the gnome tried to skim 400gp off a payment to Baskerville and so he's on the board tonight.
\item Baskerville is amused by anyone coming to talk to him and listens to whatever they have to say. However if it doesn't involve smashing, killing, or getting rich, he laughs at the intrusion and waves them away. If anyone in the party looks tough, he challenges them to arm wrestle for fun and wants to bet 20 gold on it.
\item If anyone wants to join the Howlin' Dogs, as long as they have renown with the Queen's Men, Baskerville offers any strong and tough looking characters a place if they can beat his Fighting Pit champion, Bloodhound.
\end{itemize}

\includegraphics[width=0.4\textwidth]{images/../5e.tools/img/adventure/DoDk/052-05-004.maps.png}
\subsubsection{Sneaky McGee's}
\begin{DndReadAloud}{}
A large neon-painted sign of a ram skull with swirling designs and patterns inside it is posted on the front of this old building. It's three storeys high and has balconies all along the top two floors, with a large wooden door that has a cobra motif on it. Inside the tavern there are several small private booths with tables and seating for four to six people. At the far end of the tavern is a small stage where some rough looking scoundrels are playing instruments and riling the place up. A long bar along one side has a row of red-velvet-topped stools contrasting against the dark almost black wood of the bar. A dwarf woman behind the bar named Reeda Blackstone is polishing glasses. She is friendly, kind-hearted, and welcoming to any who come through the door. Everyone who orders a drink from her tips her one gold and says "Thank you Reeda." Above the bar is a wall of swords and daggers, all stuck in place in an array of metal and leather. The entire wall is just thousands of blades. Sitting at the end of the bar is Veronica Venom and her gang.

\end{DndReadAloud}


\begin{itemize}
\item People in the tavern will tell you that Reeda may seem nice, but she's a ruthless killer and used to be one of the most notorious assassins in Westemär before she retired. Everyone respects her, and she demands respect.
\item If the characters order anything from Reeda without tipping her a gold piece , or dont say "Thank you Reeda," she grabs them and asks "Are you new in town honey? You seem lost." before motioning for Veronica to come over. This usually grabs everyone in the tavern's attention. If the characters don't smooth-talk their way out, there could be blood.
\item Veronica ignores anyone who approaches her that she didn't invite. If she notices anyone who looks stealthy and like they have a killer instinct, she may ask one of her gang members to go invite that character and that character alone over. She offers the character a spot in her gang if they can defeat her champion, Mamba Jones.
\end{itemize}

\DndQuote%
{}
{''Never turn up your nose to a good meal in Drakkenheim, you never know what your next one will have to be...''}
{Veo Sjena}
\subsubsection{Padlock Tavern}
\begin{DndReadAloud}{}
The door to this tavern is covered top to bottom in keys of various shapes and sizes. Gold, silver, bronze, copper, some rusted, some shining. All manner of keys adorn the door. Upon entering the tavern one can see the walls are covered in padlocks of equally varying shapes and sizes. A bunch of round tables fill the centre of this building. And although two stories tall, there is no upstairs save for a small balcony and a spiral staircase going up to it. A man sits up on the balcony in an animal fur coat with a large purple mohawk. The tavern is filled with leather clad bandits with spike studded armour and chains hanging from all over their gear. There is a lot of hollering and it seems a fight is breaking out every couple minutes somewhere in the place. There is a small corner bar with a shelf of various bottles of alcohol spanning both storeys from floor to ceiling. A long rickety ladder transcends the shelves on a wheeled system so it can slide back and forth across the many types of alcohol. Three goblins work the bar. One on the ladder, the other below catching and tossing bottles, and one making the drinks and sliding them down the bar to whoever catches them. People don't seem to ask questions, the goblins seem to make whatever mix of booze suits them and toss it down the bar to whoever catches it. Occasionally a glass slides right off the bar and smashes on the wall to which everyone in the tavern hollers and cheers.

\end{DndReadAloud}


\begin{itemize}
\item Blackjack Mel can be found near the end of the bar with a few of his trusted thieves. He has a private booth that he uses to discuss business.
\item Leon the Swine doesnt leave his balcony, but if anyone wants to talk to him they are welcome to come do so. He demands that anyone who speaks to him calls him the Swine Prince. He refers to the characters as his "Little Piggies."
\item To join the Pins \& Clubs all you have to do is kill someone. Leon explains that they're the biggest gang in Drakkenheim. Other gangs claim to be smarter, stronger, braver, more cunning. But no one has the numbers of the Pins \& Clubs. So to keep things proper, ya gotta go to one of the other taverns and kill someone from a different gang. Come back with their right hand and Leon will let ya join. He gestures to a box in the corner of the balcony overflowing with right hands.
\end{itemize}

\includegraphics[width=0.4\textwidth]{images/../5e.tools/img/adventure/DoDk/053-05-005.inn-keeper.png}
\subsubsection{Old Imperial Public House}
\begin{DndReadAloud}{}
This long and large building still stands almost entirely intact and seems in pretty good condition. Its blue-painted wooden exterior still holds up to only moderate chipping. The large sign outside reads "The Old Imperial" with a coat of arms on the door. Inside the tavern is fitted with beautiful hand carved tables and chairs, and an ornately carved maple bar stretches down the length of the house with a massive wall of wines and spirits behind it. There is a raised section to the bar up a few small steps where Rose Carver and her gang sit around a large table and drink. The bar is tended by a large dragonborn named Israel. There are several more heavily clad and armed folk by the bar chatting with Israel.

\end{DndReadAloud}


\begin{itemize}
\item Israel takes in new recruits for the fighting pits downstairs. If you want to fight you have to register with him. You need a crew, a name, and a catch phrase. The audience pays more when they have a catchy team to cheer for. So make it good. Once you sign up, if you are new, you "gotta do the rounds." Meaning fight all other champions.
\item Rose Carver is looking for new recruits. She likes people who are strong-willed and cunning in battle. Not brutes or overconfident muscle heads as she calls them, but tough folk looking to make a difference, one way or another.
\item To join you just need to agree to give Rose Carver five percent of any jobs you do under her guidance. If you agree, you get the symbol and join her ranks.
\end{itemize}

\subsubsection{Smee's Palace}
\begin{DndReadAloud}{}
The largest tap house in Buckledown Row, this large building used to be a popular spot for local musical shows and plays. The large stage at the back of the tavern still stands, and often you can find a group of bards here playing songs or occasionally a play being performed. The place is filled with elegant gold-accented wooden-carved features, and there is a large chandelier with what seems like over one hundred candles. Each table is illuminated by red wax candles in golden candelabras. The place feels the fanciest of the local establishments. The server is a tall elven woman with white hair and piercing hazel eyes named Tilda. She speaks elegantly and softly to any who come to her. Even though she is soft spoken, one can hear her clear as day through the noise of the ongoings of the tavern. Sitting at the table closest to the stage is Christian Lament and his gang.

\end{DndReadAloud}


\begin{itemize}
\item There is a staircase behind the bar that leads down to the fighting pits and people are constantly coming and going from below. You can hear the sound of cheering amid the thuds and clangs of combat.
\item Christain Lament hates being interrupted if there is a show on. He is a huge fan of music and theatre, and especially musical theatre. Any interruption causes him to immediately respond "Someone have this one killed please," to which usually no one does anything, but if the characters persist Christan will sigh and turn to them.
\item Christian is uninterested in any proposals from the characters unless they cater to his massive ego and explain the need for someone more cunning, smart, or handsome than any of them. Pleading to Christan's vanity is a sure way to change his demeanor towards you.
\item If one of the characters is charismatic and cunning with their words he may ask them if they have what it takes to be a Wounded Heart. If they want to join his gang they need to defeat Dorian Dare in the fighting pits, then Christain will allow them to join the gang.
\end{itemize}

\subsubsection{Fighting Pits}
Below the taverns of Buckledown Row, the basements and sewers have been restructured into a makeshift fighting arena. They are crowded with people either there to fight, or betting and cheering on the combatants.

\subparagraph{Infirmary}
A human \textbf{priest} by the name of Patty Tomkins is working here with her gnome assistant Marv to heal up some unfortunate fighters who didn't win in the pits. The two of them have some crude medical equipment and a few spells to try and patch up any who are worse for wear.


\begin{itemize}
\item She is willing to cast \textit{cure wounds} on the party if they need healing and asks for tips in exchange. She knows the Queen's private fighting pits are down the sewers some ways but has never been there. She's very tired of how injured people get in the pits and wishes they were just better at fighting.
\end{itemize}

\subparagraph{Meeting Room}
A few small tables and chairs are set up in this otherwise barren stone room. There is a \textbf{bandit captain}, 2 thugs, and 5 bandits scheming up plans in here at any given time.

\subparagraph{Prisoners Stocks}
A few people are locked in some rickety old cells here. They are watched over by a rough looking bugbear \textbf{gladiator} named Ludwig. He says that the prisoners broke the trust of their gangs and so are forced to be the early rounds of the pit's competitions. Basically, they buffer the champions with a few punching bags. Once the punching bag is spent, their gang leader can decide if they are allowed back in, or if they will be killed.

\subparagraph{Fighting Pit}
The cleared dirt-covered floor is caked in blood stains, mud, and other fluids that are best not asked about. A central square is marked off by makeshift wooden barricades, around which about forty bandits gather waiting for the next match. A big human \textbf{veteran} named Bull is refereeing the fights. Next to him, a halfling \textbf{commoner} with a moustache named Rocky is taking bets and wagers on the next fight.

\subparagraph{Fence}
Tig is a young human \textbf{spy} who looks no older than fifteen. He sits on one end of a table in this room and is extremely excited and welcoming to anyone who enters. He exclaims "Welcome to Tig's Shop!" to all who enter. Tig is accompanied by a half-orc \textbf{gladiator} named Tony who stands quietly nearby in case of any problems. Tig sells a few select stolen items (choose three common magic items every time the characters show up as well as 1d4 health potions and 1d4 \textit{aqua delerium}). Tig is also willing to buy any stolen goods from the party if they have any.

\subsubsection{Pit Fighters}
When characters join the fighting pits they may fight individually or as a team. If they fight one on one they only fight the champion of that team, or a tag team with two champions at a time - if they are doing tag teams, a \textbf{thug} joins Mamba Jones. Work your way down the list until all champions are defeated.


\begin{itemize}
\item Glassjaw Jake, a human \textbf{bandit} who looks pretty beat up already
\item Smiley Pete, a dwarf \textbf{thug} with no teeth
\item Dorian Dare, a human \textbf{veteran} with the Wounded Hearts symbol on his shoulder
\item Bloodhound, a Hobgoblin \textbf{gladiator} with the Howlin' Dogs symbol on her shoulder
\item Mamba Jones, a halfling \textbf{assassin} with the Catacomb Cobras symbol on her shoulder
\end{itemize}

\subsubsection{Pit Fight Rules}

\begin{itemize}
\item Fights can either be one-on-one, or a tag team.
\item Tag team matches pit equally-sized teams against each other.
\item In a tag team, fighters "tag in" and "tag out" by visibly tapping each other. Whoever is tagging in must be wholly outside the ring, whoever is tagging out must be wholly within.
\item You must stay in the ring. If you aren't in the ring at the end of your turn, you lose.
\item If you know magic, you can use it. You can use whatever gear you've got.
\item Your team loses if any member outside the ring is caught casting spells or otherwise influencing anyone in the ring.
\item Team members who are tagged out are allowed to cast spells on each other.
\item No poison! We mean it! If you get caught with it, we'll make you drink it all.
\item No potions! If anyone on your team drinks a potion or is caught under the effects of one, your team loses, and we're taking all your potions. The only exception is using a potion to save the life of a knocked out teammate.
\item No invisibility or darkness stuff! It's no fun if the audience can't see you fighting!
\item If you get knocked unconscious, you're out. The pit master will take you off the field and stabilize you so you don't die.
\item If someone gets taken out, you can send in another member to replace them.
\item If it's your first night in the pits, you have to fight.
\item There's no rule against killing your opponent, but generally speaking this is frowned on.
\end{itemize}

\includegraphics[width=0.4\textwidth]{images/../5e.tools/img/adventure/DoDk/054-05-006.fighting-pits.png}
\begin{DndSidebar}[float=!b]{Developments}
If the party is discovered as having ulterior motives or trying to cheat, they get chased out by the gangs and will be much more recognizable if they return. However, if the party was able to join the fighting pits and defeat the champions, they get invited to speak to Lies and Deceit who offer to take them to the Queen's Royal Pits to face the final champion. If they win, they gain an audience with the Queen. They might surmise and uncover that the Queen's chambers are through a secret sewer passage. If they attempt to uncover this passage for themselves, or if Lies and Deceit are escorting them, this brings them to the Court of Thieves.



\DndQuote%
{}
{''A half-baked plan to make your way through the streets of Drakkenheim will lead even the strongest adventurers to bite the big one.''}
{Sebastian Crowe}


\end{DndSidebar}

\section{Chapel of Saint Brenna}
Chapels devoted to the Sacred Flame form the heart of many neighbourhoods in Drakkenheim. These were places of religious worship, weddings, community gatherings, and funerals, as each stands over a small crypt where common folk interred the cremated remains of the dead. However, many are also the resting place for the mummified bodies of clerics and paladins of the Sacred Flame, and some even contain reliquaries for holy artefacts. The Chapel of Saint Brenna holds a particularly valuable relic whose location was unknown until recently: the \textit{Sceptre of Saint Vitruvio}. Now three factions seek this relic: the Followers of the Falling Fire, the Knights of the Silver Order, and the Queen's Men.

\includegraphics[width=0.4\textwidth]{images/../5e.tools/img/adventure/DoDk/055-05-007.chapel.png}
\subsection{Overview}
In this adventure, characters explore the dungeon crypts under the Chapel of Saint Brenna to recover the \textit{Sceptre of Saint Vitruvio}. Complicating matters, at least two other factions have dispatched strike teams of their own to claim the \textit{Sceptre}. If you wish to really challenge the characters, you could add a third! When the characters arrive, one faction team is already searching the upper level of the chapel, and another rival faction will arrive on the scene later to contest the claim. It's up to the characters to decide how to deal with these faction agents: they could fight them, trick them, sneak past them, or negotiate with them in some way.

The puzzles in the crypt invite player characters to use their own creativity to develop out-of-the-box solutions, as there is no "correct" way to overcome them. However, if characters rouse the remains of the interred clerics and paladins in the process, they'll have an interesting problem on their hands as the undead will rejuvenate when destroyed unless laid to rest again properly.

\subsubsection{Adventure Hooks}
\subparagraph{The Sceptre of Saint Vitruvio}
Recently, a group of rival adventurers found letters and missives in another ruin that indicated the \textit{Sceptre of Saint Vitruvio} was buried beneath the Chapel of Saint Brenna. Unfortunately, these adventurers disagreed over how to use this information and went their separate ways. They each brought the discovery to the attention of different factions. Now, the race is on to recover the Sceptre first! While each faction could send their own strike team to the chapel, whichever faction has the closest relationship to your party seeks out the characters' aid instead.

A messenger from one of the factions seeks the characters out in Emberwood Village, urgently requesting them to meet with one of the following faction lieutenants:


\begin{itemize}
\item Flamekeeper Ophelia Reed at the Chapel of Saint Ardenna
\item Nathaniel Flint at the Hendrix Farm
\item Blackjack Mel at the Skull \& Sword Taphouse
\end{itemize}

The lieutenant gets right down to business.

The characters need to act now, as other factions are already sending agents of their own.

\begin{DndReadAloud}{}

\begin{description}
\item[Ophelia Reed.]
"We cannot allow the Sceptre of Saint Vitruvio to fall into the hands of thieves or cultists. Help us retrieve it, and when you meet me back at Camp Dawn, we'll see you are rewarded faithfully for your efforts."
\end{description}

\end{DndReadAloud}

\begin{DndReadAloud}{}

\begin{description}
\item[Nathaniel Flint.]
"We see little value in worldly gold and silver, but if you help us with this task, you would be champions of our faith. You must bring the Sceptre to Lucretia Mathias herself, in the hallowed halls of Saint Selina's Monastery. There she will grant you a great boon."
\end{description}

\end{DndReadAloud}

\begin{DndReadAloud}{}

\begin{description}
\item[Blackjack Mel.]
"The Sceptre is exceptionally valuable to the right buyer... but if you ain't got the right connections you'll never move it on your own. You'll get your cut, but if we get wind of you trying to double deal us by selling it yourself, it's gonna be bad news for you."
\end{description}

\end{DndReadAloud}

\subsection{Interactions}
Whichever factions the characters \textit{aren't} working with send their own strike teams to the Chapel of Saint Brenna. The Silver Order strike team always arrives first, and the Followers of the Falling Fire arrive just as the characters are leaving the site with the sceptre. The Queen's Men can either arrive first (if the characters are working for the Silver Order) or show up if the characters take a short rest (if the characters are working for the Followers of the Falling Fire).

\subparagraph{Followers of the Falling Fire}
Katja Braun, a human \textbf{cult fanatic} leads 8 human cultists. She is a thin middle-aged woman who believes she is guided by righteous purpose.

\textbf{Knights of the Silver Order}. Ser Bryce Landry, a human \textbf{knight} leads 4 human guards. This clumsy and lanky warrior insists on doing things "properly" and refuses to damage the chapel or defile the crypts in any way. Each has a \textbf{warhorse} that they hitch to the trees outside the chapel.

\subparagraph{Queen's Men}
Rorsh, a human \textbf{bandit captain} with horrible burns all over his body leads 8 human bandits. They were promised 1,000 gp for the job, but are always looking out for a better deal.

Consult the \textbf{Roleplaying Traits} in the \textbf{Faction Dossier} for ideas on how to roleplay these faction strike teams, and the \textbf{Strike Teams} section for details on their combat tactics.

\subsection{Area Details}
\includegraphics[width=0.4\textwidth]{images/../5e.tools/img/adventure/DoDk/056-map-5.03-chapel-of-st-brena.png}
\includegraphics[width=0.4\textwidth]{images/../5e.tools/img/adventure/DoDk/057-map-5.03-chapel-of-st-brena-player.png}
\subsubsection{Chapel Grounds}
A deteriorating chapel lies atop a small hillock in the middle of a desolate residential borough. Despite heavy damage, its domed slate roof rises above the surrounding townhouses. Two cylinder-shaped towers flank the arched double-door entrance. One tower has completely collapsed into a mound of stone rubble roughly fifteen feet high, leaving a great breach leading into the chapel at ground level. The other tower remains standing up to the second storey. Piles of debris and stone surround the tower's base. Elegant arched windows filled with stained glass mosaics decorate the chapel, but most are shattered and broken.

\subsubsection{[1] Crematory Garden}
\begin{DndReadAloud}{}
Behind the chapel is a small terrace with stone benches and a sundial, and a short footpath connects the chapel to a walled garden to the east. Two soil plots within served as sites for funeral pyres. A melancholic statue of a plate-armoured woman holds vigil over these crematory mounds, clasping a silver rod in both hands over her chest.

\end{DndReadAloud}


\begin{itemize}
\item A metal hatch behind the statue is locked with a padlock. The key is long gone, but the lock can be picked or broken with a \textbf{DC 20 Strength} or \textbf{Dexterity} check. If opened, it leads to a steep staircase leading down to the Crematory Passage (Area 6).
\item DC 10 Religion. This statue depicts Saint Tarna, the founder of the faith of the Sacred Flame.
\item DC 10 History or Religion. The Faith of the Sacred Flame cremates the dead and keeps their ashes and bones within chapel crypts.
\item The statue holds a silver \textit{immovable rod}. It is one inch in diameter and fifteen inches long. The button on the rod is engaged. If released, the rod slides out of the statue's hands. The rod is perfectly smooth save for a wedge-shaped divot on the bottom edge. The rod functions as a handle for the winches in the undercroft (Area 4).
\item If the Queen's Men strike team arrives first, they are encountered here. They've mistaken the \textit{immovable rod} in the crematory garden for the true \textit{Sceptre of Saint Vitruvio}. They suspect (incorrectly) that some kind of trap will be triggered when the button is pressed. They're trying to figure out how to disarm it, and failing, since there isn't actually any trap here.
\end{itemize}

\subsubsection{[2] Chapel}
\begin{DndReadAloud}{}
The double doors open with a loud creak to a short carpeted hallway decorated with candelabras. The hallway leads directly into the chapel sanctuary. Puddles of brackish rainwater pool in the basin created by the collapsed tower roof.

At the heart of the chapel is a three-tiered stone altar set with a large bronze brazier. Petrified wood and ashes rest within. A few prayer benches and a small altar box lay nearby. Worshippers of the Sacred Flame once gathered here, encircling the brazier hand-in-hand to worship.

Dirt, dust, and debris cover most of the floor, here and there bits of cracked paint from the floor mosaic are visible beneath. Painted on the walls between the windows are fading frescoes of the saints and martyrs of the Sacred Flame. The three alcoves each hold a large statue marred by time and weather. The northern alcove depicts the legendary Saint Tarna holding steadfast against a massive two-headed demon with long tentacles for arms. The eastern alcove features a stoic statue of Saint Vitruvio, the Patron Saint of Drakkenheim. The western statue depicts Saint Brenna, for whom the chapel is named. She is shown performing the miracle of resurrection.

Scattered about the room are ruined texts, iron braziers, and bronze lamps, some of which still hang from the timbers on the roof. The room reeks of bitumen.

\end{DndReadAloud}


\begin{itemize}
\item Two Alchemist's Fire (flask) can be scavenged from the brazier basin.
\item The books and texts here are common copies of the holy books and hymns of the Sacred Flame, and contain no rare knowledge or information.
\item The altar box contains a silver ewer worth 25 gp, Spell Scroll of \textit{bless} and \textit{faerie fire}, two \textit{spell scrolls} of \textit{continual flame}, a silver holy symbol of the Sacred Flame, three vials of anointed oils (counts as Holy Water (flask)), and a block of incense.
\item The back door has rotted off its hinges, and collapses with a loud clatter unless carefully opened.
\item The puddles are slowly draining through a crack in the floor.
\item The statue of Saint Tarna bears the inscription \textit{"Saint Tarna holds silver truth that opens the way to the healing light of the Sacred Flame."} This refers to the \textit{immovable rod} held by the statue outside in the Crematory Garden.
\end{itemize}

\begin{DndReadAloud}{}
A stone spiral staircase leads down to the undercroft. Halfway down the stairs is an iron portcullis gate that bars the way. A lock is set into the gate.

\end{DndReadAloud}


\begin{itemize}
\item The gate key is missing. Instead, the gate can be broken down with a \textbf{DC 20 Strength} check or the lock picked with one minute of work and a \textbf{DC 15 Dexterity (thieves tools)} check.
\item A faint sound of crackling flame and swirling dust can be heard if one quietly steps down the stairs, with the occasional yelp, cackle, or grumbling in Primordial. These noises are caused by the mephits dwelling in the undercroft.
\item If the Silver Order strike team arrives first, they're encountered here. They hitch their horses to a tree outside, watched over by a single \textbf{guard}. Ser Landry and the others explore the site. They've set down their weapons by the door, and are looking for the key to open the gate to the undercroft. Unfortunately, they don't realize it's totally gone, and they'll waste hours searching because they refuse to "defile" the chapel by breaking down the doors.
\end{itemize}

\subsubsection{[3] Undercroft}
\begin{DndReadAloud}{}
Beneath the chapel is a crypt. Urns, clay pots, and painted vases are kept in hundreds of alcoves cut into the walls. Smaller niches in the three central pillars hold human skulls with half-melted candles resting atop them like crowns. The vaulted ceiling is ten feet above, with several pillars supporting the chapel foundations. Iron chandeliers and bronze candelabras are set throughout, many candles still burning with flickering flames in purple, green, and octarine. A peculiar breeze sweeps through this chamber in strange directions, stirring up small dust devils and causing the candlelight to shudder erratically. Two alcoves contain sturdy oak doors bound with iron lattices, both leading east.

Dripping water leaks down from a crack in the ceiling in the south end of the chamber.

\end{DndReadAloud}


\begin{itemize}
\item Ashes, bone fragments, and sometimes skulls are placed in urns, wrapped in linen, or contained in small chests, then interred in the niches in this undercroft. They bear inscriptions marking the names of the deceased.
\item The water is leaking into floor cracks in the earth.
\end{itemize}

\subsubsection{Ashley, Dusty, Smokey, and Cindy}
\begin{DndReadAloud}{}
"Stop there, meatbag! You're trespassing! You better have brought us something nice, fleshie!"

\end{DndReadAloud}

This chamber is now home to several mephits; two steam mephits, a \textbf{dust mephit}, and a \textbf{magma mephit}. These mephits were originally dispatched by the Amethyst Academy to scout the ruins fifteen years ago, but gave up on their mission and made the undercroft their new home. They quite enjoy their surroundings that are filled with materials and elements similar to their own compositions. They will suffer no trespassers in their "home" unless they are given a suitable gift: a gemstone of some kind. They have little understanding of the wider world, but enjoy inflicting capricious cruelty upon human "fleshies". However, they are terrified of anyone who wears the garb of the Amethyst Academy, and begin crying and ask for forgiveness from the character.

\subsubsection{[4A and 4B] Winch Chambers}
\begin{DndReadAloud}{}
At the centre of this small round chamber is a stout metal post three feet high and about one foot thick. There is a one-inch diameter cylindrical notch about four inches deep cut just below the topmost edge.

\end{DndReadAloud}


\begin{itemize}
\item The posts function as winches connected to the Tomb Doors in the levels below. Winch A opens the Tomb Doors marked "A". Turning Winch "B" opens the Tomb Doors marked "B".
\item Unfortunately, there are no handles attached to the winches, which makes turning them to open the doors exceptionally difficult.
\item The \textit{immovable rod} borne by the statue in the crematory garden outside functions as a handle.
\item When a lever is set into the post, a character can use an action to turn the winch and open the doors. A counterweight system unwinds the winches and closes the doors after one round, unless the winch is held in place (requiring an action) or the \textit{immovable rod} is activated.
\item A substitute could be fashioned if the original can't be located. A wooden lever has a fifty percent chance of breaking each round it is in use, but a metal lever works perfectly (but needs to be held in place)
\item A character working with the right tools could jury rig a stopper or mechanism to hold the winches open with one hour of work and a \textbf{DC 15 Intelligence} check.
\item The winch can still be raised without the lever. A character can use an action and make a successful \textbf{DC 20 Strength} check, which raises it enough for a person to crawl underneath. This check must be repeated each round thereafter to hold the doors open.
\item If the winch is released, the doors close with a crash, dealing 3d6 bludgeoning damage to anyone underneath them. It takes an action to close the doors safely using the winch.
\item The doors could be held open onced raised. Doing so requires an action and a \textbf{DC 23 Strength} check each round.
\item Characters could move a suitably durable object to brace the doors open, but the doors are heavy and could break the brace. Roll a die each round, choosing a larger die for a more durable object (1d4 for a weak one, 1d12 for a strong one, ect). On a one, the bracing object breaks or snaps and the doors slam shut.
\end{itemize}

As there is only one handle between the two winches, access to the Reliquary Chamber (Area 10) normally requires opening one set of doors with one winch, closing it, then moving the handle to the other winch to open the other set of doors, then repeating the process to allow an escape. This means either splitting the party or creating an improvised solution to open all doors simultaneously.

Your players might expect there to be a second immovable rod missing somewhere, and will likely spend some time searching for it. This is intentional, but if you are concerned your players might become unreasonably fixated on finding a second rod, you could place one in the clutches of the \textbf{mummy} in area 9.

\subsubsection{[5] Purification Fountain}
\begin{DndReadAloud}{}
This round chamber holds a circular stone font of fresh holy water. An inscription on the edge of the basin reads "May the Sacred Flame keep vigil over the souls of our revered dead."

The long curved wall of this chamber is separated into five large panels, each painted with murals. Each section shows three Flamekeepers performing the funeral rites for a slain paladin.

\end{DndReadAloud}


\begin{itemize}
\item The first shows the body being prepared: one priest washes the body, another wraps it in linens, the last collects the paladin's weapon, shield, and helm.
\item The second mural shows two priests placing the body on a tomb slab, while the third places a holy symbol in its hands and a gemstone upon each eyelid.
\item In the third mural, the priests anoint the body with holy water or sacred oils.
\item In the fourth mural, the priests place and light the candles around the body.
\item The final mural shows one priest reading from a scroll while the others join hands and recite a hymn while standing beside the body.
\end{itemize}

\begin{DndReadAloud}{}
The opposite short curved wall contains several worn shelves and alcoves filled with vials, tools, dusty scrolls, and thin linens.

\end{DndReadAloud}


\begin{itemize}
\item The shelves contain equipment and tools used to prepare the body, six vials of sacred oils, and Spell Scroll of \textit{lesser restoration}, \textit{dispel magic}, and \textit{remove curse}.
\item Amongst the vials here are two Potion of Healing.
\item 2d6 vials of Holy Water (flask) can be filled from the font.
\item DC 12 History or Religion. The Faith of the Sacred Flame treats the remains of clerics and paladins as holy relics. They believe these faithful souls might one day need their earthly bodies again, so they are preserved or mummified rather than cremated. These murals depict this process.
\end{itemize}

\subsubsection{[6] Crematory Entrance}
\begin{DndReadAloud}{}
Leaves, grass, dirt, and debris fill the floor in this stone chamber.

\end{DndReadAloud}


\begin{itemize}
\item The stairs lead to the crematory above. No one has used this passage in ages.
\item The oaken ironbound door is stuck, but can be pushed open with an action.
\end{itemize}

\includegraphics[width=0.4\textwidth]{images/../5e.tools/img/adventure/DoDk/058-05-008.sacred-aberation.png}
\subsubsection{[7] North Crypts}
\begin{DndReadAloud}{}
The stone door to this chamber has no apparent handles or hinges. It is engraved with tracery and images of angels and saints laid to rest. An inscription reads: "Compassion and temperance sustain the light; greed and hatred bring darkness."

\end{DndReadAloud}


\begin{itemize}
\item The stone door to this chamber is raised using the winch in Area 4b.
\end{itemize}

\begin{DndReadAloud}{}
Beyond the door is a hoop-shaped chamber with vaulted ceilings. The chamber is illuminated by dozens of candles placed about six stone sarcophagus slabs set in semicircular alcoves along the walls. Atop each lays the skeletal remains of clerics and paladins of the Sacred Flame. Tattered cloth covers their bones, a silk sash lain over their eyes. Most rest with arms crossed over their bodies holding a torch in one hand and candle snuffer in the other, but one has a shield, another a helm, and a greatsword etched with celestial runes and glowing with a soft blue light is laid atop another.

\end{DndReadAloud}


\begin{itemize}
\item Each corpse is a \textbf{skeleton} that lies in state. Last rites performed over them decades ago hold back the reanimating magic of the Haze. They do not rise unless their remains or possessions are disturbed, in which case the nearby candles snuff out, and they attack the grave robber.
\item The skeletons have the following additional trait:
\begin{DndReadAloud}{}
\section{Rejuvenation}
If the skeleton is destroyed, it regains all its hit points in ten minutes, unless its remains are placed upon the stone slab, surrounded with lit candles, sprinkled with holy water, and a short prayer uttered over them, or a \textit{gentle repose} or \textit{remove curse} is cast upon its remains.

\end{DndReadAloud}

\item Treasure. Each skeleton has gemstones placed in their eye sockets: two citrines worth 50 gp each. They also clutch a gold or silver holy symbol of the Sacred Flame worth 25 gp. Finally, the magical greatsword here is called \textit{Twilight}. It is a \textit{+1 greatsword} that sheds dim illumination to fifteen feet at all times. A creature wielding it can speak and understand Celestial. The other weapons or armour are wasting away and worthless.
\item DC 15 Arcana or Religion. The Haze often causes dead bodies to reanimate spontaneously. Divine rites performed over these remains may prevent this from occurring.
\item If the door was opened recently and then closed, the characters can detect signs of newly disturbed dust and grinding stone with a DC 15 Investigation check. The door must have been opened and closed again quite recently, as the dust is still settling.
\end{itemize}

\begin{DndReadAloud}{}
To the south lies a stone door engraved with an image of a majestic dragon that reads:

"No darkness is eternal while our faith fuels the Flame."

\end{DndReadAloud}

This door opens to the Hidden Reliquary (Area 10). It's raised using the winch in Area 4a.

\subsubsection{[8] South Crypts}
\begin{DndReadAloud}{}
The stone steps lead down to an ancient cobweb-filled crypt that is little more than a roughly hewn stone passage. The irregular ceiling is ten feet high.

\end{DndReadAloud}


\begin{itemize}
\item Just as in the North Crypts (Area 7), there are five stone slabs in this chamber bearing the remains of clerics and paladins of the Sacred Flame. See above for the description.
\item Treasure. Each skeleton has gemstones placed in their eye sockets: two citrines worth 50 gp each. They also clutch a gold or silver holy symbol of the Sacred Flame worth 25 gp. Unlike the north crypts there are no magical weapons left behind here.
\end{itemize}

At the end of the passage is the stone door to the Tomb of Saint Brenna (Area 9).

\subsubsection{[9] Tomb of Saint Brenna}
\begin{DndReadAloud}{}
This heavy stone door is painted with the countenance of a serene Flamekeeper in flowing robes. Beneath lies the inscription "Saint Brenna, be you surrounded in light eternal. Rest as one with the Sacred Flame."

\end{DndReadAloud}


\begin{itemize}
\item The stone door to this chamber is raised using the winch in Area 4b.
\end{itemize}

\begin{DndReadAloud}{}
Within this cavernous chamber lies a single stone slab surrounded by candles. Upon it rests the mummified body of Saint Brenna. However, leaking water from above has doused the light. Nevertheless, the chamber is illuminated by several standing candelabras surrounded by clay pots and canopic jars.

\end{DndReadAloud}


\begin{itemize}
\item The preserved corpse of Saint Brenna lies atop a stone sarcophagus at the centre of the room. Leaking water from the damage above has put out the candles surrounding Saint Brenna, awakening her \textbf{mummy}. Since her body is waterlogged, she isn't vulnerable to fire damage like other mummies. She has the following additional trait:
\begin{DndReadAloud}{}
\section{Rejuvenation}
If the mummy is destroyed, it regains all its hit points in ten minutes, unless its remains are placed upon the stone slab, surrounded with lit candles, sprinkled with holy water, and a short prayer uttered over them, or a \textit{remove curse} is cast upon its remains.

\end{DndReadAloud}

\item Saint Brenna rises if the door to her chamber is opened, and begins wandering the crypts, droning softly while misquoting scripture. She is gripped by madness and despair and attacks any living creatures. If she passes near the other skeletons interred within, she spends an action to rouse them to follow her. She doesn't leave the undercroft.
\end{itemize}

\begin{DndReadAloud}{}
To the north lies a stone door that bears the engraved countenance of Saint Vitruvio, and is inscribed with the words: "None are forever lost to shadow, candlelight glows brightly in the deepest night."

\end{DndReadAloud}


\begin{itemize}
\item The door leads to the Hidden Reliquary (Area 10). It's opened using the winch in Area 4a.
\end{itemize}

\includegraphics[width=0.4\textwidth]{images/../5e.tools/img/adventure/DoDk/059-05-009.sceptre-of-saint-vitruvio.png}
\subsubsection{[10] Hidden Reliquary}
Two doors lead into the reliquary: one from the Tomb of Saint Brenna (Area 9 and opened from Area 4a) and one from the North Crypts (Area 6 and opened from Area 4b).

\begin{DndReadAloud}{}
In the middle of this circular chamber is a large unlit pyre and several piles of ash and charred wood. A statue of Saint Vitruvio is here with two unlit candles set before him. He is an armoured warrior clad in plate wearing a decorated helm. At his right is sheathed a great blade, and a large shield rests at his left side. He holds his hands forwards and upwards towards the heavens, bearing the Sceptre of Saint Vitruvio upon them.

\end{DndReadAloud}

A character who touches the \textit{Sceptre of Saint Vitruvio} hears an ethereal voice in their mind that speaks the following:

\begin{DndReadAloud}{}
Darkness leaves a heavy burden on the heart. The light brings freedom from the weight of sin. If you would take up my scepter, it is not enough to merely bear the light. You must illuminate yourself.

\end{DndReadAloud}

While within the crypt, the \textit{Sceptre of Saint Vitruvio} effectively weighs seven hundred fifty pounds. However, the Sceptre weighs only two pounds when held or carried by a character who is shedding dim or bright light (either magical or nonmagical), or a character who is attuned to it.

\includegraphics[width=0.4\textwidth]{images/../5e.tools/img/adventure/DoDk/060-05-011.nathaniel-portrait.png}

\begin{itemize}
\item A character could use a class feature or spell to do so, such as \textit{light}, \textit{faerie fire}, \textit{continual flame}, or \textit{daylight}. Spell scrolls of each can be found elsewhere in the chapel. If such spells can normally only target an object, an object worn by a character is suitable.
\item A character could also set themselves on fire. If they do, they take 1d6 fire damage at the start of each of their turns.
\item A held or carried light source, such as a torch, candle, or lantern doesn't fulfill the requirements.
\item A character can attempt to lift or drag the sceptre, or move it using suitable magic (as long as the magic can lift objects weighing over seven hundred fifty pounds)
\item A character who lights the candles and the pyre in the room, then speaks a prayer of the Sacred Flame before the statue becomes illuminated, shedding bright light in a ten foot radius for one hour.
\item A bard, cleric, druid, or paladin can spend a short rest to attune to the sceptre. If the characters do so, this is the perfect moment for a faction strike team to arrive!
\end{itemize}

\begin{DndSidebar}[float=!b]{Developments}
\textbf{When the characters} return with the \textit{Sceptre of Saint Vitruvio}, they are rewarded with a faction boon. An uncommon magic item from the respective faction is appropriate, plus either several consumable items or a payment of 250 gold per player character.




\begin{itemize}
\item Characters working for the Silver Order are invited to Camp Dawn to meet Theodore Marshal, and characters working for the Falling Fire are invited to meet \textbf{Lucretia Mathias} at Champion's Gate. The faction leaders thank the characters for retrieving a valuable relic, and implore them to pledge their faith and service to their factions before turning the characters back over to a faction lieutenant for further instructions.
\item If Theodore Marshal takes a liking to the characters, he might arrange to set them up with mounts from the stables.
\item Characters who impress \textbf{Lucretia Mathias} and show genuine interest in joining her faith might be invited to keep the \textit{Sceptre of Saint Vitruvio}. She won't yet invite them to take the Sacrament of the Falling Fire, but foresees a great destiny for them.
\item Characters working with the Queen's Men are invited to Buckledown Row. Blackjack Mel would like to show them the fighting pits, and he might even be able to make arrangements to meet the \textbf{Queen of Thieves}.
\end{itemize}



How the characters dealt with the other faction agents encountered during this adventure could potentially influence future interactions with those factions. It's possible the characters use guile, negotiation, and stealth to navigate this scenario without making enemies. If the characters slay all the members of a strike team, there might not be anyone who can actually report back to identify that the characters were the ones responsible. Nevertheless, if faction members survive and escape, or evidence emerges implicating the characters, it could potentially sour future interactions. However, the factions still regard the characters as neutral mercenaries at this stage of the campaign. Faction agents clash with adventurers all the time in the ruins, so becoming outright enemies of either faction should only occur if the characters incite serious altercations and violence by keeping the sceptre for themselves. If they simply retrieve and return the sceptre to whichever faction hired them, it arises ire from the other factions, but not retribution.



Characters who decide to keep the sceptre become a target for both the Followers of the Falling Fire and the Knights of the Silver Order. Both factions entreat the characters to turn the sceptre over to them, offering rewards in exchange. However, if the characters don't accept these offers, the factions dispatch strike teams to attack them. The first strike team sent won't be ordered to slay the characters, but they will rough them up badly and take the scepter. If this strike team is defeated or killed, the characters might risk further escalation with one of the factions. Characters who incur the ire of the Silver Order or the Falling Fire then turn the sceptre over to one of the other factions, have made an enemy of that faction.



\end{DndSidebar}

\section{Eckerman Mill}
\begin{DndReadAloud}{}
Eckerman Mill is a dilapidated windmill that rests atop a lonely hill northwest of Drakkenheim. Its stone base is covered in moss and ivy, and the rotting wooden walls are riddled with holes. The entire structure creaks in the wind. The windmill blades still spin softly in the breeze, but the machinery within is damaged beyond repair and the millstone no longer turns. Scattered about the mill are the leftovers from countless expeditions. A well-used bonfire pit surrounded by a few well-worn logs for seating and several tattered tents are set up beside the mill. Abandoned here are a few bedrolls, battered cookware, a rusty chopping axe, and a waterlogged stack of wood.

\end{DndReadAloud}


\begin{itemize}
\item The Mill is often used as a meeting place for adventurers venturing into Drakkenheim from the northwestern side of the Drann River. Since Emberwood Village lies several miles to the south on the opposite side of the river, many use Eckerman Mill as a makeshift campsite as it is much closer to the city proper.
\item Occasionally, faction agents who wish to avoid the eyes and ears of Emberwood Village arrange clandestine meetings here. In particular, the location is favoured by field surveyors dispatched by the Amethyst Academy.
\end{itemize}

\includegraphics[width=0.4\textwidth]{images/../5e.tools/img/adventure/DoDk/061-05-010.eckerman-mill.png}
\subsection{Mill Graffiti}
\begin{DndReadAloud}{}
The mill is scrawled over with chalk graffiti and carvings, transforming the walls into an impromptu message board. There are several roughly engraved messages along the walls. Most list the names of dozens of would-be heroes who have ventured into the ruins, though many names are crossed out to mark those who did not return. Curiously, a few groups have not marked off a single casualty... perhaps they all made it out alive, or more likely, none survived. In addition to these memorials, along the wall one can find many dire warnings, cryptic riddles, poignant questions, crude maps to alleged treasure, and amateur drawings of horrible monsters encountered in the ruins.

\end{DndReadAloud}


\begin{itemize}
\item Invite your player characters to write their own messages or names on the walls of Eckerman Mill, to follow the tradition of those who came before.
\item You can distribute one rumour to characters who examine the writings of Eckerman Mill to represent what they learn from reading the scrawled messages and riddles on the walls.
\end{itemize}

\DndQuote%
{}
{''This mill used to be my favourite stomping ground, now it's full of names of those who died. I gotta say, the novelty is wearing off quick.''}
{Sebastian Crowe}
\section{Rat's Nest Tavern}
\DndQuote%
{}
{''My colleagues have an experiment in mind that requires a large delerium crystal. Smaller shards won't suffice. Unfortunately, such specimens are quite rare in the Outer City ruins. However, we've got a lead...''}
{}
A small chunk of the meteor broke off and crashed into the Rat's Nest Tavern in the Sprawl, utterly destroying most of the building. However, an intact delerium crystal was left behind. Unfortunately, a \textbf{ratling} colony has created a burrow under the taphouse basement and moved the meteorite into their warrens. The ratlings have grown bold enough to launch raids throughout the ruins and take several captives in the process.

\includegraphics[width=0.4\textwidth]{images/../5e.tools/img/adventure/DoDk/062-05-012.rat-nest-tavern.png}
\subsection{Overview}
In this adventure, the characters travel to the Rat's Nest Tavern to obtain a delerium crystal or rescue a captive prisoner. How the characters accomplish either is up to them: they can fight, sneak past, or trick the ratlings, but the creatures will viciously defend their lair by responding with ambushes and traps of their own.

\subsubsection{Adventure Hooks}
\subparagraph{Good Prospects}
(Amethyst Academy) River contacts the characters for a clandestine meeting. She wants them to retrieve a delerium crystal from the Rat's Nest. A crystal of this size is normally worth about 1,000 gp, but she offers to pay 1,250 gp for it.


\begin{itemize}
\item River's intel is based on rumours the tavern was destroyed by a breakaway piece of the original meteor. Since her intel is based on speculation, she offers the characters a \textit{potion of healing} each as an advance in the event that the characters turn up nothing.
\item River cannot disclose the nature of the experiment:
\end{itemize}

\begin{DndReadAloud}{}
"Your continued discretion ensures continued employment."

\end{DndReadAloud}

\subparagraph{Missing in Action}
(Hooded Lanterns) Ansom Lang approaches the characters in Emberwood Village. His sister, Petra Lang, was captured when ratlings ambushed them during a supply run last night. He needs the characters to rescue Petra before it's too late. He offers 250 gold pieces as a reward for her safe return.


\begin{itemize}
\item Against orders to retreat, Ansom tracked the ratlings to the Rat's Nest. He was driven off when the ratlings wounded him with a contaminated sling bullet. Ansom explains that ratlings often keep prisoners to have a "fresh meal" later. He's positive Petra is still alive.
\item Ansom warns the characters about the dangers of contamination. He can't help the characters personally as he's recovering from contamination himself. He advises them that if they get contaminated in the ruins, Flamekeeper Hanna in Emberwood Village can help them with her magic.
\item Finally, he promises to put in a good word with the Lord Commander for the characters' help. However, he emphasises this is a personal job - he's already risking insubordination:
\end{itemize}

\begin{DndReadAloud}{}
"Amongst the Hooded Lanterns, our survival doctrine says we don't risk lives on 'futile' rescue attempts. We all knew the dangers when we signed up, but she's my sister."

\end{DndReadAloud}

\subsection{Interactions}
For the most part, the ratlings in the burrow are hostile to trespassers of any sort. They know who lives in the burrow and who does not, and can recognize each other by scent. They lure intruders into traps or ambushes where the ratlings can use their numbers to their advantage. The ratlings are most active at night, when two-thirds leave the burrow on foraging raids. During the daylight hours, most ratlings are sleeping in their dens except for a few sentries.

\subsubsection{The Rat Prince}
The ratlings in the burrow are led by a large and ferocious ratling who self-styles himself as the \textbf{Rat Prince}. He rules because he's the biggest, strongest, and killed the last ratling who called themselves the Rat Prince. He believes once his colony has enough food, they will be able to grow large and multiply. Then all the other ratlings will bow down to him for sure!


\begin{description}
\item[Personality Trait.]
The Rat Prince bullies the others like a despot. When speaking, the Rat Prince repeats words two or three times for emphasis. "Yes yes yes..."
\item[Ideal.]
The Rat Prince loves killing, eating, breeding, and sleeping. Pretty much in that order.
\item[Bond.]
The Rat Prince wears a makeshift crown and claims he's the ruler of all the ratlings in Drakkenheim. This is far from true, but good luck convincing the Rat Prince otherwise.
\item[Flaw.]
Logic makes the Rat Prince confused and angry, unless that logic involves food.
\end{description}

Caught alone, the Rat Prince flees or begs for his life, offering anything he can think of to convince characters to let him go free.

\subsubsection{Squeaks the Seer}
This albino ratling \textbf{warlock of the rat god} is the spiritual leader of the colony.


\begin{description}
\item[Personality Trait.]
She makes up fantastic stories about the Rat God's wrath and cunning.
\item[Ideal.]
The Rat God wants ratlings to eat everything they can. It gave her magic powers to help the ratlings get food and get the "shiny" (delerium). One day it will eat the world. The other ratlings all think this sounds fantastic.
\item[Bond.]
She thinks collecting the shiny (delerium) pleases the Rat God, and leads cannibalistic rituals in its name. Ratlings love eating, so they love these rituals.
\item[Flaw.]
She thinks she can talk to the Rat God through the delerium crystal. She keeps the stone in her lair (Area 8) and becomes enraged if other ratlings touch it.
\end{description}

\subsubsection{Befriending the Ratlings}
In the original \textit{Dungeons of Drakkenheim} campaign, the Rat Prince surrendered to my player characters. Using deception to convince the Rat Prince they were ghosts dispatched by the Rat God, the Rat Prince befriended the characters and became a recurring NPC. The player characters continued to protect and hide the ratlings, but once the gnolls were driven out of the city, nothing was holding back a growing ratling population. Left completely unchecked, their numbers swelled to gargantuan proportions that threatened to overtake the entire city!

\subsection{Area Details}
\includegraphics[width=0.4\textwidth]{images/../5e.tools/img/adventure/DoDk/063-map-5.04-ratling-burrow.png}
\includegraphics[width=0.4\textwidth]{images/../5e.tools/img/adventure/DoDk/064-map-5.04-ratling-burrow-player.png}
\subsubsection{Scrag Lane}
\begin{DndReadAloud}{}
The Rat's Nest Tavern lies on the mud-slick and rubble-strewn Scrag Lane south of Shepherd's Way, nestled between ramshackle tenements and crumbling flats. The original tavern has been almost entirely obliterated, and the neighbouring buildings are heavily damaged. What remains are stone foundations piled with debris. The first-floor walls stand about chest high, a few supporting timbers splayed outward like a burst ribcage. The rooftop and upper levels have entirely collapsed into a mist-covered crater filled with rubble. The original tavern sign juts out obliquely from a building on the opposite side of the street. It is engraved with a smiling, obese rat laying on a plate of food.

\end{DndReadAloud}


\begin{itemize}
\item Characters who search the streets for nearby sewer entrances may discover the hidden entrance to the ratling burrow through the sewers.
\end{itemize}

\subsubsection{Taproom Ruins}
\begin{DndReadAloud}{}
Inside what remains of the tavern walls, the wooden floorboards of the ground level have almost entirely collapsed into a large gaping hole. Tumbled brick piles, ashen cinders, burned shingles, splintered wooden planks, support beams, fractured furniture, cracked mugs, broken glass, random bits of bone and cutlery are scattered everywhere.

\end{DndReadAloud}


\begin{itemize}
\item The hole opens to the basement ten feet below.
\item Investigation or Survival (DC 10). Animal-like scratchings, gnaw marks, patches of coarse fur, muddy tracks, and rat droppings may be found by examining the area.
\end{itemize}

\subsubsection{[1] Cratered Basement}
\begin{DndReadAloud}{}
The stone walls of the basement are scorched and blackened. Rubble and debris surround a muddy crater fifteen feet wide and five feet deep. A faint glimmer of octarine light twinkles within.

\end{DndReadAloud}


\begin{itemize}
\item Examining the crater for the source of the glow turns up 1d6 delerium chips (worth 10 gold each). There's a distinct rocky divot about a foot wide just off the centre of the crater, as if a large hunk of rock was dug out and excavated. If there was a meteorite here, it's gone now.
\item The ratlings removed a chunk of rock fused with a delerium crystal from the crater and dragged it into the Seer's Den (Area 8). A Survival (DC 10) check reveals drag marks, muddy tracks, and clawed footprints of several beastly creatures indicating something heavy was removed from the crater recently.
\item Investigation (DC 10), or proficiency with mason's tools/stone working. The breached wall was not caused by the meteor impact. The passage was dug relatively recently.
\end{itemize}

\subsubsection{Ratling Burrow}
These details are common to most areas within the ratling burrows:

\subparagraph{Passages}
The tunnelsare dug out from packed earth, clay, loose dirt, and bits of stone. The cramped passages are five to ten feet wide, and rarely more than six feet high. Medium-sized creatures must stoop throughout. The tunnels weave and bend left and right, but also up and down.

\subparagraph{Chambers}
The larger rooms are eight to ten feet in height. They are mostly dry past the entrance, but unrecognizable filth, bone fragments, scraps of fur and cloth, and bits of leaves and grasses lay underfoot.

\subparagraph{Light}
The ratlings rely on their darkvision for sight, so there are no light sources within.

\subparagraph{Sounds and Smells}
The air underground is humid. It reeks of rat urine and feces. The sounds of squeaking and chattering ratlings as they scrape and scurry about can be heard down the tunnels.

\subsubsection{[2] Trapped Entrance Passage}
\begin{DndReadAloud}{}
This muddy passage is roughly eight feet wide. It tunnels steeply down ten feet before veering sharply left.

\end{DndReadAloud}


\begin{description}
\item[Perception (DC 10).]
The wet mud is slippery, and could lead to a nasty fall.
\end{description}

\subparagraph{Mudslide Trap}
Characters without a climbing speed travelling down the entrance passage must succeed on a DC 12 Dexterity saving throw or slip on the muddy slide. Characters using a rope gain advantage on the saving throw. Characters who fail this save fall prone and immediately slide down the passage, where they are deposited onto the \textbf{pit trap} in the Lookout's Burrow (Area 3).

\subsubsection{[3] Lookout's Burrow}
\begin{DndReadAloud}{}
The sound of soft rustling and chittering rats fills this muddy earthen den.

\end{DndReadAloud}


\begin{itemize}
\item 2 ratling guttersnipes and 6 ratlings are on watch. Each carries a crossbow and a few prepared Alchemist's Fire (flask).
\end{itemize}

\subparagraph{Pit Trap}
A spiked pit trap ten feet deep is dug in the ground at the base of the Trapped Entrance Passage (Area 2). The hole is covered by a large cloth anchored on the pit's edge and camouflaged with mud and debris. Perception (DC 10) spots the pit. Anyone who moves onto the cloth falls into the pit onto the sharpened wooden spikes at the bottom, and takes 5 (2d4) piercing damage from the spikes, and 5 (1d10) damage from the fall. The edges of the pit are slippery, making escape difficult. Climbing out requires a DC 10 Athletics check. A creature who fails by 5 or more falls back onto the spikes.


\begin{itemize}
\item A \textbf{swarm of rats} lives at the bottom of the pit, and begins devouring anyone who falls within. The swarm attacks any creature in the pit as food.
\end{itemize}

\subsubsection{[4] Ratling Nests}
\begin{DndReadAloud}{}
This tiny dugout cavity contains a cup-shaped nest made from long grasses, hay, mud, twigs, and leaves, and is lined with shredded bits of discarded cloth, fur scraps, bits of leather, paper, and small glimmering trinkets.

\end{DndReadAloud}


\begin{itemize}
\item 8 ratlings and 1 ratling guttersnipe are present in each nest cluster, scattered across the individual nest cavities.
\item Ratlings hoard small trinkets as personal treasure, plus a delerium sliver or a pouch of 10 (3d6) gold. One nest contains a \textit{potion of healing}.
\end{itemize}

\subsubsection{[5] Food Caches}
\begin{DndReadAloud}{}
Haphazardly splayed about this rancid-smelling den are all manner of recently-scavenged dead animals, body parts, and other unsavoury stores.

\end{DndReadAloud}


\begin{itemize}
\item Living meat stays well-preserved, so occasionally prisoners are tied up here for a fresh meal later.
\item Treasure. The north cache contains \textit{thieves' tools}, a \textit{poisoner's kit}, an \textit{herbalism kit}, \textit{alchemist's supplies}, and two Potion of Healing.
\item The south cache holds a still-living captive prisoner - Petra Lang of the Hooded Lanterns. She is unarmed, wounded (with 20 hp remaining), and has expended all her spell slots. However, she's willing to help fight her way out if provided a weapon. Characters can recognize her as a Hooded Lantern lieutenant by her cloak and garb. If she survives, she asks the characters to help her get back to Shepherd's Gate. She can make the trip on her own, but wants the characters to meet her comrades.
\end{itemize}

\subsubsection{[6] Idol of the Rat God}
\begin{DndReadAloud}{}
This chamber contains a bloated effigy of a horrifically obese demonic rat made from rubble, mud, and rat feces, set with two pieces of delerium for eyes. Offerings are scattered about the statue's feet—skulls, bone, daggers, bits of rotting meat, and smaller delerium fragments. Before it lies a blood-soaked altar made from a simple stone slab.

\end{DndReadAloud}


\begin{itemize}
\item 6 ratlings attend the shrine. They flee for help if intruders arrive here.
\item Treasure. The eyes of the idol are delerium fragments (worth 100 gold each). 1d6 delerium chips are scattered about the altar. Several Spell Scroll are placed at the altar here—\textit{bane}, \textit{faerie fire}, and \textit{inflict wounds}.
\end{itemize}

\subsubsection{[7] Gathering Den}
\begin{DndReadAloud}{}
Bales of hay and animal skins surround a glowing bonfire in this large foetid chamber.

\end{DndReadAloud}


\begin{itemize}
\item The Rat Prince is often found here carousing with six other ratlings.
\item The ratlings meet here to discuss their nefarious and ill-conceived plans. There are several sets of tools in this room—including \textit{tinker's tools} and Alchemist's Supplies, which the ratlings use to create makeshift weapons and improvised explosives.
\item This burrow tunnel opens into the sewers. The ratlings flee here and scatter if they are wounded or cannot defend their lair.
\end{itemize}

\subsubsection{[8] Seer's Den}
\begin{DndReadAloud}{}
In this spacious chamber is a large and sumptuous basket nest. Woven over the nest is a canvas made from stitched together scraps of fine cloth, twigs, and choice leaves. Before the nest is a rough stone plinth upon which is placed a hunk of silvery rock about a foot wide. A protruding delerium crystal is embedded in its surface.

\end{DndReadAloud}


\begin{itemize}
\item Squeaks the Seer keeps her nest here. She has an \textbf{imp} familiar who remains in rat form. In its true form, it resembles a disgusting winged and hairless rat.
\item The delerium crystal is here. It's still stuck in a rocky chunk: a combination of stone from the tavern basement, melted earth, and meteoric rock. The Ratlings dislodged it and brought it here because they couldn't break the crystal from the rock themselves. The chunk of rock weighs about fifty pounds.
\end{itemize}

\DndQuote%
{}
{''This place was a dive before the meteor. Infested with rats, terrible ale. Honestly, maybe it's been improved?''}
{Sebastian Crowe}
\DndQuote%
{}
{''I knew it was a dive bar but I didn't think I would fall head first into an actual rats' nest.''}
{Pluto Jackson}
\begin{DndSidebar}[float=!b]{Developments}
Characters who successfully explore the Rat's Nest might trigger the following developments:



\subsection{The Delerium Crystal}
River upholds her end of the bargain faithfully. If the characters were swift and professional in their dealings with her, she gives them a single dose of \textit{aqua delerium} as an additional reward. She encourages them to take a few days off to recuperate, promising to contact them shortly (in 1d6 days) with another job. On the other hand, River is much less impressed if the characters fail to recover the delerium crystal. She still offers to buy any found delerium from the characters, but decides to wait and reassess the characters' capabilities before offering them further work.

Characters who aren't already in contact with the Amethyst Academy still have a valuable find on their hands, which could put them in contact with any faction. The Hooded Lanterns know the delerium crystal is worth a small fortune, but they're happy to let the characters keep it as an extra reward for their efforts. Ansom or Petra could suggest the characters sell the crystal in Emberwood Village to Fairweather Trades and Exports. Another NPC might suggest the characters ask around the village for contacts with the Amethyst Academy or the Queen's Men to sell the crystal to them. Alternatively, the characters might encounter a knight of the Silver Order, who encourages the characters to destroy the delerium. Finally, a pilgrim of the Falling Fire may claim the crystal has a holy purpose, and try to convince the characters to bring it to Lucretia Mathias at Saint Selina's Monastery as an offering.



\subsection{Rescuing Petra Lang}
If Petra Lang survives, she invites the characters to return with her to Shepherd's Gate. Her comrades are thrilled at her safe return, and Lord Commander Elias Drexel takes note of the characters. He won't meet them personally just yet, but instead makes arrangements for Ansom or Petra to contact them for another mission within 7 (2d6) days.

Ansom reminds the characters that their arrangement was a personal matter, and technically not official Hooded Lanterns business. As such, this mission alone won't earn the characters free passage through Shepherd's Gate. However, should the characters come to the gate in an emergency, the Hooded Lanterns will help them. In such an event, they'll overlook the tariffs "just this one time" and let them through.

If the characters fail to rescue Petra, Ansom believes them to be utterly incompetent. He hesitates to involve them further in Hooded Lanterns affairs, and bitterly mistrusts their capabilities in any future situations.

If Petra dies, but the characters weren't yet working with the Hooded Lanterns, it's up to them whether or not they reveal her fate. Regardless, her comrades are shaken by her disappearance and death. In grief, the hard-hearted Lord Commander never speaks of her again. Ansom presses on, but it feels like he's missing a part of himself.



\end{DndSidebar}

\section{Reed Manor}
\DndQuote%
{}
{''Oscar Yoren had irreconcilable differences with the Academy. I'm not privy to all the details, but the Directorate cut support for his research. When Yoren absconded with several volumes on conjuration, transmutation, and necromancy, we discovered he'd taken up unsavoury methods for acquiring "raw materials'' instead. He was declared malfeasant, and sought refuge in Drakkenheim. He's survived in the city ruins for about a decade since. We don't understand how he has managed to do so, but we want to find out. He knows something about the Haze.''}
{River}
Reed Manor stands in a now-dilapidated upper-class residential neighbourhood on the north side of Drakkenheim, just off College Road. The fate of the original occupants is unknown, but it is now occupied by the malfeasant wizard \textbf{Oscar Yoren} and his apprentices.

\includegraphics[width=0.4\textwidth]{images/../5e.tools/img/adventure/DoDk/065-05-013.reed-manor.png}
\subsection{Overview}
In this adventure, the characters encounter \textbf{Oscar Yoren}. Yoren has developed valuable research about the Haze and delerium, but the mage is not necessarily willing to share his discoveries freely. Most notably, Yoren has developed a formula for brewing \textit{aqua expurgo}, a potion that offers protection from contamination. Characters must decide how they'll deal with Oscar Yoren: they may bargain with him to obtain his research, steal his work through stealth and trickery, or attack the mage and destroy his experiments. Depending on the characters' actions, Oscar Yoren could become a recurring character in the campaign.

\subsubsection{Adventure Hooks}
\subparagraph{Stolen Homework}
River hires the characters to obtain Oscar Yoren's Research. She offers the characters a reward of 250 gold, and one magic item crafted by the Amethyst Academy—a \textit{pearl of power} made from a delerium shard.


\begin{itemize}
\begin{DndReadAloud}{}
"We located Oscar Yoren by trailing a group of expelled students who sought him out. They may be working as his apprentices now."

\end{DndReadAloud}

\begin{DndReadAloud}{}
"We don't care what you have to do to obtain this information, but you must bring back something tangible: a spellbook, his journals, or samples of his work. In the worst case scenario... all we require is his brain."

\end{DndReadAloud}

\end{itemize}

\subsubsection{Potions Inspectors}
\DndQuote%
{}
{''For several years now, we've been buying potions from a man named Oscar Yoren. He's a bit of a strange recluse, but his work is reliable and he doesn't seem to be bothering anyone. His potions are much cheaper than what the Amethyst Academy charges. Normally, he sends one of his apprentices to Shepherd's Gate every Monday with a new shipment, but they didn't arrive last week. The week before we got half the normal order. Worse, we've discovered that several of the potions we previously received simply didn't work, and a few were even contaminated. Something's going on, and we'd like you to investigate.''}
{Petra Lang}
Lieutenant Petra Lang (or Ansom Lang) contacts the characters in Emberwood Village. She offers the characters 250 gold pieces for the task. Furthermore, the Hooded Lanterns will offer characters a weapon from their stockpile: either a \textit{+1 longsword} or a \textit{+1 longbow}.


\begin{itemize}
\begin{DndReadAloud}{}
"Our past dealings with Yoren have generally been positive, but it's possible that he cleans up his operation whenever we come around to check on him."

\end{DndReadAloud}

\begin{DndReadAloud}{}
"We're worried about these faulty and contaminated potions, and we need you to make sure Oscar Yoren isn't cutting corners."

\end{DndReadAloud}

\begin{DndReadAloud}{}
"However, if something happened to him, we need to know. We still need Oscar Yoren's potions. We don't have access to a lot of our own magic, and we're stretched pretty thin right now, especially with several sick Lanterns."

\end{DndReadAloud}

\begin{DndReadAloud}{}
"He lives in an old house on the north end of town, I guess he uses some kind of magic to keep himself healthy. We think he might have been kicked out of the Amethyst Academy, but that's their business, not ours."

\end{DndReadAloud}

\end{itemize}

If the player characters don't take this mission, Petra dispatches a group of scouts instead. They show up while the characters are exploring the manor.

\subsubsection{The Buyout}
\DndQuote%
{}
{''The greencloaks have a supplier who's been making potions for 'em at a pretty steep discount. He's some sort of cast-off from the Academy. A guy like that can be a valuable asset... my employer would like to acquire this asset. That's where you come in. Make him an offer he can't refuse.''}
{Blackjack Mel}
Blackjack Mel invites the characters to meet him at the Skull \& Sword Taphouse in Emberwood Village. He wants the characters to convince Oscar Yoren to break ties with the Hooded Lanterns and start providing potions for the Queen's Men instead. He offers the characters 250 gold as payment, plus some "fancy new shoes"—\textit{boots of striding and springing}.


\begin{itemize}
\begin{DndReadAloud}{}
"Rough him up, but don't kill him. He's gotta be able to work."

\end{DndReadAloud}

\begin{DndReadAloud}{}
"We want the same price he's giving the greencloaks, but he has to cut them off. As long as he delivers the goods, the Queen of Thieves can provide all the protection he needs. If he refuses, we'll break his legs and steal his work."

\end{DndReadAloud}

\end{itemize}

\subsection{Interactions}
\subsubsection{Oscar Yoren}
This malfeasant wizard fled the Amethyst Academy ten years ago. Oscar Yoren has since taken up residence on the north side of Drakkenheim in Reed Manor, and survived on the outskirts of the city through his grim intellect, ruthless self-interest, and willingness to practice unsavoury forms of alchemy and necromancy. Several other outcast mages have since joined him as apprentices. See Reed Manor for details on Oscar Yoren's discoveries with delerium.

Oscar Yoren is a short, wide man in his mid-fifties. He wears slovenly leather breeches and a filthy black overcoat. He has shoulder-length greasy black hair and a thick beard of wiry unshaven hair creeping down his flabby neck. His skin is covered in sores and pasty blotches, with bulging blood veins under his neck. His sunken bloodshot eyes glare with a purple-green light. He has yellow teeth and fingernails.


\begin{description}
\item[Personality Trait.]
I talk down to anyone who isn't as intelligent as I am - which is everyone. I demand attention when I'm speaking, and make thinly-veiled threats when I don't get my way.
\item[Ideal.]
The Amethyst Academy doesn't understand my methods, but I cannot unlock the secrets of delerium if I am hindered by petty morality and useless ethics. If I had proper facilities and resources, unimaginable power would be mine.
\item[Bond.]
My apprentices are useful as long as they are loyal and don't question me, but disposable the moment they aren't. Only my own survival matters.
\item[Flaw.]
I hold vindictive grudges against everyone who ever wronged me. I'll have my revenge someday, but I fear acting against those with more power or resources than me. I must hide for now, as I'm terrified the Academy will send assassins to kill me and steal my brilliant research.
\end{description}

\textbf{Oscar Yoren} uses the game statistics of a human \textbf{mage}. He has darkvision 60 ft, 72 hit points, and resistance to necrotic and poison damage. He has the following spells prepared:

\begin{DndReadAloud}{}
Cantrips (at will) - \textit{chill touch}, \textit{mage hand}, \textit{ray of frost}

1st level (4 slots) - \textit{false life}, \textit{mage armor}, \textit{magic missile}, \textit{shield}

2nd level (3 slots) - \textit{invisibility}, \textit{misty step}, \textit{web}

3rd level (3 slots) - \textit{animate dead}, \textit{counterspell}, \textit{lightning bolt}

4th level (3 slots) - \textit{blight}, \textit{polymorph}

5th level (1 slot) - \textit{cloudkill}

\end{DndReadAloud}

Oscar Yoren's \textit{spellbook} contains all the spells he has prepared, plus \textit{alarm}, \textit{find familiar}, \textit{glyph of warding}, \textit{contact other plane}, \textit{dispel magic}, \textit{dimension door}, \textit{magic circle}, and \textit{planar binding}.

Thanks to long-term gradual exposure to contamination over the past decade, Oscar Yoren and his apprentices have gained a limited resistance to the Haze. They aren't immune to contamination, but they may rest in Drakkenheim normally. Oscar Yoren is working on how to reproduce these results, but his current method requires several years of risky and intense self-exposure to delerium.

Oscar Yoren and his apprentices are secretive, but not outwardly hostile if characters approach them peacefully. Instead, the apprentices offer to sell the characters potions. They admit characters into the dining room to talk, but not further into the manor. The apprentices do business in a curt yet cordial manner (see "Potion Brewery" below).

When he meets with the characters personally, Oscar Yoren bluntly demands to know who sent them. He staunchly refuses to work with anyone who admits to working with the Amethyst Academy or the Silver Order, and asks them firmly to leave. Otherwise, he sees an opportunity. He plans to use the characters to recover delerium, monster parts, and arcane reagents for him so he can advance his research, and then conveniently dispose of them once they become troublesome.

Oscar Yoren uses his research as leverage, revealing some of his discoveries about delerium as proof of his knowledge (see the sidebar on "Oscar Yoren's Research"). In exchange for ensuring his freedom, sparing his life, and making fair payment in coin and delerium, he offers to create \textit{aqua expurgo} for the characters. He overstates the efficacy of \textit{aqua expurgo}, but does not reveal its side effects. He scoffs at the idea that anyone else could create the potion. Yoren does not willingly give up his writings or formulas under any circumstances. Instead, he lies, claiming the most important secrets are "in his head."

If the characters agree to work with him, he explains that eldritch lilies are needed to create the potion, and that the plants may only be found in Queen's Park (see chapter 7). If the characters return with the flowers, he makes good on his end of the deal as long as he is paid and unmolested.

Oscar Yoren is ambivalent about working for either the Queen of Thieves or the Hooded Lanterns, but explains that he's devoting all his attention to his new research. He's only interested in continuing to work with either if they're willing to cough up extra money and muscle to help him make more \textit{aqua expurgo}.

\subparagraph{Potion Brewery}
Oscar Yoren and his apprentices brew and sell magical potions to fund their work. They sell uncommon and rare potions for twenty five percent off the normal price. Most of their potions are cheaply made with contaminated ingredients. When consumed, roll 1d20. On a one, the creature consuming the potion gains one level of contamination. If provided the right ingredients, they can also create \textit{aqua delerium} and \textit{aqua expurgo} (see Appendix D).

\subparagraph{Yoren's Apprentices}
Several expelled Academy students have since joined him as apprentices: Marco, Gemma, Taryn, and Bolter. Each is a \textbf{hedge mage} and has a familiar: a \textbf{bat}, a \textbf{cat}, an \textbf{owl}, and a \textbf{rat}, respectively.

\subsection{Area Details}
\includegraphics[width=0.4\textwidth]{images/../5e.tools/img/adventure/DoDk/066-map-5.05-reed-manor.png}
\includegraphics[width=0.4\textwidth]{images/../5e.tools/img/adventure/DoDk/067-map-5.05-reed-manor-player.png}
\subsubsection{[1] Manor Grounds}
\begin{DndReadAloud}{}
Reed Manor is a run-down two-storey stone house surrounded by a rusting wrought-iron fence and an overgrown yard filled with shrubberies and brambles. The long, thin windows are boarded over from the outside, but light flickers within the first floor. Several chimneys rise up from the slate-shingle rooftop, a light trail of smoke billowing from each. A collapsed small outbuilding, perhaps once a stable or shed, borders the yard. Sweeping around the estate grounds, there are two doors leading into the manor from the back. Both have been nailed shut with rotting planks.

A short cobblestone pathway leads through the yard up to a small set of steps leading to the central double doors. A large, round iron knocker hangs on each. Before the doors is a fountain. The central decorative statue has crumbled away into the basin, which is filled with brackish rainwater and algae. A singular obese toad basks beside the water, croaking loudly.

\end{DndReadAloud}


\begin{itemize}
\item There is a secret entrance to the lower level hidden in the bushes. Any physical inspection within five feet discovers it, but it's sealed with an \textit{arcane lock} spell.
\end{itemize}

\subsubsection{Groundskeepers}
\begin{DndReadAloud}{}
Outside the manor are two hulking figures wearing heavy brown overcoats, and both have hoods drawn up over their bandaged faces. One stands listlessly at the front gate, while the other lumbers in a patrol around the manor grounds.

\end{DndReadAloud}


\begin{itemize}
\item A pair of ogre zombies patrol the grounds. One stands at the gate. If approached, they lurch around and growl "What... you... want..." then wait for a reply.
\item The zombies respond to any mention of "potions" by replying "follow... this... way..." and escort characters to the front door. They knock loudly three times, which signals Marco and Gemma to answer the door and talk business.
\item Otherwise, the zombies reply "You... not... welcome... here... go... away..."
\item After issuing the warning, the zombies attack one round later unless they are ordered to stand down by Oscar Yoren or one of his apprentices. They're ordered to cry out loudly when they do, which automatically alerts the manor occupants that something is amiss. The zombies only leave the manor grounds to pursue a foe making ranged attacks against them.
\item The toad is Yoren's \textbf{quasit} familiar, Drovik. If it detects trouble, it swims down the fountain pipes to alert Oscar.
\end{itemize}

\subsubsection{[2-8] Manor House}
\begin{DndReadAloud}{}
\subparagraph{[2] Entrance Foyer}
This once well-appointed foyer is covered in leaves and mud tracked in from the city streets. A decaying tapestry depicting the Reed Family Crest hangs along the walls.

\end{DndReadAloud}

\begin{DndReadAloud}{}
\subparagraph{[3] Alchemical Workshop}
What was once a sitting room is now a makeshift alchemist's workshop. Furniture draped in dust-covered cloth has been pushed aside to make space for two crowded tables set with alembics, burners, distillers, and other equipment for brewing potions. A bubbling cauldron hangs in the fireplace. Strings of herbs dangle from the ceiling, and jars of ingredients are scattered about the shelves alongside racks of empty or half-filled vials. An acrid scent fills the room.

\end{DndReadAloud}


\begin{itemize}
\item Marco, Bolter, Gemma, and Taryn work here during the day.
\item Treasure. Enough tools to assemble an Herbalism Kit, a \textit{poisoner's kit}, and Alchemist's Supplies can be scavenged from the equipment here. The array of raw materials are suitable for spells or other concoctions worth 250 gold total. The cauldron contains enough liquid to fill 1d6 Potion of Healing.
\end{itemize}

\begin{DndReadAloud}{}
\subparagraph{[4] Study}
This room smells of moldy books. A mahogany desk in one corner is covered in crumpled paper and open books, with a torn velvet armchair set before it. The two candelabras don't look like they've been lit in ages. Tall bookshelves are packed with tomes, statues of a variety of fantastic creatures, and the odd portrait or knick-knack. Sliding wooden ladders mounted on rails allow easy access to higher shelves.

\end{DndReadAloud}


\begin{itemize}
\item Most of the books are ruined, but some remain in decent condition.
\item DC 10 Investigation. There are signs of traffic in this room despite the disuse. A clear pathway has been left in the room, allowing access to the far bookshelf.
\item Any physical inspection of the bookshelf reveals that it slides open to the north, revealing a staircase to the chambers below.
\end{itemize}

\begin{DndReadAloud}{}
\subparagraph{[5] Dining Room}
A warm fire burns in the large fireplace, the crackling flames sending the soft smell of wood smoke throughout the room. A black marble mantelpiece above the fireplace has a framed family portrait mounted above it. Lining the wood-paneled walls are broken wooden cabinets that contain cracked dishes, silverware, and empty candlesticks. Iron chandeliers hang from the dark mahogany ceiling above a long dining table carved from redwood. Chairs surround the table, with a tall and elaborate chair at each end. The remains of several meals lie scattered about the table.

\end{DndReadAloud}

\begin{DndReadAloud}{}
\subparagraph{[6] Kitchen}
The kitchen is well-stocked and kept in good order, suggesting frequent use. A disgusting stew simmers in a pot on the hearth. Several bottles of rank alcohol are scattered about.

\end{DndReadAloud}

\begin{DndReadAloud}{}
\subparagraph{[7] Bedrooms}
Upstairs are several slovenly and filthy bedrooms, Each contains a stinking, unmade bed, wardrobes filled with torn and mismatched clothing, a soot-stained fireplace, half-melted candles, and a personal chest.

\end{DndReadAloud}


\begin{itemize}
\item The chests in each room hold the apprentices' belongings—their spellbooks and personal trinkets totalling 50 gp each, and an uncommon \textit{potion} or \textit{spell scroll}.
\item Oscar Yoren uses the larger master bedroom, which is also in total disarray. He's locked the room with an \textit{arcane lock} and routinely sets up an \textit{alarm} spell when he sleeps here. Most of his valuable items are kept in the lab below, but he maintains some personal treasure here: a \textit{potion of invisibility} and a Spell Scroll of \textit{dimension door} under his pillow for emergencies.
\end{itemize}

\begin{DndReadAloud}{}
\subparagraph{[8] Cellar}
This dry and well-stocked cellar contains several weeks of provisions, alcohol, and other supplies. Accessed via a floor hatch in the kitchen, the cellar is wholly separated from the Secret Lab.

\end{DndReadAloud}

\subsubsection{[9-12] Secret Lab}
\begin{DndReadAloud}{}
\textit{The stone stairs lead down to a well-traveled stone hallway.}

\end{DndReadAloud}

\begin{DndReadAloud}{}
\subparagraph{[9] Oscar Yoren's Lab}
Two horrifically mutilated and mutated corpses are splayed out upon dissection tables in the middle of this nightmarish laboratory. Attached to the ceiling is a ramshackle mechanical construction consisting of a dozen metallic arms with adjustable joints and knobs: some end in clamps, blades, syringes, magnifying glasses, lenses, or lanterns. Shelves in this room are stocked with an assortment of occult tomes, vials of alchemical reagents, sealed jars containing disfigured body parts suspended in vile preserving fluids, and glass bowls holding glowing delerium shards.

A large writing desk is littered with scrawled diagrams and stacks of notebooks, and other tables are overflowing with an assortment of alchemical equipment. The stone floor is covered with faded bloodstains and discolored splotches.

\end{DndReadAloud}


\begin{itemize}
\item The main doors of the lab are protected via \textit{alarm} and \textit{glyph of warding} spells that trigger when anyone other than Oscar Yoren or his apprentices open the door.
\item The corpses of a \textbf{haze husk} and a \textbf{delerium dreg} lie on two operating tables in this room. Their identities are unknown.
\item The bookshelves contain tomes on anatomical lore, magical phenomena, curses, diseases, and poisons. Many would be quite valuable to the right buyer.
\item Treasure. In addition to his spellbook and personal possessions, Oscar Yoren keeps several \textit{potions} and Spell Scroll in this room.
\item \textbf{Oscar Yoren} works here during the day.
\end{itemize}

\subparagraph{[10] Well Cistern}
This pool of water is open to the well above. Oscar Yoren keeps the door locked, but there's a hole in the bottom of the door large enough for his familiar to sneak through.

\subparagraph{[11] Meditation Chamber}
In this chamber is a large stone statue depicting the goat-headed Demon Lord of the Undead, set before a pentagram-shaped magic circle on the floor. Several candles and candelabras illuminate the chamber with a sickly arcane glow.


\begin{itemize}
\item Oscar Yoren uses this chamber to cast \textit{contact other plane}, \textit{planar binding}, and \textit{scrying} spells.
\item He hopes one day the Demon Lord of the Undead might teach him the secrets to becoming a lich, and in the meantime wishes to conjure up fiendish servants using delerium crystals.
\item Treasure. A \textit{spell scroll} of \textit{contact other plane} and a pouch containing 300 gp worth of \textit{delerium dust} suitable for use as a material component for a spell. In an iron brazier, Oscar Yoren has a delerium crystal he is planning to use as a material component for the \textit{planar binding} spell.
\end{itemize}

\subparagraph{[12] Materials Storage}
Bits of bone and mummified flesh are scattered about this grim and dusty cellar room. Twelve zombies are kept in this chamber. The zombies are under the control of Oscar Yoren and obey his verbal commands. Occasionally, Oscar locks living prisoners here, in which case he seals the door with an \textit{arcane lock}. Currently, however, the chamber isn't locked and the zombies can respond to Oscar Yoren's call.

\begin{DndSidebar}[float=!b]{Developments}
Characters might discover the following information or trigger these events after exploring Reed Manor and interacting with Oscar Yoren.



\subsection{Oscar Yoren's Research}
The following information can be uncovered by interrogating Oscar Yoren or reading his writings:


\begin{description}
\item[Observations on delerium growth.]
"In truth, the crystal is unlike a mineral or gemstone... I have seen the crystals grow like fungus on a dead tree. The process takes some time; weeks, months, years even. I only discovered this by closely observing wild samples near my new home. They do not grow outside the Haze. I wonder if this means there are new delerium deposits in the ruins... perhaps there is more now than when the meteor originally fell."
\item[Notes on using delerium crystals.]
"Just as one cannot make a mast from a thousand saplings, a thousand small shards of delerium are nothing compared to a single large geode. Bigger crystals contain exponentially more arcane power. Fracturing the rock is a wasteful endeavor, as it causes valuable samples to lose their great magical potential. Delerium crystals and geodes are ideally suited for great spells and ambitious magical constructions. However, it has been difficult to find these specimens in the outer ruins. Nevertheless, small twigs are good for kindling, to set a fire, to keep one warm... and so the small shards make excellent alchemical reagents and arcane inks."
\item[On the alchemical properties of Haze-grown flora.]
"The warping effects of delerium have caused new plants and animals to flourish in the city. A half-mad straggler returned from Queen's Park bearing a wondrous flower - an eldritch lily. The pollen and petals soak in contamination to fuel its own germination. Combined with other rare ingredients, I believe it may prove an ideal remedy against contamination."
\item[Aqua expurgo formula.]
Oscar Yoren has developed a formula for \textit{aqua expurgo}, a magical potion that protects against contamination. However, it requires eldritch lilies from Queen's Park Garden. He needs more samples to refine the formula before he can start producing doses of the potion.
\end{description}

Oscar Yoren's research is extremely valuable. Chief amongst them is his formula for \textit{aqua expurgo}. This potent defense against the Haze will permit characters to explore much more dangerous areas in the city. The Hooded Lanterns, Queen's Men, and Amethyst Academy would deeply desire this new potion so they can advance their plans inside Drakkenheim.

The player characters themselves can study Oscar Yoren's research. Any character capable of casting 3rd level spells or higher can learn to make \textit{aqua expurgo} themselves (see the magic items section for details). However, a key ingredient for the potion is an eldritch lily, which only blooms in Queen's Park. If characters want to create this potion, they'll need to mount an expedition to the Inner City to collect this key ingredient.


\begin{itemize}
\item If the Amethyst Academy obtains the formula, they'll begin producing the potion on a periodic basis. The main constraint is the supply of eldritch lilies. This prevents the Amethyst Academy from manufacturing vast quantities of \textit{aqua expurgo}. They'll ask characters to obtain samples, and if possible, seeds or full plants. This way they can attempt to cultivate the flower themselves.
\item The Hooded Lanterns might suggest turning over the research to the Amethyst Academy as a possible inroad to an alliance. They're also willing to work with Oscar Yoren and continue providing him with protection and materials.
\item The Queen of Thieves wants to ensure that these secrets remain known only to her for the time being. She may ask characters to kidnap Oscar Yoren for her so she can force him into making the potion for her own use.
\item The Silver Order is divided over the potion: it's useful, but involves making use of contaminated flora and delerium. They defer to the characters on whether or not to use or destroy it, but don't produce it themselves.
\item The Followers of the Falling Fire possess their own means of guarding against contamination, so they have no interest in the research.
\end{itemize}

If Oscar Yoren survives, but was attacked or mistreated by the characters (or later discovers they deceived him), he will plot his revenge against them. He won't act on these plans until he obtains eldritch lilies, and tries to convince characters to bring him more delerium and other resources, citing the needs for it to further his own work. In truth, he starts using these items to conjure demons, create undead, and conspire against the characters.

If Oscar Yoren dies, the Amethyst Academy approaches the characters about buying his brain and manuscripts (even if the characters were not working for the Academy originally).

If Oscar Yoren is killed and his research is destroyed or abandoned, the formula for creating \textit{aqua expurgo} is lost. However, a player character or another group could also develop the formula again later in the campaign by studying samples of eldritch lilies.



\end{DndSidebar}

\includegraphics[width=0.4\textwidth]{images/../5e.tools/img/adventure/DoDk/068-05-014.oscar.png}
\section{Shrine of Morrigan}
\begin{DndReadAloud}{}
North of Castle Drakken lies a forested ravine known as the Overlook. Deep within these woods lies a copse of verdant trees caught in a persistent autumnal state. Bone-chimes hang from the gnarled branches of blood-red leaves, and even the tree bark is a vibrant crimson-brown. A small clearing lies at the heart of the grove, where five ancient menhir stones decorated in spiralling patterns form a semicircle around an altar of stone stained with ancient blood. Opposite is a barrow mound carved from the rock.

\end{DndReadAloud}


\begin{itemize}
\item This is the Shrine of Morrigan, a place dedicated to an ancient and shadowy old god. Also known as the Phantom Queen, this elusive entity revels in war and trickery, and holds power over life and death.
\end{itemize}

\subsection{Overview}
Few in Drakkenheim have the power to return the dead to life. Characters who do not have the ability to revive the dead by their own ability must seek out alternative solutions. The Shrine of Morrigan is one such place.

\includegraphics[width=0.4\textwidth]{images/../5e.tools/img/adventure/DoDk/069-05-015.shrine.png}
\subsubsection{Adventure Hook}
\subparagraph{Raise the Fallen}
In the event a player character dies in the ruins, Old Zoya, Elias Drexel, the Queen of Thieves, and Eldrick Runeweaver could suggest characters seek out the shrine. Their information is limited to rumours and legends:


\begin{itemize}
\item Rumours tell of an old hermit who dwells in the wooded ravine north of Castle Drakken.
\item The elf is apparently a priest of some sort, and serves an old god of life and death.
\item He can perform a ritual that brings back the dead, but it bears a terrible blood price.
\end{itemize}

\DndQuote%
{}
{''There are many who have lost things in Drakkenheim... friends and family, trinkets and truths...I just hope you don't end up among them in your quest to find yours.''}
{Veo Sjena}
\subsection{Interactions}
\subsubsection{Eoghan Ghostweaver}
This ancient elvish hermit tends the shrine. \textbf{Eoghan Ghostweaver} is a priest of Morrigan, and keeps the old faith of his goddess. He has a sallow face with skin stretched thin over his cheekbones, is garbed in black feathered robes, and wields a \textit{staff of the woodlands}. He uses the game statistics of an elf \textbf{archmage} with the following changes:


\begin{itemize}
\item He has an Intelligence score of 15 and a Wisdom score of 20.
\item He is proficient in Medicine, Nature, Survival and Perception.
\item He speaks Elven and Druidic
\begin{DndReadAloud}{}
\section{Spellcasting}
Eoghan is an 18th-level spellcaster. His spellcasting ability is Wisdom (spell save DC 17, +9 to hit with spell attacks). He has the following druid spells prepared:

Cantrips (at will): \textit{druidcraft}, \textit{mending}, \textit{poison spray}, \textit{produce flame}

1st level (4 slots): \textit{cure wounds}, \textit{entangle}, \textit{faerie fire}, \textit{speak with animals}

2nd level (3 slots): \textit{animal messenger}, \textit{beast sense}, \textit{hold person}

3rd level (3 slots): \textit{conjure animals}, \textit{purge contamination}*, \textit{sleet storm}

4th level (3 slots): \textit{ice storm}, \textit{polymorph}, \textit{wall of fire}

5th level (3 slots): \textit{commune with nature}, \textit{cone of cold}, \textit{greater restoration}

6th level (1 slot): \textit{conjure fey}, \textit{heal}

7th level (1 slot): \textit{resurrection}

8th level (1 slot): \textit{feeblemind}

9th level (1 slot): \textit{foresight}

*spells found in Appendix D

\end{DndReadAloud}

\item He gruffly rebuffs those who approach him, but then becomes curious about why they have come.
\item Three of the trees here are actually treants who defend Eoghan if he is attacked.
\item The druid dwells in the barrow mound behind the shrine.
\end{itemize}

\subsubsection{Life Trade}
If characters inquire about returning the dead to life, he explains that such a feat is possible for him, but the price is high indeed:


\begin{itemize}
\item Morrigan is a covetous deity. In order to cast \textit{resurrection}, his god demands a living sacrifice. However, not any life will do: it must be a humanoid creature with a level or challenge rating equal to or higher than the character who will be resurrected.
\item In addition, a diamond is needed. Delerium could be used as a substitute, but any creature revived using delerium as a material component gains one randomly determined form of Drakkenheim Madness (see Appendix C).
\item Finally, the hermit demands an offering to satisfy himself. Company, a great story, an act of service, a magic item, fine wine - basically, give him a nice gift that would please a grandfather.
\end{itemize}

\section{Smithy on the Scar}
\DndQuote%
{}
{''We can not allow these mining operations to continue. They are funneling these dangerous rocks out of the city and sending them all over Westemär and beyond. We need to stop the spread of this corruption before it's too late.''}
{Knight-Captain Theodore Marshal}
South of the city walls lies an industrial neighborhood known as the Spokes. The region was heavily damaged during the meteor shower, and a large trench cuts through the district. Nicknamed the "Scar", this rough ravine was created by a massive boulder ejected from the initial meteor impact that collided with the walls and ricocheted south. This rogue projectile tore open the ground and levelled many buildings as it skid across the city, only narrowly missing the Spokes Smithy.

This once-humble smithy has been transformed into a fortified mining camp by an intrepid and entrepreneurial family of dwarven prospectors. The four Ironhelm siblings arrived in Drakkenheim a few months ago, and used their wondrous mechanical weaponry to lay claim to the rich deposits of delerium found within the Scar. While this has made them one of the most successful delerium extraction teams in Drakkenheim, they've made a loose agreement with the Queen of Thieves for protection, and now find themselves caught in the crosshairs of all five factions.

\includegraphics[width=0.4\textwidth]{images/../5e.tools/img/adventure/DoDk/070-05-016.smithy.png}
\subsection{Overview}
In this adventure, one or more of the factions dispatch the player characters to deal with the Ironhelm dwarves. Each faction has different goals here, so the party may be asked to drive off the dwarves or negotiate with them. Regardless, it's up to the player characters to decide how exactly to carry out these orders: they might offer their assistance to the dwarves, attack them, trick them, or bargain with them.

\subsubsection{Adventure Hooks}
\subsubsection{Altering the Deal}
\DndQuote%
{}
{''Don't harm one hair in their mangy beards. The miners are as valuable as the delerium itself. Don't break their legs - just break their toys. The dwarves want to know what their money buys them, so show them. You can let them know they owe us twenty percent from here on out.''}
{The Queen of Thieves}
When the Ironhelms arrived in Drakkenheim, they immediately clashed with the bandit gangs who had previously been picking over the delerium in the Scar. After wresting control of the claim, they cut a deal with the Queen of Thieves: the dwarves could mine the Scar unmolested, so long as they paid fifteen percent of their profits to the Queen of Thieves each week for "protection". However, the Ironhelms have since acquired new weaponry, and no longer feel they need to pay such a large cut for protection that has never actually materialised. The Queen of Thieves is furious, and charges the characters with stealing, sabotaging, or destroying the dwarves' gunpowder weapons to make sure they'll have to depend on her services. However, she doesn't want the dwarves' mining operations otherwise disturbed. Cracking a few skulls and roughing them up is fine, but if multiple dwarves are killed, the characters won't be rewarded.

\subsubsection{An Exclusive Contract}
\DndQuote%
{}
{''The Ironhelm dwarves extract more delerium than anyone. We just want to put it to good use.''}
{River}
The Amethyst Academy wants to buy every last bit of delerium the Ironhelms collect. River explains they are willing to pay the market rate plus cover all the dwarves' expenses in return. She emphasises that the Ironhelms must cut ties with whoever else they're currently supplying (which would include their contract with the Queen of Thieves). She asks the player characters to approach the Ironhelms about this deal, and find out if the dwarves have any additional conditions or demands. If the characters can see to those demands as well, so much the better.

\subsubsection{The Hooded Lanterns Offer}
\DndQuote%
{}
{''The dwarves have weaponry, defences, and tactics that we could use to better fortify our camp, and make stronger pushes into the Inner City.''}
{Lord Commander Elias Drexel}
The Hooded Lanterns want to collaborate with the Ironhelms for mutual benefit. If the dwarves are willing to craft guns and weapons for the Hooded Lanterns, the faction is willing to send their soldiers to guard the Spokes Smithy, and let the Ironhelms take up residence with them in Drakkenheim Garrison. Elias Drexel sends the player characters as envoys to negotiate with the Ironhelms and show them the sincerity of their offer. He wants the characters to find out if the dwarves have anything else they want or need to make the deal work.

\subsubsection{The Silver Order}
The Silver Order has resolved to shut down the Ironhelms entirely. Theodore Marshal impresses that he'd rather not see the dwarves killed, but the mining simply must not continue. He charges the characters to convince the Ironhelms to abandon Drakkenheim. If they refuse, the characters must do what is necessary to drive them out: perhaps by destroying their supplies, the smithy fortifications and weapons, or even somehow making the Scar too dangerous for the dwarves to continue mining it. If the characters manage to do so without slaying any of the dwarves, Theodore Marshal promises them an additional boon, but will accept whatever solution the characters invent.

\subsubsection{The Followers of the Falling Fire}
\DndQuote%
{}
{''Our most humble and defenseless pilgrims could readily collect the holy stones found within the Scar. Not only are these dwarves defiling our sacred crystals with their shovels and mining picks, their greed has caused them to murder many innocent faithful.''}
{Lucretia Mathias}
The Followers of the Falling Fire are aghast at the dwarven strip mine, and want the operation to stop entirely. The delerium in the Scar must instead be put to a sacred purpose. Worse, the dwarves have killed several Falling Fire pilgrims who came to the Scar seeking delerium, as the devotees stubbornly refused to back down when threatened by the Ironhelms. Unfortunately, the dwarven guns are simply too much for their pilgrims to overcome on their own, and so the Followers of the Falling Fire need the player characters to be their instrument of wrath.

\subsection{Interactions}
The Ironhelm siblings have hired a company of dwarven veterans and miners to help carry out their operation. In total, the four Ironhelm siblings lead an expedition of roughly thirty individuals. The expedition consists of the four Ironhelm siblings, 8 dwarven veterans, and 12 dwarven thugs. The Ironhelms and oversee the work, while the veterans protect the campsite while the thugs work mining the scar.

The dwarves work in half-day shifts to avoid the effects of contamination and the Haze. At any given time, two of the expedition leaders are present at the camp, along with half the dwarves. The others are resting in Emberwood Village, where they have rooms at the Red Lion Hotel. Characters could opt to engage with the Ironhelms at either the Spokes Smithy or in Emberwood Village. However, because the dwarves are split up, all four need to meet personally to discuss any deal they're presented with first. If the characters want to negotiate with the dwarves, they'll have to wait until a shift changes (when all the dwarves are in the city).

\subsubsection{The Ironhelm Siblings}
The Ironhelm Siblings have relocated from their ancestral home in the Glimmer Mountains to stake a new claim in Drakkenheim. Their old gold mine has run dry, and so delerium is their new commodity.

The siblings share the following roleplaying characteristics:


\begin{description}
\item[Ideal.]
If you work hard and deal fairly, you're sure to make something of yourself in this world.
\item[Bond.]
We have to carry on the family business.
\item[Flaw.]
Words are worth horse spit. We don't trust what anyone says, and judge them only by their actions.
\end{description}

\subsubsection{Reginald Ironhelm}
This dwarf \textbf{mage} wears a trench coat inlaid with gold buttons and has a satchel full of mechanical bits. He wears a bandolier full of bullets, a monocle over his right eye, and has many tools and trinkets hanging from his belt. His bulbous nose and beady eyes accent his huge moustache and mutton chops but, unlike many dwarves, sports a clean shaven chin. His black hair is always slicked back to keep it off his brow and face.


\begin{description}
\item[Personality Trait.]
Reginald is all about the numbers - he trusts facts, figures, and proven results. He is not the type to make a risky decision on a hunch. He prefers calculated and educated results.
\end{description}

\subsubsection{Mordecai Ironhelm}
A dwarf \textbf{assassin} garbed in black cloak and trenchcoat, he has several daggers at the ready along his coat, and his belt contains several vials of poison. He has long black hair, a scar down the side of his face, and a black goatee that trails down to his belt buckle.


\begin{description}
\item[Personality Trait.]
Mordecai is always looking for a better deal or a way to sweeten the pot. He has a charisma to him, even if he is a little rough around the edges. He is blunt with his remarks, but very good at reading people. He knows when to push, and when to put on the charm.
\end{description}

\includegraphics[width=0.4\textwidth]{images/../5e.tools/img/adventure/DoDk/071-05-017.gertrude.png}
\subsubsection{Gertrude Ironhelm}
A dwarf \textbf{gladiator} who wields two golden axes, she has dark brown hair laced into three separate braided ponytails. She has long sideburns that she lets grow down, and a very stern look on her face.


\begin{description}
\item[Personality Trait.]
Gertrude has always trusted her instincts. She listens to her gut and it's gotten her this far. She doesn't mind playing rough and aggressive if it helps her and her family accomplish their goals. She is highly protective of her friends and family.
\end{description}

\subsubsection{Elouise Ironhelm}
A dwarf \textbf{priest} who keeps the ancestral dwarven faith. She is dressed in common grey robes, complemented by thick leather strappings and a breast plate. She has a mechanical trinket hanging around her neck and a belt full of various hammers and tools. She has long brown hair flowing down to her knees, big bright eyes, and a kind yet cunning look upon her face.


\begin{description}
\item[Personality Trait.]
Elouise respects tradition, and wants to preserve her family's ancestral ways. She constantly reminds her siblings they need to stick together. She is the kindest of the siblings and will always try to find a means to peaceful resolution if possible.
\end{description}

\subsubsection{Negotiating with the Ironhelms}
\begin{DndReadAloud}{}
"Promises are worth as much as pig shite. Your actions have currency with us, though. Show your mettle and we can entertain your offer. Protecting this place at night is the hardest part of the job. Prove you'll help us by doing a night shift. Then we'll talk about your little offer."

\end{DndReadAloud}

The dwarves currently have three active deals:


\begin{itemize}
\item They pay fifteen percent of their current earnings to the Queen of Thieves as part of her protection racket.
\item They spend half of their earnings on supplies, equipment upkeep, wages, accommodations, and other business expenses.
\item They sell their delerium at a bulk rate to Fairweather Trades and Exports in Emberwood Village, and so only get about eighty percent of its market value.
\end{itemize}

The Ironhelms find the offers made by the Amethyst Academy and the Hooded Lanterns intriguing. Nevertheless, the dwarves drive a hard bargain on principle, and explain that if they break their agreement with the Queen of Thieves, they're concerned she will act against them more aggressively. As such, the Ironhelms won't accept any agreement that doesn't include proven assurance of protection in the future. They need proof the faction can protect their mining operation from monsters, as well as anyone sent by the Queen of Thieves. They reason that since the faction dispatched the player characters to negotiate with them, then the characters should be the ones who provide this protection... assuming they are up to the task. They ask the characters to protect the Smithy during a night shift to demonstrate their capabilities.

\subsubsection{Sabotaging the Ironhelms}
\begin{DndReadAloud}{}
"Take your demands and shove them up your arse. There's money to be made here and we're not leaving. If you wish to contest our claim, you're free to do so before the barrel of a gun."

\end{DndReadAloud}

Convincing the Ironhelm siblings to simply abandon their operation is nearly impossible. They have invested too much time and energy into fortifying the location and building a functioning business by prospecting the Scar. They laugh in the faces of those who try to tell them to pack up and leave. Instead, characters will need to resort to other tactics to get the dwarves to give up peacefully.

Characters working for the Queen of Thieves, the Silver Order, or the Falling Fire might be looking for ways to disrupt, drive off, or sabotage the Ironhelms without necessarily slaying them. A few possibilities include one or more of the following:


\begin{itemize}
\item Luring additional monsters to attack the Spokes Smithy or occupy the Scar
\item Stealing or breaking their gunpowder weapons
\item Ruining the forge or the dwarves' supplies
\item Stealing or sabotaging the dwarves' mining equipment or wagons
\item Convincing the dwarven hirelings to desert the Ironhelms
\item Getting the dwarves kicked out of the Red Lion Hotel, or getting them blacklisted by the merchants in Emberwood Village
\item Destroying the defenses around the smithy or even the entire building itself
\end{itemize}

Finally, slaying the dwarves is always an option, though a rather murderous one. It's possible that if one or more of the Ironhelm siblings dies, the others might be convinced to give up.

\subsection{Area Details}
The Spokes Smithy is a roofed outdoor blacksmith shop with a small office. It's been heavily fortified and overlooks the Scar.

\includegraphics[width=0.4\textwidth]{images/../5e.tools/img/adventure/DoDk/072-map-5.06-spokes-smithy.png}
\includegraphics[width=0.4\textwidth]{images/../5e.tools/img/adventure/DoDk/073-map-5.06-spokes-smithy-player.png}
\subsubsection{The Scar}
\begin{DndReadAloud}{}
The Scar is three hundred feet long, twenty to thirty feet wide, and ten feet deep. The melted earth has turned to black glass and glittering obsidian sand, and the trench is now among the richest delerium veins outside the city centre.

\end{DndReadAloud}


\begin{itemize}
\item Most of the delerium found here are chips and fragments fused into the earth, and it requires diligent work to extract.
\end{itemize}

\subsubsection{[1-6] Spokes Smithy}
\begin{DndReadAloud}{}
The Spokes Smithy is a relatively small blacksmith shop. It has its covered forge, a small office, a yard that now functions as a campsite, and a series of barricades recently constructed by the Ironhelms to secure the location. They have setup three gun platforms around the perimeter, each a features a massive cannon of dwarven construction..

A thunderous staccato and a whirring mechanical sound echo across the ravine. As the noise dies down, it's replaced by a bellowing voice calling out "All clear!" in dwarven and the sound of clinking hammers and picks against stone. A field of monstrous corpses is splayed out all around: dead ratlings, mutated dregs, and other slain horrors. "Stand down." Calls the bellowing voice "Any beastie gets a bullet, but if you're no monster, you may pass."

\end{DndReadAloud}


\begin{itemize}
\item The Day shift is taken by Reginald, Elouise, 6 dwarf thugs, and 4 dwarf veterans.
\item The Night shift is taken by Gertrude, Mordecai, 6 dwarf thugs, and 4 dwarf veterans
\end{itemize}

\subparagraph{[1] Forge}
An impressive furnace and anvil rests in a covered workshop lined with hanging tools and equipment. The Ironhelms use the forge currently to manufacture bullets for their weapons and fix up their other armour and mining gear. Mordecai can often be found in the forge during night shifts making sure ammo is stocked and ready.


\begin{itemize}
\item \textit{Treasure}: There is a set of \textit{smith's tools} and a large piece of meteoric iron here that could be worked into a weapon or piece of armour by a skilled smith.
\end{itemize}

\subparagraph{[2] Smithy's Office}
This small office has a desk and a small iron safe. There is an old bookshelf along one wall lined with design plans, blueprints, and recipes for various forge-made gear. A book on the desk details the numbers regarding the prospecting and selling of delerium. Reginald can be found in the office during the day crunching numbers.


\begin{itemize}
\item \textit{Treasure}: The safe in the room is locked (DC 20). It contains four gold bars each worth 1,000gp along with a box containing two delerium crystals.
\end{itemize}

\subparagraph{[3] Campsite}
Between the smithy and the scar is a makeshift camp that houses a small bonfire, a few large chunks of debris repurposed into seating, a pair of tents for those needing a short rest or privacy, and an old cooking pot that hangs over the fire. During the day Elouise can often be found here cooking a travelers stew and aiding anyone who has suffered injuries.


\begin{itemize}
\item \textit{Treasure}: Next to one of the tents is a small satchel containing Keoghtom's Ointment and a Spell Scroll of \textit{purge contamination}.
\end{itemize}

\subparagraph{[4] East Barricade}
The larger of the two main barricades covers the open side of the building and exposed campsite and is composed of a small stone fence reinforced by wooden spikes and barbed wire. A makeshift iron gate sits in the middle and several burlap bags of sand and dirt are piled atop the wall where a row of muskets can be perched. Behind the wall are two boxes each containing eight sticks of dynamite that the dwarves occasionally use in the Scar, or to defend the camp. Gertrude can be found here most of the time during the night shift.

\subparagraph{[5] West Barricade}
The West Barricade is much the same as the east but covers less ground and doesn't have a gate. Its purpose is to block the open gap between the long end of the smithy and the Scar.

\subparagraph{[6] Gun Platforms}
Around the smithy are small platforms recently constructed, with a ladder that goes down to the camp below. The platform has a reinforced barricade wall built along the side facing out towards the city and a mounted turret has been bolted to the platform here. Upon each platform is a cannon. There is a stack of cannonballs and several barrels of gunpowder next to each.

\includegraphics[width=0.4\textwidth]{images/../5e.tools/img/adventure/DoDk/074-05-018.anvil.png}
\subsection{Working the Night Shift}
During the night, the characters have four events. Roll 1d8 for each event to determine which hour of the night the attack occurs: if you get a duplicate result, bad luck for the players! In such a case, the two attacks happen 10d6 minutes apart.

The dwarves are reluctant to let characters use the gun themselves, but can be convinced by a character who makes a good argument or who displays sophisticated engineering knowledge (DC 18). If the character insists on more illumination, they experience an additional encounter during the evening. If the players take the night shift, the normal night shift takes a break and the characters are accompanied by one of the Ironhelm siblings as well as 4 dwarf thugs and 2 dwarf veterans.

The dwarves rely on their darkvision at night. They remind characters that at night, light sources will attract more monsters. The only light source the dwarves allow at night is their bonfire.

The mists of the Haze recede at night, and a billowing corona of octarine light appears above the crater. This creates dim illumination in the area of the camp and the Scar, but creatures and objects further than 120 feet away are totally obscured unless the characters have features or traits that allow them to see further in darkness.

\subsubsection{Events}
\subparagraph{Swarm of the Walking Dead}
As night falls and the glow of the delerium clusters that pepper the Scar cast an octarine light across the lonely smithy, the sounds of shuffling and moans begin to fill the air - distant at first, but getting closer, and coming from all directions. 12 haze husks emerge on the first round, and 1d6 more haze husks appear each round afterwards for the next minute.

\subparagraph{Haze Hulks}
A series of loud wails erupt through the streets and the sounds of lumbering limbs cascade towards the smithy - something big is coming. Three haze hulks emerge fromthe haze and charge towards the smithy.

\subparagraph{Chimera}
A \textbf{chimera} swoops down to attack. "We're out of ammo!" rings out from the gun platform. None of the firearms nor the cannon are usable for any further encounters. It's at this moment that a terrible screech is heard from above. 1d4 rounds in, 10 haze husks emerge from the haze as well.

\subparagraph{Ratling Horde}
Attracted by the sounds and chaos, a horde of 20 ratlings emerge in hopes of scavenging off the dead corpses littering the field.

\begin{DndSidebar}[float=!b]{Developments}
\subsection{Accepting the Offer}
If the characters successfully defend the site during the night shift, the Ironhelms are happy to accept an offer made by the Amethyst Academy or the Hooded Lanterns. As a result, characters have won a new ally. However, Reginald and Elouise tell the characters that it's now their responsibility to help the dwarves out if they need defenders in the future. They'll contact the characters via a \textit{sending} spell if they ever need them.

At a pivotal moment in the campaign, the Queen of Thieves (or another rival faction) might time an attack against the dwarves with another crucial event. If you want to up the stakes during an otherwise straightforward mission, the Ironhelms call for aid when the players are already stretched for time. Alternatively, they might request aid immediately after an intense battle when the characters are low on resources.

The Amethyst Academy and the Hooded Lanterns reward characters with a faction boon for their excellent work. At your option, the Ironhelm dwarves might be willing to teach the characters how to use gunpowder weapons, and give them each a gun and ammunition. Later in the campaign, the dwarves could be called up to provide weapons for a major offensive as well.



\subsection{The Queen's New Deal}
If the characters sabotage the dwarven weapons but don't otherwise cause major disruptions to the dwarven mining operation, they begrudgingly acquiesce to the Queen of Thieves demands. She rewards characters with a faction boon, and lets characters keep any guns or ammunition they might have absconded with (or remarks on how foolish the characters were if they did not purloin the weapons for themselves).

However, in the event the sabotage went too far or any of the Ironhelm siblings were killed, the Queen of Thieves is greatly displeased and informs the characters that they are now expected to make up the difference in profits. She quotes an outrageous figure, saying the characters have cost her thousands of gold. She expects them to adequately make up the difference in a timely fashion, and uses this fact to blackmail them into doing her dirty work for as long as possible.



\subsection{Driving out the Dwarves}
If the dwarven mining operations are extensively sabotaged, the Ironhelm siblings are financially ruined. The cost of buying new equipment, hiring new labourers, and building new defenses is simply too much to bear. Combined with the threat of further retribution from another faction, this is sufficient to convince them to leave.

However, if it seems like the dwarves could recuperate or repair whatever damages the player characters caused, they might seek out aid from another faction to protect them from future attacks. The next time the characters arrive, a strike team from a rival faction is on site as additional muscle.



\subsection{Wiping out the Dwarves}
Should the characters decide to simply slay the dwarves, they're in for a bloody battle. If the tide of battle turns against them, Reginald and Elouise contact their siblings outside Drakkenheim via a \textit{sending} spell. Since the dwarves work in two shifts, the characters will have to deal with both groups. If one group is slain, the other does not exact immediate vengeance - instead, the survivors seek out help from the Hooded Lanterns, the Amethyst Academy, or the Queen of Thieves. Each faction is more than willing to send a strike team of their own to help the Ironhelm dwarves.

The next time a rival faction sends out a strike team against the party, add the remaining Ironhelm siblings and 1d4 dwarf veterans to the strike team.



\end{DndSidebar}

\section{Stick's Ferry}
\includegraphics[width=0.4\textwidth]{images/../5e.tools/img/adventure/DoDk/075-05-019.sticks-ferry.png}
The major bridges spanning the Drann River are all located within the Inner City. Before Drakkenheim was destroyed, crossing the river via these bridges required merchants and travellers to pay the gate tax to enter the city proper, or otherwise make a detour fifteen miles south or west to other river crossings. During this time, an old boatman named Pixis Stick made his living illicitly ferrying folk across the Drann River on the southern side of town using his rowboat. Often simply called the "Ferryman", he enjoyed meeting the urchins and ne'er-do-wells who inhabited Drakkenheim, and swapping stories and gossip with penny-pinching merchants.

When the meteor fell, the Ferryman became a \textbf{haze husk}, one of the tortured and mindless undead creatures that drift aimlessly about the city. However, he retained a strange semblance of his past memory, and is not hostile to passersby. Instead, the Ferryman takes paying passengers back and forth across the Drann River just as he did in life. He now only communicates in low moans and grunts, but the Ferryman still seems to enjoy company, conversation, and of course, gold.

\subparagraph{Rejuvenation}
If the Ferryman is destroyed, it regains all its hit points in one hour unless a \textit{remove curse} spell is cast on its remains.

\subsection{Taking the Ferry}
The Ferryman inexplicably rows up to the water's edge whenever characters arrive at either bank of the river on the southwest outskirts of Drakkenheim. If the characters approach him, the Ferryman grunts, gestures to his boat, and holds out his hand. He waits until each person who enters his rowboat gives him anything resembling a gold piece. Once every passenger has paid their fare, the Ferryman takes them across the water. He only transports characters to the opposite side of the river, and cannot be commanded to do anything else. He growls if someone else tries to take the oars from him, and only fights to defend himself or protect his boat.

\subsection{Ferryman's Treasure}
When alone, the Ferryman mindlessly returns to his old stash to hide his gold. He keeps a duck-shaped buoy in the middle of the river, which connects to several ropes attached to sacks of gold hidden below the water. All told, the Ferryman's treasure consists of 2,374 gold pieces placed in about a dozen satchels.

\DndQuote%
{}
{''Stick may not be much of a talker, but I'd take riding with the undead over risking a dip in that water any day!''}
{Veo Sjena}
\DndQuote%
{}
{''So we shoot undead on sight... unless they drive boats?''}
{Sebastian Crowe}
\chapter{Inside the Walls of Drakkenheim}
\DndDropCapLine{T}{}he impressive city walls encircle Drakkenheim both north and south of the Drann River. Although a section of city walls was destroyed by the meteor, the crater left behind is widely considered impassable due to the extreme contamination present. The walls are otherwise intact, impeding access to the Inner City for adventurers and faction agents alike except via five gatehouses.


Many baroque architectural flairs decorate the walls. The battlements feature leering demonic faces, claw-like spikes, draconic carvings, and dutiful winged gargoyles that hold up machicolations. Yet these grotesque statues are no mere ornamentation. When Drakkenheim was a living city, these magical guardians could spring to life under the monarch's will. Now, they serve the twisted whims of the Haze in mockery of their final order to safeguard the city, and viciously attack any who dare trespass over the walls!

\includegraphics[width=0.4\textwidth]{images/../5e.tools/img/adventure/DoDk/076-06-001.inside-the-walls.png}
\section{Overview}
Characters should reach at least 5th or 6th level before they attempt getting over the city walls. Even still, constantly battling animated guardians each time they travel into the Inner City will prove extremely draining on the character's resources, so they'll want to secure passage through one of the five city gates to access the city centre long-term. However, the gates are occupied by faction agents and monsters alike:


\begin{itemize}
\item Champion's Gate is nominally controlled by the Followers of the Falling Fire.
\item College Gate is now the site of a terrible dimensional rift.
\item King's Gate has been occupied by the Troll King.
\item Shepherd's Gate is defended by the Hooded Lanterns.
\item Temple Gate has become a makeshift stronghold for a massive garmyr warband.
\end{itemize}

It's up to the player characters to clear out one of these locations, make a deal with their faction occupants, or invent another solution. Rival adventurers are quick to point out that using the sewers to sneak into the city isn't a smart alternative, as far more dangerous monsters lurk there.

\subsection{Adventure Hooks}
\subparagraph{Over the Top}
The characters should learn about the hazards of the city walls through rumours or their own experiences, and search for solutions by their own initiative. The city walls can be introduced as a complication for any other adventure hook for the Inner City locations, as faction lieutenants and leaders note that any mission to the city centre will be much easier if the characters have a reliable way in and out.

\subparagraph{The Battle of Temple Gate}
At a key moment in the campaign, the Knights of the Silver Order will mount a major offensive to seize Temple Gate from the garmyr occupying it. They'll invite collaboration from the player characters as well as other factions in this joint effort.

\subsection{Encounters}
Characters who do try to pass the walls may trigger one of the following encounters:

\subsubsection{Wall Gargoyles}
7 (2d6) Wall Gargoyles animate and attack characters who dare attempt any of the following:


\begin{itemize}
\item Climbing more than halfway up the walls or towers
\item Flying within one hundred fifty feet of the city walls horizontally, or over them at any height
\item Lingering on the walls or towers for more than ten minutes
\end{itemize}

The gargoyles attack directly, viciously, and without remorse. If they reduce their quarry to 0 hit points, they spend their next turn tearing apart the creature's body before moving on to another target. The gargoyles pursue their quarry into the city to any distance until they or their quarry are destroyed, but give up pursuing characters more than three hundred feet \textit{outside} the city walls. However, Wall Gargoyles do not attack characters attuned to a \textit{Seal of Drakkenheim} (see Appendix D).

\subsubsection{Tower Dragons}
Astute characters will note that there is no barrier spanning the section of the city where the walls meet the Drann River. This section of the city is not undefended, however. Instead of gargoyles, two Tower Dragons perched upon the bastion towers flanking the river attack anyone crossing into the city via the river. They employ similar tactics to the Wall Gargoyles as described above.

\subsection{Area Details}
Described first are the general features common to the walls, towers, and gates, followed by specific details on the inhabitants of each gate:

\subsubsection{Walls}
\subparagraph{Walls}
The stone masonry walls are forty feet high and ten feet thick with a machicolated battlement facing outward. Murder holes and overhanging artillery posts feature throughout, and areas of the wall near gates have wooden hoardings or shingled rooftops for additional defenses.

\subparagraph{Towers}
Sixty-foot-tall towers thirty feet in diameter are spaced throughout the wall. Some are topped with battlements, others have spired rooftops. Each tower is decorated with a majestic bronze dragon. The largest of these towers are built where the walls meet the Drann River.

\subsubsection{Gatehouses}
There are five barbican-style gatehouses that lead into the city proper. Each is composed of two squat bastion towers sixty feet tall and forty feet in diameter. The towers flank a fifteen-foot-wide double portcullis gate. While the gatehouses themselves have largely similar construction, their current occupants vary.

Each gatehouse has four levels, each approximately fifteen feet high. Outer walls are four to five feet thick. The interior floors are wooden, but the central spiral stair is stone. The arrow slit windows are only a few inches wide, and they provide \textit{three-quarters cover} from those outside and below (+5 AC). Regular oaken doors are bound with wrought-iron and feature heavy locks and crossbars.

\includegraphics[width=0.4\textwidth]{images/../5e.tools/img/adventure/DoDk/077-map-6.01-gate-houses.png}
\subparagraph{Gate Passage}
The gate tunnel is approximately fifteen feet wide, twenty five feet tall, and twenty five feet long. At each end are a wrought-iron portcullis and massive wooden doors bound with steel. The doors can be held shut via a massive bar that slides into place like a lever.

\subparagraph{Tower Stairs}
At the core of each tower is a stone spiral staircase. The staircase is five feet wide and has landings on each level.

\subparagraph{Ground Level}
The gatehouse can be entered from a barred doorway at a landing atop a ten foot high sloping stairway on one of the bastion towers. The base level of the left tower is not accessible from the outside. These chambers are typically used as sleeping quarters or a gaol.

\subparagraph{Winch Level}
It takes a team of operators to raise or lower the gate, requiring four successful DC 12 Strength checks, but the winch can be released to cause the gate to slam shut immediately.

\subparagraph{Battlement Level}
Accessed via doors. A creature taking a firing position behind the battlements has \textit{three-quarters cover} from attackers below the walls (+5 AC).

\subparagraph{Tower Battlements}
Accessed via hatches. Ballista, catapults, or cannons were often deployed here.

\subsubsection{Champion's Gate}
This gate was struck by raining debris from the meteor and has mostly collapsed. One of the towers has toppled over into several buildings below. The gate itself is splintered and shattered. The entire passage is strewn with large rocks, chunks of worked stone, splintered wood and other scraps, debris, and refuse. There is a thin walkway cleared through the rubble and one of the towers has been cleared up to the entrance. Just inside the gate are a few tents gathered around a small makeshift campsite. The camp is huddled near the one available entrance into the gate tower. Twenty or so pilgrims of the Falling Fire can be found here resting and recuperating before going further on their journey. They occupy the small camp and some use the tower for discussions, meetings, or for the wounded who need support and healing. The gate is watched over by a Falling Fire \textbf{chaplain} named Berman Weber.


\begin{itemize}
\item The gate is almost entirely collapsed after the meteor strike, and passing through is \textit{difficult terrain}.
\item The area around this gate is heavily contaminated and covered by the Deep Haze. Characters passing through the gate must succeed on a DC 15 constitution saving throw or suffer 10 (3d6) necrotic damage and gain one level of Contamination.
\item All the followers of the Falling Fire based here have taken the sacrament and are immune to the Haze.
\item When new followers pass through on their pilgrimage they are offered quick support and healing and sent on their way to the craters edge.
\item The first time that the characters pass through this gate, \textbf{Lucretia Mathias} is here. She foresaw their coming and talks to them about joining her. On any given occasion afterwards there is a fifty percent chance that Lucretia herself is here to greet the new followers who have come to take the sacrament and knowingly awaiting the characters to implore them to join the Falling Fire.
\item If the Falling Fire become enemies of the characters, Lucretia Mathias is no longer here.
\item The gate is guarded by 5 cult fanatics and 2 \textbf{zealots} who protect the flock. The rest of the Falling Fire members camping here are cultists. If the Falling Fire are enemies, there are twice as many of the above troops, and Lucretia Mathias summons two devas who perch on either tower and guard the gate.
\item Occasionally Nathaniel Flint can be found here escorting a group of 6d6 commoners.
\end{itemize}

\subsubsection{Interactions}
The Falling Fire are welcoming but wary of wanderers who enter their camp. The fanatics and zealots question anyone approaching about their business. If they seem to have good intentions, or are just passing through, they ask if any would like to speak to Berman Weber, and offer them seats by the fire to rest for "the long road ahead".


\begin{itemize}
\item Around the campfire everyone shares stories of their call to the Falling Fire, they speak highly of Lucretia Mathias and the coming age of heroes. Many seem excited for their pilgrimage and openly discuss and debate the Falling Fire with the characters.
\item Berman Weber is willing to speak to characters and implores them to follow Lucretia Mathias. He discusses fate and the chance of their meeting, and uses that to push the player characters to take the sacrament. He believes and is very open in expressing that the age of heroes is coming, and no doubt that the characters are indeed, destined to come with them. If the characters refuse Berman expresses that this is fine, and he will see them again when the time is right.
\item If any hostility breaks out, Berman and the guards attempt to defend their flock. The cultists will defend themselves as well but they flee if the battle seems to not be going in their favour. If the Falling Fire abandons the gate, they send more people to secure it again 1d4 days later.
\end{itemize}

\subsubsection{College Gate}
College Gate lays wide open. The streets before it are littered with ruins mostly burned to ash. Nothing but the foundations of buildings can be found here. The gate itself is covered in strange shimmering spiderwebs. They cover the entire gate, giving it the look of a broken pane of glass with webbed cracks weaving throughout it. Approaching the gate reveals a strange anomaly. Through the webs one can catch glimpses of alternate timelines and realities - one might see people walking the streets beyond, or the bustle of a living city; others see horrible demonic things enslaving humans; still others see a lush overgrown Drakkenheim covered in fungus and plant life, while other glimpses see a frozen or burning wasteland. Through each of the frames of the shattered glass one can see all conceivable futures and pasts for the city.

To look upon it is both mesmerizing and madness-inducing, as the fractured realities begin to warp one's mind.


\begin{itemize}
\item Any creature who looks upon the gate from within sixty feet of it must succeed on a DC15 Intelligence saving throw or take 2d12 psychic damage, are incapacitated, and suffer an immediate Drakkenheim Madness. They remain incapacitated until someone snaps them out of it, or they take damage.
\item Touching the webs causes 7 (2d6) phase spiders to emerge, hungry to consume any creatures caught in their shifting realities. They will attack characters who are going mad first.
\item Burning away the webs removes them entirely, but they will reappear in 1d6 hours.
\end{itemize}

\subsubsection{King's Gate}
This gate is occupied by a dozen trolls. Three trolls wait outside the gate to collect the toll. Two stay inside to work the winch, and two stand atop each tower to use the ballista. Three more wander the walls and surrounding area. The Troll King and his loyal captain are inside the left tower on the second level where they hold court over this new kingdom the Troll King has claimed.


\begin{itemize}
\item The Troll King, Zaffod, a huge \textbf{troll} with two heads and three arms, sits in his makeshift throne room on the second level of the north tower. One of his heads goes by "Zaf" the other goes by "Fod," and they are constantly arguing over who gets to wear the crown. He has a hard time fitting through the doors so he rarely leaves the room. He has a large belly that hangs over his knees, and wears a necklace of human skulls. Those who do not pay the toll, or are otherwise problematic for the trolls, are brought before the king where he decides their fate. Their fate is almost always to be "boiled in a large pot and fed to the troops."
\item Zaffod wears a circlet crown that is a \textit{headband of intellect}.
\item The Trolls of the gate are somewhat intellectual and although gruff, still present themselves as a jovial and agreeable group.
\item The three trolls at the gate demand a Troll Toll from anyone who hopes to enter Drakkenheim. It's reasonable—it will only cost you an arm and a leg! The trolls first bargain that others may pass as long as they give up one of them to be eaten. They beg for whoever is the biggest and meatiest of the group. They can be bargained with and are relatively easily persuaded, however they are indignant if offered livestock or animals, they respond "You are what you eat! Do you think us animals? We are civilized folk and so only dine on civilized folk!"
\item If characters are reduced to 0 hit points they are knocked unconscious by the trolls, brought before the king tied up and bound, and presented as a feast.
\item They have a makeshift prison on the lowest level of the gatehouse where they keep prisoners that are next up to be cooked and eaten. They take the characters one at a time to a massive cauldron they keep over a burning fire and prepare to cook and eat them.
\item The trolls occupying the gate must squeeze to use the stairs and doors, so they rarely use the interior spaces except to sleep.
\end{itemize}

\DndQuote%
{}
{''The easiest way into the city might be going over the walls. What's the worst that could happen?''}
{Sebastian Crowe}
\subsubsection{Shepherd's Gate}
The Hooded Lanterns control Shepherd's Gate. A dozen veterans and thirty scouts are stationed here at all times. During the day, this force is commanded by Gate-Captain Jacob Slovak, and by Gate-Captain Raine Highlash at night (both human urban rangers).

The entire force is rotated on twelve hour shifts so the defenders may rest outside the Haze. The guard force changes over each day at 8 AM and 8 PM. They maintain a nearby shed where they keep six fresh warhorses. This is a brutal assignment, so reserve forces are often called up to take a shift to give the main guards extra time off to recover.


\begin{description}
\item[Siege Defenses.]
The rangers have installed additional defenses all around the gatehouse. Wooden hoardings are built over the walls and tower battlements, and they've dug ditches lined with wooden stakes below the walls. Finally, a wooden barricade supported by sandbags has been erected on the road. Several ballistas are deployed atop each tower.
\end{description}

\subsubsection{Interactions}
They have runners and several horses at the ready, as well as flare guns and signal horns.


\begin{itemize}
\item They offer to escort adventurers to the City Guard Barracks, which is only a few hundred feet away in the Midden.
\item There is a ten percent chance that Elias Drexel can be found here when the characters are passing through.
\item The Hooded Lanterns double the amount of troops encountered here if they've made enemies with the player characters or another faction.
\end{itemize}

\subsubsection{Temple Gate}
Occupied by garmyr who have built a shanty-town encampment around the gate. Nearly a hundred garmyr dwell here at any given time, and the warbands are like a hornet's nest.


\begin{itemize}
\item 70 (20d6) garmyr dwell here with 14 (4d6) worgs
\item 10 (3d6) garmyr berserkers and 5 (2d4) hell hounds
\end{itemize}

\subsection{Battle of Temple Gate}
At a pivotal moment in the campaign, the Silver Order announces their plans to take Temple Gate, forcing the garmyr back into the city and securing a position for their force to gain a foothold in the city. Regardless of the characters' allegiances, this battle proves to be a major event for all factions, and is the first act in a series of events that lead to a confrontation at the Cathedral of Saint Vitruvio.

The Battle of Temple Gate might occur after the player characters have achieved the following:


\begin{itemize}
\item Met the faction leaders and started forming alliances with one or more of them
\item Explored the Cosmological Clocktower and recovered one of the \textit{Seal of Drakkenheim} there, plus visited at least two other locations in the Inner City
\end{itemize}

Following these developments, Knight-Captain Theodore Marshal reasons it's time to make a push for Temple Gate. Depending on who the characters are most closely allied with, this could result in many different outcomes. Whatever faction the characters are closest with will inform them as described below.

\subparagraph{Silver Order}
If the player characters are allied with the Silver Order, Theadore Marshal summons the characters to Camp Dawn to devise a battle plan together. He proposes the player characters spearhead an assault to open the gates for the cavalry, while the Silver Order forces hold off any garmyr reinforcements. The Silver Order aims to take the gate, but want to avoid damaging or destroying it in the process so it can be used as a forward base. However, he's worried about sabotage from the other factions, and asks the characters to act as envoys to either secure aid from the Hooded Lanterns, or check in on the other factions to ensure they won't interfere.

\subparagraph{Hooded Lanterns}
The Hooded Lanterns have tried on multiple occasions to push back the garmyr, but the Lord of the Feast has proven a formidable foe. Lord Commander Elias Drexel is open to working with the Silver Order on a joint operation, and proposes a pincer strike with the Hooded Lanterns working inside the city, and the Silver Order engaging in a frontal assault. However, the Lord Commander is concerned about a foreign military force occupying ground in Drakkenheim, and wants assurance from the Silver Order that they aren't here for conquest.

\subparagraph{Queen's Men}
The garmyr are a problem for the Queen's Men, too, and have caused more than enough injuries and death amongst her gangs. While she has ulterior motives of her own, the Queen of Thieves wants to see the Silver Order accomplish their goals. She believes that holding Temple Gate will likely stretch the Silver Order to their limit once they raise the ire of the Lord of the Feast, and hopes that a victory in this battle will lead to the embarrassing defeat of the Silver Order from attrition.

\subparagraph{Amethyst Academy}
Eldrick Runeweaver is curious but concerned about the Silver Order's plans, and wants to ascertain their goals. He knows the Silver Order are keen to destroy delerium, which is a problem for the Academy. However, he's doubtful destroying all the delerium in the city is even possible, so he wants to know if the Silver Order has an ace-in-the-hole. Taking Temple Gate and destroying the monsters there is in the best interest of everyone, but the Academy doesn't want the Silver Order getting too comfortable with their place in the city.

\subparagraph{Followers of the Falling Fire}
Lucretia Mathias is concerned about the Silver Order's plan, and doesn't want them to gain a footing so close to her flock. She foresees the Silver Order will use this momentum to force a war on the Falling Fire and attempt to drive them from the city as if they were no better than the garmyr. Lucretia wants to find a way to prevent the Silver Order from taking the gate, and instead have the characters lure the Lord of The Feast out of the cathedral and into open war with the Silver Order. This may give them time to claim the cathedral for themselves while the two forces whittle each other down.

\DndQuote%
{}
{''Many adventurers think they can take the 'easy' way into the city by going over the walls...Funny enough I never see them come back that way...or come back at all.''}
{Veo Sjena}
\includegraphics[width=0.4\textwidth]{images/../5e.tools/img/adventure/DoDk/078-06-002.temple-gate.png}
\subsubsection{War At The Gate}
The characters can participate in the battle of Temple Gate in any way they choose, regardless of their alliances. The Silver Order offers griffon riders, cavalry, and a band of knights to aid in any way. The players can be creative with how they take on Temple Gate. Regardless, it will be a difficult encounter.

If the player characters have made enemies of the Silver Order at this point, there are also options to convince any of the other factions to help ensure a failure at the battle. There are many options for crafty players to swing the battle the way they want it to go.

\includegraphics[width=0.4\textwidth]{images/../5e.tools/img/adventure/DoDk/079-06-003.crossbow.png}
The force at the gate consists of 100 garmyr and their \textbf{chimera}. They have sentries on the walls armed with crossbows, and two ballistas atop each tower always manned by several garmyr. They keep the portcullis on the gate closed and keep all doors and rooms well-guarded. Even with the battle raging, characters will have to fight their way through the garmyr to open the portcullis and ensure the Silver Order's forces can take the gate.


\begin{itemize}
\item The Silver Order tasks the players with getting the gate open for them.
\item They offer to fly them in on griffons, or put them on horses for the first charge to try and get them as close to the gate as possible.
\item There are two winches inside the towers that need to be opened for the gate to be usable. The players must open both.
\item Once the gate is open, the Silver Order are able to drive out the garmyr in 1d6 rounds as their troops flood the gatehouses and clear out any monsters who do not flee.
\end{itemize}

\begin{DndSidebar}[float=!b]{Developments}
After the events at Temple Gate, there are several directions the factions might push to go.



\subsection{If The Silver Order Succeeds}
If the gate is taken, the Silver Order are able to hold it, but can not push any further without the aid of the characters. They use the gate as a forward command post and their presence draws the eyes of the other factions towards their goals and objectives. The next move for the Silver Order is to try to convince the characters to help them defeat the Lord of the Feast and take the cathedral. Without the characters' help this task is impossible.

Lucretia Mathias needs to make it to the cathedral before the Silver Order does, and pushes the matter on the characters if they are allies. Meanwhile the Queen of Thieves wants the characters to take the cathedral and obtain many items that she would like for herself. Depending on the relationship between the factions at this point, it is up to the characters to convince the Hooded Lanterns or Amethyst Academy if continuing to help the Silver Order is of any value to them.



\subsection{If The Silver Order Fails}
Perhaps the events did not go well, or perhaps the player characters found themselves opposed to the Silver Order's goals. If the Silver Order fails to take the gate, the Lord of the Feast claims the victory, putting many Silver Order knights on a gruesome display hanging over the city gate. The Silver Order maintain their position in Camp Dawn, but unless aided by the characters, fail to gain a footing in the city. If no help is offered, the Silver Order call for reinforcements from Elyria to help them take the city by force and begin to plan another attack that will take months to come to fruition.


\begin{itemize}
\item If the characters have entered the Inner City to explore it, the Silver Order will announce their plans to take Temple Gate.
\item If the player characters do not intervene, whether or not the Silver Order takes Temple Gate is up to you - we recommend whatever option creates the greatest complication for the Characters.
\item Unless the player characters intervene, the Hooded Lanterns and the Silver Order join forces to take Temple Gate, however lacking the collaboration of the characters, this alliance disintegrates soon after, leaving the Silver Order unable to push further although they maintain control of the gate.
\item If the characters have not interacted much with either the Hooded Lanterns or the Silver Order, but are already exploring the Inner City, choose either the Hooded Lanterns or the Silver Order and have them ask the characters if they will help them form an alliance for the taking of the gate. Whether or not this works is up for the characters to decide.
\end{itemize}



\end{DndSidebar}

\section{Cosmological Clocktower}
The Cosmological Clocktower is built on the Market Square Plaza at the heart of Drakkenheim. Created at the behest of an eccentric monarch, this mechanical wonder mysteriously stopped one hundred eleven years ago, and its function as a public belltower was replaced by Saint Vitruvio's Cathedral. Many early expedition teams into the ruins used the clocktower as a landmark and bivouac; remarkably the tower attic lies just above the Haze. Yet the highest spires of Drakkenheim are no safer than the streets below, especially since a bloodthirsty flock of harpies took to roost in the tower.

\includegraphics[width=0.4\textwidth]{images/../5e.tools/img/adventure/DoDk/080-06-004.cosmological-clocktower.png}
\subsection{Overview}
In this adventure, the characters assault the Cosmological Clocktower to wrest control of the strategic outpost. Not only will they be rewarded for their efforts with their own base of operations in the city, they may also find a \textit{Seal of Drakkenheim} here.

\subsubsection{Adventure Hook}
\subparagraph{Timely Rest}
A long-standing rumour holds that the Cosmological Clocktower is one of the few "safe" places to rest within Drakkenheim. These rumours are indeed true: the tower attic juts out just above the Haze, thus the area is a safe place for characters to take a long rest. Every faction leader knows the value of the clocktower as a base of operations within the city. Once the characters begin developing a positive and collaborative relationship with any faction leader, they'll suggest the characters claim it from the vicious harpies who dwell there.

\subsection{Interactions}
It's unknown if these harpies were the result of strange mutations wrought upon birds or humans, perhaps both are true. Nevertheless these capricious, cruel, and territorial monsters have little interest in bargaining or negotiations. Nearly 40 harpies live in the tower, though at any given time, half are out hunting. They are led by the \textbf{Crimson Countess}, a red-plumed harpy matriarch:


\begin{description}
\item[Personality Trait.]
I screech blood-curdling war cries and perform aerial stunts as intimidation.
\item[Ideal.]
Any who walk the earth are prey for those who soar the sky.
\item[Bond.]
I collect trophies from my kills to remind my flock who is the mightiest.
\item[Flaw.]
I am aggressively territorial, envious of my rivals, and covetous of my trinkets.
\end{description}

\DndQuote%
{}
{''Birds constantly circle this tower, but recently the birds seem much, much bigger.''}
{Sebastian Crowe}
\subsection{Area Details}
\includegraphics[width=0.4\textwidth]{images/../5e.tools/img/adventure/DoDk/081-map-6.02-cosmological-clocktower.png}
\subsubsection{Market Square Plaza}
The market square is the largest open courtyard in the city.

The square consists of two plazas, separated by a cluster of banks, vaults, guildhalls, and businesses. The ruins and remains of high town houses with storefronts surround the square. Scattered across the plaza are the remains of Drakkenheim's central open air market: rows of stalls, canvas tents, pavilions, baskets, crates, wagons, barrels, wheelbarrows, carts, and bones creating a tattered display of the chaos of that day left behind. The distant sounds of a lullaby, chimes, and faraway bells fills the air.

Market Square Plaza bears the scars of many battles: makeshift barricades, scorch marks, broken weapons, armour scraps, and a few scattered corpses. Many rats scurry about, crows pick over the remains, and a thin, silvery-purple mist hangs in the air. Looking up, you can see the Cosmological Clocktower piercing through the fog.


\begin{itemize}
\item Lingering in the market square is dangerous. Characters who do not take steps to conceal their passage are noticed and attacked by 10 (3d6) harpies who fly out from the tower.
\end{itemize}


\begin{description}
\item[Drakkenheim Guildhall.]
This building is not far from the clocktower. Formerly the administrative hub for trade and industry in Drakkenheim, this larger open hall is a gathering place for the city guild leaders and merchants, and contains offices, vaults, and a large auditorium that served as council chambers for the guild. However, the once-great vaulted roof has entirely collapsed, leaving mostly only the outer shell of the great building behind. Perhaps several impressive vaults lay buried under heaps of wreckage and rubble, but it would take significant effort to excavate the damage.
\end{description}

\subsubsection{Clocktower Facade}
At the centre of the square is the clocktower. The three-hundred-foot tall square tower is about fifty feet wide at the base, with rounded turret corners and walls supported by great buttresses. The front facade is a mechanical wonder. Near the top of the tower is a massive thirty-foot-wide clock face made from glass, the numerals and hands filigreed with silver, brass, and gold. Above are several crenellated balconies, and the roof is a steep spire topped with a wrought-iron weathervane.

Bordering the central clock is an embossed metal and glass dial that depicts the transit of the sun and the phases of the moon through each month. From this dial, one can read not only the time but the date, as it shows a stylized interpretation of the seasons and constellations. Extending down the clock face is a further system of circular glass plates, gears, and wheels that illustrate in mechanical precision the grand cosmic dance of the planets and stars as they hurtle through the astral void. Multicoloured glass lenses superimposed over the planet depict the metaphysical orbit of the afterworlds and side-realms and their relation to the mortal world. The rest of the tower is decorated with elegant carved statues depicting angels, saints, demons, devils, elementals, and the various creatures of the worlds beyond. On some levels are rows of gothic windows with trefoil arches filled with stained-glass mosaics.

The entire mechanism is frozen, stuck on the 11th day of the 11th month, 111 years ago, at 11:11 AM, the phases of the planes resting between the Elemental Chaos and Dreamland.

On the ground level, a set of stone stairs leads to the copper double doors.

\subsubsection{Clocktower Interior}
\begin{DndReadAloud}{}
The tower interior is a dramatic contrast to the elegant outer facade. The walls are exposed brick, and a rickety wooden staircase ascends the interior supported by a wrought-iron scaffold. The decorative exterior windows allow some light to pass into the tower; many along the way are shattered and broken. The stone floor is covered in dust, bone, debris, and avian feces. Gears connected by chains, stacks of heavy dangling counterweights on pulleys, and pendulum mechanisms hang down the northern wall. High above can be seen the bells.

\end{DndReadAloud}

\subparagraph{Stairs}
These creaking wooden stairs snake along the inner walls of the tower for two hundred feet, connecting the ground floor to the belfry gantries above. Each individual flight rises sharply twenty feet to a small corner landing and sports a thin metal railing.


\begin{itemize}
\item The harpies make their nests in the landings along the staircase. They enter and exit the tower through the broken windows and never use the front doors. There may be as many as 10 (3d6) harpies roosting throughout the stairs at any time.
\end{itemize}

\includegraphics[width=0.4\textwidth]{images/../5e.tools/img/adventure/DoDk/082-06-005.egg.png}
\subparagraph{Belfry}
The staircase ends at a wooden gantry clinging to the outer walls of the tower. This narrow walkway overlooks thirteen massive bronze bells hung in the middle of the tower belfry. Each ranges from three to five feet wide and weighs several tonnes. The bells are yoked to steel girders and do not swing. Instead, the bells are struck by exterior hammers fitted to each and controlled by the machinery above. As each bell sounds in a distinct tone, this mechanical system could play many unique bellsongs when the clock tower still functioned.


\begin{itemize}
\item It's a two-hundred-foot drop to the floor below. A spiral staircase in the northwestern parapet ascents twenty feet to the Clockworks above.
\item 1d4 manticores, consorts to the \textbf{Crimson Countess}, make nests in niches in the belfry.
\end{itemize}

\subsubsection{Clockworks}
\begin{DndReadAloud}{}
A complex mechanism of gears, counterweights, and clockwork dominates this chamber. Several thin windows overlook the surrounding cityscape. The floor features several steel grates and wooden hatches that allow access to the machinery and bells below. A strange granite obelisk is set in place of the main clock pendulum; the core of the "ticking" mechanism that controls the clock. Three empty sockets lined with delicate copper tracery are cut into the stone. A fine elven script inlaid with thin lines of silver and gold are etched into the surrounding surface. Several pennies are precisely balanced atop it. A system of pulleys, winches, and levers connects it to the rest of the clockwork.

A set of shelves here contain several notched metal discs, each kept in a paper folder bearing the title of a song. When the clock tower functioned, these discs could be loaded into the receptacle to determine which song was played by the bells.

Several small and curious mechanical automatons wander listlessly around this chamber, examining and tinkering with the mechanical components. They are shaped like platonic polyhedrons with spindly mechanical limbs, disturbingly realistic glass eyes, and a mechanical mouthpiece.

\end{DndReadAloud}


\begin{itemize}
\item There are 1d6 automatons here. If their game statistics are required, use those for \textbf{animated armor}, except their size is small. If interacted with, they seem addled and confused. When they speak, they repeatedly utter "zero" and "one" in long-winded combinations. The remains of a larger automaton are splayed about the chamber, it appears like it exploded from the inside.
\item The elven runes describe in precise and exacting detail several mathematical calculations involving the passage of time, the planets, the seasons, and the planes.
\item The stone itself is quite lightweight; a character proficient in Arcana or engineering-related fields could deduce that such a stone would normally not be sufficient to accurately drive such an intricate clock mechanism.
\item Examining the stone further with \textit{detect magic} or \textit{identify} spells determines it is a magic item. It is attuned to time and space itself, and could vary its density and composition to act as the perfect pendulum. However, it is missing an appropriate power source.
\end{itemize}

The harpies rarely come to this level. They don't use the stairs, and fly directly into the Rookery or through the broken windows in the Bell Tower.

\subsubsection{Rookery}
\begin{DndReadAloud}{}
A disgusting amalgamation of bird nest and boudoir. Macabre decorations of bone, skulls, and body parts are presented alongside scavenged finery from the ruins: pillows, sheets, veils, and the like. All are filthy, torn, and ruined. Gibbet cages with skeletal remains of several prisoners hang from the balconies.

\end{DndReadAloud}


\begin{itemize}
\item This is the lair of the \textbf{Crimson Countess}.
\item Treasure. The Crimson Countess has amassed many trinkets and trophies, most of them skulls and bones. She wields a \textit{javelin of lightning}, and her favourite charm is one of the \textit{Seals of Drakkenheim}: the \textit{Steward's Seal}. It can be found in her nest alongside the following note:
\end{itemize}

\begin{DndReadAloud}{}
My child—Take this treasure of mine and run. Hide it in a safe place outside the castle, then flee the city. No matter what happens, know that I will always love you.—Papa

\end{DndReadAloud}

\subsubsection{Spire Attic}
\begin{DndReadAloud}{}
The former clockmaster's apartment has seen better days, but it seems that at some point adventurers may have made camp here.

\end{DndReadAloud}


\begin{itemize}
\item 1d6 uncommon \textit{potions} and a Keoghtom's Ointment were left behind here.
\item This area is above the Haze, and characters may safely take a long rest here.
\item Searching through the bookshelves and documents uncovers a letter to the clockmaster (see below)
\end{itemize}

\includegraphics[width=0.4\textwidth]{images/../5e.tools/img/adventure/DoDk/083-06-006.harpies.png}
\begin{DndSidebar}[float=!b]{Developments}
Not all the harpies who dwell in the clocktower may be present during the characters' invasion. Characters who clear out the clocktower and wish to use it as a base of operations may find themselves dealing with counter-attacks from the former occupants for several weeks afterwards, or even from the Crimson Countess' former rivals.



Taking the clocktower is a major campaign milestone, and the other factions are jealous of the accomplishment. A rival faction might attempt to find the characters here or even attempt to wrest control of the tower. While the characters can rely on the tower as a defensible place to rest, once the factions learn the characters use this place as a respite, they might dispatch strike teams to occupy the tower while the characters are absent. The characters would do well to make a good alliance to help them keep this holding secure, as if left abandoned other monsters and faction forces are quick to seize the opportunity to control this valuable strategic location. Nevertheless, once the characters have taken adequate steps to firmly defend their new holding, it is a critical base of operations for their future explorations.



\subsection{The Steward's Seal}
If any of the factions learn the characters have found one of the \textit{Seals of Drakkenheim}, it's a major development. Locating this missing item ignites hope amongst the Hooded Lanterns of a realm restored, and attracts the intense interest of the Queen of Thieves. Any faction lieutenant can recognize the seal on sight, and if they see characters wearing it, they'll make overtures to speak to their faction leader, who can explain the significance and powers of the item.



\subsection{Letter to the Clockmaster}
\begin{DndReadAloud}{Archmage Adriana Modiera, Amethyst Academy}
Thank you for writing to me. I know you and your family have toiled hard for generations to keep the clocktower, despite the main mechanism failing a century ago. I must admit, I long believed the clocktower was a frivolous bit of architecture, based on a long-debunked cosmological theory. Nevertheless, a recent discovery upon the Isle of Skye has ignited my curiosity.

The original machinery dates beyond antiquity, and many key components appear similar to devices found in several ancient elvish sites in Skye, though I cannot speak to their original purpose. Examining the clocktower mechanism indicates several delicate parts are heavily damaged and others are now missing, including the power source. We are confident components recovered from Skye could be used to repair it.

These components are due to arrive at the Inscrutable Tower on September 13th. Please make arrangements for our mages and artificers to begin work on repairs the following week. Unfortunately, you will need to vacate the premises for the foreseeable future. We will draw up arrangements for the Amethyst Academy to lease the building.

Best Regards,

\end{DndReadAloud}


\begin{itemize}
\item September 13th is three days before the meteor struck Drakkenheim.
\item The missing power source was in fact a delerium crystal that existed on earth long before the meteor fell. Such samples are extraordinarily rare, but they do exist. The parts can be found in the Director's Office of the Inscrutable Tower.
\item What happens if characters secure the parts and fix the clock tower is up to you. It could reveal a secret about the meteor, or unlock an important clue to a personal quest.
\end{itemize}



\end{DndSidebar}

\section{The Crater}
\DndQuote%
{}
{''No one has ever been there and returned. We don't really know what we're up against, do we?''}
{Lieutenant Petra Lang}
The meteor crashed upon the South Ward of Drakkenheim. Once the industrial hub of the city, little remains here save blasted ruins and a massive mist-filled crater. The edge of the crater is a long-sought destination for prospectors and pilgrims alike for the vast deposits of delerium found there. Few have ventured deeper within, however. Shrouded deep within the most dangerous and contaminated region in all of Drakkenheim, madness and ruin fester in the pulsing Delerium Heart.

\includegraphics[width=0.4\textwidth]{images/../5e.tools/img/adventure/DoDk/084-06-007.the-crater.png}
\subsection{Overview}
Depending on their objectives, player characters planning an expedition to the crater might only aim to reach the Crater's Edge, or they might push onwards through the Crater Basin to the centre of the crater. Either way, they'll need potent tools to protect themselves against contamination.

\subsubsection{Adventure Hooks}
\subparagraph{Survey the Impact Site}
(Amethyst Academy, Silver Order, or Hooded Lanterns). If the characters discover an extraordinarily effective defense against contamination, such as the \textit{neutralizing field} spell or \textit{hazewalker plate}, a faction leader may suggest using these means to launch a mission to the crater. The goal is simple: find out what's in there. Faction leaders can offer the following warnings:


\begin{itemize}
\item Mages have tried to study the crater using \textit{scrying} spells and divination magic, but were driven mad in the process. Those who tried teleporting within were never seen again.
\item Eldritch lightning surrounds the air above the crater, and creatures struck by it are disintegrated.
\item Lucretia Mathias made the journey herself, but she does not speak of what she discovered there. She cryptically remarks that it cannot be explained or understood, and must be seen.
\end{itemize}

\subparagraph{Pilgrimage of Falling Fire}
(Falling Fire) Characters who demonstrate their faith in the teachings of Lucretia Mathias are invited to take the pilgrimage to the Crater's Edge. There, each must find a delerium fragment in the shape of their soul, and bring it to the Chapel of Saint Gresha, where the pilgrim may partake in the Sacrament of Falling Fire.

\subsection{Interactions}
\begin{DndReadAloud}{}
"Set a fire in your belly, warm your heart, and shield yourself in light for the road ahead."

\end{DndReadAloud}

\textbf{Lucretia Mathias}, Nathaniel Flint, and other Followers of the Falling Fire often lead pilgrimages of cultists and commoners to the Crater's Edge, and can frequently be encountered in this region. Any random encounters in this region may be replaced with a Falling Fire strike team, or a pilgrimage of 10 (3d6) cultists led by a \textbf{chaplain}.

\subsubsection{Pilgrimage of Falling Fire}
Only characters who are devout and sincere in their commitment to the Followers of the Falling Fire are invited to complete a pilgrimage to the Crater's Edge, where they will participate in a ritual known as the \textit{Sacrament of Falling Fire} at the Chapel of Saint Gresha. Characters can participate on their own or as part of another group of Falling Fire pilgrims. While outsiders rarely witness the rite itself, the Followers of the Falling Fire will tolerate a respectful observer in the hopes that seeing the ritual might assuage any remaining doubt in their heart.

\subsubsection{A Fragment in the Shape of a Soul}
When they arrive at the Crater's Edge, pilgrims must individually search for a delerium fragment that reflects an inner aspect of their own heart. Each attempt takes one hour, after which the character makes an ability check using their highest ability score. They find a suitable fragment on a check result of 20 or higher, and should appropriately describe and roleplay how this check guides them towards a delerium fragment. Most normal pilgrims must visit the crater several times before they successfully find such a fragment.

\subsubsection{Completing the Pilgrimage}
Once a pilgrim finds a fragment, they must bring it to the Chapel of Saint Gresha (see below) where they take the \textit{Sacrament of the Falling Fire}.

\subsection{Area Details}
The entire South Ward is suffused in the Deep Haze (see chapter 5).

\subsubsection{South Ward Ruins}
\begin{DndReadAloud}{}
Few structures are standing in the South Ward of Drakkenheim. Little remains save demolished foundations, massive chunks of blasted rock, and piles of rubble and debris. The thick fog of the Deep Haze blankets the wasteland. The direction of the devastation is unmistakable, however. Stone, brick, and timbers on any walls or pillars facing the crater itself appear melted rather than burned, and in some cases, pieces of rock or wood have partially turned into glass. Occasionally, you see the unmistakable glimmer of delerium fused into the ground or walls of a tumbledown building. The air here is cold, the mists billow about in hypnotic swirling patterns. A peculiar smell fills the air, it's acrid and pungent as if some strange material were burning nearby. Once in a while you hear the sound of crumbling rock or the creaking collapse of a distant building.

\end{DndReadAloud}


\begin{itemize}
\item Check for random encounters as normal as characters travel through this region.
\end{itemize}

\DndQuote%
{}
{''I've had some horrible nightmares in my day. This place was unfathomably worse.''}
{Sebastian Crowe}
\subsubsection{Crater's Edge}
\begin{DndReadAloud}{}
You stand upon the crater's edge. Here, the blasted ruins of the city end upon a rocky precipice that descends into a mist-filled basin. A plume of thick swirling fog rises from the centre of the crater. Crackling eldritch lighting arcs through the air into the clouds above, and a bitterly cold wind howls about. Jutting out from the rocky chasm are countless sparkling delerium fragments.

\end{DndReadAloud}


\begin{itemize}
\item The cliff edge descends one hundred fifty feet to the crater basin floor on a sharp slope.
\item It's impossible to see the crater basin floor due to the thick mists.
\item DC 15 Arcana. A character who studies the fog and makes a check realizes the Haze here is significantly thicker and likely far more contaminated than normal. You can reveal to such characters that Constitution saving throws against contamination will likely be needed for every 10 minutes within the Haze inside the crater basin.
\item DC 15 Nature or Survival. Examining the swirling mist at the centre notes the powerful lighting arcing through the fog. You can reveal to characters the details of the \textit{eldritch lightning} in the Crater Basin (see below).
\end{itemize}

\subsubsection{Scaling the Cliff}
Three successful DC 15 Strength (Athletics) checks are needed to carefully scale up or down the edge. Characters may use equipment or spells to aid their climb - spells such as \textit{feather fall} or \textit{spider climb} allow automatic success.


\begin{itemize}
\item A character takes 10 (3d6) bludgeoning damage on a failed check as they tumble and slip down the sharp rock face.
\item A character who fails three checks before reaching three successes has a terrible fall. They take 17 (5d6) bludgeoning and 17 (5d6) piercing damage.
\end{itemize}

Alternatively, characters can Search the Ruins (see chapter 5) to locate gentler slopes or a pathway where easier descent is possible.

If you'd like to add extra tension to this moment, consider having a lone \textbf{grotesque gargant} climb out of the crater and attack while characters are looking out over the basin, scaling the cliff, or searching for a way down!

\subsubsection{Chapel of Saint Gresha}
\begin{DndReadAloud}{}
Along the northwest side of the crater rests the Chapel of Saint Gresha. Nothing remains of the domed rooftop. Instead, the bricks composing the circular chapel float in the air, deconstructed in perfect geometric order. Perhaps if the magic holding them aloft were broken, they would collapse back into their proper place. Inside the chapel walls, a lighted stone brazier stands upon a tiered marble dais. The entropic flame burns in an unnatural purple-red hue, the coals within are delerium shards. The mosaic on the floor is cracked and broken. Several platforms may have once held statues, but little else remains save blasted rock.

\end{DndReadAloud}


\begin{itemize}
\item \textbf{Lucretia Mathias} can be encountered here at dawn each day, when she performs the \textit{Sacrament of the Falling Fire} for up to 12 pilgrims. Nathaniel Flint is sometimes here leading a pilgrimage as well.
\item \textbf{Saint Gresha} (see below) performs the rite each day at dusk.
\item 12 neutral good wights guard the chapel. Each wears ancient half-plate armour. These are the sacred remains of former paladins interred in the chapel, and are not haze wights.
\end{itemize}

\subsubsection{Saint Gresha}
Gresha is a neutral good \textbf{arcane wraith}. She was called back in spirit by Lucretia Mathias to tend her namesake chapel. Her Spellcasting trait is replaced with the following:

\begin{DndReadAloud}{}
At will: \textit{guidance}, \textit{light}, \textit{sacred flame}, \textit{thaumaturgy}

3/day each: \textit{bless}, \textit{flame strike}, \textit{guiding bolt}, \textit{healing word}, \textit{prayer of healing}, \textit{spirit guardians}

1/day each: \textit{divine word}, \textit{harm}, \textit{heal}, \textit{holy aura}, \textit{Sacrament of the Falling Fire}*

*new spell described in Appendix D

\end{DndReadAloud}

If Saint Gresha is destroyed, Lucretia Mathias calls her spirit back again as soon as possible.

\subsubsection{Crater Basin}
\begin{DndReadAloud}{}
The crater basin is a forlorn and desolate place. The raw earth is a deep purple-brown colour. Black glass chips and rocky gravel crunch underfoot mixed with bits and pieces of unrecognizable masonry and flecks of delerium. Here and there you encounter peculiar vegetation: creeping moss that resembles exposed sinew or swelling tissue, with leaves like fleshy membranes and roots like spreading blood vessels. The strange plants cling to clusters of delerium and around pools of sludge.

\end{DndReadAloud}

\subparagraph{Deep Haze}
Within the crater, the Deep Haze is exceptionally thick and nearly impenetrable: vision beyond 30 feet is totally obscured and blocks most of the light from the sky above, so areas within are dimly lit and \textit{lightly obscured}. At the end of every ten minutes spent within, characters must succeed on a DC 20 Constitution saving throw or suffer 10 (3d6) necrotic damage and gain one level of contamination.

\subparagraph{Eldritch Lightning}
Sheets of flashing energy and crackling eldritch lightning erupt through the mist around the crater, which extends to the clouds above. Any airborne creature flying higher than sixty feet above the Crater Basin rolls 1d6 at the end of its turn: on a one, it is struck by the lightning and must succeed on a DC 20 Dexterity saving throw or take 10d6 lightning damage and 10d6 force damage. If this damage reduces the creature to 0 hit points, it and everything it is wearing or carrying are reduced to black ashes and scattered into the wind. The creature can be restored to life only by means of a \textit{true resurrection} or a \textit{wish} spell.

\subparagraph{Twisted Space}
A creature who teleports into or out of the Crater Basin is ripped apart by eldritch magic and takes damage as if they were struck by eldritch lightning (see above)

\subsubsection{Finding the Delerium Heart}
Characters seeking the Delerium Heart must Search the Ruins for one hour, as described in "Navigating Drakkenheim" in chapter 5. In addition to making Investigation checks, characters may use Arcana, Insight, Religion, Nature, or Survival to help navigate the basin. Success requires a DC 20 check result, with any check result of 25 or higher granting an additional success. What the characters find each attempt is based on their number of successes:


\begin{DndTable}{cX}

Number of Successes & Result\\
2 successes or fewer & Nothing\\
3 successes & Massive delerium cluster\\
4 or more successes & Delerium Heart\\
\end{DndTable}

Two or more failed check results triggers one of the random encounters below. As characters are exposed to a powerful version of the Deep Haze here, they will need some means to gain immunity to contamination, otherwise they'll have to make six Constitution saving throws each time they attempt searching the crater.

\subsection{Random Encounters in the Crater Basin}
Check for random encounters as normal as characters explore the basin, using this table to generate encounters.


\begin{DndTable}{cX}

1d6 & Random Encounter\\
1 & 1 \textbf{crater wurm}\\
2 & 1d4 grotesque gargants\\
3 & 4 (1d6) protean abominations and 10 (3d6) gibbering mouthers\\
4 & 35 (10d6) delerium dregs and 1d6 haze hulks\\
5 & 1d6 contaminated elementals\\
6 & 1d4 arcane wraiths and 1d6 warp witches\\
\end{DndTable}

\subsubsection{Massive Delerium Clusters}
\begin{DndReadAloud}{}
Unfathomably large delerium geode clusters protrude from the earth throughout the crater basin. Most grow as tall as a human, some as big as ogres or even larger. Viscous opalescent liquid perspires from the crystals, forming pools around it.

\end{DndReadAloud}


\begin{itemize}
\item There are a living haze, \textbf{animated delerium sludge}, walking delerium cluster, and an \textbf{entropic flame} lurking around the cluster and the pool.
\item The pools are two or three feet deep and anywhere from ten to forty feet in diameter.
\item Characters can find 1 delerium geode, 1d6 delerium crystals, and 2d6 delerium shards that are easy enough to extract. Harvesting the massive crystals in the field is impossible with portable equipment, however.
\item If damaged or struck these massive geodes release an arcane anomaly.
\end{itemize}

\subsubsection{Delerium Heart}
\begin{DndReadAloud}{}
The air is humid and filled with the overpowering smell of ozone here. Thrumming eldritch energy makes your hair stand on end, and a low humming sound painfully rings in your ears. Suddenly, the mists clear somewhat, revealing an unfathomably colossal delerium shard seventy feet high and thirty feet wide. Pieces of rock and meteoric iron cling to the sides; it glows brightly with blinding octarine light and visibly perspires globs of prismatic delerium sludge that forms a pool around it. The edges of the crystal intersect at incomprehensible angles, making its geometry shift despite its apparent stillness. It gives you a headache to look at it. Mounds of upturned earth, huge chunks of meteoric iron, and massive delerium geodes emerge like islands from the pond.

\end{DndReadAloud}


\begin{itemize}
\item This is all that remains of the core of the meteor. The mere sight of its impossible geometry is enough to drive a creature mad. The first time a humanoid creature witnesses the Delerium Heart, it must succeed on a DC 20 Intelligence saving throw. On a failed save, it is affected as if by a \textit{feeblemind} spell.
\item The delerium sludge pool is four feet deep (see Appendix D).
\item A creature who ends its turn within one hundred twenty feet of the Delerium Heart must succeed on a DC 20 Constitution saving throw or suffer 10 (3d6) necrotic damage and gain one level of contamination. A creature who fails this saving throw by five or more instead takes 35 (10d6) necrotic damage and gains 1d4 levels of contamination.
\item Treasure. A \textit{pristine delerium geode} is found here with 10 (3d6) delerium crystals. This extremely rare crystal formation is perfectly symmetrical, notably unique, and readily extracted by any character looking to collect delerium around the Heart.
\item Characters must find their way back from the Delerium Heart in the same manner they got here, by searching the ruins for the Crater's Edge.
\item A creature who approaches within ten feet of the Delerium Heart may examine it closely. It's hollow: suspended within a quicksilver fluid are embryonic fleshy masses connected with thin veins to translucent, pulsing heart-like organs. A creature who \textit{touches} the Delerium Heart has a sudden vision:
\end{itemize}

\begin{DndReadAloud}{}
The vision lasts only a moment, but feels like a thousand years. A falling star collides with an unfamiliar world. A prismatic mist slowly creeps from the impact. Massive sharp-edged crystals burst from the surface of the world, and the people there descend into chaos and madness. The corrupt sphere violently ruptures apart, sending a hail of comets through the cosmos until one looms over this world.

\end{DndReadAloud}

\subsubsection{Destroying the Delerium Heart}
Characters can attempt to destroy the Delerium Heart:


\begin{itemize}
\item It has AC 25 and 500 hit points. It regains 50 hit points each round on initiative count 20 as long as it has at least one hit point remaining.
\item The Heart is immune to cold, lightning, necrotic, piercing, poison, psychic, and slashing damage, as well as bludgeoning damage from non-magical attacks. It is resistant to fire damage. It is vulnerable to thunder damage.
\item Once per round on initiative count 20, if the Delerium Heart took damage in the previous round, it unleashes an arcane anomaly. The resounding noise attracts nearby creatures each round, attracting creatures using the random encounter table above.
\end{itemize}

When destroyed, the Delerium Heart explodes, dealing 10d6 necrotic damage, 10d6 radiant damage, 10d6 psychic damage, 10d6 piercing damage, and 10d6 thunder damage to all creatures within five hundred feet. This explosion causes the shard to fracture into thirteen massive geodes that are sent hurtling throughout the city, which each grow into a Delerium Heart within a year if they aren't destroyed. The embryonic mass within the Heart pours into the sludge, and quickly grows into a horrific abomination: a \textbf{Tarrasque}.

\DndQuote%
{}
{''Not as impressive as the Canyon of Caspia but many more monsters ripe for sport.''}
{Pluto Jackson}
\begin{DndSidebar}[float=!b]{Developments}
\DndQuote%
{}
{''This dire news only confirms what we already assumed to be true. Drakkenheim is no more.''}
{Archwizard Eldrick Runeweaver}


Characters who return expecting the factions to be united by their report will be disappointed.



\subparagraph{Hooded Lanterns}
Elias Drexel is horrified by the characters' report, which is a despair-inducing confirmation that restoring Drakkenheim will never be a simple matter, and a realization that he will need the help of another faction to accomplish it. The Lord Commander isn't even sure he should reveal the characters' report to the rank-and-file soldiers at large. He's right. Mutiny in the ranks erupts at the news as the soldiers argue about the merits of working with the Silver Order, the Followers of the Falling Fire, or the Amethyst Academy. If the characters have made an enemy of any of these factions, a few loudly and defiantly proclaim the characters have doomed them all. Unless the characters and Elias Drexel can deliver another major victory soon, the unrest festers and the Hooded Lanterns disintegrate into fractious warbands.

\subparagraph{Amethyst Academy}
Eldrick Runeweaver is fascinated by the news of the Delerium Heart, and openly remarks that unfathomable arcane power could be harnessed from it (and thereby prevent whatever calamity affected the alien world from happening here). After several days, he returns to the characters and asks them if they might be prepared to mount another expedition to set up a magical arcane relay he's designed. If the characters didn't bring back the pristine delerium geode, he also asks them to obtain it.

\subparagraph{Falling Fire}
Once characters return from the crater's heart, Lucretia Mathias can explain the meaning of any visions they beheld. She then reveals the purpose of the Sacrament: by storing enough faithful souls in delerium, she believes the substance itself will become an earthly vessel for the Sacred Flame, transforming the scourge of the universe into its savior. The sanctified delerium will spread light through the cosmos. When asked about the fate of this world, Lucretia Mathias remarks that the Sacred Flame has always required mortals make sacrifices for the cause of righteousness. Whether or not her claims are true cannot be determined until the events are allowed to occur.

\subparagraph{Silver Order}
Theodore Marshal sees no further argument at this point: Drakkenheim must be burned to the ground. The fate of the world rests upon it. Doing so isn't going to be simple, and now the mission is clear: the Silver Order must resurrect Argonath to use his dragonfire to destroy the Delerium Heart.

\subparagraph{Queen's Men}
The Queen of Thieves can only react with laughter to this knowledge. The fact that such a horrific threat lies in the heart of the city is music to her ears. As long as it's there, no one will ever truly be able to control Drakkenheim - except her. The Queen of Thieves aims to weaponize this knowledge to sow doubt and discord in the other factions, especially the Hooded Lanterns. Whatever the characters saw in those visions isn't going to happen anytime soon, anyways.



\end{DndSidebar}

\section{Kleinberg Estate}
Deep in the noble district of Drakkenheim, just south of Castle Drakken, a walled off community known as the Kleinberg Estate stands. Amongst the ruins and debris, several of the large mansions have survived the devastation of Drakkenheim. One such estate has become lively with activity in recent years. Figures have been seen coming and going from it, which has spurred whispers and rumours within the taverns of Emberwood village. Conversations between forlorn adventurers and sickly survivors speak of a mage working out of the estate who has means to help those who are suffering with contamination. The validity of these claims is yet to be proven, and so those seeking help best be wary, as more often than not, people who have ventured deeper in search of a cure seldom return.

\includegraphics[width=0.4\textwidth]{images/../5e.tools/img/adventure/DoDk/085-06-008.kleinberg-estate.png}
\subsection{Overview}
There may be many reasons the player characters choose to uncover the Pale Man.

\subsubsection{Adventure Hooks}
\subparagraph{Search for a Cure}
Characters seeking the means to reverse a monstrous transformation might seek out the Pale Man, as rumours abound of the secretive mage who can help with mutations and contamination.

\subparagraph{The Academy Deal}
The Academy might hire the party to find the Pale Man and uncover what secret spells and arcane concoctions he has created. River asks the players to retrieve his work and bring it back to them with any relevant information or samples. The Academy is willing to offer the Pale Man a deal for his cooperation, but fully expect him to decline. They do not care if the Pale Man lives, as long as his work is recovered.

\subparagraph{The Witch Hunt}
The Silver Order might hire the players to find and kill the Pale Man - they want his work burned, and any delerium in his possession destroyed.

\subsection{Interactions}
The minions that work for the Pale Man are more social than normal dregs, and should be portrayed as having quirky, weird, and insane personalities. They direct characters to Raven, who handles all business deals. None of them know much about the Pale Man's research or experiments, and just express relief at the aid he has offered them.

\subsubsection{Friedrich von Lichten, the Pale Man}
Friedrich Von Lichten ambitiously explored the mutagenic effects of delerium after Drakkenheim's destruction, and was among the first to fully describe the effects of eldritch contamination upon the human body. He was ejected from the Amethyst Academy when it was discovered he was using apprentices as test subjects. So grave was his trespass, that the Academy notified the Silver Order of his existence, putting a death bounty on his head. Both The Silver Order and the Amethyst Academy have been searching for him since.


\begin{description}
\item[Personality Trait.]
Witness my power, understand that I am the future. Some may call me a god, and I would not correct them.
\item[Ideal.]
We must embrace the Haze, and allow contamination to change us, make us better.
\item[Bonds.]
The Amethyst Academy lacks vision. I will show them the true power of delerium.
\item[Flaw.]
It is not enough that I merely outsmart my foes, they must be shown their weakness and failure.
\end{description}

Friedrich von Lichten has fully embraced contamination and transformed into a monster, and is now the \textbf{Pale Man}. He has maintained his intellect, and drapes his hunched and withered figure in a black cloak. Beneath his paper-thin skin are wriggling and writhing things. His lidless eyes are sunken, his teeth are prominent, yellowing and decaying. One can just make out a sickly green tongue that darts out when he speaks. He talks in a stunted, almost exasperated tone.

He presents his estate as a safehouse for the lost and forgotten.

The Pale Man sows the rumours that he can help contaminated people so he can use them as test subjects for his experiments. When he meets a transformed or contaminated creature, he presents himself as a helpful expert, offering to treat any individual left in his care. He never actually cures contamination, but sometimes helps other transformed monsters become somewhat lucid so he can make better use of them as minions. The Pale Man spends most of his time in the repurposed master bedroom of the estate he uses as his laboratory, and gives orders only to his most trusted thralls.

\includegraphics[width=0.4\textwidth]{images/../5e.tools/img/adventure/DoDk/086-06-009.henry.png}
\subsubsection{Mutated Thralls}
These three haze hulks are the most intelligent and loyal of the Pale Man's servants. They came to him willingly and embraced his alterations. The Pale Man allowed them to keep their lives and minds, and in exchange they serve him.

\textbf{Raven} a gutbuster haze hulk, is a portly woman with blue skin, white glowing eyes, and thick stringy black hair wearing a burlap cloak with a hood. She is extremely cautious of visitors, but is often the first to greet and question people coming to the estate.

\subparagraph{Henry}
Henry, a hunter haze hulk, is the most intelligent of the three. He wears black scraps of leather and has a long sallow face and large black eyes.

\subparagraph{Grunt}
Grunt, a juggernaut haze hulk, is the biggest of the three, with beady eyes on his tiny bald head and a massive hulking body. He is usually found in the estate house or the stables working as the personal guard for the Pale Man.

\subparagraph{Lisa}
Lisa is the Pale Man's favourite creation, transformed by contamination into a \textbf{grotesque gargant}. He used her as the conduit for many of his early attempts to siphon contamination, transforming her into the monster she now is.

\DndQuote%
{}
{''There are monsters around every corner of this city, but the real monsters are in our minds. No, really, there are monsters that can crawl inside your head, watch out for those ones.''}
{Veo Sjena}
\subsection{Area Details}
\includegraphics[width=0.4\textwidth]{images/../5e.tools/img/adventure/DoDk/087-map-6.03-kleinberg-estate.png}
\includegraphics[width=0.4\textwidth]{images/../5e.tools/img/adventure/DoDk/088-map-6.03-kleinberg-estate-player.png}
\subsubsection{Estate Grounds}
\begin{DndReadAloud}{}
A grand mansion stands atop a rolling hill in the middle of this expansive walled estate. The steep terracotta roofs and grey stone walls are wracked with decay now. The windows are shattered and boarded over, the white-paint trim chipping away. Downhill from the main house is a large stable. Entering the grounds from the main gate, characters can see a stubby fieldstone tower with a small cottage abutting it, but the upper levels of the tower have collapsed. A small guardpost sits beside the creaking wrought-iron gates hanging from the fieldstone walls that surround the estate. Several dying trees dot the lawns, and a brackish pond rests west of the manor house. Several hooded figures wander the property. Some are watering flowers, others lounging about the grounds, while one is throwing rocks in the pond.

\end{DndReadAloud}


\begin{itemize}
\item The property walls are ten feet tall and crested with spikes.
\item Almost all the windows are boarded up. Aside from a small window on the stables, each is purposely blocked from prying eyes.
\item 10 (3d6) other delerium dregs are scattered around the courtyard tending to menial tasks or entertaining themselves.
\item Guard House. A small stone guard house sits next to the front gate. A small cloaked creature with a big toothy grin sits mumbling to himself and playing a game of cards alone.
\item This \textbf{delerium dreg} greets any who approach, asking their business. If characters seem peaceful, it will direct them to see Raven at the stables.
\end{itemize}

\subsubsection{Tower House}
\begin{DndReadAloud}{}
Sickly moss and blighted ivy creep up from the ground and wrap around this tower. There is a small red-roofed cabin off the base, but the upper level of the tower has crumbled away.

\end{DndReadAloud}

\begin{DndReadAloud}{}
\subparagraph{Servant's Quarters}
There are several little creatures working here, they all seem to be reading through notes and relaxing at a small table. The rooms are filled with cots and bunk beds for the various dregs that inhabit the grounds to sleep and rest. There are only a few rooms so they are jammed with various sleeping surfaces. Everything from beds crudely piled on top of one another, or drawers pulled out and stuffed with hay and pillows. The dregs seemed pleased to sleep just about anywhere and the rooms show a disorganized cluttered mess of their living arrangements.

\end{DndReadAloud}


\begin{itemize}
\item 5 delerium dregs are lounging in this room.
\item Henry is keeping watch of the minions.
\item Scrawled on the wall here are several bits of graffiti that seem to be directed at the servants that used to live here. Most notably are the words "Don't forget your candle."
\end{itemize}

\begin{DndReadAloud}{}
\subparagraph{Collapsed tower}
The walls are mostly intact on this second level, but the ceiling is long gone and the upper floor is covered in debris. A few empty cages are placed here.

\end{DndReadAloud}


\begin{itemize}
\item Dregs use these cages as holding for misbehaving companions.
\end{itemize}

\begin{DndReadAloud}{}
\subparagraph{Stables}
Any sign that these stables were once used to house horses is gone, the inside stripped away and made into a makeshift workstation. The stables act as both a storage compartment for finished products of delerium dust, but also housing for the many dregs who reside in the estate. Several of the stalls are lined with hay, dirty blankets, and buckets. The front of the stables by the entrance have been made into a sort of shop where Raven sits at a small table with various potions and alchemical ingredients lined for display.

\end{DndReadAloud}


\begin{itemize}
\item 30 delerium dregs are found here either sleeping, performing menial tasks, or polishing the black coach.
\item Raven often oversees products, shipping and sales that all happen out of the stable.
\item Grunt sits telling stories to several of the delerium dregs resting here.
\item Secret Entrance. At the back of the stable is an old hatch that opens to a ladder down to the old tunnels below the manor.
\end{itemize}

\begin{DndReadAloud}{}
\subparagraph{Black Coach}
Parked in the stable is a beautiful and well-kept black coach with red curtained windows.

\end{DndReadAloud}


\begin{itemize}
\item On rare occasions when the Pale Man travels the streets, he has his haze hulks pull him in this coach. He insists it be polished and cleaned daily.
\end{itemize}

\subsubsection{Mansion Ground Level}
\begin{DndReadAloud}{}
The mansion is a large two-storey property with east and west wings on either side of an entrance foyer accessed by two decorated double doors. Above the doors is a round stained glass window featuring a patterned radial design.

\end{DndReadAloud}


\begin{itemize}
\item All windows are boarded up except for one on the front right that looks into one of the wretched drawing rooms.
\end{itemize}

\begin{DndReadAloud}{}
\subparagraph{Foyer}
Entering the house by the front door you come into an open room. A huge golden chandelier holds one hundred candles that flicker light and shadows around the foyer. Straight ahead a large staircase heads up to a second floor balcony overlooking the foyer. A huge circular stained glass window fills the room with coloured light. To the left and right are doors that lead to either of the wings, while flanking the grand staircase are two statues of majestic horses carved out of marble. The walls are lined with brilliant artwork depicting sunsets on fields, galloping wild horses, a snow covered valley with a small cottage and other beautiful scenic pictures that almost always include horses.

\end{DndReadAloud}


\begin{description}
\item[Trapped Door.]
The east door on the upper level is a heavy wooden door depicting a carving of winged beings soaring up towards the top of the door, and grotesque emaciated-looking creatures below reaching up to try and grab the fleeing angels. The door is \textit{Arcane Lock} and only the password "As Above, So Below" will allow one to access the door. The Pale Man uses this door to access the balcony over the foyer to give orders to his minions
\end{description}

\begin{DndReadAloud}{}
\subparagraph{Wretched Drawing Rooms}
These once well-appointed drawing rooms are now covered in cobwebs and broken furniture, dismantled bookshelves, and dirty animal fur rugs. The rooms have been stripped of any valuables or usable goods.

\end{DndReadAloud}


\begin{itemize}
\item There are 5 (2d4) delerium dregs in each of these rooms.
\end{itemize}

\begin{DndReadAloud}{}
\subparagraph{Decrepit Dining Room}
Food and dirt cover the various walls and surfaces. Paintings are hanging crooked, while some have been broken. There is a long wooden feasting table with four chairs on either side and a grand throne-like seat at the far end. Faces are drawn over hanging portraits and crude words are scrawled across the wall in food substances or possibly blood.

\end{DndReadAloud}


\begin{itemize}
\item There are 8 delerium dregs eating. They know very well that this section of the manor is off limits to guests.
\item The food looks relatively unappetizing, a grey sludge with bits of what might be meat, but the creatures here seem keen to enjoy it.
\end{itemize}

\begin{DndReadAloud}{}
\subparagraph{Kitchen}
The kitchen is a mess with dirty pots and pans stacked almost to the ceiling. There are knives stuck in the walls, food and slop everywhere, and stains on nearly every surface. Two creatures in little makeshift chefs hats and aprons work diligently to make some sort of bland soup for the residents of the house, occasionally throwing in scraps they find on the floor into a big boiling pot.

\end{DndReadAloud}


\begin{itemize}
\item 2 delerium dregs are here and will attack anyone who interrupts them.
\end{itemize}

\begin{DndReadAloud}{}
\subparagraph{Lounge}
A couple couches and armchairs are pushed around an old fireplace. The fireplace is not lit, and has five candelabras set on its mantle with four candles placed in all but one. The furniture has been placed for convenience and conversation and less a proper attempt at outfitting a room.

\end{DndReadAloud}


\begin{itemize}
\item There are 5 (2d4) delerium dregs here.
\item Secret Entrance. There is a hatch on the ground in the corner that opens to an old stone staircase that goes down into the lower tunnels.
\item Secret Door. If a candle is placed in the empty candelabra on the fireplace mantle it opens the back of the fireplace into a small staircase that goes up to the second level.
\end{itemize}

\subsubsection{Mansion Upper Level}
\begin{DndReadAloud}{}
\subparagraph{Abandoned Bedchambers}
These once sumptuous bedchambers are now disheveled and foetid dens. Clothing has been pulled out of the wardrobes to make disgusting nests upon the shredded and broken beds.

\end{DndReadAloud}


\begin{itemize}
\item 2 (1d4) delerium dregs inhabit each bedchamber. Some dress up in the fine clothing that remains, or play with the childrens toys abandoned within.
\end{itemize}

\begin{DndReadAloud}{}
\subparagraph{Lab}
These adjacent rooms have been stripped of their previous furniture and refashioned into laboratories. There are various forms of equipment here. A table in the back has a strange hammer crafted from delerium and a crate of shards and crystals that are seemingly being smashed into dust. There are alchemical supplies and filtering devices that are used to refine the dust at another station. A third station has herbs and chemical formulas and measuring equipment for concocting arcane delerium dust as well as various delerium infused potions. It looks as though another station is for procuring delerium ink for spell scrolls and other uses.

\end{DndReadAloud}


\begin{itemize}
\item 10 (3d6) assorted delerium dregs work around the clock in this chamber at the various tables. Some have goggles or other scraps of lab gear.
\item Explosion. If combat occurs, one of the dregs hurls a half finished delerium concoction at the players, which functions as a \textit{bottled comet}.
\end{itemize}

\begin{DndReadAloud}{}
\subparagraph{Master Bedroom}
Lined with shelves of books, scrolls, folders, documents, and scattered parchments, this room appears to function as a library as well. Several shelves house glass jars with floating objects inside - a jar of eyeballs, a severed hand, an entire human head with one side of it dissolved and melted, several fleshy tentacles, and many other horrors. An old dust-covered bed in the corner appears unused. The desk is covered in various jars of octarine substances as well as an old book bound in red leather and a quill and inkwell of glowing octarine ink.

\end{DndReadAloud}


\begin{itemize}
\item Secret entrance. The bookshelf on the back wall next to the bed has a large blue book titled "the Secret Stairs" - when pulled, it opens to the staircase leading down to the lounge.
\item Amongst the shelves one can find 1d4 magical potions, half of which are Potion of Healing, and 1d6 uncommon Spell Scroll.
\item The spell book on the desk contains the Pale Man's work on the new spells he has created. It is also infused with a \textit{glyph of warding} to cast \textit{dimension door} when triggered. The Pale Man will grab his book and flee to his nest if need be.
\end{itemize}

\begin{DndReadAloud}{}
\subparagraph{Study}
As you step inside you see a mess of paperwork, tables, and several drawings stuck on the wall with red string connecting various sketches to locations in Drakkenheim. A large red circle on a map of the city indicates a focus on the Inscrutable Tower.

\end{DndReadAloud}


\begin{itemize}
\item The Pale Man has put together a rough plan to seize the Inscrutable Tower for himself. The plans seem incomplete and so far most of his research involves how to get into the tower. He has two current ideas; the first note states the front door will only open for a living Academy member wearing \textit{academy rings}, the second is research for creating flying minions to fly into the breach. Beyond this, he has not developed the means to control the Nexus.
\end{itemize}

\subsubsection{Underground Lair}
\begin{DndReadAloud}{}
Heading into the basement reveals a damp and dank cave. There is a thick moisture in the air and an acrid smell mixed with the scent of hot garbage. It's hard to pinpoint where the smells are coming from but it's unpleasant to say the least. There are a few torches along the cavern edge that seem to light the path ahead.

\end{DndReadAloud}


\begin{itemize}
\item The lair is accessed from either the stables or the manor lounge. Originally it was a passage to get to the stables and feed the horses during heavy winters without having to go outside.
\item Since the Pale Man has moved in, he has expanded the caves below to suit his needs.
\item The entire underground section is protected by the \textit{guards and wards} spell. The minions of the house are immune to its effects but any guests or intruders are not.
\item Once entered, the doors to this area vanish into the cave wall. They are still there and it is merely an illusion caused by the \textit{guards and wards} spell.
\end{itemize}

\begin{DndReadAloud}{}
\subparagraph{Collapsed Cavern}
It seems that an early attempt to expand the tunnels failed. There is a dreg crushed under a pile of rubble here, and there is no way past the rubble.

\end{DndReadAloud}


\begin{itemize}
\item The crushed dreg is clutching a candle and a vial of \textit{aqua delerium}.
\end{itemize}

\begin{DndReadAloud}{}
\subparagraph{Holding Cells}
This area contains a barred jail with a small locked door. The walls and floor are caked with refuse and debris. The cell is filled with wriggling piles of flesh. The flesh has various eyes and mouths and occasional limbs or extremities protruding from them. As you approach the cell, the fleshy mounds begin to quiver and a chorus of sorrowful moans echo through the chamber.

\end{DndReadAloud}


\begin{itemize}
\item Two dregs wander the tunnels down here and their job consists of feeding the "unfortunates" as they call them, and occasionally taking one out of the pen to feed Lisa.
\item If either of the two dregs feel threatened or feel the need to aid in an attack, they will run to pull a lever on the wall that opens the pen, allowing the "unfortunates" to wriggle out and consume anything in their path
\item There are 10 gibbering mouthers that can be released from the pen.
\end{itemize}

\begin{DndReadAloud}{}
\subparagraph{Trapped Room}
A wrong turn due to the guards and wards spell, or a misguided search of the corridors and hallways might land you in this room. Once entered, a wooden gate slams over the entrance and the room releases a cloudkill effect.

\end{DndReadAloud}

\begin{DndReadAloud}{}
\subparagraph{Lisa's Den}
This small alcove is covered in various bones, skeletons, bits of old meat and blood soaked scraps of clothing.

\end{DndReadAloud}


\begin{itemize}
\item Lisa resides here. Unless she hears commotion elsewhere in the tunnels or is retrieved by one of the dregs, she stays here eating.
\end{itemize}

\begin{DndReadAloud}{}
\subparagraph{The Nest}
A cave-in at the entrance makes this room very hard to get to. However, this room is the most important room for the Pale Man. Covered floor to ceiling with dripping and oozing delerium sludge and caked with thousands of insectoid looking eggs, this room seems to resemble a sort of nest. Hanging from the ceiling is a large fleshy cocoon-like sack, pulsing gently with a sort of heartbeat.

\end{DndReadAloud}


\begin{itemize}
\item If the Pale Man is reduced to 0 hit points he returns to this nest to recuperate.
\item Destroying the nest means the Pale Man is unable to regenerate and can be killed easily.
\item The cocoon has 100hp and an AC of 12. It is vulnerable to fire, but immune to poison, necrotic, and psychic damage.
\item The entire room is covered in delerium sludge.
\end{itemize}

\begin{DndSidebar}[float=!b]{Developments}
If the Pale Man is destroyed, his minions become confused and terrified. They flee into the city, and slowly revert back into mindless monsters.



\subsection{Pale Man's Research}
The following information can be gathered by either speaking with the Pale Man or studying his writing:


\begin{itemize}
\begin{DndReadAloud}{}
\textbf{\textit{Monstrous Transformations}.} Sufficient eldritch contamination blossoms into a wondrous and self-sustaining mutagenic transformation, after which it cannot simply be purged away. Reversing this is no simple matter. Of course, since the Great Spell may grant any mortal wish, it could be used to restore the previous form of a transformed creature. However, the Academy Directorate jealously guards knowledge of this spell, and I am sure the Archmages would not cast this spell for anyone but themselves. Nevertheless, I believe it may be possible to draw the contaminants out of one body and into that of another.

\end{DndReadAloud}

\item \textbf{\textit{Contaminated Spells}.} The Pale Man knows the following spells, which are also contained in his spellbooks: \textit{delerium blast}, \textit{delerium orb}, \textit{forced evolution}, \textit{conjure the deep haze}, \textit{purge contamination}, \textit{ray of contamination}, and \textit{warp bolt}.
\end{itemize}

The Pale Man's research provides critical clues needed for characters to cure a monstrous transformation. Not only does it reveal that the \textit{wish} spell can accomplish this task, it also suggests researching the \textit{siphon contamination} spell, which can be accomplished by a character who knows the \textit{forced evolution} and \textit{purge contamination} spells.



\subsection{Plight of the Pale Man}
In the event that the Pale Man survives and retreats to his nest to regenerate, the characters may think they have defeated him. The Pale man will plot his revenge and eventually attempt to take the Inscrutable Tower for himself.

The Academy is happy if the Pale Man is killed and his work retrieved. If the Pale Man survives, the Academy expresses the urgency that control of the tower needs to be obtained as soon as possible. The Silver Order is satisfied if the Pale Man was killed or brought in and his estate and research destroyed.



\end{DndSidebar}

\section{Old Town Cistern}
Beneath Slaughterstone Square is a large underground reservoir filled with the waterlogged corpses of the Executioner's many victims. As a result, the macabre chamber has become the favoured scavenging ground for a group of aquatic delerium dregs who serve an eldritch creature known as the Duchess. This nightmarish entity lurks deep below Drakkenheim's forsaken waterways, slowly amassing contaminated thralls while fostering dark dreams to usurp control of the surface.

\includegraphics[width=0.4\textwidth]{images/../5e.tools/img/adventure/DoDk/089-06-010.old-town-cistern.png}
\subsection{Overview}
Your characters might simply stumble into the cistern while exploring the sewers, or by taking the smuggler's passage from the Black Ivory Inn. The cistern lies directly below Slaughterstone Square, so characters may even seek refuge here from the Executioner. You should review that location before reading this section further.

\subsubsection{Adventure Hooks}
\subparagraph{An Invitation from the Duchess}
While exploring the river, sewers, or another aquatic location, the characters encounter a curious group of 12 aquatic delerium dregs and 2 chuuls who are not outwardly hostile. Two step forward peacefully with trepidation, introducing themselves as Blip-Bloop and Slurpy. They implore characters to help them find a suitable gift to bring back to their mistress, who they refer to as the "Duchess". If questioned, they can convey the following information:

\begin{DndReadAloud}{}

\begin{itemize}
\item "Our Duchess will show you great hospitality, I'm sure of it! If you would bring her a gift or perform a service for her, she could reward you greatly, I'm sure of it! What is it you desire? Knowledge! Magic! Health! Our Duchess is as mighty as she is benevolent, I'm sure of it!"
\end{itemize}

\end{DndReadAloud}

\begin{DndReadAloud}{}

\begin{itemize}
\item "She would love a magic trinket, or perhaps a great crystal, of course! I'm sure of it!"
\end{itemize}

\end{DndReadAloud}

\begin{DndReadAloud}{}

\begin{itemize}
\item "The Duchess is ruler of the deep water, all that flows below is her domain! I'm sure of it! She shelters and protects us, and in turn we are her loyal retainers, I'm sure of it!
\end{itemize}

\end{DndReadAloud}

\begin{DndReadAloud}{}

\begin{itemize}
\item "She dwells near a great pool that collects the sky-water. There is a terrible machine upon the surface, but we know a hidden way there by water! I'm sure of it!"
\end{itemize}

\end{DndReadAloud}

If the characters refuse at this first meeting, or attack, it's possible they could encounter this group again (or similar ones) during their future adventures in the city. Don't force this hook upon the player characters: this is something they should investigate on their own initiative!

\subsubsection{Alternate Hooks}
It's easy to miss this adventure site, but since it contains an important \textit{Seal of Drakkenheim}, you may need to troubleshoot. If your players don't take the bait, consider the following:


\begin{itemize}
\item One of the faction leaders asks characters to investigate the following rumour: "The bodies of many slain adventurers and faction agents are said to rest in the cistern beneath Slaughtstone Square. Who knows how many magic items are lost there?"
\item An encounter with aquatic dregs leaves behind a clue related to a character's personal quest, indicating the aquatic dregs possess a lost magic item or missing person (who is now transformed into one of the aquatic dregs or even the Duchess' attendants).
\end{itemize}

\subsection{Interactions}
The Duchess is an otherworldly abomination who holds many mutated aquatic delerium dregs in thrall, using them to gather delerium and lure unsuspecting explorers to its lair to contaminate.

\subsubsection{The Duchess}
The Duchess is an \textbf{aboleth} of unfathomable origin. Whether it was brought here by the unstable rifts of the Haze, or the result of some mutated origin, none can say. It uses its knowledge, telepathy, and mind-reading abilities to make deceptive offers that appeal to characters' deepest desires, manipulating them into embracing contamination so they will become its thralls.


\begin{description}
\item[Personality Trait.]
I rule my subjects like a benevolent matron, and eagerly invite guests to my household so they can serve my will forever.
\item[Ideal.]
The destiny of the Haze cannot be stopped, all will become contaminated and corrupt, all will bask within the eldritch embrace of distant stars.
\item[Bond.]
I have heard the clarion call of the Amalgamation in Castle Drakken. Soon it will open a dimensional rift. I must seek servants and magic so I may seize it for myself.
\item[Flaw.]
I grow angry and impatient when mortals do not regard me with fear and veneration.
\end{description}

The Duchess is deeply in tune with the magic of the Haze. Though it does not know the origins of the meteor, the Duchess is aware that the Haze will eventually overtake the world and can sense the dimensional rift within Castle Drakken (see chapter 9). It is amassing delerium, magic items, and contaminated thralls with the hopes of one day seizing the rift and ruling the surface. It dimly understands the \textit{Seals of Drakkenheim} may be useful in this regard, and so wants to acquire those items as well.

\subsubsection{Duchess' Attendants}
These three wretches are the favoured thralls of the Duchess. Called the Cardinal, the Chamberlain, and the Chancellor, each is a former human transformed by contamination. They use the game statistics for a \textbf{sea hag} with their creature type changed to aberration. The attendants use their powers to appear as humans garbed in flowing blue robes to lure and deceive visitors, but their mutations are more horrific than the aquatic delerium dregs, whom they regard as their subservient kin. Their original memories and personalities have been nearly obliterated by madness, and they now loyally serve the Duchess' strange agenda. They tend the Monument of the Duchess in the Bathhouse Ruins. The Chancellor is in truth the transformed Royal Chancellor of Drakkenheim, Balrin Barthes, and still wears the \textit{Chancellor's Crest}.


\begin{description}
\item[Personality Trait.]
We are the attendants of the Duchess. Come see our customs and partake of our strange ways, and all will be revealed. The water holds great secrets; drink deeply of its knowledge.
\item[Ideal.]
Our blessed offerings help the Duchess in her quest to open the way and the gate.
\item[Bond.]
We escaped the castle but did not make it further. The Duchess has rescued us.
\item[Flaw.]
We do not remember who we once were, but the Duchess tells us we must make preparations to one day rise from the waters to walk upon the surface once more.
\end{description}

\subsection{Area Details}
\includegraphics[width=0.4\textwidth]{images/../5e.tools/img/adventure/DoDk/090-map-6.04-old-town-cistern.png}
\includegraphics[width=0.4\textwidth]{images/../5e.tools/img/adventure/DoDk/091-map-6.04-old-town-cistern-player.png}
The water in these areas is contaminated. Characters who end their turn swimming or submerged in contaminated water must succeed on a DC 10 Constitution saving throw or gain one level of contamination. A character who drinks contaminated water gains one level of contamination.

\subsubsection{[1] Underground Cistern}
This subterranean water reservoir is an eighty-foot-tall cylinder with a domed ceiling made from watertight concrete and flagstones. At the top is a five-foot-wide opening that leads to the receptacle fountain in Slaughterstone Square. A narrow five-foot-wide stone walkway is built twenty feet above the cylinder's base and rings a thirty-foot-diameter basin.

\subparagraph{[1a] Basin}
The reservoir below is almost full, the water line is only two feet below the edge of the walkway. The basin floor is filled with dozens of waterlogged corpses belonging to assorted members of the Hooded Lanterns, Queen's Men, and rival adventurers who ran afoul of the Executioner; their bodies tossed down here by the berserk construct. While most mundane equipment here is ruined, characters can collect a satchel with 2d6 delerium fragments inside, and a sack containing coins, gems, and other valuables worth 500 gp. If you wish, you may place 1d4 common magic items or a single uncommon magic item here.


\begin{itemize}
\item Two chuuls and 12 aquatic delerium dregs pick over the treasures. If the characters found the cistern on their own, these are the group described in the Adventure Hooks section.
\end{itemize}

\subparagraph{[1b] Aqueduct Passage}
Water is siphoned from the basin into this underground aqueduct; the middle of the ten foot diameter opening is roughly in line with the walkway. A barred grate can be removed with a DC 15 Strength check. Beyond is a barrel-shaped passage flooded to the hip. After it travels two hundred feet, there is a break in the ceiling that opens to the Bathhouse Ruins (below).

\subparagraph{[1c] Sewer Canal}
The slippery steps descend this wide stairway to the sewer system at large. A battered and broken gate once kept it closed. Where these canals lead is up to you.

\subparagraph{[1d] Slaughterstone Stair}
This narrow spiral staircase ascends sixty feet to the surface, and ends at a small landing with a short ladder leading to a heavy wrought-iron drainage cover three feet wide. Opening the cover from below requires an action and a DC 15 Strength check. It opens to Slaughterstone Square.

\subparagraph{[1e] Smuggler's Tunnel}
This wide spiral staircase ascends twenty feet then ends at a small landing leading to a roughly hewn arched tunnel four feet wide and eight feet tall. The passage snakes and twists for some distance before connecting to the basement of the Black Ivory Inn (see Outer City Ruins).

\subsubsection{[2] Bathhouse Ruins}
\begin{DndReadAloud}{}
This partially-buried rectangular chamber appears to be a bathhouse of the sort built in ancient antiquity. A pillared aisle borders a central pool set with a fine circular fountain decorated with statues of nymphs basking around a long-forgotten god. Bas-relief carvings along the walls depict primordial scenes of storms, shipwrecks, and vengeful aquatic monsters. At one end of the pool is a semicircular alcove where a massive delerium crystal as tall as a human juts through the floor. It is choked by a mass of thick pulsing veins and perspiring strange black nodules laced with prismatic colour. The glowing stone illuminates the room. Opposite the pool is a wide arched canal from which water flows into the chamber. The air is cold and damp, and the sound of running water echoes through the chamber. The vaulted ceiling overhead is cracked and buckling.

\end{DndReadAloud}


\begin{itemize}
\item The statue depicts the old god Nodens, god of sea and storms.
\item Three sea hags - The Chancellor, the Chamberlain, and the Cardinal - are here. They may also be encountered with 10 (3d6) aquatic delerium dregs and possibly a \textbf{chuul}.
\item They are not initially hostile, and instead invite characters to drink contaminated water they have mixed with delerium dust, touch the delerium geode, and then communicate with the Duchess telepathically. The process will undoubtedly bring characters to the brink of a monstrous transformation.
\item \textbf{\textit{Delerium}.} This is a massive delerium geode cluster. It's exceptionally difficult to remove, but growing around it are 2d6 delerium shards that can be more easily harvested. The Duchess' attendants touch the geode when they communicate telepathically with the Duchess, and refer to it as "the Monument of the Duchess."
\item \textbf{\textit{Treasure}.} The Chancellor has prominently pinned the \textit{Chancellor's Crest} upon their robes. This valuable item is one the \textit{Seals of Drakkenheim}. Amongst the urns, bowls, and sacks kept here are 1d4 uncommon Spell Scroll or potions.
\end{itemize}

\subsubsection{[3] Hydra's Lair}
\begin{DndReadAloud}{}
This chamber is the collapsed ruin of a once-magnificent circular bath.

\end{DndReadAloud}


\begin{itemize}
\item A \textbf{hydra} dwells here in thrall to the Duchess to serve as a personal guardian. It rarely leaves this chamber except to devour worthless minions.
\end{itemize}

\subsubsection{[4] Underground Lake}
\begin{DndReadAloud}{}
This massive cavern is awash in ocatrine light, illuminated by hundreds of tiny delerium crystals embedded in the cavern walls. Water flows from cracks between hanging stalactites, filling the large underground lake with dark water that shimmers and reflects the delerium's sickly light. The air is heavy with moisture, and mucus-like slime washes up on the bank, thick globs of scum cling to the stones. The smallest sound echoes through the chamber with an eerie reverb.

\end{DndReadAloud}


\begin{itemize}
\item The underground lake begins shallow but quickly reaches a depth of thirty feet. The stalactite-covered cavern ceiling is only fifteen feet above the surface.
\item This is the lair of the Duchess. The creature hides beneath the dark water. If it senses hostile intruders, it attempts to lure them close to the banks so it can pull them into the contaminated water easily.
\item If attacked, the Duchess sends out a psychic call to its underlings, and 4 \textbf{merrow} and 12 aquatic delerium dregs arrive five rounds later to help protect the Duchess.
\item \textbf{\textit{Delerium}.} The cavern walls contain countless small sparkling flakes of delerium, but with several hours of excavation characters can unearth 10 (3d6) delerium shards and 1d4 delerium crystals. Characters inspecting the cavern walls can detect that the delerium crystals are embedded in the rock. There's simply no way these shards came from the meteor, they're too far underground.
\item Below the water's surface is a wide underground river that connects to the Drann River.
\item A mound of loose dirt and sand is not visible above the water's surface, but obvious to any who dive underwater. Digging the mound out takes one minute of work, and contains the Duchess's personal treasure. It contains 1d4 delerium crystals, 1d6 watertight tubes each containing a rare \textit{spell scroll}, 3d6 other gemstones or treasures worth 250 gp each, and a rare magic item.
\end{itemize}

\begin{DndSidebar}[float=!b]{Developments}
Regardless of whether or not the characters can learn anything from the Duchess, slaying this abomination puts an end to a terrible scourge before it festers and worsens. Left unchecked, the Duchess could become a potential interloper who interferes in future events. Alternatively, the Duchess may decide to pursue its interests elsewhere, leaving Drakkenheim via the contaminated Drann River. It travels downstream towards the ocean and spreads its corrupting influence into the wider world.



\subsection{The Chancellor's Crest}
Characters who recover the \textit{Chancellor's Crest} have a major find. As one of the missing \textit{Seals of Drakkenheim}, it's one of the keys to the city. Should the Hooded Lanterns and the Queen's Men discover the characters possess it, they'll make overtures to meet them. Both Elias Drexel and the Queen of Thieves use this as an opportunity to recruit the characters to their faction's cause. If the characters turn down the offer, these factions begin drawing up plans to take the \textit{Chancellor's Crest} by stealth, trickery, or even force if necessary.



\end{DndSidebar}

\section{Queen's Park}
Resting in the northern heart of the city between the Inscrutable Tower and Castle Drakken, Queen's Park was built to honor the Queens of Drakkenheim who ruled as monarchs, but those who were royal consorts have also been commemorated here. Over time, more monuments to other well-known women who defined the history of Drakkenheim were constructed here as well. Commonfolk and nobles alike spoke of fondly losing themselves in Queen's Park as they wandered its elegant beauty on summer afternoons. Now covered by the Deep Haze, one might indeed become lost here forever.

\includegraphics[width=0.4\textwidth]{images/../5e.tools/img/adventure/DoDk/092-06-011.queen-harden.png}
\subsection{Overview}
An expedition to Queen's Park, originally seeking key ingredients for \textit{aqua expurgo}, may eventually lead characters to a shocking discovery regarding the fate of a royal family member.

\subsubsection{Adventure Hooks}
\subparagraph{Gathering Eldritch Lilies}
Characters who find the research of Oscar Yoren or hear certain rumours may learn of the magical flowers growing in Queen's Park said to have resistive properties against contamination. Those who research further can learn the following about them:


\begin{itemize}
\item An eldritch lily is a pink-purple flower that grows in pools of contaminated water within natural areas covered by the Deep Haze. The plants have a strange odor, and emit a magical ringing sound audible within ten feet. They register a distinct magical aura when inspected with \textit{detect magic}. A single eldritch lily is worth 100 gp. The petals are potent alchemical reagents, and can be used as a component for restorative spells and many potions.
\end{itemize}

\subsection{Interactions}
While characters explore Queen's Park, they can learn the tragic fate of Queen Lenore von Kessel, and even locate the former royal consort.

The day the meteor fell, chaos erupted in Castle Drakken as the staff and royal family attempted to escape the Haze. At the behest of the royal steward, Queen Lenore's retainers helped her escape Castle Drakken via a secret waterway that exits at Queen's Park. Hindered by Lenore's baggage train, the royal entourage became lost in the spreading mist, and were transformed into monsters. They have remained here ever since, driven to madness with little remaining of their memories.

\subsubsection{Lost Lenore}
\begin{DndReadAloud}{}
This melancholic creature stands six feet tall with long, thin limbs. Her chalky white skin is cracked and rigid like stone, and her hair is a writhing mass of thick wet tendrils that end in toothy maws. She wears an emerald necklace and a distressed black ball gown. She covers her hair with a shawl and her face with a cracked porcelain mask.

\end{DndReadAloud}

\textbf{Lenore von Kessel}, the Queen-Consort of Drakkenheim and wife of Ulrich IV has become a monster that uses the game statistics of a \textbf{medusa}. Replace the normal \textit{Petrifying Gaze} trait and actions with the following:


\begin{itemize}
\item \textbf{\textit{Radiant Gaze}.} When a creature that can see Lenore's eyes starts its turn within 30 feet of her, Lenore can force it to make a DC 14 Constitution saving throw if she isn't incapacitated and can see the creature. The creature takes 22 (4d10) radiant damage on a failed save, or half as much on a successful one. In addition, creatures who fail this saving throw gain one level of contamination.
\item Unless surprised, a creature can avert its eyes to avoid the saving throw at the start of its turn. If the creature does so, it can't see Lenore until the start of its next turn, when it can avert its eyes again. If the creature looks at Lenore in the meantime, it must immediately make the save.
\item If Lenore sees herself reflected on a polished surface within 30 feet of her and in an area of bright light, she becomes blinded and incapacitated for 1d6 hours.
\item \textbf{\textit{Tentacle Hair}.}\textit{Melee Weapon Attack}: +5 to hit, reach 5 ft., one creature. \textit{Hit}: 4 (1d4 + 2) piercing damage plus 14 (4d6) necrotic damage. The target must succeed on a DC 14 Constitution saving throw or gain one level of contamination.
\end{itemize}

Transformed by the Haze, Lenore is a mockery of her former self and abhorrently reacts to her own appearance. She dimly remembers seeking shelter in her grotto, but has no recollection of how she escaped the castle.


\begin{description}
\item[Personality Trait.]
I speak softly in a perpetual daydream, rhapsodically imagining myself to be in a tragic fairytale while casting others in fictional roles.
\item[Ideal.]
The world is a mirror held up to my own glorious reflection.
\item[Bond.]
I obsessively crave beauty, decadence, and relish in art of all kinds, and especially love my garden and my lilies. They are all I have left.
\item[Flaw.]
I fly into an enraged fury at the sight of my own likeness or the names of my lost family members.
\end{description}

Any character proficient in History can instantly recognize Lenore despite her mutations. Before her transformation, Lenore was extraordinarily vain and narcissistic, though was uncharastically kind and gentle to her servants. She wore outrageously ornamented gowns, hats, makeup, and jewelry, including a famous golden necklace set with eleven emeralds. She loved her children only as extensions of herself, and absolutely despised her husband.

\subsubsection{Queen's Retinue}
The transformed creatures that once were Lenore's personal attendants each use the game statistics for a \textbf{dryad}. Now called the handmaidens, they resemble disheveled and misshapen figures clad in moss-covered petticoats, their hair like slimy grass and skin like knobbly bark.


\begin{itemize}
\item Just as they did in their former lives, these unfortunate wretches attend, obey, and protect Lenore, but have only fragmented memories of their own.
\item They act avoidant, timid, and servile, but are resolute in their duties. They use their abilities to mislead and prevent intruders from reaching the royal grotto while wandering the garden searching for delerium and eldritch lilies to bring back to please Lenore.
\item If interacted with, characters may be able to coax the handmaidens into leading them to the royal grotto by succeeding on three DC 15 Charisma checks, perhaps portraying themselves as valiant knights here to help or conjuring some other fable. Otherwise, characters can follow or chase them back to the royal grotto.
\end{itemize}

In addition, many of the haze wights wandering the garden are the undead remnants of the Queen's personal guard, and still bear the crests and livery of House von Kessel.

\subsection{Area Details}
\subsubsection{Entering Queen's Park}
\begin{DndReadAloud}{}
The labyrinthian footpaths of Queen's Park saunter by baroque-style flower beds in geometrical layouts, down winding paths through forested groves, and over ornamented bridges that cross artificial streams and ponds. The lawns and plazas are decorated with triumphal arches, marble fountains, statuary gardens, and hedge mazes, but the park has become overgrown, desolate, and blanketed by thick glimmering mists. Hardy vines cling to the statues like leeches, and strange flowers in unnatural hues bloom in the gardens. A potent pollen hangs in the air - a smell between musk and sweet perfume. It tingles the nostrils, and irritates the eyes slightly. The garden is quiet, the sounds of footsteps linger loudly. There is no sign of a single animal, not even a bird or rodents, but the plant life flourishes despite the lack of sun and influx of rain.

\end{DndReadAloud}


\begin{itemize}
\item Queen's Park Garden is entirely suffused by the Deep Haze (see chapter 5)
\end{itemize}

\subsubsection{Navigating the Garden}
Queen's Park Garden trades the streets, alleys, and ruined buildings of the city proper for footpaths, garden trails, flowerbeds, monuments, and forested groves. As such, characters move through the region just as they do the city streets.

\subsubsection{Searching the Garden}
Characters looking for an eldritch lily must attempt Searching the Ruins for one hour, as described in "Navigating Drakkenheim" in chapter 5. In addition to making Investigation checks, characters may use Arcana, History, Nature, or Survival to help search the gardens. Success requires a DC 15 check result, with any check result of 20 or higher granting an additional success. What the characters find each attempt is based on their number of successes:


\begin{DndTable}{cX}

Number of Successes & Result\\
1 success or fewer & Nothing\\
2 successes & Brackish Pool\\
3 successes & Eldritch Lily Pond\\
4 or more successes & Royal Grotto\\
\end{DndTable}

Two or more failed check results triggers one of the random encounters below. Don't forget that characters are exposed to the Deep Haze each hour they spend here!

\DndQuote%
{}
{''Do not, I repeat DO NOT mess with the flower gardens unless you want a lawn gnome's boot to your face.''}
{Sebastian Crowe}
\DndQuote%
{}
{''I've never been into botany but I found those lilies enchanting.''}
{Pluto Jackson}
\subsubsection{Random Encounters in Queen's Park}
Check for random encounters as normal as characters explore Queen's Park, using this table to generate encounters.


\begin{DndTable}{cX}

1d8 & Random Encounter\\
1 & Turned Around\\
2 & Garden Gnomes\\
3 & Weeping Willow\\
4 & Wight Knights\\
5 & The Three Angels\\
6 & Ghost Lights\\
7 & Monument Of The Queen's Men\\
8 & Handmaiden's Grove\\
\end{DndTable}

\subparagraph{Turned Around}
The characters find themselves at the edge of the garden near one of the many entrances into the grounds. They can make out rows of buildings and streets just beyond a high and moss covered rocky wall that surrounds the garden. They must have gotten turned around somewhere back there.

\subparagraph{Garden Gnomes}
A thin path winds through a grassy clearing. The grass is lush and green, perhaps the most vibrant grass within ten miles of the city. On either side of the path is a beautiful decorated floral garden. The gardens are lined with fluorescent purple flowers accented with hints of orange, yellow, or green. They are not eldritch lilies, but to the untrained eye they may appear to be. A DC 20 Nature check reveals them to be false lilies. Four small garden gnomes stand within the gardens, one on either end of the two floral beds. If the players step onto the grass or pick any of the flowers they immediately hear a voice screech out to them: "KEEP. OFF. THE. GRASS!"


\begin{itemize}
\item The garden gnomes are insane gnome berserkers, who spring to life and attack any who interfere with their lush garden.
\end{itemize}

\subparagraph{Weeping Willow}
A sad-looking willow tree hangs before a vine-covered stone bridge built across a gentle babbling creek at the base of a small waterfall. Purple and yellow vibrant flowers bloom all across the bridge's thick vines. The tree sways in an eerie breeze, while fastened to its trunk a small elegant rowboat bobs gently in the creek.


\begin{itemize}
\item As players approach the tree, it awakens as a mad \textbf{treant}.
\item Abandoned in the boat is a small hairpin with the crest of the von Kessel family upon it, as well as a book titled "The Gardens of Skye" that details exotic lilies and other flora. Inside the cover of the book is an inscription that reads "To my dear mother Lenore, From Katarina."
\item Behind the waterfall is a ten foot wide underground canal that stretches off into darkness. There are no walkways, and one thousand feet down the canal, the water becomes thick delerium sludge that is home to 1d4 animated delerium sludges. The canal gently bends and snakes upward, eventually ending in the cistern of the royal palace of Castle Drakken (see chapter 9).
\end{itemize}

\subparagraph{Wight Knights}
Six haze wights wander the garden, and hungrily attack any living creatures they encounter. They wear the livery of House von Kessel.


\begin{itemize}
\item One of the wights has a letter on his person that bears the seal of the royal steward (see the "Steward's Letter" in the Developments section).
\item With the information from the note, the characters are granted advantage on their next "Searching the Garden" roll.
\end{itemize}

\subparagraph{The Three Angels}
Three statues of weeping angels stand facing the middle of this clearing. The central statue is significantly larger than the other two, and is choked in thick vines with pulsing spores. Between them is an old stone podium entangled in overgrowth.

\includegraphics[width=0.4\textwidth]{images/../5e.tools/img/adventure/DoDk/093-06-012.lenore.png}

\begin{itemize}
\item Removing the plants reveals a map of the gardens engraved upon the stone surface. Any interaction with the statues or the map causes the three statues to spring to life: the smaller two become gargoyles, and the larger one uses the stat block of a \textbf{vrock}.
\item Examining the map gives characters advantage on their next "Searching The Garden" roll.
\end{itemize}

\subparagraph{Ghost Lights}
5 (2d4) will-o-wisps float within the trees of this thick wooded area. If followed, they attempt to lure characters into the woods towards a grove ringed by beautiful glowing fluorescent flowers. A dead mage lays in the middle of the floral patch.


\begin{itemize}
\item Any who inspect the body cause a vengeful \textbf{arcane wraith} to emerge that is immediately hostile. The will-o-wisps attack along with the wraith.
\item Upon the corpse is an uncommon magic item, as well as 1d4 rare \textit{potions} or Spell Scroll.
\end{itemize}

\subparagraph{Monument of The Queen's Men}
A monument of the original Queen's Men rests here. It depicts the regiment of castle guards dedicated to protecting the Queens of Drakkenheim. The monument has been defaced with graffiti and crude drawings left by outlaws of the current Queen's Men. An old burrow behind the monument used to be a campsite but is now home to 2 hungry trolls.


\begin{itemize}
\item The trolls can be bargained with, offering directions to the royal grotto, where they found lots of glowing flowers. In exchange, they want to eat one of the characters' limbs.
\item If the trolls are bargained with, characters have advantage on their next "Searching the Garden" roll.
\end{itemize}

\subparagraph{Handmaidens Grove}
Nestled between hedgerows is a beautiful topiary garden with bushes trimmed into elegant designs depicting mythological creatures. Tending the garden are 7 (2d6) of Lenore's handmaidens (described above).

\subsubsection{Brackish pool}
\begin{DndReadAloud}{}
This small pond is covered in a thin layer of algae that floats upon the glimmering surface of the water. Amongst the aquatic plant life are three large yellow lilies.

\end{DndReadAloud}


\begin{itemize}
\item Examining the water with a DC 10 Investigation check determines that it is contaminated.
\item The three flowers are hypnotic eldritch blossoms that attempt to lure characters into the pond.
\item If the water is disturbed, an \textbf{animated delerium sludge} emerges and attacks.
\end{itemize}

\subsubsection{Eldritch Lily Pond}
\begin{DndReadAloud}{}
This small artificial pond is surrounded by prismatic cherry blossom trees. Reeds spring up around the muddy bank, and throngs of algae grow beneath the water's surface, which is blanketed with clusters of lily pads. Blooming from the water are several large lily flowers with brightly scintillating pink-purple petals. They emit a sweet strange smell and a melodic arcane hum.

\end{DndReadAloud}


\begin{itemize}
\item Characters find this location with three successful checks when they Search the Garden (see above).
\item Standing in the contaminated water are 5 (2d4) of the Queen's handmaidens (see Interactions above), gathering the lilies with care to bring back to Lenore in the royal grotto.
\item There are 5 (2d4) eldritch lilies blooming in the pond, but also 2 (1d4) hypnotic eldritch blossoms (the blossoms ignore the Queen's handmaidens).
\end{itemize}

\subsection{Royal Grotto}
Characters who search for eldritch lilies and achieve four successes find the royal grotto.

\includegraphics[width=0.4\textwidth]{images/../5e.tools/img/adventure/DoDk/094-map-6.05-royal-grotto-tiered-garden.png}
\subsubsection{Tiered Garden}
\begin{DndReadAloud}{}
This regal garden is decorated with opulent stone staircases with pillared railings that ascend between the tiers. The main garden is highlighted by two large shimmering pools, bordered with tiled stone, while multiple lilies float upon the glistening water, pulsing gently with an eldritch glow. Several gnarled and beautiful cherry blossom trees accent the tiers of the garden in an eternal state of bloom, while petals of hot pink, electric blue, and neon violet fall softly, covering the entire area in waves of colour.

The bubbling of fountains and rustling of wind through trees add a peaceful ambiance to the garden, which is decorated with many statues of a tall woman with elegant limbs wearing outrageous dresses and costumes. The face of each statue has been chiseled off . Centred on the upper tier is a statue of the same figure with a lily in her hand, and like the others the face of the statue has been carved away. Light marble walls separate the tiers and are etched with relief carvings of faces of princesses, queens, maidens and many others who served the von Kessels. Moss, ivy, and unnaturally bright flowers cling to many of the statues and walls. Nestled between the staircases and pools is a set of large beaten brass doors, colour shifted and oxidized by years of wear. The doors have a floral patterned border and two massive handles shaped like lilies. These doors are the entrance to the grotto below.

\end{DndReadAloud}


\begin{itemize}
\item Entering the tiered garden awakens 2 shambling mounds and 2d4 hypnotic eldritch blossoms that attempt to lure characters into the contaminated pools.
\item A DC 10 History Check by a player proficient in the skill can determine that these statues are all of Lenore von Kessel
\end{itemize}

\DndQuote%
{}
{''To take a stroll through Queen's Park gardens is to walk the paths of royals past...but even the quickest of trips can leave you one sandwich short of a picnic.''}
{Veo Sjena}
\includegraphics[width=0.4\textwidth]{images/../5e.tools/img/adventure/DoDk/095-map-6.06-royal-grotto.png}
\subsubsection{[1] Grotto Entrance}
\begin{DndReadAloud}{}
Descending the stairs leads to a circular vaulted chamber decorated floor to ceiling with mosaics. The geometric patterns depict magical aquatic humanoids such as merfolk, nymphs, and water weirds delightfully relaxing in fantastically extravagant yet everyday activities: playing sports, enjoying meals, taking baths, and cavorting with terrestrial lovers. The chamber is comfortably cool and filled with a pleasant moist odor; the sound of gurgling water echoes down the hall.

In the middle of the chamber is a circular fountain set with a nymph-like figure holding a trident and a vase. Water pours from the opening to fill the fountain with crystal clear water. Around and within the fountain are arranged painted ceramic vases, glass collectables, and loose coins.

The mosaics extend along short hallways leading to the east and west, ending in smaller rooms of similar construction. Arches set with barred iron gates are cut into the walls in three places. Each leads to a roughly hewn and stepped cavernous passage that descends gently while twisting sharply south.

\end{DndReadAloud}


\begin{itemize}
\item A character proficient in mason's tools or with stonecunning can tell the cavernous passages are completely artificial, and not a naturally-formed cavern.
\item 10 (3d6) handmaidens appear using their \textit{tree stride} power in the roots within the following chambers and make a last ditch attempt to drive away intruders.
\end{itemize}

\subsubsection{[2] Wellspring}
\begin{DndReadAloud}{}
Two yawning stone pipes disgorge streams of sweet-smelling water into a pool in this cavernous chamber. The water is luminescent with a pinkish-purple sheen. Thick roots from the trees above have broken through the walls to sip the collected water. The pond empties gently down a passage that bends southwest.

\end{DndReadAloud}


\begin{itemize}
\item The pool is three feet deep, and the water within is contaminated. The pipe openings are about a foot wide.
\end{itemize}

\subsubsection{[3] Statuary Chamber}
\begin{DndReadAloud}{}
The stepped caverns converge in a rectangular chamber. Part of the chamber is finished stonework, the rest remains roughly hewn. Several statues have been carved directly from the rock. The finished chamber contains twenty statues depicting Lenore von Kessel lining the long walls. A large statue of the royal family sits at the end of this chamber, depicting King Ulrich, Queen Lenore, and their two daughters. The faces of any statue depicting Queen Lenore have been sheared off.

\end{DndReadAloud}


\begin{itemize}
\item An inscription on the family monument reads "...for my lovely Lenore, on the anniversary of our royal union."—King Ulrich von Kessel IV
\end{itemize}

\subsubsection{[4] Underground Pond}
\begin{DndReadAloud}{}
An underground pond fills most of this large artificial cavern, filled by a flowing stream to the west. The water glows softly with prismatic light, and across the surface are dozens of purple-pink eldritch lilies. The great roots from the trees above have broken through the ceiling in places to drink deeply from the lambent water.

An exquisitely upholstered velvet chaise lounge with gold-painted armrests and plush pillows is set atop several overlapping woven rugs. Scattered around the furnishings lies the extensive baggage train of a fabulously wealthy noblewoman. A mahogany chest of drawers is overflowing with pompous gowns, fashionable dresses, multicoloured wigs, oversized hats, and fantastically impractical shoes. An oaken coffer is brimming with elaborate jewellery, and several smaller containers, suitcases, and baskets are overstuffed with more clothing, artwork, gems, and coins.

Several portraits and paintings are arranged about the chamber. Many paintings show a noblewoman wearing a golden necklace set with eleven emeralds, but within each the face has been torn off. One painting shows the woman posing with a noble family: a young teen man, two girls - one perhaps eight, the other no more than five - and a portly and bearded man wearing a crown.

\end{DndReadAloud}


\begin{description}
\item[Lenore.]
resides here. If anyone touches the eldritch lilies or provokes her flaws, she tears away her veil and attacks while prismatic light emanates from her eyes, nose, and mouth.
\item[Treasure.]
Queen Lenore's emerald necklace is worth 10,000 gp. However, the necklace is famous. Anyone with the coinage to \textit{buy} it knows that it would be immediately recognized as belonging to the royal family. Lenore wears a diamond-set \textit{ring of protection} bearing the royal seal. The other costumes and accoutrements here are mostly ruined, but 1d4 pieces worth 250 gp are salvageable. 1d6 delerium crystals and 2d6 shards can be harvested here.
\item[Eldritch Lilies.]
Lenore's handmaidens have collected 35 (10d6) eldritch lilies and gathered them in the pool here, but they don't normally grow in this location.
\end{description}

\begin{DndSidebar}[float=!b]{Developments}
Characters may have come to Queen's Park seeking flowers, but could leave with a royal problem.



\subsection{Lenore}
Lenore can be kept tranquil or pacified by characters who play into her delusions, or temporarily blinded if she sees her own face. It may be possible to capture her or bring her out of the Garden, but as a contaminated creature she'll die in a few days if taken outside the Haze.

The former Queen is a powerful symbol of the old monarchy. As such, should Elias Drexel and the Hooded Lanterns learn of her survival, they demand Lenore be placed under their custody in the Drakkenheim Garrison. If the characters accept, the Hooded Lanterns reward the characters with an appropriate boon for their deed. This act alone cements the support of the Hooded Lanterns, and they'll look favourably upon an alliance with another faction the characters might be working alongside. However, should the characters refuse, the Hooded Lanterns cast them as royal kidnappers. Elias Drexel expends every effort to rescue Lenore and bring the party to justice.

If Lenore was killed and the characters inform Hooded Lanterns, the soldiers are devastated, but don't seek retribution. They understand the characters did what they had to do against a contaminated monster. Even still, this blunder means Elias Drexel no longer believes the characters have the best interests of Drakkenheim in mind. Unless they prove otherwise through greater deeds, the Hooded Lanterns become ambivalent about working with the party in the future.

\subparagraph{Curing Lenore}
Unfortunately, the Hooded Lanterns don't possess the means to reverse Lenore's transformation, and must turn to the other factions for help:


\begin{itemize}
\item The Silver Order makes an honest effort to use cleric spells to cure Lenore, but this fails. As such, they suggest a swift and merciful death is the most honorable course for a contaminated creature such as Lenore. They demand the Hooded Lanterns euthanize her as a condition for any future alliance.
\item The Amethyst Academy offers to help find a way to reverse her transformation in return for help retaking the Inscrutable Tower. They suggest the characters seek out the Pale Man to acquire his research. If these works are turned over to the Academy, the mages will aid the characters in developing the \textit{siphon contamination} spell.
\item The Queen of Thieves suspects a wish from the Crown of Westemär could be used to cure Lenore, but reveals this information carefully. Naturally, if the Hooded Lanterns keep Lenore in the Drakkenheim Garrison, the Queen's Men stage a kidnapping attempt for ransom!
\item The Followers of the Falling Fire don't care much for the old monarchy, but Lucretia Mathias offers to call upon her divine powers to restore Lenore to her human form. This requires a great show of faith: Lucrecia Mathias only agrees to perform the miracle if Elias Drexel and the Hooded Lanterns leadership agree to join the Followers of the Falling Fire by taking the Sacrament.
\end{itemize}

\subparagraph{Queen-Regent}
If her transformation is somehow reversed, Lenore has the statistics of a human \textbf{noble}. While she could possibly take up the throne as a Queen-Regent, this is an extremely dubious claim. Lenore was wed into House von Kessel, and so her political authority comes from her spouse and children. Worse, while Lenore adores the trappings of power and wealth, she loathes actual political responsibility. She wields her influence for personal gain only. Still, she could prove instrumental in confirming the identity of her offspring, but she doesn't know where they are or what happened to them.



\subsection{Eldritch Lilies and Aqua Expurgo}
Creating one dose of \textit{aqua expurgo} requires alchemical tools and eight hours of work by a character who knows the formula and can cast spells of 3rd level or higher. The recipe calls for one eldritch lily flower, 150 gp in alchemical reagents, and two ounces of blood taken from a contaminated monster.

Oscar Yoren and his apprentices can produce the potion given the materials, as can the Amethyst Academy, the Hooded Lanterns, and the Followers of the Falling Fire if they are given the formula.

\subparagraph{Cultivating an Eldritch Lily Garden}
Rather than continually revisiting the garden to acquire more flowers, inventive characters might propose creating a garden of their own. Characters who carefully collect two dozen complete eldritch lily plants - not just the flowers - can begin their own plot. They'll need a location within the Haze with ample contaminated water, as eldritch lily plants taken out of the Haze wilt and die within twenty four hours. The strongholds of the Followers of the Falling Fire and the Hooded Lanterns, as well as Reed Manor, are suitable sites for a safely accessible lily garden, but the Amethyst Academy has the means to create one at their remote strongholds as well.

This ensures a consistent supply of reagents for crafting \textit{aqua expurgo}, and working with a faction or Oscar Yoren allows each character to receive one dose at the outset of any mission. Characters who wish to create additional doses must either do so themselves or provide extra gold and materials to ensure the garden isn't overfarmed.



\subsection{The Steward's Letter}
\begin{DndReadAloud}{Johann Eisner, Royal Steward}
Please excuse the haste of this missive, and the messenger who bears it. Our present circumstances are both dire and desperate.

I have charged this man with ensuring your safety during this calamity. He is a valiant knight, and will escort you and your handmaids out of Drakkenheim. It is imperative you flee the castle immediately via the garden passage.

Do not return to the royal apartments, nor take the main castle gate.

I beg you, your Majesty, do not wait for anyone else, nor burden yourself with trinkets and treasures, nor tarry in the gardens! Escape the city with all speed.

I promise, my queen, I shall not leave myself without the king and your son. I pray that your daughters remain safe. We shall all meet again in health and mirth soon.

Yours in faithful service,

\end{DndReadAloud}



\end{DndSidebar}

\section{Rose Theatre}
The Rose Theatre was the latest in a long line of playhouses in Drakkenheim that typically remained for a few decades before horribly burning down. The three story circular theatre with a thrust stage and open-air auditorium was once a hub for popular entertainment in the city. Throngs of common folk attended bawdy comedies, exaggerated and bloody histories, and dramatic tragedies written by highly competitive bards. Now it lays abandoned within the Haze, devoid of any laughter, music, or applause.

\includegraphics[width=0.4\textwidth]{images/../5e.tools/img/adventure/DoDk/096-06-013.rose-theatre.png}
\subsection{Overview}
In this adventure, characters explore the mysterious Rose Theatre, which is oddly sheltered from the city ruins by an invisible barrier, guarded by uncontaminated elemental creatures and arcane constructs. Here, the characters encounter Ryan Greymere, a talented \textbf{archmage} who abandoned the Amethyst Academy. Though she has been branded a malfeasant wizard, Ryan Greymere doesn't practice contaminated magic, unlike Oscar Yoren and the Pale Man. Instead, she's crafted a potent suit of magic armour that protects against contamination.

It's up to the player characters how to deal with her. They can slay her to claim her magic items and seize her spellbooks, but could gain much by negotiating with her. They might even be able to broker an agreement to admit her back into the Amethyst Academy.

\subsubsection{Adventure Hooks}
\subparagraph{Dangerous Minds}
Ryan Greymere has been in the sights of the Amethyst Academy for some time. Although she was a dear friend of Eldrick Runeweaver, she was also a fierce rival. They want to know what she is up to in the ruins and what her research has uncovered about delerium.

\subparagraph{Misguided Mage}
Knights of the Silver Order would like to see her brought to justice. As a practitioner of magic who has abandoned the Edicts of Lumen, the Silver Order believes it is their duty to destroy her research and bring Ryan Greymere in for imprisonment, but they are just as happy if the players take her down.

\subsection{Interactions}
Ryan Greymere aggressively responds to intruders into the Rose Theatre, but intends to incapacitate trespassers so she can find out why they are there. Since she is hiding alone in the Inner City under magical defenses, she expects anyone who has penetrated the magical barrier is an Amethyst Academy assassin sent to kill her. However, she'll stay her hand if intruders attempt parley, but she'll try to flee before she can be captured.

\subsubsection{Ryan Greymere}
\includegraphics[width=0.4\textwidth]{images/../5e.tools/img/adventure/DoDk/097-06-014.ryan.png}
Once highly regarded within the Amethyst Academy, Ryan Greymere's trailblazing research led to many early discoveries about delerium, and she designed many of the procedures used by the Academy to make magic items using the crystals. However, she shocked the Directorate by presenting a dramatic plan for the Amethyst Academy to overtly seize control of Drakkenheim by force. Such action would grossly overstep the Edicts of Lumen, but Greymere flippantly replied that the realms and clergy were in no place to actually enforce the Edicts anymore. As a result, she was branded malfeasant, disavowed by the Academy, and most insultingly, the Directorate took credit for her discoveries. She fled into Drakkenheim, and now plans on gaining access to the Inscrutable Tower and continuing her research in private.

Ryan Greymere is a human woman in her mid-forties. Armoured head to toe and wearing white robes, her spectral appearance has garnered her the moniker the "Iron Banshee".


\begin{description}
\item[Personality Trait.]
Get to the point on a topic that excites me, or don't waste my time.
\item[Ideal.]
Only wizards can unlock the mysteries of delerium and protect the world from the Haze, but we cannot do that if we are hindered by the Edicts of Lumen.
\item[Bond.]
I greatly miss my friends and colleagues in the Amethyst Academy, but I'm furious about the way I was treated by the Directorate.
\item[Flaw.]
I'm not wasting my time with boot-licking fools looking to restore a worthless monarchy.
\end{description}

Ryan Gremere is a human \textbf{archmage} who wears \textit{hazewalker plate}, Winged Boots, and carries a \textit{wand of lightning bolts}.

In addition to the spells she has prepared, her spellbook contains \textit{purge contamination}, \textit{ride the rifts}, and \textit{octarine spray}.

\subparagraph{Survival in the Haze}
Ryan Greymere uses the \textit{Mordenkainen's Magnificent Mansion} spell each day to conjure a sanctuary where she can rest. She doesn't keep living quarters in the theatre.

\subparagraph{Bound Elementals}
Ryan Greymere's elementals were conjured and bound to her service via \textit{planar binding} and unquestionably obey her.

\subsection{Area Details}
\subsubsection{Approach}
\begin{DndReadAloud}{}
The Rose Theatre is a ring-shaped three-storey wooden playhouse with the central auditorium open to the sky.

\end{DndReadAloud}

\subsubsection{Force Barrier}
The whole theatre is protected by a hemispherical \textit{wall of force} generated by a massive delerium geode. The barrier encloses the entire building, its perimeter about twenty feet outside the walls and about fifty feet above the rooftops of the theatre.

Nothing can physically pass through the wall. It is immune to all damage and can't be dispelled by \textit{dispel magic}. A \textit{disintegrate} spell destroys the wall instantly, however it reforms ten minutes later. The wall also extends into the Ethereal Plane, blocking ethereal travel through the wall. Teleportation magic can bypass the wall, such as a \textit{dimension door} or \textit{misty step} spell.

\subsubsection{Playhouse}
The playhouse is a circular open air auditorium with a rectangular stage. Most performances were held here during daylight hours to use natural lighting.

\subparagraph{The Stage}
This large rectangular wooden platform is built at the north end of the auditorium, connected to the backstage house by exits leading into the wings. Two large wooden pillars support an overhanging rooftop that extends from the backstage house.


\begin{itemize}
\item Ryan Greymere has set up her arcane lab upon the stage.
\end{itemize}

\subparagraph{Balcony Galleries}
These wooden and pillared galleries are set with several raked rows of wooden benches, separated by evenly-spaced aisles. At the centre is a special box seat for the wealthiest patrons, and many disintegrating cushions, pamphlets, bits of old food and pitchers are left behind.

\subparagraph{Auditorium}
This gravel-filled courtyard has no seating: commonfolk stood here to watch performances instead. It is open to the air above.


\begin{itemize}
\item An \textbf{air elemental} and a \textbf{water elemental} patrol here.
\end{itemize}

\subparagraph{Barrier Generator}
This complex arcane apparatus set up on the stage powers the force barrier. Anchored by a delerium geode, it requires a constant supply of delerium to maintain. It has AC 20 and 100 hit points. If destroyed, the force barrier fails, triggering a random arcane anomaly.

\subsubsection{Backstage}
The backstage house is a rectangular construction built into the overall circular shape of the entire theatre that contains the storage and administrative offices of the theatre.

\subparagraph{The Wings}
Mouldering red curtains are hung over the doorways leading into this narrow carpeted hallway. During performances actors made their entrances and exits via this corridor; it connects the stage to the green rooms backstage. There are many shelves holding old props, several cracked mirrors, and old lamps back here. Actors have left behind notes, love-letters, and graffiti.

\subparagraph{Green Room}
Actors waiting for their cue rested here during a performance, and socialized here between performances. Bottles of wine, flower bouquets, fluttering pages from play scripts, and notebooks filled with poems, sonnets and bawdy drawings are scattered about.

\subparagraph{Wardrobes}
Contains an array of costumes. Many appear fantastically decorated, but even casual inspection reveals the gemstones and gold trim are glass beads and painted fakery.

\subparagraph{Dressing Room}
Actors donned their costumes and makeup in these rooms. Several broken mirrors and drawers of powder, wigs, and makeup are still here.

\subparagraph{Offices}
The writers and managers used these rooms. Writing desks, old playbills, and records of the ticket sales can be found here. A lockbox holds the funds, totalling 1,000 gp.

\begin{DndSidebar}[float=!b]{Developments}
As long as the player characters agree to help Ryan Greymere she is willing to grant them access to her research. Characters who either obtain her notes, or work with her can find the following.



\subsection{Ryan Greymere's Research}
Characters can learn the following from speaking with Ryan Greymere or studying her writing:


\begin{description}
\item[Formation of Delerium Crystals.]
"Magic attracts more magic. Under extraordinarily rare circumstances, magical energies crystallize into delerium. The stones slowly draw more raw magic into the world from distant planes, possible realities, and alternate dimensions. The rate is exponential. Considerably more leakage occurs around vast quantities of delerium. These chaotic energies are inherently destabilizing. Unfocused magic conjures up paradoxical improbabilities at random, leaving behind lingering eldritch contamination. Uncontained, these magical forces manifest as the Haze. More delerium forms within, and so the process accelerates."
\item[Emergence of Dimensional Rifts.]
"As it expands, the Haze both transforms the landscape, but it also conjures magical energy from the Space Between Worlds. I'm concerned that in areas where magic was already concentrated or where the borders between worlds are thin, it could result in catastrophic dimensional rifts."
\item[Possible containment.]
"Delerium presents an existential danger, but also opportunity. The concentration of delerium in Drakkenheim could provide an infinite supply of magical energy if directly harnessed, rather than simply crafting items individually. I've found promising results so far at collecting the energy given off by the stones - the problem is storage and containment. The abjurations required are just outside my wheelhouse. If it could be accomplished at-scale, not only would it be possible to fuel endless magical creations, it might actually slow the growth of the Haze, even hold it at equilibrium."
\item[Hazewalker Plate and Bottled Comet Schematics.]
Ryan Gremere's research notes include the formula and blueprints to create these items (See Appendix D).
\end{description}

Ryan Greymere keeps her research within a \textit{Leomund's Secret Chest}. If she's killed, the contents become forever lost on the Ethereal Plane.



\subsection{Ryan Greymere's Wishes}
Ryan Greymere wants to be allowed to continue her research in Drakkenheim and is upfront with characters about her relationship with the Academy, and her disdain for the monarchy. She scoffs at the idea of restoring the royal bloodline, and feels the world is better off without kings and queens. If player characters decide to help Ryan Greymere, she asks that they bring her a delerium geode for each character who wants \textit{Hazewalker Plate}. She will construct the armour for the characters but she requires access to the Inscrutable Tower.

\subparagraph{Rejoining the Academy}
Ryan Greymere holds no grudges against Eldrick Runeweaver or River, but is very open with her opinions on the Academy Directorate. However, she understands that a lot of their decisions are forced on them by the Edicts of Lumen, and regards the Silver Order and the oppression of mageborn responsible for limiting her potential to help the world and solve the issues of Drakkenheim. It may be possible for characters to convince the Academy to set aside their feud and have Ryan Greymere join once more. Ultimately, she would most like to be elected Archmage of Drakkenheim, but she might also collaborate with Eldrick Runeweaver. In either case, she proves to be a valuable asset, and her research and inventions could greatly advance the Academy's goals. Unfortunately, Ryan Greymere's fierce political opinions may complicate attempts to have the Amethyst Academy collaborate with other factions.



\subsection{Crafting Hazewalker Plate}
Building a suit of \textit{Hazewalker Plate} requires magically-enchanted smithing tools and a well-equipped arcane forge (such facilities can only be found within the Inscrutable Tower or other distant Amethyst Academy strongholds). It takes thirty days of work by a character who knows the schematics and can cast spells of 7th level or higher. The recipe calls for one delerium geode, 10,000 gp in alchemical reagents, and five ingots of refined meteoric iron in addition to enough conventional steel and materials for making either half-plate or plate. Ryan Greymere can craft a suit given facilities and materials, as can the Amethyst Academy.



\end{DndSidebar}

\section{Saint Vitruvio's Cathedral}
Saint Vitruvio was a mighty paladin of the Sacred Flame, the patron saint of Drakkenheim. Along with the ancient gold dragon Argonath, this righteous warrior saved Drakkenheim from a marauding army led by chromatic dragons that descended upon Westemär from the eastern vale centuries ago. Argonath gave their life in that battle, and Saint Vitruvio subsequently died slaying the half-dragon warlock behind the terrible scourge. Their bones and equipment have been kept as holy relics since, as Saint Vitruvio prophesied the gold dragon could be resurrected should Drakkenheim ever fall to darkness.

This gargantuan Cathedral of the Sacred Flame was built to house their remains. However, it is a relatively recent construction, only completed a little over a century ago to replace another large chapel that was torn down during the expansion of the city walls. Long before, the site was home to several crypts and crematory gardens, and the dead are still interred in the catacombs below the contemporary building. Following its completion, worshippers would travel for miles around to make prayers and offerings to the Sacred Flame here. After the meteor fell, these hallowed halls became a charnel house - the lair of the monstrous \textbf{Lord of the Feast}.

\includegraphics[width=0.4\textwidth]{images/../5e.tools/img/adventure/DoDk/098-06-015.cathedral.png}
\subsection{Overview}
In this adventure, the player characters must spearhead an assault or infiltration into Saint Vitruvio's Cathedral, where they will face off against a terrifying garmyr warband. Then characters may explore the catacombs below to recover the relics within. The cathedral is one place that can bring all the factions together, but in the aftermath their disagreements blossom into open conflict.

\subsubsection{Adventure Hooks}
\subparagraph{Reclaim the Cathedral}
Once characters have secured passage through a city gate, obtained at least one Seal of Drakkenheim, claimed the clocktower, and otherwise proven their competence and trustworthiness to either the Hooded Lanterns, the Silver Order, the Followers of the Falling Fire, or the Queen of Thieves, they are approached by a faction leader about leading an assault on the cathedral.


\begin{itemize}
\item For the Followers of the Falling Fire and the Silver Order, the cathedral is an important spiritual landmark. Not only are many important holy relics kept here, but also the magical brazier of the Sacred Flame could be re-ignited. Its protective aura might even hold back the Haze, at least around the cathedral.
\item Both the Queen of Thieves and the Hooded Lanterns wish to access the House von Kessel vaults located beneath the cathedral. This important archive may include copies of the king's will or other documents outlining the line of succession, and possibly other political secrets.
\item Finally, the Saint Vitruvio's Phylactery is believed to be lost here. As one of the six \textit{Seals of Drakkenheim}, its recovery is essential if there is to one day be a new monarch... or someone wishes to use the power of the throne. Except for the Amethyst Academy, every faction wants to recover it.
\item The Amethyst Academy doesn't have a personal interest in recovering the cathedral. However, they may throw their support behind one of the other factions to garner favour with them later.
\end{itemize}

\includegraphics[width=0.4\textwidth]{images/../5e.tools/img/adventure/DoDk/099-06-016.mutation.png}
\subsection{Interactions}
A direct attack by the player characters alone is likely folly, so they should seek aid from the factions to defeat the Lord of the Feast.

\subsubsection{Faction Alliance}
Many factions can accomplish key objectives at the cathedral, which the characters can use to facilitate cooperation between them. This is one of the few moments in the campaign where all the factions are willing to pool their resources and work together with the player characters, so long as the following conditions are met by the characters:


\begin{itemize}
\item Demonstrated their trustworthiness and full loyalty to the faction's ideals and goals.
\item Acquired at least one \textit{Seal of Drakkenheim}.
\item Secured the clocktower and Temple Gate.
\item Created an alliance with two other factions.
\end{itemize}

The factions must keep a reserve force and a garrison for their current holdings, so each can only work with a few strike teams for this operation. Plus, large numbers of people moving in the city will attract more monsters, and quickly. They propose an assault mission by the player characters supported by a handful of faction strike teams.


\begin{itemize}
\item Elias Drexel and Theodore Marshal are both willing to personally join the battle, and the other factions offer up boons in advance to help the characters.
\item The Inner City is filled with monsters, so moving material, troops, or large traps into the city over time is difficult and limited to what individual people can carry and set up in the field.
\item Once the cathedral is taken, the factions will post troops there on rotating shifts to defend it, but this will leave their forces stretched quite thin as a result.
\end{itemize}

\subsubsection{Planning the Attack}
The player characters should be instrumental in planning the attack on the cathedral and devising a way to bring down the Lord of the Feast with the faction leaders. A decisive blow must be struck. The factions can't fight a battle of attrition in the streets or hold ground in the city because their troops will quickly become contaminated. The factions can provide a few doses of \textit{aqua expurgo} or cast \textit{purge contamination} to heal an individual strike team, but none of the factions can afford to do this en masse.

\includegraphics[width=0.4\textwidth]{images/../5e.tools/img/adventure/DoDk/100-06-017.contaminated-hand.png}

\begin{itemize}
\item The best intelligence and scouting reports the factions have indicates that the garmyr number in the hundreds, perhaps even as many as two thousand. While slaying the Lord of the Feast won't suddenly cause the garmyr to disappear, it should break any cohesion between the warbands.
\item If the characters are having trouble coming up with a plan, one of the faction leaders might suggest trying to draw out the Lord of the Feast with bait of some kind. Then, faction strike teams can lure away the bulk of the garmyr with a diversion, while the player characters take down the Lord of the Feast (possibly with a faction leader). Embrace your player's creativity in coming up with what might be suitable bait and what might be a good diversion!
\end{itemize}

\subsubsection{Lord of the Feast}
Game statistics for this monstrous creature can be found in Appendix A.


\begin{description}
\item[Personality Trait.]
I live for the hunt, and enjoy nothing more than fresh meat.
\item[Ideal.]
It is the right of the strong to consume the weak.
\item[Bond.]
Drakkenheim is the perfect hunting ground for an alpha predator such as myself.
\item[Flaw.]
I am gripped with bloodlust that only subsides when I rest in the cathedral.
\end{description}

\subsection{Area Details}
\includegraphics[width=0.4\textwidth]{images/../5e.tools/img/adventure/DoDk/101-map-6.07-st-vitruvios.png}
\includegraphics[width=0.4\textwidth]{images/../5e.tools/img/adventure/DoDk/102-map-6.07-st-vitruvios-player.png}
\subsubsection{Cathedral Plaza}
\begin{DndReadAloud}{}
Broad cobblestone streets wrap around an immense baroque cathedral in this august square. The structure dominates the plaza, completely dwarfing the nearby buildings with a two-hundred-fifty-foot-tall spired central dome that looms large over the gently sloping terracotta rooftops of the surrounding townhouses. The main dome is flanked by four stocky towers meticulously decorated with religious iconography and angelic sculptures. Between each are semicircular apses lined with stained glass windows fifty feet high, thus the building has a roughly square footprint. The outer facade has rows of alcoves containing life-size statues of valiant paladins and devoted clerics. Broad steps lead up to the columned entrance portico where a massive rose window rests above three sets of golden double doors embossed with heavenly visages.

Splayed out before the majestic building is a sprawling and ramshackle camp of makeshift hovels, smoking pyres, and improvised barricades. Tall wooden stakes are driven into the ground along the streets, displaying macabre trophies of bones, ribcages, skulls, and body parts. Guttural yelps and furious howls echo down the streets, and the overpowering scent of decaying offal, feces, and wet animal fur wafts about the haze-filled courtyard.

\end{DndReadAloud}


\begin{itemize}
\item 70 (20d6) garmyr and 14 (4d6) worgs devour their latest kills in the encampment.
\end{itemize}

\subsubsection{Cathedral Interior}
\begin{DndReadAloud}{}
The Cathedral of Saint Vitruvio houses a vast interior space. A monumental rotunda topped with a two-hundred-fifty-foot-high dome forms the central sanctuary where the Brazier of the Sacred Flame is set. The ribbed roof features an elaborate fresco depicting an angelic host descending from on high, lifting the wretched and the lost towards a divine light symbolized by a blazing sun. A row of stout yawning windows rings the base of the dome, bordered by multifaceted mirrors that create angular beams of light leaving spotlights on the floor and walls around the central altar.

Four halls extend in each cardinal direction like wheel spokes from the rotunda. Colossal marble columns decorated with cherubs line the aisles down each and support high vaulted ceilings one hundred twenty feet overhead. Murals upon the ceilings of the naves and aisles depict scenes of saints, angels, and martyrs performing righteous acts of mercy, compassion, and sacrifice. Each hall ends in a semicircular apse set with arched stained glass windows fifty feet high, except the southern hall that serves as the main entrance, above which is a rose window.

Unnatural hues from outside flood this once holy place through the many windows, and cast the chambers in an unsettling and otherworldly pallor. The sanctuary has been defiled by the countless carcasses of the garmyr's many slaughters. Broken bones and severed limbs are strewn about carpets of stretched skin and furry hide. Debris and refuse are caked into the cracks and crevices, vibrant red blood smears the floors and walls. Flensing blades, barbed chains, and meat hooks dangle from the torch sconces. Unrecognizable forms of hanging meat decorate the elaborate chandeliers above. At the heart of the cathedral is a tiered platform upon which rests a five-foot-wide granite brazier. Instead of a glorious golden flame, piled within is a mountain of blood-soaked skulls.

\end{DndReadAloud}


\begin{itemize}
\item During the day, the \textbf{Lord of the Feast} rests here with 14 (4d6) garmyr berserkers, and 10 (3d6) hell hounds. The great beast hunts nightly from dusk till dawn when only 7 (2d6) berserkers, and 4 (1d6) hell hounds remain behind.
\item A mutated garmyr \textbf{chaplain} wearing the tattered robes of a Flamekeeper conducts blasphemous rites before the altar of skulls. It carries the Saint Vitruvio's Phylactery.
\end{itemize}

\subparagraph{Catacomb Passages}
Two arched doorways with oak doors open to staircases that descend down into the cathedral catacombs. Both doors are brightly illuminated by beams of light created by the reflecting mirrors in the dome above.

\subsubsection{Sidechambers}
Small groups of garmyr have claimed these rooms as their own, and each is now in a state of disarray and gore.

\subparagraph{Bell Towers}
The central dome of Saint Vitruvio's is anchored by four great bell towers. The lower floors are great galleries where statues and artworks are displayed, with the belfrys accessed by narrow staircases and ladders.

\subparagraph{Chandlery}
This western side room is a workshop for making the many candles, holy oils, and flammable fluids used in the holy ceremonies of the Sacred Flame.

\subparagraph{Sacristy}
These eastern rooms are where the Flamekeepers donned their vestments, prepared for their sermons, conducted personal meetings, and meditated privately. A small loft houses a few cots used by the nightkeepers, but most of the clergy lived at Saint Selina's Monastery.

\subparagraph{Vestibules}
Since the cathedral could be packed with thousands of worshippers on high holy days, these additional entrances helped ensure the flood of people could easily enter and exit.

\subsection{Cathedral Catacombs}
Beneath the cathedral dome lies an ancient burial crypt for former paladins and Flamekeepers who pledged their life in service to the Sacred Flame. The Lord of the Feast and his warpack leave these hallowed halls abandoned and forgotten, instead seeking fresh meat in the city ruins.

\includegraphics[width=0.4\textwidth]{images/../5e.tools/img/adventure/DoDk/103-06-018.cathedral-catacombs.png}
\subsubsection{Light Beam Puzzle}
The treasures of the Cathedral of Saint Vitruvio are sealed behind warded doors of stone, which slide open only when a beam of heavenly light is shone into the aperture lenses upon them. A beam of light originates above each of the Ascendant Archangel statues (detailed below), then characters must direct the beams down the halls by bouncing the light from mirrored shields found throughout the catacombs. Solving the puzzle requires placing the varying numbers of shields in specific positions then facing them in different directions to open the doors.


\begin{itemize}
\item The locations marked 'X' on the map indicate the approximate position where a character can hold a shield to bounce the beams, which then reflects at a ninety-degree angle towards another glyph in line with the previous glyph (think "connect the dots"). Characters can hold the shields themselves or prop them up with objects if needed.
\end{itemize}

\subparagraph{Aperture Doors}
These stone doors are eighteen inches thick and made of heavy granite. There are no seams, hinges, handles, or locks upon these doors save an embedded bronze aperture set with a thick glass lens at the centre. Peering through provides a distorted view of the chamber beyond. Engravings on each door depict resolute paladins and righteous heroes of the Sacred Flame.


\begin{itemize}
\item The magically warded doors are immune to all damage, and magic can't alter their forms.
\item A door slides open once a Heavenly Beam (see below) is shone into the aperture, and remains open for eight hours thereafter.
\end{itemize}

\subparagraph{Heavenly Beams}
Light is collected above the cathedral dome and shines down upon the Ascendant Archangels, where it becomes a clearly visible magical beam of heavenly light. These beams do not harm creatures or damage objects, and travel in straight lines. They can only be reflected off mirrored shields found in the crypts or those borne by the Ascendant Archangels. Any other reflective surface causes the beams to fizzle out.

\subparagraph{Mirrored Shields}
Six silver shields with brilliant reflective surfaces are hidden throughout the catacombs. Each weighs twenty pounds and is entirely impractical for actual defense, and don't have any other properties beyond reflecting the heavenly beams. Either side of the shield can reflect the heavenly beams, obeying the laws of reflection, although solving the puzzles in the crypt only requires the beams be bounced at hard ninety-degree angles.

\subparagraph{Holy Wards}
\textit{Hallow} spells ward the House von Kessel vault, the tomb of Saint Vitruvio, the Vault of Ignacious, and Argonath's Rest. Creatures can't teleport into these areas.

\includegraphics[width=0.4\textwidth]{images/../5e.tools/img/adventure/DoDk/104-map-6.08-cathedral-catacombs.png}
\includegraphics[width=0.4\textwidth]{images/../5e.tools/img/adventure/DoDk/105-map-6.08-cathedral-catacombs-player.png}
\subsubsection{Alcoves of the Ascendant Archangels}
Descending either set of stairs from the cathedral sanctuary above leads to similar passageways:

\begin{DndReadAloud}{}
The stairs descend into a short corridor with a semicircular alcove at one end. The opposite way leads to a crossed intersection of arched passageways heading in each cardinal direction. Inside the alcove is a magnificent ten-foot-tall marble statue of an angelic figure set upon a pedestal. The angel holds a weapon in one hand, and in the other bears a round silver shield polished to a mirror-like sheen. A column of light shines down from a hole in the ceiling above, illuminating the sculpture in a beam of heavenly light.

\end{DndReadAloud}


\begin{itemize}
\item The statues depict the Archangels Gabrielle (east alcove) and Michael (west), who hold a spear and sword respectively. Characters proficient in the Religion skill automatically recognize them; they are the angels who guided Saint Tarna down the path of a paladin.
\item The shields are magically secured to each statue. However, the angels' arms are articulated. A character can use an action to move and rotate the shield into a forty five degree angle in relation to the floor, which sends the beam of heavenly light straight down the passageway directly away from the alcoves.
\item Characters can place a mirrored shield in the intersection to bounce light north or south.
\item An inscription on each statue reads "The light opens the way of truth."
\end{itemize}

\subsubsection{[1-3] Tombs}
\begin{DndReadAloud}{}
These long rooms are lined with ancient stone slabs laid out with the mummified remains of clerics and paladins of the Sacred Flame, the bodies simply adorned with a single trinket or weapons from their former lives. Each clutches a holy symbol in their hands, and has smooth stones placed upon their eyes. Surrounding each stone platform are dozens of candles in various shapes and sizes. The tombs here are untouched and covered in thick dust. The stone slab floor and low ceilings paired with the many tombs and niches lined with skulls ever watching those who enter make even these large and spacious rooms feel claustrophobic and unsettling.

\end{DndReadAloud}

\subparagraph{Mummified Clergy}
Each chamber houses the mummified remains of clerics and paladins of the Sacred Flame. Although the dead are normally cremated, the faithful believe these blessed souls may one day return to their earthly bodies in a desperate time when the Sacred Flame enact its will through corporeal vessels. Their bodies are wrapped in linen cloth, and dressed with bits of clothing, armour, or weapons they used in life. Inscriptions, tapestries, and painted murals decorate each sarcophagus niche, telling the life and deeds of the dead. Finally, 10 (3d6) candles are arranged around each sarcophagus. A holy rite allows the candles to remain lit eternally without being consumed, but they can be doused as normal.


\begin{itemize}
\item Each is a \textbf{mummy} that lies in state. Last rites performed over them decades ago hold back the reanimating magic of the Haze. They do not rise unless their remains or possessions are disturbed, in which case the nearby candles snuff out, and they attack as unhinged and violent undead. They have the following additional trait:
\item \textbf{\textit{Rejuvenation}.} If the mummy is destroyed, it regains all its hit points in ten minutes, unless its remains are placed upon the stone slab, surrounded with lit candles, sprinkled with holy water, and a short prayer uttered over it, or a \textit{gentle repose} or \textit{remove curse} is cast upon its remains.
\item \textbf{\textit{Treasure}.} Each mummy has gemstones placed over their eye sockets worth 100 gp each. They also clutch a gold or silver holy symbol of the Sacred Flame worth 50 gp.
\end{itemize}

\subsubsection{[1] Western Tombs}
The candles of this great hall are lit with ever-burning flames that cause a dance of dim light across the various tombs and niches throughout. Shadows cast upon the walls and floor flitter with an ambience of warmth and welcome.


\begin{itemize}
\item The door to this chamber opens immediately when the arms of Michael in the west Ascendant Archangel alcove are switched to a forty-five-degree angle.
\item The mummies here remain dormant unless disturbed.
\item One of the niches is empty, but there is a small circular aperture filled with glass on the wall. Bouncing light into it from a mirrored shield causes the wall to slide open into a dusty room, which contains 1d4 rare Spell Scroll from the cleric spell list and another mirrored shield.
\end{itemize}

\subsubsection{[2] Southwest Tombs}
The candles of this hall have all gone out, casting the room in a cold darkness. The deafening silence of the hall of dead is eerie. Upon one of the ancient tombs lies a shining shield that just catches the light drifting down from the stairs. The shield is clutched in the hands of the mummified remains of a great saint.


\begin{itemize}
\item Opening the aperture door to this chamber requires using one mirror shield; bouncing light from the west Ascendant Archangel alcove, placed at the west intersection.
\item All the candles in this tomb have gone out, and the mummies interred here are mad. They rise and attack if any enter the chamber, and the mirrored shield clatters to the ground.
\item The mummies must be laid to rest properly to keep them from rising again (see above).
\end{itemize}

\subsubsection{[3] Eastern Tombs}
Several of the candles in this room have gone out except for three of the stone slabs still illuminated in the dim light. At the back of the room a statue of a great Flamekeeper sits clutching a shining shield that reflects the candlelight of the room.

\includegraphics[width=0.4\textwidth]{images/../5e.tools/img/adventure/DoDk/106-06-019.crystal.png}

\begin{itemize}
\item The door to this chamber opens immediately when the arms of Gabrielle in the east Ascendant Archangel alcove are switched to a forty-five-degree angle.
\item The mummies rise and attack if someone takes the mirrored shield, except the three with candles still lit around them who remain in state.
\end{itemize}

\subsubsection{[4] Founder's Tomb}
\begin{DndReadAloud}{}
This semicircular room is lined with rows upon rows of niches each with a single skull looking out into the centre of the room. A burning brazier rests in the middle that casts a soft glow highlighting the semi-precious stones or coins placed in the eye sockets of each of the skulls. Along the flat wall of the chamber are several rows of shelves that lie between two stone-carved archways that frame a portion of stone wall with nothing on it. The shelves are lined with small statues of various animals including a brass toad, a copper cat, a silver bird, an ivory ox, a granit bear, a gold horse, and an iron eagle.

The brazier within this chamber has an inscription carved into the stone around its base that reads "The bodies of our first Flamekeepers rest here under the watchful gaze of the dead. Ere a living soul should seek their guiding light, know that one who sees with clear eyes may open the way. Take up their shining truth with silver wings."

\end{DndReadAloud}


\begin{itemize}
\item There is one skull that has no gems placed in its eyes. A small switch within the left eye socket causes the wall of the right archway to slide open.
\item A switch on the bottom of the silver bird statue causes the left archway to open.
\item Within each open alcove is a stone slab with a gloriously decorated body adorned with holy symbols and burning candles. Each clutches a shimmering mirror shield.
\item The aperture door in this chamber can be opened using a single shield and bouncing the light from the east Ascendant Archangel alcove at the east intersection.
\end{itemize}

\subsubsection{[5] Pool of the Mourning Angel}
This circular room has a domed roof painted with frescoes of knights carrying holy flames upon their hands and marching in a spiraling formation towards a central glowing light. A large marble font of water rests within the centre of this room with a tragic looking statue perched on the far end. Water flows from the statue's eyes, down its face, falling gently into the pool below.

This statue depicts an angelic man looking up to the heavens with a look of sorrowful anguish upon his crying face. His hands rest an inch above the water, fingers curled up as if he begs "why!" to the heavens. Along the base of the fountain are four evenly spaced iron handles.


\begin{itemize}
\item A mirror shield can be placed upon the statue's hands, which can be rotated by grabbing the iron handles under the water. The light can then be reflected into the Scriptorium.
\end{itemize}

\subsubsection{[6] Scriptorium}
\begin{DndReadAloud}{}
The scriptorium is a large square room lined with wooden shelves from floor to ceiling. The shelves are overflowing with assorted texts, papers, scrolls, and documents all loosely organized by date. In the middle of the room are four large wooden tables flanked by rows of wooden stools. Several old tomes and scrolls lay upon the tables along with inkwells and quills. Ink stains the tables and one of the inkwells has toppled over causing the ink to spill to the floor and run through the cracks within the stone.

\end{DndReadAloud}


\begin{itemize}
\item The documents all appear to be writings of holy texts, copies of famous literary works and historical stories of importance to the Sacred Flame.
\item Amongst the shelves are 2d6 rare Spell Scroll from the cleric spell list, including a \textit{scroll} of \textit{raise dead}.
\item Floating within the room is a \textbf{wraith} and 2 (1d4) warp witches.
\item Opening the Vault of Ignacious requires all six mirror shields to direct both heavenly beams into the statue. Three are placed on each side.
\end{itemize}

\subsubsection{[7] House von Kessel Vault}
\begin{DndReadAloud}{}
Two long alcoves meet a small passage at the intersection of this L-shaped chamber. At the ends are large stone statues of armoured guardians each holding a kite shield before them. A large stone plaque is polished to a mirror sheen between them, opposing a narrow passage that ends in a bas-relief carving of two dragons rampant, each with an aperture lens in their mouth.

\end{DndReadAloud}


\begin{itemize}
\item Opening the vault requires bouncing light from the east intersection, then bouncing the light \textit{upwards} at the bottom of the shaft, then bouncing it east again at the top of the shaft. Since the shield would need to be positioned in midair, it's up to the characters to come up with a clever way to do this - but they'll need three mirror shields however they do it.
\item \textbf{\textit{Lineage Records}.} This family tree with written descriptions of each person, marriage agreements, dowries, and adoption decrees is magically preserved and bears royal seals.
\item \textbf{\textit{The Royal Will}.} What it says is up to you, but it could describe Ulrich IV's intentions if he were to die without a known heir.
\item \textbf{\textit{Blood Phylacteries}.} This minor magical item is a glass vial that keeps the blood stored within preserved indefinitely. Once sealed, the stopper holds celestial runes naming the person whose blood is contained within, and can only be opened by the use of \textit{universal solvent}. If a creature who is a direct blood relative holds the vial, the vial glows and runes appear within describing the nature of their relation (parent and child, siblings, cousins, ect.).
\item Vials of blood taken from Ulrich IV and his last six predecessors are kept here.
\end{itemize}

\DndQuote%
{}
{''Blessed is the lair of Ignacius! A wonderful gift for an impressive warrior.''}
{Pluto Jackson}
\subsubsection{[8] Vault of Ignacious}
\begin{DndReadAloud}{}
This circular vaulted chamber is ribbed with strange bronze pipes and niches filled with candles. Hovering in place above a central dias in a column of flame is a glimmering silver blade. The chamber is otherwise unadorned, but filled with nearly blinding light and sweltering heat from the crackling flames.

\end{DndReadAloud}


\begin{itemize}
\item This sword is \textit{Ignacious, the Sword of Burning Truth}. Whoever grips this blade is judged by the sword, which only accepts a righteous wielder.
\begin{DndReadAloud}{}
Those who take up the blade and are unworthy: "Thou art rough in spirit and brittle in soul, knave. Though perhaps if you prove to me your might, I shall take your coal-heart and iron limbs and temper them into a righteous instrument of blazing steel!"

\end{DndReadAloud}

\item The column of flame then collapses into a burning \textbf{deva} wielding the blade, who challenges the characters to battle. Ignacious will allow itself to be wielded by any who defeat the deva, believing it can turn them to the cause of righteousness.
\end{itemize}

\subsubsection{[9] Tomb of Saint Vitruvio}
\begin{DndReadAloud}{}
A marble colonnade forms a narrow aisle around this circular chamber, each pillar set with a lighted torch. Murals upon the walls depict the deeds and battles of a valiant paladin, and a life-sized statue of the man stands upon a circular podium in the centre of the room.

He is clad in plate armour and wears a resplendent helm; a shield upon his pack, sword strapped to his belt, a sceptre at his side, and a phylactery around his neck. His outstretched arms extend fully to his sides, and his palms face outwards. Cavities in the statue at the right thigh, jaw, and within several knuckles hold the few remaining bones of Saint Vitruvio, each etched with delicate prayers. Glass spheres rest in the statue's eye sockets, and within the palms of the statue's hands are stigmata-like holes, each filled with a small lens. Hundreds of lit candles are set around the statue and about the room, and a mirrored shield is fixed on the wall to the south.

\end{DndReadAloud}


\begin{itemize}
\item Accessing Saint Vitruvio's Tomb requires using two mirror shields to bounce light at the intersections, but can be done from either the east or west side.
\item Upon the wall is a golden shield. Behind the shield is a large circular aperture filled with glass. Surrounding it is a moulding depicting Saint Vitruvio riding upon Argonath.
\item Vitruvio's statue can be rotated, but also lowered and raised by pushing on it as an action. It can be moved into a position such that its palms are in line with the apertures on the doors, and its eyes are in line with the mirrored shield on the south wall.
\item If the heavenly beams of light from both Ascendant Archangel alcoves are reflected into the lenses on both of the statue's hands at the same time, a single thick beam shoots from the statue's eyes. The resulting beam strikes the aperture on the south wall, which then rumbles open revealing a wide spiral staircase leading to Argonath's Rest. Accomplishing this requires four mirror shields, one positioned at each intersection.
\end{itemize}

\subsubsection{[10] Argonath's Rest}
\begin{DndReadAloud}{}
The spiral staircase descends down a wide stone chimney into the earth, and eventually opens into a great natural cavern. Suffocating sulphurous-smelling vapours saturate the chamber. Resting in a pool of bubbling oil are the ancient remains of a colossal dragon. The great golden plates of its armoured scales still cling to its hulking skeletal remains, though the rest of its flesh has long since decomposed.

\end{DndReadAloud}


\begin{itemize}
\item A \textit{forbiddance} spell wards this chamber; creatures can't teleport into the area.
\item An inscription reads: \textit{"This was my brother in battle, Argonath. Let none disturb his deserved rest unless great darkness engulfs Drakkenheim at the end of days."}
\item Carrying an open flame in this chamber is a \textit{very} bad idea.
\item These are the bones of Argonath, an \textbf{ancient gold dragon}. The dragon can be restored to life by casting \textit{resurrection} using a spell slot and providing the \textit{Diamond of Mount Kadath} or the \textit{pristine delerium geode} as a material component. Thereafter, the dragon slumbers here unless commanded by one who wields all the relics of Saint Vitruvio.
\item The chimney by the staircase meets the cathedral floor above. If Argonath is resurrected, he exits this chamber by flying up the chimney and bursting through the floor.
\end{itemize}

\begin{DndSidebar}[float=!b]{Developments}
Slaying the Lord of the Feast and conquering the cathedral provokes an immediate response from all the factions.



\subsection{Council at the Cathedral}
Shortly after the cathedral is retaken, the factions make overtures under the banner of armistice to hold a meeting there. Each faction leader (and their lieutenants) is present.


\begin{itemize}
\item Eldrick Runeweaver sends his \textit{simulacrum}, his \textbf{homunculus}, and River. They bring a \textit{potion of invisibility} and a Spell Scroll of \textit{dimension door}, and escape immediately if conflict breaks out.
\item The Queen of Thieves sends a decoy, controlled via \textit{dominate person}. This patsy is disguised both physically \textit{and} using a \textit{seeming} spell, and the Queen of Thieves directly controls this unfortunate captive and makes them behave as if they were actually her in person. She coordinates with her spies and assassins to launch an assassination strike against a vulnerable faction leader or the player characters. Blackjack Mel does not attend, however.
\item Elias Drexel, Theodore Marshal, and Lucretia Mathias arrive with a large contingent of their own faction troops nearby, ready to storm the cathedral if violence breaks out.
\end{itemize}

All the faction leaders are keenly aware that violence could break out at any moment, but in fact, none of them are actually willing to \textit{start} a fight here. During the discussion, faction leaders should bring up the following key questions:


\begin{itemize}
\item What are the true motives of the player characters? Whose side are they on, really?
\item Which faction should control the cathedral?
\item What will be done about the royal documents?
\item Who will hold any \textit{Seals of Drakkenheim} discovered so far? Who should rule Drakkenheim?
\item How should delerium be used?
\end{itemize}

During this meeting, the roleplaying traits of the factions, their goals, and the faction leaders should be fully presented. Draw bright lines in the conflict, and don't be afraid to put the characters on the spot: each faction leader should address each player character directly at least once with a difficult question of some kind.

\subparagraph{Brutal Betrayal (Optional)}
If the player characters have greatly antagonized one of the factions before this meeting, that faction either secretly mobilizes several elite strike teams on standby or plants powerful explosives under the cathedral prior to the meeting. They also orchestrate an escape plan for their leader (which in the case of the Academy or the Queen of Thieves isn't needed since they are only present by proxy). During the meeting, the leader demands to be given either all of the Seals of Drakkenheim, the Relics of Saint Vitruvio, control of the cathedral, or any other insane ultimatum appropriate for the faction. If their demands aren't met, their forces attack or they detonate the explosives, which create a massive explosion akin to a \textit{meteor swarm} that fills the cathedral and collapses the roof and dome.



\subsection{Reigniting the Brazier}
Lucretia Mathias and Ophelia Reed both offer to reignite the \textit{Brazier of the Sacred Flame} and reconsecrate the cathedral. This requires some ashes or bone fragments from a saint: either those of Saint Vitruvio himself or the gold dragon, Argonath, which can be recovered from the cathedral catacombs. If re-lit, the cathedral brazier fills the area with effects similar to \textit{hallow} and \textit{forbiddance}, and becomes a sanctuary in the city against the Haze where characters may rest normally.

If the cathedral was bombed, the brazier is destroyed and can't be reignited.



\subsection{Aftermath}
The cathedral remains hotly contested ground: the Followers of the Falling Fire and the Silver Order will try to take it from each other after these meetings regardless of how the negotiations turned out. If the Hooded Lanterns or the Queen of Thieves aren't able to gain the information in the House von Kessel vault, they'll launch counterstrikes as well to obtain them.

Any faction who holds onto the cathedral can establish a shaky safe haven there, but doing so stretches the faction's forces thin. Nevertheless, the player characters can use it as a base of operations. Ultimately, holding the cathedral for more than 14 (4d6) weeks proves untenable due to constant attrition, unless the player characters take action to drive out enemy factions and complete their allies' major objectives. By taking the cathedral the player characters have a major opportunity, but they must take advantage of it before it's too late.



\end{DndSidebar}

\section{Slaughterstone Square}
Slaughterstone Square is an important crossroads in Drakkenheim - the intersecting streets here lead north to Castle Drakken, east to King's Gate, south to a bridge crossing the Drann River, and west to Market Square Plaza. Centuries ago the square was the core of the Old Town, home to the main marketplace and simply known as Flagstone Square. As the city grew it became a civic centre and eventually, a place for public executions. Its present moniker has never been more appropriate. The \textbf{Executioner} still occupies the plaza: a hulking clockwork knight designed to mete out death sentences. Contaminated by the Haze, the beserk machine now ferociously eviscerates any who dare pass through Slaughterstone Square.

\subsection{Overview}
Slaughterstone Square lies above the cistern, so we recommend you read this section first if your characters are seeking that location.

\subsubsection{Adventure Hooks}
Rumour claims a vast collection of treasure and magic items may be found here, left behind by the Executioner's countless victims. Many adventurers make the fatal last mistake of approaching Slaughterstone Square for this reason alone. However, your characters may simply happen upon this forsaken place by accident or on a foolish lark. Alternatively, some bold characters may be travelling through the square to reach the cistern, or possibly arrived via the smuggler's passage connected to the Black Ivory Inn.

\subsubsection{A Deadly Threat}
No matter why they came here, your characters should probably run. The Executioner is not a foe to be defeated, but rather an obstacle to be avoided (or perhaps an asset to use should the opportunity emerge to lure an unsuspecting foe here!). It is a reminder of the callous and uncontrollable dangers present in Drakkenheim. Embrace clever solutions your characters invent to escape or bypass this threat, but should your players engage the Executioner directly in combat, \textit{do not pull any punches}. They have made an incredibly foolish mistake by confronting a well-telegraphed deathtrap, and should be served their just desserts!

\DndQuote%
{}
{''When planning a trip to Slaughterstone Square, always remember: you dont have to outrun the Executioner, just your slowest party member...''}
{Veo Sjena}
\subsection{Area Details}
\subsubsection{Slaughterstone Square}
\begin{DndReadAloud}{}
Slaughterstone Square is a large octagonal plaza paved with broad flagstones with a squat, bowl-shaped granite fountain thirty feet wide and five feet tall at its centre. The tightly-packed surrounding townhouses rise up three storeys to high-steepled rooftops with wrought-iron cresting. Many feature narrow balconies and tall dour-looking windows on their front facades overlooking the square.

The building composition resembles stern faces held in silent judgement, as if all who look upon them are found wanting. Pools of coagulated blood stain the streets around several mangled corpses.

\end{DndReadAloud}


\begin{itemize}
\item The \textbf{Executioner} stands utterly motionless beside the fountain, blood dripping from its axe.
\item It jolts to life and immediately charges any creature who attacks, casts a spell, or approaches within one hundred twenty feet of it. It attacks indiscriminately, and only pursues creatures beyond the square if it continues to be attacked.
\item The Executioner has formed a bizarre habit for gathering the bodies of its slain victims and hurling their remains into the fountain. Blood and body parts slosh around the fountain, and then are deposited into the cistern.
\item The fountain collects rainwater in a flowing spiral that drains into a five-foot diameter hole at the centre. The hole leads directly into the cistern below, where the water fills up the reservoir. There was once a wrought-iron grate covering the hole, but only a few protruding pieces of sharp iron are left there now.
\item Amongst the few corpses the Executioner hasn't yet dropped into the square may be found damaged adventuring equipment and weapons. If you wish, 1d6 uncommon \textit{potions} and an uncommon magic item of your choice may be found here.
\item A heavy sewer cover beside the fountain can be lifted from this side with an action, proper leverage and a DC 20 Strength check. It opens to a stairway into the Old Town Cistern.
\end{itemize}

\includegraphics[width=0.4\textwidth]{images/../5e.tools/img/adventure/DoDk/107-06-020.slaughterstone-square.png}
\subsubsection{Magistrate's Court}
\begin{DndReadAloud}{}
The largest building in Slaughterstone Square is a somber and formal-looking manor house three storeys tall that takes up the entire block. Pointed windows with broken glass line each floor. Broad stone steps lead up to the double doors entrance, above which hang banners bearing the crest of the city itself.

\end{DndReadAloud}


\begin{itemize}
\item This building served as both a courthouse and offices for the civilian functionaries such as the Lord Mayor and the city council. Many public records are housed in the archives here, so characters on a personal quest looking for information about people who lived in Drakkenheim could find such documentation here.
\end{itemize}

\chapter{Faction Strongholds}
\DndDropCapLine{D}{}uring their adventures in Drakkenheim, characters may explore these five mighty faction strongholds in a variety of ways. Early in the campaign, they may interact on friendly terms with faction members at these locations, but once characters reach 9th level or higher they might consider hostile missions into enemy faction strongholds.


\includegraphics[width=0.4\textwidth]{images/../5e.tools/img/adventure/DoDk/108-07-001.camp-dawn.png}
\section{Interaction}
Most commonly, player characters can visit a faction stronghold to meet with faction members and speak with the faction leader. Characters who are cooperatively working with a friendly faction may be invited to use the stronghold as a base of operations.

\subsection{Infiltration}
Characters might also attempt sneaking into a faction stronghold for a specific purpose, such as rescuing a prisoner, sabotaging resources or defenses, spying on the faction, stealing an important item, or even assassinating the faction leader. Every faction expects this to happen and has taken reasonable and intelligent precautions that should provide exciting challenges for stealthy characters. Factions that acquire key items in the campaign, such as the \textit{Relics of Saint Vitruvio} or the \textit{Seals of Drakkenheim} store them in secure locations usually guarded personally by the faction leader, set up extra patrols, and may ask faction spellcasters to use protective spells such as \textit{glyph of warding} to keep the items safe.

\subsection{Assault}
At some point in the campaign, the player characters may lead an invasion on an enemy stronghold in an all-out attack to drive an enemy faction out of Drakkenheim. The stakes should be appropriately high for such an exciting and decisive moment; a stronghold assault is a win-or-die play. You should give the players the opportunity to take the lead in coming up with a battle plan with their allies, and put the full resources (and tricks) of each faction front-and-centre.

However, given the troops involved, most factions simply cannot hide when they muster forces for an all-out attack on another faction stronghold unless specific and extensive action is taken to conceal it. This gives the besieged faction time to prepare its defenses and take adequate precautions, so you should increase the numbers of NPCs present, place faction strike teams in key locations, and secure the chamber of the faction leader and the main entrances to each stronghold.

During an all-out assault, it's best to have the player characters face down elite faction strike teams, spearhead an attack on the main gates, launch a commando raid to sabotage an important defense, or have a dramatic showdown with a faction lieutenant or the faction leader, while their faction allies engage the rank-and-file to hold off reinforcements. The focus of any faction stronghold assault should be on heroic action, not a cumbersome mass combat scenario (unless your players \textit{really} enjoy such scenarios).

\subsection{Adventure}
The Inscrutable Tower and the Court of Thieves are two notable exceptions: both behave more like a typical adventure site and are presented accordingly for exploration.

\section{Camp Dawn}
A few miles southwest of Drakkenheim is a recently constructed fortified encampment. Just outside the furthest edge of the Haze, here the Knights of the Silver Order have established their base of operations for their sacred mission in Drakkenheim. \textbf{Knight-Captain Theodore Marshal} has set up his command post, and runs the camp with strong military discipline grounded in fiery faith.

\subsection{Interactions}
Camp Dawn can accommodate five hundred soldiers, support staff, and stockpiled supplies, as well as stabling hundreds of draft animals, warhorses, and flying mounts. Not everyone here is a warrior, however: many are labourers, cooks, and clergy that assist the fighting force and run the day-to-day tasks around camp. Nevertheless, dozens of Silver Order strike teams are resting, training, or resupplying here at any one time, in addition to the posted guards in the areas described below, representing two-thirds of the Silver Order's total forces in Drakkenheim.

The encampment operates under the same military standard one would expect from a front-line fortress on the eve of battle against a dangerous and duplicitous foe. The knights maintain a detailed roster of their personnel and equipment, with many checks and balances in place to follow up on missing patrols or misplaced gear.

\subsubsection{Visiting the Camp}
Any visitors approaching the gates are required to state their business. Those granted entry must disarm, and are escorted by a Silver Order \textbf{knight} or several guards wherever they go within the camp. The Silver Order is hospitable to those who are wounded or need help, but always treats any guests with caution. Allies of the Silver Order are recognized by the gate guards and may freely move about the camp, and may even be offered their own tent. Those who wish to speak with the Knight-Captain are escorted to the Command Post.

\subsubsection{Defences}
The encampment is in a perpetual state of readiness, and the rank-and-file respond to any attack, intrusion, or infiltration by quickly banding together and raising a general alarm. Every single guard on duty carries a horn, and specific calls are used to alert the camp upon the discovery of a fire, sabotage, slain guards, intruders, attackers, or all clear. These can be heard throughout the camp, and 1d4 knights rush to the site in short order, with many more to follow.

\DndQuote%
{}
{''An impressive armada designed to invade a city. I do not believe they know what they're up against.''}
{Pluto Jackson}
\subsection{Area Details}
\begin{DndReadAloud}{}
Camp Dawn has a tall wooden palisade wall surrounded with a ditch filled with water and sharpened stakes. At each corner is anchored a defended guard tower. Entry into the camp is by four gates—one at each cardinal direction. Each gate is flanked by towers, with a wooden bridge leading over the ditch to the wooden doors. Inside, the camp is an array of activities. Most of the enclosure is filled with neatly organized rows of large canvas tents, grouped around small cooking fires and supply carts.

\end{DndReadAloud}


\begin{itemize}
\item The main walls of the camp are constructed by double-layered wooden logs. Each log is twelve to fourteen feet long and ends in a spike. Stone and gravel have been piled along the base of the wall for reinforcement.
\item A ditch fifteen feet deep and ten feet wide has been dug around the entire perimeter like a moat. The trench is filled with sharpened wooden spikes and knee deep water. There is a sturdy wooden bridge over the trench leading to the camp gates.
\item Along the interior of the walls is a rampart made of wooden scaffolding. 1d4 guards patrol between each tower, each carrying heavy crossbows.
\end{itemize}

\includegraphics[width=0.4\textwidth]{images/../5e.tools/img/adventure/DoDk/109-map-7.01-camp-dawn.png}
\includegraphics[width=0.4\textwidth]{images/../5e.tools/img/adventure/DoDk/110-map-7.01-camp-dawn-player.png}
\subsubsection{[1] Supply Wagons}
There is a constant flow of wagons entering and leaving the camp. More arrive each day to keep the encampment freshly stocked with food, water, weapons, equipment, and other necessities. The ministry of the Sacred Flame funds the Knights of the Silver Order, so little gold or currency travels directly through the camp.


\begin{itemize}
\item Wagons are usually protected by a \textbf{knight} or 2d4 guards, and driven by a \textbf{commoner}.
\end{itemize}

\subsubsection{[2] Guard Towers}
Two towers flank each entrance into the camp, and one reinforces each corner. The square structures are built from thick wooden logs with ladders ascending from the ground. The defended platform rises about thirty feet off the ground with wooden stanchions and a pointed wooden rooftop.


\begin{itemize}
\item Each tower has 2 knights and 2d4 guards equipped with horns and crossbows.
\end{itemize}

\subsubsection{[3] Rookery Towers}
A round wood and brick tower sits in the southwest of the camp. Large arched windows circle around its exterior, and a wooden platform perches out from each. Ropes and ladders provide access to the tower. This tower has several nests for the griffon mounts of the Silver Order.


\begin{itemize}
\item Two dozen adult griffons are found here at any given time.
\item A \textbf{cavalier} and 1d4 knights stand guard.
\end{itemize}

\subsubsection{[4] Stables}
A long wooden building has been constructed along one of the walls of the fort. Two barn doors on each end open to a central aisle with rows of fenced stalls on each side, each filled with hay and a water trough.


\begin{itemize}
\item About 40 warhorses are housed here at any given time.
\item A staff of a dozen stablehands (human commoners) tends the horses, and 2 knights and 2d4 guards patrol nearby.
\end{itemize}

\subsubsection{[5] Infirmary}
A long short square wooden building in the northwest of the camp acts as a hospice for injured knights. The building has a warm fire inside and many cots and tables. 2d6 acolytes and a \textbf{chaplain} are taking care of the injured. At any given time, 1d6 knights are resting here recovering from contamination and other injuries.


\begin{itemize}
\item A small stockpile of 10 (3d6) Potion of Healing are kept here, along with 1d4 scrolls of \textit{lesser restoration}, \textit{prayer of healing}, \textit{purge contamination}, \textit{greater restoration} each.
\end{itemize}

\subsubsection{[6] Pyre of the Sacred Flame}
A large pyre burns ever brightly at the heart of Camp Dawn. This consecrated bonfire is lit by a \textit{continual flame} and emits a \textit{hallow} spell effect. The flame is surrounded by several wooden benches and prayer mats. This is not a place for meals or cooking. It is a place of worship; a shrine constructed in the field so the faithful knights may contemplate their duty to the Sacred Flame.


\begin{itemize}
\item Each morning, High Flamekeeper Ophelia Reed leads prayers and hymns around the pyre, a ceremony that usually lasts about an hour. Everyone not on duty attends these gatherings. Theodore Marshal is often in attendance.
\item During the rest of the day, another \textbf{priest} and 2d4 acolytes tend the pyre, and often others in the camp will stop here to make their own personal prayers.
\end{itemize}

\subsubsection{[7] Knight's Tents}
The northern half of the camp comprises eighty tents in neatly organized rows. Each tent has four bunks for sleeping staff or soldiers. Each keeps their belongings in satchels or small footlockers within. Most are occupied during the night, but empty during the day. Each cluster of tents keeps their own stoves or cooking fires.

Amongst all the tents, there are all manner of martial weapons, ammunition, suits of medium armour, and shields. Each tent typically sleeps one \textbf{knight} and the guards of their retinue.

\subsubsection{[8] Flamekeeper's Tents}
These larger tents are set up next to the infirmary. Each is decorated with silver trim and is equipped with one to four cots, a small stove, a desk with a chair, and a chest for belongings.


\begin{itemize}
\item These tents are used by the clergy of the Sacred Flame who reside in the camp: chaplains and priests, acolytes, and High Flamekeeper Ophelia Reed.
\item 1d4 knights are posted as guards for the Flamekeepers at all times.
\end{itemize}

\subsubsection{[9] Command Post}
The command post is a large open tent that sits just off the pyre. Inside is a large meeting table that can sit most officers in the regiment, with a map of Drakkenheim unrolled upon it. Several wooden figures are set on the map showing plans and tactics for military operations.

\includegraphics[width=0.4\textwidth]{images/../5e.tools/img/adventure/DoDk/111-07-002.pyre-priest.png}

\begin{itemize}
\item 4 cavaliers stand guard at the command post at all times, but typically the tent is abuzz with activity. Theodore Marshal is usually here during the day, conferring with several chaplains, priests, knights, Silver Order strike teams, or Ophelia Reed.
\end{itemize}

\subparagraph{Knight-Captain's Tent}
This large tent in the southeast of the camp is decorated with lavish designs and silver filigree. The symbol of the Sacred Flame hangs on a banner next to the entrance. Inside are animal fur rugs, a large wardrobe, a locked chest, a simple bed, and a seldom-used writing desk.


\begin{itemize}
\item This is the private tent of Theodore Marshal. 2 cavaliers are assigned to guard the Knight-Captain at all times.
\end{itemize}

\subsubsection{Stockpile}
About five hundred feet from the main camp is a wooden shed that has been constructed in a clear field. Inside the shed are several barrels and crates filled with Alchemist's Fire (flask) and oil. The shed is reinforced and locked, and clearly visible from the guard towers.


\begin{itemize}
\item The shed is manned by 1d4 knights and 2d6 guards patrol nearby at all times.
\end{itemize}

\DndQuote%
{}
{''Their whole camp is made of wood, and filled with flammable components - some even strap this stuff to their back. What's the worst that could happen?''}
{Sebastian Crowe}
\section{Court of Thieves}
Hidden deep under the city lies the stronghold of the Queen's Men, the Court of Thieves. Occasionally referred to as the Undertavern, it is the personal lair of the Queen of Thieves. While it could be near Buckledown Row, the placement of the Court of Thieves is up to you.

\includegraphics[width=0.4\textwidth]{images/../5e.tools/img/adventure/DoDk/112-07-003.court-of-thieves.png}
\subsection{Interactions}
Finding the Court of Thieves is a challenge all its own. Its location is a closely-kept secret only known to a small number of Queen's Men agents on a need-to-know basis: even Blackjack Mel himself doesn't \textit{really} know where it's found.

\DndQuote%
{}
{''I used to know my way around, but the streets have changed. Is that even possible?''}
{Sebastian Crowe}
\subsubsection{Visiting the Court of Thieves}
Named members of elite Queen's Men strike teams travel here occasionally. As such, following, interrogating, or tricking the Queen of Thieves' own strike teams is the best method for locating the court. Important prisoners and treasures seized by the Queen's Men agents eventually make their way back here. The player characters might be able to engineer an opportunity to lure out a strike team to track back to the court by using tantalizing bait.

Regardless, those who come to parley with the Queen of Thieves are presented with a fiendish puzzle before they are granted an audience (see the "Queen's Riddle" below).

\subparagraph{Queen's Champions}
A widely-circulating rumour claims that the Queen of Thieves grants a personal favour to any who defeat the arena champion, and so characters who made a name for themselves fighting in the rings of Buckledown Row might be invited to the Court of Thieves to test their mettle in the arena instead.

\subsubsection{Defences}
The Court of Thieves focuses on security through obscurity. A Queen's Men strike team is usually on guard here. Devilish traps protect the inner sanctum of the Queen of Thieves.

\subsection{Area Details}
\includegraphics[width=0.4\textwidth]{images/../5e.tools/img/adventure/DoDk/113-map-7.02-theives-throne-room.png}
\includegraphics[width=0.4\textwidth]{images/../5e.tools/img/adventure/DoDk/114-map-7.02-theives-throne-room-player.png}
\subsubsection{Aquifer Junction}
This circular aqueduct passage is inhabited by three giant crocodiles. They've been lured here deliberately to ward off snoops or intruders. They don't answer to any of the Queen's Men, but any of the faction members who regularly come here know to give them a wide berth.

The walkway along the sewer canal here is heavily damaged and has partially collapsed into the waterway. Huge blocks of stone and gravel are piled up, but the wall itself is completely intact.


\begin{itemize}
\item A DC 15 Investigation check determines the walkway isn't collapsed: it's an intersection heading east blocked off with wreckage.
\item The wall here is an illusion created by a permanent \textit{major image} spell cast by the Queen of Thieves. Creatures and objects simply pass through to a wooden staircase in the Undertavern.
\item Excavation equipment or magic would be needed to remove the wreckage.
\end{itemize}

\subsubsection{Undertavern}
\begin{DndReadAloud}{}
This wide sewer passage has been transformed into a makeshift tavern. The main canal is drained, relatively dry, and blocked with debris in several places. Wooden planks hammered and lashed across the canal form bridges and decks above, accessible through ladder-like stairways. The broad upper sewer walkways have been cleared out to make room for about a dozen round tables with stools set around each. Hanging lanterns are strung up throughout to light the chamber, and a few coal stoves provide meager warmth. An impressive bar has been constructed, lined with barstools from which a curiously dapper mechanical barman serves ale and mead. At the far end of the tavern is a small band of ruffians performing lively music upon an improvised stage. Several rogues and thieves mingle about.

\end{DndReadAloud}


\begin{itemize}
\item The canal continues east to another intersection, which heads north into a fog-shrouded tunnel and further east to a collapsed tunnel.
\item One or two Queen's Men strike teams are watchfully relaxing here most of the time. They aren't hostile to visitors, even enemies of the Queen of Thieves. Instead, they invite characters to share a drink with them and chat. Characters who inquire about seeing the Queen of Thieves are told to speak with the Barman.
\item The Barman is a mechanical creature with the game statistics of a medium-sized \textbf{shield guardian}. He can serve and mix drinks with aplomb. He doesn't speak, but any mention of the Queen of Thieves prompts him to deliver a programmed response via the \textit{magic mouth} spell (see the "Queen's Riddle" below).
\end{itemize}

\subparagraph{Kitchen}
Several wooden panels have been set up here to create a separate kitchen. Inside are several well-stocked pantries, and a pig roasts over a fire pit. Anyone is free to use the kitchen, there's no designated cook.

\subparagraph{Secret Doors}
Seams in the walls are noticeable by a creature within five feet who makes a DC 20 Intelligence (Investigation) check. Knocking on the wall causes the door to slide open, revealing spiral staircases to the Arena Level or the Prison.

\subsubsection{Poisoned Passage}
\begin{DndReadAloud}{}
The canal exits the tavern into a ten-foot-wide arched tunnel shrouded by a menacing green vapour.

\end{DndReadAloud}


\begin{itemize}
\item The mist is magical, but whether or not it is harmless is impossible to determine (see below). If dispersed by wind, it reappears one round later.
\item NPCs in the bar explain that beyond lies the audience chamber of the Queen of Thieves. They claim the mist contains a powerful toxin that will render anyone who isn't inoculated against it completely unconscious and unresponsive.
\item Those who wish to see the Queen uninvited must solve the Queen's Riddle given by the Barman first. Any who solve the riddle will become inoculated against the toxic mist, and may pass through safely.
\item Queen's Men agents and strike teams in the bar have been "inoculated" already as they've proven their loyalty to the Queen of Thieves.
\item What the NPCs don't know is that the mist is actually harmless, and the potion they drank was merely a placebo. However, since this is what they \textit{believe to be true}, they aren't lying. They do not know the answer to the riddle or how the riddle works.
\end{itemize}

\subsubsection{The Queen's Riddle}
Recite the following carefully:

\begin{DndReadAloud}{}
The barman gestures to the mist-filled passageway and speaks:

"The Queen of Thieves bids any who seek an audience drink to her true name."

\end{DndReadAloud}

He then sets the following drinks upon the table before you:


\begin{itemize}
\item A bottle of red wine....
\item a cask of yellow mead...
\item a jug of white milk...
\item a pot of black coffee...
\item and several empty drinking cups.
\end{itemize}

\begin{DndReadAloud}{}
Then the barman speaks again:

\end{DndReadAloud}

\begin{DndReadAloud}{}
"It's a drink more pure than water, and more deadly than poison. When you finish your drink, it fills your cup, but you'll die if it's all you have to drink."

\end{DndReadAloud}

\subsubsection{Solution}
The entire riddle is a misdirection:


\begin{itemize}
\item The liquids are non-magical, but each contains an undetectable reactive toxin.
\item A creature who drinks any of the liquids then walks into the mist activates the toxin.
\item The toxin causes the creature to fall unconscious for 8 hours or until they take damage, with no saving throw.
\item Unconscious characters are stripped and locked in the prison cells below.
\end{itemize}

The answer to the riddle is "nothing", but speaking the answer prompts no response from the Barman. If someone drinks, he gestures towards the mist. The correct solution is to simply ignore all the drinks, and walk through the mist. It's not even necessary to drink from the empty cups.

\subsubsection{Thieves' Throne Room}
\begin{DndReadAloud}{}
This sewer junction has been blocked off, drained, and cleaned to create a throne room. Most pillars are repurposed pipes and a large sewer grate rests in a depression in the centre of the room. The brick walls still bear the stains and discolourations of sewer water, though the scent of perfume fills the air. Fine tapestries and benches with well-worn but choice scavenged pillows are positioned along the walls, and hanging braziers illuminate the room. At the far side of the chamber is a large circular alcove illuminated by lanterns, upon which is an impressive stone throne. A throng of chattering and laughing guests lounge about the room, including all manner of ne'er do wells, dashing rogues, and sly charmers.

\end{DndReadAloud}

\subparagraph{Courtiers}
Most people in this chamber aren't real, and are instead the result of many carefully designed \textit{programmed illusion} spells. This illusory model audience is perfectly programmed to play along with the Queen of Thieves' charades: cavorting like affable flunkies, gasping dramatically, uproariously laughing, sneering, barking cruel insults, clapping enthusiastically, and otherwise performing as needed.


\begin{itemize}
\item A Queen's Men strike team is positioned here whenever the Queen of Thieves holds an actual audience.
\end{itemize}

\subparagraph{Throne}
The Queen of Thieves is seldom here in person, and instead appears here via the \textit{project image} power of her resplendent desk in the office (Area X).

\subparagraph{Trap Grate}
This iron grate is an obvious trap. A button on the armrest of the throne and a hidden lever on one of the pillar pipes cause it to snap open suddenly. Any creature standing on it is deposited into the Arena Floor, where they land in a pool of deep water and suffer 4d6 damage from the fall.

\subparagraph{Secret Doors}
A depression in the wall here is noticeable with a DC 20 Intelligence (Investigation) check. Knocking on the wall causes the door to slide open to the guest rooms.

\subparagraph{Guest Rooms}
Offered to "guests" of the Queen of Thieves, these well-appointed bedchambers are in fact sufficiently insulated from the Haze thanks to their underground position. Depending on the nature of her company, the Queen of Thieves may place poisoned drinks and food here, or seal the doors with \textit{Arcane Lock} and \textit{Glyph of Warding}.

\subsubsection{Rogue's Gallery}
\begin{DndReadAloud}{}
This circular room may have once been a cistern. Upon the walls are now dozens of immaculate paintings and artworks. Several passages lead from this room to other exits. One has a metallic door covered in stunningly complex mechanisms and locks. Beside it is a wooden wine cabinet. The other is rather nondescript.

\end{DndReadAloud}


\begin{description}
\item[Treasure.]
There are forty three paintings in this room. Seven are recognizable as authentic masterpieces by some of the greatest artists who ever lived and are worth 2,500 gp each, the rest are worth 250 gp each. Garnering the full value for each painting would require hosting an auction to attract a host of wealthy collectors, however.
\item[Wine Cabinet.]
One of the alcoves of this room contains a wooden cabinet that holds twelve wine glasses and twenty two bottles of fine wine. Half contain real wine; two bottles are actually Potion of Gaseous Form disguised using \textit{Nystul's Magic Aura} to appear nonmagical. The rest are poisoned. A creature who drinks from one must succeed on a DC 17 Constitution saving throw, taking 35 (10d6) poison damage on a failed save, or half as much damage on a successful one. Only the Queen of Thieves knows which is which.
\item[Shackled Door.]
At the bottom of a short stair is an impressive door covered in a massive array of mechanical gears, locks, keyholes, buttons, latches, and handles. Each is frustratingly complex and may be disarmed or picked with a DC 20 Dexterity (thieves' tools) check, which either reveals yet another compartment with even more locks or resets a previously opened lock. There's no real logic to the system, it's a useless distraction as the door leads nowhere. However, one of the keyholes is actually a thin tube that leads to the office, but only a creature physically small enough to enter the keyhole could detect it.
\end{description}

\subsubsection{Queen's Gauntlet}
\begin{DndReadAloud}{}
Before you is a perfectly clean, featureless stone passage.

\end{DndReadAloud}


\begin{itemize}
\item This long hallway snakes around before ending in an \textit{Arcane Lock} door (DC 30 to open).
\item This devilish trap was devised by the Queen of Thieves to defend her personal sanctum; the entire hall is lined with a dynamic mechanical and magical system. Roll initiative when characters move down this hallway. Each round on initiative count 20 and 10 (losing initiative ties) the Queen's Gauntlet activates one of the following defences, which spring out from anywhere along the walls, floor, or ceiling:
\end{itemize}

\subparagraph{Multiattack}
The gauntlet makes two \textbf{whirling blade} or \textbf{poisoned darts} attacks, each against different targets within range.


\begin{description}
\item[Whirling Blade.]
\textit{Melee Weapon Attack}: +5 to hit, reach 5 ft., one target. \textit{Hit}: 14 (4d6) slashing damage.
\item[Poisoned Darts.]
\textit{Ranged Weapon Attack}: +5 to hit, range 30/120 ft., one target. \textit{Hit}: 4 (1d8) piercing damage plus 10 (3d6) poison damage.
\end{description}

\subparagraph{Battering Ram (Recharge 5-6)}
A section of the wall springs open. \textsl{Melee Weapon Attack:} +9 to hit, reach 5 ft., one target. Hit: 20 (3d10 + 4) bludgeoning damage. If the target is a creature, it must succeed on a DC 17 Strength saving throw or be flung 20 feet away from the battering ram and knocked prone.

\subparagraph{Snare (Recharge 5-6)}
A spinning bolas launches from a panel on the wall to ensnare one creature within 10 feet of any wall. The target must succeed on a successful DC 17 Dexterity saving throw or be restrained by sticky barbed wire. A restrained creature can escape by using an action to make a successful DC 17 Dexterity or Strength check.

\subparagraph{Flame Jet (Recharge 5-6)}
A bronze nozzle pops from the wall shooting gouts of flame in a 15-foot cone. Each creature in that area must succeed on a DC 17 Dexterity saving throw, taking 27 (8d6) fire damage on a failed save, or half as much damage on a successful one.

\subparagraph{Pit Trap (5/day)}
A ten-foot by ten-foot trapdoor in the floor snaps open anywhere in the hallway more than five feet from another open trap door, leading to a fifty-foot-deep pit. Creatures standing on the floor in the area must succeed on a DC 17 Dexterity saving throw or fall into the pit. On a successful save the creature leaps clear on the side of the pit (it can choose which). The bottom of the pit is different each time this ability activates:


\begin{description}
\item[Featureless.]
There's nothing at the bottom of this pit.
\item[Acid Pool.]
The bottom of the pit is filled with a five-foot-deep pool of acid. A creature who starts its turn in the acid takes 10 (3d6) acid damage.
\item[Spikes.]
The pit is lined with sharp spikes. A creature falling into the pit takes an additional 5d6 piercing damage.
\item[Contaminated Water.]
The bottom of the pit is filled with contaminated water ten feet deep. A creature that ends its turn there must succeed on a DC 17 Constitution saving throw or gain one level of Contamination.
\item[Liquid Hot Magma.]
Ten feet at the bottom of the pit is filled with lava. A creature who enters the pit or who starts its turn there takes 10d10 fire damage.
\end{description}

\subparagraph{Portcullis (5/day)}
A heavy iron gate slams down from the ceiling above. Creatures below can make a DC 17 Dexterity saving throw to leap clear on either side. A creature who fails this saving throw lands prone. Those who fail the save by 5 or more are instead crushed under the portcullis. The creature suffers 5d6 bludgeoning and 5d6 piercing damage, and becomes restrained under the portcullis until freed. A creature can lift the portcullis by using an action to make a successful DC 17 Strength check. It may also be attacked and damaged; it has AC 20, 30 hit points, and is immune to psychic and poison damage.

\subparagraph{Nerve Gas (2/day)}
The hallway releases a deadly neurotoxin. Each creature in the hall must succeed on a DC 17 Wisdom saving throw against this magical gas. On a failed save, the target takes 14 (4d6) psychic damage at the start of each of its turns. In addition, it can't use reactions, its speed is halved, and can take either an action or a bonus action on its turn, not both. These effects last for one minute. A target can repeat the saving throw at the end of each of its turns, ending the effect on itself on a success.

\subsubsection{Countermeasures}
Characters can use whatever means at their disposal to defend themselves, or simply take the Dash or Dodge actions and hope for the best, or try the following:

\subparagraph{Perception}
On its turn, a creature can spend an action to examine the area around it for signs of traps with a successful DC 17 Perception check. On a successful check, either the next saving throw it makes against the trap automatically succeeds, or the next attack roll made against it by the trap automatically misses.

\subparagraph{Disarm}
A creature can attempt to disable one of the gauntlet's systems. It chooses which action it wants to disarm and makes a DC 22 Dexterity (thieves' tools) check. On a successful check, the mechanism is disabled for one minute. On a failed check, the creature accidentally triggers the mechanism.

\DndQuote%
{}
{''Respect where it is due although the crown belongs to the youngest as the older generation has been corrupted by greed and madness.''}
{Pluto Jackson}
\subsubsection{The Office}
\begin{DndReadAloud}{}
Beyond the door is a tidy and luxurious office paneled with wood and bookshelves. Three bronze candelabras illuminate the chamber with soft ember light. A mahogany desk sits in the middle of the room with a sumptuous upholstered chair behind and two fine wooden chairs before it. Atop are several neatly stacked papers and journals, a lamp, and a box of fine fountain pens. Several large framed mirrors hang on the wall opposite the desk and to the left and right side of the room. A model globe rests in one corner, and in the other is a petrified wingback chair fixed to the floor.

\end{DndReadAloud}

A mental \textit{alarm} spell notifies the Queen of Thieves whenever someone enters this chamber.

\subparagraph{Bookshelves}
Gothic novels, play scripts, and other fiction involving dastardly criminals, tricksters, and other villains are found here. Written in the margins of each text are commentaries left by a genre-savvy reader about the many mistakes made by the villain in the story. Occasionally, particularly clever reversals made by the heroes are circled in red ink.

\subparagraph{Mirrors}
These magical mirrors are positioned such that anyone seated at the desk may see behind themselves and observe others in the room from every angle. In addition, the reflections of invisible creatures and objects are visible in the mirrors, which also reveal the true nature of any illusions or shapechangers. Hidden behind one of the mirrors is the end of a thin tube that connects to the Shackled Door in the Rogue's Gallery.

\subparagraph{Resplendent Desk}
Kept in the desk drawers are comprehensive dossiers on each of the player characters, faction leaders, and lieutenants. Not only do these notes detail virtually everything about their capabilities and exploits thus far, within may be secrets unknown even to the player characters themselves. The desk and chair are magical; a creature sitting at the desk on the chair can use an action to cast \textit{project image} to a range of 300 feet.

\subparagraph{Petrified Wingback}
A creature who sits on this stone chair is immediately teleported as if by \textit{dimension door} to an identical-looking chamber (see below). In that moment, the teleported creature sees an illusion of any other creatures in the room being disintegrated, while those creatures see an illusion of the teleported creature having its flesh painfully seared away, leaving behind nothing but a blackened skeleton that crumbles into ash if touched. Looking at the skeleton through the mirrors in the room reveals it to be an illusion.

\subparagraph{Globe}
This inaccurate but meticulously detailed globe opens at the equator, revealing a liquor cabinet with seventeen bottles of fine spirits and eight tumblers. Only four contain real alcohol; two bottles are actually Potion of Gaseous Form disguised using \textit{Nystul's Magic Aura} to appear nonmagical. The rest are poisoned. A creature who drinks from one must succeed on a DC 17 Constitution saving throw, taking 35 (10d6) poison damage on a failed save, or half as much damage on a successful one. Only the Queen of Thieves knows which is which.

\subsubsection{Secret Sanctum}
These chambers are deep under the earth, and almost entirely inaccessible except via teleportation. They are outside the Haze, and creatures can rest here normally.

This secret room is completely identical to the office in every way, and sitting on the wingback again returns a creature to the original office. However, the doorway exits to a short hall with two impressively secure iron doors. Each door is \textit{Arcane Lock} and warded with a \textit{symbol} spell. The rooms beyond are both protected by \textit{Mordenkainen's Private Sanctum}:

\subparagraph{Bedchamber}
A \textit{mirage arcane} spell has decorated this room to appear like the penthouse apartment in Castle Drakken with a magnificent view of the city as it was twenty years ago. The only real objects in the room are a comfortable bed, a table upon which the Queen of Thieves keeps her spellbook, and a \textit{decanter of endless water}. If you wish, a trinket that provides a revealing clue about her true identity is found here, and the \textit{mirage arcane} can be adjusted to better reflect the real person behind the Queen of Thieves. Hidden beneath a loose tile under the bed is a Spell Scroll of \textit{teleport} she uses to escape if she ever becomes trapped here.

\subparagraph{Queen's Vault}
This vault contains 10,000 gp worth of bullion, coins, and gems, neatly stacked, a \textit{helm of brilliance}, and one other rare magic item of your choice. If the Queen of Thieves has acquired any Seals of Drakkenheim, Relics of Saint Vitruvio, or other important items, they are kept here (unless she is using them herself).

\subsubsection{Arena}
\begin{DndReadAloud}{}
This natural cavernous chamber is roughly fifty feet high and supported by stone pillars. Perhaps it was a water reservoir, drained by the blockages created by the Queen's Men. Only a small pool of water remains here now, sloshing against a narrow sewer pipe.

\end{DndReadAloud}

When word gets out a fight is going down, the Queen's Men rush from the Undertavern down here using the secret staircase to watch the fight.

\subparagraph{Drainpipe}
Just above the waterline of the pool is a two-foot-wide drainpipe that used to siphon from the reservoir. The pipe travels one hundred feet and exits to the wider sewer system.

\subparagraph{Bleachers}
This narrow walkway forty feet above the arena floor sports a rickety wooden railing and dozens of benches.

\subparagraph{Queen's Balcony}
This overhanging section of the bleachers is reserved for the Queen of Thieves to watch the fights. A special throne is reserved for her, but she usually appears here using a \textit{project image} spell.

\subparagraph{Fighter Chamber}
Gladiators invited to fight in the arena get ready here. Weapon racks hold at least one of each martial weapon, several shields, and a few breastplates and helmets.

\subparagraph{Big Linda's Den}
This bone-strewn chamber is where \textbf{Big Linda} (see below) sleeps most of the time.

\subsubsection{Gladiators}
Should the characters find themselves volunteering to fight in the arena, a rival adventuring party or an evenly-matched Queen's Men strike team make for suitable foes. The no-holds barred contest merely stipulates that combatants who are knocked unconscious or surrender are pulled from the field. If the characters defeat their opponents, the Queen of Thieves invites them to face her champion, \textbf{Big Linda}.

\subsubsection{Big Linda}
\DndQuote%
{}
{''Big Linda has proven well worth the lives I spent to capture her. She is the prize champion of my fighting pit, and a loyal friend. I am very much excited for you to meet her. Also, before you ask, yes she bites''}
{The Queen Of Thieves}
\includegraphics[width=0.4\textwidth]{images/../5e.tools/img/adventure/DoDk/115-07-004.big-linda.png}
\textbf{Big Linda} resembles a dragon without wings, and has knobbly flesh accented by shimmering scales that stretch along its back. The monstrous beast has a gaping maw of razor-sharp teeth and a tail of razor-sharp spines, but dishevelled and useless forearms. Big Linda uses the game statistics of a \textbf{tyrannosaurus rex} with the following additional traits that increase its challenge rating to 9. Big Linda has been permanently charmed by the Queen of Thieves and (usually) obeys her spoken commands. The creature emerges from its den in response to the Queen's whistle call.


\begin{itemize}
\item Immunity to acid damage and blindsight 30 feet.
\item \textbf{\textit{Acid Spew (Recharge 5–6}).} Big Linda spews acid in a 30 foot cone. Each creature in that line of fire must succeed on a DC 15 Dexterity saving throw, taking 49 (11d8) acid damage on a failed save, or half as much damage on a successful one.
\end{itemize}

\subsubsection{The Queen's Request}
Characters who do defeat Big Linda in an arena battle may ask a single favour of the Queen of Thieves. Much like a capricious genie, she will grant the request, but will twist it to her advantage.

\subsubsection{Prison Block}
\includegraphics[width=0.4\textwidth]{images/../5e.tools/img/adventure/DoDk/116-07-005.thief.png}
If any player characters or faction agents are kidnapped by the Queen's Men, they are stripped to their skivvies and brought here. The Queen of Thieves takes any important or choice items belonging to captives and stores them in her vault, and the minions who made the capture get their pick of the rest. As such, there's no convenient place where captured characters can recover lost equipment, as it's quickly dispersed amongst the ranks of the Queen's Men. When they release prisoners, the Queen's Men don't return stolen equipment.

\subparagraph{Guard Room}
This dingy chamber has a rickety wooden card table and containers with manacles and chains. A bugbear \textbf{gladiator} named Mulg is chief jailer and is stationed here with a detail of four scoundrels.

\subparagraph{Interrogation Rooms}
These dingy chambers have a swinging lamp overhead and a single chair. Ropes, chains, pliers, and torture implements are kept within each. Any screams resound from within and travel unsettlingly down the halls. A \textbf{doppelganger} named Thought-stealer carries out interrogations for the Queen of Thieves.

\subparagraph{Cells}
These wretched cells stink of sweat and waste. They are locked (DC 20 to open). Keys are kept by Mulg, Thought-stealer, and the Queen of Thieves. Rolf Wagner, a rank-and-file member of the Hooded Lanterns, is kept here. He was captured months ago and the Queen of Thieves has been using his identity to spy on the Hooded Lanterns.

\section{Drakkenheim Garrison}
Not far from Shepherd's Gate lies the Drakkenheim Garrison. This building was the former stronghold of the city guard, and so made a perfect base of operations for the Hooded Lanterns. The Garrison takes up an entire city block; fortified stone walls extend from a squat round keep and enclose several barracks, stables, workshops, and storehouses. Deep beneath the garrison is a well-secured stockade prison.

The Hooded Lanterns secured a major victory just last year when they drove out the monsters occupying the garrison. When the greencloaks discovered that the prison cells were sufficiently shielded and insulated from the Haze, Lord Commander Elias Drexel ordered it repurposed into living quarters then began relocating his forces here. Since then, the regiment has expended considerable resources repairing the defences. It has been a gruelling, dangerous, and time-consuming operation, but now the renewed fortress is the perfect strongpoint for launching future incursions deep into the city.

\includegraphics[width=0.4\textwidth]{images/../5e.tools/img/adventure/DoDk/117-07-006.garrison.png}
\subsection{Interactions}
The Drakkenheim Garrison is a well-fortified base and meant to withstand infiltration, attacks, and sabotage. It is by far one of the most secure locations in the city besides Castle Drakken itself. Fully stocked, it could house and feed hundreds of defenders for weeks. However, since the Haze forces the occupants to rest in the underground prison and rapidly contaminates food and water stores, only about 300 soldiers occupy the base at any one time. Moreover, the force must still be constantly rotated to ensure troops stay fresh and well-rested, as field conditions are not unlike a castle under siege. Not everyone stationed here is a warrior, however, and a few especially brave cooks, smiths, and other personnel have come to aid the efforts of the Hooded Lanterns here.

\subsubsection{Visiting the Garrison}
While the Hooded Lanterns are open to new recruits, they regard those who approach their holdings within the city with suspicion, expecting either scavengers looking for help and handouts or worse, saboteurs looking to infiltrate their base. Nevertheless, they are willing to help injured and contaminated adventurers evacuate the city, and those seeking either permission to pass through Shepherd's Gate or an audience with \textbf{Lord Commander Elias Drexel} are invited to enter the garrison and escorted by a detail of 1d4 urban rangers to the Commander's Keep.

\subsubsection{Defences}
The Garrison is perpetually in a state of readiness, and able to respond to any alarm for attack, intrusions, or infiltration. The scouts and urban rangers on guard duty are each equipped with signal flares that may be loaded with various smoke colours to call for help or retreat, warn about monsters, flag intruders, and give the all clear. In addition, the Garrison Towers are each equipped with bellowing alarm bells that can be heard throughout the city. When rung, resting Hooded Lanterns spring into action and patrols are recalled from the streets to defend the garrison.

\subsection{Area Details}
\includegraphics[width=0.4\textwidth]{images/../5e.tools/img/adventure/DoDk/118-map-7.03-garrison.png}
\includegraphics[width=0.4\textwidth]{images/../5e.tools/img/adventure/DoDk/119-map-7.03-garrison-player.png}
\begin{DndReadAloud}{}
A small fortress looms over the streets, its towers jutting out above the impenetrable walls and thick iron portcullis. Several circular bastion towers topped with ballistas reinforce the walls. One can make out from a distance several green-cloaked figures with crossbows manning the walls and towers.

\end{DndReadAloud}

\subparagraph{Walls}
The masonry walls of the Garrison are thirty feet high and five feet thick. A wooden scaffold lines the inside of the walls with ladders and stairs ascending from the inner courtyard.


\begin{itemize}
\item Each segment of the wall is patrolled by 1d6 scouts at all times.
\end{itemize}

\subparagraph{Towers}
These fifty-foot-tall towers are ten to twenty feet in diameter. Atop is a crenellated battlement accessible via a wooden hatch and fixed with a ballista. The towers are accessed from inside the garrison by a short wooden staircase, and within each tower are four wooden floors connected by sturdy ladders. Arrowslits, windows, and stockpiles of ammunition are on each level, but the topmost floor has an arched window facing the courtyard set with a large bell that sounds in case of an attack. An \textbf{urban ranger} and 2d4 scouts are posted to each tower.

\subparagraph{Gates}
Two gates lead into the garrison; the south is flanked by bastion towers and the other is set into the north wall. Each is an archway fixed with ironbound and reinforced double doors that may be fastened closed with a bolted brace. Two urban rangers stand guard on the ground outside each gate.

\subsubsection{Training Yard}
This courtyard rests within the walls and has been fashioned with a fenced off area that comprises the majority of it. The ground here is muddy and worn. The far end of the courtyard is lined with sandbags and targets for archery, and a few scarecrow dummies set up for weapons training. Along the left side of the yard a small wooden platform has been erected where Ansom or Petra usually oversee training.


\begin{itemize}
\item During the day 4d6 scouts, and 2d4 urban rangers are training here.
\end{itemize}

\subparagraph{Stables}
Off to the side of the training yard is a small stable that is no longer in use for housing horses. Instead, this building has been repurposed into cells to keep prisoners short term if the need arises.

\subparagraph{Well}
This well draws from the city aqueducts, and so the water is contaminated. As a result, the Hooded Lanterns have blocked it off with rubble and boarded over the opening.

\subparagraph{Storehouse}
A small storehouse is constructed along the southern wall of the yard. This small wooden shed is primarily used to store raw materials such as iron, wood, and various tools that may be used for general purposes around the yard. Due to the impact of the Haze, this storehouse can no longer be used for storing cured meats, water, or any other items that could spoil due to the contamination.

\DndQuote%
{}
{''Survivors fueled by hope. A familiar feeling of battle readiness hangs in the air.''}
{Pluto Jackson}
\subsubsection{Armoury Tower}
This four storey tower rests at the southwest corner of the garrison. The ground floor houses a small table with four chairs and a deck of cards where guards often take breaks. The levels above are equipped with stockpiles of weapons, armour, shields, ammunition, and materials that are ready for use.


\begin{itemize}
\item 4 urban rangers and twelve scouts occupy this tower at any time.
\item \textbf{Smithy}. At the base of the armoury tower is a small open workstation and blacksmith forge where gear and ammunition is manufactured or maintained.
\item 2 commoners and scouts can be found here at any time working on repairs and maintenance on gear.
\item Modrin Engers, a human \textbf{priest} devoted to the old god of craftsmen works the smithy. He also helps out in the infirmary with healing, as he knows the \textit{purge contamination} spell.
\end{itemize}

\subsubsection{Barracks}
Three large buildings inside the garrison make up the barracks. The ground floor of each wooden building consists of a mess hall, kitchen, infirmary, and a storeroom. The second level of each is a long room for bunk beds and footlockers. However, the Haze renders the bunks unusable for rest and recuperation. One of the buildings has been entirely boarded up and the furniture moved into the stockade. The other has been similarly emptied of furniture, and repurposed to store construction materials, spare parts for siege equipment, and barricades in case repairsor extra fortifications are needed.

\subparagraph{Mess Hall}
The north barracks building still operates as a common hall for the troops, and is often filled with off-duty staff and rangers taking their meals.

During the day, the barracks area contains 5d6 scouts, 2d6 urban rangers, and 3d6 commoners working in the kitchen, relaxing, or enjoying meals in the hall.

\subsubsection{Lord Commander's Keep}
The largest structure in the garrison is a squat four storey stone keep, sixty feet tall and forty feet wide. It features a broad battlement level upon which a catapult has been constructed. \textbf{Lord Commander Elias Drexel} keeps his personal quarters and coordinates the operations of the Hooded Lanterns here. A set of broad steps lead up to the fortified entrance doors, and each floor is connected by a square staircase on the southern edge of the tower.


\begin{itemize}
\item 6 urban rangers and 12 scouts form the security detail for the keep, with two posted outside the front entrance at all times, and about half manning the battlements.
\end{itemize}

\subparagraph{Reception Hall}
This stern entrance hallway is supported by several stone pillars and features a large hearth. The walls bear tapestries depicting the coats-of-arms of previous Lord Commanders.

\subparagraph{War room}
Taking up the second level of the keep is a meeting room with a large round table at the centre. A high wooden chair is set before the coat of arms of Drakkenheim, and many smaller chairs are placed around the table. A highly detailed and complete map of Drakkenheim is spread out on the table before a panoply of scout's reports, mission briefs, dossiers, and supply lists.


\begin{itemize}
\item During the day, Eliss Drexel holds important discussions and devises plans and operations for his goals in Drakkenheim. The busy man is almost always meeting with a Hooded Lanterns strike team, or Petra and Ansom Lang.
\end{itemize}

\subparagraph{Offices}
There are a few small rooms on this floor with desks equipped with inks and papers, as well as a common area lined with bookshelves, documents, maps, and other miscellaneous works and writings. Much of the bureaucratic work of the city guard was carried out here, such as soldiers' pay and schedules, but now such formalities are largely dispensed with so the room is rather disused.

\subparagraph{Commander's Quarters}
The top level of the tower is a regal bed chamber with large windows, a writing desk, wardrobe, and a thick fur rug. The quarters appear untouched and have gone unused for some time. Elias Drexel hasn't bothered to use these quarters since he returned to Drakkenheim, and keeps a bunk in the stockade with his troops rather than use this Haze-filled chamber.

\DndQuote%
{}
{''The Hooded Lanterns saved my bacon when the city fell, I would have been toast living on the streets of Drakkenheim alone! These loyal survivors strive to keep unprepared adventurers from the dangers beyond the walls, if only to make sure they can survive to pay their due.''}
{Veo Sjena}
\subsubsection{Stockade}
The stairs down from the keep stop at a small office, then descend further into the dungeons.

Beneath the garrison is a small office and a large dungeon lined with cells. This entire area has been repurposed into the new barracks of the Hooded Lanterns, housing their entire force for rest and recuperation.

\subparagraph{Jailer's Office}
At the bottom of the stairs is a small office with an iron gated door on the far side leading to the cell blocks. A small desk and chair are set next to it. An \textbf{urban ranger} can be found on duty here at all times

\subparagraph{Cell Blocks}
Each cell block consists of long hallways with a dozen individual cells on each side accessed via a sturdy iron door with a barred window. Originally, the cells were used for one or two prisoners, but now bunks have been installed within each to sleep four soldiers. Footlockers under each bunk hold the soldiers' personal belongings. Only the Lord Commander has a private cell.

\subparagraph{Infirmary}
Instead of being used as bunks, several cells in this block were repurposed into a spartan infirmary equipped with Healer's Kit and medical supplies. Furthermore, a few of the cells here are still used for this original purpose if the need arises to house prisoners for longer than a day.


\begin{itemize}
\item Cal Grice, a human \textbf{druid} serves the Hooded Lanterns as a medic. They know the \textit{purge contamination} spell.
\item The Hooded Lanterns keep a stockpile of 10 (3d6) Potion of Healing and 1d4 jars of Keoghtom's Ointment here.
\end{itemize}

\section{Inscrutable Tower}
Alongside the distant Paradox Castle and the Enigma Ziggurat, the Inscrutable Tower is one of three widely known strongholds of the Amethyst Academy. Though there are five more hidden throughout the world, this is unquestionably the greatest of them all. The loss of the tower was a nearly crippling blow for the otherwise mighty and resourceful Academy, though the Directorate has taken great pains to conceal this uncomfortable fact.

This enigmatic structure was a mystery even before a meteor struck it. None can say exactly when it appeared in Drakkenheim, but most think that it must have been some time after the Edicts of Lumen. Yet there is no record of its construction, which seems incongruent given the massive tower would have taken decades to build.

\includegraphics[width=0.4\textwidth]{images/../5e.tools/img/adventure/DoDk/120-07-007.meteor-impact.png}
In truth, the Inscrutable Tower is the ancient and legendary hidden fortress of the Sorcerer Kings. This extraplanar dungeon was built within the roiling realm of chaotic possibility known as the Space Between Worlds in far antiquity by powerful and insane sorcerers to hide away their magical knowledge, but was left forgotten and abandoned for a thousand years. When the mages of the Amethyst Academy located it several hundred years ago, they tapped into the immense power of its Arcane Nexus to shift the whole structure into the mortal world, where it now stands in Drakkenheim. Since then, the Amethyst Academy gradually expunged most records and memories of this accomplishment to rehabilitate the tower's evil and ominous history. Yet there are many secrets about this place unknown to even the Amethyst Academy, including whatever magical forces now hold the tower together even after a meteor smashed through it.

\subsection{Overview}
Unlike the other faction strongholds described herein, the Inscrutable Tower is not occupied by its controlling faction. Instead, the Amethyst Academy aims to regain control of it. Though it will take them months, if not years, to refurbish and repair the tower once this is done, it could cement the Academy's controlling influence over Drakkenheim.

\subsubsection{Adventure Hooks}
\subparagraph{Seize the Tower}
Securing control of the Inscrutable Tower is a make-or-break mission for the Amethyst Academy, but also its enemies. If the Academy is able to regain control over the tower and begin repairs, it will allow them to enact more ambitious actions to seize control of Drakkenheim. On the other hand, the Amethyst Academy can be organizationally paralyzed if another faction occupies the tower and takes the \textit{Inscrutable Staff}.

\subsection{Interactions}
\subsubsection{Interloper}
Several malfeasant mages are found throughout Drakkenheim: \textbf{Oscar Yoren}, the \textbf{Pale Man}, and Ryan Greymere. If any of them have survived to this point in the campaign, choose one of them to make their own incursion into the Inscrutable Tower to seize it for themselves (you might want to have Oscar Yoren "level up" to an \textbf{archmage}). They gather up some appropriate minions of their own: wights, haze hulks, or summoned \textbf{elementals} (respectively) and are encountered at the breach in the library, plumbing the arcane knowledge within. The malfeasant has managed to claim the \textit{Inscrutable Staff} for themselves, and are searching for an alternate means to control the Nexus as they cannot do so alone.

\includegraphics[width=0.4\textwidth]{images/../5e.tools/img/adventure/DoDk/121-07-008.academy-tower.png}
Alternatively, if the players have made an enemy of the Amethyst Academy, \textbf{Eldrick Runeweaver} sends his \textit{simulacrum} here instead under similar circumstances. If his simulacrum was already destroyed, he arrives in person with an \textbf{iron golem} and an Amethyst Academy strike team.

Finally, if neither seems like a viable option, an \textbf{efreeti} has managed to escape the Deep Cells and conjured several fire elementals and azers to serve it. The creature hasn't fully escaped the binding magic holding it to the Inscrutable Tower, and is looking through the books for another means of escape.

\subsubsection{Random Encounters}
Check for random encounters each time characters enter rooms in the Inscrutable Tower:


\begin{DndTable}{cX}

1d6 & Random Encounter\\
1 & Unfortunate Apprentices. Delerium dregs. Each can cast two wizard cantrips at-will.\\
2 & Restless Furniture. \textbf{Animated armor}, flying swords, and \textbf{rug of smothering}\\
3 & Residual Magic. 10 (3d6) Will-o-wisps\\
4 & Disembodied Mages. 5 (2d4) Arcane wraiths\\
5 & Elements Unbound. 1d6 elementals of any kind (fire, air, water, or earth)\\
6 & Malfunctioning Protectors. 1d4 malfunctioning shield guardians. There's a 50\\% chance each still is storing a spell within.\\
\end{DndTable}

\includegraphics[width=0.4\textwidth]{images/../5e.tools/img/adventure/DoDk/122-07-010.the-tower-academy.png}
\subsection{Area Details}
Three-quarters of the Inscrutable Tower is composed of a massive library (detailed below) that spans most levels. The remaining floors in this unfathomably large structure have a sprawling maze-like layout with roughly a dozen rooms connected by a twisting set of corridors. There is no consistency in the floorplan from floor to floor, and the configuration of the rooms is often illogical and may even defy the spatial geometry. Dormitories might be beside highly dangerous laboratories, or storage rooms for ingredients used by a workshop kept several floors apart. You can construct layouts as you see fit by using the various rooms described in the Residence Floors and Research Floors. The Lobby, Great Hall, Library, Director's Office, Nexus, and Deep Cells are singular specific locations in the tower, however.

\subparagraph{Exotic Materials}
The walls, halls, and furnishings found in the Inscrutable Tower are made from obsidian, dark steel, granite, black marble, and glass.

\subparagraph{Magical Amenities}
Magic provides lush comforts throughout the Inscrutable Tower. Fireplace hearths don't have chimneys, instead using special \textit{continual flame} spells that radiate pleasant but otherwise harmless heat. Most rooms are illuminated with persistent \textit{dancing lights}. Bound \textit{Unseen Servant} tend the halls, keeping rooms tidy. Magical fonts produce water via the \textit{Create or Destroy Water} spell. A creature can focus their attention (as an action) on a light source, heat source, or water source they can see to cause it to turn on or off, or adjust its brightness or intensity. The parameters are such that these can't become so hot or cold to harm a creature, or so bright as to blind them. However, each time these are manipulated, roll 1d20. On a 1, a random arcane anomaly is triggered due to the erratic magic that grips the tower.

\subparagraph{Magical Wards}
Powerful magic protects the Inscrutable Tower. The exterior walls and doors are magically reinforced and are immune to all damage (except meteors). The doors and walls extend into the ethereal plane, blocking ethereal travel through them, and magic can't alter their forms. Creatures outside the tower can't teleport inside, nor can they teleport within the tower rooms unless they have a direct line of sight to their destination.

\subsubsection{Tower Exterior}
\begin{DndReadAloud}{}
The Inscrutable Tower of the Amethyst Academy is an immensely tall structure that reaches a dizzying height of twelve hundred feet. The shimmering obsidian walls are perfectly smooth; there is not a single window or embellishment along its outer face save a deep vertical groove cut into the pinnacle. Seen at a distance the tower appears relatively narrow, though it is roughly one hundred fifty feet wide and anchored by sharply arched buttresses at the base. The tower has been smashed apart roughly two-thirds up its height. Remarkably, the top of the tower hovers above in utter defiance of gravity. Between the upper and lower sections is a gaping breach with pieces of floating debris locked in the air; frozen in time at the moment a meteor soared through it. A keen-eyed observer can see fragments of delerium, broken bookshelves, and scattered tomes bobbing about the wreckage.

\end{DndReadAloud}

\subsubsection{Academy Plaza}
\begin{DndReadAloud}{}
The cobblestone streets break here, a campus park surrounding the Inscrutable Tower. A forest of delerium crystals have grown up around the tower that surrounds a small artificial lake. Overlooking the pond is a large stone plinth surrounded by arcane arches.

\end{DndReadAloud}

The area around the Inscrutable Tower is covered by the Deep Haze.


\begin{itemize}
\item Etched into the stone is a great \textit{teleportation circle}.
\item The remains of a pair of stone golems who once guarded this site.
\item 1d4 protean abominations and 2d4 gibbering mouthers inhabit the plaza.
\end{itemize}

\subsubsection{Tower Entrance}
\begin{DndReadAloud}{}
A broad staircase of seventeen steps bordered by a sloping ramp leads up to the singular entrance of the Inscrutable Tower. The dark steel doors are engraved with the sigil of the Amethyst Academy: a lidless arcane eye set within an octagram. There are no handles or hinges upon the door.

\end{DndReadAloud}


\begin{itemize}
\item A character attuned to a set of \textit{academy rings} who places their hands upon the door causes a set of glowing arcane runes to appear. After a short set of arcane clicks and whirring noises, there is a flash of light and the doors slide open with a screech.
\item Since \textit{academy rings} are made for a specific individual, if the player characters don't possess a set of their own, they will need either an ally or a living prisoner wearing them to gain entry into the Inscrutable Tower. Characters can also enter via the breach itself, assuming they can fly or climb there.
\end{itemize}

\subsubsection{Lobby}
\begin{DndReadAloud}{}
The front doors lead into a massive lobby that fills the tower base. The ceiling is nearly fifty feet high and supported by obsidian pillars etched with intricate sigils, and the floor is reflective marble decorated with intricate arcane patterns. Rows of tall, thin vents along the outer walls form transparent openings allowing a view of Drakkenheim outside. On two opposite sides of the lobby are three sets of metallic double doors, and railingless spiral staircases of metal and glass wrap upward around the pillars on each corner.

Throughout the rest of the room are chest-high stone dividing walls that form semi-private meeting areas set with metallic chairs and glass-topped tables. Bobbling motes of spectral lights drift through the air, illuminating the chamber with soft arcane light.

\end{DndReadAloud}


\begin{itemize}
\item Four erratic shield guardians linger in this chamber and attack any who enter. Two hold spells: one \textit{wall of fire}, the other \textit{ice storm}, and both are as if cast by a \textbf{mage}.
\end{itemize}

\subsubsection{Elevators}
\begin{DndReadAloud}{}
Four of the six sets of metal double doors are closed, one has slid partially ajar. The last is erratically slamming shut and jolting open at random. Looking beyond is a ten-foot-wide shaft that bristles and hums with arcane energy. Once in a while there is a thunderous crack of lightning that resounds up and down the shaft.

\end{DndReadAloud}


\begin{itemize}
\item The shafts extend up the length of the tower, but are broken by the breach. Four of these elevators still function, however, and can provide access to the lower portion of the tower.
\item Once every 1d6 rounds a \textit{lightning bolt} (DC 16) shoots up the length of the two broken shafts.
\item The closed doors magically open when any creature approaches within ten feet of them, and a prismatic platform of magical energy forms within the shaft, above which float glimmering arcane runes arranged vertically: from bottom to top, the runes read "Lobby", "Great Hall", "Residences", "Research Levels", "Library", and "Director's Office".
\item A creature may step onto the platform and speak where it wishes to go, and the platform ascends with stomach-churning speed to that level.
\item There are multiple runes for the residences, research levels, and library indicating the dozens of levels for each, thus the speaker must specify individual floors. The levels past "63" don't function, so travel via the elevators to the upper portions of the library, the director's office, or the nexus isn't possible.
\end{itemize}

\subsubsection{Great Hall}
\begin{DndReadAloud}{}
This resplendent dining hall prominently features an elevated platform where a long marble-topped dinner table with high chairs faces the room. Before it are three long rows of polished granite-topped tables with dark steel chairs. Spectral chandeliers and candles of magical energy hover above, illuminating the room. Along the walls are tall portraits of archmages and professors of the Amethyst Academy.

\end{DndReadAloud}


\begin{itemize}
\item The portraits are magical, and the figures within come to life and \textit{counterspell} anyone who casts a spell in this room.
\item The sleek dinnerware are all magically animated, and trundle into their place at the dinner tables with food from the kitchen, then whisk away to the scullery to clean themselves.
\item The large kitchen is immaculately clean and organized, featuring well-made iron ovens, stoves, steel-topped food preparation tables, and organized glass cabinets for the dinnerware and utensils. The cookware is similarly animated, preparing food and tending the kitchen on its own.
\item The well-stocked larder opens into a vast featureless extradimensional space containing prodigious food stores upon neatly organized metal shelves, all magically preserved. The kitchen cookware daintily retrieves ingredients as needed.
\end{itemize}

\subsubsection{Residence Floors}
Throughout the Inscrutable Tower are numerous apartments and quarters for the mages:

\subparagraph{Apprentices' Dormitory}
These rooms contain four single bunks set with personal chests and nightstands. There is a shelf over each bed for personal journals with an iron lantern hanging from it. Clusters of these rooms are often found around a small common room with a tiled hearth, glass-topped table, and shelved books on mundane topics.

\subparagraph{Wizard's Apartments}
These well-appointed apartments are sealed with \textit{arcane lock} spells. Most often, entry is keyed to the original (long dead) occupant only. Each consists of several rooms:


\begin{itemize}
\item A common room to hold meetings and take meals.
\item A workroom, summoning chamber, mediation room, or scrying pool.
\item A small study and private library to conduct personal research.
\item A private bedroom and water closet.
\end{itemize}

A permanent \textit{unseen servant} tends each apartment, keeping it clean and tidy. Left behind in some apartments are Spellbook containing 10 (3d6) wizard spells of 5th level or lower, plus 1d4 uncommon Spell Scroll or potions.

\subsubsection{Research Floor}
Other floors of the Inscrutable Tower are dedicated to the study of magic and the manufacture of magical items:

\subparagraph{Alchemy Labs}
These chambers reek of chemicals, and contain cauldrons for creating potions or brewing alchemical concoctions. Cupboards and shelves contain material components for transmutation spells.

\subparagraph{Arcane Workshops}
These rooms are used to craft magic items.


\begin{description}
\item[Golem Workshop.]
Making golems. Partially complete golem.
\item[Scriptorium.]
Scribing scrolls.
\item[Surgeon's Chamber.]
Not common, but there is one such room in this tower to provide biological and necromantic research.
\item[Spellforge.]
A pedestal before a "firing range." "Target dummies."
\item[Elemental Forge.]
Fabrication Chamber.
\end{description}

\subparagraph{Chamber of Illusions}
These square rooms have scalar lines in a one-foot grid along the ceilings and walls. The otherwise entirely featureless void is the perfect place to practice illusion magic.

\subparagraph{Copy Room}
These rooms each contain 1d6 bound imps chained to writing desks. These miserable creatures are tasked with making copies of any book or writing given to them.

\subparagraph{Materials Storage}
Secure chambers sealed with \textit{arcane lock} and \textit{alarm} spells. Spell components, inks, herbs, gemstones and other rare arcane ingredients, mundane materials such as stones, metals, leather, wood, cloth, parchment, and paper.


\begin{description}
\item[Morgue.]
Bodies used for biological and necromantic research.
\item[Menagerie.]
Contains animals.
\end{description}

\subparagraph{Seminar Rooms}
Used for classes.


\begin{itemize}
\item A handful of larger lecture halls are twice the size, suitable for groups of forty to fifty. Larger gatherings are held in the Assembly Hall, however.
\end{itemize}

\subparagraph{Summoning Chamber}
Specially prepared with materials needed for summoning and binding extraplanar creatures. Some have \textit{Magic Circle} and \textit{dimensional shackles}.


\begin{itemize}
\item One of the summoning chambers contains a woman sitting in the middle of a magic circle. The circle is lined with salt. She begs players to break the salt barrier so she can leave. She argues that evil mages trapped her here and she has been stuck for years. She will plead with the players and offer them aid, information, wealth, or anything they want to hear. If the seal is broken she instantly transforms into a \textbf{bone devil}, and attempts to kill the party and flee.
\end{itemize}

\subparagraph{Scrying Pools}
These private and secure round meditation chambers contain a mirror, pool of water, or other font needed for divination spells.

\subsubsection{Library}
\begin{DndReadAloud}{}
The Library of the Inscrutable Tower is a vast and dizzying place. Staircases zig-zag from wall to wall, weaving up the height of the tower between walkways of black marble with railings of dark steel and massive floating bookshelves of speckled granite and glass lined up like airborne dominoes. The stairs and shelves rumble lowly as they shift positions to converge at suspended balconies and galleries where obsidian-topped tables, plush couches, and amber spheres of glowing light are arranged into reading nooks and studying desks. Throughout, many stairs and platforms are nauseatingly oriented sideways, some even completely upside-down and backwards.

Every conceivable surface stores books. The balconies, furniture, walls, and even the steps of each stair all have cavities, niches, and shelves containing tomes of varying shapes, sizes, materials, and colours, some in neat and tightly organized rows, others in disheveled piles and scattered masses of pages.

\end{DndReadAloud}


\begin{itemize}
\item The library comprises an ever-changing and dynamic layout. It is one hundred fifty feet wide on each side.
\item The elevators stop at the ground level of the library and rise no higher. Before the damage was done, secret doors behind bookshelves along the walls opened for easier access, but these are broken now.
\end{itemize}

\subparagraph{Stairs}
Most staircases are thirty to fifty foot long sections, each five feet wide. The stairs connect and disconnect to each other and the platforms by some unfathomable pattern. For any given stair section, roll 1d20 each round, and on a 1, it moves.

\subparagraph{Platforms}
The various platforms have a variety of equilateral shapes - square, triangular, hexagonal - and range from ten to thirty feet in diameter. The ground level of the library is arranged into reading nooks and study desks like the other platforms.

\subparagraph{Gravity}
Gravity is subjective here. Characters can walk normally on staircases and platforms, even sideways or upside-down. Creatures levitate in midair instead of falling, but those lacking a flying speed must move by pushing or pulling against a fixed object or surface within reach (such as a railing or bookcase), which allows them to move as if they were climbing. Unattended objects float around randomly.

\subparagraph{Soaring Intellect}
A creature can fly and hover within the library with a speed equal to ten times its Intelligence modifier.

\subparagraph{Tomes}
Herein are books on every conceivable topic, from artworks, history, natural philosophy, magical matters, mathematics, and more. Among the books are introductory and advanced texts, exhaustive commentaries on those texts, poignant commentaries on those commentaries, scathing refutations of that commentary, a rebuttal to refutations of the commentaries, treatises that purportedly reinvent the topics altogether, and many works of complete fiction and utter fantasy.

\subparagraph{Spellbooks}
The library contains multiple tomes detailing every spell from the sorcerer, warlock, or wizard spell list in the \textit{Core Rules} (and many more beyond) with the notable exceptions of \textit{wish}, \textit{demiplane}, \textit{clone}, \textit{simulacrum}, \textit{gate}, \textit{plane shift}, and the new spells described in Appendix D.

\subsubsection{Shattered Breach}
\begin{DndReadAloud}{}
Looking above from the library, characters can see the breach in the tower. Visible are the remains of that section of the library, where the stairs, shelves, and platforms still shift and move, but in a field of wrecked debris. A crackingly storm of eldritch lightning erupts between the gently floating debris, cast adrift gently in space, defying gravity.

\end{DndReadAloud}


\begin{itemize}
\item The breach is a one-hundred-fifty-foot-wide gap in the library. Traversing the breach requires flight, but there are several stray shelves, platforms, and stairs creatures can rest upon.
\item Chunks of delerium and meteoric iron hover about the air, interacting with the magic within.
\item Gravity remains broken in the breach, and creatures may still fly via \textit{soaring intellect}.
\item A creature who ends its turn in the breach rolls 1d6. On a 1, it's struck by eldritch lightning. It must succeed on a DC 15 Dexterity saving throw or take 10d6 lighting damage, or half as much on a successful save.
\item Casting a spell in the breach automatically triggers an arcane anomaly.
\end{itemize}

\includegraphics[width=0.4\textwidth]{images/../5e.tools/img/adventure/DoDk/123-07-009.old-book.png}
\DndQuote%
{}
{''I've been to a couple Academy strongholds, none of them felt welcoming, this one feels less so.''}
{Sebastian Crowe}
\subsubsection{Director's Office}
\begin{DndReadAloud}{}
Behind an impressive set of double doors is the director's audience chamber. Opposite the entrance to this circular room is a sumptuous obsidian throne flanked by iron candelabras. The sigil of the Amethyst Academy decorates the wall overhead, and the marble floor depicts a complex planar astrolabe. Overhead, a chandelier of steel and copper rings gently rotates, showing the movement of the celestial spheres. On either side of the throne is an arched iron door.

\end{DndReadAloud}


\begin{itemize}
\item This room is located on one of the topmost floors of the tower, just below the Nexus.
\item The archmages of the Amethyst Academy keep personal chambers within the Inscrutable Tower merely as a formality: all maintain private \textit{Demiplane} as their primary abodes. Instead, this room serves as a meeting place to conduct business with others.
\end{itemize}

\begin{DndReadAloud}{}
The left door opens outside the tower itself. As the door opens, a glimmering balcony of magical force takes shape that rings the entire tower with a narrow walkway five feet wide. There is a heaving lurch, as panels in the outer wall swing open below. A great telescope of glass and bronze emerges from within, supported by a metallic crane that moves it into position over the balcony. A chair is built into the side of the telescope, upon which are several intricate switches and levers.

\end{DndReadAloud}


\begin{itemize}
\item Controls on the telescope allow it to be moved and rotated around the tower.
\item Several runes on the telescope play back a record of past observations. The last rune indicates a record was made the day the meteor struck, indicating that the archmage only became aware of the meteor in the final moments before it struck.
\end{itemize}

\begin{DndReadAloud}{}
The right door opens into a smaller circular chamber where a round granite slab hovers three feet off the ground. Splayed out over it are all manner of missives, magical implements, and occult instruments. Glass shelves holding arcane components and books line the walls.

\end{DndReadAloud}


\begin{itemize}
\item \textbf{\textit{Treasure}.} This is a tuning fork keyed to the mortal plane, the material focus for a \textit{plane shift} spell. Scrolls of \textit{teleport} and \textit{forcecage}, and a very rare magic item can be found here.
\item A sealed chest in this chamber contains the components needed to repair the clocktower, including an ancient (pre-meteor) delerium crystal in a preserved glass sphere.
\end{itemize}

\subsubsection{Nexus}
\begin{DndReadAloud}{}
A gigantic and incomprehensible arcane apparatus whirs, clicks, and hums upon an obsidian plinth at the centre of this chamber. Revolving rings of brass, rotating golden plates embossed with diagrams of magical formula, whirling crystal lenses, spinning bronze rings etched with markings bearing out the movement of the planes, receptacles bearing lambent glass orbs of various colours and sizes, and delicate brass instruments orbit around a pulsing sphere of collecting energy. It glows with a strange prismatic light as it erratically convulses and emits spasmodic arcs of arcane plasma that send shudders and shakes through the room.

Surrounding the machine are eight massive cushioned chairs held aloft by articulated bronze arms. Levers, buttons, switches, and plates that glow with magical runes are built into the armrests, and an array of telescoping apertures and lenses are affixed to the headrests with jointed metal rods. Seated within each is a warped and barely recognizable lifeless body.

The ceiling is a massive dome made from interlocking glass panes in organic shapes, each which gazes upon a different alien sky. Booming lightning crackles above. The rest of the room is ringed with shelves holding jars of spell components and eldritch curiosities.

\end{DndReadAloud}


\begin{itemize}
\item This is the topmost level of the Inscrutable Tower.
\item The nexus contains a dimensional rift that has become unstable due to the damage caused by the meteor. The nexus must be rebooted to stabilize the dimensional rift.
\item Rebooting the nexus requires each of the eight chairs be disengaged and shut down. A creature seated in the chair can use an action to manipulate the apparatus. Doing so takes an action and DC 20 Arcana, Investigation, or Sleight of Hand check. Checks can also be made with tinker's tools or thieves' tools.
\item On a failed check, the creature suffers 4d6 psychic damage from backlash. Failure by 5 or more causes the creature to become stunned for 1 round.
\item Each attempt immediately triggers an arcane anomaly.
\item After the first chair is reset, on the following round on initiative count 20, 1d6 random extraplanar creatures are conjured by the erratic interference of the machine being partially on and partially off.
\item Once all eight chairs are successfully disengaged and shut down, the nexus powers down and the dimensional rift becomes stable.
\end{itemize}

\includegraphics[width=0.4\textwidth]{images/../5e.tools/img/adventure/DoDk/124-07-011.academy-tower-cut-out.png}
\subsubsection{Deep Cells}
\begin{DndReadAloud}{}
The floor of this forty-foot-wide obsidian chamber contains only a twenty-foot-wide cylindrical shaft lined with rows of hexagonal cells arrayed in a honeycomb.

\end{DndReadAloud}


\begin{itemize}
\item This is the bottom chamber of the Inscrutable Tower, but space warps around itself, and it can only be accessed from a stairway opposite the Director's Office that seems to impossibly lead back down the entire tower.
\item The cylindrical shaft descends one thousand feet into the earth before ending in an immobile \textit{sphere of annihilation}. All in all, there are over eighty thousand cells here.
\item Two stone golems stationed here as guards have gone berserk, and attack any who enter the chamber.
\item A \textit{wall of force} once covered the opening of the shaft, but the magic has failed. Fortunately, most of the other defensive features of the cells remain intact:
\end{itemize}

\subparagraph{Cells}
Each cell is a hexagonal cavity about ten feet wide and ten feet deep. An \textit{antimagic field} fills each, and either a \textit{wall of force} or a \textit{prismatic wall} covers the opening. An arcane rune identifying the contents of the cell hovers in the air before the wall. A destroyed wall reforms one round later.

\subparagraph{Keys}
The magical wards upon each cell are individually keyed to a specific person, a password, and a token, and are manipulated by providing two of the keys and then speaking a command word within ten feet of the corresponding cell. The command words and their effects are as follows:


\begin{description}
\item[Secure.]
Raises the \textit{Wall of Force} and activates the \textit{antimagic field}.
\item[Release.]
Lowers the \textit{wall}.
\item[Disengage.]
Deactivates the \textit{antimagic field}.
\item[Discharge.]
Lowers the \textit{wall} and deactivates the \textit{antimagic field}, then causes a blast of force to violently expel the contents of the cell into the shaft, where they fall into the \textit{sphere of annihilation} below.
\end{description}

Each token is a magical trinket and could be almost any kind of object, but each always bears an invisible arcane rune that identifies the corresponding cell. Thousands of such objects are scattered about the Inscrutable Tower.

The exact contents of the cells are best left to the imagination, but they could contain all manner of extraplanar entities, fiends, and magical artefacts. Given the variety of creatures and objects stored within, it isn't uncommon for unique specimens to have specific safeguards.

\begin{DndSidebar}[float=!b]{Developments}
\subsection{The Tower Reclaimed}
While it will be some time before the tower can be repaired and brought to a fully operational state again, accomplishing this major objective allows the Academy to move on to the next stage of their plans for Drakkenheim.


\begin{itemize}
\item The Amethyst Academy will permit a player character who is a sorcerer, warlock, or wizard that can cast spells of 6th level or higher and opts to formally join the Amethyst Academy to carry the \textit{Inscrutable Staff} on a mission to Castle Drakken that requires it.
\item The Amethyst Academy offers three very rare magic items in exchange for the \textit{Inscrutable Staff} to groups who wish to keep it. If this offer is refused, they may attempt to steal it back.
\end{itemize}



\subsection{Stolen Tower}
Simply holding the Inscrutable Tower and the staff hostage is enough to stay the hand of the Amethyst Academy in the short term, who are terrified of losing such a precious resource. A cunning foe could extort the Academy for all its worth by doing so. Although Eldrick Runeweaver might enact a last-ditch effort by launching a few schemes to reclaim the tower, if he's defeated, the distant Academy Directorate takes the long view and resorts to appeasement.

\DndQuote%
{}
{''This is why we don't build anything tall in Caspia. The bigger they are the harder they fall...or float.''}
{Pluto Jackson}


\end{DndSidebar}

\section{Saint Selina's Monastery}
This cloistered refuge lies on a small hill overlooking the Drann River in the South Ward of the city. Originally built when Drakkenheim only occupied the north bank, the city streets were built up around the monastery grounds over time. It served as the primary residence for the clergy of the Sacred Flame who tended the cathedrals and chapels throughout Drakkenheim. A few hundred acolytes and priests dwelt in humble comfort within its austere walls and living cells.

Lucretia Mathias and her close followers began gathering here shortly after they completed the first pilgrimages of the Falling Fire. The area is completely covered by the Deep Haze and the nearby waters are contaminated, but those who complete the sacrament are immune to such hazards and can rest in Drakkenheim normally. As such, the otherwise lightly defended abbey is an ideal sanctuary for the most devoted adherents of the Falling Fire. Roughly a hundred sanctified members of the Falling Fire dwell within the monastery proper now, but hundreds more occupy the squalid campsite that has sprung up immediately outside the monastery.

\DndQuote%
{}
{''I swear I can point to this place on a map, but actually getting there is another story.''}
{Sebastian Crowe}
\includegraphics[width=0.4\textwidth]{images/../5e.tools/img/adventure/DoDk/125-07-012.saint-selina-monastery.png}
\subsection{Interactions}
Followers of the Falling Fire live in a tight community and know each other quite well. Everyone here has taken the sacrament.

\subsubsection{Visiting the Monastery}
The Followers of the Falling Fire do not permit anyone who has not completed the Sacrament of the Falling Fire to enter the monastery for any reason. Those who wish to meet with Lucretia Mathias may meet her outside the gates or in the surrounding campsite.

\subsubsection{Defences}
The Followers of the Falling Fire are largely disorganized and live casually within the monastery, safe with the knowledge that few dare approach their sanctuary within the Deep Haze. Nevertheless, a group of six cult fanatics and berserkers are stationed at each of the front gates and passages into the cloisters.

\subparagraph{Blessings of Protection}
Lucretia Mathias has permanently warded the monastery cloisters, tower, and chapel with \textit{hallow} and \textit{forbiddance} spells. As such, when a creature who is not a Follower of the Falling Fire enters the area, it must succeed on a DC 20 Charisma saving throw or become frightened while it remains within. On a success, the creature ignores this effect until it leaves the monastery. These effects cannot be dispelled while the Brazier of the Falling Fire remains lit in the chapel.

\subsection{Area Details}
Saint Selina's Monastery is a broad building composed of two abutting square cloisters with a sixty-foot-tall fortified tower rising between them, and an adjoining chapel of the Sacred Flame. Most of the plastered sandstone building is only twenty feet high with a sloping terracotta roof.

\includegraphics[width=0.4\textwidth]{images/../5e.tools/img/adventure/DoDk/126-map-7.04-saint-selenas.png}
\includegraphics[width=0.4\textwidth]{images/../5e.tools/img/adventure/DoDk/127-map-7.04-saint-selenas-player.png}
\subsubsection{Monastery Garden}
\subparagraph{Monastery Walls}
Rings of tiered lawns and a ten-foot-high fieldstone wall surrounds the grounds. Two wrought-iron gates allow access up a gravel path towards the monastery entrances.


\begin{itemize}
\item 4 cult fanatics and a \textbf{berserker} are posted at each gate.
\end{itemize}

\subparagraph{Gardens}
These were once lush and peaceful gardens, but the plans have now withered and turned to dust. The trees are dead and plants petrified.

\subparagraph{Follower's Tents}
A makeshift encampment of dozens of simple canvas tents fills the garden surrounding Saint Selina's Monastery.


\begin{itemize}
\item During the day, 10d10 sanctified cultists gather around meager campfires set up amongst the dingy tents and stinking hovels here, but at night hundreds huddle close to sleep in truly squalid conditions.
\end{itemize}

\subsubsection{Cloisters}
Each cloister consists of a ten-foot-wide pillared aisle surrounding a square courtyard. Stone benches and long-faded frescoes decorate the tiled path underneath the awning roof.

\subparagraph{Entrance Passages}
A ten-foot-wide covered passageway barred with heavy oaken doors allows entry into each cloister. They are kept open during the day, and barred at night.


\begin{itemize}
\item 4 berserkers led by a converted \textbf{chaplain} keep watch at each entrance at all times.
\end{itemize}

\subparagraph{Courtyard}
The centre of each cloister is a broad lawn criss-crossed with gravel paths. Tall long-dead petrified trees stretch up above the monastery rooftops, and the grass is now merely dead scrub. At the centre of the north courtyard is a covered well that draws water from the city aqueducts, which is surprisingly still functional. The centre of the other is a circular fountain decorated with an angelic statue.


\begin{itemize}
\item 10d6 cultists are praying or meditating in each courtyard.
\end{itemize}

\subparagraph{Monk's Cells}
Along three sides of each cloister is a row of a dozen cells. Accessed via small rickety wooden doors that are seldom locked, these practical stone chambers are tidy and unadorned. Each is no more than ten feet in width and fifteen feet long, and contains two straw mattresses with rough spun sheets, a writing table, a bucket, a satchel with an extra set of clothing for each occupant, and a few candlesticks. Two small square windows look out from each. One row of cells each have a backdoor that leads to individual small disused vegetable gardens ringed by a four-foot-high wall.

\subparagraph{Scriptorium}
The Followers of the Falling Fire have set up two printing presses salvaged from the wreckage of the South Ward in this chamber, which was once the monastery library and scriptorium. They have disposed of the ruined books which once filled the shelves to make space for print moulds, sheets of parchment, bottles of ink, and other printmaking and bookbinding supplies.


\begin{itemize}
\item 12 cult fanatics work around the clock producing printed copies of the \textit{Testament of the Falling Fire} or pamphlets describing the proclamations of Lucretia Mathias. The finished copies are neatly placed in nearby crates and wagons for transport into the wider world.
\end{itemize}

\subparagraph{Common Hall}
Inside this large hall are three rows of long tables set with wooden benches. Disciples of the Falling Fire prepare simple meals of broth and gruel in the kitchen.

\includegraphics[width=0.4\textwidth]{images/../5e.tools/img/adventure/DoDk/128-07-013.abbess-tower-vasburg.png}
\subsubsection{Abbess' Tower}
This sixty-foot-tall tower has four creaking wooden floors. One is now used as accommodations for Lucretia Mathias herself, the others are used for meditation or meetings. Much like her followers, she keeps austere living conditions with a simple straw bed and a writing desk. The rest of the floors contain the scrolls and theological texts moved out of the scriptorium.


\begin{itemize}
\item When she is not personally leading a pilgrimage to the Crater's Edge or delivering sermons in the courtyards, \textbf{Lucretia Mathias} spends her time in the tower composing new writing. She rests here each night.
\item 1d4 very rare Spell Scroll from the cleric spell list are kept here.
\item 6 cavaliers, paladins of the Sacred Flame who converted to the Falling Fire, keep watch over the tower as the personal protectors for Lucretia Mathias. At any given time two are sleeping whilst the others are on patrol, and the other two are posted here.
\end{itemize}

\subsubsection{Chapel of Saint Selina}
This chapel abuts the main cloister and is accessed from the courtyard. It is a round building with a domed roof.

\subparagraph{Chapel}
The brass brazier at the heart is alight with \textit{continual flame}. So long as this brazier remains lit, the \textit{hallow} and \textit{forbiddance} effects warding the abbey cannot be dispelled.

\subparagraph{Crypts}
When a sanctified member of the Followers of the Falling Fire dies, the disciples collect their delerium soulstones and store them in this crypt with the remains of Saint Selina.


\begin{itemize}
\item Lucretia Mathias has summoned two planetars to act as crypt guardians here.
\item The Followers of the Falling Fire keep the \textit{Helm of Patron Saints} here.
\end{itemize}

\chapter{Castle Drakken}
\DndDropCapLine{L}{}ooming over Drakkenheim stands the mighty Castle Drakken. Once a royal palace, military fortress, and administrative hub, the stronghold is now a grim monument to madness and horror. Any who wish to discover what ultimately happened to King Ulrich IV or unearth conclusive evidence to prove a royal claim will find what they are seeking here. Those who come bearing the \textit{Seals of Drakkenheim} may even be able seat a new ruler upon the \textit{Throne of Drakkenheim}, if they can defeat the deadly creatures conjured by the Haze and the former castle guardians left unhinged by its influence.


An unspeakable abomination known as the \textbf{Amalgamation} has taken root in the throne room. However, its physical form is a mere corporeal manifestation of a multidimensional entity that extends into multiple places and planes at once. Feeding on the powerful magic within the foundations of Castle Drakken, the Amalgamation anchors the first major dimensional rift opened by the Haze, which slowly draws extraplanar terrors into the mortal world. Eventually the rift will expand, tearing apart the castle and leaving behind a sickening hole in reality.

\includegraphics[width=0.4\textwidth]{images/../5e.tools/img/adventure/DoDk/129-08-001.throne-room-amalgamation.png}
\section{Overview}
Characters should be at least 11th level before taking on Castle Drakken, and may reach 13th level or higher before defeating all the challenges within. Read this chapter carefully: Castle Drakken is a sprawling and dynamic environment with many interconnected interactions. There is no set order in which characters must progress through the castle! High level characters have many abilities including flight and teleportation, and the following scenarios are designed with these capabilities in mind. Characters lacking these powers may find Castle Drakken significantly more difficult.

\subsection{Adventure Hooks}
\subparagraph{Seize the Throne}
The political future of Westemär rests upon successfully reclaiming Castle Drakken. Though their motives differ, both the Hooded Lanterns and the Queen of Thieves wish to gain entry to the throne room, obtain the \textit{Crown of Westemär}, and install a new monarch using the Seals of Drakkenheim. Meanwhile, the Amethyst Academy, the Followers of the Falling Fire, and the Silver Order all understand that supporting a new monarch means they can gain immense influence over what is done about the delerium in Drakkenheim. If you have more than one potential candidate for the throne, you might have a serious political drama on your hands.

\subparagraph{Investigate the Dimensional Rift}
The dimensional rift signals the dangerous progression of the Haze, but characters might first discover a similar rift within the Inscrutable Tower. In this case, Eldrick Runeweaver may theorize that other such rifts could form within the Haze, especially around areas where structures were infused with powerful magic already. Castle Drakken is such a place, so the possibility of a similar effect happening there simply cannot be ignored.

\subparagraph{Treasures of Castle Drakken}
The Followers of the Falling Fire and the Silver Order need two great treasures found inside Castle Drakken to achieve their final goals. First is the Diamond of Mount Kadath, the crowning jewel of the royal treasure vault. The other is a Relic of Saint Vitruvio: the The Shield of Sacred Flame. Though believed to be in the castle chapel or the royal treasure vaults, the relic is now actually in the royal apartments.

\section{Interactions}
Whether ally or enemy, the factions play their best cards now. In addition, players might find unexpected aid among the horrific creatures who inhabit the castle.

\subsection{Allied Support}
Even the mightiest player characters will need assistance from one or more factions to overcome the monsters and challenges present in Castle Drakken. Factions only place their full resources behind an expedition to Castle Drakken to allied characters who have achieved the following:


\begin{itemize}
\item Demonstrated their trustworthiness and full loyalty to the faction's ideals and goals.
\item Acquired the \textit{Seals of Drakkenheim} and located a potential heir to the throne.
\item Secured the cathedral and/or the Inscrutable Tower.
\item Defeated at least one other rival faction and/or created an alliance with two other factions.
\item Acquired reliable means to protect themselves against contamination.
\end{itemize}

Launching a full-scale siege on Castle Drakken is far beyond the capabilities of any faction. Instead, they propose an elite strike team led by the player characters. Before committing their forces, the faction leaders request the characters thoroughly scout the castle grounds and determine what lurks within the palace, especially the throne room and the royal apartments. Each wants to develop a clear strategy, and notes that securing Castle Drakken in a single expedition is unlikely.

Once characters return from Castle Drakken with a well-developed report that confirms what lies in the throne room, the factions should offer their greatest boons and finest agents to help players put together a team to take the palace. This force could reasonably include a faction leader, one or more faction lieutenants, a group of rival adventurers, and a faction strike team. Many faction leaders can also offer key insight and knowledge about the castle:


\begin{itemize}
\item The \textbf{Queen of Thieves} and Elias Drexel intimately know the layout of Castle Drakken, including the location of any secret passages. They can answer questions about the former inhabitants of the castle and provide specific details about how the \textit{Seals of Drakkenheim} are used to crown a new monarch.
\item \textbf{Eldrick Runeweaver} can provide characters with excellent information about the magical defenses of Castle Drakken. He can study the nature of the dimensional rift and accurately assess what is required to close it down, given a full report.
\item \textbf{Lucretia Mathias} performs powerful divinations to discern the future. Characters may ask her anything about the castle, but she replies in truthful yet cryptic clues in the form of prophetic statements.
\item Theodore Marshal can't offer much information about Castle Drakken, but he commands the finest air force amongst the factions. He can offer characters advice on how to take the castle by the air, noting that there's no reason for the characters to walk in the front gate if they can fly up to the balconies!
\end{itemize}

\subsection{Enemy Interlopers}
Allied faction leaders should warn the player characters that counterstrikes and interference from the other factions are almost certain. When enemy factions discover the player characters are exploring Castle Drakken, they see it as an opportunity to defeat them once and for all. They begin enacting their more dangerous and underhanded schemes, and enemy faction leaders will personally lead an elite strike team into Castle Drakken to confront the player characters.

\subsection{Neutral Parties}
The time to remain on the sidelines has passed. If by some chance the characters kept relatively neutral interactions with one or more of the factions, the leaders of those factions should contact the player characters to arrange a meeting once they learn they are successfully exploring Castle Drakken. The faction leader offers their support, but demands their goals be factored into any plans the player characters might have for the future of Drakkenheim. The outcome of this final interaction will determine if that faction emerges as an ally or enemy.

\subsection{Missing Seals}
By this point in the adventure the player characters should have located the \textit{Seals of Drakkenheim} and learned of their significance. However, if the characters missed any of the seals still lost in the ruins, a group of rival adventurers might find them instead. If enemy faction leaders still possess one or more of the seals, you can have these NPCs reveal this fact during a tense negotiation as a bargaining chip. The players might make a deal, or plot to take the seal by force or stealth. Either way, you should ensure all six seals are "in play" before the player characters head to Castle Drakken.

\subsection{The Royal Steward}
Johann Eisner was the royal steward of Drakkenheim. As the Haze engulfed the castle, he attempted to rescue Ulrich IV, but failed. He fled to the crypts where he became trapped. Now, his body and mind are fused to the Foundation Pillar in the Monarchal Crypt. As a result, he can now tap into the magic within the castle. He dutifully acts as custodian of this damned place to contain the horrors within. However, the steward has endured this torturous experience for fifteen years. It is only a matter of time before he falters and fully succumbs to madness.

Initially, the steward attempts to bar characters' entry to the royal palace and drive them out of the castle. Nevertheless, if characters make an earnest attempt to communicate with him, the steward responds cautiously. Once he believes the characters are capable and well-intentioned, the steward implores them to rescue him from the crypts. Then, he shares his knowledge and plans (see "Eternal Custodian" in the Developments section for details.)


\begin{description}
\item[Personality.]
I pay exceptional attention to the finest details in my work.
\item[Ideal.]
A loyal steward discreetly manages the affairs of state and household unseen.
\item[Bond.]
I pledged my life to serve my liege, my homeland, and my family.
\item[Flaw.]
Others often fail to live up to my exacting standards. I personally perform many tasks I don't trust anyone else to handle properly.
\end{description}

The steward's mind inhabits Castle Drakken itself. He has Intelligence, Wisdom, and Charisma 16, and a +11 bonus on History checks. He speaks, reads, and understands common plus four other languages, and can see and hear anything within sixty feet of the Foundation Pillar. He has the following additional traits:

\subparagraph{Custodian's Eye}
The steward creates an invisible, immobile magical sensor anywhere within Castle Drakken that hovers in the air for one minute. He can see and hear through the sensor as if he was there, and may cast spells as though he were in the sensor's space. The sensor dissipates if it takes any damage.

\subparagraph{Innate Spellcasting}
The steward's innate spellcasting ability is Intelligence (spell save DC 15). When he casts a spell of 1st level or higher, roll 1d6. On a 1, it triggers an arcane anomaly. He can innately cast the following spells, requiring no components:


\begin{itemize}
\item At will: \textit{mage hand} (the hand is invisible), \textit{thaumaturgy}, \textit{unseen servant}.
\item 3/day: \textit{animate objects}, \textit{dispel magic}, \textit{fear}, \textit{telekinesis}, \textit{wall of force}
\item 1/day: \textit{forcecage}, \textit{reverse gravity}, \textit{symbol}
\end{itemize}

When using his magic, the steward may only leave cryptic messages of one or two words.

\section{Sealed Fate}
At the DM's discretion, characters who possess a signet ring or token of some kind once owned by a member of the royal family (such as the \textit{ring of protection} worn by Lenore) can also bypass castle defenses as if the item were a \textit{Seal of Drakkenheim}.

\section{Castle Grounds}
Castle Drakken is a sprawling palatial fortress built upon a two-hundred-foot-tall stepped plateau surrounded by rocky cliffs. The castle grounds are separated into two baileys by an internal dividing wall: the Lower Bailey contains administrative buildings, a military barracks, and servants quarters. The Upper Bailey grounds are elevated thirty feet above the Lower Bailey, feature the castle gardens, a chapel, and access to the royal palace of Castle Drakken. The palace itself perches atop the highest rocky pinnacle approximately two hundred fifty feet above the city proper, its imposing silhouette visible everywhere in the city below.

\includegraphics[width=0.4\textwidth]{images/../5e.tools/img/adventure/DoDk/130-map-8.01-castle-drakken.png}
\includegraphics[width=0.4\textwidth]{images/../5e.tools/img/adventure/DoDk/131-map-8.01-castle-drakken-player.png}
\subsection{Arcane Wards}
Generations ago, House von Drakken commissioned the Amethyst Academy to weave a powerful magical ward upon the castle grounds and palace that remains in place today:


\begin{itemize}
\item Doors, walls and windows are magically reinforced, and are immune to fire, poison, and psychic damage, as well as non-magical bludgeoning, piercing, and slashing damage from attacks. They are resistant to all other damage. The doors, walls and windows extend into the ethereal plane, blocking ethereal travel through them.
\item A radiant purple glow emanates from every window blocking line of sight from outside.
\item Creatures outside castle buildings can't teleport into any interior areas, nor can they teleport within the castle rooms unless they have a direct line of sight to their destination. A character attuned to a \textit{Seal of Drakkenheim} ignores this effect.
\end{itemize}

\subsection{Tower Dragons}
The bronze dragons upon the towers that peer down with serpentine necks are no mere decorations. Castle Drakken is defended in a similar manner as the city walls themselves (see chapter 7), but the guardians are much more aggressive and hostile. 2 (1d4) tower dragons and 10 (3d6) wall gargoyles animate and attack characters who dare attempt any of the following:


\begin{itemize}
\item Travel up the roadway past the Cliff Gate on foot.
\item Climb more than halfway up the cliffs, walls, or towers.
\item Fly within three hundred feet of the outer walls horizontally, or over them at any height.
\item Enter the bailey courtyards with teleportation magic.
\end{itemize}

The guardians relentlessly pursue their quarry within the castle grounds, but do not chase characters more than three hundred feet away from the castle grounds. They do not attack characters attuned to a \textit{Seal of Drakkenheim}.

\subsection{Outer Walls}
Stone masonry walls forty feet high and ten feet thick ring the clifftop, anchored by several bastion towers each sixty feet tall. Like the walls that surround the city, they are richly decorated with architectural details and menacing gargoyles.

\subsection{[1] Gates}
A gently sloped cobblestone road twenty feet wide and lined with stepped battlements winds up the two-hundred-foot-tall cliff face from the Inner City towards Castle Drakken. Three gates are set along its path: The road begins from the Lower Gate, passes through the Cliff Gate halfway up the hill, then finally terminates at the Upper Gate that stands atop the cliff.

Each gate features a pair of bastion towers, double portcullis, and a drawbridge. Unfurled and torn banners bearing the coat-of-arms of House von Kessel hang from each tower. Though more ornamented than the city gates, their internal layout is quite similar.

\subparagraph{[1a] Lower Gate}
The lower gate of Castle Drakken rests at the base of the castle cliff. A broad city street meets the gate, and the surrounding neighbourhood is made up of noble estates and wealthy homes now occupied by ghosts and mutated dregs.


\begin{itemize}
\item The last regiment of the royal guard held the gate when the meteor fell. The petrified forms of the royal guard and a panicked throng of common folk are arranged here in a horrific tableaux. Their features are melted and warped, and several animate into 10 (3d6) gibbering mouthers with flesh like wet concrete when approached.
\item Several guards who held fast against the gates have merged into it as a huddled mass. Their limbs are elongated and their faces stretched in odd otherworldly shapes that spring to life as 2 (1d4) ropers. The gate will not open unless they are destroyed.
\end{itemize}

\subparagraph{[1b] Cliff Gate}
Two stone golems defend this gate.

\subparagraph{[1c] Upper Gate}
These gates are built into the main castle walls and are sealed shut, opening only for one who bears a Seal of Drakkenheim.

\DndQuote%
{}
{''It stands over a city lost to chaos. I intend to make it stand again for something.''}
{Pluto Jackson}
\subsection{Random Encounters}
Monsters wander the castle grounds. Check for random encounters each hour characters spend exploring Castle Drakken as described in chapter 5 using the table below. In addition, check for random encounters when characters enter certain buildings or rooms, as indicated by the area description:


\begin{DndTable}{cX}

1d8 & Random Encounter\\
1 & Castle Guardians. 2 (1d4) tower dragons and 10 (3d6) wall gargoyles (ignore this result for interior locations)\\
2 & Formless Horrors. 3 (1d6) protean abominations and 10 (3d6) gibbering mouthers\\
3 & Spectral Housekeepers. 3 (1d6) warp witches\\
4 & Wretched Workers. 5 (2d4) haze hulks and 10 (3d6) delerium dregs\\
5 & Royal Guards. 10 (3d6) haze wights leading 14 (4d6) haze husks\\
6 & Roving Demons. A \textbf{glabrezu} with 3 (1d6) vrocks\\
7 & Wandering Devils. A \textbf{horned devil} with 10 (3d6) bearded devils\\
8 & Illusions of Grandeur. All the furniture and objects in the room are mimics\\
\end{DndTable}

\subsection{Deep Haze}
The entire castle is blanketed by a thick and insidious manifestation of the Deep Haze. Blinking motes of purple and red light bounce through the air, and flakes of delerium flutter like ash within.

\subsection{Delerium Deposits}
Massive clusters of delerium shards have burst through the earth throughout the castle grounds. Crystals creep up the outer walls like a cancerous moss, surrounded by thick gobs of congealed Haze and bubbling delerium sludge. Corpses writhe and wriggle within, mouths and eyes forming at random.


\begin{itemize}
\item Though the massive clusters are too large for field extraction, characters can pry loose 1d4 crystals or 2d6 delerium shards from each.
\item This delerium was not left by the meteor; it has since grown in the Haze.
\end{itemize}

\subsection{Endless Horror}
Monsters inhabiting Castle Drakken are rejuvenated each night at midnight by the Haze, with replacements conjured for any banished or wholly destroyed creatures. This effect ends if the \textbf{Amalgamation} in the palace throne room is destroyed. However, if the dimensional rift is not then sealed (see the throne room) the Amalgamation itself rejuvenates 1d10 days later, and the monsters return shortly thereafter.

\subsection{[2] Lower Bailey Areas}
The dusty and gravel-filled lower bailey has a fine bricked road leading through it from the front gate to the bailey gate. The ground is ashen and cracked, delerium clusters jutting from the earth. Check for a random encounter each time the characters enter one of the buildings below:

\subparagraph{[2a] Armoury}
This fortified warehouse houses cannons, siege engines, and esoteric weaponry, as well as a massive stockpile of conventional armour and armaments.

\subparagraph{[2b] Barracks}
In the middle of the bailey is a fortified building with several crenellated parapets and a high steepled roof. Outside are unfurled banners displaying the sigil of the royal guard. The ground floors contain officers quarters, a mess hall, and kitchen, while the upper floors are comfortable if spartan accommodations for the royal guards. The topmost level has private chambers for the guard captain.

\subparagraph{[2c] Carriage House}
This large garage houses a dozen royal carriages, wagons, and vehicles.

\subparagraph{[2d] Servant's Lodge}
This bending building housed accomodations for the common castle staff. It also features a large kitchen, bakery, well, and laundry facilities. The adjoining tower serves as a granary and dedicated storehouse for castle provisions including grain and preserved foods.

\subparagraph{[2e] Stables}
This massive L-shaped stable once sheltered dozens of horses. It is now a charnel house of bones and decaying leather; the noble steeds within were consumed by the castle's monstrous residents long ago. The loft above the stables contained a residence for the master of horses and the royal stablehands. The adjoining tower is the castle aerie, which was roost for the royal giant owls. In addition to a possible random encounter, six manticores live here now.

\subparagraph{[2f] Workshops}
A large outdoor forge, with several anvils and workstations. An adjoining shed contains raw materials: wood, iron ingots, finished nails, and all manner of tools.

\subparagraph{[2g] Bailey Gate}
The gate passage covers a curved ramp leading to the Upper Bailey.

\subsection{Prince of Carnage}
A wingless \textbf{balor} known as the Prince of Carnage has made the courtyard his territory, its body covered in lacerations and exposed raw muscle. Conjured by the Haze and bound to the castle grounds, the Prince of Carnage is an insane demon who believes himself to be a royal prince. He furiously attacks any who trespass the lower bailey, barring the way to the upper bailey in delusional claim over the ruined realm.

\subsection{Servant's Passage}
A prominent set of doors in the basement of the \textbf{servant's lodge} opens into a narrow stone passage that leads under the Bailey Gate and connects to the palace proper. The passage isn't concealed: it was used daily by the castle staff to get to work in the palace, especially in the winter. However, the doors of the passage were sealed with an \textit{arcane lock}. They open for anyone bearing a \textit{Seal of Drakkenheim} or who speaks the password "Giggleshorts".

Characters interacting with the warp witches and delerium dregs found in the castle grounds may be able to learn the password from them.

\subsection{[3] Upper Bailey Areas}
The regal Upper Bailey is a stark contrast to the Lower Bailey. Fields of dead grass, large contaminated pools, dried-out fountains, baroque statuary covered in ash, and twisted trees line a fine brick roadway leading to the palace entrance.

\subparagraph{[3a] Overlook Tower}
The north tower connects to a small exterior doorway large enough for humans to pass through. It allows access via a leisurely (if exhausting) walk down the cliff via a graceful stone staircase that leads to the ravine below. Pillared bannisters, exquisite statues, luscious fountains, and stone benches decorate the many porches, overlooking platforms, and balconies along the way.

\subparagraph{[3b] Knight's Tower}
This tower manor has small, if comfortable, accommodations for the knights in service to the monarch. Many levels contain several individual rooms, whereas others are whole apartments for particularly favoured champions.

\subparagraph{[3c] Castle Gardens}
A field large enough to set up a great outdoor tent for garden parties, the gardens contain fountains, statues, hedgerows, and decrepit flowerbeds.

\subsection{Castle Chapel}
This small chapel of the Sacred Fire is dedicated to an obscure and ancient saint and lavishly decorated for its size with angelic sculptures. However, the crypts beneath are unfinished, the additions intended as a reliquary for the The Shield of Sacred Flame left incomplete.


\begin{itemize}
\item Roll or choose a random encounter the first time characters enter the chapel.
\item An altarbox left beside the brazier contains 1d4 Potion of Superior Healing, plus scrolls of \textit{greater restoration}, \textit{heal}, and \textit{raise dead}. Inside is a letter from the steward:
\end{itemize}

\DndQuote%
{}
{''Flamekeeper Maresa—My apologies for the construction delays. Lest you doubt the piety of our liege, rest assured his majesty has resolved to keep the shield of Saint Vitruvio over his mantleplace so he may closely meditate upon the Flame nightly until the new reliquary is complete.''}
{Johann Eisner, Royal Steward Royal Palace Areas}
Though outwardly imposing, the interior rooms of Castle Drakken contain luxurious chambers featuring grandeur beyond reckoning. These sumptuous rooms have been darkly twisted by eldritch forces; once lively halls shifted into destatured tones and sickly hues. Characters can enter the palace through the main door, the servant's passages, or by flying onto the balconies.

\includegraphics[width=0.4\textwidth]{images/../5e.tools/img/adventure/DoDk/132-08-002.metal-box.png}
\section{Approach}
\begin{DndReadAloud}{}
The royal palace of Castle Drakken perches atop the highest pinnacle overlooking the city below. The heart of the citadel is a massive central keep connected to four great round towers via soaring bridges, which rise above two large stone chateaus each bedecked with parapets. The imperial building is topped with slate rooftops of midnight blue and elongated spires sharp like spear points. The outer walls feature immaculate stone moldings, overhanging balconies, flying buttresses, and tall arched windows with intricate stained glass designs. Upon the battlements are leering gargoyles and roosting dragons cast in bronze.

The main entrance is a fortified gate flanked by squat square towers. Engraved into each of the heavy iron double doors is a dragon rampant set with gold filigree. Above is a circular rose window fixed with delicate tracery. Beams of sickly purple light radiate from within. Gusts of sweeping wind kick up the smell of grave-dirt and ancient stone, and thunder crackles through the air.

\end{DndReadAloud}

\subsection{Minazorond}
An \textbf{ancient bronze dragon} named \textbf{Minazorond} is eternally bound to serve as the guardian of Castle Drakken. He roosts atop the Stair Tower of the royal palace with his mind awash in nightmares. The dragon now believes the monstrous creatures dwelling in the castle to be its proper inhabitants, and perceives humanoid creatures as trespassers to be destroyed. Minazorond awakes with a thunderous roar that can be heard for miles away, and attacks if intruders approach within one hundred fifty feet of the royal palace by any means. The dragon does not speak, nor respond to reason or parley. He does not pursue his quarry beyond the castle grounds. However, Minazorond is too large to enter buildings or the palace, and cannot pursue creatures who flee inside.

\textbf{Minazorond} possesses the additional traits of a \textbf{tower dragon}, and may be rebuked by a creature bearing a \textit{Seal of Drakkenheim} in a similar manner. Only a \textit{wish} spell can cure his madness: he interacts as a dutiful, stoic, and wise protector thereafter.

\section{Main Entrance}
\subparagraph{Front Gates}
The grand double doors are closed and barred, and the portcullis lowered. A creature within thirty feet can use an action to present a Seal of Drakkenheim, which magically opens the doors.

\subparagraph{Vestibule}
The entrance hallway of the royal palace is adorned with grand tapestries proudly displaying the House von Kessel arms. Tight rows of arched columns form aisles on either side of this room. Clerestory windows allow light to illuminate the room.

\subparagraph{Guardrooms}
These unadorned stone chambers have narrow arrowslit windows. Weapon racks with bows, crossbows, and arrows, and spartan bunks furnish each. A ladder allows access between three levels with wooden floors and a winch room to operate the portcullis.


\begin{itemize}
\item The first time characters enter the palace, a force of 12 haze wights stands guard at the entrance. Thereafter, check for random encounters each time characters pass through the main entrance.
\end{itemize}

\subsection{Palace Interior Features}
The following details are common to all rooms in the royal palace except the basement, kitchen, and hidden passages.

\subparagraph{Chambers and Corridors}
Castle Drakken is lavishly and ostentatiously decorated. The castle rooms are each individually themed around a different pantheon of mythological heroes, gods, and monsters depicted in scenes from their greatest stories. You are encouraged to choose your favourite characters (even those from real-world myth) and place their likenesses throughout the castle art.

\subparagraph{Amenities}
Running water fountains, grand fireplace hearths, and coal stoves once provided comfort throughout the castle. The pipes and chimneys are now choked with contaminated water and Haze-infused soot.

\subparagraph{Artwork}
Massive painted murals and framed canvases adorn the walls and ceilings, and lifelike statues of carved marble rest on pedestals in many rooms. However, the once expressive and idealistic fairy-tale scenes now have disturbing undertones, bizarre contrasting colours, distorted features, warped perspectives, and oblique angles.

\subparagraph{Ceilings}
Most rooms and corridors have vaulted ceilings thirty feet high. All are decorated with frescoes and carved mouldings. Some truly elegant rooms feature a dome painted with a stunning circular panorama.

\subparagraph{Doors}
These paired doorways of reinforced wood have embossed panels and ornamentations.

\subparagraph{Floors}
Tiles of marble, porcelain, or stone are laid out in intricate arrangements. Many rooms also have large woven carpets bearing geometric patterns and golden tassels and trim.

\includegraphics[width=0.4\textwidth]{images/../5e.tools/img/adventure/DoDk/133-08-003.royal-ring.png}
\subparagraph{Furnishings}
These elaborate antiques are upholstered with fine velvet and painted with gold.

\subparagraph{Lighting}
Brass and gold chandeliers strung with crystals and kept alight by \textit{continual flame} illuminate each chamber.

The Haze has transformed the once soft and golden firelight into unsettling hues of crimson, purple, and green.

\subparagraph{Secret Doors}
Castle Drakken contains many secret doors. These may be found with a DC 20 Investigation check. Each wall section rotates on a central axis; pushing them open takes an action and a DC 10 Strength check. They automatically close afterwards.

\subparagraph{Walls}
Sculpted borders and baseboards in marble, ivory, or painted wood frame the walls. Many feature mascarons and leering faces with draconic designs. The rare areas without windows or paintings often feature intricate patterned tapestries or wallpaper.

\subparagraph{Windows}
Paned glass windows are decorated with tracery and stained glass, and hung with curtains of silk and sable.

\subsection{Looting the Castle}
The artwork and furnishings in each room of Castle Drakken are extraordinarily valuable. Any given chamber contains 1d6 objects worth 250 gp each and 1d4 objects worth 750 gp each. However, transporting and fencing these items may be a challenge.

\subsection{Chateau Halls}
Ceilings in the castle chateau reach a height of forty feet, though the attic above extends the roof twenty feet higher beyond.

\subparagraph{Central Passage}
Artwork and statues in this columned hall depict the past monarchs of Drakkenheim. Most belong to the original ruling family, House von Drakken, with only the most recent ones belonging to House von Kessel. Stairs on either end connected to the chateau basement, beside sets of double doors lead to the palace entrance and the antechamber.

\subparagraph{Banquet Hall}
This room features two lines of feasting tables set with tablecloths and golden candlesticks. Porcelain plates, silver utensils, and napkins provide table settings to accommodate a hundred people.

\subparagraph{Dining Room}
This dining room is richly decorated with far more lavish settings than the banquet hall, as it seats only twenty four important guests.

\subparagraph{Drawing Room}
This well-appointed salon contains several comfortable wingback chairs, settees, and side tables. A delicate cabinet is stocked with exquisite spirits and fine tobacco.

\subparagraph{Grand Ballroom}
This splendid long hall has grand mirrors opposite the windows, filling the walls at each end and incorporated into the ceiling, creating the illusion of a vast and almost endless interior space. Gold-painted pillars, glass and silver chandeliers, and flowing curtains inlaid with fine embroidery complete the room. Upholstered benches and dancing statues rest in alcoves on each side of the room. A small platform at one end has a collection of chairs, music stands, and antique stringed instruments.


\begin{itemize}
\item Hundreds of intangible spirits of long-dead nobles manifest in these chambers, playing out the last masquerade ball held in Castle Drakken. Meanwhile, three dozen haze husks carry out their duties as serving staff, bringing spoiled wine and foetid hor'dourves to the "guests". Characters passing through the chamber are treated like familiar and celebrated guests, and invited to share their heroic tales and partake in the dance. Characters can pass through the chateau unopposed so long as they humor the spirits; but should they shatter the delusions, the spirits angrily vanish save for 10 (3d6) hostile warp witches.
\item Check for a random encounter each time characters pass through the chateau thereafter.
\end{itemize}

\subsection{Chateau Basement}
\subparagraph{Kitchen}
This gigantic kitchen could prepare a banquet for hundreds, and contains several ovens, stoves, and large preparation tables. It reeks of offal, but is disturbingly tidy otherwise. Cabinets and shelves with spices, cooking utensils, plates, and spoiled stores line the walls.


\begin{itemize}
\item The three renowned castle chefs - Frier, Olivier, and Ramsay - still remain in the kitchen. They have transformed into hideous creatures imbued with magic power (use the game statistic of a \textbf{night hag} but their creature type is aberration). The three chefs despise each other yet feed off their rivalry, and assume anyone entering the kitchen who isn't openly wearing a \textit{Seal of Drakkenheim} is an underling they can boss around. The chefs furiously reprimand intruders for being "late", order them to get to work, then berate and insult the quality of any tasks performed. They attack if they aren't obeyed.
\item Each time characters enter the kitchen, check for a random encounter. Any creatures that might be encountered are working (poorly) in the kitchen and being abused by the chefs.
\end{itemize}

\subparagraph{Larders}
These three chambers contained the food supplies for the palace upon and within shelves and containers, almost all of which were eaten or contaminated. One stored dry goods at normal room temperature, but the other two are magically enchanted with temperatures of 40 and -180 fahrenheit.

\subparagraph{Basement Storeroom}
These dusty storerooms contain extra furnishings, textiles, and utensils used during parties and galas held in the palace, neatly stacked and cleanly organized. Behind an upturned table is a secret door to a short passage connected to the Dungeon Vaults.

\subparagraph{Wine Cellar}
Racks containing hundreds of bottles of fine wine and spirits organized in a byzantine fashion line this cellar. One wine rack conceals a secret door opening to a long passage that leads to the Bureaucratic Archives in the Steward's Tower.

\subparagraph{Scullery}
This room smells of soap, and has large basins for washing laundry connected to taps of running water.

\subparagraph{Privvies}
This washchamber contains rows of privvies used by castle guests.

\subparagraph{Side Entrance}
These broad steps include a gently sloping ramp that leads outside the palace.

\subparagraph{Servants Staircase}
This spiral staircase has an integrated dumbwaiter elevator. It leads directly to the lowest level of the royal apartments.

\subparagraph{Servant's Passage}
This door is \textit{Arcane Lock} and opened via the servant's password. The narrow tunnel leads to the Servant's Lodge in the Lower Bailey.

\subparagraph{Water Spring}
Water flows into this large pool from several narrow pipes connected to the city aqueducts. The once-fresh water is contaminated, inhabited by two animated delerium sludges.

\subparagraph{Canal}
Water from the cistern pool flows into a downward-sloping canal with a small cleated pier for mooring a small boat or raft. A portcullis bars passage further; the canal terminates at another spring in Queen's Park. The gate opens only for a character bearing a Seal of Drakkenheim.

\subsection{Throne Room Antechamber}
This ceiling of this magnificent antechamber rises sixty feet overhead, the top bordered by arched clerestory windows that allow natural light into the chamber. Two sets of broad marble steps with pillared banisters wrap around the room and converge at a central columned balcony thirty feet above the ground floor. Golden double doors each etched with the design of a dragon rampant upon a tower are set atop and below the balcony passage below it, and broad landings halfway up each staircase connect to halls heading north and south.


\begin{itemize}
\item Check for a random encounter each time characters pass through this chamber.
\item The steward uses his spells such as \textit{wall of force} to prevent characters from entering the throne room.
\item This chamber is the ground level of the central tower, but does not directly connect to the royal apartments.
\end{itemize}

\includegraphics[width=0.4\textwidth]{images/../5e.tools/img/adventure/DoDk/134-map-8.02-castle-drakken-towers.png}
\includegraphics[width=0.4\textwidth]{images/../5e.tools/img/adventure/DoDk/135-map-8.02-castle-drakken-towers-player.png}
\subsection{Stair Tower}
A brilliant chandelier extends down the height of this hollow one-hundred-fifty-foot-tall tower. Thousands of glass crystals with frames of silver and gold are arranged in patterns of constellations and stars. A wide and impressive spiral staircase glides up and down the length of the tower lined by a carved balustrade, exiting to the tower roof. There are four landings along the length of the stair: three connected to arched bridgeways. The stairs also descend to a lower level, where the sparkling chandelier light reflects from a large fountain at the tower's base.


\begin{itemize}
\item The stairs connect to the Ambassador's Tower, the Steward's Tower, and the Royal Apartments, the Antechamber, and the Vaults.
\end{itemize}

\subparagraph{Dragon Roost}
The stairs continue and exit the top of the tower to the tower roof, which is a flat platform surrounded by battlements. \textbf{Minazorond} uses this place as a roost.

\subsection{Conservatory Tower}
Beyond the tower base, this one-hundred-fifty-foot-tall tower narrows considerably to a width of twenty feet as the spiral staircase makes its way upward.

\subparagraph{Library}
The musty smell of ancient books wafts through this room, lined floor to ceiling with shelves with a sliding ladder rail. This impressive collection includes rare works of poetry, literature, and history. A loft above contains a study where the monarch's children were instructed by a tutor.

\subparagraph{Mage's Chamber}
An \textit{Arcane Lock} door on the middle level opens only for members of the Amethyst Academy. Inside is a private suite inhabited by 4 arcane wraiths, the insane remnants of the last Academy representatives posted to Castle Drakken. Amongst their belongings are a rare magic item plus 1d4 rare potions and 1d6 rare Spell Scroll of your choice.

\subparagraph{Solarium}
The topmost level is a fine glass-paned sitting room with perhaps one of the finest views of the landscape north of Drakkenheim. Elegant silver furniture with spiraling design and plush pillows provides a place for quiet relaxation.

\subsection{Steward's Tower}
This one-hundred-foot-tall square tower connects to the Stair Tower via an arched and covered bridge. It can be accessed from the outside via the balconies and battlements.

\subparagraph{Private Quarters}
The top two levels of the tower are the private quarters for the steward and his family, consisting of a drawing room, personal study, and three bedrooms. The furnishing is comparatively modest next to the ostentatious decoration of the rest of the castle.

\subparagraph{Steward's Office}
The walls of this otherwise ordinary room are covered floor to ceiling in cabinets, bookshelves, chests of drawers, and a still-ticking grandfather clock. A massive oak desk covered in heavy tomes and papers dominates the room.


\begin{itemize}
\item Upon the desk is an unsent letter written in the steward's hand and bearing his seal addressed to Ulrich IV. It contains revealing details about an illegitimate heir to House von Kessel, including their name and possible location. The exact contents are up to you to decide based on how you've arranged the royal succession.
\end{itemize}

\subparagraph{Bureaucratic Archives}
This cramped room has narrow aisles between tall rows of drawer cabinets and shelves. They contain detailed records of the city; including property deeds, noble titles, tax records, construction permits, royal decrees, diaries, succession laws, and years of correspondence and records. Any key information or clues helpful to the characters, factions, or their personal quests could conceivably be found here, but finding any specific document or information kept here would take days of searching, possibly even weeks. However, making a DC 25 Intelligence (Investigation) check allows a character to find one specific document with a one-hour search by piecing together how the material here is actually organized.

A secret door in the archives leads to a rough tunnel that connects to the wine cellar in the Chateau Kitchen.

\subsection{Ambassador's Tower}
This tower is one hundred twenty feet tall and connects to the Stair Tower via a bridge.

\subparagraph{Reception Chamber}
This luxurious lounge served as a meeting area for more informal conversations between dignitaries.

\subparagraph{Guest Apartments}
The remaining levels of this tower contain several fine suites for important visitors and political hostages of the royal family with private bedrooms and common rooms.

\subsubsection{The Emissary's Entourage}
A scheming and sauve \textbf{erinyes} known as the Emissary occupies the Ambassador's Tower with several other devils: two horned devils and a bodyguard of 12 bearded devils. The devils desperately wish to leave the mortal world and return to Hell (without being destroyed in the process). They believe that if they were to acquire mortal souls or a \textit{Seal of Drakkenheim}, their hellish overlord will relent and pull them back to Hell. They tell the following diabolical lies:


\begin{itemize}
\item "We were summoned here by the will of the throne to receive an important message for our master. We are bound here until one who sits the throne gives us our mission."
\item "Ulrich IV lives. His horrific transformation may be undone by destroying the corrupted pillar in the crypts, but a powerful guardian protects it now."
\item "The steward is in fact a powerful lich. He conspired with demons to summon the meteor. He has locked himself in his tower. Do not listen to anything he says."
\item "The Prince of Carnage is Leonard von Kessel."
\end{itemize}

\subsubsection{The Devil's Bargain}
\begin{DndReadAloud}{}
"Our master is the Lord of Hell. All secrets are known to them. If you offer us your mortal soul or one of the Seals of Drakkenheim, we can contact them. Then, you may ask them three questions, and you will receive a truthful reply in kind."

\end{DndReadAloud}

The Lord of Hell possesses true knowledge, and the deal is genuine. Characters could learn about the following:


\begin{itemize}
\item Identify a royal heir or where proof of lineage may be found.
\item Locate another \textit{Seal of Drakkenheim} or a \textit{Relic of Saint Vitruvio}.
\item Provide information needed to complete a personal quest.
\item Gain knowledge about one or more new spells in this book.
\item Explain the true origin of the meteor and the effects of the Haze.
\end{itemize}

A character who offers their mortal soul withers and dies immediately, and can't be raised unless their soul is freed. If a \textit{Seal of Drakkenheim} is offered to the Emissary, it disappears in a puff of brimstone when placed in her hand, lost in the Lower Planes. Then the devils answer the first three questions they are asked truthfully, and are subsequently transported back to Hell.

\subsection{Royal Apartments}
The two-hundred-foot-tall central tower of the palace contains the royal apartments. It is accessible by the Servant's Stair from the Basement Kitchens, or by a bridge from the Stair Tower. The royal apartments are spread over four floors. Each has a similar layout with a central hall and six rooms, connected by a staircase that runs across the floors.

\subparagraph{Attendant's Room}
These small rooms and closets were used by servants waiting on the royal family, and are connected via the Servant's Stair to the Basement Kitchen. They contain fine shelves and cabinets used to store linens, cleaning supplies, cutlery, a small pantry, wine rack, and a cot for servants working the night shift.

\subparagraph{Baths}
Two rooms in the apartments contain large, elegant porcelain clawfoot bathtubs. Levers beside gold-painted spouts cause water to fill the tub when a stopper is placed over the drain.

\subparagraph{Bedchambers}
There are a dozen bedchambers across the four floors of the royal apartments. Each has a large sumptuous bed with a satin canopy, silk sheets, plush pillows, and soft velvet carpets, and a hearth or heating stove. The carved headboards have gold filigree and embedded pearls. Candelabras, chandeliers, and nightstands are made of gold or silver.

\subparagraph{Dining Room}
The royal family took most meals in this elegant dining room. The carved dining table has a glass top and is a work of art in its own right.

\subparagraph{Dressing Rooms}
About half the bedchambers have an adjoining dressing room appointed with broad wardrobes, dressing tables, and large framed mirrors. Lockboxes with wigs, jewellry, trinkets, and rotting makeup are scattered throughout.

\subparagraph{Water Closets}
Three curious rooms in the apartment floors each contain a porcelain bowl set with a cushioned hoop-shaped seat and attached to a tank of water. Pulling a nearby metal chain causes water from the tank to flood the bowl and carry waste away into the castle plumbing. The water-closet on the penthouse level is plated with gold.

\subparagraph{Sitting Room}
Velvet couches and two great leather wingback chairs flank either side of an impressive fireplace. Thick carpets with intricate designs are laid out beneath, and a square polished marble table with streaks of gold is set in the middle of the room. A warm fire illuminates the room.

\subparagraph{Penthouse}
The topmost floor is the private penthouse used by the monarch and ringed by battlements with an impressive view. It includes an impressive bed, and the rest of the level includes its own bathroom, water closet, short hallway, dressing room, and study. The Regent (see below) has made this chamber his personal lair, and it's now in a state between lavish and squalor. The The Shield of Sacred Flame is here.

\subparagraph{Royal Study}
This seldom used chamber nevertheless has an impressive desk and bookshelves lined with aspirational reading material.

\subparagraph{Rooftop Garden}
The battlements have observational parapets and flower boxes filled with lavenders that housed a royal apiary and falcon's roost. These majestic creatures have been slaughtered, and now eight hypnotic eldritch blossoms occupy the flowerbeds. 1d4 vrocks may also be fluttering around the spire.

\subparagraph{Escape Chute}
The fireplace in the sitting room and the penthouse bedchamber has a false back. When pressed inward and slid upward, it opens into a narrow drop chute with a metal sliding pole. Sliding down the pole leads to the Dungeon Vault.

\subsubsection{The Regent's Retinue}
A gluttonous \textbf{nalfeshnee} known as the Regent conducts depraved and debauched acts in the royal apartments with his retinue of demons: a \textbf{marilith}, 1d4 \textbf{glabrezu}, 1d6 vrocks, and dozens of servile quasits. These creatures can be encountered anywhere in the royal apartments, and relish in transforming the regal appointments into a place of gory grandeur. They have no desire to leave Castle Drakken. Instead, they want to gain control of the dimensional rift to summon more demons. They believe they can accomplish this by stealing mortal souls and drawing power from holy relics. They tell the following demonic lies:


\begin{itemize}
\item "We were summoned here by the will of the throne, and are bound to keep the castle. Only one who sits upon the throne may discharge us from our servitude."
\item "Ulrich IV lives. His mind is bound to the disfigured flesh before the throne, and his will alone contains the Haze. If you destroy him, it will consume the world."
\item "The meteor was in fact a diabolical vessel that crashed here. The devils arrived in Drakkenheim from a portal to Hell in the crater."
\item "The Rat Prince is Leonard von Kessel."
\end{itemize}

\subsubsection{Demonic Sacrifice}
\begin{DndReadAloud}{}
"The king is not strong enough to defeat the Haze on his own. However, we will show you how to help him, but it requires you sacrifice your mortal life... and give up the Relics of Saint Vitruvio."

\end{DndReadAloud}

The demons use the guise of civility to lure characters into a trap where they can use them as a living sacrifice or steal one of the \textit{Relics of Saint Vitruvio} from them. Though they explain the rite will break the power of the Amalgamation, it's a total lie. Instead, the demons gleefully murder any characters who submit to them and use the relic in their profane summoning rites.

\DndQuote%
{}
{''Always wanted to see the inside of Castle Drakken when I was young. Now not so much.''}
{Sebastian Crowe}
\includegraphics[width=0.4\textwidth]{images/../5e.tools/img/adventure/DoDk/136-08-004.bronze-dragon.png}
\section{Throne Room}
\begin{DndReadAloud}{}
This sovereign chamber has a grandiose vaulted ceiling sixty feet overhead and is ringed by magnificent columns supporting a stately balcony halfway up the height of the room. Throughout are countless detailed architectural ornamentations of gold inlaid with emeralds that accent Intricate painted murals depicting regal scenes of glorious conquest. Fleshy tendrils entwine the pillars, twitching membranes stretch across the room, and massive delerium crystals burst through the marble tiled floor. The tall arched windows that lined the walls and galleries have burst inward; whatever intricate designs they once bore shattered amongst pieces of sharp glass floating in the air around each. Instead, the openings look out upon bizarre vistas.

\end{DndReadAloud}


\begin{itemize}
\item The throne room is accessed via the royal antechamber. The rooftop attic rises twenty feet above that outside, and from outside the building, the windows appear intact and opaque.
\item Stumbling around the balcony are 14 (4d6) delerium dregs in shredded courtier's outfits.
\end{itemize}

\subsection{The Throne of Drakkenheim}
\begin{DndReadAloud}{}
Upon a three-tiered marble platform surrounded by a balustrade and flanked by warped statues of mighty knights is the Throne of Drakkenheim. This imperial throne is fixed into the structure of Castle Drakken, and fashioned from granite and embellished with wrought-iron, marble, and golden filigree. Its broad claw-shaped armrests are embedded with gemstones. The cushions are upholstered in rich hues of deep green.

A colossal dragon skull adorns the top of the chair, and skeletal wings rise up from the throne. These once framed a breathtaking circular rose window decorated with prismatic tracery, but instead forms a turbulent eldritch portal beyond which lies an unfathomable landscape.

\end{DndReadAloud}


\begin{itemize}
\item Above the throne is the dimensional rift, which leads to the Space Between Worlds. This churning chaos is the source of all magic and metaphysically connected to every other dimension, plane, and reality. Characters may enter the Space Between Worlds through this portal, but are struck by \textit{feeblemind} (DC 20) in the process. Those who fail this save are driven mad. See the "Developments" section for details on what lies beyond.
\item The throne is a magical artefact in its own right and indestructible while the castle stands.
\item While seated on the throne, a character attuned to the \textit{Crown of Westemär} may use an action to cast the following spells:
\begin{DndReadAloud}{}
At-will—\textit{command} (DC 20)

1/day each—\textit{globe of invulnerability}, \textit{mind blank} (self only), \textit{stoneskin} (self only)

3/day—\textit{sending}

\end{DndReadAloud}

\item The \textit{Throne of Drakkenheim} is irreversibly contaminated whilst the Haze remains. A humanoid creature seated on the throne can't regain hit points, and each time a creature activates its powers, they must succeed on a DC 20 Constitution saving throw or take 10 (3d6) necrotic damage and gain one level of Contamination.
\end{itemize}

\includegraphics[width=0.4\textwidth]{images/../5e.tools/img/adventure/DoDk/137-map-8.03-castle-drakken-main.png}
\includegraphics[width=0.4\textwidth]{images/../5e.tools/img/adventure/DoDk/138-map-8.03-castle-drakken-main-player.png}
\subsubsection{The Amalgamation}
\begin{DndReadAloud}{}
Before the throne is a gibbering amalgamation of fleshy tumors, rotting bones, and rancid layers of bloated fat suspended between the ceiling and floor by stalky strands of gnarled tissue. The mass is covered in bursting sores, shuddering mucus nodules, and weeps bile and delerium sludge, The creature mutters in alien tongues from six yawning mouths, piping and sputtering in discordant musical tones, and lazily shifting eyes of massive size open and close throughout its bulk. Thick trunk-like tentacles reach out from its hideous stinking form to squirm and writhe around the chamber. Several appendages are wrapped around pulsing human organs: lungs, liver, kidney, a heart, a brain, and dangling amongst them is the jewelled Crown of Westemär. A palpable aura of contamination emanates from the beast.

\end{DndReadAloud}


\begin{itemize}
\item The \textbf{Amalgamation} and 2d6 gibbering mouthers are here, and slavishly attack any who enter the throne room.
\item While the physical remains of Ulrich IV have become part of the Amalgamation, his mind and soul are forever annihilated.
\item If the amalgamation is destroyed, it returns with all its hit points after 1d10 days unless the dimensional rift is closed.
\end{itemize}

\DndQuote%
{}
{''Even my fondest childhood memories of Castle Drakken can't shake the feeling of dread when the shadow looms heavy over the haze.''}
{Veo Sjena}
\section{Dungeons of Drakkenheim}
The dungeon vaults are accessed through the Stair Tower. Instead of the gilded finery and impressive paintings of the upper chambers, the dungeons are decorated with solemn statuary, mosaic brickwork, detailed bas-relief carvings, and ribbed and vaulted ceilings twenty feet overhead.

\subparagraph{Pool of Tears}
The base of the Stair Tower houses a glimmering marble fountain. Three mischievous gargoyles are bound to the fountain, nicknamed Cheeky, Rump, and Bazoo. They are friendly to passersby and speak common if addressed. They don't know much about what's happened in the castle, but many years ago they saw the steward run down into the dungeons. He never returned.

\subparagraph{Escape Chute Base}
This dusty chimney is the base for the escape chimneys in the Royal Apartments. The sliding pole ends about four feet from the ground.

\subparagraph{Heroes' Hall}
The statues lining this hallway honor adventurers and heroes who performed great deeds for the royal family, which should be subtly described as having the likenesses of your players' past characters (if appropriate). Two can be pushed aside easily for access to the secret passages.

\subparagraph{Vault Watchers}
Guarding the way to the royal vaults and crypts are two stone golems. The steward once attempted to take control of them, but ended up scrambling the arcane commands that bound them to the castle. They haltingly stumble around their posts and attack anyone who enters.

\includegraphics[width=0.4\textwidth]{images/../5e.tools/img/adventure/DoDk/139-map-8.04-castle-drakken-dungeons.png}
\includegraphics[width=0.4\textwidth]{images/../5e.tools/img/adventure/DoDk/140-map-8.04-castle-drakken-dungeons-player.png}
\subsection{Royal Vaults}
The wealth of Westemär rests behind these large, tightly sealed iron doors: there are no handles, locks, hinges, or seams. Each chamber is warded with \textit{Mordenkainen's Private Sanctum}; the doors are \textit{Arcane Lock} and immune to all damage. The vault doors slide open when touched by one who wears the \textit{Crown of Westemär}, or three characters attuned to the \textit{Seals of Drakkenheim}.


\begin{itemize}
\item The royal vaults contain an incredible treasure of nearly 35,000 gold pieces in heavy gold bullion and silver talents, neatly stacked.
\item A gemstone collection of 6 emeralds, 4 diamonds, 7 rubies, and 8 sapphires, worth 1,000 gp each. Resting as the crowning jewel amongst these is the \textit{Diamond of Mount Kadath}.
\item A collection of magically notarized, preserved, and sealed bonds, credit notes, financial agreements, and property deeds involving staggering sums and distant foreign powers. These pieces of paper are the \textit{true} wealth of House von Kessel.
\item Jewelry boxes containing a sapphire amulet, diamond necklace, sapphire tiara, and a ruby-encrusted circlet each worth 2,500 gp.
\item The final chamber houses one very rare magic item and one rare magic item of your choice, and a Spell Scroll of \textit{resurrection}.
\end{itemize}

\subsection{Monarchal Crypt}
\begin{DndReadAloud}{}
The royal mausoleum is a fifty-foot-high chamber supported by two dozen rectangular columns. Gargantuan delerium crystals have breached through the walls and floors, and thick delerium sludge fills the hall to the knee's height, bubbling and hissing like a foul stew. The air is cold and stale.

Upon the sides of each pillar are five decorated niches roughly eight feet tall each. They are placed one atop another and separated by ornate borders. Some are empty, but many others hold a life-sized sculpture of a former monarch of Drakkenheim. Each is cast with a regal countenance and idealized features; each holds an object symbolizing their legacy and is flanked by candelabras. Horrifyingly, the outer shells of the statues are crumbling away, revealing exposed muscle and weeping tissue beneath. Some murmur from mouths half-stone and half-flesh. Rays of prismatic light illuminate the chamber from a bulbous form that floats and bobs between the pillars.

\end{DndReadAloud}


\begin{itemize}
\item The \textbf{Gravekeeper} attacks if it detects any trespassers in this chamber.
\item Delerium sludge floods the floor - characters within must succeed on a DC 18 Constitution saving throw at the start of their turn. They take 42 (12d6) necrotic damage on a failed saving throw, and half as much on a successful one. In addition, characters who fail the saving throw gain one level of contamination
\item A plaque below each statue names each monarch with dates of their life and rule, beginning with King Vladimir von Drakken and ending with Queen Helena von Kessel II, Ulrich IV's mother. There are roughly sixty monarchs interred here across five distinct family dynasties. with reigns ranging from decades to days. Behind each statue is a small sealed sepulcher that holds their skulls and ashes. The outer walls are lined with covered sepulchral drawers holding the remains of other royal family members.
\end{itemize}

\includegraphics[width=0.4\textwidth]{images/../5e.tools/img/adventure/DoDk/141-08-005.doundation-pillar.png}
\subsubsection{Foundation Pillar}
\begin{DndReadAloud}{}
A large alcove on the left-hand side of the chamber features a heavily damaged granite pillar. Unlike the others, it is cylindrical, etched with arcane runes, and appears on the brink of collapse. However, the deep cracks are held together by pulsing tendrils of living stone emerging from a peculiar lump of melted rock latched to the pillar like a leech.

\end{DndReadAloud}


\begin{itemize}
\item The Foundation Pillar is the power source for the castle wards. It is magically connected to the \textit{Seals of Drakkenheim}, the \textit{Throne of Drakkenheim}, and the \textit{Crown of Westemär}. These magic items lose their magical properties if the Foundation Pillar is destroyed.
\item The dimensional rift is drawing out its immense magic to fuel its growth. If the steward wasn't holding the pillar together to slow this process, the rift would destructively spill open and destroy the castle.
\item Johann Eisner fled here after failing to save Ulrich IV when the Haze overtook Castle Drakken fifteen years ago, He was about to undergo a monstrous transformation when he laid his hands upon the foundation pillar. He has become fused to it as a result, and this accident granted him limited control over its magic. He can't be freed without slaying him, except by a \textit{wish} spell. A creature within five feet of the foundation pillar can see his face merged within the melted rock. The steward can speak in a halting whisper to characters here.
\item The foundation pillar has AC 20 and 250 hit points. It is immune to fire, poison, and psychic damage, and bludgeoning, piercing, and slashing damage from non-magical weapon attacks.
\item If Eisner is freed from the pillar while the dimension rift remains open, the foundation pillar crumbles in twenty four hours, and Castle Drakken is destroyed by the rift.
\end{itemize}

\DndQuote%
{}
{''The power of a single seal is multiplied when the noble few take up the cause to bring unity and hope to this crumbling city. If only that was all that's needed to restore what we've lost here...''}
{Veo Sjena}
\begin{DndSidebar}[float=!b]{Developments}
Possible outcomes for events directly involving Castle Drakken are summarized below, while the next chapter, the "Fate of Drakkenheim" details possible conclusions for the campaign as a whole.



\end{DndSidebar}

\section{Closing the Dimensional Rift}
The rift can be sealed in three ways:


\begin{itemize}
\item A \textit{wish} spell (such as one used from the \textit{Crown of Westemär}).
\item Investing the power of the \textit{Seals of Drakkenheim} into the steward (see Eternal Custodian).
\item Physically travelling into the rift and destroying the Amalgamation again on the other side.
\end{itemize}

The Haze will eventually tear open more rifts as it expands beyond the city, but this is the only one in Drakkenheim... for now.

\section{Beyond the Dimensional Rift}
The Space Between Worlds is a roiling non-place. The rift manifests on the other side as an abrupt tear in the peripheral perception. A pernicious beam of multihued light stretches across the effulgent sky, marking a distant place beyond an impossible horizon where the Amalgamation's multitudinous manifestations converge. There is no meaningful concept of time in the Space Between Worlds - it is beyond the perception of mortal creatures - so reaching that place is a journey of eons that unfolds in only 1d10 minutes of time in the mortal world. Nevertheless, should an intrepid band make the journey and destroy the Amalgamation here as well, the dimensional rift closes one round of apparent time later. Any characters left behind in the Space Between Worlds must find another means back, but the scope of such adventures is beyond this book.

\section{Eternal Custodian}
Johann Eisner believes Drakkenheim is forever cursed, and has resigned himself to be its jailer. When he meets the characters, Eisner explains how it might be possible for him to close the dimensional rift if he is invested with the power of the \textit{Seals of Drakkenheim}:


\begin{itemize}
\item Accomplishing this task requires each of the six \textit{Seals of Drakkenheim} and the \textit{Crown of Westemär} to be pressed against the foundation pillar.
\item This will attune Eisner to the Seals, but won't grant him the full powers of the throne or the crown since a coronation ritual isn't performed. He is only able to pacify the wall gargoyles and tower dragons so they don't attack humanoid creatures anymore, but instead attack monsters who come near. However, he is able to use this power to close the dimension rift.
\item If Eisner dies or another creature attunes to the crown, the rift reopens.
\end{itemize}

\section{Coronation Ceremony}
This special ceremony takes two hours to perform and requires seven participants. During this time, one participant (the "heir") must sit upon the \textit{Throne of Drakkenheim} wearing the \textit{Crown of Westemär}. The other six participants must each be attuned to a different \textit{Seal of Drakkenheim}. The ritual fails immediately if any participant becomes incapacitated, moves more than thirty feet away from the throne, or if the heir does not remain seated on the throne for the duration.

One after the other, each participant attuned to a \textit{Seal} must approach the throne, pledge their allegiance, and offer their lifelong service to the new monarch. Each then makes a DC 20 ability check using an ability score that corresponds to their \textit{Seal of Drakkenheim}:


\begin{itemize}
\item \textit{Lord Commander's Badge} (Strength)
\item \textit{Steward's Seal} (Constitution)
\item \textit{Spymaster's Signet} (Dexterity)
\item Inscrutable Staff (Intelligence)
\item Saint Vitruvio's Phylactery (Wisdom)
\item \textit{Chancellor's Crest} (Charisma)
\end{itemize}

The designated heir responds to each ability check by making an appropriate check of their own using the same ability score. While all 12 checks must be attempted to finish the ceremony, only 8 successful check results are needed to successfully complete it and attune the heir to the crown. Encourage each character participating in the ritual to roleplay out their vow, and the response of the monarch to each.

Since the throne is contaminated, each failed ability check causes the designated heir to immediately take 14 (4d6) psychic damage. In addition, they must succeed on a DC 20 Constitution saving throw or take an additional 10 (3d6) necrotic damage and gain one level of contamination. Failure by 5 or more causes the heir to take 35 (10d6) damage and gain 1d4 levels of contamination instead. Remember, a creature can't regain hit points while seated on the throne! If the heir is transformed into a monster or reduced to 0 hit points, they immediately transform into another \textbf{Amalgamation}.

\subparagraph{NOTE}
Not every Seal needs to be worn by a player character: the characters might choose one or more faction leaders or other key non-player characters to wear the Seals. In such cases, these NPCs should offer their hopes for the future, express their gratitude to the player characters, and celebrate their accomplishments. You might even allow the players to make dice rolls for them.

\section{Be careful what you wish for..}
No matter how carefully worded, it's well beyond the power of \textit{wish} to instantaneously or fully cleanse Drakkenheim of all contamination, delerium, and the Haze. A wish can restore the original form of a single creature who suffered a monstrous transformation, however.

\section{Lost Seals}
Characters foolish enough to offer a \textit{Seal of Drakkenheim} to the devils have effectively doomed any hope for restoring the throne properly. Perhaps such characters could mount an expedition to Hell to retrieve the \textit{Seal} on an adventure of your own creation!

\includegraphics[width=0.4\textwidth]{images/../5e.tools/img/adventure/DoDk/142-08-006.delirium-crystal.png}
\chapter{The Fate of Drakkenheim}
\DndDropCapLine{T}{}he player characters define their adventures in Drakkenheim, so there's no telling exactly how their story might end. Each character has a personal quest to complete, but there are two major questions that should be resolved at the end of the campaign: \textit{What should be done about delerium? Who will rule Drakkenheim}?


\includegraphics[width=0.4\textwidth]{images/../5e.tools/img/adventure/DoDk/143-09-001.journey-starts.png}
\section{Resolving Personal Quests}
When your characters finish their personal quests, they should ponder on the following:

\subparagraph{Find a Cure}
What was the cost of the cure? Did the character discover a new spell like \textit{siphon contamination}, or did they use a \textit{wish} from the \textit{Crown of Westemär}?

\subparagraph{Claim the Throne}
If successfully accomplished, the new monarch can make three \textit{Wish} using the crown. What will they do with this power?

\subparagraph{Apocalyptic Vision}
After seeing the vision at the Delerium Heart, what did the character decide to do? Why?

\subparagraph{Faction Aspirant}
The character emerged as a mighty champion of their faction. Perhaps they should take up the mantle as the new faction leader?

\subparagraph{Score to Settle}
What was the price of revenge?

\subparagraph{Arcane Secrets}
This character has unlocked unimaginable arcane power. What will they do with this knowledge? What dark new spells will they discover next?

\subparagraph{Blighted Landscape}
What discoveries did this character make? How did their knowledge help them resolve the conflicts in Drakkenheim?

\subparagraph{Family Matter}
Was the character reunited with their loved one, or did they pay a personal price in this conflict?

\subparagraph{Overwhelming Debt}
A debt once paid need not be entirely forgiven. What other favours and obligations might the character's debtors call in later?

\subparagraph{Lost Heirloom}
What mysterious powers and strange properties might the heirloom have?

\subparagraph{Personal Pilgrimage}
What is the character's place in the divine order? Did they come to save these places, or see them for the last time?

\subparagraph{Monster Slayer}
What songs will the bards sing of this character's exploits?

\section{Defeated Factions}
Factions can be directly defeated in myriad ways. At the very least, this requires removing its leader, routing the rank-and-file members, and driving them out of their stronghold. However, other inventive solutions are possible. Perhaps characters help the Followers of the Falling Fire convert the Knights of the Silver Order to believe the words of Lucretia Mathias by revealing the meteor's origin. Maybe they present irrefutable evidence to the Amethyst Academy that convinces the mages delerium is too dangerous to use responsibly. Working with the Queen of Thieves, characters might weave a tangled conspiracy and control the Hooded Lanterns from the shadows. How your players might accomplish such dramatic reversals is up to them to attempt and you to adjudicate!

Even in defeat, the factions still possess distant strongholds, fresh recruits, vast resources, and additional allies outside Drakkenheim. Nevertheless, retreating from Drakkenheim sends the faction reeling. Before they can return and exact vengeance, they will need at least a year to regroup, appoint new leaders, and determine a new strategy. You might decide to end the campaign in this moment of victory with the spark of future conflict on the horizon.

Alternatively, you might continue the campaign into higher levels. If you do, give characters time to make plans of their own, pursue other adventures, and simply take some well-deserved downtime first. While such adventures are beyond the scope of this book, consider the following scenarios:


\begin{itemize}
\item The High Paladin of the Silver Order dispatches a new crusading force to lay siege to Drakkenheim. They leverage the full military might of Elyria and seek support from any religious nobles in Westemär to form a powerful new coalition.
\item The Directorate of the Amethyst Academy sets grand designs in motion to reassert their dominance over delerium. They devise potent spells, artefacts of terrible power, and cunning manipulations to exact revenge against the player characters.
\item A mysterious figure takes up the mantle of the Queen of Thieves, weaving a larger web of intrigue. Perhaps the Queen of Thieves herself was merely an elaborate ruse for another shadowy mastermind?
\item Inspired by the martyrdom of Lucretia Mathias, a new wave of disciples and acolytes rouse the faithful into a militant pilgrimage. A new wave of religious conflict threatens to set the continent ablaze.
\item Outraged that an incomptetent nobility stood idle while the Hooded Lanterns were defeated, the people of Westemär unite under a charismatic rebel who promises a democratic future.
\end{itemize}

\section{Epilogues}
Below are a few possible outcomes and suggestions about how a Drakkenheim campaign may end, and the aftermath that follows. Each offers a dramatic conclusion to the campaign, and an epilogue that presents an uncertain future. Not all of these are mutually exclusive, either! Your players could devise an inventive solution that combines several different possibilities.

\subsection{Restoring the Throne}
Though the world is indeed threatened by cosmic forces, these issues may seem distant and aloof compared to the pressing political turmoil that presently engulfs the continent. Whoever manages to wrest control of Drakkenheim will assume vast influence, enough to dictate the balance of power. While the Queen of Thieves plots the means to take advantage of this situation for her own gain, the Hooded Lanterns desperately hope they might locate a potential heir to the throne and secure the means to crown them. Who this heir might be is up to you to decide: consult the "Long Lost Heir" section in chapter 2 for details.

\subsubsection{Proving the Claim}
Finding the heir is one thing; but \textit{proving} that the claimant has royal lineage is another. An heir might be identified by a broken locket, birthmark, signet ring, or some other document, but this must be substantiated with something from the ruins. There are three places such evidence might be found:


\begin{itemize}
\item The vaults of Saint Vitruvio's Cathedral.
\item The steward's archives in Castle Drakken.
\item The private office of the Queen of Thieves.
\end{itemize}

You decide the exact nature of this evidence. It could include a detailed family bible outlining the order of succession, a royal will specificing the next heir, vials of preserved royal blood that can be used in a magic ritual to confirm relation, an unsent letter that grants trueborn status to an illegitimate child, a scandalous document that exposes infidelity, or the other half of a broken locket. Regardless, it should be something tangible and irrefutably authentic. If characters haven't yet found an heir, these documents could also reveal their identity and location.

\subparagraph{Faction Support}
A claimant to the throne will almost certainly need the support of two or more of the factions to secure the political and military force needed to win over the nobility. This doesn't necessarily have to be the Hooded Lanterns, but they are a safe bet.

\subsubsection{Coronation}
The \textit{Crown of Westemär} and the \textit{Throne of Drakkenheim} are more than just political concepts: they're also powerful magic items. However, unlocking the full abilities of these items requires the six \textit{Seals of Drakkenheim} be used to perform a special coronation ritual in Castle Drakken (see chapter 9). Unfortunately, the corruption of the \textit{Throne of Drakkenheim} by the Haze means that there's a good chance a potential claimant to the throne may die in the process if they aren't strong enough to survive - unless your players make cunning use of protective spells!

Technically speaking, these artefacts aren't actually \textit{required} for a new monarch to reign over Westemär, but holding those items conclusively places a potential claimant in the seat of political power. A rightful ruler seated upon the \textit{Throne of Drakkenheim} will be able to unite the nobility of Westemär once more, but one without them will have a much more difficult time commanding the respect and cooperation of the wider nobility.

\subsubsection{New Dynasty}
Elias Drexel and the Hooded Lanterns believe that restoring both Drakkenheim and the Realm of Westemär \textit{requires} a scion of House von Kessel to sit on the \textit{Throne of Drakkenheim}. However, observant characters may notice that at least four other families have ruled Drakkenheim prior to House von Kessel. As such, another noble house could make the claim provided they possess the wealth, political influence, and military might needed to secure power. This family then establishes a new royal dynasty to rule Westemär.

Such a task would require demonstrating \textit{some} connection to the last dynasty (as per above), the political support of at least three factions, as well as holding the power of the royal treasures. Especially intrepid heroes might not like the idea of a monarchy at all. If your characters have ambitions to set up a new form of government, embrace it, but don't make it easy for them!

\subsubsection{Queen's Castle}
If the Queen of Thieves successfully obtains the \textit{Crown of Westemär}, she'll use its power to cement her rule over the ruins. If you opted to portray the Queen of Thieves as Katarina von Kessel, she uses her first wish to close the dimensional rift, the second to restore the life of her lost sister, and the final to erase all knowledge of her magical abilities. She then reveals her true identity so she may legitimately rule Drakkenheim. Whether or not Katarina von Kessel is a good ruler is uncertain, but she proves a shrewd politician. Undoubtedly, her political ambitions will not stop with Westemär alone.

\subsection{Reborn Flames}
The Knights of the Silver Order insist that destroying all the delerium in Drakkenheim must be done to avert an apocalypse, but whether accomplishing such a task is even possible is uncertain. Conversely, the Followers of the Falling Fire vehemently believe delerium offers the key to salvation, but this divine upheaval will require thousands make pilgrimage to Drakkenheim. Both factions think they can achieve their goals by resurrecting a legendary harbinger of the Sacred Flame.

Beneath the Cathedral of Saint Vitruvio lie the bones of Argonath, an ancient gold dragon. Although the wyrm has been dead for over five hundred years, Saint Vitruvio proclaimed that in the city's darkest hour, a faithful soul could cast \textit{resurrection} on the dragon's remains to bring it back to life. However, a much more powerful gemstone is required for this purpose: either the \textit{Diamond of Mount Kadath}, a priceless gem located in the royal vault of Castle Drakken, or the \textit{pristine delerium geode} taken from the Delerium Heart in the crater.

Once resurrected, Argonath only heeds the words of one who bears all the \textit{Relics of Saint Vitruvio}. Otherwise, he slumbers.


\begin{itemize}
\item Saint Vitruvio's Phylactery, found in Saint Vitruvio's Cathedral
\item \textit{Ignacious, the Sword of Burning Truth} sealed beneath the cathedral catacombs.
\item The Shield of Sacred Flame found in the royal apartments of Castle Drakken
\item \textit{Helm of Patron Saints} located in the crypts of Saint Selina's Monastery
\item \textit{Sceptre of Saint Vitruvio}, located in the Chapel of Saint Brenna
\end{itemize}

\subsubsection{Purged with Dragonfire}
If resurrected by the Silver Order, Argonath helps the knights burn Drakkenheim to the ground. Destroying all the delerium within the ruins is a monumental undertaking, but just might be possible with the power of an ancient gold dragon. Argonath's flames deal radiant damage, not fire damage.

A suitable climax could see the player characters riding Argonath into the crater alongside the Knights of the Silver Order, where they battle to destroy the Delerium Heart. If they are successful, the Haze will gradually disappear over the next two decades. The change becomes evident when many monsters grow timid and the mists slowly recede. Smaller pockets of delerium continue to emerge, however, necessitating constant cleanup efforts. The Silver Order staunchly maintains control over the city until the work is done, though other factions such as the Amethyst Academy and the Followers of the Falling Fire will almost certainly continue opposing their efforts. Once every trace of contamination is eradicated, people may return to Drakkenheim to rebuild. Unfortunately, unless a new monarch is firmly established in another city (see "Restoring the Throne" above), the political outcome for Westemär is grim. The nation disintegrates into a dozen warring provinces long before Drakkenheim itself is restored.

\subsubsection{Darkness and Salvation}
Should Argonath be resurrected by the Followers of the Falling Fire, the worshipers use the dragon to drive out the Knights of the Silver Order and the Amethyst Academy. Thereafter, Argonath keeps vigil over Drakkenheim, protecting future pilgrims who come to join the Falling Fire. Lucretia Mathias will die a few months before her 100th birthday, but before her death, she writes one final testament. It affirms the deepest shadows are yet to come, but the Sacred Flame will spark again from a moment of utter darkness.

Provided charismatic disciples such as the player characters rise to take her place and lead the Followers of the Falling Fire, the nascent religion spreads. New converts flow into Drakkenheim to take the sacrament - first hundreds, then thousands, then more. The Falling Fire continues to face intense persecution from the mainstream clergy, but the resurrection of Argonath by a group of supposed heretics shakes the foundations of the faith in the eyes of the laity. In time, millions come seeking a place in the new divine order. When these individuals die, their golden delerium soulstones are gathered in the crater. However, Drakkenheim is never truly rebuilt; the ruins come to be seen as a holy place.

The Followers of the Falling Fire do nothing to prevent new delerium formations, and so the Haze spreads unabated. The faithful stay the course as the world slowly descends into chaos over several decades. Sanctified devotees may survive within the Haze, but monsters grow great and terrible as the whole earth is swallowed by the mists. Even mighty Argonath eventually falls. When the world begins shaking apart, the final Follower of the Falling Fire completes the sacrament before a vast assemblage of golden delerium crystals containing millions of righteous souls. In that moment, Lucretia Mathias' prophecies might come to pass, or they might not. Either way, the world undergoes a new apotheosis in a brilliant corona of multicoloured light.

\subsection{Sequestered City}
The Amethyst Academy wants to harness the arcane power of delerium, and believes they can contain the dangers involved. Eldrick Runeweaver has designed a monstrously powerful version of the \textit{prismatic wall} spell, one large enough to form a containment dome enclosing the entire city. This is a closely guarded secret, and he will not reveal this plan until every piece is in place.

Such a dramatic display of magical power is a major political risk for the Amethyst Academy. They won't carry out this plan without support from a new monarch of Drakkenheim, or at least with sufficient backing from whatever authority remains in Westemär. The archmages claim the barrier is the best hope for eventually restoring Drakkenheim: it will allow the Amethyst Academy to carefully extract the delerium so they can put it to proper use constructing magic items, which they will happily sell to anyone who can afford them. Once all the delerium is quarried, the Amethyst Academy claims they will lower the barrier and help rebuild the city.

The spell requires tremendous arcane power to set in motion and maintain:


\begin{itemize}
\item The Inscrutable Tower Nexus acts as the fulcrum of the spell. The \textit{Inscrutable Staff} is needed to reactivate and control the Nexus for this purpose.
\item Characters must install a suitable local power source in the Nexus: a \textit{pristine delerium geode} taken from the Delerium Heart in the crater. Thereafter, a new delerium crystal must be supplied daily to maintain the spell.
\item A special arcane siphon must be installed at the Delerium Heart to channel its power to bolster the spell.
\item 8 arcane spellcasters who can cast spells of 5th level or higher must participate in the ritual. At least one must be capable of casting 9th level spells.
\end{itemize}

Naturally of course, the player characters are called upon to enact these plans and participate in casting the spell, which might prompt a last-ditch sabotage attempt from an enemy faction. Once cast, the spell is indeed powerful enough to halt expansion of the Haze, effectively turning Drakkenheim into a greenhouse where the Amethyst Academy may harvest and study delerium at their sole discretion. Thereafter, only Academy mages and their designated allies are permitted through the barrier. A few scavengers and looters find secret tunnels and hidden means of entry, and occasionally, monsters escape via the river or underground passages. The city remains a dangerous place, but one largely quarantined within the colossal prismatic barrier. However, rather than strip-mine delerium, the Academy carefully cultivates it. The barrier remains indefinitely, because the Academy intentionally allows more delerium to grow within to maintain their supply. The Amethyst Academy uses its newfound wealth and influence to clandestinely reshape the political destiny of the continent, and possibly even unravel the Edicts of Lumen.

Nevertheless, should the prismatic barrier fall through sabotage or negligence, the results would be appropriately catastrophic. The feedback force shatters the Inscrutable Tower and causes the Haze to violently explode outward for hundreds of miles around. The volatile magic greatly accelerates the expansion of the Haze, and the world is destroyed only a few years later.

\subsection{Anarchy in Drakkenheim}
\includegraphics[width=0.4\textwidth]{images/../5e.tools/img/adventure/DoDk/144-09-002.wizard-portrait.png}
In the midst of ongoing conflict, it's possible none rule the ruins save monsters and mayhem. Perhaps several of the plans above fail, no credible heir survives, and the player characters and factions decide that Drakkenheim is best abandoned. A new ruler might decide to establish another city as the capital of Westemär. Unless they possess both political savvy and military strength, Westemär will likely fragment into several smaller realms as a result. Indeed, this is the most likely course unless the player characters intervene, though one that unfolds over decades of intrigue, betrayal, and war.

By that time, the situation in Drakkenheim may become even more dire as the Haze grows larger and larger. In about a decade, containment will become almost impossible when more dimensional rifts and greater horrors begin to awaken within the crater. A few may callously conclude that the destruction is still centuries away, and so the Haze is a problem for future generations to inherit. Other powerful and wealthy recluses might selfishly plot to save themselves by fleeing to other worlds. Perhaps, a new generation of hopeful heroes may come together and try to save the world before it is too late.

\DndQuote%
{}
{''The city is not a place people come without purpose. But some people's purpose can only be found here.''}
{Sebastian Crowe}
\appendix
\chapter{Monsters}
Herein are described several new monsters that may be found in Drakkenheim. This appendix also includes many creatures based on monsters found in the \textit{Core Rules} with additional traits and new descriptions to better suit this dark fantasy setting. However, any 5e-compatible monster could conceivably stalk the cursed streets of Drakkenheim as a twisted product of eldritch contamination, or magically conjured from distant netherworlds by the Haze. Use these modifications as inspiration to adapt your favourite monsters from other sources into twisted abominations in Drakkenheim.

\includegraphics[width=0.4\textwidth]{images/../5e.tools/img/adventure/DoDk/145-10-001.alleyway.png}
\section{Mechanical Features}
All non-humanoid creatures encountered in Drakkenheim possess the following traits. These do not change the monsters' challenge ratings.


\begin{description}
\item[Fully Contaminated.]
This creature is immune to contamination, and has advantage on saving throws against contaminated spells.
\item[Haze Adaptation.]
This creature may benefit from a long rest finished within the Haze. It can see twice as far through the Haze mists and Deep Haze fog.
\item[Haze Dependence.]
This creature gains one level of exhaustion for every twenty four hours it spends outside the Haze. These exhaustion levels cannot be removed while the creature is outside the Haze.
\item[Delerium Healing.]
As an action, a monster can touch a delerium shard to regain 10 (3d6) hit points. Once a monster has used a delerium shard in this way, the shard can't be used in this manner again for twenty four hours.
\end{description}

These benefits do not apply to creatures originating outside the Haze or who are recruited or conjured by the characters, such as familiars and summoned monsters.

\section{Roleplaying Traits}
All monsters inhabiting Drakkenheim are irrevocably gripped by otherworldly madness. Many become predatory and violent as a result, but a few are lucid. However, even these creatures possess an unstable consciousness that drifts between psychosis and psychopathy. They do not understand the origin of their malevolent predilections. All monsters have the following roleplaying traits:


\begin{description}
\item[Personality Trait.]
I'm irritable or outright hostile around humanoid creatures. Logic and reason confuse and frustrate me.
\item[Ideals.]
I feel a strange compulsion to contaminate or devour humanoid creatures.
\item[Bonds.]
I'm obsessed with delerium crystals and wish to bask in their warm glow.
\item[Flaws.]
I have one or more forms of indefinite madness, and often have trouble remembering what I was doing even a few moments ago.
\end{description}

A creature's reaction to other monsters in the Haze depends on its own nature. It may be ambivalent, avoidant, hostile, or even servile.

\section{Monster Modifications}

\begin{itemize}
\begin{multicols}{3}
\item \textbf{Arcane Wraith}
\item \textbf{Crater Wurm}
\item \textbf{Garmyr}
\item \textbf{Grotesque Gargant}
\item \textbf{Haze Husk}
\item \textbf{Haze Wight}
\item \textbf{Tower Dragon}
\item \textbf{Wall Gargoyle}
\item \textbf{Warp Witch}
\end{multicols}

\end{itemize}

\section{New Monsters}
\DndQuote%
{}
{''Who knows how many of these twisted abominations were once our steadfast comrades, friends, and loved ones. When you see these lost souls, you must not hesitate.''}
{Elias Drexel}

\begin{itemize}
\begin{multicols}{3}
\item \textbf{Delerium Dreg}
\item \textbf{Haze Hulk}
\item \textbf{Protean Abomination}
\item \textbf{Ratling}
\item \textbf{Warlock of the Rat God}
\item \textbf{Animated Delerium Sludge}
\item \textbf{Living Deep Haze}
\item \textbf{Entropic Flame}
\item \textbf{Walking Delerium Geode}
\item \textbf{Hypnotic Eldritch Blossom}
\end{multicols}

\end{itemize}

\section{Legendary Creatures}

\begin{itemize}
\begin{multicols}{3}
\item \textbf{Amalgamation}
\item \textbf{Crimson Countess}
\item \textbf{Executioner}
\item \textbf{Gravekeeper}
\item \textbf{Lord of the Feast}
\item \textbf{Pale Man}
\item \textbf{Rat Prince}
\end{multicols}

\end{itemize}

\chapter{Nonplayer Characters}

\begin{itemize}
\begin{multicols}{3}
\item \textbf{Lord Commander Elias Drexel}
\item \textbf{Eldrick Runeweaver}
\item \textbf{Knight-Captain Theodore Marshal}
\item \textbf{Lucretia Mathias}
\item \textbf{Queen of Thieves}
\item \textbf{Cavalier}
\item \textbf{Chaplain}
\item \textbf{Hedge Mage}
\item \textbf{Scoundrel}
\item \textbf{Urban Ranger}
\end{multicols}

\end{itemize}

\chapter{Contamination}
Characters exploring Drakkenheim may be exposed to eldritch pollutants within the Haze or arcane radiation emitted by delerium. These hazards cause a new condition called contamination, which inflicts both debilitating symptoms and otherworldly mutations.

Contamination is measured in six levels. An effect can give a creature one or more levels of contamination, as specified in the effect's description. If an already contaminated creature suffers another effect that causes further contamination, its current level of contamination increases by the amount specified in the effect's description.

In addition, each time a character finishes a long rest while they have one or more contamination levels, roll 1d20. On a 1, they gain a contamination level.

Upon reaching contamination level 6 or higher, the creature transforms as described in the "Monstrous Transformation" section.

\section{Contamination Symptoms}
A contaminated creature suffers the symptoms from its current level as well as all lower levels.


\begin{DndTable}{cX}

Level & Symptoms\\
1 & None.\\
2 & Hit points regained by expending hit dice halved.\\
3 & No hit points regained at the end of a long rest.\\
4 & Damage dealt by weapon attacks and spells halved.\\
5 & Incapacitated.\\
6 & Monstrous Transformation!\\
\end{DndTable}

\section{Mutations}
In addition to suffering symptoms, each time a character gains a contamination level, it rolls 1d6. If the result is equal to or less than the character's current contamination level the creature gains a mutation. The Game Master may choose the mutation, or determine it randomly by rolling 1d20 on the Mutations Table. If the result is the same as a mutation affecting the creature currently, the GM chooses a different one or rolls again.

\begin{DndTable}[header=Mutations]{cX}

1d20 & Mutations\\
1 & Rampant Mutation! Roll twice, ignoring this result on subsequent rolls.\\
2 & Rasping. Your vocal cords warp, and you may only speak in a halting gurgle. If you have 4 or more contamination levels, your tongue rots and falls out, and you can no longer speak.\\
3 & Wasting. Your fingernails, teeth, and toenails start falling out. 2d6 fall out for each contamination level you have gained.\\
4 & Rotting. Your lips, nose, and ears blacken and wither. If you reach 4 or more contamination levels, they rot and fall off. You can still speak and hear, however.\\
5 & Molting. Painful blisters, welts, and multicoloured lesions appear all over your skin, which burst and peel off painfully, exposing the raw sinew underneath. Once you reach 4 contamination levels, your skin entirely sloughs off.\\
6 & Shedding. Each time you gain a contamination level, some of your hair falls out in patches. Once you reach 4 or more contamination levels, all hair on your body completely falls out.\\
7 & Lambent Glow. You emit a dim octarine glow to a range of 10 feet. If you have 4 or more contamination levels, you instead emit bright light to a range of 30 feet.\\
8 & Ocular Tumors. An eyeball opens somewhere on your body for each contamination level you have gained. If you have 4 or more contamination levels, you can see in all directions.\\
9 & Spiked Growths. At the start of each of your turns, you deal 5 (1d10) piercing damage to any creature you are grapple.\\
10 & Aquatic Adaptation. You sprout fish-like fins and gills. You gain a swimming speed equal to your land speed and can breathe underwater. If you have 4 or more contamination levels, you can only breathe underwater; but can hold your breath outside water for up to 1 hour.\\
11 & Amorphous Form. Your bones and organs become gelatinous. You can move through a space as narrow as 6 inches wide without squeezing.\\
12 & Chitinous Skin. Shell-like growths appear all over your body, giving you a +1 bonus to AC. If you have four or more contamination levels, this bonus increases to +2.\\
13 & Cyclopean Vision. Your eyes merge into a single central eye which can emit an energy beam as a ranged spell attack using your Intelligence modifier for the attack roll. If it hits, it deals 2d6 radiant damage.\\
14 & Spatial Displacement. You can cast \textit{misty step} once for each contamination level you have gained. You regain these uses when you finish a long rest.\\
15 & Tentacled Limb. One of your arms becomes a fleshy tentacle. When you make a melee attack on your turn, increase your reach by 5 feet.\\
16 & Spider Climb. You gain a climb speed equal to your walking speed. You can climb difficult surfaces, including upside down on ceilings, without needing to make an ability check.\\
17 & Whispering Voices. You gain telepathy to a range of 10 feet, but other people hear it as their own voice. If you have 4 or more contamination levels, the range extends to 60 feet.\\
18 & Belly Maw. A toothy mouth appears on your stomach, which you can use to make unarmed strikes. If you hit with it, you deal piercing damage equal to 1d6 + your Strength modifier.\\
19 & Eyeless Sight. Your eyes become milky orbs, and you gain blindsight to a range of 10 feet. If you have 4 or more contamination levels, your eyes rot out, and dim octarine light burns in the sockets. Your blindsight increases to 30 feet, but you are blind beyond this radius.\\
20 & Arcane Blood. Gain an additional spell slot of the highest level you can cast (to a maximum of 5th level). If you don't have spell slots, your hit point maximum increases by an amount equal to your level. Double with 4 or more contamination levels.\\
\end{DndTable}

\section{Removing Contamination}
Only specific magic can remove contamination levels. Creatures do not recover naturally. The \textit{purge contamination} spell can remove contamination, but it leaves affected characters exhausted. A \textit{heal} spell removes all contamination levels and mutations from a contaminated character.

An effect which removes a contamination level also removes one randomly determined mutation. All contamination symptoms end and all mutations are removed if a creature's contamination level is reduced below 1. Skin, hair, fingernails, and toenails lost to mutations regrow normally once contamination is removed. However, a \textit{regenerate} spell or similar magic is needed to restore any other body parts (such as teeth, limbs, or eyes) lost as a result of contamination levels. Appendages or limbs which develop as a result of mutations wither and fall off when contamination is removed, and other warped body parts are restored to their original form.

\section{Death and Dying while Contaminated}
When a humanoid creature with any contamination levels dies, it animates as a \textbf{haze husk} 24 hours later. A creature with six or more class levels or hit dice rises as a \textbf{haze wight} instead.

Being raised from the dead reduces a creature's contamination level by one.

\section{Monstrous Transformation!}
A creature who reaches contamination level 6 or higher triggers a transformation into a horrific monster controlled by the Game Master. Once triggered, the transformation finishes in 1 round. It is thereafter permanent. Game Masters are encouraged to narrate a suitably gory description of the transformation.

The Game Master chooses the creature's new form, which is most often an aberration or monstrosity of some kind such as a \textbf{delerium dreg}, \textbf{haze hulk}, or \textbf{gibbering mouther} (see Appendix A). The new form usually has a challenge rating similar to that of the original creature, but under appropriate circumstances you can choose one that is much higher or much lower.

The creature's game statistics are entirely replaced by those of the chosen monster. The creature assumes the new form's hit points and hit dice. All contamination levels are removed, though the GM may grant additional traits to the new form similar to any beneficial mutations gained from contamination levels.

The GM determines what remains of the original creature's personality and memories, if anything. Regardless, the creature is invariably driven mad by the transformation and falls under the GM's control.

\subsection{Reversing the Transformation}
The means of reversing a monstrous transformation are completely unknown at the outset of the campaign. A monstrous transformation is widely regarded as permanent, though many still search the ruins for answers. Those who do may eventually be able to research the \textit{siphon contamination} spell.

A \textit{wish} spell or similar magic can restore a single transformed creature to its previous form, removing all madness and contamination in the process. This is considered a stressful use of the \textit{wish} spell, and thus there is a 33 percent chance that the caster will be unable to cast \textit{wish} ever again if they use the spell in this manner.

The transformation is otherwise irreversible by any means short of divine intervention. Nevertheless, at the Game Master's discretion, it may be possible to temporarily alleviate a transformed creature's innate madness. However, an intelligent monster who receives long-term treatment of any sort may eventually conspire to contaminate or subtly corrupt its caretakers.

\begin{DndTable}[header=Madness]{cX}

1d10 & Madness\\
1 & "I wish I didn't have all these useless organs inside me."\\
2 & "The contamination is a blessing which will transform me into a wondrous creature."\\
3 & "The monsters are civilians trying to live a peaceful life! We need to protect them!"\\
4 & "I must wear a flesh-coat made from my slain enemies to gain their strength!"\\
5 & "My companions died in the ruins. I'm sorry friends, you are merely ghosts haunting me, you aren't real. Stop trying to talk to other people."\\
6 & "I need to eat everything I can find. It's probably going to be my last meal."\\
7 & "Drakkenheim is so beautiful at night. I could spend forever wandering the streets by moonlight. We should go tonight! Let's go every night!"\\
8 & "Don't you fools get it?! If you die in Drakkenheim, you die IN REAL LIFE!!!!"\\
9 & "A sinister cabal of disembodied hands is plotting against me."\\
10 & "I must go into the ruins and kill. Rip and tear, until it is done."\\
\end{DndTable}

\chapter{Delerium, Magic \& Spells}
\includegraphics[width=0.4\textwidth]{images/../5e.tools/img/adventure/DoDk/171-13-001.ride.png}
\section{Delerium}
Delerium is a magical mineral left behind by the meteor. It appears in geode clusters of translucent, sharp-edged crystals which reflect octarine light. The eldritch stones softly hum in dissonant tones, and glow brightly at night or when exposed to magic. Deposits are found throughout Drakkenheim, often fused into stone streets and buildings like crystalline moss.

A typical delerium fragment is about the size of a finger. Crystals may be fist-sized or slightly larger, and geodes may be as big as a pumpkin. Massive clusters might grow taller than a human.


\begin{DndTable}{cccccc}

Size & Market Value & Weight & AC & HP & Extraction\\
Chip & 10 gold & 1/4 lbs & 15 & 5 & 1 action\\
Fragment & 100 gold & 1/2 lbs & 17 & 10 & 1 minute\\
Shard & 500 gold & 1 lbs & 19 & 15 & 5 minutes\\
Crystal & 1,000 gold & 2 lbs & 21 & 20 & 30 minutes\\
Geode & 5,000 gold & 20 lbs+ & 23 & 25 & 1 hour\\
Massive Cluster & Priceless & 8000 lbs+ & 25 & 50 & 7 days\\
\end{DndTable}

\DndQuote%
{}
{''"The real meat and potatoes of our city's problem... delerium may be worth its weight in tuna, but is it worth all the trouble that follows?''}
{Veo Sjena}
\subsection{Delerium Properties}
All delerium samples have the following traits regardless of size:


\begin{itemize}
\item Immune to necrotic, poison, and psychic damage; as well as bludgeoning, piercing, and slashing damage from nonmagical weapons.
\item Resistant to acid, cold, fire, and lightning damage; as well as slashing and piercing damage from magical weapons.
\item Vulnerable to bludgeoning damage from magical weapons.
\item Unless contained within an \textit{antimagic field}, delerium shatters and crumbles into worthless ash when reduced to zero hit points. However, delerium geodes release a random \textbf{arcane anomaly} when destroyed (see below).
\end{itemize}

\subsubsection{Delerium Hazards}
When a humanoid creature touches delerium without protective gear (such as gloves or tongs) for the first time on their turn, or ends their turn in bodily contact with delerium, they must succeed on a DC 10 Constitution saving throw or take 1d6 necrotic damage and gain one level of contamination (see Appendix C).

\subsubsection{Delerium Harvesting}
Delerium found in Drakkenheim is usually fused into the ground or stone buildings, and may be carefully extracted using handheld mining equipment such as shovels, picks, hammers, and chisels. Consult the table to determine how long it takes characters to harvest a shard. Lacking proper equipment, extraction takes ten times as long. Massive clusters are impossible to extract without heavy equipment or powerful magic.

\subsection{Uses for Delerium}
Delerium has vast arcane potential. Beyond trade, crystals are used for several purposes:

\subparagraph{Spell Component}
Delerium may be used as an arcane focus or material component for any spell on the sorcerer, warlock, or wizard spell list. When casting a spell which requires a costly material component, a spellcaster may instead use delerium of equivalent value to the required cost.

\subparagraph{Magic Items}
Delerium is an exceptional material for creating magical items of all kinds. Several new magic items are described herein, but any existing magic item described in the Core Rules could be made with delerium. Fragments, shards, crystals, and geodes may be used to make uncommon, rare, very rare, and legendary items respectively. At the GM's discretion, characters who can cast 5th level spells or higher may learn techniques to craft magic items of their own during their downtime.

\subparagraph{Stable Delerium}
The process for crafting magic items with delerium renders the crystals stable. Unless otherwise specified, characters do not take damage nor risk contamination when they touch or handle stable delerium. Stable delerium is damaged and destroyed in the same manner as a normal magic item of its kind.

\subparagraph{Delerium Dust}
Delerium dust is made by grinding the crystals against one another to create a fine powder. This milling process is extraordinarily hazardous unless performed within an \textit{antimagic field}. Delerium dust can be used in alchemy, as a spell component, a reagent for brewing potions, or mixed into inks for Spell Scroll.

\subparagraph{Improvised Weapons \& Ammunition}
Delerium fragments may be used as ammunition for a sling or fashioned into a makeshift club. Such weapons and ammunition deal an extra 1d6 necrotic damage, and non-contaminated creatures struck must succeed on a DC 10 Constitution saving throw or gain one level of contamination. Delerium used as ammunition in this way is destroyed.

\DndQuote%
{}
{''The essence of delerium is intoxicating, like looking upon raw power. I swear you can feel it draw you in. A lesser man might be drawn to the contamination, embracing it fully as a gift.''}
{Sebastian Crowe}
\DndQuote%
{}
{''It's heavy and dangerous like some of the best men I know, including me.''}
{Pluto Jackson}
\begin{DndTable}[header=Arcane Anomalies]{cX}

1d20 & Arcane Anomalies\\
1 & Gravity breaks within a 100-foot radius area for 1 hour. Creatures levitate in midair, and must move by pushing or pulling against a fixed object or surface within reach (such as a wall or a ceiling), which allows them to move as if they were climbing. Unattended objects float around randomly.\\
2 & The nearest creature is affected by a \textit{Tasha's Hideous Laughter} spell (Save DC 15) but instead of laughing, the creature repeats unfathomable combinations of syllables and words.\\
3 & Time skips a beat. Creatures within 60 feet experience a palpable feeling of vertigo followed by a powerful sensation of deja vu and are stunned for 1 round (No save).\\
4 & The nearest creature becomes unstuck in time. It is affected by the \textit{blink} spell for 1 minute. Instead of vanishing into the Ethereal Plane, the creature vanishes into a sliver of time in its past or possible future.\\
5 & A prismatic burst of energy erupts in a 20-foot radius. Creatures in the area must succeed on a DC 15 Constitution saving throw or take 8d6 radiant damage and become blinded for 1 round. The smell of ozone fills the area, and wood transforms into glass.\\
6 & Echoes of possible realities are briefly visible for 1 minute. When a creature within 60 feet is hit by an attack, a faint vision of the creature being killed by that attack appears.\\
7 & Discordant music fills the mind of all creatures within 30 feet, who are affected as if by \textit{Otto's Irresistible Dance} (Save DC 15).\\
8 & A section of stone, water, air, or energy becomes an appropriate contaminated elemental.\\
9 & An extraplanar creature is summoned and remains for 1 hour. The DM either chooses the creature or determines it randomly. It is friendly to the creature who triggered the anomaly.\\
10 & All humanoid corpses within 120 feet animate as hostile haze husks. The shrieking undead beg frantically for forgiveness as they rip apart the living.\\
11 & The shadows of 1d6 random creatures in the area animate and try to kill them. Once destroyed, the creatures don't cast a shadow for 24 hours.\\
12 & All creatures within 60 feet become invisible for 1 min or until they attack or cast a spell.\\
13 & Tendrils of life flow from the nearest creature to others. It must succeed on a DC 15 Constitution saving throw or take 8d8 necrotic damage, half on a success. The three nearest creatures within 60 feet each regain hit points equal to the damage taken.\\
14 & A \textit{hypnotic pattern} (save DC 15) appears. It creates scintillating impossible colours in shapes which are simply wrong. Creatures incapacitated by the spell weep uncontrollably.\\
15 & A \textit{Evard's Black Tentacles} spell appears in the area for 1 hour (Escape DC 15).\\
16 & Objects within 60 feet come to life for the next hour, as if affected by the \textit{animate objects} spell. They mumble awful truths, but are not otherwise hostile.\\
17 & The nearest creature is polymorphed into an \textbf{awakened shrub} for 1 hour or until reduced to 0 hp (No save).\\
18 & Time slows down for up to six randomly determined creatures within 120 feet of the anomaly. They are affected by the \textit{slow} spell for 1 minute. (Save DC 15).\\
19 & Time speeds up for one randomly determined creature within 60 feet of the anomaly. They are affected by the \textit{haste} spell for 1 minute.\\
20 & A bowl of flowers and a very surprised aquatic mammal appear 100 feet in the air. "Oh no, not again..." thinks the flowers.\\
\end{DndTable}

\section{Spells}
This section contains new spells found in the Drakkenheim campaign. While many monsters and NPCs know these spells, \textbf{they are not available to player characters at the outset of the campaign}. Instead, characters must either study the spell from a found spellbook or scroll, locate a mentor to teach them the spell, or research the spells themselves!

\subsection{Learning the Spells}
Once they've found a spellbook or teacher, a character must spend downtime and gold to master the spell before they may cast it. For each level of the spell, the character must spend two days of downtime and 100 gp practicing the spell. Once complete, the character may prepare the spell, add it to their spellbook, or replace a spell they already know with the new spell, as appropriate for the spellcasting features of their class.

\subsection{Researching New Spells}
\textit{Contamination immunity}, \textit{contaminated power}, \textit{neutralizing field}, and \textit{siphon contamination} are totally undiscovered at the outset of the campaign, and must be researched by the player characters themselves. The spell descriptions outline the prerequisites needed to successfully research the spell. Once these requirements are met, a character must spend downtime and gold to complete the research: for each level of the spell, the character must spend 5 days of downtime and 250 gp. Once complete, the character may prepare the spell, add it to their spellbook, or replace a spell they already know with the new spell, as appropriate for the spellcasting features of their class.

\subsection{Secret Spells}
In the world of Drakkenheim, several powerful spells found in the \textit{Core Rules} are significantly rarer than normal: \textit{astral projection}, \textit{gate}, \textit{plane shift}, \textit{simulacrum}, \textit{true resurrection}, and \textit{wish}. At the Game Master's option, player characters must obtain knowledge of these spells in the same manner as other spells presented in this book.

Non-player characters do not have access to these spells unless specifically indicated in their description or game statistics. Furthermore, the material components for these spells are exceptionally scarce and should be treated as consumable magic items of very rare (for 7th and 8th level spells) or legendary rarity (for 9th level spells).

\subsection{Contaminated Spells}
Several spells include the "contaminated" descriptor beside their school of magic. Such spells always require delerium as a material component, and cause the caster to gain contamination when cast, as described in the spell description of each spell. The spellcaster can't prevent or ignore gaining this contamination level by any means, otherwise a contaminated spell simply fails.

\subsection{New Spells}

\begin{itemize}
\begin{multicols}{3}
\item \textit{Conjure the Deep Haze}
\item \textit{Contaminated Power}
\item \textit{Contamination Immunity}
\item \textit{Delerium Blast}
\item \textit{Delerium Orb}
\item \textit{Forced Evolution}
\item \textit{Neutralizing Field}
\item \textit{Octarine Spray}
\item \textit{Purge Contamination}
\item \textit{Ray of Contamination}
\item \textit{Ride the Rifts}
\item \textit{Sacrament of the Falling Fire}
\item \textit{Siphon Contamination}
\item \textit{Warp Bolt}
\end{multicols}

\end{itemize}

\section{Magic Items}

\begin{itemize}
\begin{multicols}{3}
\item \textit{Aqua Delerium}
\item \textit{Aqua Expurgo}
\item \textit{Bottled Comet}
\item \textit{Comet Smasher}
\item \textit{Delerium Crystal Focus}
\item \textit{Delerium-Forged Blade}
\item \textit{Flame Lance}
\item \textit{Hardened Delerium-Tipped Arrows}
\item \textit{Hazewalker Plate}
\item \textit{Purging Rod}
\item \textit{Refined Delerium Dust}
\item \textit{Skymetal Shield}
\item \textit{Skymetal Staff}
\item \textit{Spellpiercing Wand}
\item \textit{Staff of Contaminated Power}
\item \textit{Starcrossed Bow}
\end{multicols}

\end{itemize}

\section{Seals of Drakkenheim}
The Seals of Drakkenheim were jointly crafted by mages of the Amethyst Academy and the clergy of the Sacred Flame, and gifted to the royal family of Drakkenheim after the Edicts of Lumen. The six members of the royal council traditionally each carried one of these potent symbols of political authority, while the \textit{Crown of Westemär} was worn by the monarch. While each faction seeks at least one of these magic items, the Hooded Lanterns and the Queen of Thieves wish to collect all of them so they can determine the political destiny of Drakkenheim.

\subsection{Crown Authority}
A character cannot attune to more than one \textit{Seal of Drakkenheim}. Guardian constructs in Drakkenheim such as wall gargoyles and tower dragons do not attack you while you are attuned to one of the \textit{Seals of Drakkenheim}. An individual guardian ignores this effect for 24 hours if you attack it or cast a spell upon it.

As an action, you may present a \textit{Seal of Drakkenheim} to rebuke the guardians. You make an opposed Charisma ability check against each \textbf{wall gargoyle} or \textbf{tower dragon} that can see or hear you within 30 feet of you. If you win the contest, the creature is rebuked for 1 minute or until it takes any damage. A rebuked creature doesn't move, and can't take actions or reactions. A creature who wins the contest is immune to your rebuke for 24 hours.


\begin{itemize}
\begin{multicols}{3}
\item \textit{Chancellor's Crest}
\item \textit{Inscrutable Staff}
\item \textit{Lord Commander's Badge}
\item \textit{Spymaster's Signet}
\item \textit{Steward's Seal}
\item \textit{Crown of Westemär}
\end{multicols}

\end{itemize}

\section{Relics of Saint Vitruvio}
These holy relics of the Sacred Flame were wielded by Saint Vitruvio and are kept in several religious sites throughout Drakkenheim. The Knights of the Silver Order and the Followers of the Falling Fire now seek these relics. Not only do both factions revere Saint Vitruvio, both wish to use these relics to help advance their goals in Drakkenheim.

A character attuned to three \textit{Relics of Saint Vitruvio} may cast the \textit{heal} spell three times per day, after which they must finish a long rest before they may do so again. In addition, that character gains the ability to attune to two additional magic items so long as those magic items are both \textit{Relics of Saint Vitruvio}.


\begin{itemize}
\begin{multicols}{3}
\item \textit{Ignacious, the Sword of Burning Truth}
\item \textit{Sceptre of Saint Vitruvio}
\item \textit{Helm of Patron Saints}
\item \textit{The Shield of Sacred Flame}
\item \textit{Saint Vitruvio's Phylactery}
\end{multicols}

\end{itemize}

\chapter{The World of Drakkenheim}
This chapter describes the wider world surrounding Drakkenheim. This material is meant to provide additional ideas for players creating characters adventuring in Drakkenheim, and useful background information to help Game Masters understand the context behind the conflicts in the ruined city. Perhaps it might also inspire future adventures!

\includegraphics[width=0.4\textwidth]{images/../5e.tools/img/adventure/DoDk/179-14-001.sharpen.png}
\section{People of the Continent}
Many different people lived in Drakkenheim. The neighbouring countryside and the realms beyond are no less diverse, though humans make up the majority of the population in the current age. In ancient days, the industrious dwarven folk built great underground empires, the alien elves raised mystical towers, while the orcish people constructed monuments and menhirs to honor their old gods. Today, their descendants and those of many other ancestries inhabit the continent, though few would associate their varied biological distinctions with any single nation, culture, or religion.

\subsection{Dwarves}
Many dwarves tell tales of ancient underground kingdoms that hoarded vast mines of gold and gems and fought great wars miles below the surface. The foundations upon which Drakkenheim were built date back to these times, and traces of dwarven architecture and engineering can still be found in the fortifications, aqueducts, and deep vaults of the city. However, the days of the dwarven realms have long passed, and the dwarven people have no great holds of their own any more. Still, many dwarves dwell in the mountains and keep the traditions of their ancestors: mining, smithing, and engineering. Many more are much less tied to these traditions now, so it is not unusual to see a small number of dwarves in just about any community.

\subsection{Elves}
No one knows when elves came to these lands, or exactly how they got here. Even the most ancient among them tell little more than vague myths about their origins. These stories claim that millennia ago, the elves came into this world through fairy realms and shadow gates from a distant and long-lost land. Why the elves left, none can say. Today, the elves live in diaspora. Only a handful of elven families dwell in even the most populous regions. Although rare, they are no longer regarded as strangers to these lands. In fact, half-elves are now much more common than those of pure elven descent. Though elves remember little of their history, they still possess great knowledge of magic. The few elven ruins found across the continent contain mysterious portals and arcane lore. The reason they were abandoned remains unknown.

\subsection{Halflings}
These smallfolk often form their own communities within settled lands, just as often as they live alongside the other people of the continent. Many live a rural life, tending farms, orchards, and fishing villages. Dozens of halfling villages surrounded Drakkenheim, the bounty of their fields feeding the city. Though these communities were not outright destroyed when the meteor struck, they have not escaped its devastation. The eldritch contamination spreading from the ruins and the threat of monstrous raiders have forced many smallfolk to abandon their homes. Their once-quaint underhill abodes and welcoming hamlets are now eerie ghost-towns.

\subsection{Humans}
Most of the population of the continent are humans, who range in all shapes, sizes, colours, skills, and beliefs. Humans tell countless tales of their origins, and tend to regard wherever they were born as home. They formed many of the great noble and royal houses of the continent, and the endless conflicts between these families have defined the political landscape of the realms. One such family, House von Kessel, ruled Drakkenheim until the meteor and years of civil war ended their regal line.

Humans have long mingled with other peoples, resulting in half-elves, half-orcs, and half-dragons. Though such people often bear distinct physical traits, such individuals are recognized for their shared human origin. Indeed, humanity has often intentionally sought to enhance their lineages. During antiquity, many powerful sorcerous dynasties infused their bloodlines with the essence of infernal and celestial entities, bequeathing their progeny with an enduring magical legacy. Though these old empires have long since fallen, their descendants have since spread throughout the world.

\subsection{Uncommon Origins}
Though their populations are sparse, many other peoples inhabit the world of Drakkenheim, including catfolk, dragonborn, gnomes, goblins, orcs, and tieflings. No matter their origin, there is no reason any type of character would be out of place arriving in Drakkenheim. Players should work with the Game Master to invent how their ancestry choice could have a place in Drakkenheim.

\section{Realms of the Continent}
The people who inhabit the world of Drakkenheim simply call it "the earth" or "the world". They live on "the continent" that is divided into several nations which people refer to as "the realms". The largest and most politically influential of these sovereign states are Caspia, Elyria, and Westemär. Neighbouring each are several petty kingdoms and dozens of city states.

\includegraphics[width=0.4\textwidth]{images/../5e.tools/img/adventure/DoDk/180-14-002.map.png}
\subsection{Westemär}
The Realm of Westemär stretches north from the Drakeclaw and Glimmer Mountains, along the Crystal Coast to northern tundras. The heartland is made up of deep dark forests surrounding wide flowing rivers, interspersed with areas of fertile marshland. The land is prosperous and rich in natural resources: iron, salt, and coal mines, fine timber, and excellent quarries. These industries supported the realm's once-great military strength and built the impressive architecture of its cities. As such, the nobility are distinguished, but materialistic and prideful, in contrast to the hard-working and pragmatic common folk.

\subsubsection{A Broken Realm}
Drakkenheim was the capital of Westemär and home to the ruling family, \textbf{House von Kessel}. After the meteor fell, the two surviving scions of House von Kessel clashed over succession in a brutal civil war that tore the nation apart. After many assassinations, accidents, and disastrous battles, the civil war ended with no clear victor. Most believe House von Kessel is now extinct, and the throne remains abandoned and empty in the ruins of Castle Drakken. The lands surrounding Drakkenheim are now contaminated, and most villages nearby are deserted.

Lacking a sovereign ruler, the nobility manages their own domains directly. They remain deeply divided based on which branch of House von Kessel they supported during the civil war. Many grievances remain, and occasional skirmishes occur between the squabbling barons, counts, and dukes. However, after nearly a decade of fighting, their coffers are empty and their blood is spent. Few possess the will or resources to continue these feuds for the time being. Many noble families possess some relation to the royal line, yet none have the political clout or military strength to press a claim upon the throne and reunify the nation.

\subsection{Elyria}
The realm of Elyria lies on a large peninsula south of the Drakeclaw Mountains. The land comprises vast rolling prairies and river basins broken up by rocky arid hills and a mountainous coast. Nestled within are pockets of imposing pinewood forests, and great deposits of precious minerals such as gold and silver. Elyrians are known for being kind-hearted and outgoing folk eager to share the bounty of their land. Most live a rural life in small villages along the rivers and coasts. Indeed, the fertile river valleys at the heart of Elyria were an early cradle of civilization, and the countryside is dotted with the ruins of ancient cities.

The Holy City of Lumen is the capital of Elyria. The Faith of the Sacred Flame was founded here, and the shining city is the seat of power for the \textbf{Divine Matriarch}. The most fervent and orthodox adherents to the faith are found amongst the Elyrian people. Beyond Lumen, there are few large cities save several bustling port towns that are hubs for international trade. The Elyrian navy is amongst the finest in the continent.

\subsubsection{A Rising Theocracy}
Originally, Elyria was a monarchy, but today, the nation is effectively a theocracy ruled by the Divine Matriarch. As leader of the Faith of the Sacred Flame, the Divine Matriarch has traditionally been regarded as a pillar of wisdom, truth, and justice, so they often mediate difficult political matters. When the Elyrian royal line went extinct, a parliament drawn from the Elyrian nobility determined that the Divine Matriarch would appoint a Lord Regent to lead them. Together, the noble parliament and the Lord Regent would serve as a ruling body for the nation.

However, in recent decades, successive Divine Matriarchs have used this power to select a line of caretaker regents who defer to the Faith on virtually all matters of state. While members of the clergy take vows of poverty that preclude them from holding actual political station, the parliament is now mostly made up of fundementalist nobles staunchly loyal to the religion, and the Lord Regent is almost always a close relative of a prominent High Flamekeeper or the Divine Matriarch. As a result, the Divine Matriarch wields \textit{de facto} political authority in Elyria.

\subsection{Caspia}
The realm of Caspia stretches in an expanse between the Glimmer Mountains and the Diamond Sea. It is a rugged and mediterranean nation, home to a boisterous and affable people led by a constantly squabbling nobility.

\subsubsection{The Kingsmoot}
Caspians elect their High King once each decade, from candidates put forth by the princes of their noble houses. The valorous deeds and reputation of the noble houses is a major factor in determining the High King. Though Caspian politics are always shifting, especially as the next Kingsmoot draws near, it has stood as a remarkably stable nation. As such, Caspians love to meddle in the politics of Elyria and Westemär. They often intermarry with the nobles of different realms to gain greater clout and influence for the next Kingsmoot.

\subsection{Wars and Conflicts}
The realms have been shaped by centuries of war, from Elyrian rebellions against ancient sorcerer-kings, to great battles between rival Caspian princes, and most recently, the civil war that ravaged Westemär.

By law and custom, none with arcane spellcasting powers may hold a noble title within the realms. Instead, the noble houses prize the martial arts. Many fighting traditions have developed: armoured phalanxes, archery, fencing, berserking, and diverse unarmed combat techniques. Duels and trial by combat are a particularly favoured means to settle disputes amongst the nobility with minimal bloodshed. The realms of Caspia, Elyria, and Westemär have each developed their own distinctive styles and practices for such contests.

\subsection{Travel and Trade}
Commerce, rather than conflict, defines everyday life for common folk in the realms. Merchant caravels and ships frequently carry cargo across the Diamond Sea from the Caspian port of Jackson's Keep or the independent city-state of Liberio. Other vessels depart from the ports of Drannsmund in Westemär and Angel's Landing in Elyria, and sail the Middle Ocean to the far west. Trade caravans from Geldstadt travel across the Twin Vales to distant lands beyond the steppes.

Magical goods and services are bought and sold throughout the continent, but only amongst the wealthy and powerful. High-ranking nobles outfit their favoured knights with magical arms and armour, and rich merchants buy spells and wondrous items to safeguard their valuable mundane wares. The manufacture and sale of arcane commodities is controlled by the Mage Guild, a branch of the \textbf{Amethyst Academy}. Under strict edicts recognized across the realms, any mage or trader dealing in magic pays an annual fee to the Mage Guild, and must follow the prices set by the guild.

A few are bold enough to skirt these requirements. As a result, smugglers, hired muscle, and rogues can make their fortune in the illicit trade of enchanted gear and forbidden magic items. This underground black market has only expanded since delerium has become such a valuable and sought-after export.

\subsection{Free Lands}
Though the realms boast mighty armies and legal customs, these have little authority outside cities and settled farmlands. Not far off the major trade roads and rivers, the wilderness remains a place ruled by the laws of nature, not the laws of mortals. The dark forest of Achtungwald, the wet bogs of the Elfmire, the massive mountains that separate the realms, the northern tundras, and the arid hills of Elyria are dangerous lands, where dragons, trolls, giants, and monsters keep their lairs.

However, these wild places are also home to many people who do not reside under the control of the realms. Living in small tribes, they keep their old ways and old faiths, led by stout-hearted warriors and druids who embrace the fury of the elements.

\subsection{The Wider World}
The continent meets the vast Middle Ocean along its western coast. Across this body of water is a distant continent. While bold and well-travelled merchants often make the voyage, few folk in the realms know anything of those antique lands.

East of Westemär are the Twin Vales, nestled between great mountain ranges. Beyond these are expansive steppelands and vast deserts that lead to parts unknown. Off the Crystal Coast, to the north, is the island of Skye, further beyond is the North Pole. South of Caspia is the Diamond Sea; across are realms that often trade with Caspia and Elyria.

\section{Magic and Spellcasters}
Magic suffuses the world of Drakkenheim, but mortal spellcasters are rare. Only about one in every thousand people can cast spells, and roughly two-thirds of these only ever master cantrips and first-level spells in their entire lifetimes. The few who can work greater magic often wield tremendous personal, social, and political power. Though rare, magic is common currency for nobility, clergy, wealthy merchants, and adventurers.

Only a handful of living mortals can cast ninth-level spells. Most are secretive recluses or legendary figures. Harnessing more magical might than most mortals can even comprehend, they often turn their attention away from earthly affairs to contemplate the cosmic order, fighting unseen battles against supernatural foes whose very nature threatens to tear the world asunder.

The common folk regard spellcasters with superstition and awe. There was a time when arcane magic users faced extreme persecution, and echoes of those prejudices remain. Few even understand the difference between sorcerers, warlocks, and wizards—they are collectively referred to as 'mages'. However, commoners can recognize spellcasters who are clerics of the Sacred Flame or mages of the Amethyst Academy by their symbols and attire. Spellcasters who are not (or do not claim to be) members of these recognized organizations are often mistrusted and feared by commoners.

\subsection{Arcane Magic}
The arcane magic of sorcerers and wizards draws upon otherworldly energies. The ability to tap into these forces is reliant on a recessive genetic trait in mortals - such individuals are called mageborn. Signs that a person is mageborn typically emerge in late childhood or early puberty, usually as random displays of wild magic. In some rare cases it does not develop until adulthood. Mortals who lack the mageborn gene cannot become sorcerers or wizards under any circumstances.

By law, the Amethyst Academy takes guardianship over any child born in the realms who manifests magical ability, whether they are a royal heir or commoner. A child who bristles against Academy training—or who was not taken in by them—most often becomes a sorcerer.

Bards and warlocks also work arcane magic, but do not necessarily depend on mageborn lineage to do so. Instead, warlocks gain arcane power through pacts with eldritch entities, and bards find supernatural inspiration to harness the Words of Creation. Regardless, many keep their powers secret.

\subsection{Divine Magic}
In the world of Drakkenheim, divine spellcasting powers are not granted by gods. Instead, clerics, druids, and paladins learn to cast spells by channeling the supernatural forces of light and shadow that underscore the cosmos. In theory, such powers are attainable by anyone, but in practice only exceptionally gifted individuals master even the lowest-level divine spells. As such, divine spellcasting is viewed as a disciplined and sacred technique attainable only by those with sufficient will and devotion, and only after years of meditation and training. Its practitioners are not considered mageborn.

Violating one's religious tenets will not cause a cleric, druid, or paladin to lose their powers. However, a personal crisis of faith brought on by some discordant change in behaviour and morality could cause an individual to lose their divine spellcasting powers. More often, a cleric, druid, or paladin who commits some grave trespass risks retribution from other members of their religion.

\subsection{Other Spellcasters}
There are many rudimentary techniques for casting a small smattering of cantrips and lesser spells, represented by rangers and subclasses or feats that grant minor spellcasting features. Such characters are viewed as merely dabblers, not mageborn.

\DndQuote%
{}
{''Warmongers that hide behind divine purpose are still warmongers. My people are not perfect but we are honest with our enemies and ourselves.''}
{Pluto Jackson}
\DndQuote%
{}
{''I've seen plenty of Elyrians set places on fire, but when I do it, it's a "crime" and I am "dangerous."''}
{Sebastian Crowe}
\section{Faith and Religions}
In the world of Drakkenheim, gods do not physically manifest in the world, intervene in mortal affairs, or speak directly with mortals. Myths and legends passed down from distant eons tell such tales, but these events have not happened in recorded history. Instead, the inspired words of prophets, oracles, and heroes laid the foundations for the world's many religions.

Powerful extraplanar entities such as archangels, demon lords, and outsiders exist and interact with mortals. Many extensively meddle in earthly matters, making pacts and offering dark counsel. Some even claim to be gods. When asked questions regarding the true nature of the cosmos or divinity, however, they speak in cryptic riddles.

\subsection{The Sacred Flame}
This is the dominant religion throughout the realms. Also known as the Divine Light, its teachings and beliefs are tightly woven into the social and political fabric of Caspia, Elyria, Westemär, and their neighbouring domains.

\subsubsection{Background}
Myths about the Sacred Flame have ancient roots, but the contemporary religion began a thousand years ago with the valiant \textbf{Saint Tarna}. Although most of her life she was a ruthless warlord, light flickered in her dark heart. Guided by angels, she learned to channel the Sacred Flame and became the first paladin. She took up the cause of righteousness and sacrificed herself to defeat an evil sorcerer-king. After her martyrdom, her disciples spread her teachings and established the Faith of the Sacred Flame.

\subsubsection{Beliefs and Tenets}
Faithful of the Sacred Flame do not worship a deity. Rather, devotees pray for guidance from a transcendent divine force that they believe is the source of all light and goodness. Symbolized as the Sacred Flame, this brilliant beacon inspires mortals to act with benevolence, honor, and righteousness. The religion brings a message of dutiful hope: even the smallest flame may light the way through the darkest night. The core tenets of the religion are conceptualized as follows:


\begin{description}
\item[The Hearth.]
spreads compassion. One must nourish the hungry; offer shelter to the cold; and give succour to those who suffer. Spiritual warmth comes when the community is gathered together before the Flame as neighbours.
\item[The Lantern.]
illuminates the way of truth. The light reveals what is hidden by darkness, and shows the way to understanding. One should not silence truth, nor speak falsehoods.
\item[The Torch.]
displays the duty of all faithful to uphold the cause of justice. One cannot abide when evil stirs; nor stand idle when corruption spreads: callousness and indifference permit selfishness, greed, and hate to multiply.
\item[The Candle.]
symbolizes hope and redemption. Show those who are in darkness the light. Offer mercy to the guilty; no shadow is absolute so long as a flickering flame remains.
\end{description}

Followers of the Sacred Flame believe that when mortals die, the divine light guides righteous spirits to the place where dawn breaks over the Shadowlands. Their luminous souls shine eternally, driving back the endless darkness that threatens all life.

\subsubsection{Holy Sites and Symbols}
A goblet of fire or a lit candle represents the Sacred Flame in holy symbols. Believers join hands in a circle around a roaring fire and sing hymns during worship rites. The holy text of the faith is called the \textit{Song of Fire}. Saint Tarna is typically depicted as a silver-haired warrior bearing a longsword. She is often shown in battle against demons and witches, sometimes astride her griffon, Aarak.

In Caspia, Elyria, and Westemär, most towns have at least one chapel devoted to the Sacred Flame. Often at the heart of their communities, these domed sanctuaries are built around a hallowed brazier set alight with a golden \textit{continual flame}. Major cathedrals to the Sacred Flame are found in large cities, and many are brilliant architectural wonders in their own right, decorated with painted murals, stained-glass windows, and statues depicting the saints and martyrs. The small shrines in most villages are usually simply an awning built over a bonfire pit.

\subsubsection{Religious Hierarchy}
Ordained priests of the Sacred Flame are known as \textbf{Flamekeepers}. These clerics of the Sacred Flame are predominantly women*; they wear flowing vestments of white, yellow, and gold. Flamekeepers take vows of poverty, eschewing all personal possessions and living on the commonwealth of the faith. Chapels and cathedrals of the Sacred Flame are ministered by a Flamekeeper, who is assisted by a few devoted acolytes. Smaller villages rely on lay ministers to keep their shrines. The highest authority amongst the priesthood of the Sacred Fire is the \textbf{Divine Matriarch}, who tends the Cathedral of Saint Tarna in the holy city of \textbf{Lumen} in the land of Elyria.

\textit{* Although this represents the formal religious hierarchy, people of any gender identity may become clerics and paladins of the Sacred Flame and rise to positions of honour and reverence}.

\subsubsection{Clerics of the Sacred Flame}
Through deep spiritual discipline and steadfast devotion, some faithful may become living vessels for the Sacred Flame. These rare individuals hold the power to channel its light in the earthly realm to work divine magic. As clerics of the Sacred Flame, they become enlightened healers, custodians of truth and knowledge, or fiery beacons of light. However, not every cleric of the Sacred Flame becomes a Flamekeeper or participates in the organization and politics of the religion.

\subsubsection{Paladins of the Sacred Flame}
Legend claims Saint Tarna was the first paladin. These warriors hold deep reverence amongst the faithful, as only truly virtuous souls can wield the Sacred Flame as she once did. Paladins of the Sacred Flame swear sacred oaths of devotion, redemption, and vigilance. Inspired by Saint Tarna's example, many join knightly orders ordained for righteous purposes. Such companies of holy warriors and their martial retainers embark on quests and crusades to destroy supernatural forces of evil and root out creatures of otherworldly darkness. The most well-known and decorated amongst these militant fellowships are the Knights of the Silver Order.

\subsection{Other Religions}
Although many people keep the Flame, there are hundreds of smaller religions with disparate beliefs, doctrines, and practices. Here are some of the most notable.

\subsubsection{The Old Faith}
This disparate religion is dedicated to a pantheon of primal deities who have whispered to mortals through nature for untold generations, such as primal Nodens, vengeful Kromac, honorable Nuada, and nurturing Danu, but there are countless more. Their beliefs and practices vary widely based on the tenets of the specific patron god, but their myths often surround the elements, the seasons, the land, and ancestor worship. Many such beliefs originated with the ancient orcish and dwarven peoples of the continent. Others are closely connected to mysterious entities who inhabit extraplanar worlds such as Dreamland, the Elemental Chaos, or the Feywild, and are often kept by the elves. Druids are particularly common adherents and priests of the Old Faith, but clerics and paladins of the Old Faith are not unknown.

\subsubsection{Shadow Cults}
Throughout the ages, various 'shadow faiths' and heretical sects have emerged surrounding the Sacred Flame, some of which embrace an opposing force called the Exalted Darkness. These cults develop their own mysteries and esoteric practices, or devote themselves to worshipping enigmatic god-like beings such as Morrigan the Phantom Queen, the Night Serpent, the Lord of the Undead, and powerful entities who inhabit the Shadowlands or the Abyss.

Some of these chthonic entities seek to devour the light and plunge the world into darkness unending. However, not all of these faiths are evil, with practitioners who view shadow and light as part of a necessary balance. Regardless, followers of the Sacred Flame condemn these beliefs, banishing their rites and practices. In the same manner that deep faith and contemplation leads followers of the Sacred Flame to become clerics and paladins, one may also channel the Shadow itself. Clerics are heralds of twilight, trickery, magic, and death, and paladins often invoke dark oaths to vengeful causes or ambitious conquest. Followers of the Shadow Faith may even become warlocks, embracing dark and fiendish patrons.

\DndQuote%
{}
{''I am told many a story about the exploits of brave Caspian warriors, and I take them in the same way I do my meals... with a grain of salt.''}
{Veo Sjena}
\DndQuote%
{}
{''The name breathes life into me. Glory to the great Caspia! I can hear the chants and cheers in the halls clearly in my head. I do everything for them.''}
{Pluto Jackson}
\section{The Edicts of Lumen}
During antiquity, mighty sorcerer-kings ruled a continent-spanning empire using their immense arcane power. The sorcerer-kings became tyrants who scarred the land with magical wars and held those without arcane magical power in bondage. The iron rule of the sorcerer-kings was ended by the clerics and paladins of the Faith of the Sacred Flame a thousand years ago. However, those born with arcane magical ability have faced discrimination ever since, and at times, the religious ministry has endorsed vicious persecution against mages. Even today, many common folk consider a mageborn child a terrible curse.

Outright conflict ended with the Edicts of Lumen, a magical treaty between the Amethyst Academy, the continental nobility, and the Faith of the Sacred Flame. Signed three centuries ago, this agreement separates the political and economic powers of clergy, mages, and nobles.

The Edicts of Lumen comprise six main articles:


\begin{itemize}
\item The \textbf{Articles of Inheritance} bar all mages from holding noble titles or owning land.
\item The \textbf{Articles of Neutrality} require the Amethyst Academy remain neutral and impartial with regards to the social, economic, and political affairs of the continent, but also grant protection under the law to all mages.
\item The \textbf{Articles of Guardianship} grant the Amethyst Academy the authority to take any child born with magical capabilities as wards to train in magic until they come of age, including mageborn of noble houses and royalty.
\item The \textbf{Articles of Enterprise} permit the Amethyst Academy to operate as a self-governing mage guild to control the sale of magical goods and arcane magical services throughout the continent, but they may not discriminate or charge their clients differently.
\item The \textbf{Articles of Malediction} disallow the Academy from teaching magical practices such as summoning demons and creating undead, and forbid mages from using magic to directly influence or control members of the nobility or clergy.
\item The \textbf{Articles of Umbrage} detail a complex system for resolving minor infractions in a peaceful manner, reparations for damages for more serious violations, and the circumstances when parties may use violent force to enforce compliance.
\end{itemize}

\section{The Cosmos}
Sages of the Amethyst Academy have put forth differing theories to explain the order of the cosmos throughout history. These included bizarre concepts such as the "great wheel", the "world axis", the "conjunction of the spheres", "planetary orbits", and "general relativity". There is no consensus; every theorist is quick to point out the glaring flaws and logical inconsistencies in competing hypotheses.

Most mortal knowledge of the planes comes from contact with extraplanar creatures via summoning magic or divination spells such as \textit{contact other plane}. Deliberate planar travel by mortals is the stuff of legends. The knowledge and material components for the \textit{plane shift} spell are jealously guarded secrets. However, there are countless ley-lines and thin places where the unwary might slip between worlds. Described below are the most well-known, but countless others exist.

\subsection{Mortal Worlds}
The mortal world is a sphere of rock that orbits a fiery sun. Several moons and planets join this cosmic dance, all suspended in an endless astral void. Far away, the sparkling stars are other suns. Who can say what distant worlds lie among them?

\subsection{Otherworlds}
These worlds are metaphysically close to the mortal world, such that some creatures can pass through the Ethereal Plane into the mortal world, and rarely vice-versa.

\subparagraph{The Ethereal Plane}
A truly liminal space, this plane is the metaphysical border between the mortal world and the other planes.

\subparagraph{Dreamland}
This strange realm is the mortal unconscious made manifest. It is a world of thoughts, dreams, and nightmares, inhabited by impossible creatures.

\subparagraph{The Feywild}
The realm of fey is sometimes called by different names, such as Tír na nÓg, Wonderland, Avalon, Elfhame, and the Elusive Realm. It is uncertain whether or not these names refer to regions within the realm or are in fact distinct worlds entirely.

\subparagraph{The Elemental Chaos}
A primordial place of raw energy and unbridled potential, the individual planes of air, earth, fire, and water are found within its churning storms and burning seas.

\includegraphics[width=0.4\textwidth]{images/../5e.tools/img/adventure/DoDk/181-14-003.gun.png}
\subsection{Afterworlds}
The afterworlds are places where mortal souls go when they die. There has been no recorded instance of a living mortal physically entering these planes and returning. However, the creatures who inhabit these worlds can be contacted or summoned using powerful spells.

\subparagraph{The Shadowlands}
A dark reflection of the mortal world held in perpetual gloom. Most souls pass into the Shadowlands on their way to the afterworlds, but some wander here forever. Why certain souls remain in this place, and others pass on is the subject of many religious myths. The energies from this place may be invoked to create undead.

\subparagraph{Eternity}
The sacred celestial afterworld where angels dwell. Little is known about this holy place, but many religions believe it is a paradise for righteous souls.

\subparagraph{The Abyss}
Interchangeably referred to as Hell, sages believe that mortal souls are transformed into demons or devils in this fiendish afterworld. These evil beings plot mayhem and conspiracy against mortals when they are not waging endless war amongst themselves. Few mortals understand the distinction between such fiendish creatures.

\DndQuote%
{}
{''There are many reasons to be here, but you need to know why and finish it. Your life is in the balance and it is not worth a meaningless pursuit.''}
{Pluto Jackson}
\begin{DndTable}[header=Historical Timeline]{cX}

Year & Event\\
Year -12 & A golden comet leads Tarna to the archangels Gabrielle and Michael, and she becomes the First Paladin.\\
Year 1 & The Martyrdom of Saint Tarna, who saved the world from the madness of sorcerer-king Xandor XIII. After her martyrdom, her followers share how to channel the Sacred Flame and war against the sorcerer-kings.\\
Circa 300 & Rule of the sorcerer-kings finally crumbles. More learn to wield the Light of the Sacred Flame. These faithful clerics organize the early religion.\\
Year 381 & Vowing tyrannical mage-lords will never rule again, the newly-established Divine Matriarch of the Sacred Flame endorses persecution against arcane spellcasters.\\
Circa 500 & Following a century of persecution, a group of surviving mages form the Amethyst Academy as a secret society to protect themselves.\\
Circa 600 & House von Drakken establishes a petty kingdom along the Drann River, where they lay the foundations of Castle Drakken.\\
Year 678 & The Knights of the Silver Order are ordained to hunt down the Amethyst Academy.\\
Circa 700 & Bloody wars break out between noble houses supported by the Amethyst Academy and houses supported by the Faith of the Sacred Flame.\\
Year 743 & The Edicts of Lumen bring an end to the wars by establishing a clear separation between mages, clergy, and nobles.\\
Circa 800 & House von Drakken goes extinct, setting off 150 years of succession crises.\\
Year 965 & House Von Kessel claims the throne of Drakkenheim and establishes the Kingdom of Westemär.\\
Circa 1000 & The monarchy of Elyria crumbles; the Divine Matriarch appoints a Lord Regent.\\
Year 1,015 & The Clocktower of Drakkenheim mysteriously stops.\\
Year 1,093 & Lucretia Mathias pens an essay predicting the early return of Tarna's Comet, and that its earthly arrival heralds the dawn of a new age for the faithful of the Divine Light. Her apocalyptic predictions are extremely unpopular within the mainstream clergy and her writings are suppressed.\\
Year 1,111 & The meteor strikes Drakkenheim on September 16th at 8:13 PM.\\
Year 1,112 & A failed attempt to reclaim Drakkenheim results in thousands of soldiers dying in the Haze. The few who return are stricken with madness.\\
Year 1,113 & Lucretia Mathias greatly expands her original work and publishes the Testament of the Falling Fire. The accuracy of her early prediction garners her a significant following, and she plans a personal pilgrimage to Drakkenheim.\\
Year 1,114 & After two more military expeditions to reclaim Drakkenheim fail disastrously, the succession crisis between the surviving members of House von Kessel boils over, setting off the Westemär civil war.\\
Year 1,115 & As the civil war rages, several groups of adventurers investigate the ruins of Drakkenheim for plunder. Unexpectedly, some return alive. They speak of horrible monsters within, and bring strange samples of a dangerous crystal.\\
Year 1,116 & Lucretia Mathias completes her pilgrimage to Drakkenheim, personally reaching the Crater's Edge. She returns to Elyria to spread her teachings.\\
Year 1,117 & The Amethyst Academy starts hiring adventurers to collect the "meteor shards" for study. They name the strange crystals delerium.\\
Year 1,119 & Mannfred von Kessel is assassinated. Cecilia von Kessel dies mysteriously. The civil war fizzles out with no clear successor.\\
Year 1,120 & Lucretia Mathias is branded a heretic and excommunicated by the Divine Matriarch of the Church of the Sacred Flame. She returns to Drakkenheim.\\
Year 1,121 & More adventurers venture into the ruins of Drakkenheim. Emberwood Village becomes a hub for treasure-seekers and cutthroats alike.\\
Year 1,122 & Belief in the Falling Fire spreads rapidly, and more pilgrims make the dangerous voyage to Drakkenheim. Few are successful.\\
Year 1,123 & The 4th Provisional Expeditionary Force to Reclaim the Capital—the Hooded Lanterns—are established and make their first forays into Drakkenheim.\\
Year 1,124 & The Hooded Lanterns finally establish a base of operations at the old Shepherd's Gate and then the Drakkenheim Garrison, but cannot make further progress into the city.\\
Year 1,125 & A small group of the Followers of the Falling Fire established a foothold in Saint Selina's Monastery to aid future pilgrims to the crater.\\
Year 1,126 & March 4—Dungeons of Drakkenheim begins.\\
\end{DndTable}

\chapter{Character Backgrounds}
\includegraphics[width=0.4\textwidth]{images/../5e.tools/img/adventure/DoDk/182-15-001.backgrounds.png}
Many types of Adventurers travel to the ruins of Drakkenheim in hopes of uncovering things once lost. Whatever their call to the city may be, there is a chance that their history is ripe with danger. Some have colourful pasts that have readied them for the dangers ahead, some have crimes and mistakes they seek to hide from in the ruins, and others are bound to those who fight for what's left of the crumbling empire.


\begin{itemize}
\begin{multicols}{3}
\item Continental Nobility
\item Mageborn
\item Devoted Missionary
\item Survivor
\item Treasure Seeker
\end{multicols}

\end{itemize}

\includegraphics[width=0.4\textwidth]{images/../5e.tools/img/adventure/DoDk/185-15-004.concept.png}
\chapter{Credits}

\begin{description}
\begin{multicols}{2}
\item[Written by:]
Monty Martin & Kelly McLaughlin
\item[Lead Design:]
Monty Martin & Kelly McLaughlin
\item[Art Direction:]
Marius Bota
\item[Writing:]
Monty Martin & Kelly McLaughlin
\item[Commentary:]
Jill Danaitis, Joe O'Gorman
\item[Editing:]
Bob Everson, Monty Martin, Kelly McLaughlin
\item[Proofreaders:]
Bob Everson, Monty Martin, Kelly McLaughlin, Hannah Peart
\item[Maps:]
Made using Dungeon Draft
\item[Playtesting:]
All of the wonderful backers who were kind enough to help us out!
\item[Graphic Design:]
Martin Hughes
\item[Interior Illustration:]
Jokubas Uogintas, Marzena Nereida Piwowar, Mihai Radu, Toni Munteanu, Andreia Ugrai, Mike Daarken, Elizabeth Peiró, Decio Junior, Derek Murphy, Gaga Turmanishvili, Erel Maatita, Vasburg, Mauro Alocci, Andre Garcia, Russell Jones, Karin Wittig, Stephen Sykes, Augusto Costa, Andrei Iacob, Oliver Medak, Chad Lewis, Humble Squid, Raven Rusch, James Patel, Bastian Wilk, Clement Blum, Karina Igorevna Pavlova, Sam Carr, Katariina Sofia Kemi, Daniel Correia, Fajri Muhammad, Cristiana Leone, Alberto Dal Lago, Quintin Gleim, Foeh, Vita, Marius Bota, Josh Orchard
\item[Cover Art:]
Toni Munteanu
\item[Product Design:]
Matthew Witbreuk, Martin Hughes, Simon Sherry, Nick Ingamells, Marius Bota, Joshua Orchard
\item[Fulfilment:]
Matthew Witbreuk, Nick Ingamells, Rosa Man
\item[Marketing:]
Hannah Peart, Kathryn Griggs, Eduardo Cabrera
\item[Special Thanks:]
Shay Bytheway, Jill Danaitis, Kierston Drier, Kyle Drier, Joe O'Gorman.
\item[Product Identity:]
The following items are hereby identified as Product Identity, as defined in the Open Game License version 1.0a, Section 1(e), and are not Open Content: All trademarks, registered trademarks, proper names (characters, deities, etc.), dialogue, plots, storylines, locations, characters, illustrations, and trade dress. (Elements that have previously been designated as Open Game Content or are in the public domain are not included in this declaration.)
\end{multicols}

\end{description}

\includegraphics[width=0.4\textwidth]{images/../5e.tools/img/adventure/DoDk/186-16-002.thank-you.png}
Beyond the incredible team who helped put this book together, the authors would like to express our deep gratitude for our viewers, subscribers, Patreon supporters, and Kickstarter backers. We would not be where we are today if it weren't for the encouragement and support you have shown for both us and \textit{Dungeons of Drakkenheim}. Whether you joined us all the way back in 2018 when we streamed our first episode or are discovering Drakkenheim for the first time with this book - however you found us, we're so glad that you did and that we've been able to share this world with you. We truly thank you from the bottom of our hearts.

Furthermore, a huge thank you to Jill Danaitis \& Joe O'Gorman - our very own Veo Sjena and Pluto Jackson! These two helped shape the world of Drakkenheim and the many many adventures within the ruins. \textit{Dungeons of Drakkenheim} truly wouldn't be what it is without them. We'd also like to thank Kyle Drier: the unseen hand of Drakkenheim! Kyle has been instrumental in helping produce and edit our livestreams and podcasts. We are so grateful to have such positive, kind, giving and wonderful friends to collaborate with and see where the road to Drakkenheim leads.

Finally, we also want to thank our incredibly patient and loving partners - Kierston Drier and Shay Bytheway - who have put up with us through long hours spent writing and producing this project. They have been with us first hand through the thick and thin of our entire careers on YouTube, and we are so grateful for their love and support over the years.

\chapter{Poster Map: City of Drakkenheim}
\includegraphics[width=0.4\textwidth]{images/../5e.tools/img/adventure/DoDk/189-map-0.01-drakkenheim.png}
\includegraphics[width=0.4\textwidth]{images/../5e.tools/img/adventure/DoDk/190-map-0.01-drakkenheim-player.png}
\includegraphics[width=0.4\textwidth]{images/../5e.tools/img/adventure/DoDk/191-map-0.01-drakkenheim-delrium.png}
\includegraphics[width=0.4\textwidth]{images/../5e.tools/img/adventure/DoDk/192-map-0.01-drakkenheim-delrium-player.png}
\chapter{Creatures}
\setlist[description]{nolistsep,listparindent=3pt,topsep=6pt,font={\bfseries\sffamily\small}}
\pagestyle{empty}

\begin{DndMonster}[float*=b,width=\textwidth + 8pt]{Aboleth}
\begin{multicols}{2}
\DndMonsterType{Large aberration, lawful evil}
\DndMonsterBasics[
armor-class = {17 (natural armor)},
hit-points = {\DndDice{18d10 + 36}},
speed = {10 ft., swim 40 ft.}
]
\DndMonsterAbilityScores[
 str = 21,
 dex = 9,
 con = 15,
 int = 18,
 wis = 15,
 cha = 18
]
\DndMonsterDetails[
saving-throws = {Con +6, Int +8, Wis +6},
skills = {History +12, Perception +10},
senses = {darkvision 120 ft., passive perception 20},
languages = {Deep Speech, telepathy 120 ft.},
challenge = 10,
]
\DndMonsterAction{Amphibious}
The aboleth can breathe air and water.

\DndMonsterAction{Mucous Cloud}
While underwater, the aboleth is surrounded by transformative mucus. A creature that touches the aboleth or that hits it with a melee attack while within 5 feet of it must make a DC 14 Constitution saving throw. On a failure, the creature is diseased for 1d4 hours. The diseased creature can breathe only underwater.

\DndMonsterAction{Probing Telepathy}
If a creature communicates telepathically with the aboleth, the aboleth learns the creature's greatest desires if the aboleth can see the creature.

\DndMonsterSection{Actions}
\DndMonsterAction{Multiattack}
The aboleth makes three tentacle attacks.

\DndMonsterAction{Tentacle}
\textsl{Melee Weapon Attack:} 9 to hit, reach 10 ft., one target. \textsl{Hit:} 12 (2d6 + 5) bludgeoning damage. If the target is a creature, it must succeed on a DC 14 Constitution saving throw or become diseased. The disease has no effect for 1 minute and can be removed by any magic that cures disease. After 1 minute, the diseased creature's skin becomes translucent and slimy, the creature can't regain hit points unless it is underwater, and the disease can be removed only by \textit{heal} or another disease-curing spell of 6th level or higher. When the creature is outside a body of water, it takes 6 (1d12) acid damage every 10 minutes unless moisture is applied to the skin before 10 minutes have passed.

\DndMonsterAction{Tail}
\textsl{Melee Weapon Attack:} 9 to hit, reach 10 ft., one target. \textsl{Hit:} 15 (3d6 + 5) bludgeoning damage.

\DndMonsterAction{Enslave (3/Day)}
The aboleth targets one creature it can see within 30 feet of it. The target must succeed on a DC 14 Wisdom saving throw or be magically charmed by the aboleth until the aboleth dies or until it is on a different plane of existence from the target. The charmed target is under the aboleth's control and can't take reactions, and the aboleth and the target can communicate telepathically with each other over any distance.

Whenever the charmed target takes damage, the target can repeat the saving throw. On a success, the effect ends. No more than once every 24 hours, the target can also repeat the saving throw when it is at least 1 mile away from the aboleth.

\DndMonsterSection{Legendary Actions}
Aboleth can take 3 legendary actions, choosing from the options below. Only one legendary action can be used at a time and only at the end of another creature's turn. Aboleth regains spent legendary actions at the start of its turn.

\begin{DndMonsterLegendaryActions}
\DndMonsterLegendaryAction{Detect}{The aboleth makes a Wisdom (Perception) check.
}
\DndMonsterLegendaryAction{Tail Swipe}{The aboleth makes one tail attack.
}
\DndMonsterLegendaryAction{Psychic Drain (Costs 2 Actions)}{One creature charmed by the aboleth takes 10 (3d6) psychic damage, and the aboleth regains hit points equal to the damage the creature takes.
}
\end{DndMonsterLegendaryActions}
\DndMonsterSection{Lair Actions}
When fighting inside its lair, an aboleth can invoke the ambient magic to take lair actions. On initiative count 20 (losing initiative ties), the aboleth takes a lair action to cause one of the following effects:


\begin{itemize}
\item The aboleth casts \textit{phantasmal force} (no components required) on any number of creatures it can see within 60 feet of it. While maintaining concentration on this effect, the aboleth can't take other lair actions. If a target succeeds on the saving throw or if the effect ends for it, the target is immune to the aboleth's phantasmal force lair action for the next 24 hours, although such a creature can choose to be affected.
\item Pools of water within 90 feet of the aboleth surge outward in a grasping tide. Any creature on the ground within 20 feet of such a pool must succeed on a DC 14 Strength saving throw or be pulled up to 20 feet into the water and knocked prone. The aboleth can't use this lair action again until it has used a different one.
\item Water in the aboleth's lair magically becomes a conduit for the creature's rage. The aboleth can target any number of creatures it can see in such water within 90 feet of it. A target must succeed on a DC 14 Wisdom saving throw or take 7 (2d6) psychic damage. The aboleth can't use this lair action again until it has used a different one.
\end{itemize}

\DndMonsterSection{Regional Effects}
The region containing an aboleth's lair is warped by the creature's presence, which creates one or more of the following effects:


\begin{itemize}
\item Underground surfaces within 1 mile of the aboleth's lair are slimy and wet and are difficult terrain.
\item Water sources within 1 mile of the lair are supernaturally fouled. Enemies of the aboleth that drink such water vomit it within minutes.
\item As an action, the aboleth can create an illusory image of itself within 1 mile of the lair. The copy can appear at any location the aboleth has seen before or in any location a creature charmed by the aboleth can currently see. Once created, the image lasts for as long as the aboleth maintains concentration, as if concentration on a spell. Although the image is intangible, it looks, sounds, and can move like the aboleth. The aboleth can sense, speak, and use telepathy from the image's position as if present at that position. If the image takes any damage, it disappears.
\end{itemize}

If the aboleth dies, the first two effects fade over the course of 3d10 days.

\end{multicols}
\end{DndMonster}



\begin{DndMonster}[float=t]{Acolyte}
\DndMonsterType{Medium humanoid (any race), any alignment}
\DndMonsterBasics[
armor-class = {10},
hit-points = {\DndDice{2d8}},
speed = {30 ft.}
]
\DndMonsterAbilityScores[
 str = 10,
 dex = 10,
 con = 10,
 int = 10,
 wis = 14,
 cha = 11
]
\DndMonsterDetails[
skills = {Medicine +4, Religion +2},
languages = {any one language (usually Common)},
challenge = 1/4,
]
\DndMonsterAction{Spellcasting}
The acolyte is a 1st-level spellcaster. Its spellcasting ability is Wisdom (spell save DC 12, 4 to hit with spell attacks). The acolyte has following cleric spells prepared:
\begin{DndMonsterSpells}
  \DndMonsterSpellLevel{\textit{light}, \textit{sacred flame}, \textit{thaumaturgy}}
   \DndMonsterSpellLevel[1][3]{\textit{bless}, \textit{cure wounds}, \textit{sanctuary}}
\end{DndMonsterSpells}
\DndMonsterSection{Actions}
\DndMonsterAction{Club}
\textsl{Melee Weapon Attack:} 2 to hit, reach 5 ft., one target. \textsl{Hit:} 2 (1d4) bludgeoning damage.

\end{DndMonster}



\begin{DndMonster}[float=t]{Air Elemental}
\DndMonsterType{Large elemental, neutral}
\DndMonsterBasics[
armor-class = {15},
hit-points = {\DndDice{12d10 + 24}},
speed = {0 ft., fly 90 ft. \textit{(hover)}}
]
\DndMonsterAbilityScores[
 str = 14,
 dex = 20,
 con = 14,
 int = 6,
 wis = 10,
 cha = 6
]
\DndMonsterDetails[
damage-resistances = {lightning; thunder; bludgeoning, piercing, slashing from nonmagical attacks},
damage-immunities = {poison},
condition-immunities = {exhaustion, grappled, paralyzed, petrified, poisoned, prone, restrained, unconscious},
senses = {darkvision 60 ft., passive perception 10},
languages = {Auran},
challenge = 5,
]
\DndMonsterAction{Air Form}
The elemental can enter a hostile creature's space and stop there. It can move through a space as narrow as 1 inch wide without squeezing.

\DndMonsterSection{Actions}
\DndMonsterAction{Multiattack}
The elemental makes two slam attacks.

\DndMonsterAction{Slam}
\textsl{Melee Weapon Attack:} 8 to hit, reach 5 ft., one target. \textsl{Hit:} 14 (2d8 + 5) bludgeoning damage.

\DndMonsterAction{Whirlwind (Recharge 4\textendash 6)}
Each creature in the elemental's space must make a DC 13 Strength saving throw. On a failure, a target takes 15 (3d8 + 2) bludgeoning damage and is flung up 20 feet away from the elemental in a random direction and knocked prone. If a thrown target strikes an object, such as a wall or floor, the target takes 3 (1d6) bludgeoning damage for every 10 feet it was thrown. If the target is thrown at another creature, that creature must succeed on a DC 13 Dexterity saving throw or take the same damage and be knocked prone.

If the saving throw is successful, the target takes half the bludgeoning damage and isn't flung away or knocked prone.

\end{DndMonster}



\begin{DndMonster}[float*=b,width=\textwidth + 8pt]{Amalgamation}
\begin{multicols}{2}
\DndMonsterType{Gargantuan aberration, chaotic evil}
\DndMonsterBasics[
armor-class = {18 (natural armor)},
hit-points = {\DndDice{30d20 + 210}},
speed = {5 ft.}
]
\DndMonsterAbilityScores[
 str = 30,
 dex = 11,
 con = 27,
 int = 22,
 wis = 18,
 cha = 20
]
\DndMonsterDetails[
saving-throws = {Str +17, Dex +7, Con +15, Int +13, Wis +11},
skills = {Arcana +13, Deception +12, Insight +11, Perception +11},
damage-immunities = {bludgeoning, piercing, slashing from nonmagical attacks},
condition-immunities = {charmed, exhaustion, frightened, paralyzed, petrified, poisoned, prone, stunned},
senses = {truesight 120 ft., passive perception 21},
languages = {all, telepathy 1 mile},
challenge = 23,
]
\DndMonsterAction{Legendary Resistance (3/Day)}
If the amalgamation fails a saving throw, it can choose to succeed instead.

\DndMonsterAction{Malignancy}
The amalgamation can't be teleported or travel to another plane against its will.

\DndMonsterAction{Immutable Form}
The amalgamation is immune to any spell or effect that would alter its form.

\DndMonsterAction{Magic Weapons}
The amalgamation's weapon attacks are magical.

\DndMonsterAction{Innate Spellcasting}
The amalgamation's innate spellcasting ability is Charisma (spell save DC 20). It can innately cast the following spells, requiring no components:
\begin{DndMonsterSpells}
  \DndInnateSpellLevel{\textit{detect thoughts}, \textit{dispel magic}}
\end{DndMonsterSpells}
\DndMonsterSection{Actions}
\DndMonsterAction{Multiattack}
The amalgamation uses its Contaminating Presence. It then makes three tentacle attacks, each of which it can replace with one use of Fling.

\DndMonsterAction{Bite}
\textsl{Melee Weapon Attack:} 17 to hit, reach 5 ft., one target. \textsl{Hit:} 23 (3d8 + 10) piercing damage. If the target is a Large or smaller creature grappled by the amalgamation, that creature is swallowed, and the grapple ends.

While swallowed, the creature is blinded and restrained, it has total cover against attacks and other effects outside the amalgamation, and it takes 42 (12d6) acid damage and gains one level of contamination at the start of each of the amalgamation's turns.

If the amalgamation takes 50 damage or more on a single turn from a creature inside it, the amalgamation must succeed on a DC 25 Constitution saving throw at the end of that turn or regurgitate all swallowed creatures, which fall prone in a space within 10 feet of the amalgamation. If the amalgamation dies, a swallowed creature is no longer restrained by it and can escape from the corpse using 15 feet of movement, exiting prone.

\DndMonsterAction{Tentacle}
\textsl{Melee Weapon Attack:} 17 to hit, reach 50 ft., one target. \textsl{Hit:} 20 (3d6 + 10) bludgeoning damage, and the target is grappled (escape DC 18). Until this grapple ends, the target is restrained. The amalgamation has twelve tentacles, each of which can grapple one target.

\DndMonsterAction{Fling}
One Large or smaller object held or creature grappled by the amalgamation is thrown up to 60 feet in a direction of the amalgamation's choice and knocked prone. If a thrown target strikes a solid surface, the target takes 3 (1d6) bludgeoning damage for every 10 feet it was thrown. If the target is thrown at another creature, that creature must succeed on a DC 18 Dexterity saving throw or take the same damage and be knocked prone.

\DndMonsterAction{Mind-breaking Whispers}
The amalgamation utters a terrible truth within the mind of one creature within 120 feet of it. That creature must succeed on a DC 20 Intelligence saving throw. On a failure, it takes 35 (10d6) psychic damage and becomes insane until it finishes a long rest. While insane, it can't take actions, can't understand what other creatures say, can't read, and can speak only in gibberish. A \textit{greater restoration} spell cast on the creature ends this effect.

\DndMonsterAction{Contaminating Presence}
Humanoid creatures within 120 feet of the amalgamation must succeed on a DC 20 Constitution saving throw or take 10 (3d6) necrotic damage and gain one level of contamination.

\DndMonsterSection{Legendary Actions}
Amalgamation can take 3 legendary actions, choosing from the options below. Only one legendary action can be used at a time and only at the end of another creature's turn. Amalgamation regains spent legendary actions at the start of its turn.

\begin{DndMonsterLegendaryActions}
\DndMonsterLegendaryAction{Tentacle Attack or Fling}{The amalgamation makes one tentacle attack or uses its Fling.
}
\DndMonsterLegendaryAction{Mind-breaking Whispers (Costs 2 Actions)}{The amalgamation uses Mind-breaking Whispers.
}
\DndMonsterLegendaryAction{Piper at the Gate (Costs 3 Actions)}{Up to five gibbering mouthers appear in unoccupied spaces within 30 feet of the amalgamation and remain until destroyed. These creatures act after the amalgamation on the same initiative count and are under its control. The amalgamation can have up to ten gibbering mouthers summoned by this ability at a time.
}
\end{DndMonsterLegendaryActions}
\end{multicols}
\end{DndMonster}



\begin{DndMonster}[float*=b,width=\textwidth + 8pt]{Ancient Bronze Dragon}
\begin{multicols}{2}
\DndMonsterType{Gargantuan dragon, lawful good}
\DndMonsterBasics[
armor-class = {22 (natural armor)},
hit-points = {\DndDice{24d20 + 192}},
speed = {40 ft., fly 80 ft., swim 40 ft.}
]
\DndMonsterAbilityScores[
 str = 29,
 dex = 10,
 con = 27,
 int = 18,
 wis = 17,
 cha = 21
]
\DndMonsterDetails[
saving-throws = {Dex +7, Con +15, Wis +10, Cha +12},
skills = {Insight +10, Perception +17, Stealth +7},
damage-immunities = {lightning},
senses = {blindsight 60 ft., darkvision 120 ft., passive perception 27},
languages = {Common, Draconic},
challenge = 22,
]
\DndMonsterAction{Amphibious}
The dragon can breathe air and water.

\DndMonsterAction{Legendary Resistance (3/Day)}
If the dragon fails a saving throw, it can choose to succeed instead.

\DndMonsterSection{Actions}
\DndMonsterAction{Multiattack}
The dragon can use its Frightful Presence. It then makes three attacks: one with its bite and two with its claws.

\DndMonsterAction{Bite}
\textsl{Melee Weapon Attack:} 16 to hit, reach 15 ft., one target. \textsl{Hit:} 20 (2d10 + 9) piercing damage.

\DndMonsterAction{Claw}
\textsl{Melee Weapon Attack:} 16 to hit, reach 10 ft., one target. \textsl{Hit:} 16 (2d6 + 9) slashing damage.

\DndMonsterAction{Tail}
\textsl{Melee Weapon Attack:} 16 to hit, reach 20 ft., one target. \textsl{Hit:} 18 (2d8 + 9) bludgeoning damage.

\DndMonsterAction{Frightful Presence}
Each creature of the dragon's choice that is within 120 feet of the dragon and aware of it must succeed on a DC 20 Wisdom saving throw or become frightened for 1 minute. A creature can repeat the saving throw at the end of each of its turns, ending the effect on itself on a success. If a creature's saving throw is successful or the effect ends for it, the creature is immune to the dragon's Frightful Presence for the next 24 hours.

\DndMonsterAction{Breath Weapons (Recharge 5\textendash 6)}
The dragon uses one of the following breath weapons.


\begin{description}
\item[Lightning Breath.]
The dragon exhales lightning in a 120-foot line that is 10 feet wide. Each creature in that line must make a DC 23 Dexterity saving throw, taking 88 (16d10) lightning damage on a failed save, or half as much damage on a successful one.
\item[Repulsion Breath.]
The dragon exhales repulsion energy in a 30-foot cone. Each creature in that area must succeed on a DC 23 Strength saving throw. On a failed save, the creature is pushed 60 feet away from the dragon.
\end{description}

\DndMonsterAction{Change Shape}
The dragon magically polymorphs into a humanoid or beast that has a challenge rating no higher than its own, or back into its true form. It reverts to its true form if it dies. Any equipment it is wearing or carrying is absorbed or borne by the new form (the dragon's choice).

In a new form, the dragon retains its alignment, hit points, Hit Dice, ability to speak, proficiencies, Legendary Resistance, lair actions, and Intelligence, Wisdom, and Charisma scores, as well as this action. Its statistics and capabilities are otherwise replaced by those of the new form, except any class features or legendary actions of that form.

\DndMonsterSection{Legendary Actions}
Ancient Bronze Dragon can take 3 legendary actions, choosing from the options below. Only one legendary action can be used at a time and only at the end of another creature's turn. Ancient Bronze Dragon regains spent legendary actions at the start of its turn.

\begin{DndMonsterLegendaryActions}
\DndMonsterLegendaryAction{Detect}{The dragon makes a Wisdom (Perception) check.
}
\DndMonsterLegendaryAction{Tail Attack}{The dragon makes a tail attack.
}
\DndMonsterLegendaryAction{Wing Attack (Costs 2 Actions)}{The dragon beats its wings. Each creature within 15 feet of the dragon must succeed on a DC 24 Dexterity saving throw or take 16 (2d6 + 9) bludgeoning damage and be knocked prone. The dragon can then fly up to half its flying speed.
}
\end{DndMonsterLegendaryActions}
\end{multicols}
\end{DndMonster}



\begin{DndMonster}[float*=b,width=\textwidth + 8pt]{Ancient Gold Dragon}
\begin{multicols}{2}
\DndMonsterType{Gargantuan dragon, lawful good}
\DndMonsterBasics[
armor-class = {22 (natural armor)},
hit-points = {\DndDice{28d20 + 252}},
speed = {40 ft., fly 80 ft., swim 40 ft.}
]
\DndMonsterAbilityScores[
 str = 30,
 dex = 14,
 con = 29,
 int = 18,
 wis = 17,
 cha = 28
]
\DndMonsterDetails[
saving-throws = {Dex +9, Con +16, Wis +10, Cha +16},
skills = {Insight +10, Perception +17, Persuasion +16, Stealth +9},
damage-immunities = {fire},
senses = {blindsight 60 ft., darkvision 120 ft., passive perception 27},
languages = {Common, Draconic},
challenge = 24,
]
\DndMonsterAction{Amphibious}
The dragon can breathe air and water.

\DndMonsterAction{Legendary Resistance (3/Day)}
If the dragon fails a saving throw, it can choose to succeed instead.

\DndMonsterSection{Actions}
\DndMonsterAction{Multiattack}
The dragon can use its Frightful Presence. It then makes three attacks: one with its bite and two with its claws.

\DndMonsterAction{Bite}
\textsl{Melee Weapon Attack:} 17 to hit, reach 15 ft., one target. \textsl{Hit:} 21 (2d10 + 10) piercing damage.

\DndMonsterAction{Claw}
\textsl{Melee Weapon Attack:} 17 to hit, reach 10 ft., one target. \textsl{Hit:} 17 (2d6 + 10) slashing damage.

\DndMonsterAction{Tail}
\textsl{Melee Weapon Attack:} 17 to hit, reach 20 ft., one target. \textsl{Hit:} 19 (2d8 + 10) bludgeoning damage.

\DndMonsterAction{Frightful Presence}
Each creature of the dragon's choice that is within 120 feet of the dragon and aware of it must succeed on a DC 24 Wisdom saving throw or become frightened for 1 minute. A creature can repeat the saving throw at the end of each of its turns, ending the effect on itself on a success. If a creature's saving throw is successful or the effect ends for it, the creature is immune to the dragon's Frightful Presence for the next 24 hours.

\DndMonsterAction{Breath Weapons (Recharge 5\textendash 6)}
The dragon uses one of the following breath weapons.


\begin{description}
\item[Fire Breath.]
The dragon exhales fire in a 90-foot cone. Each creature in that area must make a DC 24 Dexterity saving throw, taking 71 (13d10) fire damage on a failed save, or half as much damage on a successful one.
\item[Weakening Breath.]
The dragon exhales gas in a 90-foot cone. Each creature in that area must succeed on a DC 24 Strength saving throw or have disadvantage on Strength-based attack rolls, Strength checks, and Strength saving throws for 1 minute. A creature can repeat the saving throw at the end of each of its turns, ending the effect on itself on a success.
\end{description}

\DndMonsterAction{Change Shape}
The dragon magically polymorphs into a humanoid or beast that has a challenge rating no higher than its own, or back into its true form. It reverts to its true form if it dies. Any equipment it is wearing or carrying is absorbed or borne by the new form (the dragon's choice).

In a new form, the dragon retains its alignment, hit points, Hit Dice, ability to speak, proficiencies, Legendary Resistance, lair actions, and Intelligence, Wisdom, and Charisma scores, as well as this action. Its statistics and capabilities are otherwise replaced by those of the new form, except any class features or legendary actions of that form.

\DndMonsterSection{Legendary Actions}
Ancient Gold Dragon can take 3 legendary actions, choosing from the options below. Only one legendary action can be used at a time and only at the end of another creature's turn. Ancient Gold Dragon regains spent legendary actions at the start of its turn.

\begin{DndMonsterLegendaryActions}
\DndMonsterLegendaryAction{Detect}{The dragon makes a Wisdom (Perception) check.
}
\DndMonsterLegendaryAction{Tail Attack}{The dragon makes a tail attack.
}
\DndMonsterLegendaryAction{Wing Attack (Costs 2 Actions)}{The dragon beats its wings. Each creature within 15 feet of the dragon must succeed on a DC 25 Dexterity saving throw or take 17 (2d6 + 10) bludgeoning damage and be knocked prone. The dragon can then fly up to half its flying speed.
}
\end{DndMonsterLegendaryActions}
\end{multicols}
\end{DndMonster}



\begin{DndMonster}[float=t]{Animated Armor}
\DndMonsterType{Medium construct, unaligned}
\DndMonsterBasics[
armor-class = {18 (natural armor)},
hit-points = {\DndDice{6d8 + 6}},
speed = {25 ft.}
]
\DndMonsterAbilityScores[
 str = 14,
 dex = 11,
 con = 13,
 int = 1,
 wis = 3,
 cha = 1
]
\DndMonsterDetails[
damage-immunities = {poison, psychic},
condition-immunities = {blinded, charmed, deafened, exhaustion, frightened, paralyzed, petrified, poisoned},
senses = {blindsight 60 ft. (blind beyond this radius), passive perception 6},
challenge = 1,
]
\DndMonsterAction{Antimagic Susceptibility}
The armor is incapacitated while in the area of an \textit{antimagic field}. If targeted by \textit{dispel magic}, the armor must succeed on a Constitution saving throw against the caster's spell save DC or fall unconscious for 1 minute.

\DndMonsterAction{False Appearance}
While the armor remains motionless, it is indistinguishable from a normal suit of armor.

\DndMonsterSection{Actions}
\DndMonsterAction{Multiattack}
The armor makes two melee attacks.

\DndMonsterAction{Slam}
\textsl{Melee Weapon Attack:} 4 to hit, reach 5 ft., one target. \textsl{Hit:} 5 (1d6 + 2) bludgeoning damage.

\end{DndMonster}



\begin{DndMonster}[float=t]{Animated Delerium Sludge}
\DndMonsterType{Large elemental, unaligned}
\DndMonsterBasics[
armor-class = {14 (natural armor)},
hit-points = {\DndDice{12d10 + 48}},
speed = {30 ft., swim 90 ft.}
]
\DndMonsterAbilityScores[
 str = 18,
 dex = 14,
 con = 18,
 int = 5,
 wis = 10,
 cha = 8
]
\DndMonsterDetails[
damage-resistances = {acid; bludgeoning, piercing, slashing from nonmagical attacks},
damage-immunities = {necrotic, poison},
condition-immunities = {exhaustion, grappled, paralyzed, petrified, poisoned, prone, restrained, unconscious},
senses = {darkvision 60 ft., passive perception 10},
languages = {Deep Speech},
challenge = 5,
]
\DndMonsterAction{Sludge Form}
The animated sludge can enter a hostile creature's space and stop there. It can move through a space as narrow as 1 inch wide without squeezing.

\DndMonsterSection{Actions}
\DndMonsterAction{Contaminated Touch}
\textsl{Melee Weapon Attack:} 7 to hit, reach 5 ft., one target. \textsl{Hit:} 10 (3d6) necrotic damage. In addition, a target hit by this attack must succeed on a DC 15 Constitution saving throw or gain one level of contamination.

\DndMonsterAction{Engulf}
Each creature in the animated sludge's space must succeed on a DC 15 Strength saving throw. On a failed save, a creature becomes grappled by the sludge if it is Large or smaller (escape DC 14). Until this grapple ends, the target is restrained. If the saving throw is successful, the target is pushed out of the sludge's space.

The sludge can grapple one Large creature or up to two Medium or smaller creatures at one time. At the start of each of the sludge's turns, each target grappled by it takes 28 (8d6) necrotic damage and must succeed on a DC 15 Constitution saving throw or gain one level of contamination.

A creature within 5 feet of the sludge can attempt to pull a creature or object out of it. This takes an action and requires a successful DC 14 Strength check.

\end{DndMonster}



\begin{DndMonster}[float=t]{Arcane Wraith}
\DndMonsterType{Medium undead, neutral evil}
\DndMonsterBasics[
armor-class = {13},
hit-points = {\DndDice{9d8 + 27}},
speed = {0 ft., fly 60 ft. \textit{(hover)}}
]
\DndMonsterAbilityScores[
 str = 6,
 dex = 16,
 con = 16,
 int = 12,
 wis = 14,
 cha = 15
]
\DndMonsterDetails[
damage-resistances = {acid; cold; fire; lightning; thunder; bludgeoning, piercing, slashing from nonmagical attacks that aren't silvered},
damage-immunities = {necrotic, poison},
condition-immunities = {charmed, exhaustion, grappled, paralyzed, petrified, poisoned, prone, restrained},
senses = {darkvision 60 ft., passive perception 12},
languages = {the languages it knew in life},
challenge = 5,
]
\DndMonsterAction{Copy Warning}
This content has been copied from another creature, but the data replacement hasn't been run. Check details.

\DndMonsterAction{Incorporeal Movement}
The wraith can move through other creatures and objects as if they were difficult terrain. It takes 5 (1d10) force damage if it ends its turn inside an object.

\DndMonsterAction{Sunlight Sensitivity}
While in sunlight, the wraith has disadvantage on attack rolls, as well as on Wisdom (Perception) checks that rely on sight.

\DndMonsterSection{Actions}
\DndMonsterAction{Life Drain}
\textsl{Melee Weapon Attack:} 6 to hit, reach 5 ft., one creature. \textsl{Hit:} 21 (4d8 + 3) necrotic damage. The target must succeed on a DC 14 Constitution saving throw or its hit point maximum is reduced by an amount equal to the damage taken. This reduction lasts until the target finishes a long rest. The target dies if this effect reduces its hit point maximum to 0.

\DndMonsterAction{Create Specter}
The wraith targets a humanoid within 10 feet of it that has been dead for no longer than 1 minute and died violently. The target's spirit rises as a \textbf{specter} in the space of its corpse or in the nearest unoccupied space. The \textbf{specter} is under the wraith's control. The wraith can have no more than seven specters under its control at one time.

\end{DndMonster}



\begin{DndMonster}[float=t]{Archmage}
\DndMonsterType{Medium humanoid (any race), any alignment}
\DndMonsterBasics[
armor-class = {12 (15 with \textit{mage armor})},
hit-points = {\DndDice{18d8 + 18}},
speed = {30 ft.}
]
\DndMonsterAbilityScores[
 str = 10,
 dex = 14,
 con = 12,
 int = 20,
 wis = 15,
 cha = 16
]
\DndMonsterDetails[
saving-throws = {Int +9, Wis +6},
skills = {Arcana +13, History +13},
damage-resistances = {damage from spells; bludgeoning, piercing, slashing (from stoneskin)},
languages = {any six languages},
challenge = 12,
]
\DndMonsterAction{Magic Resistance}
The archmage has advantage on saving throws against spells and other magical effects.

\DndMonsterAction{Spellcasting}
The archmage is an 18th-level spellcaster. Its spellcasting ability is Intelligence (spell save DC 17, 9 to hit with spell attacks). The archmage can cast \textit{disguise self} and \textit{invisibility} at will and has the following wizard spells prepared:
\begin{DndMonsterSpells}
  \DndInnateSpellLevel{\textit{disguise self}, \textit{invisibility}}
  \DndMonsterSpellLevel{\textit{fire bolt}, \textit{light}, \textit{mage hand}, \textit{prestidigitation}, \textit{shocking grasp}}
   \DndMonsterSpellLevel[1][4]{\textit{detect magic}, \textit{identify}, \textit{mage armor}*, \textit{magic missile}}
   \DndMonsterSpellLevel[2][3]{\textit{detect thoughts}, \textit{mirror image}, \textit{misty step}}
   \DndMonsterSpellLevel[3][3]{\textit{counterspell}, \textit{fly}, \textit{lightning bolt}}
   \DndMonsterSpellLevel[4][3]{\textit{banishment}, \textit{fire shield}, \textit{stoneskin}*}
   \DndMonsterSpellLevel[5][3]{\textit{cone of cold}, \textit{scrying}, \textit{wall of force}}
   \DndMonsterSpellLevel[6][1]{\textit{globe of invulnerability}}
   \DndMonsterSpellLevel[7][1]{\textit{teleport}}
   \DndMonsterSpellLevel[8][1]{\textit{mind blank}*}
\end{DndMonsterSpells}
\DndMonsterSection{Actions}
\DndMonsterAction{Dagger}
\textsl{Melee or Ranged Weapon Attack:} 6 to hit, reach 5 ft. or range 20/60 ft., one target. \textsl{Hit:} 4 (1d4 + 2) piercing damage.

\end{DndMonster}



\begin{DndMonster}[float=t]{Assassin}
\DndMonsterType{Medium humanoid (any race), any non-good alignment}
\DndMonsterBasics[
armor-class = {15 (studded leather armor)},
hit-points = {\DndDice{12d8 + 24}},
speed = {30 ft.}
]
\DndMonsterAbilityScores[
 str = 11,
 dex = 16,
 con = 14,
 int = 13,
 wis = 11,
 cha = 10
]
\DndMonsterDetails[
saving-throws = {Dex +6, Int +4},
skills = {Acrobatics +6, Deception +3, Perception +3, Stealth +9},
damage-resistances = {poison},
languages = {Thieves' cant plus any two languages},
challenge = 8,
]
\DndMonsterAction{Assassinate}
During its first turn, the assassin has advantage on attack rolls against any creature that hasn't taken a turn. Any hit the assassin scores against a Surprise creature is a critical hit.

\DndMonsterAction{Evasion}
If the assassin is subjected to an effect that allows it to make a Dexterity saving throw to take only half damage, the assassin instead takes no damage if it succeeds on the saving throw, and only half damage if it fails.

\DndMonsterAction{Sneak Attack (1/Turn)}
The assassin deals an extra 14 (4d6) damage when it hits a target with a weapon attack and has advantage on the attack roll, or when the target is within 5 feet of an ally of the assassin that isn't incapacitated and the assassin doesn't have disadvantage on the attack roll.

\DndMonsterSection{Actions}
\DndMonsterAction{Multiattack}
The assassin makes two shortsword attacks.

\DndMonsterAction{Shortsword}
\textsl{Melee Weapon Attack:} 6 to hit, reach 5 ft., one target. \textsl{Hit:} 6 (1d6 + 3) piercing damage, and the target must make a DC 15 Constitution saving throw, taking 24 (7d6) poison damage on a failed save, or half as much damage on a successful one.

\DndMonsterAction{Light Crossbow}
\textsl{Ranged Weapon Attack:} 6 to hit, range 80/320 ft., one target. \textsl{Hit:} 7 (1d8 + 3) piercing damage, and the target must make a DC 15 Constitution saving throw, taking 24 (7d6) poison damage on a failed save, or half as much damage on a successful one.

\end{DndMonster}



\begin{DndMonster}[float=t]{Awakened Shrub}
\DndMonsterType{Small plant, unaligned}
\DndMonsterBasics[
armor-class = {9},
hit-points = {\DndDice{3d6}},
speed = {20 ft.}
]
\DndMonsterAbilityScores[
 str = 3,
 dex = 8,
 con = 11,
 int = 10,
 wis = 10,
 cha = 6
]
\DndMonsterDetails[
damage-vulnerabilities = {fire},
damage-resistances = {piercing},
languages = {one language known by its creator},
challenge = 0,
]
\DndMonsterAction{False Appearance}
While the shrub remains motionless, it is indistinguishable from a normal shrub.

\DndMonsterSection{Actions}
\DndMonsterAction{Rake}
\textsl{Melee Weapon Attack:} 1 to hit, reach 5 ft., one target. \textsl{Hit:} 1 (1d4 - 1) slashing damage.

\end{DndMonster}



\begin{DndMonster}[float=t]{Azer}
\DndMonsterType{Medium elemental, lawful neutral}
\DndMonsterBasics[
armor-class = {17 (natural armor)},
hit-points = {\DndDice{6d8 + 12}},
speed = {30 ft.}
]
\DndMonsterAbilityScores[
 str = 17,
 dex = 12,
 con = 15,
 int = 12,
 wis = 13,
 cha = 10
]
\DndMonsterDetails[
saving-throws = {Con +4},
damage-immunities = {fire, poison},
condition-immunities = {poisoned},
languages = {Ignan},
challenge = 2,
]
\DndMonsterAction{Heated Body}
A creature that touches the azer or hits it with a melee attack while within 5 feet of it takes 5 (1d10) fire damage.

\DndMonsterAction{Heated Weapons}
When the azer hits with a metal melee weapon, it deals an extra 3 (1d6) fire damage (included in the attack).

\DndMonsterAction{Illumination}
The azer sheds bright light in a 10-foot radius and dim light for an additional 10 feet.

\DndMonsterSection{Actions}
\DndMonsterAction{Warhammer}
\textsl{Melee Weapon Attack:} 5 to hit, reach 5 ft., one target. \textsl{Hit:} 7 (1d8 + 3) bludgeoning damage, or 8 (1d10 + 3) bludgeoning damage if used with two hands to make a melee attack, plus 3 (1d6) fire damage.

\end{DndMonster}



\begin{DndMonster}[float=t]{Balor}
\DndMonsterType{Huge fiend (demon), chaotic evil}
\DndMonsterBasics[
armor-class = {19 (natural armor)},
hit-points = {\DndDice{21d12 + 126}},
speed = {40 ft., fly 80 ft.}
]
\DndMonsterAbilityScores[
 str = 26,
 dex = 15,
 con = 22,
 int = 20,
 wis = 16,
 cha = 22
]
\DndMonsterDetails[
saving-throws = {Str +14, Con +12, Wis +9, Cha +12},
damage-resistances = {cold; lightning; bludgeoning, piercing, slashing from nonmagical attacks},
damage-immunities = {fire, poison},
condition-immunities = {poisoned},
senses = {truesight 120 ft., passive perception 13},
languages = {Abyssal, telepathy 120 ft.},
challenge = 19,
]
\DndMonsterAction{Death Throes}
When the balor dies, it explodes, and each creature within 30 feet of it must make a DC 20 Dexterity saving throw, taking 70 (20d6) fire damage on a failed save, or half as much damage on a successful one. The explosion ignites flammable objects in that area that aren't being worn or carried, and it destroys the balor's weapons.

\DndMonsterAction{Fire Aura}
At the start of each of the balor's turns, each creature within 5 feet of it takes 10 (3d6) fire damage, and flammable objects in the aura that aren't being worn or carried ignite. A creature that touches the balor or hits it with a melee attack while within 5 feet of it takes 10 (3d6) fire damage.

\DndMonsterAction{Magic Resistance}
The balor has advantage on saving throws against spells and other magical effects.

\DndMonsterAction{Magic Weapons}
The balor's weapon attacks are magical.

\DndMonsterSection{Actions}
\DndMonsterAction{Multiattack}
The balor makes two attacks: one with its longsword and one with its whip.

\DndMonsterAction{Longsword}
\textsl{Melee Weapon Attack:} 14 to hit, reach 10 ft., one target. \textsl{Hit:} 21 (3d8 + 8) slashing damage plus 13 (3d8) lightning damage. If the balor scores a critical hit, it rolls damage dice three times, instead of twice.

\DndMonsterAction{Whip}
\textsl{Melee Weapon Attack:} 14 to hit, reach 30 ft., one target. \textsl{Hit:} 15 (2d6 + 8) slashing damage plus 10 (3d6) fire damage, and the target must succeed on a DC 20 Strength saving throw or be pulled up to 25 feet toward the balor.

\DndMonsterAction{Teleport}
The balor magically teleports, along with any equipment it is wearing or carrying, up to 120 feet to an unoccupied space it can see.

\end{DndMonster}



\begin{DndMonster}[float=t]{Bandit}
\DndMonsterType{Medium humanoid (any race), any non-lawful alignment}
\DndMonsterBasics[
armor-class = {12 (\textit{leather armor})},
hit-points = {\DndDice{2d8 + 2}},
speed = {30 ft.}
]
\DndMonsterAbilityScores[
 str = 11,
 dex = 12,
 con = 12,
 int = 10,
 wis = 10,
 cha = 10
]
\DndMonsterDetails[
languages = {any one language (usually Common)},
challenge = 1/8,
]
\DndMonsterSection{Actions}
\DndMonsterAction{Scimitar}
\textsl{Melee Weapon Attack:} 3 to hit, reach 5 ft., one target. \textsl{Hit:} 4 (1d6 + 1) slashing damage.

\DndMonsterAction{Light Crossbow}
\textsl{Ranged Weapon Attack:} 3 to hit, range 80/320 ft., one target. \textsl{Hit:} 5 (1d8 + 1) piercing damage.

\end{DndMonster}



\begin{DndMonster}[float=t]{Bandit Captain}
\DndMonsterType{Medium humanoid (any race), any non-lawful alignment}
\DndMonsterBasics[
armor-class = {15 (studded leather armor)},
hit-points = {\DndDice{10d8 + 20}},
speed = {30 ft.}
]
\DndMonsterAbilityScores[
 str = 15,
 dex = 16,
 con = 14,
 int = 14,
 wis = 11,
 cha = 14
]
\DndMonsterDetails[
saving-throws = {Str +4, Dex +5, Wis +2},
skills = {Athletics +4, Deception +4},
languages = {any two languages},
challenge = 2,
]
\DndMonsterSection{Actions}
\DndMonsterAction{Multiattack}
The captain makes three melee attacks: two with its scimitar and one with its dagger. Or the captain makes two ranged attacks with its daggers.

\DndMonsterAction{Scimitar}
\textsl{Melee Weapon Attack:} 5 to hit, reach 5 ft., one target. \textsl{Hit:} 6 (1d6 + 3) slashing damage.

\DndMonsterAction{Dagger}
\textsl{Melee or Ranged Weapon Attack:} 5 to hit, reach 5 ft. or range 20/60 ft., one target. \textsl{Hit:} 5 (1d4 + 3) piercing damage.

\DndMonsterSection{Reactions}
\DndMonsterAction{Parry}
The captain adds 2 to its AC against one melee attack that would hit it. To do so, the captain must see the attacker and be wielding a melee weapon.

\end{DndMonster}



\begin{DndMonster}[float=t]{Bat}
\DndMonsterType{Tiny beast, unaligned}
\DndMonsterBasics[
armor-class = {12},
hit-points = {\DndDice{1d4 - 1}},
speed = {5 ft., fly 30 ft.}
]
\DndMonsterAbilityScores[
 str = 2,
 dex = 15,
 con = 8,
 int = 2,
 wis = 12,
 cha = 4
]
\DndMonsterDetails[
senses = {blindsight 60 ft., passive perception 11},
challenge = 0,
]
\DndMonsterAction{Echolocation}
The bat can't use its blindsight while deafened.

\DndMonsterAction{Keen Hearing}
The bat has advantage on Wisdom (Perception) checks that rely on hearing.

\DndMonsterSection{Actions}
\DndMonsterAction{Bite}
\textsl{Melee Weapon Attack:} 0 to hit, reach 5 ft., one creature. \textsl{Hit:} 1 piercing damage.

\end{DndMonster}



\begin{DndMonster}[float=t]{Bearded Devil}
\DndMonsterType{Medium fiend (devil), lawful evil}
\DndMonsterBasics[
armor-class = {13 (natural armor)},
hit-points = {\DndDice{8d8 + 16}},
speed = {30 ft.}
]
\DndMonsterAbilityScores[
 str = 16,
 dex = 15,
 con = 15,
 int = 9,
 wis = 11,
 cha = 11
]
\DndMonsterDetails[
saving-throws = {Str +5, Con +4, Wis +2},
damage-resistances = {cold; bludgeoning, piercing, slashing from nonmagical attacks that aren't silvered},
damage-immunities = {fire, poison},
condition-immunities = {poisoned},
senses = {darkvision 120 ft., passive perception 10},
languages = {Infernal, telepathy 120 ft.},
challenge = 3,
]
\DndMonsterAction{Devil's Sight}
Magical darkness doesn't impede the devil's darkvision.

\DndMonsterAction{Magic Resistance}
The devil has advantage on saving throws against spells and other magical effects.

\DndMonsterAction{Steadfast}
The devil can't be frightened while it can see an allied creature within 30 feet of it.

\DndMonsterSection{Actions}
\DndMonsterAction{Multiattack}
The devil makes two attacks: one with its beard and one with its glaive.

\DndMonsterAction{Beard}
\textsl{Melee Weapon Attack:} 5 to hit, reach 5 ft., one creature. \textsl{Hit:} 6 (1d8 + 2) piercing damage, and the target must succeed on a DC 12 Constitution saving throw or be poisoned for 1 minute. While poisoned in this way, the target can't regain hit points. The target can repeat the saving throw at the end of each of its turns, ending the effect on itself on a success.

\DndMonsterAction{Glaive}
\textsl{Melee Weapon Attack:} 5 to hit, reach 10 ft., one target. \textsl{Hit:} 8 (1d10 + 3) slashing damage. If the target is a creature other than an undead or a construct, it must succeed on a DC 12 Constitution saving throw or lose 5 (1d10) hit points at the start of each of its turns due to an infernal wound. Each time the devil hits the wounded target with this attack, the damage dealt by the wound increases by 5 (1d10). Any creature can take an action to stanch the wound with a successful DC 12 Wisdom (Medicine) check. The wound also closes if the target receives magical healing.

\end{DndMonster}



\begin{DndMonster}[float=t]{Berserker}
\DndMonsterType{Medium humanoid (any race), any chaotic alignment}
\DndMonsterBasics[
armor-class = {13 (\textit{hide armor})},
hit-points = {\DndDice{9d8 + 27}},
speed = {30 ft.}
]
\DndMonsterAbilityScores[
 str = 16,
 dex = 12,
 con = 17,
 int = 9,
 wis = 11,
 cha = 9
]
\DndMonsterDetails[
languages = {any one language (usually Common)},
challenge = 2,
]
\DndMonsterAction{Reckless}
At the start of its turn, the berserker can gain advantage on all melee weapon attack rolls during that turn, but attack rolls against it have advantage until the start of its next turn.

\DndMonsterSection{Actions}
\DndMonsterAction{Greataxe}
\textsl{Melee Weapon Attack:} 5 to hit, reach 5 ft., one target. \textsl{Hit:} 9 (1d12 + 3) slashing damage.

\end{DndMonster}



\begin{DndMonster}[float=t]{Big Linda}
\DndMonsterType{Huge beast, unaligned}
\DndMonsterBasics[
armor-class = {13 (natural armor)},
hit-points = {\DndDice{13d12 + 52}},
speed = {50 ft.}
]
\DndMonsterAbilityScores[
 str = 25,
 dex = 10,
 con = 19,
 int = 2,
 wis = 12,
 cha = 9
]
\DndMonsterDetails[
skills = {Perception +4},
damage-immunities = {acid},
senses = {blindsight 30 ft., passive perception 14},
challenge = 9,
]
\DndMonsterSection{Actions}
\DndMonsterAction{Multiattack}
Big Linda makes two attacks: one with its bite and one with its tail. It can't make both attacks against the same target.

\DndMonsterAction{Bite}
\textsl{Melee Weapon Attack:} 10 to hit, reach 10 ft., one target. \textsl{Hit:} 33 (4d12 + 7) piercing damage. If the target is a Medium or smaller creature, it is grappled (escape DC 17). Until this grapple ends, the target is restrained, and Big Linda can't bite another target.

\DndMonsterAction{Tail}
\textsl{Melee Weapon Attack:} 10 to hit, reach 10 ft., one target. \textsl{Hit:} 20 (3d8 + 7) bludgeoning damage.

\DndMonsterAction{Acid Spew (Recharge 5\textendash 6)}
Big Linda spews acid in a 30 foot cone. Each creature in that line of fire must succeed on a DC 15 Dexterity saving throw, taking 49 (11d8) acid damage on a failed save, or half as much damage on a successful one.

\end{DndMonster}


\clearpage

\begin{DndMonster}[float=t]{Black Pudding}
\DndMonsterType{Large ooze, unaligned}
\DndMonsterBasics[
armor-class = {7},
hit-points = {\DndDice{10d10 + 30}},
speed = {20 ft., climb 20 ft.}
]
\DndMonsterAbilityScores[
 str = 16,
 dex = 5,
 con = 16,
 int = 1,
 wis = 6,
 cha = 1
]
\DndMonsterDetails[
damage-immunities = {acid, cold, lightning, slashing},
condition-immunities = {blinded, charmed, deafened, exhaustion, frightened, prone},
senses = {blindsight 60 ft. (blind beyond this radius), passive perception 8},
challenge = 4,
]
\DndMonsterAction{Amorphous}
The pudding can move through a space as narrow as 1 inch wide without squeezing.

\DndMonsterAction{Corrosive Form}
A creature that touches the pudding or hits it with a melee attack while within 5 feet of it takes 4 (1d8) acid damage. Any nonmagical weapon made of metal or wood that hits the pudding corrodes. After dealing damage, the weapon takes a permanent and cumulative −1 penalty to damage rolls. If its penalty drops to −5, the weapon is destroyed. Nonmagical ammunition made of metal or wood that hits the pudding is destroyed after dealing damage. The pudding can eat through 2-inch-thick, nonmagical wood or metal in 1 round.

\DndMonsterAction{Spider Climb}
The pudding can climb difficult surfaces, including upside down on ceilings, without needing to make an ability check.

\DndMonsterSection{Actions}
\DndMonsterAction{Pseudopod}
\textsl{Melee Weapon Attack:} 5 to hit, reach 5 ft., one target. \textsl{Hit:} 6 (1d6 + 3) bludgeoning damage plus 18 (4d8) acid damage. In addition, nonmagical armor worn by the target is partly dissolved and takes a permanent and cumulative −1 penalty to the AC it offers. The armor is destroyed if the penalty reduces its AC to 10.

\DndMonsterSection{Reactions}
\DndMonsterAction{Split}
When a pudding that is Medium or larger is subjected to lightning or slashing damage, it splits into two new puddings if it has at least 10 hit points. Each new pudding has hit points equal to half the original pudding's, rounded down. New puddings are one size smaller than the original pudding.

\end{DndMonster}



\begin{DndMonster}[float=t]{Bone Devil}
\DndMonsterType{Large fiend (devil), lawful evil}
\DndMonsterBasics[
armor-class = {19 (natural armor)},
hit-points = {\DndDice{15d10 + 60}},
speed = {40 ft., fly 40 ft.}
]
\DndMonsterAbilityScores[
 str = 18,
 dex = 16,
 con = 18,
 int = 13,
 wis = 14,
 cha = 16
]
\DndMonsterDetails[
saving-throws = {Int +5, Wis +6, Cha +7},
skills = {Deception +7, Insight +6},
damage-resistances = {cold; bludgeoning, piercing, slashing from nonmagical attacks that aren't silvered},
damage-immunities = {fire, poison},
condition-immunities = {poisoned},
senses = {darkvision 120 ft., passive perception 12},
languages = {Infernal, telepathy 120 ft.},
challenge = 9,
]
\DndMonsterAction{Devil's Sight}
Magical darkness doesn't impede the devil's darkvision.

\DndMonsterAction{Magic Resistance}
The devil has advantage on saving throws against spells and other magical effects.

\DndMonsterSection{Actions}
\DndMonsterAction{Multiattack}
The devil makes three attacks: two with its claws and one with its sting.

\DndMonsterAction{Claw}
\textsl{Melee Weapon Attack:} 8 to hit, reach 10 ft., one target. \textsl{Hit:} 8 (1d8 + 4) slashing damage.

\DndMonsterAction{Sting}
\textsl{Melee Weapon Attack:} 8 to hit, reach 10 ft., one target. \textsl{Hit:} 13 (2d8 + 4) piercing damage plus 17 (5d6) poison damage, and the target must succeed on a DC 14 Constitution saving throw or become poisoned for 1 minute. The target can repeat the saving throw at the end of each of its turns, ending the effect on itself on a success.

\end{DndMonster}



\begin{DndMonster}[float=t]{Bugbear}
\DndMonsterType{Medium humanoid (goblinoid), chaotic evil}
\DndMonsterBasics[
armor-class = {16 (\textit{hide armor})},
hit-points = {\DndDice{5d8 + 5}},
speed = {30 ft.}
]
\DndMonsterAbilityScores[
 str = 15,
 dex = 14,
 con = 13,
 int = 8,
 wis = 11,
 cha = 9
]
\DndMonsterDetails[
skills = {Stealth +6, Survival +2},
senses = {darkvision 60 ft., passive perception 10},
languages = {Common, Goblin},
challenge = 1,
]
\DndMonsterAction{Brute}
A melee weapon deals one extra die of its damage when the bugbear hits with it (included in the attack).

\DndMonsterAction{Surprise Attack}
If the bugbear surprises a creature and hits it with an attack during the first round of combat, the target takes an extra 7 (2d6) damage from the attack.

\DndMonsterSection{Actions}
\DndMonsterAction{Morningstar}
\textsl{Melee Weapon Attack:} 4 to hit, reach 5 ft., one target. \textsl{Hit:} 11 (2d8 + 2) piercing damage.

\DndMonsterAction{Javelin}
\textsl{Melee or Ranged Weapon Attack:} 4 to hit, reach 5 ft. or range 30/120 ft., one target. \textsl{Hit:} 9 (2d6 + 2) piercing damage in melee or 5 (1d6 + 2) piercing damage at range.

\end{DndMonster}



\begin{DndMonster}[float=t]{Cat}
\DndMonsterType{Tiny beast, unaligned}
\DndMonsterBasics[
armor-class = {12},
hit-points = {\DndDice{1d4}},
speed = {40 ft., climb 30 ft.}
]
\DndMonsterAbilityScores[
 str = 3,
 dex = 15,
 con = 10,
 int = 3,
 wis = 12,
 cha = 7
]
\DndMonsterDetails[
skills = {Perception +3, Stealth +4},
challenge = 0,
]
\DndMonsterAction{Keen Smell}
The cat has advantage on Wisdom (Perception) checks that rely on smell.

\DndMonsterSection{Actions}
\DndMonsterAction{Claws}
\textsl{Melee Weapon Attack:} 0 to hit, reach 5 ft., one target. \textsl{Hit:} 1 slashing damage.

\end{DndMonster}



\begin{DndMonster}[float=t]{Cavalier}
\DndMonsterType{Medium humanoid, any alignment}
\DndMonsterBasics[
armor-class = {19 (splint armor)},
hit-points = {\DndDice{12d8 + 36}},
speed = {30 ft.}
]
\DndMonsterAbilityScores[
 str = 18,
 dex = 10,
 con = 16,
 int = 10,
 wis = 12,
 cha = 15
]
\DndMonsterDetails[
saving-throws = {Wis +4, Cha +5},
skills = {Animal Handling +4, Athletics +7, Perception +4, Persuasion +5},
languages = {Any two},
challenge = 5,
]
\DndMonsterAction{Bravery}
The cavalier has advantage on saving throws against being frightened.

\DndMonsterAction{Light's Fury}
The cavalier deals an extra 9 (2d8) radiant damage when it hits with a weapon attack.

\DndMonsterAction{Skilled Rider}
While mounted, the cavalier may force an attack targeted at their mount to target them instead.

\DndMonsterAction{Spellcasting}
The cavalier is an 8th-level spellcaster that uses Charisma as its spellcasting ability (spell save DC 13; 5 to hit with spell attacks). They have the following paladin spells prepared:
\begin{DndMonsterSpells}
   \DndMonsterSpellLevel[1][4]{\textit{bless}, \textit{command}, \textit{cure wounds}}
   \DndMonsterSpellLevel[2][3]{\textit{branding smite}, \textit{find steed}}
\end{DndMonsterSpells}
\DndMonsterSection{Actions}
\DndMonsterAction{Multiattack}
The cavalier makes two melee attacks.

\DndMonsterAction{Lance}
\textsl{Melee Weapon Attack:} 7 to hit (with disadvantage against a target within 5 ft.), reach 10 ft., one target. \textsl{Hit:} 10 (1d12 + 4) piercing damage plus 9 (2d8) radiant damage.

\DndMonsterAction{Javelin}
\textsl{Melee or Ranged Weapon Attack:} 7 to hit, reach 5 ft. or range 30/120 ft., one target \textsl{Hit:} 7 (1d6 + 4) piercing damage plus 9 (2d8) radiant damage.

\end{DndMonster}



\begin{DndMonster}[float=t]{Chaplain}
\DndMonsterType{Medium humanoid, any alignment}
\DndMonsterBasics[
armor-class = {18 (\textit{chain mail})},
hit-points = {\DndDice{9d8 + 18}},
speed = {30 ft.}
]
\DndMonsterAbilityScores[
 str = 13,
 dex = 10,
 con = 14,
 int = 11,
 wis = 17,
 cha = 13
]
\DndMonsterDetails[
saving-throws = {Wis +6, Cha +4},
skills = {Insight +6, Persuasion +4, Religion +3},
languages = {Any two},
challenge = 6,
]
\DndMonsterAction{Devotion}
The chaplain has advantage on saving throws against being charmed or frightened.

\DndMonsterAction{Spellcasting}
The chaplain is a 9th-level spellcaster that uses Wisdom as its spellcasting ability (spell save DC 12; 4 to hit with spell attacks). They have the following cleric spells prepared:
\begin{DndMonsterSpells}
  \DndMonsterSpellLevel{\textit{light}, \textit{mending}, \textit{sacred flame}, \textit{spare the dying}}
   \DndMonsterSpellLevel[1][4]{\textit{bless}, \textit{guiding bolt}, \textit{healing word}}
   \DndMonsterSpellLevel[2][3]{\textit{lesser restoration}, \textit{prayer of healing}, \textit{spiritual weapon}}
   \DndMonsterSpellLevel[3][3]{\textit{dispel magic}, \textit{purge contamination}, \textit{spirit guardians}}
   \DndMonsterSpellLevel[4][3]{\textit{banishment}, \textit{divination}}
   \DndMonsterSpellLevel[5][1]{\textit{flame strike}}
\end{DndMonsterSpells}
\DndMonsterSection{Actions}
\DndMonsterAction{Spear}
\textsl{Melee or Ranged Weapon Attack:} 3 to hit, reach 5 ft. or range 20/60 ft., one target. \textsl{Hit:} 4 (1d6 + 1) piercing damage or 5 (1d8 + 1) damage when used in two hands to make a melee attack.

\end{DndMonster}



\begin{DndMonster}[float=t]{Chimera}
\DndMonsterType{Large monstrosity, chaotic evil}
\DndMonsterBasics[
armor-class = {14 (natural armor)},
hit-points = {\DndDice{12d10 + 48}},
speed = {30 ft., fly 60 ft.}
]
\DndMonsterAbilityScores[
 str = 19,
 dex = 11,
 con = 19,
 int = 3,
 wis = 14,
 cha = 10
]
\DndMonsterDetails[
skills = {Perception +8},
senses = {darkvision 60 ft., passive perception 18},
languages = {understands Draconic but can't speak},
challenge = 6,
]
\DndMonsterSection{Actions}
\DndMonsterAction{Multiattack}
The chimera makes three attacks: one with its bite, one with its horns, and one with its claws. When its fire breath is available, it can use the breath in place of its bite or horns.

\DndMonsterAction{Bite}
\textsl{Melee Weapon Attack:} 7 to hit, reach 5 ft., one target. \textsl{Hit:} 11 (2d6 + 4) piercing damage.

\DndMonsterAction{Horns}
\textsl{Melee Weapon Attack:} 7 to hit, reach 5 ft., one target. \textsl{Hit:} 10 (1d12 + 4) bludgeoning damage.

\DndMonsterAction{Claws}
\textsl{Melee Weapon Attack:} 7 to hit, reach 5 ft., one target. \textsl{Hit:} 11 (2d6 + 4) slashing damage.

\DndMonsterAction{Fire Breath (Recharge 5\textendash 6)}
The dragon head exhales fire in a 15-foot cone. Each creature in that area must make a DC 15 Dexterity saving throw, taking 31 (7d8) fire damage on a failed save, or half as much damage on a successful one.

\end{DndMonster}



\begin{DndMonster}[float=t]{Chuul}
\DndMonsterType{Large aberration, chaotic evil}
\DndMonsterBasics[
armor-class = {16 (natural armor)},
hit-points = {\DndDice{11d10 + 33}},
speed = {30 ft., swim 30 ft.}
]
\DndMonsterAbilityScores[
 str = 19,
 dex = 10,
 con = 16,
 int = 5,
 wis = 11,
 cha = 5
]
\DndMonsterDetails[
skills = {Perception +4},
damage-immunities = {poison},
condition-immunities = {poisoned},
senses = {darkvision 60 ft., passive perception 14},
languages = {understands Deep Speech but can't speak},
challenge = 4,
]
\DndMonsterAction{Amphibious}
The chuul can breathe air and water.

\DndMonsterAction{Sense Magic}
The chuul senses magic within 120 feet of it at will. This trait otherwise works like the \textit{detect magic} spell but isn't itself magical.

\DndMonsterSection{Actions}
\DndMonsterAction{Multiattack}
The chuul makes two pincer attacks. If the chuul is grappling a creature, the chuul can also use its tentacles once.

\DndMonsterAction{Pincer}
\textsl{Melee Weapon Attack:} 6 to hit, reach 10 ft., one target. \textsl{Hit:} 11 (2d6 + 4) bludgeoning damage. The target is grappled (escape DC 14) if it is a Large or smaller creature and the chuul doesn't have two other creatures grappled.

\DndMonsterAction{Tentacles}
One creature grappled by the chuul must succeed on a DC 13 Constitution saving throw or be poisoned for 1 minute. Until this poison ends, the target is paralyzed. The target can repeat the saving throw at the end of each of its turns, ending the effect on itself on a success.

\end{DndMonster}



\begin{DndMonster}[float=t]{Cloaker}
\DndMonsterType{Large aberration, chaotic neutral}
\DndMonsterBasics[
armor-class = {14 (natural armor)},
hit-points = {\DndDice{12d10 + 12}},
speed = {10 ft., fly 40 ft.}
]
\DndMonsterAbilityScores[
 str = 17,
 dex = 15,
 con = 12,
 int = 13,
 wis = 12,
 cha = 14
]
\DndMonsterDetails[
skills = {Stealth +5},
senses = {darkvision 60 ft., passive perception 11},
languages = {Deep Speech, Undercommon},
challenge = 8,
]
\DndMonsterAction{Damage Transfer}
While attached to a creature, the cloaker takes only half the damage dealt to it (rounded down). and that creature takes the other half.

\DndMonsterAction{False Appearance}
While the cloaker remains motionless without its underside exposed, it is indistinguishable from a dark leather cloak.

\DndMonsterAction{Light Sensitivity}
While in bright light, the cloaker has disadvantage on attack rolls and Wisdom (Perception) checks that rely on sight.

\DndMonsterSection{Actions}
\DndMonsterAction{Multiattack}
The cloaker makes two attacks: one with its bite and one with its tail.

\DndMonsterAction{Bite}
\textsl{Melee Weapon Attack:} 6 to hit, reach 5 ft., one creature. \textsl{Hit:} 10 (2d6 + 3) piercing damage, and if the target is Large or smaller, the cloaker attaches to it. If the cloaker has advantage against the target, the cloaker attaches to the target's head, and the target is blinded and unable to breathe while the cloaker is attached. While attached, the cloaker can make this attack only against the target and has advantage on the attack roll. The cloaker can detach itself by spending 5 feet of its movement. A creature, including the target, can take its action to detach the cloaker by succeeding on a DC 16 Strength check.

\DndMonsterAction{Tail}
\textsl{Melee Weapon Attack:} 6 to hit, reach 10 ft., one creature. \textsl{Hit:} 7 (1d8 + 3) slashing damage.

\DndMonsterAction{Moan}
Each creature within 60 feet of the cloaker that can hear its moan and that isn't an aberration must succeed on a DC 13 Wisdom saving throw or become frightened until the end of the cloaker's next turn. If a creature's saving throw is successful, the creature is immune to the cloaker's moan for the next 24 hours.

\DndMonsterAction{Phantasms (Recharges after a Short or Long Rest)}
The cloaker magically creates three illusory duplicates of itself if it isn't in bright light. The duplicates move with it and mimic its actions, shifting position so as to make it impossible to track which cloaker is the real one. If the cloaker is ever in an area of bright light, the duplicates disappear.

Whenever any creature targets the cloaker with an attack or a harmful spell while a duplicate remains, that creature rolls randomly to determine whether it targets the cloaker or one of the duplicates. A creature is unaffected by this magical effect if it can't see or if it relies on senses other than sight.

A duplicate has the cloaker's AC and uses its saving throws. If an attack hits a duplicate, or if a duplicate fails a saving throw against an effect that deals damage, the duplicate disappears.

\end{DndMonster}



\begin{DndMonster}[float=t]{Commoner}
\DndMonsterType{Medium humanoid (any race), any alignment}
\DndMonsterBasics[
armor-class = {10},
hit-points = {\DndDice{1d8}},
speed = {30 ft.}
]
\DndMonsterAbilityScores[
 str = 10,
 dex = 10,
 con = 10,
 int = 10,
 wis = 10,
 cha = 10
]
\DndMonsterDetails[
languages = {any one language (usually Common)},
challenge = 0,
]
\DndMonsterSection{Actions}
\DndMonsterAction{Club}
\textsl{Melee Weapon Attack:} 2 to hit, reach 5 ft., one target. \textsl{Hit:} 2 (1d4) bludgeoning damage.

\end{DndMonster}



\begin{DndMonster}[float=t]{Crater Wurm}
\DndMonsterType{Gargantuan monstrosity, unaligned}
\DndMonsterBasics[
armor-class = {18 (natural armor)},
hit-points = {\DndDice{15d20 + 90}},
speed = {50 ft., burrow 30 ft.}
]
\DndMonsterAbilityScores[
 str = 28,
 dex = 7,
 con = 22,
 int = 1,
 wis = 8,
 cha = 4
]
\DndMonsterDetails[
saving-throws = {Con +11, Wis +4},
senses = {blindsight 30 ft., tremorsense 60 ft., passive perception 9},
challenge = 15,
]
\DndMonsterAction{Copy Warning}
This content has been copied from another creature, but the data replacement hasn't been run. Check details.

\DndMonsterAction{Tunneler}
The worm can burrow through solid rock at half its burrow speed and leaves a 10-foot-diameter tunnel in its wake.

\DndMonsterSection{Actions}
\DndMonsterAction{Multiattack}
The worm makes two attacks: one with its bite and one with its stinger.

\DndMonsterAction{Bite}
\textsl{Melee Weapon Attack:} 14 to hit, reach 10 ft., one target. \textsl{Hit:} 22 (3d8 + 9) piercing damage. If the target is a Large or smaller creature, it must succeed on a DC 19 Dexterity saving throw or be swallowed by the worm. A swallowed creature is blinded and restrained, it has total cover against attacks and other effects outside the worm, and it takes 21 (6d6) acid damage at the start of each of the worm's turns.

If the worm takes 30 damage or more on a single turn from a creature inside it, the worm must succeed on a DC 21 Constitution saving throw at the end of that turn or regurgitate all swallowed creatures, which fall prone in a space within 10 feet of the worm. If the worm dies, a swallowed creature is no longer restrained by it and can escape from the corpse by using 20 feet of movement, exiting prone.

\DndMonsterAction{Tail Stinger}
\textsl{Melee Weapon Attack:} 14 to hit, reach 10 ft., one creature. \textsl{Hit:} 19 (3d6 + 9) piercing damage, and the target must make a DC 19 Constitution saving throw, taking 42 (12d6) poison damage on a failed save, or half as much damage on a successful one.

\end{DndMonster}



\begin{DndMonster}[float=t]{Crimson Countess}
\DndMonsterType{Medium monstrosity, chaotic evil}
\DndMonsterBasics[
armor-class = {17 (natural armor)},
hit-points = {\DndDice{20d8 + 40}},
speed = {20 ft., fly 60 ft.}
]
\DndMonsterAbilityScores[
 str = 10,
 dex = 21,
 con = 15,
 int = 13,
 wis = 14,
 cha = 19
]
\DndMonsterDetails[
saving-throws = {Con +5, Wis +5},
skills = {Acrobatics +11, Deception +7, Perception +8},
damage-immunities = {lightning, thunder},
senses = {darkvision 120 ft., passive perception 16},
languages = {Abyssal},
challenge = 8,
]
\DndMonsterAction{Aerial Acrobat}
On each of her turns, the Crimson Countess can use a bonus action to take the Dash, Disengage, or Dodge action.

\DndMonsterAction{Keen Sight}
The Crimson Countess has advantage on Wisdom (Perception) checks that rely on sight.

\DndMonsterSection{Actions}
\DndMonsterAction{Multiattack}
The Crimson Countess can use her Frightful Screech. She then makes two attacks with her talons or one with her javelin.

\DndMonsterAction{Talons}
\textsl{Melee Weapon Attack:} 8 to hit, reach 5 ft., one target. \textsl{Hit:} 11 (1d10 + 5) slashing damage.

\DndMonsterAction{Frightful Screech}
Each creature of the Crimson Countess' choice that is within 120 feet of her and not deafened must succeed on a DC 15 Wisdom saving throw or become frightened for 1 minute. If a creature's saving throw is successful or the effect ends for it, the creature is immune to the effects of the Crimson Countess' Frightful Screech for the next 24 hours.

While frightened by this effect, a creature must take the Dash action and move away from the Crimson Countess by the safest available route on each of its turns, unless there is nowhere to move. If the creature ends its turn in a location where it doesn't have line of sight to the Crimson Countess, the creature can repeat the saving throw. On a successful save, the effect ends for that creature.

\DndMonsterAction{Javelin}
\textsl{Melee or Ranged Weapon Attack:} 8 to hit, reach 5 ft. or ranged 20/60 feet, one target. \textsl{Hit:} 9 (1d6 + 5) piercing damage, plus 14 (4d6) lightning damage.

\DndMonsterAction{Thunder Blast (Recharge 5\textendash 6)}
The Crimson Countess unleashes a thunderous cry in a 15-foot cube. Each creature in the cube must succeed on a DC 15 Constitution saving throw, taking 14 (4d6) thunder damage on a failed save, or half as much damage on a successful one. Creatures who fail this saving throw are pushed 10 feet and knocked prone.

\end{DndMonster}



\begin{DndMonster}[float=t]{Cult Fanatic}
\DndMonsterType{Medium humanoid (any race), any non-good alignment}
\DndMonsterBasics[
armor-class = {13 (\textit{leather armor})},
hit-points = {\DndDice{6d8 + 6}},
speed = {30 ft.}
]
\DndMonsterAbilityScores[
 str = 11,
 dex = 14,
 con = 12,
 int = 10,
 wis = 13,
 cha = 14
]
\DndMonsterDetails[
skills = {Deception +4, Persuasion +4, Religion +2},
languages = {any one language (usually Common)},
challenge = 2,
]
\DndMonsterAction{Dark Devotion}
The fanatic has advantage on saving throws against being charmed or frightened.

\DndMonsterAction{Spellcasting}
The fanatic is a 4th-level spellcaster. Its spellcasting ability is Wisdom (spell save DC 11, 3 to hit with spell attacks). The fanatic has the following cleric spells prepared:
\begin{DndMonsterSpells}
  \DndMonsterSpellLevel{\textit{light}, \textit{sacred flame}, \textit{thaumaturgy}}
   \DndMonsterSpellLevel[1][4]{\textit{command}, \textit{inflict wounds}, \textit{shield of faith}}
   \DndMonsterSpellLevel[2][3]{\textit{hold person}, \textit{spiritual weapon}}
\end{DndMonsterSpells}
\DndMonsterSection{Actions}
\DndMonsterAction{Multiattack}
The fanatic makes two melee attacks.

\DndMonsterAction{Dagger}
\textsl{Melee or Ranged Weapon Attack:} 4 to hit, reach 5 ft. or range 20/60 ft., one creature. \textsl{Hit:} 4 (1d4 + 2) piercing damage.

\end{DndMonster}



\begin{DndMonster}[float=t]{Cultist}
\DndMonsterType{Medium humanoid (any race), any non-good alignment}
\DndMonsterBasics[
armor-class = {12 (\textit{leather armor})},
hit-points = {\DndDice{2d8}},
speed = {30 ft.}
]
\DndMonsterAbilityScores[
 str = 11,
 dex = 12,
 con = 10,
 int = 10,
 wis = 11,
 cha = 10
]
\DndMonsterDetails[
skills = {Deception +2, Religion +2},
languages = {any one language (usually Common)},
challenge = 1/8,
]
\DndMonsterAction{Dark Devotion}
The cultist has advantage on saving throws against being charmed or frightened.

\DndMonsterSection{Actions}
\DndMonsterAction{Scimitar}
\textsl{Melee Weapon Attack:} 3 to hit, reach 5 ft., one creature. \textsl{Hit:} 4 (1d6 + 1) slashing damage.

\end{DndMonster}



\begin{DndMonster}[float=t]{Delerium Dreg}
\DndMonsterType{Medium aberration, unaligned}
\DndMonsterBasics[
armor-class = {11},
hit-points = {\DndDice{5d8}},
speed = {30 ft.}
]
\DndMonsterAbilityScores[
 str = 15,
 dex = 13,
 con = 11,
 int = 7,
 wis = 10,
 cha = 5
]
\DndMonsterDetails[
senses = {darkvision 60 ft., passive perception 10},
languages = {Common},
challenge = 1/2,
]
\DndMonsterAction{Misty Camouflage}
A delerium dreg has advantage on Dexterity (Stealth) checks to hide in any area obscured by mist or fog.

\DndMonsterSection{Actions}
\DndMonsterAction{Spear}
\textsl{Melee or Ranged Weapon Attack:} 4 to hit, reach 5 ft. or range 20/60 ft., one target. \textsl{Hit:} 5 (1d6 + 2) piercing damage, or 6 (1d8 + 2) piercing damage if used with two hands to make a melee attack.

\end{DndMonster}



\begin{DndMonster}[float=t]{Deva}
\DndMonsterType{Medium celestial, lawful good}
\DndMonsterBasics[
armor-class = {17 (natural armor)},
hit-points = {\DndDice{16d8 + 64}},
speed = {30 ft., fly 90 ft.}
]
\DndMonsterAbilityScores[
 str = 18,
 dex = 18,
 con = 18,
 int = 17,
 wis = 20,
 cha = 20
]
\DndMonsterDetails[
saving-throws = {Wis +9, Cha +9},
skills = {Insight +9, Perception +9},
damage-resistances = {radiant; bludgeoning, piercing, slashing from nonmagical attacks},
condition-immunities = {charmed, exhaustion, frightened},
senses = {darkvision 120 ft., passive perception 19},
languages = {all, telepathy 120 ft.},
challenge = 10,
]
\DndMonsterAction{Angelic Weapons}
The deva's weapon attacks are magical. When the deva hits with any weapon, the weapon deals an extra 4d8 radiant damage (included in the attack).

\DndMonsterAction{Magic Resistance}
The deva has advantage on saving throws against spells and other magical effects.

\DndMonsterAction{Innate Spellcasting}
The deva's spellcasting ability is Charisma (spell save DC 17). The deva can innately cast the following spells, requiring only verbal components:
\begin{DndMonsterSpells}
  \DndInnateSpellLevel{\textit{detect evil and good}}
  \DndInnateSpellLevel[1e]{\textit{commune}, \textit{raise dead}}
\end{DndMonsterSpells}
\DndMonsterSection{Actions}
\DndMonsterAction{Multiattack}
The deva makes two melee attacks.

\DndMonsterAction{Mace}
\textsl{Melee Weapon Attack:} 8 to hit, reach 5 ft., one target. \textsl{Hit:} 7 (1d6 + 4) bludgeoning damage plus 18 (4d8) radiant damage.

\DndMonsterAction{Healing Touch (3/Day)}
The deva touches another creature. The target magically regains 20 (4d8 + 2) hit points and is freed from any curse, disease, poison, blindness, or deafness.

\DndMonsterAction{Change Shape}
The deva magically polymorphs into a humanoid or beast that has a challenge rating equal to or less than its own, or back into its true form. It reverts to its true form if it dies. Any equipment it is wearing or carrying is absorbed or borne by the new form (the deva's choice).

In a new form, the deva retains its game statistics and ability to speak, but its AC, movement modes, Strength, Dexterity, and special senses are replaced by those of the new form, and it gains any statistics and capabilities (except class features, legendary actions, and lair actions) that the new form has but that it lacks.

\end{DndMonster}



\begin{DndMonster}[float=t]{Djinni}
\DndMonsterType{Large elemental, chaotic good}
\DndMonsterBasics[
armor-class = {17 (natural armor)},
hit-points = {\DndDice{14d10 + 84}},
speed = {30 ft., fly 90 ft.}
]
\DndMonsterAbilityScores[
 str = 21,
 dex = 15,
 con = 22,
 int = 15,
 wis = 16,
 cha = 20
]
\DndMonsterDetails[
saving-throws = {Dex +6, Wis +7, Cha +9},
damage-immunities = {lightning, thunder},
senses = {darkvision 120 ft., passive perception 13},
languages = {Auran},
challenge = 11,
]
\DndMonsterAction{Elemental Demise}
If the djinni dies, its body disintegrates into a warm breeze, leaving behind only equipment the djinni was wearing or carrying.

\DndMonsterAction{Innate Spellcasting}
The djinni's innate spellcasting ability is Charisma (spell save DC 17, 9 to hit with spell attacks). It can innately cast the following spells, requiring no material components:
\begin{DndMonsterSpells}
  \DndInnateSpellLevel{\textit{detect evil and good}, \textit{detect magic}, \textit{thunderwave}}
  \DndInnateSpellLevel[3e]{\textit{create food and water} (can create wine instead of water), \textit{tongues}, \textit{wind walk}}
  \DndInnateSpellLevel[1e]{\textit{conjure elemental} (\textbf{air elemental} only), \textit{creation}, \textit{gaseous form}, \textit{invisibility}, \textit{major image}, \textit{plane shift}}
\end{DndMonsterSpells}
\DndMonsterSection{Actions}
\DndMonsterAction{Multiattack}
The djinni makes three scimitar attacks.

\DndMonsterAction{Scimitar}
\textsl{Melee Weapon Attack:} 9 to hit, reach 5 ft., one target. \textsl{Hit:} 12 (2d6 + 5) slashing damage plus 3 (1d6) lightning or thunder damage (djinni's choice).

\DndMonsterAction{Create Whirlwind}
A 5-foot-radius, 30-foot-tall cylinder of swirling air magically forms on a point the djinni can see within 120 feet of it. The whirlwind lasts as long as the djinni maintains concentration (as if concentration on a spell). Any creature but the djinni that enters the whirlwind must succeed on a DC 18 Strength saving throw or be restrained by it. The djinni can move the whirlwind up to 60 feet as an action, and creatures restrained by the whirlwind move with it. The whirlwind ends if the djinni loses sight of it.

A creature can use its action to free a creature restrained by the whirlwind, including itself, by succeeding on a DC 18 Strength check. If the check succeeds, the creature is no longer restrained and moves to the nearest space outside the whirlwind.

\end{DndMonster}



\begin{DndMonster}[float=t]{Doppelganger}
\DndMonsterType{Medium monstrosity (shapechanger), neutral}
\DndMonsterBasics[
armor-class = {14},
hit-points = {\DndDice{8d8 + 16}},
speed = {30 ft.}
]
\DndMonsterAbilityScores[
 str = 11,
 dex = 18,
 con = 14,
 int = 11,
 wis = 12,
 cha = 14
]
\DndMonsterDetails[
skills = {Deception +6, Insight +3},
condition-immunities = {charmed},
senses = {darkvision 60 ft., passive perception 11},
languages = {Common},
challenge = 3,
]
\DndMonsterAction{Shapechanger}
The doppelganger can use its action to polymorph into a Small or Medium humanoid it has seen, or back into its true form. Its statistics, other than its size, are the same in each form. Any equipment it is wearing or carrying isn't transformed. It reverts to its true form if it dies.

\DndMonsterAction{Ambusher}
In the first round of a combat, the doppelganger has advantage on attack rolls against any creature it Surprise.

\DndMonsterAction{Surprise Attack}
If the doppelganger surprises a creature and hits it with an attack during the first round of combat, the target takes an extra 10 (3d6) damage from the attack.

\DndMonsterSection{Actions}
\DndMonsterAction{Multiattack}
The doppelganger makes two melee attacks.

\DndMonsterAction{Slam}
\textsl{Melee Weapon Attack:} 6 to hit, reach 5 ft., one target. \textsl{Hit:} 7 (1d6 + 4) bludgeoning damage.

\DndMonsterAction{Read Thoughts}
The doppelganger magically reads the surface thoughts of one creature within 60 feet of it. The effect can penetrate barriers, but 3 feet of wood or dirt, 2 feet of stone, 2 inches of metal, or a thin sheet of lead blocks it. While the target is in range, the doppelganger can continue reading its thoughts, as long as the doppelganger's concentration isn't broken (as if concentration on a spell). While reading the target's mind, the doppelganger has advantage on Wisdom (Insight) and Charisma (Deception, Intimidation, and Persuasion) checks against the target.

\end{DndMonster}



\begin{DndMonster}[float=t]{Draft Horse}
\DndMonsterType{Large beast, unaligned}
\DndMonsterBasics[
armor-class = {10},
hit-points = {\DndDice{3d10 + 3}},
speed = {40 ft.}
]
\DndMonsterAbilityScores[
 str = 18,
 dex = 10,
 con = 12,
 int = 2,
 wis = 11,
 cha = 7
]
\DndMonsterDetails[
challenge = 1/4,
]
\DndMonsterSection{Actions}
\DndMonsterAction{Hooves}
\textsl{Melee Weapon Attack:} 6 to hit, reach 5 ft., one target. \textsl{Hit:} 9 (2d4 + 4) bludgeoning damage.

\end{DndMonster}



\begin{DndMonster}[float=t]{Druid}
\DndMonsterType{Medium humanoid (any race), any alignment}
\DndMonsterBasics[
armor-class = {11 (16 with \textit{barkskin})},
hit-points = {\DndDice{5d8 + 5}},
speed = {30 ft.}
]
\DndMonsterAbilityScores[
 str = 10,
 dex = 12,
 con = 13,
 int = 12,
 wis = 15,
 cha = 11
]
\DndMonsterDetails[
skills = {Medicine +4, Nature +3, Perception +4},
languages = {Druidic plus any two languages},
challenge = 2,
]
\DndMonsterAction{Spellcasting}
The druid is a 4th-level spellcaster. Its spellcasting ability is Wisdom (spell save DC 12, 4 to hit with spell attacks). It has the following druid spells prepared:
\begin{DndMonsterSpells}
  \DndMonsterSpellLevel{\textit{druidcraft}, \textit{produce flame}, \textit{shillelagh}}
   \DndMonsterSpellLevel[1][4]{\textit{entangle}, \textit{longstrider}, \textit{speak with animals}, \textit{thunderwave}}
   \DndMonsterSpellLevel[2][3]{\textit{animal messenger}, \textit{barkskin}}
\end{DndMonsterSpells}
\DndMonsterSection{Actions}
\DndMonsterAction{Quarterstaff}
\textsl{Melee Weapon Attack:} 2 to hit (4 to hit with shillelagh), reach 5 ft., one target. \textsl{Hit:} 3 (1d6) bludgeoning damage, 4 (1d8) bludgeoning damage if wielded with two hands, or 6 (1d8 + 2) bludgeoning damage with \textit{shillelagh}.

\end{DndMonster}


\clearpage

\begin{DndMonster}[float=t]{Dryad}
\DndMonsterType{Medium fey, neutral}
\DndMonsterBasics[
armor-class = {11 (16 with \textit{barkskin})},
hit-points = {\DndDice{5d8}},
speed = {30 ft.}
]
\DndMonsterAbilityScores[
 str = 10,
 dex = 12,
 con = 11,
 int = 14,
 wis = 15,
 cha = 18
]
\DndMonsterDetails[
skills = {Perception +4, Stealth +5},
senses = {darkvision 60 ft., passive perception 14},
languages = {Elvish, Sylvan},
challenge = 1,
]
\DndMonsterAction{Magic Resistance}
The dryad has advantage on saving throws against spells and other magical effects.

\DndMonsterAction{Speak with Beasts and Plants}
The dryad can communicate with beasts and plants as if they shared a language.

\DndMonsterAction{Tree Stride}
Once on her turn, the dryad can use 10 feet of her movement to step magically into one living tree within her reach and emerge from a second living tree within 60 feet of the first tree, appearing in an unoccupied space within 5 feet of the second tree. Both trees must be large or bigger.

\DndMonsterAction{Innate Spellcasting}
The dryad's innate spellcasting ability is Charisma (spell save DC 14). The dryad can innately cast the following spells, requiring no material components:
\begin{DndMonsterSpells}
  \DndInnateSpellLevel{\textit{druidcraft}}
  \DndInnateSpellLevel[3e]{\textit{entangle}, \textit{goodberry}}
  \DndInnateSpellLevel[1e]{\textit{barkskin}, \textit{pass without trace}, \textit{shillelagh}}
\end{DndMonsterSpells}
\DndMonsterSection{Actions}
\DndMonsterAction{Club}
\textsl{Melee Weapon Attack:} 2 to hit (6 to hit with shillelagh), reach 5 ft., one target. \textsl{Hit:} 2 (1d4) bludgeoning damage, or 8 (1d8 + 4) bludgeoning damage with shillelagh.

\DndMonsterAction{Fey Charm}
The dryad targets one humanoid or beast that she can see within 30 feet of her. If the target can see the dryad, it must succeed on a DC 14 Wisdom saving throw or be magically charmed. The charmed creature regards the dryad as a trusted friend to be heeded and protected. Although the target isn't under the dryad's control, it takes the dryad's requests or actions in the most favorable way it can.

Each time the dryad or its allies do anything harmful to the target, it can repeat the saving throw, ending the effect on itself on a success. Otherwise, the effect lasts 24 hours or until the dryad dies, is on a different plane of existence from the target, or ends the effect as a bonus action. If a target's saving throw is successful, the target is immune to the dryad's Fey Charm for the next 24 hours.

The dryad can have no more than one humanoid and up to three beasts charmed at a time.

\end{DndMonster}



\begin{DndMonster}[float=t]{Dust Mephit}
\DndMonsterType{Small elemental, neutral evil}
\DndMonsterBasics[
armor-class = {12},
hit-points = {\DndDice{5d6}},
speed = {30 ft., fly 30 ft.}
]
\DndMonsterAbilityScores[
 str = 5,
 dex = 14,
 con = 10,
 int = 9,
 wis = 11,
 cha = 10
]
\DndMonsterDetails[
skills = {Perception +2, Stealth +4},
damage-vulnerabilities = {fire},
damage-immunities = {poison},
condition-immunities = {poisoned},
senses = {darkvision 60 ft., passive perception 12},
languages = {Auran, Terran},
challenge = 1/2,
]
\DndMonsterAction{Death Burst}
When the mephit dies, it explodes in a burst of dust. Each creature within 5 feet of it must then succeed on a DC 10 Constitution saving throw or be blinded for 1 minute. A blinded creature can repeat the saving throw on each of its turns, ending the effect on itself on a success.

\DndMonsterAction{Innate Spellcasting (1/Day)}
The mephit can innately cast \textit{sleep}, requiring no material components. Its innate spellcasting ability is Charisma.
\begin{DndMonsterSpells}
  \DndInnateSpellLevel{\textit{sleep}}
\end{DndMonsterSpells}
\DndMonsterSection{Actions}
\DndMonsterAction{Claws}
\textsl{Melee Weapon Attack:} 4 to hit, reach 5 ft., one creature. \textsl{Hit:} 4 (1d4 + 2) slashing damage.

\DndMonsterAction{Blinding Breath (Recharge 6)}
The mephit exhales a 15-foot cone of blinding dust. Each creature in that area must succeed on a DC 10 Dexterity saving throw or be blinded for 1 minute. A creature can repeat the saving throw at the end of each of its turns, ending the effect on itself on a success.

\end{DndMonster}



\begin{DndMonster}[float=t]{Earth Elemental}
\DndMonsterType{Large elemental, neutral}
\DndMonsterBasics[
armor-class = {17 (natural armor)},
hit-points = {\DndDice{12d10 + 60}},
speed = {30 ft., burrow 30 ft.}
]
\DndMonsterAbilityScores[
 str = 20,
 dex = 8,
 con = 20,
 int = 5,
 wis = 10,
 cha = 5
]
\DndMonsterDetails[
damage-vulnerabilities = {thunder},
damage-resistances = {bludgeoning, piercing, slashing from nonmagical attacks},
damage-immunities = {poison},
condition-immunities = {exhaustion, paralyzed, petrified, poisoned, unconscious},
senses = {darkvision 60 ft., tremorsense 60 ft., passive perception 10},
languages = {Terran},
challenge = 5,
]
\DndMonsterAction{Earth Glide}
The elemental can burrow through nonmagical, unworked earth and stone. While doing so, the elemental doesn't disturb the material it moves through.

\DndMonsterAction{Siege Monster}
The elemental deals double damage to objects and structures.

\DndMonsterSection{Actions}
\DndMonsterAction{Multiattack}
The elemental makes two slam attacks.

\DndMonsterAction{Slam}
\textsl{Melee Weapon Attack:} 8 to hit, reach 10 ft., one target. \textsl{Hit:} 14 (2d8 + 5) bludgeoning damage.

\end{DndMonster}



\begin{DndMonster}[float=t]{Efreeti}
\DndMonsterType{Large elemental, lawful evil}
\DndMonsterBasics[
armor-class = {17 (natural armor)},
hit-points = {\DndDice{16d10 + 112}},
speed = {40 ft., fly 60 ft.}
]
\DndMonsterAbilityScores[
 str = 22,
 dex = 12,
 con = 24,
 int = 16,
 wis = 15,
 cha = 16
]
\DndMonsterDetails[
saving-throws = {Int +7, Wis +6, Cha +7},
damage-immunities = {fire},
senses = {darkvision 120 ft., passive perception 12},
languages = {Ignan},
challenge = 11,
]
\DndMonsterAction{Elemental Demise}
If the efreeti dies, its body disintegrates in a flash of fire and puff of smoke, leaving behind only equipment the efreeti was wearing or carrying.

\DndMonsterAction{Innate Spellcasting}
The efreeti's innate spellcasting ability is Charisma (spell save DC 15, 7 to hit with spell attacks). It can innately cast the following spells, requiring no material components:
\begin{DndMonsterSpells}
  \DndInnateSpellLevel{\textit{detect magic}}
  \DndInnateSpellLevel[3e]{\textit{enlarge/reduce}, \textit{tongues}}
  \DndInnateSpellLevel[1e]{\textit{conjure elemental} (\textbf{fire elemental} only), \textit{gaseous form}, \textit{invisibility}, \textit{major image}, \textit{plane shift}, \textit{wall of fire}}
\end{DndMonsterSpells}
\DndMonsterSection{Actions}
\DndMonsterAction{Multiattack}
The efreeti makes two scimitar attacks or uses its Hurl Flame twice.

\DndMonsterAction{Scimitar}
\textsl{Melee Weapon Attack:} 10 to hit, reach 5 ft., one target. \textsl{Hit:} 13 (2d6 + 6) slashing damage plus 7 (2d6) fire damage.

\DndMonsterAction{Hurl Flame}
\textsl{Ranged Spell Attack:} 7 to hit, range 120 ft., one target. \textsl{Hit:} 17 (5d6) fire damage.

\end{DndMonster}



\begin{DndMonster}[float*=b,width=\textwidth + 8pt]{Eldrick Runeweaver}
\begin{multicols}{2}
\DndMonsterType{Medium humanoid (human), lawful neutral}
\DndMonsterBasics[
armor-class = {17 (\textit{robe of the archmagi})},
hit-points = {\DndDice{20d8 + 60}},
speed = {30 ft.}
]
\DndMonsterAbilityScores[
 str = 8,
 dex = 10,
 con = 16,
 int = 22,
 wis = 19,
 cha = 15
]
\DndMonsterDetails[
saving-throws = {Str +5, Dex +6, Con +9, Int +12, Wis +10, Cha +8},
skills = {Arcana +16, History +16, Perception +16, Persuasion +14},
damage-resistances = {damage from spells},
senses = {truesight 120 ft., passive perception 26},
languages = {all},
challenge = 17,
]
\DndMonsterAction{Special Equipment}
Eldrick Runeweaver wears a \textit{robe of the archmagi}, carries a \textit{staff of power}, and wears a \textit{ring of spell storing}. He carries 2d6 very rare spell scrolls and potions.

\DndMonsterAction{Legendary Resistance (3/Day)}
If Eldrick Runeweaver fails a saving throw, he can choose to succeed instead.

\DndMonsterAction{Magic Resistance}
Eldrick Runeweaver has advantage on saving throws against spells and other magical effects.

\DndMonsterAction{Spellcasting}
Eldrick Runeweaver is a 20th-level spellcaster. His spellcasting ability is Intelligence (spell save DC 19, 11 to hit with spell attacks). He can cast shield and invisibility at will and has the following wizard spells prepared:
\begin{DndMonsterSpells}
  \DndMonsterSpellLevel{\textit{fire bolt}, \textit{mage hand}, \textit{minor illusion}, \textit{prestidigitation}, \textit{ray of frost}}
   \DndMonsterSpellLevel[1][4]{\textit{mage armor}, \textit{magic missile}, \textit{thunderwave}}
   \DndMonsterSpellLevel[2][3]{\textit{detect thoughts}, \textit{levitate}, \textit{shatter}}
   \DndMonsterSpellLevel[3][3]{\textit{counterspell}, \textit{dispel magic}, \textit{lightning bolt}}
   \DndMonsterSpellLevel[4][3]{\textit{banishment}, \textit{dimension door}, \textit{ice storm}}
   \DndMonsterSpellLevel[5][3]{\textit{Bigby's hand}, \textit{scrying}, \textit{wall of force}}
   \DndMonsterSpellLevel[6][2]{\textit{chain lightning}, \textit{globe of invulnerability}}
   \DndMonsterSpellLevel[7][2]{\textit{forcecage}, \textit{teleport}}
   \DndMonsterSpellLevel[8][1]{\textit{maze}}
\end{DndMonsterSpells}
\DndMonsterSection{Actions}
\DndMonsterAction{Staff of Power}
\textsl{Melee Weapon Attack:} 6 to hit, reach 5 ft., one creature. \textsl{Hit:} 10 (1d6 + 1) bludgeoning damage.

\DndMonsterSection{Legendary Actions}
Eldrick Runeweaver can take 3 legendary actions, choosing from the options below. Only one legendary action can be used at a time and only at the end of another creature's turn. Eldrick Runeweaver regains spent legendary actions at the start of its turn.

\begin{DndMonsterLegendaryActions}
\DndMonsterLegendaryAction{Cantrip}{Eldrick Runeweaver casts a cantrip.
}
\DndMonsterLegendaryAction{Teleport}{Eldrick Runeweaver teleports 60 feet to an unoccupied space he can see.
}
\DndMonsterLegendaryAction{Cast a Spell (Costs 3 actions)}{Eldrick Runeweaver casts a spell using a spell slot.
}
\end{DndMonsterLegendaryActions}
\end{multicols}
\end{DndMonster}



\begin{DndMonster}[float=t]{Entropic Flame}
\DndMonsterType{Large elemental, unaligned}
\DndMonsterBasics[
armor-class = {13},
hit-points = {\DndDice{12d10 + 12}},
speed = {30 ft.}
]
\DndMonsterAbilityScores[
 str = 10,
 dex = 17,
 con = 13,
 int = 5,
 wis = 10,
 cha = 5
]
\DndMonsterDetails[
damage-resistances = {bludgeoning, piercing, slashing from nonmagical attacks},
damage-immunities = {fire, poison},
condition-immunities = {exhaustion, grappled, paralyzed, petrified, poisoned, prone, restrained, unconscious},
senses = {darkvision 60 ft., passive perception 10},
languages = {Deep Speech},
challenge = 5,
]
\DndMonsterAction{Fire Form}
The entropic flame can move through a space as narrow as 1 inch wide without squeezing. A creature that touches the flame or hits it with a melee attack while within 5 feet of it takes 5 (1d10) fire damage. In addition, the flame can enter a hostile creature's space and stop there. The first time it enters a creature's space on a turn, that creature takes 5 (1d10) fire damage and catches fire; until someone takes an action to douse the fire, the creature takes 5 (1d10) fire damage at the start of each of its turns.

\DndMonsterAction{Illumination}
The entropic flame sheds bright light in a 30-foot radius and dim light for an additional 30 feet.

\DndMonsterSection{Actions}
\DndMonsterAction{Multiattack}
The entropic flame makes three hurl flame attacks. Before or after making each attack, the entropic flame can teleport up to 10 feet to a space it can see. It may teleport into a space occupied by another creature.

\DndMonsterAction{Hurl Flame}
\textsl{Ranged Weapon Attack:} 6 to hit, range 120 ft., one target. \textsl{Hit:} 10 (3d6) fire damage.

\DndMonsterAction{Touch}
\textsl{Melee Weapon Attack:} 6 to hit, reach 5 ft., one target. \textsl{Hit:} 10 (2d6 + 3) fire damage. If the target is a creature or a flammable object, it ignites. Until a creature takes an action to douse the fire, the target takes 5 (1d10) fire damage at the start of each of its turns.

\end{DndMonster}



\begin{DndMonster}[float=t]{Eoghan Ghostweaver}
\DndMonsterType{Medium humanoid (elf), any alignment}
\DndMonsterBasics[
armor-class = {12 (15 with \textit{mage armor})},
hit-points = {\DndDice{18d8 + 18}},
speed = {30 ft.}
]
\DndMonsterAbilityScores[
 str = 10,
 dex = 14,
 con = 12,
 int = 15,
 wis = 20,
 cha = 16
]
\DndMonsterDetails[
saving-throws = {Int +6, Wis +9},
skills = {Arcana +13, History +13, Medicine +9, Nature +6, Perception +9, Survival +9},
damage-resistances = {damage from spells; bludgeoning, piercing, slashing (from stoneskin)},
languages = {Druidic, Elvish},
challenge = 12,
]
\DndMonsterAction{Magic Resistance}
Eoghan has advantage on saving throws against spells and other magical effects.

\DndMonsterAction{Spellcasting}
Eoghan is an 18th-level spellcaster. His spellcasting ability is Wisdom (spell save DC 17, 9 to hit with spell attacks). He has the following druid spells prepared:
\begin{DndMonsterSpells}
  \DndMonsterSpellLevel{\textit{druidcraft}, \textit{mending}, \textit{poison spray}, \textit{produce flame}}
   \DndMonsterSpellLevel[1][4]{\textit{cure wounds}, \textit{entangle}, \textit{faerie fire}, \textit{speak with animals}}
   \DndMonsterSpellLevel[2][3]{\textit{animal messenger}, \textit{beast sense}, \textit{hold person}}
   \DndMonsterSpellLevel[3][3]{\textit{conjure animals}, \textit{purge contamination}*, \textit{sleet storm}}
   \DndMonsterSpellLevel[4][3]{\textit{ice storm}, \textit{polymorph}, \textit{wall of fire}}
   \DndMonsterSpellLevel[5][3]{\textit{commune with nature}, \textit{cone of cold}, \textit{greater restoration}}
   \DndMonsterSpellLevel[6][1]{\textit{conjure fey}, \textit{heal}}
   \DndMonsterSpellLevel[7][1]{\textit{resurrection}}
   \DndMonsterSpellLevel[8][1]{\textit{feeblemind}}
\end{DndMonsterSpells}
\DndMonsterSection{Actions}
\DndMonsterAction{Dagger}
\textsl{Melee or Ranged Weapon Attack:} 6 to hit, reach 5 ft. or range 20/60 ft., one target. \textsl{Hit:} 4 (1d4 + 2) piercing damage.

\end{DndMonster}



\begin{DndMonster}[float=t]{Erinyes}
\DndMonsterType{Medium fiend (devil), lawful evil}
\DndMonsterBasics[
armor-class = {18 (\textit{plate armor})},
hit-points = {\DndDice{18d8 + 72}},
speed = {30 ft., fly 60 ft.}
]
\DndMonsterAbilityScores[
 str = 18,
 dex = 16,
 con = 18,
 int = 14,
 wis = 14,
 cha = 18
]
\DndMonsterDetails[
saving-throws = {Dex +7, Con +8, Wis +6, Cha +8},
damage-resistances = {cold; bludgeoning, piercing, slashing from nonmagical attacks that aren't silvered},
damage-immunities = {fire, poison},
condition-immunities = {poisoned},
senses = {truesight 120 ft., passive perception 12},
languages = {Infernal, telepathy 120 ft.},
challenge = 12,
]
\DndMonsterAction{Hellish Weapons}
The erinyes's weapon attacks are magical and deal an extra 13 (3d8) poison damage on a hit (included in the attacks).

\DndMonsterAction{Magic Resistance}
The erinyes has advantage on saving throws against spells and other magical effects.

\DndMonsterSection{Actions}
\DndMonsterAction{Multiattack}
The erinyes makes three attacks.

\DndMonsterAction{Longsword}
\textsl{Melee Weapon Attack:} 8 to hit, reach 5 ft., one target. \textsl{Hit:} 8 (1d8 + 4) slashing damage, or 9 (1d10 + 4) slashing damage if used with two hands, plus 13 (3d8) poison damage.

\DndMonsterAction{Longbow}
\textsl{Ranged Weapon Attack:} 7 to hit, range 150/600 ft., one target. \textsl{Hit:} 7 (1d8 + 3) piercing damage plus 13 (3d8) poison damage, and the target must succeed on a DC 14 Constitution saving throw or be poisoned. The poison lasts until it is removed by the \textit{lesser restoration} spell or similar magic.

\DndMonsterSection{Reactions}
\DndMonsterAction{Parry}
The erinyes adds 4 to its AC against one melee attack that would hit it. To do so, the erinyes must see the attacker and be wielding a melee weapon.

\end{DndMonster}



\begin{DndMonster}[float*=b,width=\textwidth + 8pt]{Executioner}
\begin{multicols}{2}
\DndMonsterType{Huge construct, lawful evil}
\DndMonsterBasics[
armor-class = {20 (natural armor)},
hit-points = {\DndDice{30d12 + 210}},
speed = {40 ft.}
]
\DndMonsterAbilityScores[
 str = 29,
 dex = 11,
 con = 25,
 int = 10,
 wis = 11,
 cha = 10
]
\DndMonsterDetails[
saving-throws = {Str +15, Con +13, Wis +6},
damage-immunities = {poison; psychic; bludgeoning, piercing, slashing from nonmagical attacks that aren't adamantine},
condition-immunities = {charmed, exhaustion, frightened, paralyzed, petrified, poisoned},
senses = {darkvision 120 ft., passive perception 10},
languages = {understands common but does not speak},
challenge = 20,
]
\DndMonsterAction{Legendary Resistance (3/Day)}
If the Executioner fails a saving throw, it can choose to succeed instead.

\DndMonsterAction{Immutable Form}
The Executioner is immune to any spell or effect that would alter its form.

\DndMonsterAction{Magic Resistance}
The Executioner has advantage on saving throws against spells and other magical effects.

\DndMonsterAction{Magic Weapons}
The Executioner's weapon attacks are magical.

\DndMonsterAction{Rejuvenation}
If the Executioner is destroyed, it returns with all its hit points in 24 hours unless a \textit{disintegrate} spell is cast on its remains.

\DndMonsterSection{Actions}
\DndMonsterAction{Guillotine Blade}
\textsl{Melee or Ranged Weapon Attack:} 15 to hit, reach 10 ft., or range 30/120 ft., one target. \textsl{Hit:} 30 (6d6 + 9) slashing damage plus 14 (4d6) necrotic damage. The target must succeed on a DC 21 Constitution saving throw or its hit point maximum is reduced by an amount equal to the necrotic damage taken. This reduction lasts until the creature finishes a long rest. The target dies if this effect reduces its hit point maximum to 0. Immediately after making a ranged attack, this weapon flies back to the Executioner's hands.

\DndMonsterAction{Necrotic Vents (Recharge 5\textendash 6)}
The Executioner vents necrotic gas in a 30-foot radius. Each creature in that area must succeed on a DC 21 Constitution saving throw, taking 45 (10d8) necrotic damage on a failed save, or half as much damage on a successful one. A creature who fails this saving throw by 5 or more gains one level of contamination.

\DndMonsterSection{Legendary Actions}
Executioner can take 3 legendary actions, choosing from the options below. Only one legendary action can be used at a time and only at the end of another creature's turn. Executioner regains spent legendary actions at the start of its turn.

\begin{DndMonsterLegendaryActions}
\DndMonsterLegendaryAction{Lumber Forward}{The Executioner moves half its speed.
}
\DndMonsterLegendaryAction{Attack (Costs 2 actions)}{The Executioner makes a melee or ranged attack with its Guillotine Blade.
}
\DndMonsterLegendaryAction{Summon Shadows. (Costs 3 actions)}{The Executioner targets up to six humanoid corpses it can see within 60 feet of it that have died violently. Each target's spirit rises as a shadow in the space of its corpse or in the nearest unoccupied space. The shadows act after the Executioner on the same initiative count, are under its control, and remain for 1 minute or until they're destroyed. They disappear if the Executioner is destroyed.
}
\end{DndMonsterLegendaryActions}
\end{multicols}
\end{DndMonster}



\begin{DndMonster}[float=t]{Fire Elemental}
\DndMonsterType{Large elemental, neutral}
\DndMonsterBasics[
armor-class = {13},
hit-points = {\DndDice{12d10 + 36}},
speed = {50 ft.}
]
\DndMonsterAbilityScores[
 str = 10,
 dex = 17,
 con = 16,
 int = 6,
 wis = 10,
 cha = 7
]
\DndMonsterDetails[
damage-resistances = {bludgeoning, piercing, slashing from nonmagical attacks},
damage-immunities = {fire, poison},
condition-immunities = {exhaustion, grappled, paralyzed, petrified, poisoned, prone, restrained, unconscious},
senses = {darkvision 60 ft., passive perception 10},
languages = {Ignan},
challenge = 5,
]
\DndMonsterAction{Fire Form}
The elemental can move through a space as narrow as 1 inch wide without squeezing. A creature that touches the elemental or hits it with a melee attack while within 5 feet of it takes 5 (1d10) fire damage. In addition, the elemental can enter a hostile creature's space and stop there. The first time it enters a creature's space on a turn, that creature takes 5 (1d10) fire damage and catches fire; until someone takes an action to douse the fire, the creature takes 5 (1d10) fire damage at the start of each of its turns.

\DndMonsterAction{Illumination}
The elemental sheds bright light in a 30-foot radius and dim light in an additional 30 feet.

\DndMonsterAction{Water Susceptibility}
For every 5 feet the elemental moves in water, or for every gallon of water splashed on it, it takes 1 cold damage.

\DndMonsterSection{Actions}
\DndMonsterAction{Multiattack}
The elemental makes two touch attacks.

\DndMonsterAction{Touch}
\textsl{Melee Weapon Attack:} 6 to hit, reach 5 ft., one target. \textsl{Hit:} 10 (2d6 + 3) fire damage. If the target is a creature or a flammable object, it ignites. Until a creature takes an action to douse the fire, the target takes 5 (1d10) fire damage at the start of each of its turns.

\end{DndMonster}



\begin{DndMonster}[float=t]{Flying Sword}
\DndMonsterType{Small construct, unaligned}
\DndMonsterBasics[
armor-class = {17 (natural armor)},
hit-points = {\DndDice{5d6}},
speed = {0 ft., fly 50 ft. \textit{(hover)}}
]
\DndMonsterAbilityScores[
 str = 12,
 dex = 15,
 con = 11,
 int = 1,
 wis = 5,
 cha = 1
]
\DndMonsterDetails[
saving-throws = {Dex +4},
damage-immunities = {poison, psychic},
condition-immunities = {blinded, charmed, deafened, frightened, paralyzed, petrified, poisoned},
senses = {blindsight 60 ft. (blind beyond this radius), passive perception 7},
challenge = 1/4,
]
\DndMonsterAction{Antimagic Susceptibility}
The sword is incapacitated while in the area of an \textit{antimagic field}. If targeted by \textit{dispel magic}, the sword must succeed on a Constitution saving throw against the caster's spell save DC or fall unconscious for 1 minute.

\DndMonsterAction{False Appearance}
While the sword remains motionless and isn't flying, it is indistinguishable from a normal sword.

\DndMonsterSection{Actions}
\DndMonsterAction{Longsword}
\textsl{Melee Weapon Attack:} 3 to hit, reach 5 ft., one target. \textsl{Hit:} 5 (1d8 + 1) slashing damage.

\end{DndMonster}



\begin{DndMonster}[float=t]{Gargoyle}
\DndMonsterType{Medium elemental, chaotic evil}
\DndMonsterBasics[
armor-class = {15 (natural armor)},
hit-points = {\DndDice{7d8 + 21}},
speed = {30 ft., fly 60 ft.}
]
\DndMonsterAbilityScores[
 str = 15,
 dex = 11,
 con = 16,
 int = 6,
 wis = 11,
 cha = 7
]
\DndMonsterDetails[
damage-resistances = {bludgeoning, piercing, slashing from nonmagical attacks that aren't adamantine},
damage-immunities = {poison},
condition-immunities = {exhaustion, petrified, poisoned},
senses = {darkvision 60 ft., passive perception 10},
languages = {Terran},
challenge = 2,
]
\DndMonsterAction{False Appearance}
While the gargoyle remains motionless, it is indistinguishable from an inanimate statue.

\DndMonsterSection{Actions}
\DndMonsterAction{Multiattack}
The gargoyle makes two attacks: one with its bite and one with its claws.

\DndMonsterAction{Bite}
\textsl{Melee Weapon Attack:} 4 to hit, reach 5 ft., one target. \textsl{Hit:} 5 (1d6 + 2) piercing damage.

\DndMonsterAction{Claws}
\textsl{Melee Weapon Attack:} 4 to hit, reach 5 ft., one target. \textsl{Hit:} 5 (1d6 + 2) slashing damage.

\end{DndMonster}



\begin{DndMonster}[float=t]{Garmyr}
\DndMonsterType{Medium humanoid (gnoll), chaotic evil}
\DndMonsterBasics[
armor-class = {15 (\textit{hide armor})},
hit-points = {\DndDice{5d8}},
speed = {30 ft.}
]
\DndMonsterAbilityScores[
 str = 14,
 dex = 12,
 con = 11,
 int = 6,
 wis = 10,
 cha = 7
]
\DndMonsterDetails[
senses = {darkvision 60 ft., passive perception 10},
languages = {Gnoll},
challenge = 1/2,
]
\DndMonsterAction{Copy Warning}
This content has been copied from another creature, but the data replacement hasn't been run. Check details.

\DndMonsterAction{Rampage}
When the gnoll reduces a creature to 0 hit points with a melee attack on its turn, the gnoll can take a bonus action to move up to half its speed and make a bite attack.

\DndMonsterSection{Actions}
\DndMonsterAction{Bite}
\textsl{Melee Weapon Attack:} 4 to hit, reach 5 ft., one creature. \textsl{Hit:} 4 (1d4 + 2) piercing damage.

\DndMonsterAction{Spear}
\textsl{Melee or Ranged Weapon Attack:} 4 to hit, reach 5 ft. or range 20/60 ft., one target. \textsl{Hit:} 5 (1d6 + 2) piercing damage, or 6 (1d8 + 2) piercing damage if used with two hands to make a melee attack.

\DndMonsterAction{Longbow}
\textsl{Ranged Weapon Attack:} 3 to hit, range 150/600 ft., one target. \textsl{Hit:} 5 (1d8 + 1) piercing damage.

\end{DndMonster}



\begin{DndMonster}[float=t]{Gelatinous Cube}
\DndMonsterType{Large ooze, unaligned}
\DndMonsterBasics[
armor-class = {6},
hit-points = {\DndDice{8d10 + 40}},
speed = {15 ft.}
]
\DndMonsterAbilityScores[
 str = 14,
 dex = 3,
 con = 20,
 int = 1,
 wis = 6,
 cha = 1
]
\DndMonsterDetails[
condition-immunities = {blinded, charmed, deafened, exhaustion, frightened, prone},
senses = {blindsight 60 ft. (blind beyond this radius), passive perception 8},
challenge = 2,
]
\DndMonsterAction{Ooze Cube}
The cube takes up its entire space. Other creatures can enter the space, but a creature that does so is subjected to the cube's Engulf and has disadvantage on the saving throw.

Creatures inside the cube can be seen but have total cover.

A creature within 5 feet of the cube can take an action to pull a creature or object out of the cube. Doing so requires a successful DC 12 Strength check, and the creature making the attempt takes 10 (3d6) acid damage.

The cube can hold only one Large creature or up to four Medium or smaller creatures inside it at a time.

\DndMonsterAction{Transparent}
Even when the cube is in plain sight, it takes a successful DC 15 Wisdom (Perception) check to spot a cube that has neither moved nor attacked. A creature that tries to enter the cube's space while unaware of the cube is Surprise by the cube.

\DndMonsterSection{Actions}
\DndMonsterAction{Pseudopod}
\textsl{Melee Weapon Attack:} 4 to hit, reach 5 ft., one creature. \textsl{Hit:} 10 (3d6) acid damage.

\DndMonsterAction{Engulf}
The cube moves up to its speed. While doing so, it can enter Large or smaller creatures' spaces. Whenever the cube enters a creature's space, the creature must make a DC 12 Dexterity saving throw.

On a successful save, the creature can choose to be pushed 5 feet back or to the side of the cube. A creature that chooses not to be pushed suffers the consequences of a failed saving throw.

On a failed save, the cube enters the creature's space, and the creature takes 10 (3d6) acid damage and is engulfed. The engulfed creature can't breathe, is restrained, and takes 21 (6d6) acid damage at the start of each of the cube's turns. When the cube moves, the engulfed creature moves with it.

An engulfed creature can try to escape by taking an action to make a DC 12 Strength check. On a success, the creature escapes and enters a space of its choice within 5 feet of the cube.

\end{DndMonster}



\begin{DndMonster}[float=t]{Giant Crocodile}
\DndMonsterType{Huge beast, unaligned}
\DndMonsterBasics[
armor-class = {14 (natural armor)},
hit-points = {\DndDice{9d12 + 27}},
speed = {30 ft., swim 50 ft.}
]
\DndMonsterAbilityScores[
 str = 21,
 dex = 9,
 con = 17,
 int = 2,
 wis = 10,
 cha = 7
]
\DndMonsterDetails[
skills = {Stealth +5},
challenge = 5,
]
\DndMonsterAction{Hold Breath}
The crocodile can hold its breath for 30 minutes.

\DndMonsterSection{Actions}
\DndMonsterAction{Multiattack}
The crocodile makes two attacks: one with its bite and one with its tail.

\DndMonsterAction{Bite}
\textsl{Melee Weapon Attack:} 8 to hit, reach 5 ft., one target. \textsl{Hit:} 21 (3d10 + 5) piercing damage, and the target is grappled (escape DC 16). Until this grapple ends, the target is restrained, and the crocodile can't bite another target.

\DndMonsterAction{Tail}
\textsl{Melee Weapon Attack:} 8 to hit, reach 10 ft., one target not grappled by the crocodile. \textsl{Hit:} 14 (2d8 + 5) bludgeoning damage. If the target is a creature, it must succeed on a DC 16 Strength saving throw or be knocked prone.

\end{DndMonster}



\begin{DndMonster}[float=t]{Gibbering Mouther}
\DndMonsterType{Medium aberration, neutral}
\DndMonsterBasics[
armor-class = {9},
hit-points = {\DndDice{9d8 + 27}},
speed = {10 ft., swim 10 ft.}
]
\DndMonsterAbilityScores[
 str = 10,
 dex = 8,
 con = 16,
 int = 3,
 wis = 10,
 cha = 6
]
\DndMonsterDetails[
condition-immunities = {prone},
senses = {darkvision 60 ft., passive perception 10},
challenge = 2,
]
\DndMonsterAction{Aberrant Ground}
The ground in a 10-foot radius around the mouther is doughlike difficult terrain. Each creature that starts its turn in that area must succeed on a DC 10 Strength saving throw or have its speed reduced to 0 until the start of its next turn.

\DndMonsterAction{Gibbering}
The mouther babbles incoherently while it can see any creature and isn't incapacitated. Each creature that starts its turn within 20 feet of the mouther and can hear the gibbering must succeed on a DC 10 Wisdom saving throw. On a failure, the creature can't take reactions until the start of its next turn and rolls a d8 to determine what it does during its turn. On a 1 to 4, the creature does nothing. On a 5 or 6, the creature takes no action or bonus action and uses all its movement to move in a randomly determined direction. On a 7 or 8, the creature makes a melee attack against a randomly determined creature within its reach or does nothing if it can't make such an attack.

\DndMonsterSection{Actions}
\DndMonsterAction{Multiattack}
The gibbering mouther makes one bite attack and, if it can, uses its Blinding Spittle.

\DndMonsterAction{Bites}
\textsl{Melee Weapon Attack:} 2 to hit, reach 5 ft., one creature. \textsl{Hit:} 17 (5d6) piercing damage. If the target is Medium or smaller, it must succeed on a DC 10 Strength saving throw or be knocked prone. If the target is killed by this damage, it is absorbed into the mouther.

\DndMonsterAction{Blinding Spittle (Recharge 5\textendash 6)}
The mouther spits a chemical glob at a point it can see within 15 feet of it. The glob explodes in a blinding flash of light on impact. Each creature within 5 feet of the flash must succeed on a DC 13 Dexterity saving throw or be blinded until the end of the mouther's next turn.

\end{DndMonster}



\begin{DndMonster}[float=t]{Glabrezu}
\DndMonsterType{Large fiend (demon), chaotic evil}
\DndMonsterBasics[
armor-class = {17 (natural armor)},
hit-points = {\DndDice{15d10 + 75}},
speed = {40 ft.}
]
\DndMonsterAbilityScores[
 str = 20,
 dex = 15,
 con = 21,
 int = 19,
 wis = 17,
 cha = 16
]
\DndMonsterDetails[
saving-throws = {Str +9, Con +9, Wis +7, Cha +7},
damage-resistances = {cold; fire; lightning; bludgeoning, piercing, slashing from nonmagical attacks},
damage-immunities = {poison},
condition-immunities = {poisoned},
senses = {truesight 120 ft., passive perception 13},
languages = {Abyssal, telepathy 120 ft.},
challenge = 9,
]
\DndMonsterAction{Magic Resistance}
The glabrezu has advantage on saving throws against spells and other magical effects.

\DndMonsterAction{Innate Spellcasting}
The glabrezu's spellcasting ability is Intelligence (spell save DC 16). The glabrezu can innately cast the following spells, requiring no material components:
\begin{DndMonsterSpells}
  \DndInnateSpellLevel{\textit{darkness}, \textit{detect magic}, \textit{dispel magic}}
  \DndInnateSpellLevel[1e]{\textit{confusion}, \textit{fly}, \textit{power word stun}}
\end{DndMonsterSpells}
\DndMonsterSection{Actions}
\DndMonsterAction{Multiattack}
The glabrezu makes four attacks: two with its pincers and two with its fists. Alternatively, it makes two attacks with its pincers and casts one spell.

\DndMonsterAction{Pincer}
\textsl{Melee Weapon Attack:} 9 to hit, reach 10 ft., one target. \textsl{Hit:} 16 (2d10 + 5) bludgeoning damage. If the target is a Medium or smaller creature, it is grappled (escape DC 15). The glabrezu has two pincers, each of which can grapple only one target.

\DndMonsterAction{Fist}
\textsl{Melee Weapon Attack:} 9 to hit, reach 5 ft., one target. \textsl{Hit:} 7 (2d4 + 2) bludgeoning damage.

\end{DndMonster}



\begin{DndMonster}[float=t]{Gladiator}
\DndMonsterType{Medium humanoid (any race), any alignment}
\DndMonsterBasics[
armor-class = {16 (studded leather armor)},
hit-points = {\DndDice{15d8 + 45}},
speed = {30 ft.}
]
\DndMonsterAbilityScores[
 str = 18,
 dex = 15,
 con = 16,
 int = 10,
 wis = 12,
 cha = 15
]
\DndMonsterDetails[
saving-throws = {Str +7, Dex +5, Con +6},
skills = {Athletics +10, Intimidation +5},
languages = {any one language (usually Common)},
challenge = 5,
]
\DndMonsterAction{Brave}
The gladiator has advantage on saving throws against being frightened.

\DndMonsterAction{Brute}
A melee weapon deals one extra die of its damage when the gladiator hits with it (included in the attack).

\DndMonsterSection{Actions}
\DndMonsterAction{Multiattack}
The gladiator makes three melee attacks or two ranged attacks.

\DndMonsterAction{Spear}
\textsl{Melee or Ranged Weapon Attack:} 7 to hit, reach 5 ft. and range 20/60 ft., one target. \textsl{Hit:} 11 (2d6 + 4) piercing damage, or 13 (2d8 + 4) piercing damage if used with two hands to make a melee attack.

\DndMonsterAction{Shield Bash}
\textsl{Melee Weapon Attack:} 7 to hit, reach 5 ft., one creature. \textsl{Hit:} 9 (2d4 + 4) bludgeoning damage. If the target is a Medium or smaller creature, it must succeed on a DC 15 Strength saving throw or be knocked prone.

\DndMonsterSection{Reactions}
\DndMonsterAction{Parry}
The gladiator adds 3 to its AC against one melee attack that would hit it. To do so, the gladiator must see the attacker and be wielding a melee weapon.

\end{DndMonster}



\begin{DndMonster}[float=t]{Gnoll}
\DndMonsterType{Medium humanoid (gnoll), chaotic evil}
\DndMonsterBasics[
armor-class = {15 (\textit{hide armor})},
hit-points = {\DndDice{5d8}},
speed = {30 ft.}
]
\DndMonsterAbilityScores[
 str = 14,
 dex = 12,
 con = 11,
 int = 6,
 wis = 10,
 cha = 7
]
\DndMonsterDetails[
senses = {darkvision 60 ft., passive perception 10},
languages = {Gnoll},
challenge = 1/2,
]
\DndMonsterAction{Rampage}
When the gnoll reduces a creature to 0 hit points with a melee attack on its turn, the gnoll can take a bonus action to move up to half its speed and make a bite attack.

\DndMonsterSection{Actions}
\DndMonsterAction{Bite}
\textsl{Melee Weapon Attack:} 4 to hit, reach 5 ft., one creature. \textsl{Hit:} 4 (1d4 + 2) piercing damage.

\DndMonsterAction{Spear}
\textsl{Melee or Ranged Weapon Attack:} 4 to hit, reach 5 ft. or range 20/60 ft., one target. \textsl{Hit:} 5 (1d6 + 2) piercing damage, or 6 (1d8 + 2) piercing damage if used with two hands to make a melee attack.

\DndMonsterAction{Longbow}
\textsl{Ranged Weapon Attack:} 3 to hit, range 150/600 ft., one target. \textsl{Hit:} 5 (1d8 + 1) piercing damage.

\end{DndMonster}



\begin{DndMonster}[float=t]{Gorgon}
\DndMonsterType{Large monstrosity, unaligned}
\DndMonsterBasics[
armor-class = {19 (natural armor)},
hit-points = {\DndDice{12d10 + 48}},
speed = {40 ft.}
]
\DndMonsterAbilityScores[
 str = 20,
 dex = 11,
 con = 18,
 int = 2,
 wis = 12,
 cha = 7
]
\DndMonsterDetails[
skills = {Perception +4},
condition-immunities = {petrified},
senses = {darkvision 60 ft., passive perception 14},
challenge = 5,
]
\DndMonsterAction{Trampling Charge}
If the gorgon moves at least 20 feet straight toward a creature and then hits it with a gore attack on the same turn, that target must succeed on a DC 16 Strength saving throw or be knocked prone. If the target is prone, the gorgon can make one attack with its hooves against it as a bonus action.

\DndMonsterSection{Actions}
\DndMonsterAction{Gore}
\textsl{Melee Weapon Attack:} 8 to hit, reach 5 ft., one target. \textsl{Hit:} 18 (2d12 + 5) piercing damage.

\DndMonsterAction{Hooves}
\textsl{Melee Weapon Attack:} 8 to hit, reach 5 ft., one target. \textsl{Hit:} 16 (2d10 + 5) bludgeoning damage.

\DndMonsterAction{Petrifying Breath (Recharge 5\textendash 6)}
The gorgon exhales petrifying gas in a 30-foot cone. Each creature in that area must succeed on a DC 13 Constitution saving throw. On a failed save, a target begins to turn to stone and is restrained. The restrained target must repeat the saving throw at the end of its next turn. On a success, the effect ends on the target. On a failure, the target is petrified until freed by the  \textit{greater restoration} spell or other magic.

\end{DndMonster}


\clearpage

\begin{DndMonster}[float*=b,width=\textwidth + 8pt]{Gravekeeper}
\begin{multicols}{2}
\DndMonsterType{Large aberration, chaotic evil}
\DndMonsterBasics[
armor-class = {18 (natural armor)},
hit-points = {\DndDice{20d10 + 100}},
speed = {0 ft., fly 30 ft. \textit{(hover)}}
]
\DndMonsterAbilityScores[
 str = 18,
 dex = 7,
 con = 21,
 int = 19,
 wis = 17,
 cha = 21
]
\DndMonsterDetails[
saving-throws = {Con +10, Wis +8, Cha +10},
skills = {Perception +13},
senses = {truesight 120 ft., passive perception 23},
languages = {all but does not speak},
challenge = 13,
]
\DndMonsterAction{Dazzling Illumination}
The Gravekeeper projects a brilliant corona of light which sheds bright light to a range of 60 feet, and causes attack rolls against it to have disadvantage. If it is hit by an attack, this trait is disrupted until the end of its next turn. This trait is also disrupted while the Gravekeeper is incapacitated or has a speed of 0.

\DndMonsterAction{Legendary Resistance (3/Day)}
If the Gravekeeper fails a saving throw, it can choose to succeed instead.

\DndMonsterSection{Actions}
\DndMonsterAction{Slam}
\textsl{Melee Weapon Attack:} 9 to hit, reach 10 ft., one target. \textsl{Hit:} 22 (4d8 + 4) bludgeoning damage plus 10 (3d6) necrotic damage.

\DndMonsterAction{Lantern Blast}
Magical octarine beams flash from the Gravekeeper's many lanterns. Each beam has a different power. Each creature within 120 feet of the Gravekeeper which is not behind total cover rolls a d6 to determine which colour ray affects it. The DC for all saving throws is 18:


\begin{description}
\item[1. Dispelling Blast.]
The targeted creature must succeed on an Intelligence saving throw or take 35 (10d6) psychic damage, and any spell or magical effect upon the target ends. A creature who fails this saving throw has their concentration broken on any spell they were actively maintaining.

\item[2. Forceful Blast.]
The targeted creature must succeed on a Strength saving throw or take 35 (10d6) force damage, be pushed 20 feet away, and knocked prone.

\item[3. Blinding Blast.]
The targeted creature must succeed on a Constitution saving throw or take 35 (10d6) radiant damage and be blinded for 1 minute. A blinded creature can repeat the saving throw on each of its turns, ending the effect on itself on a success.

\item[4. Contamining Blast.]
The targeted creature must succeed on a Constitution saving throw or take 35 (10d6) necrotic damage and gain one level of contamination.

\item[5. Igniting Blast.]
The targeted creature must succeed on a Dexterity saving throw or take 35 (10d6) fire damage, and catches fire; until someone takes an action to douse the fire, the target takes 35 (10d6) fire damage at the start of each of its turns.

\item[6. Freezing Blast.]
The targeted creature must succeed on a Constitution saving throw or take 35 (10d6) cold damage, and it can't take a reaction until the end of its next turn. Moreover, on its next turn, it must choose whether it gets a move, an action, or a bonus action; it gets only one of the three. On a successful save, the target takes half as much damage and suffers none of the blast's other effects.

\end{description}

\DndMonsterSection{Legendary Actions}
Gravekeeper can take 3 legendary actions, choosing from the options below. Only one legendary action can be used at a time and only at the end of another creature's turn. Gravekeeper regains spent legendary actions at the start of its turn.

\begin{DndMonsterLegendaryActions}
\DndMonsterLegendaryAction{Slam}{The Gravekeeper makes a slam attack.
}
\DndMonsterLegendaryAction{Glide}{The Gravekeeper moves its speed without provoking opportunity attacks.
}
\DndMonsterLegendaryAction{Lantern Blast (Costs 3 actions)}{The Gravekeeper uses its Lantern Blast.
}
\end{DndMonsterLegendaryActions}
\DndMonsterSection{Lair Actions}
When fighting inside its lair, the Gravekeeper can invoke the ambient magic to take lair actions. On initiative count 20 (losing initiative ties), the Gravekeeper can take one lair action to cause one of the following effects:


\begin{itemize}
\item A 10 foot radius cylinder of delerium sludge spouts up from the floor and raises 30 feet into the air before bubbling back down. Any creature caught in this spout must succeed on a DC 18 Constitution saving throw. They take 42 (12d6) necrotic damage on a failed saving throw, and half as much on a successful one. In addition, characters who fail the saving throw gain one level of contamination.
\item Walls within 120 feet of the Gravekeeper sprout grasping, fleshy arms. Each creature of the Gravekeepers choice that starts its turn within 10 feet of such a wall must succeed on a DC 15 Dexterity saving throw or be grappled. Escaping requires a successful DC 15 Strength (
\item Athletics
\item ) or Dexterity (
\item Acrobatics
\item ) check.
\item A torch or lantern hanging in the room emits one of the Gravekeeper's lantern blasts.
\end{itemize}

The Gravekeeper can't repeat an effect until they have all been used, and it can't use the same effect two rounds in a row.

\end{multicols}
\end{DndMonster}



\begin{DndMonster}[float=t]{Griffon}
\DndMonsterType{Large monstrosity, unaligned}
\DndMonsterBasics[
armor-class = {12},
hit-points = {\DndDice{7d10 + 21}},
speed = {30 ft., fly 80 ft.}
]
\DndMonsterAbilityScores[
 str = 18,
 dex = 15,
 con = 16,
 int = 2,
 wis = 13,
 cha = 8
]
\DndMonsterDetails[
skills = {Perception +5},
senses = {darkvision 60 ft., passive perception 15},
challenge = 2,
]
\DndMonsterAction{Keen Sight}
The griffon has advantage on Wisdom (Perception) checks that rely on sight.

\DndMonsterSection{Actions}
\DndMonsterAction{Multiattack}
The griffon makes two attacks: one with its beak and one with its claws.

\DndMonsterAction{Beak}
\textsl{Melee Weapon Attack:} 6 to hit, reach 5 ft., one target. \textsl{Hit:} 8 (1d8 + 4) piercing damage.

\DndMonsterAction{Claws}
\textsl{Melee Weapon Attack:} 6 to hit, reach 5 ft., one target. \textsl{Hit:} 11 (2d6 + 4) slashing damage.

\end{DndMonster}



\begin{DndMonster}[float=t]{Grotesque Gargant}
\DndMonsterType{Huge giant, neutral}
\DndMonsterBasics[
armor-class = {17 (natural armor)},
hit-points = {\DndDice{11d12 + 55}},
speed = {40 ft.}
]
\DndMonsterAbilityScores[
 str = 23,
 dex = 15,
 con = 20,
 int = 10,
 wis = 12,
 cha = 9
]
\DndMonsterDetails[
saving-throws = {Dex +5, Con +8, Wis +4},
skills = {Athletics +12, Perception +4},
senses = {darkvision 60 ft., passive perception 14},
languages = {Giant},
challenge = 7,
]
\DndMonsterAction{Copy Warning}
This content has been copied from another creature, but the data replacement hasn't been run. Check details.

\DndMonsterAction{Stone Camouflage}
The giant has advantage on Dexterity (Stealth) checks made to hide in rocky terrain.

\DndMonsterSection{Actions}
\DndMonsterAction{Multiattack}
The giant makes two greatclub attacks.

\DndMonsterAction{Greatclub}
\textsl{Melee Weapon Attack:} 9 to hit, reach 15 ft., one target. \textsl{Hit:} 19 (3d8 + 6) bludgeoning damage.

\DndMonsterAction{Rock}
\textsl{Ranged Weapon Attack:} 9 to hit, range 60/240 ft., one target. \textsl{Hit:} 28 (4d10 + 6) bludgeoning damage. If the target is a creature, it must succeed on a DC 17 Strength saving throw or be knocked prone.

\DndMonsterSection{Reactions}
\DndMonsterAction{Rock Catching}
If a rock or similar object is hurled at the giant, the giant can, with a successful DC 10 Dexterity saving throw, catch the missile and take no bludgeoning damage from it.

\end{DndMonster}



\begin{DndMonster}[float=t]{Guard}
\DndMonsterType{Medium humanoid (any race), any alignment}
\DndMonsterBasics[
armor-class = {16 (\textit{chain shirt})},
hit-points = {\DndDice{2d8 + 2}},
speed = {30 ft.}
]
\DndMonsterAbilityScores[
 str = 13,
 dex = 12,
 con = 12,
 int = 10,
 wis = 11,
 cha = 10
]
\DndMonsterDetails[
skills = {Perception +2},
languages = {any one language (usually Common)},
challenge = 1/8,
]
\DndMonsterSection{Actions}
\DndMonsterAction{Spear}
\textsl{Melee or Ranged Weapon Attack:} 3 to hit, reach 5 ft. or range 20/60 ft., one target. \textsl{Hit:} 4 (1d6 + 1) piercing damage, or 5 (1d8 + 1) piercing damage if used with two hands to make a melee attack.

\end{DndMonster}



\begin{DndMonster}[float=t]{Harpy}
\DndMonsterType{Medium monstrosity, chaotic evil}
\DndMonsterBasics[
armor-class = {11},
hit-points = {\DndDice{7d8 + 7}},
speed = {20 ft., fly 40 ft.}
]
\DndMonsterAbilityScores[
 str = 12,
 dex = 13,
 con = 12,
 int = 7,
 wis = 10,
 cha = 13
]
\DndMonsterDetails[
languages = {Common},
challenge = 1,
]
\DndMonsterSection{Actions}
\DndMonsterAction{Multiattack}
The harpy makes two attacks: one with its claws and one with its club.

\DndMonsterAction{Claws}
\textsl{Melee Weapon Attack:} 3 to hit, reach 5 ft., one target. \textsl{Hit:} 6 (2d4 + 1) slashing damage.

\DndMonsterAction{Club}
\textsl{Melee Weapon Attack:} 3 to hit, reach 5 ft., one target. \textsl{Hit:} 3 (1d4 + 1) bludgeoning damage.

\DndMonsterAction{Luring Song}
The harpy sings a magical melody. Every humanoid and giant within 300 feet of the harpy that can hear the song must succeed on a DC 11 Wisdom saving throw or be charmed until the song ends. The harpy must take a bonus action on its subsequent turns to continue singing. It can stop singing at any time. The song ends if the harpy is incapacitated.

While charmed by the harpy, a target is incapacitated and ignores the songs of other harpies. If the charmed target is more than 5 feet away from the harpy, the target must move on its turn toward the harpy by the most direct route. It doesn't avoid opportunity attacks, but before moving into damaging terrain, such as lava or a pit, and whenever it takes damage from a source other than the harpy, a target can repeat the saving throw. A creature can also repeat the saving throw at the end of each of its turns. If a creature's saving throw is successful, the effect ends on it.

A target that successfully saves is immune to this harpy's song for the next 24 hours.

\end{DndMonster}



\begin{DndMonster}[float=t]{Haze Hulk}
\DndMonsterType{Large aberration, unaligned}
\DndMonsterBasics[
armor-class = {15 (natural armor)},
hit-points = {\DndDice{10d10 + 40}},
speed = {40 ft.}
]
\DndMonsterAbilityScores[
 str = 21,
 dex = 10,
 con = 18,
 int = 7,
 wis = 10,
 cha = 6
]
\DndMonsterDetails[
senses = {darkvision 60 ft., passive perception 10},
languages = {Common},
challenge = 4,
]
\DndMonsterSection{Actions}
\DndMonsterAction{Multiattack}
The haze hulk makes two slam attacks.

\DndMonsterAction{Slam}
\textsl{Melee Weapon Attack:} 7 to hit, reach 5 ft., one target. \textsl{Hit:} 14 (2d8 + 5) bludgeoning damage.

\end{DndMonster}



\begin{DndMonster}[float=t]{Haze Husk}
\DndMonsterType{Medium undead, neutral evil}
\DndMonsterBasics[
armor-class = {8},
hit-points = {\DndDice{3d8 + 9}},
speed = {20 ft.}
]
\DndMonsterAbilityScores[
 str = 13,
 dex = 6,
 con = 16,
 int = 3,
 wis = 6,
 cha = 5
]
\DndMonsterDetails[
saving-throws = {Wis +0},
damage-immunities = {poison},
condition-immunities = {poisoned},
senses = {darkvision 60 ft., passive perception 8},
languages = {understands all languages it spoke in life but can't speak},
challenge = 1/4,
]
\DndMonsterAction{Copy Warning}
This content has been copied from another creature, but the data replacement hasn't been run. Check details.

\DndMonsterAction{Undead Fortitude}
If damage reduces the zombie to 0 hit points, it must make a Constitution saving throw with a DC of 5 + the damage taken, unless the damage is radiant or from a critical hit. On a success, the zombie drops to 1 hit point instead.

\DndMonsterSection{Actions}
\DndMonsterAction{Slam}
\textsl{Melee Weapon Attack:} 3 to hit, reach 5 ft., one target. \textsl{Hit:} 4 (1d6 + 1) bludgeoning damage.

\end{DndMonster}



\begin{DndMonster}[float=t]{Haze Wight}
\DndMonsterType{Medium undead, neutral evil}
\DndMonsterBasics[
armor-class = {14 (studded leather armor)},
hit-points = {\DndDice{6d8 + 18}},
speed = {30 ft.}
]
\DndMonsterAbilityScores[
 str = 15,
 dex = 14,
 con = 16,
 int = 10,
 wis = 13,
 cha = 15
]
\DndMonsterDetails[
skills = {Perception +3, Stealth +4},
damage-resistances = {necrotic; bludgeoning, piercing, slashing from nonmagical attacks that aren't silvered},
damage-immunities = {poison},
condition-immunities = {exhaustion, poisoned},
senses = {darkvision 60 ft., passive perception 13},
languages = {the languages it knew in life},
challenge = 3,
]
\DndMonsterAction{Copy Warning}
This content has been copied from another creature, but the data replacement hasn't been run. Check details.

\DndMonsterAction{Sunlight Sensitivity}
While in sunlight, the wight has disadvantage on attack rolls, as well as on Wisdom (Perception) checks that rely on sight.

\DndMonsterSection{Actions}
\DndMonsterAction{Multiattack}
The wight makes two longsword attacks or two longbow attacks. It can use its Life Drain in place of one longsword attack.

\DndMonsterAction{Life Drain}
\textsl{Melee Weapon Attack:} 4 to hit, reach 5 ft., one creature. \textsl{Hit:} 5 (1d6 + 2) necrotic damage. The target must succeed on a DC 13 Constitution saving throw or its hit point maximum is reduced by an amount equal to the damage taken. This reduction lasts until the target finishes a long rest. The target dies if this effect reduces its hit point maximum to 0.

A humanoid slain by this attack rises 24 hours later as a \textbf{zombie} under the wight's control, unless the humanoid is restored to life or its body is destroyed. The wight can have no more than twelve zombies under its control at one time.

\DndMonsterAction{Longsword}
\textsl{Melee Weapon Attack:} 4 to hit, reach 5 ft., one target. \textsl{Hit:} 6 (1d8 + 2) slashing damage, or 7 (1d10 + 2) slashing damage if used with two hands.

\DndMonsterAction{Longbow}
\textsl{Ranged Weapon Attack:} 4 to hit, range 150/600 ft., one target. \textsl{Hit:} 6 (1d8 + 2) piercing damage.

\end{DndMonster}



\begin{DndMonster}[float=t]{Hedge Mage}
\DndMonsterType{Medium humanoid, any alignment}
\DndMonsterBasics[
armor-class = {12 (15 with \textit{mage armor})},
hit-points = {\DndDice{4d8 + 4}},
speed = {30 ft.}
]
\DndMonsterAbilityScores[
 str = 10,
 dex = 14,
 con = 13,
 int = 15,
 wis = 11,
 cha = 11
]
\DndMonsterDetails[
skills = {Arcana +4},
languages = {Any two},
challenge = 1,
]
\DndMonsterAction{Spellcasting}
The hedge mage is a 4th-level spellcaster that uses Intelligence as its spellcasting ability (spell save DC 12; 4 to hit with spell attacks). The hedge mage has the following wizard spells prepared:
\begin{DndMonsterSpells}
  \DndMonsterSpellLevel{\textit{chill touch}, \textit{mage hand}, \textit{shocking grasp}}
   \DndMonsterSpellLevel[1][4]{\textit{false life}, \textit{mage armor}, \textit{magic missile}}
   \DndMonsterSpellLevel[2][3]{\textit{Melf's acid arrow}, \textit{web}}
\end{DndMonsterSpells}
\DndMonsterSection{Actions}
\DndMonsterAction{Dagger}
\textsl{Melee or Ranged Weapon Attack:} 4 to hit, reach 5 ft. or range 20/60 ft., one target. \textsl{Hit:} 4 (1d4 + 2) piercing damage.

\end{DndMonster}



\begin{DndMonster}[float=t]{Hell Hound}
\DndMonsterType{Medium fiend, lawful evil}
\DndMonsterBasics[
armor-class = {15 (natural armor)},
hit-points = {\DndDice{7d8 + 14}},
speed = {50 ft.}
]
\DndMonsterAbilityScores[
 str = 17,
 dex = 12,
 con = 14,
 int = 6,
 wis = 13,
 cha = 6
]
\DndMonsterDetails[
skills = {Perception +5},
damage-immunities = {fire},
senses = {darkvision 60 ft., passive perception 15},
languages = {understands Infernal but can't speak it},
challenge = 3,
]
\DndMonsterAction{Keen Hearing and Smell}
The hound has advantage on Wisdom (Perception) checks that rely on hearing or smell.

\DndMonsterAction{Pack Tactics}
The hound has advantage on an attack roll against a creature if at least one of the hound's allies is within 5 feet of the creature and the ally isn't incapacitated.

\DndMonsterSection{Actions}
\DndMonsterAction{Bite}
\textsl{Melee Weapon Attack:} 5 to hit, reach 5 ft., one target. \textsl{Hit:} 7 (1d8 + 3) piercing damage plus 7 (2d6) fire damage.

\DndMonsterAction{Fire Breath (Recharge 5\textendash 6)}
The hound exhales fire in a 15-foot cone. Each creature in that area must make a DC 12 Dexterity saving throw, taking 21 (6d6) fire damage on a failed save, or half as much damage on a successful one.

\end{DndMonster}



\begin{DndMonster}[float=t]{Homunculus}
\DndMonsterType{Tiny construct, neutral}
\DndMonsterBasics[
armor-class = {13 (natural armor)},
hit-points = {\DndDice{2d4}},
speed = {20 ft., fly 40 ft.}
]
\DndMonsterAbilityScores[
 str = 4,
 dex = 15,
 con = 11,
 int = 10,
 wis = 10,
 cha = 7
]
\DndMonsterDetails[
damage-immunities = {poison},
condition-immunities = {charmed, poisoned},
senses = {darkvision 60 ft., passive perception 10},
languages = {understands the languages of its creator but can't speak},
challenge = 0,
]
\DndMonsterAction{Telepathic Bond}
While the homunculus is on the same plane of existence as its master, it can magically convey what it senses to its master, and the two can communicate telepathically.

\DndMonsterSection{Actions}
\DndMonsterAction{Bite}
\textsl{Melee Weapon Attack:} 4 to hit, reach 5 ft., one creature. \textsl{Hit:} 1 piercing damage, and the target must succeed on a DC 10 Constitution saving throw or be poisoned for 1 minute. If the saving throw fails by 5 or more, the target is instead poisoned for 5 (1d10) minutes and unconscious while poisoned in this way.

\end{DndMonster}



\begin{DndMonster}[float=t]{Horned Devil}
\DndMonsterType{Large fiend (devil), lawful evil}
\DndMonsterBasics[
armor-class = {18 (natural armor)},
hit-points = {\DndDice{17d10 + 85}},
speed = {20 ft., fly 60 ft.}
]
\DndMonsterAbilityScores[
 str = 22,
 dex = 17,
 con = 21,
 int = 12,
 wis = 16,
 cha = 17
]
\DndMonsterDetails[
saving-throws = {Str +10, Dex +7, Wis +7, Cha +7},
damage-resistances = {cold; bludgeoning, piercing, slashing from nonmagical attacks that aren't silvered},
damage-immunities = {fire, poison},
condition-immunities = {poisoned},
senses = {darkvision 120 ft., passive perception 13},
languages = {Infernal, telepathy 120 ft.},
challenge = 11,
]
\DndMonsterAction{Devil's Sight}
Magical darkness doesn't impede the devil's darkvision.

\DndMonsterAction{Magic Resistance}
The devil has advantage on saving throws against spells and other magical effects.

\DndMonsterSection{Actions}
\DndMonsterAction{Multiattack}
The devil makes three melee attacks: two with its fork and one with its tail. It can use Hurl Flame in place of any melee attack.

\DndMonsterAction{Fork}
\textsl{Melee Weapon Attack:} 10 to hit, reach 10 ft., one target. \textsl{Hit:} 15 (2d8 + 6) piercing damage.

\DndMonsterAction{Tail}
\textsl{Melee Weapon Attack:} 10 to hit, reach 10 ft., one target. \textsl{Hit:} 10 (1d8 + 6) piercing damage. If the target is a creature other than an undead or a construct, it must succeed on a DC 17 Constitution saving throw or lose 10 (3d6) hit points at the start of each of its turns due to an infernal wound. Each time the devil hits the wounded target with this attack, the damage dealt by the wound increases by 10 (3d6). Any creature can take an action to stanch the wound with a successful DC 12 Wisdom (Medicine) check. The wound also closes if the target receives magical healing.

\DndMonsterAction{Hurl Flame}
\textsl{Ranged Spell Attack:} 7 to hit, range 150 ft., one target. \textsl{Hit:} 14 (4d6) fire damage. If the target is a flammable object that isn't being worn or carried, it also catches fire.

\end{DndMonster}



\begin{DndMonster}[float=t]{Hydra}
\DndMonsterType{Huge monstrosity, unaligned}
\DndMonsterBasics[
armor-class = {15 (natural armor)},
hit-points = {\DndDice{15d12 + 75}},
speed = {30 ft., swim 30 ft.}
]
\DndMonsterAbilityScores[
 str = 20,
 dex = 12,
 con = 20,
 int = 2,
 wis = 10,
 cha = 7
]
\DndMonsterDetails[
skills = {Perception +6},
senses = {darkvision 60 ft., passive perception 16},
challenge = 8,
]
\DndMonsterAction{Hold Breath}
The hydra can hold its breath for 1 hour.

\DndMonsterAction{Multiple Heads}
The hydra has five heads. While it has more than one head, the hydra has advantage on saving throws against being blinded, charmed, deafened, frightened, stunned, and knocked unconscious.

Whenever the hydra takes 25 or more damage in a single turn, one of its heads dies. If all its heads die, the hydra dies.

At the end of its turn, it grows two heads for each of its heads that died since its last turn, unless it has taken fire damage since its last turn. The hydra regains 10 hit points for each head regrown in this way.

\DndMonsterAction{Reactive Heads}
For each head the hydra has beyond one, it gets an extra reaction that can be used only for opportunity attacks.

\DndMonsterAction{Wakeful}
While the hydra sleeps, at least one of its heads is awake.

\DndMonsterSection{Actions}
\DndMonsterAction{Multiattack}
The hydra makes as many bite attacks as it has heads.

\DndMonsterAction{Bite}
\textsl{Melee Weapon Attack:} 8 to hit, reach 10 ft., one target. \textsl{Hit:} 10 (1d10 + 5) piercing damage.

\end{DndMonster}



\begin{DndMonster}[float=t]{Hypnotic Eldritch Blossom}
\DndMonsterType{Medium plant, chaotic evil}
\DndMonsterBasics[
armor-class = {5},
hit-points = {\DndDice{4d8}},
speed = {5 ft.}
]
\DndMonsterAbilityScores[
 str = 3,
 dex = 1,
 con = 10,
 int = 1,
 wis = 3,
 cha = 5
]
\DndMonsterDetails[
condition-immunities = {blinded, deafened, frightened},
senses = {blindsight 30 ft., passive perception 6},
challenge = 1,
]
\DndMonsterAction{False Appearance}
While the Hypnotic Eldritch Blossom remains motionless, it is indistinguishable from an ordinary flower.

\DndMonsterSection{Actions}
\DndMonsterAction{Contaminated Nectar}
The Hypnotic Eldritch Blossom targets one creature it can see within 10 feet of it. The target must succeed on a DC 10 Constitution saving throw or gain one level of contamination.

\DndMonsterAction{Hypnotic Pollen (recharges on a short or long rest)}
The Hypnotic Eldritch Blossom releases a cloud of glittering pollen out to a radius of 30 feet. Creatures in the area must succeed on a DC 10 Wisdom saving throw or become charmed by the Hypnotic Eldritch Blossom for 1 minute. While charmed, the creature is incapacitated and must use all its movement on its turns to get as close to the Hypnotic Eldritch Blossom as possible.

The effect ends for an affected creature if it takes any damage or if someone else uses an action to shake the creature out of its stupor.

\DndMonsterSection{Reactions}
\DndMonsterAction{Spell Reflection}
If the hypnotic eldritch blossom makes a successful saving throw against a spell which targets only it and does not have an area of effect, or a spell attack misses it, the blossom can choose another creature (including the spellcaster) it can see within 30 feet of it. The spell targets the chosen creature instead of the hypnotic eldritch blossom. If the spell forces a saving throw, the chosen creature makes its own save. If the spell was an attack, the attack roll is rerolled against the chosen creature.

\end{DndMonster}



\begin{DndMonster}[float=t]{Imp}
\DndMonsterType{Tiny fiend (devil), lawful evil}
\DndMonsterBasics[
armor-class = {13},
hit-points = {\DndDice{3d4 + 3}},
speed = {20 ft., fly 40 ft.}
]
\DndMonsterAbilityScores[
 str = 6,
 dex = 17,
 con = 13,
 int = 11,
 wis = 12,
 cha = 14
]
\DndMonsterDetails[
skills = {Deception +4, Insight +3, Persuasion +4, Stealth +5},
damage-resistances = {cold; bludgeoning, piercing, slashing from nonmagical attacks not made with silvered weapons},
damage-immunities = {fire, poison},
condition-immunities = {poisoned},
senses = {darkvision 120 ft., passive perception 11},
languages = {Infernal, Common},
challenge = 1,
]
\DndMonsterAction{Shapechanger}
The imp can use its action to polymorph into a beast form that resembles a rat (speed 20 ft.), a raven (20 ft., fly 60 ft.), or a spider (20 ft., climb 20 ft.), or back into its true form. Its statistics are the same in each form, except for the speed changes noted. Any equipment it is wearing or carrying isn't transformed. It reverts to its true form if it dies.

\DndMonsterAction{Devil's Sight}
Magical darkness doesn't impede the imp's darkvision.

\DndMonsterAction{Magic Resistance}
The imp has advantage on saving throws against spells and other magical effects.

\DndMonsterSection{Actions}
\DndMonsterAction{Sting (Bite in Beast Form)}
\textsl{Melee Weapon Attack:} 5 to hit, reach 5 ft., one target. \textsl{Hit:} 5 (1d4 + 3) piercing damage, and the target must make a DC 11 Constitution saving throw, taking 10 (3d6) poison damage on a failed save, or half as much damage on a successful one.

\DndMonsterAction{Invisibility}
The imp magically turns invisible until it attacks, or until its concentration ends (as if concentration on a spell). Any equipment the imp wears or carries is invisible with it.

\end{DndMonster}



\begin{DndMonster}[float=t]{Invisible Stalker}
\DndMonsterType{Medium elemental, neutral}
\DndMonsterBasics[
armor-class = {14},
hit-points = {\DndDice{16d8 + 32}},
speed = {50 ft., fly 50 ft. \textit{(hover)}}
]
\DndMonsterAbilityScores[
 str = 16,
 dex = 19,
 con = 14,
 int = 10,
 wis = 15,
 cha = 11
]
\DndMonsterDetails[
skills = {Perception +8, Stealth +10},
damage-resistances = {bludgeoning, piercing, slashing from nonmagical attacks},
damage-immunities = {poison},
condition-immunities = {exhaustion, grappled, paralyzed, petrified, poisoned, prone, restrained, unconscious},
senses = {darkvision 60 ft., passive perception 18},
languages = {Auran, understands Common but doesn't speak it},
challenge = 6,
]
\DndMonsterAction{Invisibility}
The stalker is invisible.

\DndMonsterAction{Faultless Tracker}
The stalker is given a quarry by its summoner. The stalker knows the direction and distance to its quarry as long as the two of them are on the same plane of existence. The stalker also knows the location of its summoner.

\DndMonsterSection{Actions}
\DndMonsterAction{Multiattack}
The stalker makes two slam attacks.

\DndMonsterAction{Slam}
\textsl{Melee Weapon Attack:} 6 to hit, reach 5 ft., one target. \textsl{Hit:} 10 (2d6 + 3) bludgeoning damage.

\end{DndMonster}



\begin{DndMonster}[float=t]{Iron Golem}
\DndMonsterType{Large construct, unaligned}
\DndMonsterBasics[
armor-class = {20 (natural armor)},
hit-points = {\DndDice{20d10 + 100}},
speed = {30 ft.}
]
\DndMonsterAbilityScores[
 str = 24,
 dex = 9,
 con = 20,
 int = 3,
 wis = 11,
 cha = 1
]
\DndMonsterDetails[
damage-immunities = {fire; poison; psychic; bludgeoning, piercing, slashing from nonmagical attacks that aren't adamantine},
condition-immunities = {charmed, exhaustion, frightened, paralyzed, petrified, poisoned},
senses = {darkvision 120 ft., passive perception 10},
languages = {understands the languages of its creator but can't speak},
challenge = 16,
]
\DndMonsterAction{Fire Absorption}
Whenever the golem is subjected to fire damage, it takes no damage and instead regains a number of hit points equal to the fire damage dealt.

\DndMonsterAction{Immutable Form}
The golem is immune to any spell or effect that would alter its form.

\DndMonsterAction{Magic Resistance}
The golem has advantage on saving throws against spells and other magical effects.

\DndMonsterAction{Magic Weapons}
The golem's weapon attacks are magical.

\DndMonsterSection{Actions}
\DndMonsterAction{Multiattack}
The golem makes two melee attacks.

\DndMonsterAction{Slam}
\textsl{Melee Weapon Attack:} 13 to hit, reach 5 ft., one target. \textsl{Hit:} 20 (3d8 + 7) bludgeoning damage.

\DndMonsterAction{Sword}
\textsl{Melee Weapon Attack:} 13 to hit, reach 10 ft., one target. \textsl{Hit:} 23 (3d10 + 7) slashing damage.

\DndMonsterAction{Poison Breath (Recharge 5\textendash 6)}
The golem exhales poisonous gas in a 15-foot cone. Each creature in that area must make a DC 19 Constitution saving throw, taking 45 (10d8) poison damage on a failed save, or half as much damage on a successful one.

\end{DndMonster}



\begin{DndMonster}[float=t]{Knight}
\DndMonsterType{Medium humanoid (any race), any alignment}
\DndMonsterBasics[
armor-class = {18 (\textit{plate armor})},
hit-points = {\DndDice{8d8 + 16}},
speed = {30 ft.}
]
\DndMonsterAbilityScores[
 str = 16,
 dex = 11,
 con = 14,
 int = 11,
 wis = 11,
 cha = 15
]
\DndMonsterDetails[
saving-throws = {Con +4, Wis +2},
languages = {any one language (usually Common)},
challenge = 3,
]
\DndMonsterAction{Brave}
The knight has advantage on saving throws against being frightened.

\DndMonsterSection{Actions}
\DndMonsterAction{Multiattack}
The knight makes two melee attacks.

\DndMonsterAction{Greatsword}
\textsl{Melee Weapon Attack:} 5 to hit, reach 5 ft., one target. \textsl{Hit:} 10 (2d6 + 3) slashing damage.

\DndMonsterAction{Heavy Crossbow}
\textsl{Ranged Weapon Attack:} 2 to hit, range 100/400 ft., one target. \textsl{Hit:} 5 (1d10) piercing damage.

\DndMonsterAction{Leadership (Recharges after a Short or Long Rest)}
For 1 minute, the knight can utter a special command or warning whenever a nonhostile creature that it can see within 30 feet of it makes an attack roll or a saving throw. The creature can add a d4 to its roll provided it can hear and understand the knight. A creature can benefit from only one Leadership die at a time. This effect ends if the knight is incapacitated.

\DndMonsterSection{Reactions}
\DndMonsterAction{Parry}
The knight adds 2 to its AC against one melee attack that would hit it. To do so, the knight must see the attacker and be wielding a melee weapon.

\end{DndMonster}



\begin{DndMonster}[float*=b,width=\textwidth + 8pt]{Knight-Captain Theodore Marshal}
\begin{multicols}{2}
\DndMonsterType{Medium humanoid (human), lawful good}
\DndMonsterBasics[
armor-class = {20 (plate armor)},
hit-points = {\DndDice{30d8 + 120}},
speed = {30 ft.}
]
\DndMonsterAbilityScores[
 str = 20,
 dex = 10,
 con = 18,
 int = 13,
 wis = 15,
 cha = 19
]
\DndMonsterDetails[
saving-throws = {Wis +7, Cha +9},
skills = {Animal Handling +7, Athletics +10, Insight +7, Persuasion +9},
condition-immunities = {charmed, frightened},
languages = {Celestial, Common},
challenge = 15,
]
\DndMonsterAction{Divine Strikes}
Theodore Marshal's weapon attacks are magical. When he hits with any weapon, the weapon deals an extra 4d8 radiant damage (included in the attack).

\DndMonsterAction{Legendary Resistance (3/Day)}
If Theodore Marshal fails a saving throw, he can choose to succeed instead.

\DndMonsterAction{Special Equipment}
Theodore Marshal wields a \textit{holy avenger longsword}.

\DndMonsterAction{Spellcasting}
Theodore Marshal is a 20th-level spellcaster. His spellcasting ability is Charisma (spell save DC 17, 9 to hit with spell attacks). He has the following paladin spells prepared:
\begin{DndMonsterSpells}
   \DndMonsterSpellLevel[1][4]{\textit{bless}, \textit{command}, \textit{cure wounds}, \textit{guiding bolt}}
   \DndMonsterSpellLevel[2][3]{\textit{aid}, \textit{branding smite}, \textit{lesser restoration}}
   \DndMonsterSpellLevel[3][3]{\textit{beacon of hope}, \textit{dispel magic}}
   \DndMonsterSpellLevel[4][3]{\textit{banishment}, \textit{death ward}}
   \DndMonsterSpellLevel[5][2]{\textit{flame strike}, \textit{hallow}}
\end{DndMonsterSpells}
\DndMonsterSection{Actions}
\DndMonsterAction{Multiattack}
Theodore Marshal makes two melee attacks.

\DndMonsterAction{Holy Avenger}
\textsl{Melee Weapon Attack:} 13 to hit, reach 5 ft., one target. \textsl{Hit:} 12 (1d8 + 8) slashing damage plus 18 (4d8) radiant damage.

\DndMonsterSection{Legendary Actions}
Knight-Captain Theodore Marshal can take 3 legendary actions, choosing from the options below. Only one legendary action can be used at a time and only at the end of another creature's turn. Knight-Captain Theodore Marshal regains spent legendary actions at the start of its turn.

\begin{DndMonsterLegendaryActions}
\DndMonsterLegendaryAction{Attack}{Theodore Marshal makes one weapon attack.
}
\DndMonsterLegendaryAction{Holy Bolt (Costs 2 actions)}{Theodore Marshal casts \textit{guiding bolt} without expending a spell slot.
}
\end{DndMonsterLegendaryActions}
\end{multicols}
\end{DndMonster}



\begin{DndMonster}[float=t]{Lenore Von Kessel}
\DndMonsterType{Medium monstrosity, lawful evil}
\DndMonsterBasics[
armor-class = {15 (natural armor)},
hit-points = {\DndDice{17d8 + 51}},
speed = {30 ft.}
]
\DndMonsterAbilityScores[
 str = 10,
 dex = 15,
 con = 16,
 int = 12,
 wis = 13,
 cha = 15
]
\DndMonsterDetails[
skills = {Deception +5, Insight +4, Perception +4, Stealth +5},
senses = {darkvision 60 ft., passive perception 14},
languages = {Common},
challenge = 6,
]
\DndMonsterAction{Copy Warning}
This content has been copied from another creature, but the data replacement hasn't been run. Check details.

\DndMonsterAction{Petrifying Gaze}
When a creature that can see the medusa's eyes starts its turn within 30 feet of the medusa, the medusa can force it to make a DC 14 Constitution saving throw if the medusa isn't incapacitated and can see the creature. If the saving throw fails by 5 or more, the creature is instantly petrified. Otherwise, a creature that fails the save begins to turn to stone and is restrained. The restrained creature must repeat the saving throw at the end of its next turn, becoming petrified on a failure or ending the effect on a success. The petrification lasts until the creature is freed by the  \textit{greater restoration} spell or other magic.

Unless Surprise, a creature can avert its eyes to avoid the saving throw at the start of its turn. If the creature does so, it can't see the medusa until the start of its next turn, when it can avert its eyes again. If the creature looks at the medusa in the meantime, it must immediately make the save.

If the medusa sees itself reflected on a polished surface within 30 feet of it and in an area of bright light, the medusa is, due to its curse, affected by its own gaze.

\DndMonsterSection{Actions}
\DndMonsterAction{Multiattack}
The medusa makes either three melee attacks—one with its snake hair and two with its shortsword—or two ranged attacks with its longbow.

\DndMonsterAction{Snake Hair}
\textsl{Melee Weapon Attack:} 5 to hit, reach 5 ft., one creature. \textsl{Hit:} 4 (1d4 + 2) piercing damage plus 14 (4d6) poison damage.

\DndMonsterAction{Shortsword}
\textsl{Melee Weapon Attack:} 5 to hit, reach 5 ft., one target. \textsl{Hit:} 5 (1d6 + 2) piercing damage.

\DndMonsterAction{Longbow}
\textsl{Ranged Weapon Attack:} 5 to hit, range 150/600 ft., one target. \textsl{Hit:} 6 (1d8 + 2) piercing damage plus 7 (2d6) poison damage.

\end{DndMonster}


\clearpage

\begin{DndMonster}[float=t]{Living Deep Haze}
\DndMonsterType{Large elemental, unaligned}
\DndMonsterBasics[
armor-class = {15},
hit-points = {\DndDice{12d10 + 24}},
speed = {0 ft., fly 90 ft. \textit{(hover)}}
]
\DndMonsterAbilityScores[
 str = 14,
 dex = 20,
 con = 14,
 int = 5,
 wis = 10,
 cha = 5
]
\DndMonsterDetails[
damage-resistances = {bludgeoning, piercing, slashing from nonmagical attacks},
damage-immunities = {necrotic, poison},
condition-immunities = {exhaustion, grappled, paralyzed, petrified, poisoned, prone, restrained, unconscious},
senses = {darkvision 60 ft., passive perception 10},
languages = {Deep Speech},
challenge = 5,
]
\DndMonsterAction{Fog Form}
The Living Deep Haze can enter a hostile creature's space and stop there. It can move through a space as narrow as 1 inch wide without squeezing.

\DndMonsterAction{Deep Haze}
A creature who enters or starts its turn in the same space as the Living Deep Haze must succeed on a DC 15 Constitution saving throw or take 10 (3d6) necrotic damage and gain one level of contamination.

\DndMonsterSection{Actions}
\DndMonsterAction{Multiattack}
The Living Deep Haze makes two slam attacks.

\DndMonsterAction{Slam}
\textsl{Melee Weapon Attack:} 8 to hit, reach 5 ft., one target. \textsl{Hit:} 14 (2d8 + 5) bludgeoning damage.

\end{DndMonster}



\begin{DndMonster}[float*=b,width=\textwidth + 8pt]{Lord Commander Elias Drexel}
\begin{multicols}{2}
\DndMonsterType{Medium humanoid (human), lawful neutral}
\DndMonsterBasics[
armor-class = {19 (\textit{+3 breastplate})},
hit-points = {\DndDice{30d8 + 120}},
speed = {30 ft.}
]
\DndMonsterAbilityScores[
 str = 18,
 dex = 15,
 con = 18,
 int = 15,
 wis = 16,
 cha = 11
]
\DndMonsterDetails[
saving-throws = {Str +9, Dex +7, Wis +8},
skills = {Athletics +9, Perception +13, Stealth +12, Survival +8},
senses = {blindsight 10 ft., passive perception 23},
languages = {Common},
challenge = 15,
]
\DndMonsterAction{Survivor}
At the start of each of his turns, Elias Drexel regains 10 hit points if he has no more than half of his hit points left. He does not gain this benefit if he has 0 hit points.

\DndMonsterAction{Brute}
A melee weapon deals one extra die of its damage when Elias Drexel hits with it (included in the attack).

\DndMonsterAction{Legendary Resistance (3/Day)}
If Elias Drexel fails a saving throw, he can choose to succeed instead.

\DndMonsterAction{Special Equipment}
Elias Drexel wields a \textit{+3 Longsword}, wears a \textit{+3 breastplate} and the \textit{Lord Commander's Badge}.

\DndMonsterAction{Spellcasting}
Elias Drexel is a 20th-level spellcaster. His spellcasting ability is Wisdom (spell save DC 16, 8 to hit with spell attacks). He knows the following ranger spells:
\begin{DndMonsterSpells}
   \DndMonsterSpellLevel[1][4]{\textit{cure wounds}, \textit{fog cloud}, \textit{hunter's mark}}
   \DndMonsterSpellLevel[2][3]{\textit{pass without trace}, \textit{protection from poison}, \textit{spike growth}}
   \DndMonsterSpellLevel[3][3]{\textit{conjure animals}, \textit{nondetection}}
   \DndMonsterSpellLevel[4][3]{\textit{freedom of movement}, \textit{stoneskin}}
   \DndMonsterSpellLevel[5][2]{\textit{tree stride} (works on buildings instead of trees)}
\end{DndMonsterSpells}
\DndMonsterSection{Actions}
\DndMonsterAction{Multiattack}
Elias Drexel uses Leadership. He then makes four attacks with either his longsword or his longbow.

\DndMonsterAction{Longsword}
\textsl{Melee Weapon Attack:} 12 to hit, reach 5 ft., one target. \textsl{Hit:} 16 (2d8 + 7) slashing damage.

\DndMonsterAction{Longbow}
\textsl{Ranged Weapon Attack:} 7 to hit, range 150/600 ft., one target. \textsl{Hit:} 7 (1d8 + 2) piercing damage.

\DndMonsterAction{Leadership (Recharges after a Short or Long Rest)}
For 1 minute, Elias Drexel can utter a special command or warning whenever a nonhostile creature that he can see within 120 feet of him makes an attack roll or a saving throw. The creature can add a d4 to its roll provided it can hear and understand Elias Drexel. A creature can benefit from only one Leadership die at a time. This effect ends if Elias Drexel is incapacitated.

\DndMonsterSection{Reactions}
\DndMonsterAction{Parry}
Elias Drexel adds 5 to his AC against one melee attack that would hit him. To do so, Elias Drexel must see the attacker and be wielding a melee weapon.

\DndMonsterSection{Legendary Actions}
Lord Commander Elias Drexel can take 3 legendary actions, choosing from the options below. Only one legendary action can be used at a time and only at the end of another creature's turn. Lord Commander Elias Drexel regains spent legendary actions at the start of its turn.

\begin{DndMonsterLegendaryActions}
\DndMonsterLegendaryAction{Attack}{Elias Drexel makes one weapon attack.
}
\DndMonsterLegendaryAction{Command}{Elias Drexel chooses one ally he can see. If that ally can see and hear him, that ally can make one weapon attack as a reaction and gains advantage on the attack roll. The attack deals an extra 10 3d6 damage on a hit.
}
\end{DndMonsterLegendaryActions}
\end{multicols}
\end{DndMonster}



\begin{DndMonster}[float*=b,width=\textwidth + 8pt]{Lord of the Feast}
\begin{multicols}{2}
\DndMonsterType{Large monstrosity, chaotic evil}
\DndMonsterBasics[
armor-class = {18 (natural armor)},
hit-points = {\DndDice{24d10 + 96}},
speed = {30 ft., climb 30 ft.}
]
\DndMonsterAbilityScores[
 str = 18,
 dex = 18,
 con = 18,
 int = 11,
 wis = 18,
 cha = 13
]
\DndMonsterDetails[
saving-throws = {Dex +9, Con +9, Wis +9, Cha +6},
skills = {Athletics +9, Perception +14, Stealth +14, Survival +9},
senses = {darkvision 120 ft., passive perception 24},
languages = {Abyssal},
challenge = 15,
]
\DndMonsterAction{Legendary Resistance (3/Day)}
If the Lord of the Feast fails a saving throw, he can choose to succeed instead.

\DndMonsterAction{Regeneration}
The Lord of the Feast regains 20 hit points at the start of his turn if he has at least 1 hit point remaining.

\DndMonsterAction{Close-Quarters Hunter}
Being within 5 feet of a hostile creature doesn't impose disadvantage on the Lord of the Feast's ranged attack rolls.

\DndMonsterAction{Predatory Senses}
The Lord of the Feast has advantage on Wisdom (Perception) checks. He can pinpoint the position of any creature within 120 feet of him that is invisible.

\DndMonsterSection{Actions}
\DndMonsterAction{Multiattack}
The Lord of the Feast makes three attacks: either one with each of its different arrow attacks, or one with his bite and two with his claws.

\DndMonsterAction{Bite}
\textsl{Melee Weapon Attack:} 9 to hit, reach 5 ft., one target. \textsl{Hit:} 15 (2d10 + 4) piercing damage plus 18 (4d8) necrotic damage. A target hit by this attack must succeed on a DC 17 Constitution saving throw or gain one level of contamination.

\DndMonsterAction{Claws}
\textsl{Melee Weapon Attack:} 9 to hit, reach 10 ft, one target. \textsl{Hit:} 11 (2d6 + 4) slashing damage.

\DndMonsterAction{Contaminating Arrow}
\textsl{Ranged Weapon Attack:} 9 to hit, range 150/600 ft., one target. \textsl{Hit:} 13 (2d8 + 4) piercing damage plus 18 (4d8) necrotic damage. A target hit by this attack must succeed on a DC 17 Constitution saving throw or gain one level of contamination.

\DndMonsterAction{Ensnaring Arrow}
\textsl{Ranged Weapon Attack:} 9 to hit, range 150/600 ft., one target. \textsl{Hit:} 13 (2d8 + 4) piercing damage. If the target is Large or smaller, it must succeed on a DC 17 Dexterity saving throw or become restrained in a magical snare for 1 minute. A restrained target can use its action to repeat the saving throw, ending the effect on itself on a success.

\DndMonsterAction{Harpoon Arrow}
\textsl{Ranged Weapon Attack:} 9 to hit, range 150/600 ft., one target. \textsl{Hit:} 13 (2d8 + 4) piercing damage. If the target is a Huge or smaller creature, it must succeed on DC 17 Strength saving throw or be pulled up to 20 feet directly toward the Lord of the Feast.

\DndMonsterSection{Legendary Actions}
Lord of the Feast can take 3 legendary actions, choosing from the options below. Only one legendary action can be used at a time and only at the end of another creature's turn. Lord of the Feast regains spent legendary actions at the start of its turn.

\begin{DndMonsterLegendaryActions}
\DndMonsterLegendaryAction{Dash}{The Lord of the Feast moves his speed without provoking opportunity attacks.
}
\DndMonsterLegendaryAction{Arrow (Costs 2 actions)}{The Lord of the Feast makes one arrow attack of his choice.
}
\DndMonsterLegendaryAction{Call of the Hunt (costs 3 actions)}{The Lord of the Feast releases a blood-curdling howl which can be heard all across Drakkenheim. Garmyr within 60 feet who can see and hear the Lord of the Feast can use their reaction to move up to half their speed and make one melee or ranged attack with advantage.
}
\end{DndMonsterLegendaryActions}
\end{multicols}
\end{DndMonster}



\begin{DndMonster}[float*=b,width=\textwidth + 8pt]{Lucretia Mathias}
\begin{multicols}{2}
\DndMonsterType{Medium humanoid (human), neutral good}
\DndMonsterBasics[
armor-class = {17 (Unarmored Defense)},
hit-points = {\DndDice{20d8}},
speed = {30 ft.}
]
\DndMonsterAbilityScores[
 str = 8,
 dex = 10,
 con = 10,
 int = 19,
 wis = 24,
 cha = 19
]
\DndMonsterDetails[
saving-throws = {Wis +12, Cha +9},
skills = {Insight +17, Perception +17, Persuasion +14, Religion +14},
condition-immunities = {charmed, frightened},
senses = {truesight 120 ft., passive perception 26},
languages = {all},
challenge = 15,
]
\DndMonsterAction{Prescient}
Lucretia Mathias can't be surprised while she is conscious.

\DndMonsterAction{Unarmoured Defense}
While Lucretia Mathias is wearing no armour and wielding no shield, her AC includes her Wisdom modifier.

\DndMonsterAction{Undeniable Faith}
Lucretia Mathias automatically succeeds on Constitution saving throws to maintain concentration on her spells. Any spell she casts is treated as if it were cast using a 9th level spell slot, regardless of the spell slot used to cast the spell.

\DndMonsterAction{Legendary Resistance (3/Day)}
If Lucretia Mathias fails a saving throw, she can choose to succeed instead.

\DndMonsterAction{Divine Grace}
Any creature who targets Lucretia Mathias with an attack or a harmful spell must first make a DC 20 Wisdom saving throw. On a failed save, the creature must choose a new target or lose the attack or spell. This effect doesn't protect Lucretia Mathias from area effects, such as the explosion of a fireball.

If Lucretia Mathias makes an attack, casts a spell that affects an enemy, or deals damage to another creature, this effect ends until the start of her next turn.

\DndMonsterAction{Spellcasting}
Lucretia Mathias is a 20th-level spellcaster. Her spellcasting ability is Wisdom (spell save DC 20, 12 to hit with spell attacks). She has the following cleric spells prepared:
\begin{DndMonsterSpells}
  \DndMonsterSpellLevel{\textit{chill touch}, \textit{guidance}, \textit{light}, \textit{sacred flame}, \textit{thaumaturgy}}
   \DndMonsterSpellLevel[1][4]{\textit{bless}, \textit{command}, \textit{guiding bolt}, \textit{healing word}}
   \DndMonsterSpellLevel[2][3]{\textit{aid}, \textit{hold person}, \textit{lesser restoration}, \textit{spiritual weapon}}
   \DndMonsterSpellLevel[3][3]{\textit{dispel magic}, \textit{sending}, \textit{speak with dead}, \textit{spirit guardians}}
   \DndMonsterSpellLevel[4][3]{\textit{banishment}, \textit{death ward}, \textit{divination}}
   \DndMonsterSpellLevel[5][3]{\textit{commune}, \textit{flame strike}, \textit{greater restoration}, \textit{scrying}}
   \DndMonsterSpellLevel[6][2]{\textit{heal}, \textit{word of recall}}
   \DndMonsterSpellLevel[7][2]{\textit{divine word}, \textit{resurrection}}
   \DndMonsterSpellLevel[8][1]{\textit{holy aura}}
\end{DndMonsterSpells}
\DndMonsterSection{Actions}
\DndMonsterAction{Pacifying Touch}
\textsl{Melee Spell Attack:} 12 to hit, reach 5 ft., one creature. \textsl{Hit:} 10 (3d6) psychic damage. A creature reduced to 0 hit points by this damage is stable at 0 hit points. The target must succeed on a DC 20 Wisdom saving throw or be incapacitated for 1 minute. This effect ends for an affected creature if it takes any damage.

\DndMonsterAction{Divine Intervention (1/ week)}
Lucretia Mathias calls on the Sacred Flame for aid. The GM chooses the nature of the intervention; the effect of any cleric spell or cleric domain spell would be appropriate.

\DndMonsterSection{Legendary Actions}
Lucretia Mathias can take 3 legendary actions, choosing from the options below. Only one legendary action can be used at a time and only at the end of another creature's turn. Lucretia Mathias regains spent legendary actions at the start of its turn.

\begin{DndMonsterLegendaryActions}
\DndMonsterLegendaryAction{Cantrip}{Lucretia Mathias casts a cantrip.
}
\DndMonsterLegendaryAction{Pacifying Touch (Costs 2 Actions)}{Lucretia Mathias uses her Pacifying Touch.
}
\DndMonsterLegendaryAction{Cast a Spell (Costs 3 Actions)}{Lucretia Mathias casts a spell using a spell slot.
}
\end{DndMonsterLegendaryActions}
\end{multicols}
\end{DndMonster}



\begin{DndMonster}[float=t]{Mage}
\DndMonsterType{Medium humanoid (any race), any alignment}
\DndMonsterBasics[
armor-class = {12 (15 with \textit{mage armor})},
hit-points = {\DndDice{9d8}},
speed = {30 ft.}
]
\DndMonsterAbilityScores[
 str = 9,
 dex = 14,
 con = 11,
 int = 17,
 wis = 12,
 cha = 11
]
\DndMonsterDetails[
saving-throws = {Int +6, Wis +4},
skills = {Arcana +6, History +6},
languages = {any four languages},
challenge = 6,
]
\DndMonsterAction{Spellcasting}
The mage is a 9th-level spellcaster. Its spellcasting ability is Intelligence (spell save DC 14, 6 to hit with spell attacks). The mage has the following wizard spells prepared:
\begin{DndMonsterSpells}
  \DndMonsterSpellLevel{\textit{fire bolt}, \textit{light}, \textit{mage hand}, \textit{prestidigitation}}
   \DndMonsterSpellLevel[1][4]{\textit{detect magic}, \textit{mage armor}, \textit{magic missile}, \textit{shield}}
   \DndMonsterSpellLevel[2][3]{\textit{misty step}, \textit{suggestion}}
   \DndMonsterSpellLevel[3][3]{\textit{counterspell}, \textit{fireball}, \textit{fly}}
   \DndMonsterSpellLevel[4][3]{\textit{greater invisibility}, \textit{ice storm}}
   \DndMonsterSpellLevel[5][1]{\textit{cone of cold}}
\end{DndMonsterSpells}
\DndMonsterSection{Actions}
\DndMonsterAction{Dagger}
\textsl{Melee or Ranged Weapon Attack:} 5 to hit, reach 5 ft. or range 20/60 ft., one target. \textsl{Hit:} 4 (1d4 + 2) piercing damage.

\end{DndMonster}



\begin{DndMonster}[float=t]{Magma Mephit}
\DndMonsterType{Small elemental, neutral evil}
\DndMonsterBasics[
armor-class = {11},
hit-points = {\DndDice{5d6 + 5}},
speed = {30 ft., fly 30 ft.}
]
\DndMonsterAbilityScores[
 str = 8,
 dex = 12,
 con = 12,
 int = 7,
 wis = 10,
 cha = 10
]
\DndMonsterDetails[
skills = {Stealth +3},
damage-vulnerabilities = {cold},
damage-immunities = {fire, poison},
condition-immunities = {poisoned},
senses = {darkvision 60 ft., passive perception 10},
languages = {Ignan, Terran},
challenge = 1/2,
]
\DndMonsterAction{Death Burst}
When the mephit dies, it explodes in a burst of lava. Each creature within 5 feet of it must make a DC 11 Dexterity saving throw, taking 7 (2d6) fire damage on a failed save, or half as much damage on a successful one.

\DndMonsterAction{False Appearance}
While the mephit remains motionless, it is indistinguishable from an ordinary mound of magma.

\DndMonsterAction{Innate Spellcasting (1/Day)}
The mephit can innately cast \textit{heat metal} (spell save DC 10), requiring no material components. Its innate spellcasting ability is Charisma.
\begin{DndMonsterSpells}
  \DndInnateSpellLevel{\textit{heat metal}}
\end{DndMonsterSpells}
\DndMonsterSection{Actions}
\DndMonsterAction{Claws}
\textsl{Melee Weapon Attack:} 3 to hit, reach 5 ft., one creature. \textsl{Hit:} 3 (1d4 + 1) slashing damage plus 2 (1d4) fire damage.

\DndMonsterAction{Fire Breath (Recharge 6)}
The mephit exhales a 15-foot cone of fire. Each creature in that area must make a DC 11 Dexterity saving throw, taking 7 (2d6) fire damage on a failed save, or half as much damage on a successful one.

\end{DndMonster}



\begin{DndMonster}[float=t]{Manticore}
\DndMonsterType{Large monstrosity, lawful evil}
\DndMonsterBasics[
armor-class = {14 (natural armor)},
hit-points = {\DndDice{8d10 + 24}},
speed = {30 ft., fly 50 ft.}
]
\DndMonsterAbilityScores[
 str = 17,
 dex = 16,
 con = 17,
 int = 7,
 wis = 12,
 cha = 8
]
\DndMonsterDetails[
senses = {darkvision 60 ft., passive perception 11},
languages = {Common},
challenge = 3,
]
\DndMonsterAction{Tail Spike Regrowth}
The manticore has twenty-four tail spikes. Used spikes regrow when the manticore finishes a long rest.

\DndMonsterSection{Actions}
\DndMonsterAction{Multiattack}
The manticore makes three attacks: one with its bite and two with its claws or three with its tail spikes.

\DndMonsterAction{Bite}
\textsl{Melee Weapon Attack:} 5 to hit, reach 5 ft., one target. \textsl{Hit:} 7 (1d8 + 3) piercing damage.

\DndMonsterAction{Claw}
\textsl{Melee Weapon Attack:} 5 to hit, reach 5 ft., one target. \textsl{Hit:} 6 (1d6 + 3) slashing damage.

\DndMonsterAction{Tail Spike}
\textsl{Ranged Weapon Attack:} 5 to hit, range 100/200 ft., one target. \textsl{Hit:} 7 (1d8 + 3) piercing damage.

\end{DndMonster}



\begin{DndMonster}[float=t]{Marilith}
\DndMonsterType{Large fiend (demon), chaotic evil}
\DndMonsterBasics[
armor-class = {18 (natural armor)},
hit-points = {\DndDice{18d10 + 90}},
speed = {40 ft.}
]
\DndMonsterAbilityScores[
 str = 18,
 dex = 20,
 con = 20,
 int = 18,
 wis = 16,
 cha = 20
]
\DndMonsterDetails[
saving-throws = {Str +9, Con +10, Wis +8, Cha +10},
damage-resistances = {cold; fire; lightning; bludgeoning, piercing, slashing from nonmagical attacks},
damage-immunities = {poison},
condition-immunities = {poisoned},
senses = {truesight 120 ft., passive perception 13},
languages = {Abyssal, telepathy 120 ft.},
challenge = 16,
]
\DndMonsterAction{Magic Resistance}
The marilith has advantage on saving throws against spells and other magical effects.

\DndMonsterAction{Magic Weapons}
The marilith's weapon attacks are magical.

\DndMonsterAction{Reactive}
The marilith can take one reaction on every turn in combat.

\DndMonsterSection{Actions}
\DndMonsterAction{Multiattack}
The marilith can make seven attacks: six with its longswords and one with its tail.

\DndMonsterAction{Longsword}
\textsl{Melee Weapon Attack:} 9 to hit, reach 5 ft., one target. \textsl{Hit:} 13 (2d8 + 4) slashing damage.

\DndMonsterAction{Tail}
\textsl{Melee Weapon Attack:} 9 to hit, reach 10 ft., one creature. \textsl{Hit:} 15 (2d10 + 4) bludgeoning damage. If the target is Medium or smaller, it is grappled (escape DC 19). Until this grapple ends, the target is restrained, the marilith can automatically hit the target with its tail, and the marilith can't make tail attacks against other targets.

\DndMonsterAction{Teleport}
The marilith magically teleports, along with any equipment it is wearing or carrying, up to 120 feet to an unoccupied space it can see.

\DndMonsterSection{Reactions}
\DndMonsterAction{Parry}
The marilith adds 5 to its AC against one melee attack that would hit it. To do so, the marilith must see the attacker and be wielding a melee weapon.

\end{DndMonster}



\begin{DndMonster}[float=t]{Medusa}
\DndMonsterType{Medium monstrosity, lawful evil}
\DndMonsterBasics[
armor-class = {15 (natural armor)},
hit-points = {\DndDice{17d8 + 51}},
speed = {30 ft.}
]
\DndMonsterAbilityScores[
 str = 10,
 dex = 15,
 con = 16,
 int = 12,
 wis = 13,
 cha = 15
]
\DndMonsterDetails[
skills = {Deception +5, Insight +4, Perception +4, Stealth +5},
senses = {darkvision 60 ft., passive perception 14},
languages = {Common},
challenge = 6,
]
\DndMonsterAction{Petrifying Gaze}
When a creature that can see the medusa's eyes starts its turn within 30 feet of the medusa, the medusa can force it to make a DC 14 Constitution saving throw if the medusa isn't incapacitated and can see the creature. If the saving throw fails by 5 or more, the creature is instantly petrified. Otherwise, a creature that fails the save begins to turn to stone and is restrained. The restrained creature must repeat the saving throw at the end of its next turn, becoming petrified on a failure or ending the effect on a success. The petrification lasts until the creature is freed by the  \textit{greater restoration} spell or other magic.

Unless Surprise, a creature can avert its eyes to avoid the saving throw at the start of its turn. If the creature does so, it can't see the medusa until the start of its next turn, when it can avert its eyes again. If the creature looks at the medusa in the meantime, it must immediately make the save.

If the medusa sees itself reflected on a polished surface within 30 feet of it and in an area of bright light, the medusa is, due to its curse, affected by its own gaze.

\DndMonsterSection{Actions}
\DndMonsterAction{Multiattack}
The medusa makes either three melee attacks—one with its snake hair and two with its shortsword—or two ranged attacks with its longbow.

\DndMonsterAction{Snake Hair}
\textsl{Melee Weapon Attack:} 5 to hit, reach 5 ft., one creature. \textsl{Hit:} 4 (1d4 + 2) piercing damage plus 14 (4d6) poison damage.

\DndMonsterAction{Shortsword}
\textsl{Melee Weapon Attack:} 5 to hit, reach 5 ft., one target. \textsl{Hit:} 5 (1d6 + 2) piercing damage.

\DndMonsterAction{Longbow}
\textsl{Ranged Weapon Attack:} 5 to hit, range 150/600 ft., one target. \textsl{Hit:} 6 (1d8 + 2) piercing damage plus 7 (2d6) poison damage.

\end{DndMonster}



\begin{DndMonster}[float=t]{Merrow}
\DndMonsterType{Large monstrosity, chaotic evil}
\DndMonsterBasics[
armor-class = {13 (natural armor)},
hit-points = {\DndDice{6d10 + 12}},
speed = {10 ft., swim 40 ft.}
]
\DndMonsterAbilityScores[
 str = 18,
 dex = 10,
 con = 15,
 int = 8,
 wis = 10,
 cha = 9
]
\DndMonsterDetails[
senses = {darkvision 60 ft., passive perception 10},
languages = {Abyssal, Aquan},
challenge = 2,
]
\DndMonsterAction{Amphibious}
The merrow can breathe air and water.

\DndMonsterSection{Actions}
\DndMonsterAction{Multiattack}
The merrow makes two attacks: one with its bite and one with its claws or harpoon.

\DndMonsterAction{Bite}
\textsl{Melee Weapon Attack:} 6 to hit, reach 5 ft., one target. \textsl{Hit:} 8 (1d8 + 4) piercing damage.

\DndMonsterAction{Claws}
\textsl{Melee Weapon Attack:} 6 to hit, reach 5 ft., one target. \textsl{Hit:} 9 (2d4 + 4) slashing damage.

\DndMonsterAction{Harpoon}
\textsl{Melee or Ranged Weapon Attack:} 6 to hit, reach 5 ft. or range 20/60 ft., one target. \textsl{Hit:} 11 (2d6 + 4) piercing damage. If the target is a Huge or smaller creature, it must succeed on a Strength contest against the merrow or be pulled up to 20 feet toward the merrow.

\end{DndMonster}



\begin{DndMonster}[float=t]{Mimic}
\DndMonsterType{Medium monstrosity (shapechanger), neutral}
\DndMonsterBasics[
armor-class = {12 (natural armor)},
hit-points = {\DndDice{9d8 + 18}},
speed = {15 ft.}
]
\DndMonsterAbilityScores[
 str = 17,
 dex = 12,
 con = 15,
 int = 5,
 wis = 13,
 cha = 8
]
\DndMonsterDetails[
skills = {Stealth +5},
damage-immunities = {acid},
condition-immunities = {prone},
senses = {darkvision 60 ft., passive perception 11},
challenge = 2,
]
\DndMonsterAction{Shapechanger}
The mimic can use its action to polymorph into an object or back into its true, amorphous form. Its statistics are the same in each form. Any equipment it is wearing or carrying isn't transformed. It reverts to its true form if it dies.

\DndMonsterAction{Adhesive (Object Form Only)}
The mimic adheres to anything that touches it. A Huge or smaller creature adhered to the mimic is also grappled by it (escape DC 13). Ability checks made to escape this grapple have disadvantage.

\DndMonsterAction{False Appearance (Object Form Only)}
While the mimic remains motionless, it is indistinguishable from an ordinary object.

\DndMonsterAction{Grappler}
The mimic has advantage on attack rolls against any creature grappled by it.

\DndMonsterSection{Actions}
\DndMonsterAction{Pseudopod}
\textsl{Melee Weapon Attack:} 5 to hit, reach 5 ft., one target. \textsl{Hit:} 7 (1d8 + 3) bludgeoning damage. If the mimic is in object form, the target is subjected to its Adhesive trait.

\DndMonsterAction{Bite}
\textsl{Melee Weapon Attack:} 5 to hit, reach 5 ft., one target. \textsl{Hit:} 7 (1d8 + 3) piercing damage plus 4 (1d8) acid damage.

\end{DndMonster}



\begin{DndMonster}[float*=b,width=\textwidth + 8pt]{Minazorond}
\begin{multicols}{2}
\DndMonsterType{Gargantuan dragon, lawful good}
\DndMonsterBasics[
armor-class = {22 (natural armor)},
hit-points = {\DndDice{24d20 + 192}},
speed = {40 ft., fly 80 ft., swim 40 ft.}
]
\DndMonsterAbilityScores[
 str = 29,
 dex = 10,
 con = 27,
 int = 18,
 wis = 17,
 cha = 21
]
\DndMonsterDetails[
saving-throws = {Dex +7, Con +15, Wis +10, Cha +12},
skills = {Insight +10, Perception +17, Stealth +7},
damage-resistances = {poison; psychic; bludgeoning, piercing, slashing from nonmagical attacks that aren't adamantine},
damage-immunities = {lightning},
condition-immunities = {charmed, exhaustion, frightened, paralyzed, petrified, poisoned},
senses = {blindsight 60 ft., darkvision 120 ft., truesight 120 ft., passive perception 27},
languages = {Common, Draconic},
challenge = 22,
]
\DndMonsterAction{Amphibious}
Minazorond can breathe air and water.

\DndMonsterAction{Constructed Nature}
Minazorond doesn't require air, food, drink, or sleep.

\DndMonsterAction{False Appearance}
While Minazorond remains motionless, he is indistinguishable from an inanimate bronze statue.

\DndMonsterAction{Faultless Tracker}
The creature(s) that caused Minazorond to animate becomes his quarry. Minazorond knows the direction and distance to his quarry as long the quarry remains within Drakkenheim or until Minazorond is destroyed.

\DndMonsterAction{Immutable Form}
Minazorond is immune to any spell or effect that would alter his form.

\DndMonsterAction{Legendary Resistance (3/Day)}
If Minazorond fails a saving throw, he can choose to succeed instead.

\DndMonsterAction{Magic Resistance}
Minazorond has advantage on saving throws against spells and other magical effects.

\DndMonsterAction{Magic Weapons}
Minazorond's weapon attacks are magical.

\DndMonsterAction{Rejuvenation}
If Minazorond is destroyed, he regains all his hit points in 1 hour unless a \textit{disintegrate} spell is cast on his remains.

\DndMonsterSection{Actions}
\DndMonsterAction{Multiattack}
Minazorond can use his Frightful Presence. He then makes three attacks: one with his bite and two with his claws.

\DndMonsterAction{Bite}
\textsl{Melee Weapon Attack:} 16 to hit, reach 15 ft., one target. \textsl{Hit:} 20 (2d10 + 9) piercing damage.

\DndMonsterAction{Claw}
\textsl{Melee Weapon Attack:} 16 to hit, reach 10 ft., one target. \textsl{Hit:} 16 (2d6 + 9) slashing damage.

\DndMonsterAction{Tail}
\textsl{Melee Weapon Attack:} 16 to hit, reach 20 ft., one target. \textsl{Hit:} 18 (2d8 + 9) bludgeoning damage.

\DndMonsterAction{Frightful Presence}
Each creature of Minazorond's choice that is within 120 feet of Minazorond and aware of him must succeed on a DC 20 Wisdom saving throw or become frightened for 1 minute. A creature can repeat the saving throw at the end of each of its turns, ending the effect on itself on a success. If a creature's saving throw is successful or the effect ends for it, the creature is immune to Minazorond's Frightful Presence for the next 24 hours.

\DndMonsterAction{Breath Weapons (Recharge 5\textendash 6)}
Minazorond uses one of the following breath weapons.

\DndMonsterAction{Lightning Breath}
Minazorond exhales lightning in a 120-foot line that is 10 feet wide. Each creature in that line must make a DC 23 Dexterity saving throw, taking 88 (16d10) lightning damage on a failed save, or half as much damage on a successful one.

\DndMonsterAction{Repulsion Breath}
Minazorond exhales repulsion energy in a 30-foot cone. Each creature in that area must succeed on a DC 23 Strength saving throw. On a failed save, the creature is pushed 60 feet away from Minazorond.

\DndMonsterAction{Change Shape}
Minazorond magically polymorphs into a humanoid or beast that has a challenge rating no higher than his own, or back into his true form. He reverts to his true form if he dies. Any equipment he is wearing or carrying is absorbed or borne by the new form (Minazorond's choice).

In a new form, Minazorond retains his alignment, hit points, Hit Dice, ability to speak, proficiencies, Legendary Resistance, lair actions, and Intelligence, Wisdom, and Charisma scores, as well as this action. His statistics and capabilities are otherwise replaced by those of the new form, except any class features or legendary actions of that form.

\DndMonsterSection{Legendary Actions}
Minazorond can take 3 legendary actions, choosing from the options below. Only one legendary action can be used at a time and only at the end of another creature's turn. Minazorond regains spent legendary actions at the start of its turn.

\begin{DndMonsterLegendaryActions}
\DndMonsterLegendaryAction{Detect}{Minazorond makes a Wisdom (Perception) check.
}
\DndMonsterLegendaryAction{Tail Attack}{Minazorond makes a tail attack.
}
\DndMonsterLegendaryAction{Wing Attack (Costs 2 Actions)}{Minazorond beats his wings. Each creature within 15 feet of Minazorond must succeed on a DC 24 Dexterity saving throw or take 16 (2d6 + 9) bludgeoning damage and be knocked prone. Minazorond can then fly up to half his flying speed.
}
\end{DndMonsterLegendaryActions}
\end{multicols}
\end{DndMonster}



\begin{DndMonster}[float=t]{Mule}
\DndMonsterType{Medium beast, unaligned}
\DndMonsterBasics[
armor-class = {10},
hit-points = {\DndDice{2d8 + 2}},
speed = {40 ft.}
]
\DndMonsterAbilityScores[
 str = 14,
 dex = 10,
 con = 13,
 int = 2,
 wis = 10,
 cha = 5
]
\DndMonsterDetails[
challenge = 1/8,
]
\DndMonsterAction{Beast of Burden}
The mule is considered to be a Large animal for the purpose of determining its carrying capacity.

\DndMonsterAction{Sure-Footed}
The mule has advantage on Strength and Dexterity saving throws made against effects that would knock it prone.

\DndMonsterSection{Actions}
\DndMonsterAction{Hooves}
\textsl{Melee Weapon Attack:} 2 to hit, reach 5 ft., one target. \textsl{Hit:} 4 (1d4 + 2) bludgeoning damage.

\end{DndMonster}



\begin{DndMonster}[float=t]{Mummy}
\DndMonsterType{Medium undead, lawful evil}
\DndMonsterBasics[
armor-class = {11 (natural armor)},
hit-points = {\DndDice{9d8 + 18}},
speed = {20 ft.}
]
\DndMonsterAbilityScores[
 str = 16,
 dex = 8,
 con = 15,
 int = 6,
 wis = 10,
 cha = 12
]
\DndMonsterDetails[
saving-throws = {Wis +2},
damage-vulnerabilities = {fire},
damage-resistances = {bludgeoning, piercing, slashing from nonmagical attacks},
damage-immunities = {necrotic, poison},
condition-immunities = {charmed, exhaustion, frightened, paralyzed, poisoned},
senses = {darkvision 60 ft., passive perception 10},
languages = {the languages it knew in life},
challenge = 3,
]
\DndMonsterSection{Actions}
\DndMonsterAction{Multiattack}
The mummy can use its Dreadful Glare and makes one attack with its rotting fist.

\DndMonsterAction{Rotting Fist}
\textsl{Melee Weapon Attack:} 5 to hit, reach 5 ft., one target. \textsl{Hit:} 10 (2d6 + 3) bludgeoning damage plus 10 (3d6) necrotic damage. If the target is a creature, it must succeed on a DC 12 Constitution saving throw or be cursed with mummy rot. The cursed target can't regain hit points, and its hit point maximum decreases by 10 (3d6) for every 24 hours that elapse. If the curse reduces the target's hit point maximum to 0, the target dies, and its body turns to dust. The curse lasts until removed by the \textit{remove curse} spell or other magic.

\DndMonsterAction{Dreadful Glare}
The mummy targets one creature it can see within 60 feet of it. If the target can see the mummy, it must succeed on a DC 11 Wisdom saving throw against this magic or become frightened until the end of the mummy's next turn. If the target fails the saving throw by 5 or more, it is also paralyzed for the same duration. A target that succeeds on the saving throw is immune to the Dreadful Glare of all mummies (but not mummy lords) for the next 24 hours.

\end{DndMonster}



\begin{DndMonster}[float=t]{Nalfeshnee}
\DndMonsterType{Large fiend (demon), chaotic evil}
\DndMonsterBasics[
armor-class = {18 (natural armor)},
hit-points = {\DndDice{16d10 + 96}},
speed = {20 ft., fly 30 ft.}
]
\DndMonsterAbilityScores[
 str = 21,
 dex = 10,
 con = 22,
 int = 19,
 wis = 12,
 cha = 15
]
\DndMonsterDetails[
saving-throws = {Con +11, Int +9, Wis +6, Cha +7},
damage-resistances = {cold; fire; lightning; bludgeoning, piercing, slashing from nonmagical attacks},
damage-immunities = {poison},
condition-immunities = {poisoned},
senses = {truesight 120 ft., passive perception 11},
languages = {Abyssal, telepathy 120 ft.},
challenge = 13,
]
\DndMonsterAction{Magic Resistance}
The nalfeshnee has advantage on saving throws against spells and other magical effects.

\DndMonsterSection{Actions}
\DndMonsterAction{Multiattack}
The nalfeshnee uses Horror Nimbus if it can. It then makes three attacks: one with its bite and two with its claws.

\DndMonsterAction{Bite}
\textsl{Melee Weapon Attack:} 10 to hit, reach 5 ft., one target. \textsl{Hit:} 32 (5d10 + 5) piercing damage.

\DndMonsterAction{Claw}
\textsl{Melee Weapon Attack:} 10 to hit, reach 10 ft., one target. \textsl{Hit:} 15 (3d6 + 5) slashing damage.

\DndMonsterAction{Horror Nimbus (Recharge 5\textendash 6)}
The nalfeshnee magically emits scintillating, multicolored light. Each creature within 15 feet of the nalfeshnee that can see the light must succeed on a DC 15 Wisdom saving throw or be frightened for 1 minute. A creature can repeat the saving throw at the end of each of its turns, ending the effect on itself on a success. If a creature's saving throw is successful or the effect ends for it, the creature is immune to the nalfeshnee's Horror Nimbus for the next 24 hours.

\DndMonsterAction{Teleport}
The nalfeshnee magically teleports, along with any equipment it is wearing or carrying, up to 120 feet to an unoccupied space it can see.

\end{DndMonster}



\begin{DndMonster}[float*=b,width=\textwidth + 8pt]{Night Hag}
\begin{multicols}{2}
\DndMonsterType{Medium fiend, neutral evil}
\DndMonsterBasics[
armor-class = {17 (natural armor)},
hit-points = {\DndDice{15d8 + 45}},
speed = {30 ft.}
]
\DndMonsterAbilityScores[
 str = 18,
 dex = 15,
 con = 16,
 int = 16,
 wis = 14,
 cha = 16
]
\DndMonsterDetails[
skills = {Deception +6, Insight +5, Perception +5, Stealth +5},
damage-resistances = {cold; fire; bludgeoning, piercing, slashing from nonmagical attacks that aren't silvered},
condition-immunities = {charmed},
senses = {darkvision 120 ft., passive perception 16},
languages = {Abyssal, Common, Infernal, Primordial},
challenge = 5,
]
\DndMonsterAction{Magic Resistance}
The hag has advantage on saving throws against spells and other magical effects.

\DndMonsterAction{Night Hag Items}
A night hag carries two very rare magic items that she must craft for herself. If either object is lost, the night hag will go to great lengths to retrieve it, as creating a new tool takes time and effort.

Heartstone: This lustrous black gem allows a night hag to become ethereal while it is in her possession. The touch of a heartstone also cures any disease. Crafting a heartstone takes 30 days.

Soul Bag: When an evil humanoid dies as a result of a night hag's Nightmare Haunting, the hag catches the soul in this black sack made of stitched flesh. A soul bag can hold only one evil soul at a time, and only the night hag who crafted the bag can catch a soul with it. Crafting a soul bag takes 7 days and a humanoid sacrifice (whose flesh is used to make the bag).

\DndMonsterAction{Innate Spellcasting}
The hag's innate spellcasting ability is Charisma (spell save DC 14, 6 to hit with spell attacks). She can innately cast the following spells, requiring no material components:
\begin{DndMonsterSpells}
  \DndInnateSpellLevel{\textit{detect magic}, \textit{magic missile}}
  \DndInnateSpellLevel[2e]{\textit{plane shift} (self only), \textit{ray of enfeeblement}, \textit{sleep}}
\end{DndMonsterSpells}
\DndMonsterSection{Actions}
\DndMonsterAction{Claws (Hag Form Only)}
\textsl{Melee Weapon Attack:} 7 to hit, reach 5 ft., one target. \textsl{Hit:} 13 (2d8 + 4) slashing damage.

\DndMonsterAction{Change Shape}
The hag magically polymorphs into a Small or Medium female humanoid, or back into her true form. Her statistics are the same in each form. Any equipment she is wearing or carrying isn't transformed. She reverts to her true form if she dies.

\DndMonsterAction{Etherealness}
The hag magically enters the Ethereal Plane from the Material Plane, or vice versa. To do so, the hag must have a heartstone in her possession.

\DndMonsterAction{Nightmare Haunting (1/Day)}
While on the Ethereal Plane, the hag magically touches a sleeping humanoid on the Material Plane. A \textit{protection from evil and good} spell cast on the target prevents this contact, as does a magic circle. As long as the contact persists, the target has dreadful visions. If these visions last for at least 1 hour, the target gains no benefit from its rest, and its hit point maximum is reduced by 5 (1d10). If this effect reduces the target's hit point maximum to 0, the target dies, and if the target was evil, its soul is trapped in the hag's soul bag. The reduction to the target's hit point maximum lasts until removed by the  \textit{greater restoration} spell or similar magic.

\end{multicols}
\end{DndMonster}



\begin{DndMonster}[float=t]{Noble}
\DndMonsterType{Medium humanoid (any race), any alignment}
\DndMonsterBasics[
armor-class = {15 (\textit{breastplate})},
hit-points = {\DndDice{2d8}},
speed = {30 ft.}
]
\DndMonsterAbilityScores[
 str = 11,
 dex = 12,
 con = 11,
 int = 12,
 wis = 14,
 cha = 16
]
\DndMonsterDetails[
skills = {Deception +5, Insight +4, Persuasion +5},
languages = {any two languages},
challenge = 1/8,
]
\DndMonsterSection{Actions}
\DndMonsterAction{Rapier}
\textsl{Melee Weapon Attack:} 3 to hit, reach 5 ft., one target. \textsl{Hit:} 5 (1d8 + 1) piercing damage.

\DndMonsterSection{Reactions}
\DndMonsterAction{Parry}
The noble adds 2 to its AC against one melee attack that would hit it. To do so, the noble must see the attacker and be wielding a melee weapon.

\end{DndMonster}



\begin{DndMonster}[float=t]{Ochre Jelly}
\DndMonsterType{Large ooze, unaligned}
\DndMonsterBasics[
armor-class = {8},
hit-points = {\DndDice{6d10 + 12}},
speed = {10 ft., climb 10 ft.}
]
\DndMonsterAbilityScores[
 str = 15,
 dex = 6,
 con = 14,
 int = 2,
 wis = 6,
 cha = 1
]
\DndMonsterDetails[
damage-resistances = {acid},
damage-immunities = {lightning, slashing},
condition-immunities = {blinded, charmed, deafened, exhaustion, frightened, prone},
senses = {blindsight 60 ft. (blind beyond this radius), passive perception 8},
challenge = 2,
]
\DndMonsterAction{Amorphous}
The jelly can move through a space as narrow as 1 inch wide without squeezing.

\DndMonsterAction{Spider Climb}
The jelly can climb difficult surfaces, including upside down on ceilings, without needing to make an ability check.

\DndMonsterSection{Actions}
\DndMonsterAction{Pseudopod}
\textsl{Melee Weapon Attack:} 4 to hit, reach 5 ft., one target. \textsl{Hit:} 9 (2d6 + 2) bludgeoning damage plus 3 (1d6) acid damage.

\DndMonsterSection{Reactions}
\DndMonsterAction{Split}
When a jelly that is Medium or larger is subjected to lightning or slashing damage, it splits into two new jellies if it has at least 10 hit points. Each new jelly has hit points equal to half the original jelly's, rounded down. New jellies are one size smaller than the original jelly.

\end{DndMonster}



\begin{DndMonster}[float=t]{Ogre}
\DndMonsterType{Large giant, chaotic evil}
\DndMonsterBasics[
armor-class = {11 (\textit{hide armor})},
hit-points = {\DndDice{7d10 + 21}},
speed = {40 ft.}
]
\DndMonsterAbilityScores[
 str = 19,
 dex = 8,
 con = 16,
 int = 5,
 wis = 7,
 cha = 7
]
\DndMonsterDetails[
senses = {darkvision 60 ft., passive perception 8},
languages = {Common, Giant},
challenge = 2,
]
\DndMonsterSection{Actions}
\DndMonsterAction{Greatclub}
\textsl{Melee Weapon Attack:} 6 to hit, reach 5 ft., one target. \textsl{Hit:} 13 (2d8 + 4) bludgeoning damage.

\DndMonsterAction{Javelin}
\textsl{Melee or Ranged Weapon Attack:} 6 to hit, reach 5 ft. or range 30/120 ft., one target. \textsl{Hit:} 11 (2d6 + 4) piercing damage.

\end{DndMonster}



\begin{DndMonster}[float=t]{Ogre Zombie}
\DndMonsterType{Large undead, neutral evil}
\DndMonsterBasics[
armor-class = {8},
hit-points = {\DndDice{9d10 + 36}},
speed = {30 ft.}
]
\DndMonsterAbilityScores[
 str = 19,
 dex = 6,
 con = 18,
 int = 3,
 wis = 6,
 cha = 5
]
\DndMonsterDetails[
saving-throws = {Wis +0},
damage-immunities = {poison},
condition-immunities = {poisoned},
senses = {darkvision 60 ft., passive perception 8},
languages = {understands Common and Giant but can't speak},
challenge = 2,
]
\DndMonsterAction{Undead Fortitude}
If damage reduces the zombie to 0 hit points, it must make a Constitution saving throw with a DC of 5 + the damage taken, unless the damage is radiant or from a critical hit. On a success, the zombie drops to 1 hit point instead.

\DndMonsterSection{Actions}
\DndMonsterAction{Morningstar}
\textsl{Melee Weapon Attack:} 6 to hit, reach 5 ft., one target. \textsl{Hit:} 13 (2d8 + 4) bludgeoning damage.

\end{DndMonster}


\clearpage

\begin{DndMonster}[float=t]{Oscar Yoren}
\DndMonsterType{Medium humanoid (human), any alignment}
\DndMonsterBasics[
armor-class = {12 (15 with \textit{mage armor})},
hit-points = {\DndDice{9d8}},
speed = {30 ft.}
]
\DndMonsterAbilityScores[
 str = 9,
 dex = 14,
 con = 11,
 int = 17,
 wis = 12,
 cha = 11
]
\DndMonsterDetails[
saving-throws = {Int +6, Wis +4},
skills = {Arcana +6, History +6},
damage-resistances = {necrotic, poison},
senses = {darkvision 60 ft., passive perception 11},
languages = {Any four languages},
challenge = 6,
]
\DndMonsterAction{Spellcasting}
Oscar Yoren is a 9th-level spellcaster. His spellcasting ability is Intelligence (spell save DC 14, 6 to hit with spell attacks). Oscar Yoren has the following wizard spells prepared:
\begin{DndMonsterSpells}
  \DndMonsterSpellLevel{\textit{chill touch}, \textit{mage hand}, \textit{ray of frost}}
   \DndMonsterSpellLevel[1][4]{\textit{false life}, \textit{mage armor}, \textit{magic missile}, \textit{shield}}
   \DndMonsterSpellLevel[2][3]{\textit{invisibility}, \textit{misty step}, \textit{web}}
   \DndMonsterSpellLevel[3][3]{\textit{animate dead}, \textit{counterspell}, \textit{lightning bolt}}
   \DndMonsterSpellLevel[4][3]{\textit{blight}, \textit{polymorph}}
   \DndMonsterSpellLevel[5][1]{\textit{cloudkill}}
\end{DndMonsterSpells}
\DndMonsterSection{Actions}
\DndMonsterAction{Dagger}
\textsl{Melee or Ranged Weapon Attack:} 5 to hit, reach 5 ft. or range 20/60 ft., one target. \textsl{Hit:} 4 (1d4 + 2) piercing damage.

\end{DndMonster}



\begin{DndMonster}[float=t]{Otyugh}
\DndMonsterType{Large aberration, neutral}
\DndMonsterBasics[
armor-class = {14 (natural armor)},
hit-points = {\DndDice{12d10 + 48}},
speed = {30 ft.}
]
\DndMonsterAbilityScores[
 str = 16,
 dex = 11,
 con = 19,
 int = 6,
 wis = 13,
 cha = 6
]
\DndMonsterDetails[
saving-throws = {Con +7},
senses = {darkvision 120 ft., passive perception 11},
languages = {Otyugh},
challenge = 5,
]
\DndMonsterAction{Limited Telepathy}
The otyugh can magically transmit simple messages and images to any creature within 120 feet of it that can understand a language. This form of telepathy doesn't allow the receiving creature to telepathically respond.

\DndMonsterSection{Actions}
\DndMonsterAction{Multiattack}
The otyugh makes three attacks: one with its bite and two with its tentacles.

\DndMonsterAction{Bite}
\textsl{Melee Weapon Attack:} 6 to hit, reach 5 ft., one target. \textsl{Hit:} 12 (2d8 + 3) piercing damage. If the target is a creature, it must succeed on a DC 15 Constitution saving throw against disease or become poisoned until the disease is cured. Every 24 hours that elapse, the target must repeat the saving throw, reducing its hit point maximum by 5 (1d10) on a failure. The disease is cured on a success. The target dies if the disease reduces its hit point maximum to 0. This reduction to the target's hit point maximum lasts until the disease is cured.

\DndMonsterAction{Tentacle}
\textsl{Melee Weapon Attack:} 6 to hit, reach 10 ft., one target. \textsl{Hit:} 7 (1d8 + 3) bludgeoning damage plus 4 (1d8) piercing damage. If the target is Medium or smaller, it is grappled (escape DC 13) and restrained until the grapple ends. The otyugh has two tentacles, each of which can grapple one target.

\DndMonsterAction{Tentacle Slam}
The otyugh slams creatures grappled by it into each other or a solid surface. Each creature must succeed on a DC 14 Constitution saving throw or take 10 (2d6 + 3) bludgeoning damage and be stunned until the end of the otyugh's next turn. On a successful save, the target takes half the bludgeoning damage and isn't stunned.

\end{DndMonster}



\begin{DndMonster}[float=t]{Owl}
\DndMonsterType{Tiny beast, unaligned}
\DndMonsterBasics[
armor-class = {11},
hit-points = {\DndDice{1d4 - 1}},
speed = {5 ft., fly 60 ft.}
]
\DndMonsterAbilityScores[
 str = 3,
 dex = 13,
 con = 8,
 int = 2,
 wis = 12,
 cha = 7
]
\DndMonsterDetails[
skills = {Perception +3, Stealth +3},
senses = {darkvision 120 ft., passive perception 13},
challenge = 0,
]
\DndMonsterAction{Flyby}
The owl doesn't provoke opportunity attacks when it flies out of an enemy's reach.

\DndMonsterAction{Keen Hearing and Sight}
The owl has advantage on Wisdom (Perception) checks that rely on hearing or sight.

\DndMonsterSection{Actions}
\DndMonsterAction{Talons}
\textsl{Melee Weapon Attack:} 3 to hit, reach 5 ft., one target. \textsl{Hit:} 1 slashing damage.

\end{DndMonster}



\begin{DndMonster}[float*=b,width=\textwidth + 8pt]{Pale Man}
\begin{multicols}{2}
\DndMonsterType{Medium aberration, neutral}
\DndMonsterBasics[
armor-class = {12 (15 with \textit{mage armor})},
hit-points = {\DndDice{16d8 + 48}},
speed = {30 ft., climb 30 ft.}
]
\DndMonsterAbilityScores[
 str = 10,
 dex = 15,
 con = 17,
 int = 20,
 wis = 15,
 cha = 12
]
\DndMonsterDetails[
saving-throws = {Con +8, Int +10, Wis +7},
skills = {Arcana +10, History +10, Perception +7},
damage-resistances = {bludgeoning, piercing, slashing from nonmagical attacks},
damage-immunities = {necrotic, poison},
condition-immunities = {charmed, frightened, paralyzed, petrified, poisoned, restrained},
senses = {darkvision 60 ft., passive perception 15},
languages = {Common and any 5 other languages},
challenge = 13,
]
\DndMonsterAction{Legendary Resistance (3/Day)}
If the Pale Man fails a saving throw, it can choose to succeed instead.

\DndMonsterAction{Insect Body}
The Pale Man can occupy another creature's space and vice versa, and the Pale Man can move through any opening large enough for a Tiny insect.

\DndMonsterAction{Crawling Flesh}
When the Pale Man is reduced to 0 hit points, it breaks apart into a swarm of insects in the same space. Unless the swarm is destroyed, the Pale Man swarm returns to its cocoon and reforms 1d6 days later.

\DndMonsterAction{Spellcasting}
The Pale Man is a 16th-level spellcaster. Its spellcasting ability is Intelligence (spell save DC 18, 10 to hit with spell attacks). The Pale Man has the following wizard spells prepared:
\begin{DndMonsterSpells}
  \DndMonsterSpellLevel{\textit{chill touch}, \textit{mage hand}, \textit{mending}, \textit{prestidigitation}, \textit{shocking grasp}}
   \DndMonsterSpellLevel[1][4]{\textit{delerium orb}*, \textit{mage armor}, \textit{magic missile}, \textit{shield}}
   \DndMonsterSpellLevel[2][3]{\textit{hold person}, \textit{invisibility}, \textit{warp bolt}*}
   \DndMonsterSpellLevel[3][3]{\textit{counterspell}, \textit{lightning bolt}, \textit{purge contamination}*}
   \DndMonsterSpellLevel[4][3]{\textit{delerium blast}*, \textit{forced evolution}*}
   \DndMonsterSpellLevel[5][3]{\textit{conjure the deep haze}*}
   \DndMonsterSpellLevel[6][1]{\textit{chain lightning}, \textit{siphon contamination}*}
   \DndMonsterSpellLevel[7][1]{\textit{octarine spray}*}
   \DndMonsterSpellLevel[8][1]{\textit{antimagic field}}
\end{DndMonsterSpells}
\DndMonsterSection{Actions}
\DndMonsterAction{Contaminating Touch}
\textsl{Melee Spell Attack:} 8 to hit, reach 10 ft., one target. \textsl{Hit:} 10 (3d6) necrotic damage, and the target must succeed on a DC 16 Constitution saving throw or receive one level of contamination.

\DndMonsterAction{Infest the Flesh (Recharge 6)}
Each creature within 30 feet of the Pale Man must succeed on a DC 16 Dexterity saving throw or take 22 (5d8) necrotic damage and be restrained by masses of swarming fleshy bugs. The affected creature takes 22 (5d8) necrotic damage at the start of each of the Pale Man's turns. The creature can repeat the saving throw at the end of each of its turns, ending the effect on itself on a success.

\DndMonsterSection{Legendary Actions}
Pale Man can take 3 legendary actions, choosing from the options below. Only one legendary action can be used at a time and only at the end of another creature's turn. Pale Man regains spent legendary actions at the start of its turn.

\begin{DndMonsterLegendaryActions}
\DndMonsterLegendaryAction{Cantrip (Costs 1 Action)}{The Pale Man casts one cantrip.
}
\DndMonsterLegendaryAction{Slither (Costs 2 Actions)}{The Pale Man moves its speed without provoking opportunity attacks.
}
\DndMonsterLegendaryAction{Contaminate (Costs 3 Actions)}{Each creature restrained by the Pale Man's Infest the Flesh must succeed on a DC 16 Constitution saving throw or gain one level of contamination.
}
\end{DndMonsterLegendaryActions}
\end{multicols}
\end{DndMonster}



\begin{DndMonster}[float=t]{Phase Spider}
\DndMonsterType{Large monstrosity, unaligned}
\DndMonsterBasics[
armor-class = {13 (natural armor)},
hit-points = {\DndDice{5d10 + 5}},
speed = {30 ft., climb 30 ft.}
]
\DndMonsterAbilityScores[
 str = 15,
 dex = 15,
 con = 12,
 int = 6,
 wis = 10,
 cha = 6
]
\DndMonsterDetails[
skills = {Stealth +6},
senses = {darkvision 60 ft., passive perception 10},
challenge = 3,
]
\DndMonsterAction{Ethereal Jaunt}
As a bonus action, the spider can magically shift from the Material Plane to the Ethereal Plane, or vice versa.

\DndMonsterAction{Spider Climb}
The spider can climb difficult surfaces, including upside down on ceilings, without needing to make an ability check.

\DndMonsterAction{Web Walker}
The spider ignores movement restrictions caused by webbing.

\DndMonsterSection{Actions}
\DndMonsterAction{Bite}
\textsl{Melee Weapon Attack:} 4 to hit, reach 5 ft., one creature. \textsl{Hit:} 7 (1d10 + 2) piercing damage, and the target must make a DC 11 Constitution saving throw, taking 18 (4d8) poison damage on a failed save, or half as much damage on a successful one. If the poison damage reduces the target to 0 hit points, the target is stable but poisoned for 1 hour, even after regaining hit points, and is paralyzed while poisoned in this way.

\end{DndMonster}



\begin{DndMonster}[float=t]{Planetar}
\DndMonsterType{Large celestial, lawful good}
\DndMonsterBasics[
armor-class = {19 (natural armor)},
hit-points = {\DndDice{16d10 + 112}},
speed = {40 ft., fly 120 ft.}
]
\DndMonsterAbilityScores[
 str = 24,
 dex = 20,
 con = 24,
 int = 19,
 wis = 22,
 cha = 25
]
\DndMonsterDetails[
saving-throws = {Con +12, Wis +11, Cha +12},
skills = {Perception +11},
damage-resistances = {radiant; bludgeoning, piercing, slashing from nonmagical attacks},
condition-immunities = {charmed, exhaustion, frightened},
senses = {truesight 120 ft., passive perception 21},
languages = {all, telepathy 120 ft.},
challenge = 16,
]
\DndMonsterAction{Angelic Weapons}
The planetar's weapon attacks are magical. When the planetar hits with any weapon, the weapon deals an extra 5d8 radiant damage (included in the attack).

\DndMonsterAction{Divine Awareness}
The planetar knows if it hears a lie.

\DndMonsterAction{Magic Resistance}
The planetar has advantage on saving throws against spells and other magical effects.

\DndMonsterAction{Innate Spellcasting}
The planetar's spellcasting ability is Charisma (spell save DC 20). The planetar can innately cast the following spells, requiring no material components:
\begin{DndMonsterSpells}
  \DndInnateSpellLevel{\textit{detect evil and good}, \textit{invisibility} (self only)}
  \DndInnateSpellLevel[3e]{\textit{blade barrier}, \textit{dispel evil and good}, \textit{flame strike}, \textit{raise dead}}
  \DndInnateSpellLevel[1e]{\textit{commune}, \textit{control weather}, \textit{insect plague}}
\end{DndMonsterSpells}
\DndMonsterSection{Actions}
\DndMonsterAction{Multiattack}
The planetar makes two melee attacks.

\DndMonsterAction{Greatsword}
\textsl{Melee Weapon Attack:} 12 to hit, reach 5 ft., one target. \textsl{Hit:} 21 (4d6 + 7) slashing damage plus 22 (5d8) radiant damage.

\DndMonsterAction{Healing Touch (4/Day)}
The planetar touches another creature. The target magically regains 30 (6d8 + 3) hit points and is freed from any curse, disease, poison, blindness, or deafness.

\end{DndMonster}



\begin{DndMonster}[float=t]{Priest}
\DndMonsterType{Medium humanoid (any race), any alignment}
\DndMonsterBasics[
armor-class = {13 (\textit{chain shirt})},
hit-points = {\DndDice{5d8 + 5}},
speed = {30 ft.}
]
\DndMonsterAbilityScores[
 str = 10,
 dex = 10,
 con = 12,
 int = 13,
 wis = 16,
 cha = 13
]
\DndMonsterDetails[
skills = {Medicine +7, Persuasion +3, Religion +5},
languages = {any two languages},
challenge = 2,
]
\DndMonsterAction{Divine Eminence}
As a bonus action, the priest can expend a spell slot to cause its melee weapon attacks to magically deal an extra 10 (3d6) radiant damage to a target on a hit. This benefit lasts until the end of the turn. If the priest expends a spell slot of 2nd level or higher, the extra damage increases by 1d6 for each level above 1st.

\DndMonsterAction{Spellcasting}
The priest is a 5th-level spellcaster. Its spellcasting ability is Wisdom (spell save DC 13, 5 to hit with spell attacks). The priest has the following cleric spells prepared:
\begin{DndMonsterSpells}
  \DndMonsterSpellLevel{\textit{light}, \textit{sacred flame}, \textit{thaumaturgy}}
   \DndMonsterSpellLevel[1][4]{\textit{cure wounds}, \textit{guiding bolt}, \textit{sanctuary}}
   \DndMonsterSpellLevel[2][3]{\textit{lesser restoration}, \textit{spiritual weapon}}
   \DndMonsterSpellLevel[3][2]{\textit{dispel magic}, \textit{spirit guardians}}
\end{DndMonsterSpells}
\DndMonsterSection{Actions}
\DndMonsterAction{Mace}
\textsl{Melee Weapon Attack:} 2 to hit, reach 5 ft., one target. \textsl{Hit:} 3 (1d6) bludgeoning damage.

\end{DndMonster}



\begin{DndMonster}[float=t]{Protean Abomination}
\DndMonsterType{Huge aberration, unaligned}
\DndMonsterBasics[
armor-class = {14 (natural armor)},
hit-points = {\DndDice{16d12 + 64}},
speed = {30 ft., climb 10 ft.}
]
\DndMonsterAbilityScores[
 str = 21,
 dex = 8,
 con = 18,
 int = 5,
 wis = 10,
 cha = 5
]
\DndMonsterDetails[
condition-immunities = {charmed, frightened, prone},
senses = {blindsight 60 ft., passive perception 10},
challenge = 6,
]
\DndMonsterAction{Absorption}
When the protean abomination kills a creature that is neither a construct nor undead, it absorbs the corpse into its body and regains 20 hit points.

\DndMonsterAction{Amorphous}
The protean abomination can move through a space as narrow as 1 inch wide without squeezing.

\DndMonsterAction{Spider Climb}
The protean abomination can climb difficult surfaces, including upside down on ceilings, without needing to make an ability check.

\DndMonsterSection{Actions}
\DndMonsterAction{Multiattack}
The protean abomination makes two slam attacks. If both attacks hit a Large or smaller target, the target is grappled (escape DC 15), and the protean abomination uses Engulf on it.

\DndMonsterAction{Slam}
\textsl{Melee Weapon Attack:} 8 to hit, reach 5 ft., one target. \textsl{Hit:} 19 (3d8 + 5) bludgeoning damage.

\DndMonsterAction{Engulf}
The protean abomination engulfs a Large or smaller creature grappled by it. Until the grapple ends the target is blinded, restrained, and unable to breathe, and it must succeed on a DC 15 Constitution saving throw at the start of each of the abomination's turns or take 19 (3d8 + 5) bludgeoning damage. If the abomination moves, the engulfed target moves with it. The abomination can have only one creature engulfed at a time.

\end{DndMonster}



\begin{DndMonster}[float=t]{Pseudodragon}
\DndMonsterType{Tiny dragon, neutral good}
\DndMonsterBasics[
armor-class = {13},
hit-points = {\DndDice{2d4 + 2}},
speed = {15 ft., fly 60 ft.}
]
\DndMonsterAbilityScores[
 str = 6,
 dex = 15,
 con = 13,
 int = 10,
 wis = 12,
 cha = 10
]
\DndMonsterDetails[
skills = {Perception +3, Stealth +4},
senses = {blindsight 10 ft., darkvision 60 ft., passive perception 13},
languages = {understands Common and Draconic but can't speak},
challenge = 1/4,
]
\DndMonsterAction{Keen Senses}
The pseudodragon has advantage on Wisdom (Perception) checks that rely on sight, hearing, or smell.

\DndMonsterAction{Magic Resistance}
The pseudodragon has advantage on saving throws against spells and other magical effects.

\DndMonsterAction{Limited Telepathy}
The pseudodragon can magically communicate simple ideas, emotions, and images telepathically with any creature within 100 feet of it that can understand a language.

\DndMonsterSection{Actions}
\DndMonsterAction{Bite}
\textsl{Melee Weapon Attack:} 4 to hit, reach 5 ft., one target. \textsl{Hit:} 4 (1d4 + 2) piercing damage.

\DndMonsterAction{Sting}
\textsl{Melee Weapon Attack:} 4 to hit, reach 5 ft., one creature. \textsl{Hit:} 4 (1d4 + 2) piercing damage, and the target must succeed on a DC 11 Constitution saving throw or become poisoned for 1 hour. If the saving throw fails by 5 or more, the target falls unconscious for the same duration, or until it takes damage or another creature uses an action to shake it awake.

\end{DndMonster}



\begin{DndMonster}[float=t]{Quasit}
\DndMonsterType{Tiny fiend (demon), chaotic evil}
\DndMonsterBasics[
armor-class = {13},
hit-points = {\DndDice{3d4}},
speed = {40 ft.}
]
\DndMonsterAbilityScores[
 str = 5,
 dex = 17,
 con = 10,
 int = 7,
 wis = 10,
 cha = 10
]
\DndMonsterDetails[
skills = {Stealth +5},
damage-resistances = {cold; fire; lightning; bludgeoning, piercing, slashing from nonmagical attacks},
damage-immunities = {poison},
condition-immunities = {poisoned},
senses = {darkvision 120 ft., passive perception 10},
languages = {Abyssal, Common},
challenge = 1,
]
\DndMonsterAction{Shapechanger}
The quasit can use its action to polymorph into a beast form that resembles a bat (speed 10 feet fly 40 ft.), a centipede (40 ft., climb 40 ft.), or a toad (40 ft., swim 40 ft.), or back into its true form. Its statistics are the same in each form, except for the speed changes noted. Any equipment it is wearing or carrying isn't transformed. It reverts to its true form if it dies.

\DndMonsterAction{Magic Resistance}
The quasit has advantage on saving throws against spells and other magical effects.

\DndMonsterSection{Actions}
\DndMonsterAction{Claw (Bite in Beast Form)}
\textsl{Melee Weapon Attack:} 4 to hit, reach 5 ft., one target. \textsl{Hit:} 5 (1d4 + 3) piercing damage, and the target must succeed on a DC 10 Constitution saving throw or take 5 (2d4) poison damage and become poisoned for 1 minute. The target can repeat the saving throw at the end of each of its turns, ending the effect on itself on a success.

\DndMonsterAction{Scare (1/Day)}
One creature of the quasit's choice within 20 feet of it must succeed on a DC 10 Wisdom saving throw or be frightened for 1 minute. The target can repeat the saving throw at the end of each of its turns, with disadvantage if the quasit is within line of sight, ending the effect on itself on a success.

\DndMonsterAction{Invisibility}
The quasit magically turns invisible until it attacks or uses Scare, or until its concentration ends (as if concentration on a spell). Any equipment the quasit wears or carries is invisible with it.

\end{DndMonster}



\begin{DndMonster}[float*=b,width=\textwidth + 8pt]{Queen of Thieves}
\begin{multicols}{2}
\DndMonsterType{Medium humanoid, neutral evil}
\DndMonsterBasics[
armor-class = {23 (canny defense)},
hit-points = {\DndDice{20d8 + 20}},
speed = {30 ft.}
]
\DndMonsterAbilityScores[
 str = 10,
 dex = 20,
 con = 13,
 int = 20,
 wis = 16,
 cha = 19
]
\DndMonsterDetails[
saving-throws = {Dex +10, Int +10},
skills = {Acrobatics +10, Arcana +10, Deception +14, Insight +8, Perception +8, Persuasion +9, Sleight Of Hand +10, Stealth +15},
condition-immunities = {charmed},
senses = {darkvision 60 ft., passive perception 18},
languages = {Common, Dwarvish, Elvish, Undercommon},
challenge = 15,
]
\DndMonsterAction{Special Equipment}
The Queen of Thieves wears an \textit{amulet of proof against detection and location} and a \textit{spymaster's signet}. She wields a \textit{+3 rapier}.

\DndMonsterAction{Canny Defense}
While the Queen of Thieves is wearing light or no armour and wielding no shield, her AC includes her Intelligence modifier.

\DndMonsterAction{Sneak Attack (1/Turn)}
The Queen of Thieves deals an extra 18 5d6 damage when she hits a target with a weapon attack and has advantage on the attack roll, or when the target is within 5 feet of an ally of the Queen of Thieves that isn't incapacitated and she doesn't have disadvantage on the attack roll.

\DndMonsterAction{Legendary Resistance (3/Day)}
If the Queen of Thieves fails a saving throw, she can choose to succeed instead.

\DndMonsterAction{Memory Thief}
Creatures are not aware when the Queen of Thieves reads their thoughts or charms them. Additionally, when a spell cast by the Queen of Thieves that causes creatures to become charmed ends, she can make the affected creatures lose their memories. They must succeed on a DC 18 Intelligence saving throw or forget up to 8 hours of the time spent charmed by her. These memories can be restored with a \textit{heal} spell or similar magic.

\DndMonsterAction{Innate Spellcasting}
The Queen of Thieves' innate spellcasting ability is Intelligence. She can innately cast the following spells, requiring no material, somatic, or verbal components:
\begin{DndMonsterSpells}
  \DndInnateSpellLevel{\textit{detect thoughts}, \textit{disguise self}, \textit{suggestion}}
\end{DndMonsterSpells}
\DndMonsterAction{Spellcasting}
The Queen of Thieves is a 20th-level spellcaster. Her spellcasting ability is Intelligence (spell save DC 18, 10 to hit with spell attacks). She has the following wizard spells prepared:
\begin{DndMonsterSpells}
  \DndMonsterSpellLevel{\textit{chill touch}, \textit{mage hand}, \textit{message}, \textit{minor illusion}, \textit{prestidigitation}}
   \DndMonsterSpellLevel[1][4]{\textit{feather fall}, \textit{mage armor}, \textit{shield}, \textit{sleep}}
   \DndMonsterSpellLevel[2][3]{\textit{mirror image}, \textit{misty step}, \textit{see invisibility}}
   \DndMonsterSpellLevel[3][3]{\textit{gaseous form}, \textit{hypnotic pattern}, \textit{major image}}
   \DndMonsterSpellLevel[4][3]{\textit{dimension door}, \textit{greater invisibility}, \textit{polymorph}}
   \DndMonsterSpellLevel[5][3]{\textit{dominate person}, \textit{scrying}, \textit{seeming}}
   \DndMonsterSpellLevel[6][2]{\textit{eyebite}, \textit{mass suggestion}}
   \DndMonsterSpellLevel[7][2]{\textit{forcecage}, \textit{project image}}
   \DndMonsterSpellLevel[8][1]{\textit{feeblemind}}
\end{DndMonsterSpells}
\DndMonsterSection{Actions}
\DndMonsterAction{+3 Rapier}
\textsl{Melee Weapon Attack:} 13 to hit, reach 5 ft., one target \textsl{Hit:} 13 (1d8 + 8) piercing damage.

\DndMonsterAction{Dagger}
\textsl{Ranged Weapon Attack:} 10 to hit, range 20/60 ft., one target. \textsl{Hit:} 7 (1d4 + 7) piercing damage.

\DndMonsterSection{Reactions}
\DndMonsterAction{Misdirection}
When she would take damage, the Queen of Thieves becomes invisible and teleports 60 feet to an unoccupied space she can see. At the same time, an illusory double of her appears where she was standing and lasts for 1 minute. The invisibility ends if the Queen of Thieves attacks or casts a spell, or after 1 minute.

\DndMonsterSection{Legendary Actions}
Queen of Thieves can take 3 legendary actions, choosing from the options below. Only one legendary action can be used at a time and only at the end of another creature's turn. Queen of Thieves regains spent legendary actions at the start of its turn.

\begin{DndMonsterLegendaryActions}
\DndMonsterLegendaryAction{Attack}{The Queen of Thieves makes one melee or ranged attack.
}
\DndMonsterLegendaryAction{Move}{The Queen of Thieves moves up to her speed without provoking opportunity attacks.
}
\end{DndMonsterLegendaryActions}
\end{multicols}
\end{DndMonster}



\begin{DndMonster}[float=t]{Rat}
\DndMonsterType{Tiny beast, unaligned}
\DndMonsterBasics[
armor-class = {10},
hit-points = {\DndDice{1d4 - 1}},
speed = {20 ft.}
]
\DndMonsterAbilityScores[
 str = 2,
 dex = 11,
 con = 9,
 int = 2,
 wis = 10,
 cha = 4
]
\DndMonsterDetails[
senses = {darkvision 30 ft., passive perception 10},
challenge = 0,
]
\DndMonsterAction{Keen Smell}
The rat has advantage on Wisdom (Perception) checks that rely on smell.

\DndMonsterSection{Actions}
\DndMonsterAction{Bite}
\textsl{Melee Weapon Attack:} 0 to hit, reach 5 ft., one target. \textsl{Hit:} 1 piercing damage.

\end{DndMonster}



\begin{DndMonster}[float=t]{Rat Prince}
\DndMonsterType{Medium monstrosity, chaotic evil}
\DndMonsterBasics[
armor-class = {15 (\textit{leather armor})},
hit-points = {\DndDice{12d8 + 12}},
speed = {30 ft., swim 30 ft., climb 30 ft.}
]
\DndMonsterAbilityScores[
 str = 12,
 dex = 19,
 con = 13,
 int = 9,
 wis = 15,
 cha = 9
]
\DndMonsterDetails[
skills = {Acrobatics +6, Athletics +3, Perception +6, Stealth +8},
senses = {darkvision 60 ft., passive perception 15},
languages = {Common},
challenge = 3,
]
\DndMonsterAction{Cunning Action}
The Rat Prince can take the Dash, Disengage, or Hide action as a bonus action on each of its turns.

\DndMonsterAction{Keen Smell}
The Rat Prince has advantage on Wisdom (Perception) checks that rely on smell.

\DndMonsterAction{Pack Tactics}
The Rat Prince has advantage on an attack roll against a creature if at least one of his allies is within 5 feet of the creature and the ally isn't incapacitated.

\DndMonsterAction{Contaminated Ammunition}
The Rat Prince carries 2d6 pieces of ammunition laced with delerium fragments. These projectiles cause 1d6 extra damage on a hit, and the target must succeed on a DC 10 Constitution saving throw or gain one level of contamination.

\DndMonsterSection{Actions}
\DndMonsterAction{Multiattack}
The Rat Prince makes three attacks: two with his short swords, and one with his contaminated bite.

\DndMonsterAction{Short Sword}
\textsl{Melee Weapon Attack:} 6 to hit, reach 5 ft., one target. \textsl{Hit:} 7 (1d6 + 4) piercing damage.

\DndMonsterAction{Contaminated Bite}
\textsl{Melee Weapon Attack:} 6 to hit, reach 5 ft., one target. \textsl{Hit:} 4 (1d4 + 4) piercing damage, plus an extra 10 (3d6) necrotic damage, and the creature must succeed on a DC 11 Constitution saving throw or gain one level of contamination.

\DndMonsterAction{Hand Crossbow}
\textsl{Ranged Weapon Attack:} 6 to hit, range 30/120 ft., one target. \textsl{Hit:} 7 (1d6 + 4) piercing damage.

\end{DndMonster}



\begin{DndMonster}[float=t]{Ratling}
\DndMonsterType{Small monstrosity, chaotic evil}
\DndMonsterBasics[
armor-class = {12},
hit-points = {\DndDice{2d6 - 2}},
speed = {30 ft., swim 30 ft., climb 30 ft.}
]
\DndMonsterAbilityScores[
 str = 7,
 dex = 15,
 con = 8,
 int = 8,
 wis = 10,
 cha = 5
]
\DndMonsterDetails[
skills = {Stealth +6},
senses = {darkvision 60 ft., passive perception 10},
languages = {Common},
challenge = 1/8,
]
\DndMonsterAction{Skulker}
The ratling can take the Hide action as a bonus action on each of its turns.

\DndMonsterAction{Keen Smell}
The ratling has advantage on Wisdom (Perception) checks that rely on smell.

\DndMonsterAction{Pack Tactics}
The ratling has advantage on an attack roll against a creature if at least one of the ratling's allies is within 5 feet of the creature and the ally isn't incapacitated.

\DndMonsterSection{Actions}
\DndMonsterAction{Bite}
\textsl{Melee Weapon Attack:} 4 to hit, reach 5 ft., one target. \textsl{Hit:} 4 (1d4 + 2) piercing damage.

\DndMonsterAction{Sling}
\textsl{Ranged Weapon Attack:} 4 to hit, range 30/120 ft., one target. \textsl{Hit:} 4 (1d4 + 2) bludgeoning damage.

\end{DndMonster}



\begin{DndMonster}[float=t]{Riding Horse}
\DndMonsterType{Large beast, unaligned}
\DndMonsterBasics[
armor-class = {10},
hit-points = {\DndDice{2d10 + 2}},
speed = {60 ft.}
]
\DndMonsterAbilityScores[
 str = 16,
 dex = 10,
 con = 12,
 int = 2,
 wis = 11,
 cha = 7
]
\DndMonsterDetails[
challenge = 1/4,
]
\DndMonsterSection{Actions}
\DndMonsterAction{Hooves}
\textsl{Melee Weapon Attack:} 5 to hit, reach 5 ft., one target. \textsl{Hit:} 8 (2d4 + 3) bludgeoning damage.

\end{DndMonster}



\begin{DndMonster}[float=t]{Roper}
\DndMonsterType{Large monstrosity, neutral evil}
\DndMonsterBasics[
armor-class = {20 (natural armor)},
hit-points = {\DndDice{11d10 + 33}},
speed = {10 ft., climb 10 ft.}
]
\DndMonsterAbilityScores[
 str = 18,
 dex = 8,
 con = 17,
 int = 7,
 wis = 16,
 cha = 6
]
\DndMonsterDetails[
skills = {Perception +6, Stealth +5},
senses = {darkvision 60 ft., passive perception 16},
challenge = 5,
]
\DndMonsterAction{False Appearance}
While the roper remains motionless, it is indistinguishable from a normal cave formation, such as a stalagmite.

\DndMonsterAction{Grasping Tendrils}
The roper can have up to six tendrils at a time. Each tendril can be attacked (AC 20; 10 hit points; immunity to poison and psychic damage). Destroying a tendril deals no damage to the roper, which can extrude a replacement tendril on its next turn. A tendril can also be broken if a creature takes an action and succeeds on a DC 15 Strength check against it.

\DndMonsterAction{Spider Climb}
The roper can climb difficult surfaces, including upside down on ceilings, without needing to make an ability check.

\DndMonsterSection{Actions}
\DndMonsterAction{Multiattack}
The roper makes four attacks with its tendrils, uses Reel, and makes one attack with its bite.

\DndMonsterAction{Bite}
\textsl{Melee Weapon Attack:} 7 to hit, reach 5 ft., one target. \textsl{Hit:} 22 (4d8 + 4) piercing damage.

\DndMonsterAction{Tendril}
\textsl{Melee Weapon Attack:} 7 to hit, reach 50 ft., one creature. \textsl{Hit:} The target is grappled (escape DC 15). Until the grapple ends, the target is restrained and has disadvantage on Strength checks and Strength saving throws, and the roper can't use the same tendril on another target.

\DndMonsterAction{Reel}
The roper pulls each creature grappled by it up to 25 feet straight toward it.

\end{DndMonster}



\begin{DndMonster}[float=t]{Rug of Smothering}
\DndMonsterType{Large construct, unaligned}
\DndMonsterBasics[
armor-class = {12},
hit-points = {\DndDice{6d10}},
speed = {10 ft.}
]
\DndMonsterAbilityScores[
 str = 17,
 dex = 14,
 con = 10,
 int = 1,
 wis = 3,
 cha = 1
]
\DndMonsterDetails[
damage-immunities = {poison, psychic},
condition-immunities = {blinded, charmed, deafened, frightened, paralyzed, petrified, poisoned},
senses = {blindsight 60 ft. (blind beyond this radius), passive perception 6},
challenge = 2,
]
\DndMonsterAction{Antimagic Susceptibility}
The rug is incapacitated while in the area of an \textit{antimagic field}. If targeted by \textit{dispel magic}, the rug must succeed on a Constitution saving throw against the caster's spell save DC or fall unconscious for 1 minute.

\DndMonsterAction{Damage Transfer}
While it is grappling a creature, the rug takes only half the damage dealt to it, and the creature grappled by the rug takes the other half.

\DndMonsterAction{False Appearance}
While the rug remains motionless, it is indistinguishable from a normal rug.

\DndMonsterSection{Actions}
\DndMonsterAction{Smother}
\textsl{Melee Weapon Attack:} 5 to hit, reach 5 ft., one Medium or smaller creature. \textsl{Hit:} The creature is grappled (escape DC 13). Until this grapple ends, the target is restrained, blinded, and at risk of suffocating, and the rug can't smother another target. In addition, at the start of each of the target's turns, the target takes 10 (2d6 + 3) bludgeoning damage.

\end{DndMonster}



\begin{DndMonster}[float=t]{Saint Gresha}
\DndMonsterType{Medium undead, neutral good}
\DndMonsterBasics[
armor-class = {13},
hit-points = {\DndDice{9d8 + 27}},
speed = {0 ft., fly 60 ft.}
]
\DndMonsterAbilityScores[
 str = 6,
 dex = 16,
 con = 16,
 int = 12,
 wis = 14,
 cha = 15
]
\DndMonsterDetails[
damage-resistances = {acid; cold; fire; lightning; thunder; bludgeoning, piercing, slashing from nonmagical attacks that aren't silvered},
damage-immunities = {necrotic, poison},
condition-immunities = {charmed, exhaustion, grappled, paralyzed, petrified, poisoned, prone, restrained},
senses = {darkvision 60 ft., passive perception 12},
languages = {The languages it knew in life},
challenge = 6,
]
\DndMonsterAction{Incorporeal Movement}
Saint Gresha can move through other creatures and objects as if they were difficult terrain. It takes 5 (1d10) force damage if it ends its turn inside an object.

\DndMonsterAction{Sunlight Sensitivity}
While in sunlight, Saint Gresha has disadvantage on attack rolls, as well as on Wisdom (Perception) checks that rely on sight.

\DndMonsterAction{Spellcasting}
Saint Gresha is an 8th-level spellcaster. Its spellcasting ability is Charisma (spell save DC 13, 5 to hit with spell attacks). Saint Gresha knows the following sorcerer spells:
\begin{DndMonsterSpells}
  \DndInnateSpellLevel{\textit{guidance}, \textit{light}, \textit{sacred flame}, \textit{thaumaturgy}}
  \DndInnateSpellLevel[3e]{\textit{bless}, \textit{flame strike}, \textit{guiding bolt}, \textit{healing word}, \textit{prayer of healing}, \textit{spirit guardians}}
  \DndInnateSpellLevel[1e]{\textit{divine word}, \textit{harm}, \textit{heal}, \textit{holy aura}, \textit{sacrament of the falling fire}*}
\end{DndMonsterSpells}
\DndMonsterSection{Actions}
\DndMonsterAction{Life Drain}
\textsl{Melee Weapon Attack:} 6 to hit, reach 5 ft., one creature. \textsl{Hit:} 21 (4d8 + 3) necrotic damage. The target must succeed on a DC 14 Constitution saving throw or its hit point maximum is reduced by an amount equal to the damage taken. This reduction lasts until the target finishes a long rest. The target dies if this effect reduces its hit point maximum to 0.

\DndMonsterAction{Create Specter}
Saint Gresha targets a humanoid within 10 feet of it that has been dead for no longer than 1 minute and died violently. The target's spirit rises as a specter in the space of its corpse or in the nearest unoccupied space. The specter is under Saint Gresha's control. Saint Gresha can have no more than seven specters under its control at one time.

\end{DndMonster}



\begin{DndMonster}[float=t]{Scoundrel}
\DndMonsterType{Medium humanoid, any alignment}
\DndMonsterBasics[
armor-class = {14 (leather armor)},
hit-points = {\DndDice{8d8 + 8}},
speed = {30 ft.}
]
\DndMonsterAbilityScores[
 str = 10,
 dex = 17,
 con = 12,
 int = 14,
 wis = 10,
 cha = 15
]
\DndMonsterDetails[
skills = {Acrobatics +5, Deception +6, Sleight Of Hand +5, Stealth +7},
languages = {Thieves' cant plus any two languages},
challenge = 3,
]
\DndMonsterAction{Cunning Action}
On each of its turns, the scoundrel can use a bonus action to take the Dash, Disengage, or Hide action.

\DndMonsterAction{Sneak Attack (1/Turn)}
The scoundrel deals an extra 14 4d6 damage when it hits a target with a weapon attack and has advantage on the attack roll, or when the target is within 5 feet of an ally of the scoundrel that isn't incapacitated and the scoundrel doesn't have disadvantage on the attack roll.

\DndMonsterAction{Spellcasting}
The scoundrel is a 3rd level spellcaster that uses Intelligence as its spellcasting ability (spell save DC 12; 4 to hit with spell attacks). It knows the following spells from the wizard spell list:
\begin{DndMonsterSpells}
  \DndMonsterSpellLevel{\textit{mage hand} (the hand is invisible), \textit{minor illusion}, \textit{prestidigitation}}
   \DndMonsterSpellLevel[1][4]{\textit{charm person}, \textit{disguise self}, \textit{sleep}}
   \DndMonsterSpellLevel[2][2]{\textit{invisibility}, \textit{suggestion}}
\end{DndMonsterSpells}
\DndMonsterSection{Actions}
\DndMonsterAction{Shortsword}
\textsl{Melee Weapon Attack:} 5 to hit, reach 5 ft., one target. \textsl{Hit:} 6 (1d6 + 3) piercing damage.

\DndMonsterAction{Hand Crossbow}
\textsl{Ranged Weapon Attack:} 5 to hit, range 30/120 ft., one target. \textsl{Hit:} 6 (1d6 + 3) piercing damage.

\end{DndMonster}



\begin{DndMonster}[float=t]{Scout}
\DndMonsterType{Medium humanoid (any race), any alignment}
\DndMonsterBasics[
armor-class = {13 (\textit{leather armor})},
hit-points = {\DndDice{3d8 + 3}},
speed = {30 ft.}
]
\DndMonsterAbilityScores[
 str = 11,
 dex = 14,
 con = 12,
 int = 11,
 wis = 13,
 cha = 11
]
\DndMonsterDetails[
skills = {Nature +4, Perception +5, Stealth +6, Survival +5},
languages = {any one language (usually Common)},
challenge = 1/2,
]
\DndMonsterAction{Keen Hearing and Sight}
The scout has advantage on Wisdom (Perception) checks that rely on hearing or sight.

\DndMonsterSection{Actions}
\DndMonsterAction{Multiattack}
The scout makes two melee attacks or two ranged attacks.

\DndMonsterAction{Shortsword}
\textsl{Melee Weapon Attack:} 4 to hit, reach 5 ft., one target. \textsl{Hit:} 5 (1d6 + 2) piercing damage.

\DndMonsterAction{Longbow}
\textsl{Ranged Weapon Attack:} 4 to hit, ranged 150/600 ft., one target. \textsl{Hit:} 6 (1d8 + 2) piercing damage.

\end{DndMonster}


\clearpage

\begin{DndMonster}[float*=b,width=\textwidth + 8pt]{Sea Hag}
\begin{multicols}{2}
\DndMonsterType{Medium fey, chaotic evil}
\DndMonsterBasics[
armor-class = {14 (natural armor)},
hit-points = {\DndDice{7d8 + 21}},
speed = {30 ft., swim 40 ft.}
]
\DndMonsterAbilityScores[
 str = 16,
 dex = 13,
 con = 16,
 int = 12,
 wis = 12,
 cha = 13
]
\DndMonsterDetails[
senses = {darkvision 60 ft., passive perception 11},
languages = {Aquan, Common, Giant},
challenge = 2,
]
\DndMonsterAction{Amphibious}
The hag can breathe air and water.

\DndMonsterAction{Horrific Appearance}
Any humanoid that starts its turn within 30 feet of the hag and can see the hag's true form must make a DC 11 Wisdom saving throw. On a failed save, the creature is frightened for 1 minute. A creature can repeat the saving throw at the end of each of its turns, with disadvantage if the hag is within line of sight, ending the effect on itself on a success. If a creature's saving throw is successful or the effect ends for it, the creature is immune to the hag's Horrific Appearance for the next 24 hours.

Unless the target is Surprise or the revelation of the hag's true form is sudden, the target can avert its eyes and avoid making the initial saving throw. Until the start of its next turn, a creature that averts its eyes has disadvantage on attack rolls against the hag.

\DndMonsterSection{Actions}
\DndMonsterAction{Claws}
\textsl{Melee Weapon Attack:} 5 to hit, reach 5 ft., one target. \textsl{Hit:} 10 (2d6 + 3) slashing damage.

\DndMonsterAction{Death Glare}
The hag targets one frightened creature she can see within 30 feet of her. If the target can see the hag, it must succeed on a DC 11 Wisdom saving throw against this magic or drop to 0 hit points.

\DndMonsterAction{Illusory Appearance}
The hag covers herself and anything she is wearing or carrying with a magical illusion that makes her look like an ugly creature of her general size and humanoid shape. The effect ends if the hag takes a bonus action to end it or if she dies.

The changes wrought by this effect fail to hold up to physical inspection. For example, the hag could appear to have no claws, but someone touching her hand might feel the claws. Otherwise, a creature must take an action to visually inspect the illusion and succeed on a DC 16 Intelligence (Investigation) check to discern that the hag is disguised.

\end{multicols}
\end{DndMonster}



\begin{DndMonster}[float=t]{Shadow}
\DndMonsterType{Medium undead, chaotic evil}
\DndMonsterBasics[
armor-class = {12},
hit-points = {\DndDice{3d8 + 3}},
speed = {40 ft.}
]
\DndMonsterAbilityScores[
 str = 6,
 dex = 14,
 con = 13,
 int = 6,
 wis = 10,
 cha = 8
]
\DndMonsterDetails[
skills = {Stealth +4},
damage-vulnerabilities = {radiant},
damage-resistances = {acid; cold; fire; lightning; thunder; bludgeoning, piercing, slashing from nonmagical attacks},
damage-immunities = {necrotic, poison},
condition-immunities = {exhaustion, frightened, grappled, paralyzed, petrified, poisoned, prone, restrained},
senses = {darkvision 60 ft., passive perception 10},
challenge = 1/2,
]
\DndMonsterAction{Amorphous}
The shadow can move through a space as narrow as 1 inch wide without squeezing.

\DndMonsterAction{Shadow Stealth}
While in dim light or darkness, the shadow can take the Hide action as a bonus action. Its stealth bonus is also improved to +6.

\DndMonsterAction{Sunlight Weakness}
While in sunlight, the shadow has disadvantage on attack rolls, ability checks, and saving throws.

\DndMonsterSection{Actions}
\DndMonsterAction{Strength Drain}
\textsl{Melee Weapon Attack:} 4 to hit, reach 5 ft., one creature. \textsl{Hit:} 9 (2d6 + 2) necrotic damage, and the target's Strength score is reduced by 1d4. The target dies if this reduces its Strength to 0. Otherwise, the reduction lasts until the target finishes a short or long rest.

If a non-evil humanoid dies from this attack, a new shadow rises from the corpse 1d4 hours later.

\end{DndMonster}



\begin{DndMonster}[float=t]{Shambling Mound}
\DndMonsterType{Large plant, unaligned}
\DndMonsterBasics[
armor-class = {15 (natural armor)},
hit-points = {\DndDice{16d10 + 48}},
speed = {20 ft., swim 20 ft.}
]
\DndMonsterAbilityScores[
 str = 18,
 dex = 8,
 con = 16,
 int = 5,
 wis = 10,
 cha = 5
]
\DndMonsterDetails[
skills = {Stealth +2},
damage-resistances = {cold, fire},
damage-immunities = {lightning},
condition-immunities = {blinded, deafened, exhaustion},
senses = {blindsight 60 ft. (blind beyond this radius), passive perception 10},
challenge = 5,
]
\DndMonsterAction{Lightning Absorption}
Whenever the shambling mound is subjected to lightning damage, it takes no damage and regains a number of hit points equal to the lightning damage dealt.

\DndMonsterSection{Actions}
\DndMonsterAction{Multiattack}
The shambling mound makes two slam attacks. If both attacks hit a Medium or smaller target, the target is grappled (escape DC 14), and the shambling mound uses its Engulf on it.

\DndMonsterAction{Slam}
\textsl{Melee Weapon Attack:} 7 to hit, reach 5 ft., one target. \textsl{Hit:} 13 (2d8 + 4) bludgeoning damage.

\DndMonsterAction{Engulf}
The shambling mound engulfs a Medium or smaller creature grappled by it. The engulfed target is blinded, restrained, and unable to breathe, and it must succeed on a DC 14 Constitution saving throw at the start of each of the mound's turns or take 13 (2d8 + 4) bludgeoning damage. If the mound moves, the engulfed target moves with it. The mound can have only one creature engulfed at a time.

\end{DndMonster}



\begin{DndMonster}[float=t]{Shield Guardian}
\DndMonsterType{Large construct, unaligned}
\DndMonsterBasics[
armor-class = {17 (natural armor)},
hit-points = {\DndDice{15d10 + 60}},
speed = {30 ft.}
]
\DndMonsterAbilityScores[
 str = 18,
 dex = 8,
 con = 18,
 int = 7,
 wis = 10,
 cha = 3
]
\DndMonsterDetails[
damage-immunities = {poison},
condition-immunities = {charmed, exhaustion, frightened, paralyzed, poisoned},
senses = {blindsight 10 ft., darkvision 60 ft., passive perception 10},
languages = {understands commands given in any language but can't speak},
challenge = 7,
]
\DndMonsterAction{Bound}
The shield guardian is magically bound to an amulet. As long as the guardian and its amulet are on the same plane of existence, the amulet's wearer can telepathically call the guardian to travel to it, and the guardian knows the distance and direction to the amulet. If the guardian is within 60 feet of the amulet's wearer, half of any damage the wearer takes (rounded up) is transferred to the guardian.

\DndMonsterAction{Regeneration}
The shield guardian regains 10 hit points at the start of its turn if it has at least 1 hit point.

\DndMonsterAction{Spell Storing}
A spellcaster who wears the shield guardian's amulet can cause the guardian to store one spell of 4th level or lower. To do so, the wearer must cast the spell on the guardian. The spell has no effect but is stored within the guardian. When commanded to do so by the wearer or when a situation arises that was predefined by the spellcaster, the guardian casts the stored spell with any parameters set by the original caster, requiring no components. When the spell is cast or a new spell is stored, any previously stored spell is lost.

\DndMonsterSection{Actions}
\DndMonsterAction{Multiattack}
The guardian makes two fist attacks.

\DndMonsterAction{Fist}
\textsl{Melee Weapon Attack:} 7 to hit, reach 5 ft., one target. \textsl{Hit:} 11 (2d6 + 4) bludgeoning damage.

\DndMonsterSection{Reactions}
\DndMonsterAction{Shield}
When a creature makes an attack against the wearer of the guardian's amulet, the guardian grants a +2 bonus to the wearer's AC if the guardian is within 5 feet of the wearer.

\end{DndMonster}



\begin{DndMonster}[float=t]{Skeleton}
\DndMonsterType{Medium undead, lawful evil}
\DndMonsterBasics[
armor-class = {13 (armor scraps)},
hit-points = {\DndDice{2d8 + 4}},
speed = {30 ft.}
]
\DndMonsterAbilityScores[
 str = 10,
 dex = 14,
 con = 15,
 int = 6,
 wis = 8,
 cha = 5
]
\DndMonsterDetails[
damage-vulnerabilities = {bludgeoning},
damage-immunities = {poison},
condition-immunities = {exhaustion, poisoned},
senses = {darkvision 60 ft., passive perception 9},
languages = {understands all languages it spoke in life but can't speak},
challenge = 1/4,
]
\DndMonsterSection{Actions}
\DndMonsterAction{Shortsword}
\textsl{Melee Weapon Attack:} 4 to hit, reach 5 ft., one target. \textsl{Hit:} 5 (1d6 + 2) piercing damage.

\DndMonsterAction{Shortbow}
\textsl{Ranged Weapon Attack:} 4 to hit, range 80/320 ft., one target. \textsl{Hit:} 5 (1d6 + 2) piercing damage.

\end{DndMonster}



\begin{DndMonster}[float=t]{Spy}
\DndMonsterType{Medium humanoid (any race), any alignment}
\DndMonsterBasics[
armor-class = {12},
hit-points = {\DndDice{6d8}},
speed = {30 ft.}
]
\DndMonsterAbilityScores[
 str = 10,
 dex = 15,
 con = 10,
 int = 12,
 wis = 14,
 cha = 16
]
\DndMonsterDetails[
skills = {Deception +5, Insight +4, Investigation +5, Perception +6, Persuasion +5, Sleight Of Hand +4, Stealth +4},
languages = {any two languages},
challenge = 1,
]
\DndMonsterAction{Cunning Action}
On each of its turns, the spy can use a bonus action to take the Dash, Disengage, or Hide action.

\DndMonsterAction{Sneak Attack (1/Turn)}
The spy deals an extra 7 (2d6) damage when it hits a target with a weapon attack and has advantage on the attack roll, or when the target is within 5 feet of an ally of the spy that isn't incapacitated and the spy doesn't have disadvantage on the attack roll.

\DndMonsterSection{Actions}
\DndMonsterAction{Multiattack}
The spy makes two melee attacks.

\DndMonsterAction{Shortsword}
\textsl{Melee Weapon Attack:} 4 to hit, reach 5 ft., one target. \textsl{Hit:} 5 (1d6 + 2) piercing damage.

\DndMonsterAction{Hand Crossbow}
\textsl{Ranged Weapon Attack:} 4 to hit, range 30/120 ft., one target. \textsl{Hit:} 5 (1d6 + 2) piercing damage.

\end{DndMonster}



\begin{DndMonster}[float=t]{Steam Mephit}
\DndMonsterType{Small elemental, neutral evil}
\DndMonsterBasics[
armor-class = {10},
hit-points = {\DndDice{6d6}},
speed = {30 ft., fly 30 ft.}
]
\DndMonsterAbilityScores[
 str = 5,
 dex = 11,
 con = 10,
 int = 11,
 wis = 10,
 cha = 12
]
\DndMonsterDetails[
damage-immunities = {fire, poison},
condition-immunities = {poisoned},
senses = {darkvision 60 ft., passive perception 10},
languages = {Aquan, Ignan},
challenge = 1/4,
]
\DndMonsterAction{Death Burst}
When the mephit dies, it explodes in a cloud of steam. Each creature within 5 feet of the mephit must succeed on a DC 10 Dexterity saving throw or take 4 (1d8) fire damage.

\DndMonsterAction{Innate Spellcasting (1/Day)}
The mephit can innately cast \textit{blur}, requiring no material components. Its innate spellcasting ability is Charisma.
\begin{DndMonsterSpells}
  \DndInnateSpellLevel{\textit{blur}}
\end{DndMonsterSpells}
\DndMonsterSection{Actions}
\DndMonsterAction{Claws}
\textsl{Melee Weapon Attack:} 2 to hit, reach 5 ft., one creature. \textsl{Hit:} 2 (1d4) slashing damage plus 2 (1d4) fire damage.

\DndMonsterAction{Steam Breath (Recharge 6)}
The mephit exhales a 15-foot cone of scalding steam. Each creature in that area must succeed on a DC 10 Dexterity saving throw, taking 4 (1d8) fire damage on a failed save, or half as much damage on a successful one.

\end{DndMonster}



\begin{DndMonster}[float=t]{Stone Golem}
\DndMonsterType{Large construct, unaligned}
\DndMonsterBasics[
armor-class = {17 (natural armor)},
hit-points = {\DndDice{17d10 + 85}},
speed = {30 ft.}
]
\DndMonsterAbilityScores[
 str = 22,
 dex = 9,
 con = 20,
 int = 3,
 wis = 11,
 cha = 1
]
\DndMonsterDetails[
damage-immunities = {poison; psychic; bludgeoning, piercing, slashing from nonmagical attacks that aren't adamantine},
condition-immunities = {charmed, exhaustion, frightened, paralyzed, petrified, poisoned},
senses = {darkvision 120 ft., passive perception 10},
languages = {understands the languages of its creator but can't speak},
challenge = 10,
]
\DndMonsterAction{Immutable Form}
The golem is immune to any spell or effect that would alter its form.

\DndMonsterAction{Magic Resistance}
The golem has advantage on saving throws against spells and other magical effects.

\DndMonsterAction{Magic Weapons}
The golem's weapon attacks are magical.

\DndMonsterSection{Actions}
\DndMonsterAction{Multiattack}
The golem makes two slam attacks.

\DndMonsterAction{Slam}
\textsl{Melee Weapon Attack:} 10 to hit, reach 5 ft., one target. \textsl{Hit:} 19 (3d8 + 6) bludgeoning damage.

\DndMonsterAction{Slow (Recharge 5\textendash 6)}
The golem targets one or more creatures it can see within 10 feet of it. Each target must make a DC 17 Wisdom saving throw against this magic. On a failed save, a target can't use reactions, its speed is halved, and it can't make more than one attack on its turn. In addition, the target can take either an action or a bonus action on its turn, not both. These effects last for 1 minute. A target can repeat the saving throw at the end of each of its turns, ending the effect on itself on a success.

\end{DndMonster}



\begin{DndMonster}[float=t]{Swarm of Rats}
\DndMonsterType{Medium beast, unaligned}
\DndMonsterBasics[
armor-class = {10},
hit-points = {\DndDice{7d8 - 7}},
speed = {30 ft.}
]
\DndMonsterAbilityScores[
 str = 9,
 dex = 11,
 con = 9,
 int = 2,
 wis = 10,
 cha = 3
]
\DndMonsterDetails[
damage-resistances = {bludgeoning, piercing, slashing},
condition-immunities = {charmed, frightened, grappled, paralyzed, petrified, prone, restrained, stunned},
senses = {darkvision 30 ft., passive perception 10},
challenge = 1/4,
]
\DndMonsterAction{Keen Smell}
The swarm has advantage on Wisdom (Perception) checks that rely on smell.

\DndMonsterAction{Swarm}
The swarm can occupy another creature's space and vice versa, and the swarm can move through any opening large enough for a Tiny rat. The swarm can't regain hit points or gain temporary hit points.

\DndMonsterSection{Actions}
\DndMonsterAction{Bites}
\textsl{Melee Weapon Attack:} 2 to hit, reach 0 ft., one target in the swarm's space. \textsl{Hit:} 7 (2d6) piercing damage, or 3 (1d6) piercing damage if the swarm has half of its hit points or fewer.

\end{DndMonster}



\begin{DndMonster}[float*=b,width=\textwidth + 8pt]{Tarrasque}
\begin{multicols}{2}
\DndMonsterType{Gargantuan monstrosity (titan), unaligned}
\DndMonsterBasics[
armor-class = {25 (natural armor)},
hit-points = {\DndDice{33d20 + 330}},
speed = {40 ft.}
]
\DndMonsterAbilityScores[
 str = 30,
 dex = 11,
 con = 30,
 int = 3,
 wis = 11,
 cha = 11
]
\DndMonsterDetails[
saving-throws = {Int +5, Wis +9, Cha +9},
damage-immunities = {fire; poison; bludgeoning, piercing, slashing from nonmagical attacks},
condition-immunities = {charmed, frightened, paralyzed, poisoned},
senses = {blindsight 120 ft., passive perception 10},
challenge = 30,
]
\DndMonsterAction{Legendary Resistance (3/Day)}
If the tarrasque fails a saving throw, it can choose to succeed instead.

\DndMonsterAction{Magic Resistance}
The tarrasque has advantage on saving throws against spells and other magical effects.

\DndMonsterAction{Reflective Carapace}
Any time the tarrasque is targeted by a \textit{magic missile} spell, a line spell, or a spell that requires a ranged attack roll, roll a d6. On a 1 to 5, the tarrasque is unaffected. On a 6, the tarrasque is unaffected, and the effect is reflected back at the caster as though it originated from the tarrasque, turning the caster into the target.

\DndMonsterAction{Siege Monster}
The tarrasque deals double damage to objects and structures.

\DndMonsterSection{Actions}
\DndMonsterAction{Multiattack}
The tarrasque can use its Frightful Presence. It then makes five attacks: one with its bite, two with its claws, one with its horns, and one with its tail. It can use its Swallow instead of its bite.

\DndMonsterAction{Bite}
\textsl{Melee Weapon Attack:} 19 to hit, reach 10 ft., one target. \textsl{Hit:} 36 (4d12 + 10) piercing damage. If the target is a creature, it is grappled (escape DC 20). Until this grapple ends, the target is restrained, and the tarrasque can't bite another target.

\DndMonsterAction{Claw}
\textsl{Melee Weapon Attack:} 19 to hit, reach 15 ft., one target. \textsl{Hit:} 28 (4d8 + 10) slashing damage.

\DndMonsterAction{Horns}
\textsl{Melee Weapon Attack:} 19 to hit, reach 10 ft., one target. \textsl{Hit:} 32 (4d10 + 10) piercing damage.

\DndMonsterAction{Tail}
\textsl{Melee Weapon Attack:} 19 to hit, reach 20 ft., one target. \textsl{Hit:} 24 (4d6 + 10) bludgeoning damage. If the target is a creature, it must succeed on a DC 20 Strength saving throw or be knocked prone.

\DndMonsterAction{Frightful Presence}
Each creature of the tarrasque's choice within 120 feet of it and aware of it must succeed on a DC 17 Wisdom saving throw or become frightened for 1 minute. A creature can repeat the saving throw at the end of each of its turns, with disadvantage if the tarrasque is within line of sight, ending the effect on itself on a success. If a creature's saving throw is successful or the effect ends for it, the creature is immune to the tarrasque's Frightful Presence for the next 24 hours.

\DndMonsterAction{Swallow}
The tarrasque makes one bite attack against a Large or smaller creature it is grappling. If the attack hits, the target takes the bite's damage, the target is swallowed, and the grapple ends. While swallowed, the creature is blinded and restrained, it has total cover against attacks and other effects outside the tarrasque, and it takes 56 (16d6) acid damage at the start of each of the tarrasque's turns.

If the tarrasque takes 60 damage or more on a single turn from a creature inside it, the tarrasque must succeed on a DC 20 Constitution saving throw at the end of that turn or regurgitate all swallowed creatures, which fall prone in a space within 10 feet of the tarrasque. If the tarrasque dies, a swallowed creature is no longer restrained by it and can escape from the corpse by using 30 feet of movement, exiting prone.

\DndMonsterSection{Legendary Actions}
Tarrasque can take 3 legendary actions, choosing from the options below. Only one legendary action can be used at a time and only at the end of another creature's turn. Tarrasque regains spent legendary actions at the start of its turn.

\begin{DndMonsterLegendaryActions}
\DndMonsterLegendaryAction{Attack}{The tarrasque makes one claw attack or tail attack.
}
\DndMonsterLegendaryAction{Move}{The tarrasque moves up to half its speed.
}
\DndMonsterLegendaryAction{Chomp (Costs 2 Actions)}{The tarrasque makes one bite attack or uses its Swallow.
}
\end{DndMonsterLegendaryActions}
\end{multicols}
\end{DndMonster}



\begin{DndMonster}[float=t]{Thug}
\DndMonsterType{Medium humanoid (any race), any non-good alignment}
\DndMonsterBasics[
armor-class = {11 (\textit{leather armor})},
hit-points = {\DndDice{5d8 + 10}},
speed = {30 ft.}
]
\DndMonsterAbilityScores[
 str = 15,
 dex = 11,
 con = 14,
 int = 10,
 wis = 10,
 cha = 11
]
\DndMonsterDetails[
skills = {Intimidation +2},
languages = {any one language (usually Common)},
challenge = 1/2,
]
\DndMonsterAction{Pack Tactics}
The thug has advantage on an attack roll against a creature if at least one of the thug's allies is within 5 feet of the creature and the ally isn't incapacitated.

\DndMonsterSection{Actions}
\DndMonsterAction{Multiattack}
The thug makes two melee attacks.

\DndMonsterAction{Mace}
\textsl{Melee Weapon Attack:} 4 to hit, reach 5 ft., one creature. \textsl{Hit:} 5 (1d6 + 2) bludgeoning damage.

\DndMonsterAction{Heavy Crossbow}
\textsl{Ranged Weapon Attack:} 2 to hit, range 100/400 ft., one target. \textsl{Hit:} 5 (1d10) piercing damage.

\end{DndMonster}



\begin{DndMonster}[float=t]{Tower Dragon}
\DndMonsterType{Large dragon, lawful good}
\DndMonsterBasics[
armor-class = {18 (natural armor)},
hit-points = {\DndDice{15d10 + 60}},
speed = {40 ft., fly 80 ft., swim 40 ft.}
]
\DndMonsterAbilityScores[
 str = 21,
 dex = 10,
 con = 19,
 int = 14,
 wis = 13,
 cha = 17
]
\DndMonsterDetails[
saving-throws = {Dex +3, Con +7, Wis +4, Cha +6},
skills = {Insight +4, Perception +7, Stealth +3},
damage-immunities = {lightning},
senses = {blindsight 30 ft., darkvision 120 ft., passive perception 17},
languages = {Common, Draconic},
challenge = 8,
]
\DndMonsterAction{Copy Warning}
This content has been copied from another creature, but the data replacement hasn't been run. Check details.

\DndMonsterAction{Amphibious}
The dragon can breathe air and water.

\DndMonsterSection{Actions}
\DndMonsterAction{Multiattack}
The dragon makes three attacks: one with its bite and two with its claws.

\DndMonsterAction{Bite}
\textsl{Melee Weapon Attack:} 8 to hit, reach 10 ft., one target. \textsl{Hit:} 16 (2d10 + 5) piercing damage.

\DndMonsterAction{Claw}
\textsl{Melee Weapon Attack:} 8 to hit, reach 5 ft., one target. \textsl{Hit:} 12 (2d6 + 5) slashing damage.

\DndMonsterAction{Breath Weapons (Recharge 5\textendash 6)}
The dragon uses one of the following breath weapons.


\begin{description}
\item[Lightning Breath.]
The dragon exhales lightning in a 60-foot line that is 5 feet wide. Each creature in that line must make a DC 15 Dexterity saving throw, taking 55 (10d10) lightning damage on a failed save, or half as much damage on a successful one.
\item[Repulsion Breath.]
The dragon exhales repulsion energy in a 30-foot cone. Each creature in that area must succeed on a DC 15 Strength saving throw. On a failed save, the creature is pushed 40 feet away from the dragon.
\end{description}

\end{DndMonster}



\begin{DndMonster}[float=t]{Treant}
\DndMonsterType{Huge plant, chaotic good}
\DndMonsterBasics[
armor-class = {16 (natural armor)},
hit-points = {\DndDice{12d12 + 60}},
speed = {30 ft.}
]
\DndMonsterAbilityScores[
 str = 23,
 dex = 8,
 con = 21,
 int = 12,
 wis = 16,
 cha = 12
]
\DndMonsterDetails[
damage-vulnerabilities = {fire},
damage-resistances = {bludgeoning, piercing},
languages = {Common, Druidic, Elvish, Sylvan},
challenge = 9,
]
\DndMonsterAction{False Appearance}
While the treant remains motionless, it is indistinguishable from a normal tree.

\DndMonsterAction{Siege Monster}
The treant deals double damage to objects and structures.

\DndMonsterSection{Actions}
\DndMonsterAction{Multiattack}
The treant makes two slam attacks.

\DndMonsterAction{Slam}
\textsl{Melee Weapon Attack:} 10 to hit, reach 5 ft., one target. \textsl{Hit:} 16 (3d6 + 6) bludgeoning damage.

\DndMonsterAction{Rock}
\textsl{Ranged Weapon Attack:} 10 to hit, range 60/180 ft., one target. \textsl{Hit:} 28 (4d10 + 6) bludgeoning damage.

\DndMonsterAction{Animate Trees (1/Day)}
The treant magically animates one or two trees it can see within 60 feet of it. These trees have the same statistics as a \textbf{treant}, except they have Intelligence and Charisma scores of 1, they can't speak, and they have only the Slam action option. An animated tree acts as an ally of the treant. The tree remains animate for 1 day or until it dies; until the treant dies or is more than 120 feet from the tree; or until the treant takes a bonus action to turn it back into an inanimate tree. The tree then takes root if possible.

\end{DndMonster}



\begin{DndMonster}[float=t]{Troll}
\DndMonsterType{Large giant, chaotic evil}
\DndMonsterBasics[
armor-class = {15 (natural armor)},
hit-points = {\DndDice{8d10 + 40}},
speed = {30 ft.}
]
\DndMonsterAbilityScores[
 str = 18,
 dex = 13,
 con = 20,
 int = 7,
 wis = 9,
 cha = 7
]
\DndMonsterDetails[
skills = {Perception +2},
senses = {darkvision 60 ft., passive perception 12},
languages = {Giant},
challenge = 5,
]
\DndMonsterAction{Keen Smell}
The troll has advantage on Wisdom (Perception) checks that rely on smell.

\DndMonsterAction{Regeneration}
The troll regains 10 hit points at the start of its turn. If the troll takes acid or fire damage, this trait doesn't function at the start of the troll's next turn. The troll dies only if it starts its turn with 0 hit points and doesn't regenerate.

\DndMonsterSection{Actions}
\DndMonsterAction{Multiattack}
The troll makes three attacks: one with its bite and two with its claws.

\DndMonsterAction{Bite}
\textsl{Melee Weapon Attack:} 7 to hit, reach 5 ft., one target. \textsl{Hit:} 7 (1d6 + 4) piercing damage.

\DndMonsterAction{Claw}
\textsl{Melee Weapon Attack:} 7 to hit, reach 5 ft., one target. \textsl{Hit:} 11 (2d6 + 4) slashing damage.

\end{DndMonster}



\begin{DndMonster}[float=t]{Tyrannosaurus Rex}
\DndMonsterType{Huge beast, unaligned}
\DndMonsterBasics[
armor-class = {13 (natural armor)},
hit-points = {\DndDice{13d12 + 52}},
speed = {50 ft.}
]
\DndMonsterAbilityScores[
 str = 25,
 dex = 10,
 con = 19,
 int = 2,
 wis = 12,
 cha = 9
]
\DndMonsterDetails[
skills = {Perception +4},
challenge = 8,
]
\DndMonsterSection{Actions}
\DndMonsterAction{Multiattack}
The tyrannosaurus makes two attacks: one with its bite and one with its tail. It can't make both attacks against the same target.

\DndMonsterAction{Bite}
\textsl{Melee Weapon Attack:} 10 to hit, reach 10 ft., one target. \textsl{Hit:} 33 (4d12 + 7) piercing damage. If the target is a Medium or smaller creature, it is grappled (escape DC 17). Until this grapple ends, the target is restrained, and the tyrannosaurus can't bite another target.

\DndMonsterAction{Tail}
\textsl{Melee Weapon Attack:} 10 to hit, reach 10 ft., one target. \textsl{Hit:} 20 (3d8 + 7) bludgeoning damage.

\end{DndMonster}



\begin{DndMonster}[float=t]{Urban Ranger}
\DndMonsterType{Medium humanoid, any alignment}
\DndMonsterBasics[
armor-class = {15 (studded leather armor)},
hit-points = {\DndDice{8d8 + 16}},
speed = {30 ft., climb 30 ft.}
]
\DndMonsterAbilityScores[
 str = 12,
 dex = 17,
 con = 14,
 int = 11,
 wis = 15,
 cha = 11
]
\DndMonsterDetails[
skills = {Athletics +5, Perception +6, Stealth +7, Survival +6},
languages = {Any two},
challenge = 3,
]
\DndMonsterAction{Spellcasting}
The urban ranger is an 8th-level spellcaster that uses Wisdom as its spellcasting ability (spell save DC 12; 4 to hit with spell attacks). It knows the following spells from the ranger spell list:
\begin{DndMonsterSpells}
   \DndMonsterSpellLevel[1][4]{\textit{cure wounds}, \textit{fog cloud}, \textit{hunter's mark}}
   \DndMonsterSpellLevel[2][3]{\textit{pass without trace}, \textit{spike growth}}
\end{DndMonsterSpells}
\DndMonsterSection{Actions}
\DndMonsterAction{Multiattack}
The urban ranger makes two melee attacks or two ranged attacks.

\DndMonsterAction{Shortsword}
\textsl{Melee Weapon Attack:} 5 to hit, reach 5 ft., one target. \textsl{Hit:} 6 (1d6 + 3) piercing damage.

\DndMonsterAction{Longbow}
\textsl{Ranged Weapon Attack:} 5 to hit, ranged 150/600 ft., one target. \textsl{Hit:} 7 (1d8 + 3) piercing damage.

\end{DndMonster}



\begin{DndMonster}[float=t]{Veteran}
\DndMonsterType{Medium humanoid (any race), any alignment}
\DndMonsterBasics[
armor-class = {17 (\textit{splint armor})},
hit-points = {\DndDice{9d8 + 18}},
speed = {30 ft.}
]
\DndMonsterAbilityScores[
 str = 16,
 dex = 13,
 con = 14,
 int = 10,
 wis = 11,
 cha = 10
]
\DndMonsterDetails[
skills = {Athletics +5, Perception +2},
languages = {any one language (usually Common)},
challenge = 3,
]
\DndMonsterSection{Actions}
\DndMonsterAction{Multiattack}
The veteran makes two longsword attacks. If it has a shortsword drawn, it can also make a shortsword attack.

\DndMonsterAction{Longsword}
\textsl{Melee Weapon Attack:} 5 to hit, reach 5 ft., one target. \textsl{Hit:} 7 (1d8 + 3) slashing damage, or 8 (1d10 + 3) slashing damage if used with two hands.

\DndMonsterAction{Shortsword}
\textsl{Melee Weapon Attack:} 5 to hit, reach 5 ft., one target. \textsl{Hit:} 6 (1d6 + 3) piercing damage.

\DndMonsterAction{Heavy Crossbow}
\textsl{Ranged Weapon Attack:} 3 to hit, range 100/400 ft., one target. \textsl{Hit:} 6 (1d10 + 1) piercing damage.

\end{DndMonster}



\begin{DndMonster}[float=t]{Vrock}
\DndMonsterType{Large fiend (demon), chaotic evil}
\DndMonsterBasics[
armor-class = {15 (natural armor)},
hit-points = {\DndDice{11d10 + 44}},
speed = {40 ft., fly 60 ft.}
]
\DndMonsterAbilityScores[
 str = 17,
 dex = 15,
 con = 18,
 int = 8,
 wis = 13,
 cha = 8
]
\DndMonsterDetails[
saving-throws = {Dex +5, Wis +4, Cha +2},
damage-resistances = {cold; fire; lightning; bludgeoning, piercing, slashing from nonmagical attacks},
damage-immunities = {poison},
condition-immunities = {poisoned},
senses = {darkvision 120 ft., passive perception 11},
languages = {Abyssal, telepathy 120 ft.},
challenge = 6,
]
\DndMonsterAction{Magic Resistance}
The vrock has advantage on saving throws against spells and other magical effects.

\DndMonsterSection{Actions}
\DndMonsterAction{Multiattack}
The vrock makes two attacks: one with its beak and one with its talons.

\DndMonsterAction{Beak}
\textsl{Melee Weapon Attack:} 6 to hit, reach 5 ft., one target. \textsl{Hit:} 10 (2d6 + 3) piercing damage.

\DndMonsterAction{Talons}
\textsl{Melee Weapon Attack:} 6 to hit, reach 5 ft., one target. \textsl{Hit:} 14 (2d10 + 3) slashing damage.

\DndMonsterAction{Spores (Recharge 6)}
A 15-foot-radius cloud of toxic spores extends out from the vrock. The spores spread around corners. Each creature in that area must succeed on a DC 14 Constitution saving throw or become poisoned. While poisoned in this way, a target takes 5 (1d10) poison damage at the start of each of its turns. A target can repeat the saving throw at the end of each of its turns, ending the effect on itself on a success. Emptying a vial of holy water on the target also ends the effect on it.

\DndMonsterAction{Stunning Screech (1/Day)}
The vrock emits a horrific screech. Each creature within 20 feet of it that can hear it and that isn't a demon must succeed on a DC 14 Constitution saving throw or be stunned until the end of the vrock's next turn.

\end{DndMonster}



\begin{DndMonster}[float=t]{Walking Delerium Geode}
\DndMonsterType{Large elemental, unaligned}
\DndMonsterBasics[
armor-class = {17 (natural armor)},
hit-points = {\DndDice{12d10 + 60}},
speed = {30 ft., burrow 30 ft.}
]
\DndMonsterAbilityScores[
 str = 20,
 dex = 8,
 con = 20,
 int = 5,
 wis = 10,
 cha = 5
]
\DndMonsterDetails[
damage-vulnerabilities = {radiant, thunder},
damage-resistances = {bludgeoning, piercing, slashing from nonmagical attacks},
damage-immunities = {necrotic, poison},
condition-immunities = {exhaustion, paralyzed, petrified, poisoned, unconscious},
senses = {darkvision 60 ft., tremorsense 60 ft., passive perception 10},
languages = {Deep Speech},
challenge = 5,
]
\DndMonsterAction{Death Burst}
If damage reduces the walking delerium geode to 0 hit points, the delerium crystals in its body explode in a burst of contaminated energy, unless the damage was force, radiant, or thunder. Each creature within 30 feet of it must succeed on a DC 15 Constitution saving throw or take 42 (12d6) necrotic damage and gain one level of contamination. In addition, this explosion triggers a random arcane anomaly (see Dungeons of Drakkenheim).

\DndMonsterAction{Earth Glide}
The walking delerium geode can burrow through nonmagical, unworked earth and stone. While doing so, the geode doesn't disturb the material it moves through.

\DndMonsterSection{Actions}
\DndMonsterAction{Shard Slam}
\textsl{Melee Weapon Attack:} 8 to hit, reach 10 ft., one target. \textsl{Hit:} 14 (2d8 + 5) bludgeoning damage. A target hit by this attack must succeed on a DC 15 Constitution saving throw or take 10 (3d6) necrotic damage and gain one level of contamination.

\end{DndMonster}



\begin{DndMonster}[float=t]{Wall Gargoyle}
\DndMonsterType{Medium elemental, chaotic evil}
\DndMonsterBasics[
armor-class = {15 (natural armor)},
hit-points = {\DndDice{7d8 + 21}},
speed = {30 ft., fly 60 ft.}
]
\DndMonsterAbilityScores[
 str = 15,
 dex = 11,
 con = 16,
 int = 6,
 wis = 11,
 cha = 7
]
\DndMonsterDetails[
damage-resistances = {bludgeoning, piercing, slashing from nonmagical attacks that aren't adamantine},
damage-immunities = {poison},
condition-immunities = {exhaustion, petrified, poisoned},
senses = {darkvision 60 ft., passive perception 10},
languages = {Terran},
challenge = 2,
]
\DndMonsterAction{Copy Warning}
This content has been copied from another creature, but the data replacement hasn't been run. Check details.

\DndMonsterAction{False Appearance}
While the gargoyle remains motionless, it is indistinguishable from an inanimate statue.

\DndMonsterSection{Actions}
\DndMonsterAction{Multiattack}
The gargoyle makes two attacks: one with its bite and one with its claws.

\DndMonsterAction{Bite}
\textsl{Melee Weapon Attack:} 4 to hit, reach 5 ft., one target. \textsl{Hit:} 5 (1d6 + 2) piercing damage.

\DndMonsterAction{Claws}
\textsl{Melee Weapon Attack:} 4 to hit, reach 5 ft., one target. \textsl{Hit:} 5 (1d6 + 2) slashing damage.

\end{DndMonster}


\clearpage

\begin{DndMonster}[float=t]{Warhorse}
\DndMonsterType{Large beast, unaligned}
\DndMonsterBasics[
armor-class = {11},
hit-points = {\DndDice{3d10 + 3}},
speed = {60 ft.}
]
\DndMonsterAbilityScores[
 str = 18,
 dex = 12,
 con = 13,
 int = 2,
 wis = 12,
 cha = 7
]
\DndMonsterDetails[
challenge = 1/2,
]
\DndMonsterAction{Trampling Charge}
If the horse moves at least 20 feet straight toward a creature and then hits it with a hooves attack on the same turn, that target must succeed on a DC 14 Strength saving throw or be knocked prone. If the target is prone, the horse can make another attack with its hooves against it as a bonus action.

\DndMonsterSection{Actions}
\DndMonsterAction{Hooves}
\textsl{Melee Weapon Attack:} 6 to hit, reach 5 ft., one target. \textsl{Hit:} 11 (2d6 + 4) bludgeoning damage.

\end{DndMonster}



\begin{DndMonster}[float=t]{Warlock of the Rat God}
\DndMonsterType{Small monstrosity, chaotic evil}
\DndMonsterBasics[
armor-class = {12 (15 with \textit{mage armor})},
hit-points = {\DndDice{6d6 + 6}},
speed = {30 ft., swim 30 ft., climb 30 ft.}
]
\DndMonsterAbilityScores[
 str = 7,
 dex = 14,
 con = 13,
 int = 13,
 wis = 11,
 cha = 15
]
\DndMonsterDetails[
skills = {Deception +4, Stealth +4},
senses = {darkvision 60 ft., passive perception 10},
languages = {Common},
challenge = 2,
]
\DndMonsterAction{Skulker}
The ratling can take the Hide action as a bonus action on each of its turns.

\DndMonsterAction{Keen Smell}
The ratling has advantage on Wisdom (Perception) checks that rely on smell.

\DndMonsterAction{Innate Spellcasting}
The warlock's innate spellcasting ability is Charisma. It can innately cast the following spells (spell save DC 12, 4 to hit with spell attacks), requiring no material components:
\begin{DndMonsterSpells}
  \DndInnateSpellLevel{\textit{eldritch blast} (2 beams), \textit{mage armor} (self only), \textit{minor illusion}, \textit{thaumaturgy}}
  \DndInnateSpellLevel[1e]{\textit{augury}, \textit{burning hands}, \textit{conjure animals} (giant rats only), \textit{faerie fire}, \textit{invisibility}, \textit{misty step}}
\end{DndMonsterSpells}
\DndMonsterSection{Actions}
\DndMonsterAction{Bite}
\textsl{Melee Weapon Attack:} 4 to hit, reach 5 ft., one target. \textsl{Hit:} 4 (1d4 + 2) piercing damage.

\end{DndMonster}



\begin{DndMonster}[float=t]{Warp Witch}
\DndMonsterType{Medium undead, any alignment}
\DndMonsterBasics[
armor-class = {11},
hit-points = {\DndDice{10d8}},
speed = {0 ft., fly 40 ft. \textit{(hover)}}
]
\DndMonsterAbilityScores[
 str = 7,
 dex = 13,
 con = 10,
 int = 10,
 wis = 12,
 cha = 17
]
\DndMonsterDetails[
damage-resistances = {acid; fire; lightning; thunder; bludgeoning, piercing, slashing from nonmagical attacks},
damage-immunities = {cold, necrotic, poison},
condition-immunities = {charmed, exhaustion, frightened, grappled, paralyzed, petrified, poisoned, prone, restrained},
senses = {darkvision 60 ft., passive perception 11},
languages = {any languages it knew in life},
challenge = 4,
]
\DndMonsterAction{Copy Warning}
This content has been copied from another creature, but the data replacement hasn't been run. Check details.

\DndMonsterAction{Ethereal Sight}
The ghost can see 60 feet into the Ethereal Plane when it is on the Material Plane, and vice versa.

\DndMonsterAction{Incorporeal Movement}
The ghost can move through other creatures and objects as if they were difficult terrain. It takes 5 (1d10) force damage if it ends its turn inside an object.

\DndMonsterSection{Actions}
\DndMonsterAction{Withering Touch}
\textsl{Melee Weapon Attack:} 5 to hit, reach 5 ft., one target. \textsl{Hit:} 17 (4d6 + 3) necrotic damage.

\DndMonsterAction{Etherealness}
The ghost enters the Ethereal Plane from the Material Plane, or vice versa. It is visible on the Material Plane while it is in the Border Ethereal, and vice versa, yet it can't affect or be affected by anything on the other plane.

\DndMonsterAction{Horrifying Visage}
Each non-undead creature within 60 feet of the ghost that can see it must succeed on a DC 13 Wisdom saving throw or be frightened for 1 minute. If the save fails by 5 or more, the target also ages 1d4 × 10 years. A frightened target can repeat the saving throw at the end of each of its turns, ending the frightened condition on itself on a success. If a target's saving throw is successful or the effect ends for it, the target is immune to this ghost's Horrifying Visage for the next 24 hours. The aging effect can be reversed with a  \textit{greater restoration} spell, but only within 24 hours of it occurring.

\DndMonsterAction{Possession (Recharge 6)}
One humanoid that the ghost can see within 5 feet of it must succeed on a DC 13 Charisma saving throw or be possessed by the ghost; the ghost then disappears, and the target is incapacitated and loses control of its body. The ghost now controls the body but doesn't deprive the target of awareness. The ghost can't be targeted by any attack, spell, or other effect, except ones that turn undead, and it retains its alignment, Intelligence, Wisdom, Charisma, and immunity to being charmed and frightened. It otherwise uses the possessed target's statistics, but doesn't gain access to the target's knowledge, class features, or proficiencies.

The possession lasts until the body drops to 0 hit points, the ghost ends it as a bonus action, or the ghost is turned or forced out by an effect like the \textit{dispel evil and good} spell. When the possession ends, the ghost reappears in an unoccupied space within 5 feet of the body. The target is immune to this ghost's Possession for 24 hours after succeeding on the saving throw or after the possession ends.

\end{DndMonster}



\begin{DndMonster}[float=t]{Water Elemental}
\DndMonsterType{Large elemental, neutral}
\DndMonsterBasics[
armor-class = {14 (natural armor)},
hit-points = {\DndDice{12d10 + 48}},
speed = {30 ft., swim 90 ft.}
]
\DndMonsterAbilityScores[
 str = 18,
 dex = 14,
 con = 18,
 int = 5,
 wis = 10,
 cha = 8
]
\DndMonsterDetails[
damage-resistances = {acid; bludgeoning, piercing, slashing from nonmagical attacks},
damage-immunities = {poison},
condition-immunities = {exhaustion, grappled, paralyzed, petrified, poisoned, prone, restrained, unconscious},
senses = {darkvision 60 ft., passive perception 10},
languages = {Aquan},
challenge = 5,
]
\DndMonsterAction{Water Form}
The elemental can enter a hostile creature's space and stop there. It can move through a space as narrow as 1 inch wide without squeezing.

\DndMonsterAction{Freeze}
If the elemental takes cold damage, it partially freezes; its speed is reduced by 20 feet until the end of its next turn.

\DndMonsterSection{Actions}
\DndMonsterAction{Multiattack}
The elemental makes two slam attacks.

\DndMonsterAction{Slam}
\textsl{Melee Weapon Attack:} 7 to hit, reach 5 ft., one target. \textsl{Hit:} 13 (2d8 + 4) bludgeoning damage.

\DndMonsterAction{Whelm (Recharge 4\textendash 6)}
Each creature in the elemental's space must make a DC 15 Strength saving throw. On a failure, a target takes 13 (2d8 + 4) bludgeoning damage. If it is Large or smaller, it is also grappled (escape DC 14). Until this grapple ends, the target is restrained and unable to breathe unless it can breathe water. If the saving throw is successful, the target is pushed out of the elemental's space.

The elemental can grapple one Large creature or up to two Medium or smaller creatures at one time. At the start of each of the elemental's turns, each target grappled by it takes 13 (2d8 + 4) bludgeoning damage. A creature within 5 feet of the elemental can pull a creature or object out of it by taking an action to make a DC 14 Strength check and succeeding.

\end{DndMonster}



\begin{DndMonster}[float=t]{Wight}
\DndMonsterType{Medium undead, neutral evil}
\DndMonsterBasics[
armor-class = {14 (studded leather armor)},
hit-points = {\DndDice{6d8 + 18}},
speed = {30 ft.}
]
\DndMonsterAbilityScores[
 str = 15,
 dex = 14,
 con = 16,
 int = 10,
 wis = 13,
 cha = 15
]
\DndMonsterDetails[
skills = {Perception +3, Stealth +4},
damage-resistances = {necrotic; bludgeoning, piercing, slashing from nonmagical attacks that aren't silvered},
damage-immunities = {poison},
condition-immunities = {exhaustion, poisoned},
senses = {darkvision 60 ft., passive perception 13},
languages = {the languages it knew in life},
challenge = 3,
]
\DndMonsterAction{Sunlight Sensitivity}
While in sunlight, the wight has disadvantage on attack rolls, as well as on Wisdom (Perception) checks that rely on sight.

\DndMonsterSection{Actions}
\DndMonsterAction{Multiattack}
The wight makes two longsword attacks or two longbow attacks. It can use its Life Drain in place of one longsword attack.

\DndMonsterAction{Life Drain}
\textsl{Melee Weapon Attack:} 4 to hit, reach 5 ft., one creature. \textsl{Hit:} 5 (1d6 + 2) necrotic damage. The target must succeed on a DC 13 Constitution saving throw or its hit point maximum is reduced by an amount equal to the damage taken. This reduction lasts until the target finishes a long rest. The target dies if this effect reduces its hit point maximum to 0.

A humanoid slain by this attack rises 24 hours later as a \textbf{zombie} under the wight's control, unless the humanoid is restored to life or its body is destroyed. The wight can have no more than twelve zombies under its control at one time.

\DndMonsterAction{Longsword}
\textsl{Melee Weapon Attack:} 4 to hit, reach 5 ft., one target. \textsl{Hit:} 6 (1d8 + 2) slashing damage, or 7 (1d10 + 2) slashing damage if used with two hands.

\DndMonsterAction{Longbow}
\textsl{Ranged Weapon Attack:} 4 to hit, range 150/600 ft., one target. \textsl{Hit:} 6 (1d8 + 2) piercing damage.

\end{DndMonster}



\begin{DndMonster}[float=t]{Will-o'-Wisp}
\DndMonsterType{Tiny undead, chaotic evil}
\DndMonsterBasics[
armor-class = {19},
hit-points = {\DndDice{9d4}},
speed = {0 ft., fly 50 ft. \textit{(hover)}}
]
\DndMonsterAbilityScores[
 str = 1,
 dex = 28,
 con = 10,
 int = 13,
 wis = 14,
 cha = 11
]
\DndMonsterDetails[
damage-resistances = {acid; cold; fire; necrotic; thunder; bludgeoning, piercing, slashing from nonmagical attacks},
damage-immunities = {lightning, poison},
condition-immunities = {exhaustion, grappled, paralyzed, poisoned, prone, restrained, unconscious},
senses = {darkvision 120 ft., passive perception 12},
languages = {the languages it knew in life},
challenge = 2,
]
\DndMonsterAction{Consume Life}
As a bonus action, the will-o'-wisp can target one creature it can see within 5 feet of it that has 0 hit points and is still alive. The target must succeed on a DC 10 Constitution saving throw against this magic or die. If the target dies, the will-o'-wisp regains 10 (3d6) hit points.

\DndMonsterAction{Ephemeral}
The will-o'-wisp can't wear or carry anything.

\DndMonsterAction{Incorporeal Movement}
The will-o'-wisp can move through other creatures and objects as if they were difficult terrain. It takes 5 (1d10) force damage if it ends its turn inside an object.

\DndMonsterAction{Variable Illumination}
The will-o'-wisp sheds bright light in a 5 to 20-foot radius and dim light for an additional number of ft. equal to the chosen radius. The will-o'-wisp can alter the radius as a bonus action.

\DndMonsterSection{Actions}
\DndMonsterAction{Shock}
\textsl{Melee Spell Attack:} 4 to hit, reach 5 ft., one creature. \textsl{Hit:} 9 (2d8) lightning damage.

\DndMonsterAction{Invisibility}
The will-o'-wisp and its light magically become invisible until it attacks or uses its Consume Life, or until its concentration ends (as if concentration on a spell).

\end{DndMonster}



\begin{DndMonster}[float=t]{Worg}
\DndMonsterType{Large monstrosity, neutral evil}
\DndMonsterBasics[
armor-class = {13 (natural armor)},
hit-points = {\DndDice{4d10 + 4}},
speed = {50 ft.}
]
\DndMonsterAbilityScores[
 str = 16,
 dex = 13,
 con = 13,
 int = 7,
 wis = 11,
 cha = 8
]
\DndMonsterDetails[
skills = {Perception +4},
senses = {darkvision 60 ft., passive perception 14},
languages = {Goblin, Worg},
challenge = 1/2,
]
\DndMonsterAction{Keen Hearing and Smell}
The worg has advantage on Wisdom (Perception) checks that rely on hearing or smell.

\DndMonsterSection{Actions}
\DndMonsterAction{Bite}
\textsl{Melee Weapon Attack:} 5 to hit, reach 5 ft., one target. \textsl{Hit:} 10 (2d6 + 3) piercing damage. If the target is a creature, it must succeed on a DC 13 Strength saving throw or be knocked prone.

\end{DndMonster}



\begin{DndMonster}[float=t]{Wraith}
\DndMonsterType{Medium undead, neutral evil}
\DndMonsterBasics[
armor-class = {13},
hit-points = {\DndDice{9d8 + 27}},
speed = {0 ft., fly 60 ft. \textit{(hover)}}
]
\DndMonsterAbilityScores[
 str = 6,
 dex = 16,
 con = 16,
 int = 12,
 wis = 14,
 cha = 15
]
\DndMonsterDetails[
damage-resistances = {acid; cold; fire; lightning; thunder; bludgeoning, piercing, slashing from nonmagical attacks that aren't silvered},
damage-immunities = {necrotic, poison},
condition-immunities = {charmed, exhaustion, grappled, paralyzed, petrified, poisoned, prone, restrained},
senses = {darkvision 60 ft., passive perception 12},
languages = {the languages it knew in life},
challenge = 5,
]
\DndMonsterAction{Incorporeal Movement}
The wraith can move through other creatures and objects as if they were difficult terrain. It takes 5 (1d10) force damage if it ends its turn inside an object.

\DndMonsterAction{Sunlight Sensitivity}
While in sunlight, the wraith has disadvantage on attack rolls, as well as on Wisdom (Perception) checks that rely on sight.

\DndMonsterSection{Actions}
\DndMonsterAction{Life Drain}
\textsl{Melee Weapon Attack:} 6 to hit, reach 5 ft., one creature. \textsl{Hit:} 21 (4d8 + 3) necrotic damage. The target must succeed on a DC 14 Constitution saving throw or its hit point maximum is reduced by an amount equal to the damage taken. This reduction lasts until the target finishes a long rest. The target dies if this effect reduces its hit point maximum to 0.

\DndMonsterAction{Create Specter}
The wraith targets a humanoid within 10 feet of it that has been dead for no longer than 1 minute and died violently. The target's spirit rises as a \textbf{specter} in the space of its corpse or in the nearest unoccupied space. The \textbf{specter} is under the wraith's control. The wraith can have no more than seven specters under its control at one time.

\end{DndMonster}



\begin{DndMonster}[float=t]{Zombie}
\DndMonsterType{Medium undead, neutral evil}
\DndMonsterBasics[
armor-class = {8},
hit-points = {\DndDice{3d8 + 9}},
speed = {20 ft.}
]
\DndMonsterAbilityScores[
 str = 13,
 dex = 6,
 con = 16,
 int = 3,
 wis = 6,
 cha = 5
]
\DndMonsterDetails[
saving-throws = {Wis +0},
damage-immunities = {poison},
condition-immunities = {poisoned},
senses = {darkvision 60 ft., passive perception 8},
languages = {understands all languages it spoke in life but can't speak},
challenge = 1/4,
]
\DndMonsterAction{Undead Fortitude}
If damage reduces the zombie to 0 hit points, it must make a Constitution saving throw with a DC of 5 + the damage taken, unless the damage is radiant or from a critical hit. On a success, the zombie drops to 1 hit point instead.

\DndMonsterSection{Actions}
\DndMonsterAction{Slam}
\textsl{Melee Weapon Attack:} 3 to hit, reach 5 ft., one target. \textsl{Hit:} 4 (1d6 + 1) bludgeoning damage.

\end{DndMonster}


\end{document}
